% The chapter input list was generated from the do Carmo docs.
% Preserve the introductory material when refreshing the chapter inputs.
\providecommand{\dcref}[1]{}      % provenance marker (no-op for PDF)
\providecommand{\entryref}[1]{}   % legacy fallback if any survive

\chapter*{About This Blueprint}

This project follows Manfredo P. do Carmo, \emph{Riemannian Geometry},
published by Birkh\"auser in 1992 \cite{doCarmo}.

The book develops the Riemannian geometry used throughout the Poincare
project: metrics and connections, geodesics and the exponential map,
curvature and Jacobi fields, completeness and comparison theorems, and global
applications relating curvature to topology. These results are background for
the formalization of the Poincare conjecture rather than a standalone proof of
the conjecture.

The blueprint is a concise, independently worded distillation, not a verbatim
copy of the book, it only retains statements and proof ideas needed for
formalization. \textbf{We recommend you to read the blueprint alongside the book}. The \emph{dcref} tag records the corresponding theorem,
definition, or other source number in do Carmo's book. 

The blueprint covers the preliminary material on differentiable manifolds and
Chapters 1--13: Riemannian metrics and connections, geodesics, curvature,
Jacobi fields, isometric immersions, completeness, constant-curvature spaces,
energy variation, comparison and index theorems, negative-curvature topology,
and the sphere theorem.

\chapter{Differentiable Manifolds}
\label{chap:dc-differentiable-manifolds}

The model example in \cite{doCarmo} is a regular surface $S\subset\mathbb{R}^3$:
for every $p\in S$, some open $U\subset\mathbb{R}^2$ and neighborhood
$V\subset\mathbb{R}^3$ admit a map
$\mathbf{x}:U\to V\cap S$ that is a differentiable homeomorphism onto $V\cap S$ and
has injective differential at every point. For overlapping parametrizations, the
transition maps are diffeomorphisms (\cref{ex:dc-ch0-4-2}). The abstract definition
below replaces the ambient $\mathbb{R}^3$ by $\mathbb{R}^n$.

Throughout, ``differentiable'' means $C^\infty$.

\section{Differentiable manifolds; tangent space}

\begin{definition}[differentiable manifold]
  \label{def:dc-ch0-2-1}
  \lean{Riemannian.DCDifferentiableManifold}
  \leanok
  \dcref{ch0:2.1}
  \source{docarmo:page-0017}
  A \emph{differentiable manifold} of dimension $n$ is a set $M$ together with a family
  of injective mappings $\mathbf{x}_\alpha:U_\alpha\subset\mathbb{R}^n\to M$ of open
  sets $U_\alpha$ of $\mathbb{R}^n$ into $M$ such that:
  1. $\bigcup_\alpha\mathbf{x}_\alpha(U_\alpha)=M$.
  2. for any pair $\alpha,\beta$ with
     $\mathbf{x}_\alpha(U_\alpha)\cap\mathbf{x}_\beta(U_\beta)=W\neq\phi$, the sets
     $\mathbf{x}_\alpha^{-1}(W)$ and $\mathbf{x}_\beta^{-1}(W)$ are open in
     $\mathbb{R}^n$ and the mappings $\mathbf{x}_\beta^{-1}\circ\mathbf{x}_\alpha$
     are differentiable.
  3. the family $\{(U_\alpha,\mathbf{x}_\alpha)\}$ is maximal relative to (1) and (2).
  The pair $(U_\alpha,\mathbf{x}_\alpha)$ (or $\mathbf{x}_\alpha$) with
  $p\in\mathbf{x}_\alpha(U_\alpha)$ is a \emph{parametrization} (or \emph{system of
  coordinates}) of $M$ at $p$; $\mathbf{x}_\alpha(U_\alpha)$ is then a \emph{coordinate
  neighborhood} at $p$. A family satisfying (1) and (2) is a \emph{differentiable
  structure} on $M$. Any structure satisfying (1)--(2) has a unique maximal extension.
\end{definition}

\begin{remark}[change of parameters as an axiom]
  \label{rem:dc-ch0-2-2}
  \uses{def:dc-ch0-2-1}
  \notready
  \dcref{ch0:2.2}
  \notready
  Condition (2) of \cref{def:dc-ch0-2-1} declares every coordinate change
  differentiable; it is the
  compatibility condition for transferring calculus from $\mathbb{R}^n$ to $M$.
\end{remark}

\begin{remark}[natural topology]
  \label{rem:dc-ch0-2-3}
  \notready
  \dcref{ch0:2.3}
  \notready
  A differentiable structure on $M$ induces a natural topology: define $A\subset M$
  to be \emph{open} if and only if $\mathbf{x}_\alpha^{-1}(A\cap\mathbf{x}_\alpha(U_\alpha))$
  is open in $\mathbb{R}^n$ for all $\alpha$. This defines a topology: $M$ and $\phi$
  are open, arbitrary unions and finite intersections are open, each coordinate
  neighborhood is open, and every $\mathbf{x}_\alpha$ is continuous. The identity
  atlas gives the usual differentiable structure on $\mathbb{R}^n$.
\end{remark}

\begin{example}[the real projective space $P^n(\mathbb{R})$]
  \label{ex:dc-ch0-2-4}
  \notready
  \dcref{ch0:2.4}
  \notready
  $P^n(\mathbb{R})$ is the set of straight lines of $\mathbb{R}^{n+1}$ through the
  origin, i.e. the quotient of $\mathbb{R}^{n+1}-\{0\}$ by
  $(x_1,\dots,x_{n+1})\sim(\lambda x_1,\dots,\lambda x_{n+1})$, $\lambda\neq0$,
  with points denoted $[x_1,\dots,x_{n+1}]$. Let
  $V_i=\{[x_1,\dots,x_{n+1}]; x_i\neq0\}$, $i=1,\dots,n+1$, and define
  $\mathbf{x}_i:\mathbb{R}^n\to V_i$ by
  $$
  \mathbf{x}_i(y_1,\dots,y_n)=[y_1,\dots,y_{i-1},1,y_i,\dots,y_n].
  $$
  Each $\mathbf{x}_i$ is bijective with $\bigcup\mathbf{x}_i(\mathbb{R}^n)=P^n(\mathbb{R})$,
  $\mathbf{x}_i^{-1}(V_i\cap V_j)$ is open, and (for $i>j$, the case $i<j$ similar)
  $$
  \mathbf{x}_j^{-1}\circ\mathbf{x}_i(y_1,\dots,y_n)=\Big(\frac{y_1}{y_j},\dots,\frac{y_{j-1}}{y_j},\frac{y_{j+1}}{y_j},\dots,\frac{y_{i-1}}{y_j},\frac{1}{y_j},\frac{y_i}{y_j},\dots,\frac{y_n}{y_j}\Big)
  $$
  is differentiable. Thus $\{(\mathbb{R}^n,\mathbf{x}_i)\}$ is a differentiable
  structure: $P^n(\mathbb{R})$ is covered by $n+1$ coordinate neighborhoods $V_i$
  (directions not in the hyperplane $x_i=0$), with \emph{inhomogeneous coordinates}
  $\big(\frac{x_1}{x_i},\dots,\frac{x_{i-1}}{x_i},\frac{x_{i+1}}{x_i},\dots,\frac{x_{n+1}}{x_i}\big)$
  corresponding to the \emph{homogeneous coordinates} $(x_1,\dots,x_{n+1})$.
\end{example}

\begin{definition}[differentiable mapping]
  \label{def:dc-ch0-2-5}
  \lean{Riemannian.DCDifferentiableAt}
  \uses{def:dc-ch0-2-1}
  \leanok
  \dcref{ch0:2.5}
  Let $M_1^n$ and $M_2^m$ be differentiable manifolds. A mapping
  $\varphi:M_1\to M_2$ is \emph{differentiable at} $p\in M_1$ if, given a parametrization
  $\mathbf{y}:V\subset\mathbb{R}^m\to M_2$ at $\varphi(p)$, there is a
  parametrization $\mathbf{x}:U\subset\mathbb{R}^n\to M_1$ at $p$ with
  $\varphi(\mathbf{x}(U))\subset\mathbf{y}(V)$ such that
  $$
  \mathbf{y}^{-1}\circ\varphi\circ\mathbf{x}:U\subset\mathbb{R}^n\to\mathbb{R}^m\qquad\text{(1)}
  $$
  is differentiable at $\mathbf{x}^{-1}(p)$. By condition (2) of \cref{def:dc-ch0-2-1}
  this is independent of the choice of parametrizations. The mapping (1) is the
  \emph{expression} of $\varphi$ in the parametrizations $\mathbf{x}$ and $\mathbf{y}$.
\end{definition}

\begin{definition}[tangent vector; tangent space]
  \label{def:dc-ch0-2-6}
  \lean{Riemannian.DCTangentSpaceAt}
  \leanok
  \dcref{ch0:2.6}
  Let $M$ be a differentiable manifold. A differentiable
  $\alpha:(-\varepsilon,\varepsilon)\to M$ is a \emph{curve} in $M$. Suppose
  $\alpha(0)=p$, and let $\mathcal{D}$ be the functions on $M$ differentiable at
  $p$. The \emph{tangent vector to the curve $\alpha$ at $t=0$} is the function
  $\alpha'(0):\mathcal{D}\to\mathbb{R}$ given by
  $$
  \alpha'(0)f=\frac{d(f\circ\alpha)}{dt}\Big|_{t=0},\qquad f\in\mathcal{D}.
  $$
  A \emph{tangent vector at $p$} is the tangent vector at $t=0$ of some curve
  $\alpha:(-\varepsilon,\varepsilon)\to M$ with $\alpha(0)=p$. The set of all such
  is denoted $T_pM$. Choosing a parametrization $\mathbf{x}:U\to M^n$ at
  $p=\mathbf{x}(0)$, with $\mathbf{x}^{-1}\circ\alpha(t)=(x_1(t),\dots,x_n(t))$,
  one obtains
  $$
  \alpha'(0)=\sum_i x_i'(0)\Big(\frac{\partial}{\partial x_i}\Big)_0.\qquad\text{(2)}
  $$
  Here $\big(\frac{\partial}{\partial x_i}\big)_0$ is the tangent vector at $p$ of
  the coordinate curve $x_i\mapsto\mathbf{x}(0,\dots,x_i,\dots,0)$. Expression (2)
  shows $T_pM$ is a vector space of dimension $n$ with \emph{associated basis}
  $\big\{\big(\frac{\partial}{\partial x_1}\big)_0,\dots,\big(\frac{\partial}{\partial x_n}\big)_0\big\}$,
  independent of the parametrization; $T_pM$ is the \emph{tangent space} of $M$ at $p$.
\end{definition}

\begin{proposition}[the differential]
  \label{prop:dc-ch0-2-7}
  \lean{Riemannian.DCDifferentialAt_curve}
  \uses{def:dc-ch0-2-6}
  \leanok
  \dcref{ch0:2.7}
  Let $M_1^n,M_2^m$ be differentiable manifolds and $\varphi:M_1\to M_2$
  differentiable. For $p\in M_1$ and $v\in T_pM_1$, choose a differentiable curve
  $\alpha:(-\varepsilon,\varepsilon)\to M_1$ with $\alpha(0)=p$, $\alpha'(0)=v$, and
  put $\beta=\varphi\circ\alpha$. Then
  $d\varphi_p:T_pM_1\to T_{\varphi(p)}M_2$ given by $d\varphi_p(v)=\beta'(0)$ is a
  linear mapping independent of the chosen curve $\alpha$.
\end{proposition}
\begin{proof}
  \sketch
  \leanok
  In parametrizations $\mathbf{x}$ at $p$ and $\mathbf{y}$ at $\varphi(p)$,
  write $\mathbf{y}^{-1}\circ\varphi\circ\mathbf{x}(q)=(y_1(x_1,\dots,x_n),\dots,y_m(x_1,\dots,x_n))$
  and $\mathbf{x}^{-1}\circ\alpha(t)=(x_1(t),\dots,x_n(t))$. Then with respect to
  the basis $\{(\frac{\partial}{\partial y_i})_0\}$,

  $$
  \beta'(0)=\Big(\sum_{i=1}^n\frac{\partial y_1}{\partial x_i}x_i'(0),\dots,\sum_{i=1}^n\frac{\partial y_m}{\partial x_i}x_i'(0)\Big),\qquad\text{(3)}
  $$
  
  which is independent of $\alpha$ and equals
  $d\varphi_p(v)=\big(\frac{\partial y_i}{\partial x_j}\big)(x_j'(0))$. Thus
  $d\varphi_p$ is linear with matrix $\big(\frac{\partial y_i}{\partial x_j}\big)$
  in the associated bases.
\end{proof}

\begin{definition}[differential]
  \label{def:dc-ch0-2-8}
  \lean{Riemannian.DCDifferentialAt}
  \uses{prop:dc-ch0-2-7}
  \leanok
  \dcref{ch0:2.8}
  The linear mapping $d\varphi_p$ defined by \cref{prop:dc-ch0-2-7} is the \emph{differential}
  of $\varphi$ at $p$.
\end{definition}

\begin{definition}[diffeomorphism; local diffeomorphism]
  \label{def:dc-ch0-2-9}
  \lean{Riemannian.DCDiffeomorph}
  \uses{def:dc-ch0-2-5}
  \leanok
  \dcref{ch0:2.9}
  $\varphi:M_1\to M_2$ is a \emph{diffeomorphism} if it is differentiable, bijective, and
  $\varphi^{-1}$ is differentiable. It is a \emph{local diffeomorphism at} $p$ if there
  exist neighborhoods $U$ of $p$ and $V$ of $\varphi(p)$ with $\varphi:U\to V$ a
  diffeomorphism. By the chain rule, if $\varphi$ is a diffeomorphism then
  $d\varphi_p:T_pM_1\to T_{\varphi(p)}M_2$ is an isomorphism for all $p$; in
  particular $\dim M_1=\dim M_2$.
\end{definition}

\begin{lemma}[differential of a diffeomorphism is an isomorphism]
  \label{lem:dc-ch0-2-9-iso}
  \lean{Riemannian.DCDiffeomorph.mfderiv_bijective}
  \uses{def:dc-ch0-2-9, def:dc-ch0-2-8}
  \leanok
  \dcref{ch0:2.9}
  If $\varphi:M_1\to M_2$ is a diffeomorphism then for every $p\in M_1$ the
  differential $d\varphi_p:T_pM_1\to T_{\varphi(p)}M_2$ is a linear isomorphism.
\end{lemma}
\begin{proof}
  \sketch
  \leanok
  Let $\psi=\varphi^{-1}$. By the chain rule (\cref{prop:dc-ch0-2-7}) applied to
  $\psi\circ\varphi=\mathrm{id}_{M_1}$ and $\varphi\circ\psi=\mathrm{id}_{M_2}$,
  $d\psi_{\varphi(p)}\circ d\varphi_p=\mathrm{id}_{T_pM_1}$ and
  $d\varphi_p\circ d\psi_{\varphi(p)}=\mathrm{id}_{T_{\varphi(p)}M_2}$. Hence
  $d\varphi_p$ is a two-sided invertible linear map, i.e. a linear isomorphism, with
  inverse $d\psi_{\varphi(p)}$.
\end{proof}

\begin{theorem}[local converse, via inverse function theorem]
  \label{thm:dc-ch0-2-10}
  \lean{Riemannian.isLocalDiffeomorphAt_of_mfderiv_equiv}
  \uses{def:dc-ch0-2-5, def:dc-ch0-2-8, def:dc-ch0-2-9}
  \leanok
  \dcref{ch0:2.10}
  Let $\varphi:M_1^n\to M_2^n$ be a differentiable mapping and $p\in M_1$ such that
  $d\varphi_p:T_pM_1\to T_{\varphi(p)}M_2$ is an isomorphism. Then $\varphi$ is a
  local diffeomorphism at $p$.
\end{theorem}
\begin{proof}
  \sketch
  \leanok
  Choose coordinate charts $\mathbf{x}_1$ at $p$ and $\mathbf{x}_2$ at $\varphi(p)$
  and let $\tilde\varphi=\mathbf{x}_2\circ\varphi\circ\mathbf{x}_1^{-1}$ be the
  coordinate representative, a map between open subsets of $\mathbb{R}^n$. Because
  $\varphi$ is differentiable, $\tilde\varphi$ is $C^\infty$ on a neighbourhood of
  $q=\mathbf{x}_1(p)$, and its Fréchet derivative $d\tilde\varphi_q$ is the coordinate
  expression of $d\varphi_p$, hence a linear isomorphism. The inverse function theorem
  in $\mathbb{R}^n$ then makes $\tilde\varphi$ a local diffeomorphism at $q$: the
  analytic core is that a $C^\infty$ map whose derivative at a point is invertible has
  a $C^\infty$ local inverse, the derivative being invertible on a whole neighbourhood
  and the inverse's smoothness propagating pointwise. Finally $\varphi=\mathbf{x}_2^{-1}
  \circ\tilde\varphi\circ\mathbf{x}_1$ near $p$ is a composition of local
  diffeomorphisms (the two charts and $\tilde\varphi$), hence a local diffeomorphism
  at $p$.
\end{proof}

\section{Immersions, embeddings, and submanifolds}

\begin{definition}[immersion; embedding; submanifold]
  \label{def:dc-ch0-3-1}
  \lean{Riemannian.DCIsImmersion, Riemannian.DCIsEmbedding}
  \uses{def:dc-ch0-2-5, def:dc-ch0-2-8}
  \leanok
  \dcref{ch0:3.1}
  % submanifold = image of an embedded inclusion (an instance of DCIsEmbedding);
  % codimension $n-m$ is a derived numeric notion, not separate Lean data.
  Let $M^m,N^n$ be differentiable manifolds. A differentiable $\varphi:M\to N$ is an
  \emph{immersion} if $d\varphi_p:T_pM\to T_{\varphi(p)}N$ is injective for all $p\in M$.
  If in addition $\varphi$ is a homeomorphism onto $\varphi(M)\subset N$ (subspace
  topology), $\varphi$ is an \emph{embedding}. If $M\subset N$ and the inclusion
  $i:M\subset N$ is an embedding, $M$ is a \emph{submanifold} of $N$. If $\varphi:M^m\to
  N^n$ is an immersion then $m\le n$; the difference $n-m$ is the \emph{codimension}.
\end{definition}

\begin{example}[a regular surface in $\mathbb{R}^3$ is a submanifold]
  \label{ex:dc-ch0-3-6}
  \notready
  \dcref{ch0:3.6}
  \notready
  A regular surface $S\subset\mathbb{R}^3$ has the differentiable structure given by
  its parametrizations $\mathbf{x}_\alpha:U_\alpha\to S$; these are embeddings of
  $U_\alpha$ into $S$, by conditions (a),(b) of the definition of regular surface.
  The inclusion $i:S\subset\mathbb{R}^3$ is differentiable (for $p\in S$ take
  $\mathbf{x}$ at $p$ and $j:V\subset\mathbb{R}^3\to V$ the identity, so
  $j^{-1}\circ i\circ\mathbf{x}=\mathbf{x}$ is differentiable); from (b) $i$ is an
  immersion and from (a) $i$ is a homeomorphism onto its image. Hence $S$ is a
  submanifold of $\mathbb{R}^3$.
\end{example}

\begin{proposition}[every immersion is locally an embedding]
  \label{prop:dc-ch0-3-7}
  \lean{Riemannian.DCIsImmersion.exists_isEmbedding_restrict}
  \uses{def:dc-ch0-3-1, lem:dc-ch0-3-7-embedding, lem:dc-ch0-3-7-normalform}
  \leanok
  \dcref{ch0:3.7}
  Let $\varphi:M_1^n\to M_2^m$, $n\le m$, be an immersion. For every $p\in M_1$
  there exists a neighborhood $V\subset M_1$ of $p$ such that the restriction
  $\varphi\,|\,V\to M_2$ is an embedding.
\end{proposition}
\begin{proof}
  \sketch
  \leanok
  A consequence of the inverse function theorem. Take coordinates
  $\mathbf{x}_1:U_1\to M_1$ at $p$, $\mathbf{x}_2:U_2\to M_2$ at $\varphi(p)$, and
  write $\tilde\varphi=\mathbf{x}_2^{-1}\circ\varphi\circ\mathbf{x}_1=(y_1(x),\dots,y_m(x))$.
  With $q=\mathbf{x}_1^{-1}(p)$, renumber so that
  $\frac{\partial(y_1,\dots,y_n)}{\partial(x_1,\dots,x_n)}(q)\neq0$. Extend to
  $\phi:U_1\times\mathbb{R}^{m-n=k}\to\mathbb{R}^m$ by
  
  $$
  \phi(x_1,\dots,x_n,t_1,\dots,t_k)=(y_1(x),\dots,y_n(x),y_{n+1}(x)+t_1,\dots,y_{n+k}(x)+t_k),
  $$
  
  which restricts to $\tilde\varphi$ on $U_1\times\{0\}$ and has
  $\det(d\phi_q)=\frac{\partial(y_1,\dots,y_n)}{\partial(x_1,\dots,x_n)}(q)\neq0$. By
  the inverse function theorem $\phi\,|\,W_1$ is a diffeomorphism onto $W_2$; with
  $\tilde V=W_1\cap U_1$, $V=\mathbf{x}_1(\tilde V)$, the restriction
  $\mathbf{x}_2\circ\tilde\varphi\circ\mathbf{x}_1^{-1}:V\to\varphi(V)$ is a
  diffeomorphism, hence an embedding. This is the normal-form argument of
  \cref{lem:dc-ch0-3-7-normalform}.

\end{proof}

\begin{lemma}[immersion normal form yields a local embedding]
  \label{lem:dc-ch0-3-7-embedding}
  \lean{Riemannian.isEmbedding_restrict_of_isImmersionAt}
  \uses{def:dc-ch0-3-1}
  \leanok
  \dcref{ch0:3.7}
  Let $\varphi:M_1\to M_2$ and $p\in M_1$. Suppose that in suitable coordinate charts
  around $p$ and $\varphi(p)$ the map $\varphi$ has the standard immersion normal form
  $u\mapsto(u,0)$ (a linear inclusion of the model space of $M_1$ into that of $M_2$).
  Then there is an open neighborhood $U\ni p$ such that the restriction $\varphi\,|\,U$ is
  a topological embedding. Hence every immersion is \emph{locally} an embedding.
\end{lemma}
\begin{proof}
  \sketch
  \leanok
  In the chart around $p$, $\varphi$ is the composite of three maps: the chart, which is a
  homeomorphism of a neighborhood of $p$ onto an open subset of the model space; the linear
  inclusion $u\mapsto(u,0)$, which is a topological embedding; and the inverse of the chart
  around $\varphi(p)$, a homeomorphism onto its image. The subspace topology on $U$ is
  therefore induced by $\varphi$ (each factor induces its topology, and inducing maps
  compose), and $\varphi\,|\,U$ is injective because the chart, the linear inclusion, and
  the codomain chart are each injective. An injective inducing map is an embedding.
\end{proof}

\begin{lemma}[local immersion normal form in normed spaces]
  \label{lem:dc-ch0-3-7-normalform}
  \lean{Riemannian.isEmbedding_restrict_of_hasFDerivAt_injective}
  \leanok
  \dcref{ch0:3.7}
  Let $g:E\to F$ be a $C^\infty$ map between finite-dimensional normed spaces, defined on an
  open set $s\ni q$, whose Fréchet derivative $L=dg_q$ is injective. Then $g$ restricts to a
  topological embedding on a neighborhood of $q$.
\end{lemma}
\begin{proof}
  \sketch
  \leanok
  Choose a complement $G$ of the range of $L$ in $F$ (finite-dimensional, hence closed and
  complemented), so that $(x,t)\mapsto Lx+t$ is a linear isomorphism $E\times G\to F$. Extend
  $g$ to $\mathrm{Ext}(x,t)=g(x)+t$ on $s\times G$; its derivative at $(q,0)$ is exactly that
  isomorphism, so by the inverse function theorem $\mathrm{Ext}$ is a homeomorphism from an
  open neighborhood of $(q,0)$ onto an open subset of $F$. The slice inclusion
  $u\mapsto(u,0)$ is a topological embedding, and $g$ agrees with $\mathrm{Ext}\circ(u\mapsto(u,0))$
  near $q$; being a composite of an embedding with a homeomorphism onto its image, $g$ is a
  topological embedding on a neighborhood of $q$.
\end{proof}

\section{Other examples of manifolds. Orientation}

\begin{example}[the tangent bundle]
  \label{ex:dc-ch0-4-1}
  \lean{Riemannian.DCTangentBundle}
  \uses{def:dc-ch0-2-1}
  \leanok
  \dcref{ch0:4.1}
  Let $M^n$ be a differentiable manifold and $TM=\{(p,v); p\in M, v\in T_pM\}$. With
  a maximal structure $\{(U_\alpha,\mathbf{x}_\alpha)\}$ on $M$, coordinates
  $(x_1^\alpha,\dots,x_n^\alpha)$ on $U_\alpha$ and bases
  $\{\frac{\partial}{\partial x_i^\alpha}\}$, define
  $\mathbf{y}_\alpha:U_\alpha\times\mathbb{R}^n\to TM$ by
  $$
  \mathbf{y}_\alpha(x_1^\alpha,\dots,x_n^\alpha,u_1,\dots,u_n)=\Big(\mathbf{x}_\alpha(x_1^\alpha,\dots,x_n^\alpha),\sum_{i=1}^n u_i\frac{\partial}{\partial x_i^\alpha}\Big).
  $$
  Then $\bigcup_\alpha\mathbf{y}_\alpha(U_\alpha\times\mathbb{R}^n)=TM$ verifies (1)
  of \cref{def:dc-ch0-2-1}, and since $\mathbf{x}_\beta^{-1}\circ\mathbf{x}_\alpha$ is
  differentiable so is $d(\mathbf{x}_\beta^{-1}\circ\mathbf{x}_\alpha)$, giving
  $$
  \mathbf{y}_\beta^{-1}\circ\mathbf{y}_\alpha(q_\alpha,v_\alpha)=\big((\mathbf{x}_\beta^{-1}\circ\mathbf{x}_\alpha)(q_\alpha),\,d(\mathbf{x}_\beta^{-1}\circ\mathbf{x}_\alpha)(v_\alpha)\big)
  $$
  differentiable, verifying (2). Thus $\{(U_\alpha\times\mathbb{R}^n,\mathbf{y}_\alpha)\}$
  is a differentiable structure of dimension $2n$ on $TM$, the \emph{tangent bundle}.
\end{example}

\begin{lemma}[the tangent bundle is a smooth manifold]
  \label{lem:dc-ch0-4-1-manifold}
  \lean{Riemannian.DCTangentBundle.isManifold}
  \uses{ex:dc-ch0-4-1}
  \leanok
  \dcref{ch0:4.1}
  The tangent bundle $TM$ of a smooth $n$-manifold $M$ carries a smooth manifold
  structure of dimension $2n$, modelled on $H\times E$.
\end{lemma}
\begin{proof}
  \sketch
  \leanok
  The charts $\mathbf{y}_\alpha$ of \cref{ex:dc-ch0-4-1} form a differentiable
  atlas: the transition maps
  $\mathbf{y}_\beta^{-1}\circ\mathbf{y}_\alpha(q,v)=\big((\mathbf{x}_\beta^{-1}\circ\mathbf{x}_\alpha)(q),\,d(\mathbf{x}_\beta^{-1}\circ\mathbf{x}_\alpha)(v)\big)$
  are smooth, being built from the smooth change of coordinates on $M$ and its
  differential. Hence $TM$ is a smooth manifold modelled on $H\times E$.
\end{proof}

\begin{example}[regular surfaces in $\mathbb{R}^n$]
  \label{ex:dc-ch0-4-2}
  \uses{ex:dc-ch0-3-6}
  \notready
  \dcref{ch0:4.2}
  \notready
  $M^k\subset\mathbb{R}^n$, $k\le n$, is a \emph{regular surface of dimension $k$} if for
  every $p\in M^k$ there is a neighborhood $V$ of $p$ in $\mathbb{R}^n$ and
  $\mathbf{x}:U\subset\mathbb{R}^k\to M\cap V$ with (a) $\mathbf{x}$ a
  differentiable homeomorphism and (b) $(d\mathbf{x})_q:\mathbb{R}^k\to\mathbb{R}^n$
  injective for all $q\in U$. If
  $\mathbf{x}:U\to M^k$, $\mathbf{y}:V\to M^k$ are two parametrizations with
  $\mathbf{x}(U)\cap\mathbf{y}(V)=W\neq\phi$ then
  $h=\mathbf{x}^{-1}\circ\mathbf{y}:\mathbf{y}^{-1}(W)\to\mathbf{x}^{-1}(W)$ is a
  diffeomorphism. Consequently, $M^k$ is a differentiable manifold of dimension
  $k$, and the inclusion $i:M^k\subset\mathbb{R}^n$ is an embedding.
\end{example}
\begin{proof}
  \sketch
  $h$ is a homeomorphism (composition of homeomorphisms). For
  $r\in\mathbf{y}^{-1}(W)$, $q=h(r)$, write
  $\mathbf{x}(u_1,\dots,u_k)=(v_1,\dots,v_n)$; by (b) suppose
  $\frac{\partial(v_1,\dots,v_k)}{\partial(u_1,\dots,u_k)}(q)\neq0$. Extend
  $\mathbf{x}$ to $F:U\times\mathbb{R}^{n-k}\to\mathbb{R}^n$,
  $F(u,t)=(v_1(u),\dots,v_k(u),v_{k+1}(u)+t_{k+1},\dots,v_n(u)+t_n)$, with
  $\det(dF_q)=\frac{\partial(v_1,\dots,v_k)}{\partial(u_1,\dots,u_k)}(q)\neq0$. By
  the inverse function theorem $F^{-1}$ exists and is differentiable on a
  neighborhood $Q$ of $\mathbf{x}(q)$; choosing $R$ with $\mathbf{y}(R)\subset Q$,
  $h\,|\,R=F^{-1}\circ\mathbf{y}\,|\,R$ is differentiable, hence $h$ is, and
  similarly $h^{-1}$. The final assertion follows from
  \cref{ex:dc-ch0-3-6}.
\end{proof}

\begin{example}[inverse image of a regular value]
  \label{ex:dc-ch0-4-3}
  \uses{ex:dc-ch0-4-2}
  \notready
  \dcref{ch0:4.3}
  \notready
  Let $F:U\subset\mathbb{R}^n\to\mathbb{R}^m$ be differentiable on an open set $U$.
  A point $p\in U$ is a \emph{critical point} if $dF_p:\mathbb{R}^n\to\mathbb{R}^m$ is
  not surjective; $F(p)$ is then a \emph{critical value}. A point $a\in\mathbb{R}^m$ not
  a critical value is a \emph{regular value} (any $a\notin F(U)$ is trivially regular; if
  a regular value of $F$ exists in $\mathbb{R}^m$ then $n\ge m$). If $a\in F(U)$ is
  a regular value, then $F^{-1}(a)\subset\mathbb{R}^n$ is a regular surface of
  dimension $n-m=k$, hence (\cref{ex:dc-ch0-4-2}) a submanifold of $\mathbb{R}^n$.
\end{example}
\begin{proof}
  \sketch
  For $p\in F^{-1}(a)$, writing $q=(y_1,\dots,y_m,x_1,\dots,x_k)$ and
  $F(q)=(f_1(q),\dots,f_m(q))$, since $dF_p$ is surjective suppose
  $\frac{\partial(f_1,\dots,f_m)}{\partial(y_1,\dots,y_m)}(p)\neq0$. Define
  $\varphi:U\to\mathbb{R}^{m+k}$ by
  $\varphi(y,x)=(f_1(q),\dots,f_m(q),x_1,\dots,x_k)$, with
  $\det(d\varphi)_p=\frac{\partial(f_1,\dots,f_m)}{\partial(y_1,\dots,y_m)}(p)\neq0$.
  By the inverse function theorem $\varphi$ is a diffeomorphism of a neighborhood
  $Q$ of $p$ onto $W$; with a cube $K^{m+k}$ centered at $\varphi(p)$ and
  $V=\varphi^{-1}(K^{m+k})\cap Q$, $\varphi$ maps $V$ diffeomorphically onto
  $K^m\times K^k$, and $\mathbf{x}:K^k\to V$,
  $\mathbf{x}(x_1,\dots,x_k)=\varphi^{-1}(a_1,\dots,a_m,x_1,\dots,x_k)$, satisfies
  (a),(b) of \cref{ex:dc-ch0-4-2}. Hence $F^{-1}(a)$ is a regular surface.
\end{proof}

\begin{definition}[orientable manifold; orientation]
  \label{def:dc-ch0-4-4}
  \lean{Riemannian.DCIsOrientable}
  \leanok
  \dcref{ch0:4.4}
  $M$ is \emph{orientable} if it admits a differentiable structure
  $\{(U_\alpha,\mathbf{x}_\alpha)\}$ such that:
  (i) for every pair $\alpha,\beta$ with
  $\mathbf{x}_\alpha(U_\alpha)\cap\mathbf{x}_\beta(U_\beta)=W\neq\phi$, the
  differential of the change of coordinates
  $\mathbf{x}_\beta^{-1}\circ\mathbf{x}_\alpha$ has positive determinant.
  Otherwise $M$ is \emph{non-orientable}. A choice of structure satisfying (i) is an
  \emph{orientation}; $M$ is then \emph{oriented}. Two structures satisfying (i) \emph{determine
  the same orientation} if their union again satisfies (i). If $M$ is orientable and
  connected there exist exactly two orientations. For a diffeomorphism
  $\varphi:M_1\to M_2$, $M_1$ is orientable iff $M_2$ is; if both are connected and
  oriented, $\varphi$ \emph{preserves} or \emph{reverses} the orientation according as the
  induced orientation coincides with the given one or not.
\end{definition}

\begin{example}[two-neighborhood criterion]
  \label{ex:dc-ch0-4-5}
  \notready
  \dcref{ch0:4.5}
  \notready
  If $M$ can be covered by two coordinate neighborhoods $V_1,V_2$ with
  $V_1\cap V_2$ connected, then $M$ is orientable: the determinant of the change of
  coordinates has constant sign on the connected $V_1\cap V_2$, and if negative one
  changes the sign of one coordinate to make it positive throughout.
\end{example}

\begin{example}[another description of projective space]
  \label{ex:dc-ch0-4-7}
  \uses{ex:dc-ch0-2-4}
  \notready
  \dcref{ch0:4.7}
  \notready
  $P^n(\mathbb{R})$ is the quotient of the unit sphere $S^n=\{p\in\mathbb{R}^{n+1};
  |p|=1\}$ by $p\sim A(p)=-p$. Introducing on $S^n\subset\mathbb{R}^{n+1}$ the
  regular-surface structure with hemisphere parametrizations
  $\mathbf{x}_i^{\pm}:U_i\to S^n$,
  $U_i=\{x_i=0,\,x_1^2+\dots+x_{i-1}^2+x_{i+1}^2+\dots+x_{n+1}^2<1\}$,
  $\mathbf{x}_i^{\pm}(x_1,\dots,\hat x_i,\dots,x_{n+1})=(x_1,\dots,\pm D_i,\dots,x_{n+1})$,
  $D_i=\sqrt{1-(x_1^2+\dots+x_{i-1}^2+x_{i+1}^2+\dots+x_{n+1}^2)}$. With the
  canonical projection $\pi:S^n\to P^n(\mathbb{R})$, $\pi(p)=\{p,-p\}$, define
  $\mathbf{y}_i=\pi\circ\mathbf{x}_i^+:U_i\to P^n(\mathbb{R})$. Then
  $\mathbf{y}_i^{-1}\circ\mathbf{y}_j=(\mathbf{x}_i^+)^{-1}\circ\mathbf{x}_j^+$ is
  differentiable, so $\{(U_i,\mathbf{y}_i)\}$ is a differentiable structure for
  $P^n(\mathbb{R})$, giving the same maximal structure as \cref{ex:dc-ch0-2-4}.
  The projective space is orientable exactly when $n$ is odd.
\end{example}

\begin{example}[discontinuous action of a group]
  \label{ex:dc-ch0-4-8}
  \uses{ex:dc-ch0-4-7}
  \notready
  \dcref{ch0:4.8}
  \notready
  A group $G$ \emph{acts on} a differentiable manifold $M$ if there is
  $\varphi:G\times M\to M$ such that (i) for each $g\in G$, $\varphi_g(p)=\varphi(g,p)$
  is a diffeomorphism with $\varphi_e=\mathrm{id}$, and (ii)
  $\varphi_{g_1 g_2}=\varphi_{g_1}\circ\varphi_{g_2}$ (writing $gp$ for $\varphi(g,p)$,
  this is $(g_1 g_2)p=g_1(g_2 p)$). The action is \emph{properly discontinuous} if every
  $p\in M$ has a neighborhood $U$ with $U\cap g(U)=\phi$ for all $g\neq e$. The
  action determines $\sim$ on $M$ ($p_1\sim p_2$ iff $p_2=gp_1$), with quotient
  $M/G$ and \emph{projection} $\pi:M\to M/G$, $\pi(p)=Gp$.
  For a properly discontinuous action, $M/G$ has a differentiable structure with
  respect to which $\pi:M\to M/G$ is a local diffeomorphism.
\end{example}
\begin{proof}
  \sketch
  For $p\in M$ choose $\mathbf{x}:V\to M$ at $p$ with $\mathbf{x}(V)\subset
  U$, $U\cap g(U)=\phi$ for $g\neq e$; then $\pi\,|\,U$ is injective, so
  $\mathbf{y}=\pi\circ\mathbf{x}:V\to M/G$ is injective and $\{(V,\mathbf{y})\}$
  covers $M/G$. Given $\mathbf{y}_1=\pi\circ\mathbf{x}_1$,
  $\mathbf{y}_2=\pi\circ\mathbf{x}_2$ with overlapping images, with $\pi_i=\pi\,|\,\mathbf{x}_i(V_i)$,
  on a suitable $W$ the change is
  $\mathbf{y}_1^{-1}\circ\mathbf{y}_2\,|\,W=\mathbf{x}_1^{-1}\circ\pi_1^{-1}\circ\pi_2\circ\mathbf{x}_2$;
  it suffices that $\pi_1^{-1}\circ\pi_2$ be differentiable at
  $p_2=\pi_2^{-1}(q)$. With $p_1=\pi_1^{-1}\circ\pi_2(p_2)$ there is $g\in G$ with
  $gp_2=p_1$, and $\pi_1^{-1}\circ\pi_2\,|\,\mathbf{x}_2(W)$ coincides with the
  diffeomorphism $\varphi_g\,|\,\mathbf{x}_2(W)$. Hence $\pi$ is a local
  diffeomorphism. The quotient in \cref{ex:dc-ch0-4-7} is obtained from
  $M=S^n$ and $G=\{A,I=A^2\}$.
\end{proof}

\section{Vector fields; brackets. Topology of manifolds}

\begin{definition}[vector field]
  \label{def:dc-ch0-5-1}
  \lean{Riemannian.VectorFieldSection, Riemannian.TangentSmoothAt, Riemannian.TangentSmoothAt.iff_coord, Riemannian.SmoothVectorField.smoothAt}
  \uses{ex:dc-ch0-4-1}
  \leanok
  \dcref{ch0:5.1}
  A \emph{vector field $X$ on $M$} associates to each $p\in M$ a vector $X(p)\in T_pM$;
  as a mapping, $X:M\to TM$. $X$ is \emph{differentiable} if $X:M\to TM$
  is. In a parametrization $\mathbf{x}:U\subset\mathbb{R}^n\to M$,
  $$
  X(p)=\sum_{i=1}^n a_i(p)\frac{\partial}{\partial x_i},\qquad\text{(4)}
  $$
  $a_i:U\to\mathbb{R}$; $X$ is differentiable iff the $a_i$ are differentiable for
  some (hence any) parametrization. Let $\mathcal{D}$ be the differentiable
  functions on $M$ and $\mathcal{F}$ the functions on $M$; thinking of
  $X:\mathcal{D}\to\mathcal{F}$,
  $$
  (Xf)(p)=\sum_i a_i(p)\frac{\partial f}{\partial x_i}(p),\qquad\text{(5)}
  $$
  $Xf$ is independent of $\mathbf{x}$, and $X$ is differentiable iff
  $Xf\in\mathcal{D}$ for all $f\in\mathcal{D}$. If $\varphi:M\to M$ is a
  diffeomorphism, $v\in T_pM$, and $f$ differentiable near $\varphi(p)$, then
  $(d\varphi(v)f)\varphi(p)=v(f\circ\varphi)(p)$.
\end{definition}

\begin{lemma}[existence and uniqueness of the bracket field]
  \label{lem:dc-ch0-5-2}
  \lean{Riemannian.exists_unique_bracketField, Riemannian.bracketField_dir}
  \leanok
  \dcref{ch0:5.2}
  Let $X,Y$ be differentiable vector fields on $M$. Then there exists a unique
  differentiable vector field $Z$ such that $Zf=(XY-YX)f$ for all
  $f\in\mathcal{D}$. This field is the \emph{bracket} $[X,Y]=XY-YX$.
  Existence, uniqueness, and smoothness are isolated in
  \cref{lem:dc-ch0-5-2-commutator,lem:dc-ch0-5-2-unique,lem:dc-ch0-5-2-smooth};
  the structural identities are treated in \cref{lem:dc-ch0-5-3-antisymm,lem:dc-ch0-5-3-jacobi}.
\end{lemma}
\begin{proof}
  \sketch
  \leanok
  \uses{lem:dc-ch0-5-2-unique,lem:dc-ch0-5-2-commutator,lem:dc-ch0-5-2-smooth}
  In a parametrization $\mathbf{x}$ at $p$ with
  $X=\sum_i a_i\frac{\partial}{\partial x_i}$,
  $Y=\sum_j b_j\frac{\partial}{\partial x_j}$,
  
  $$
  XYf=\sum_{i,j}a_i\frac{\partial b_j}{\partial x_i}\frac{\partial f}{\partial x_j}+\sum_{i,j}a_i b_j\frac{\partial^2 f}{\partial x_i\partial x_j},
  $$
  
  and symmetrically for $YXf$; the second-order terms cancel, giving
  
  $$
  Zf=XYf-YXf=\sum_{i,j}\Big(a_i\frac{\partial b_j}{\partial x_i}-b_i\frac{\partial a_j}{\partial x_i}\Big)\frac{\partial f}{\partial x_j},
  $$
  
  proving uniqueness. For existence, define $Z_\alpha$ on each
  $\mathbf{x}_\alpha(U_\alpha)$ by this expression; by uniqueness $Z_\alpha=Z_\beta$
  on overlaps, so $Z$ is defined over all of $M$. The coordinate formula shows
  that it is differentiable.
\end{proof}

\begin{lemma}[existence characterization: the bracket acts as the commutator of derivations]
  \label{lem:dc-ch0-5-2-commutator}
  \lean{Riemannian.mfderiv_mlieBracket_eq_commutator}
  \uses{lem:dc-ch0-5-2}
  \leanok
  \dcref{ch0:5.2}
  Let $X,Y$ be smooth vector fields on a boundaryless manifold $M$ and $f$ a smooth
  real function. Then the bracket field $[X,Y]$ acts on $f$ as the
  commutator of the derivations $X$ and $Y$:
  $$
  [X,Y]f = X(Yf)-Y(Xf), \qquad\text{i.e.}\qquad
  \mathrm{d}f_p([X,Y]_p) = \mathrm{d}(\mathrm{d}f\,Y)_p(X_p)-\mathrm{d}(\mathrm{d}f\,X)_p(Y_p).
  $$
  Thus $[X,Y]$ realizes the derivation $(XY-YX)$.
\end{lemma}
\begin{proof}
  \sketch
  \leanok
  In a chart around $p$, the boundaryless hypothesis identifies the model set with the
  model space. Pull back $X,Y$ to this chart and apply the Euclidean commutator identity
  $\mathrm{d}f([V,W])=\mathrm{d}(\mathrm{d}f\,W)(V)-\mathrm{d}(\mathrm{d}f\,V)(W)$;
  the chain rule transports the result back to $M$.
\end{proof}

\begin{lemma}[uniqueness: a vector field is determined by its action on functions]
  \label{lem:dc-ch0-5-2-unique}
  \lean{Riemannian.tangentSection_eq_of_forall_mfderiv}
  \uses{lem:dc-ch0-5-2}
  \leanok
  \dcref{ch0:5.2}
  A tangent vector at $p$, and hence a vector field on $M$, is determined by its
  action on differentiable real functions: if $Z_1,Z_2$ are sections of $TM$ with
  $Z_1f=Z_2f$ for every $f\in\mathcal{D}$, then $Z_1=Z_2$; in particular at most one
  field acts as $(XY-YX)$.
\end{lemma}
\begin{proof}
  \sketch
  \leanok
  It suffices to prove the pointwise statement: if $v\in T_pM$ satisfies
  $\mathrm{d}f_p(v)=0$ for every $f$ differentiable at $p$, then $v=0$
  (apply it to $Z_1(p)-Z_2(p)$ using linearity of $\mathrm{d}f_p$). For a
  continuous linear functional $L$ on the model space $E$, the coordinate function
  $f_L=L\circ\mathbf{x}$
  is differentiable at $p$, and since $\mathrm{d}\mathbf{x}_p=\mathrm{id}$ on
  $T_pM$ one has $\mathrm{d}(f_L)_p(v)=L(v)$. Thus $L(v)=0$ for all such $L$, and
  continuous functionals separate points of a normed space, so $v=0$.
\end{proof}

\begin{lemma}[smoothness of the bracket field]
  \label{lem:dc-ch0-5-2-smooth}
  \lean{Riemannian.DCLieBracket_smoothAt}
  \uses{lem:dc-ch0-5-2}
  \leanok
  \dcref{ch0:5.2}
  If $X$ and $Y$ are smooth vector fields, then their bracket $[X,Y]$ is a smooth
  vector field.
\end{lemma}
\begin{proof}
  \sketch
  \leanok
  In each parametrization, write
  $X=\sum_i a_i\frac{\partial}{\partial x_i}$ and
  $Y=\sum_j b_j\frac{\partial}{\partial x_j}$. By
  \cref{lem:dc-ch0-5-2}, the bracket has coefficients
  $\sum_i(a_i\partial_i b_j-b_i\partial_i a_j)$. These coefficients are smooth
  because the coefficients of $X$ and $Y$ are smooth, so the local expressions
  define a smooth vector field on each chart. Since the bracket was constructed
  uniquely, the local smooth fields agree on overlaps and glue to a smooth global
  vector field.
\end{proof}

\begin{lemma}[antisymmetry of the bracket]
  \label{lem:dc-ch0-5-3-antisymm}
  \lean{Riemannian.DCLieBracket_antisymm}
  \uses{lem:dc-ch0-5-2}
  \leanok
  \dcref{ch0:5.3}
  For differentiable vector fields $X$ and $Y$ one has $[X,Y]=-[Y,X]$.
\end{lemma}
\begin{proof}
  \sketch
  \leanok
  By the defining identity of \cref{lem:dc-ch0-5-2},
  $[X,Y]f=XYf-YXf$ and $[Y,X]f=YXf-XYf$ for every differentiable function $f$.
  Hence $[X,Y]f=-[Y,X]f$ for all such $f$, and uniqueness in
  \cref{lem:dc-ch0-5-2} gives $[X,Y]=-[Y,X]$.
\end{proof}

\begin{lemma}[Jacobi identity in Leibniz form]
  \label{lem:dc-ch0-5-3-jacobi}
  \lean{Riemannian.DCLieBracket_jacobi}
  \uses{lem:dc-ch0-5-2}
  \leanok
  \dcref{ch0:5.3}
  For differentiable vector fields $X,Y,Z$ one has
  $[X,[Y,Z]]=[[X,Y],Z]+[Y,[X,Z]]$.
\end{lemma}
\begin{proof}
  \sketch
  \leanok
  Evaluate both sides on a differentiable function $f$. Expanding each bracket
  as a commutator of derivations gives
  $XYZf-XZYf-YZXf+ZYXf$ on the left and the same four terms on the right after
  cancellation of the repeated middle terms. By uniqueness in
  \cref{lem:dc-ch0-5-2}, the two vector fields are equal.
\end{proof}

\begin{lemma}[Jacobi identity in cyclic form]
  \label{lem:dc-ch0-5-3-jacobi-cyclic}
  \lean{Riemannian.DCLieBracket_jacobi_cyclic}
  \uses{lem:dc-ch0-5-3-jacobi, lem:dc-ch0-5-3-antisymm}
  \leanok
  \dcref{ch0:5.3}
  For differentiable vector fields $X,Y,Z$ one has the cyclic Jacobi identity
  $[[X,Y],Z]+[[Y,Z],X]+[[Z,X],Y]=0$.
\end{lemma}
\begin{proof}
  \sketch
  \leanok
  Anticommutativity (\cref{lem:dc-ch0-5-3-antisymm}) rewrites each outer bracket as
  $[[X,Y],Z]=-[Z,[X,Y]]$, so the cyclic sum equals
  $-\big([Z,[X,Y]]+[X,[Y,Z]]+[Y,[Z,X]]\big)$. The Leibniz-form identity
  $[X,[Y,Z]]=[[X,Y],Z]+[Y,[X,Z]]$ of \cref{lem:dc-ch0-5-3-jacobi}, together with
  $[X,Z]=-[Z,X]$, makes the parenthesized derivation sum vanish, giving $0$.
\end{proof}

\begin{lemma}[linearity of the bracket]
  \label{lem:dc-ch0-5-3-bilinear}
  \lean{Riemannian.DCLieBracket_linear_left}
  \uses{lem:dc-ch0-5-2}
  \leanok
  \dcref{ch0:5.3}
  For smooth vector fields $X,Y,Z$ and scalars $a,b\in\mathbb{R}$ one has
  $[aX+bY,Z]=a[X,Z]+b[Y,Z]$.
\end{lemma}
\begin{proof}
  \sketch
  \leanok
  Apply additivity of the bracket in its first slot to the sum $aX+bY$, then pull
  out the constants $a$ and $b$ by homogeneity in the first slot. Both facts follow
  from the local commutator expression of \cref{lem:dc-ch0-5-2}.
\end{proof}

\begin{lemma}[Leibniz rule for the bracket]
  \label{lem:dc-ch0-5-3-leibniz}
  \lean{Riemannian.DCLieBracket_leibniz}
  \uses{lem:dc-ch0-5-2}
  \leanok
  \dcref{ch0:5.3}
  For smooth vector fields $X,Y$ and smooth functions $f,g$ one has
  $[fX,gY]=fg[X,Y]+fX(g)Y-gY(f)X$, where $X(g)=dg_p(X_p)$ denotes the derivative
  of $g$ along $X$.
\end{lemma}
\begin{proof}
  \sketch
  \leanok
  Expand the bracket by the first-slot Leibniz identity
  $[fX,W]=f[X,W]-W(f)X$ with $W=gY$, then the second-slot Leibniz identity
  $[X,gY]=g[X,Y]+X(g)Y$; combining and using linearity of the differential
  $dg_p$ gives $[fX,gY]=fg[X,Y]+fX(g)Y-gY(f)X$.
\end{proof}

\begin{proposition}[properties of the bracket]
  \label{prop:dc-ch0-5-3}
  \lean{Riemannian.DCLieBracket_properties}
  \uses{lem:dc-ch0-5-3-antisymm, lem:dc-ch0-5-3-jacobi-cyclic, lem:dc-ch0-5-3-bilinear, lem:dc-ch0-5-3-leibniz}
  \leanok
  \dcref{ch0:5.3}
  If $X,Y,Z$ are differentiable vector fields on $M$, $a,b\in\mathbb{R}$, and $f,g$
  differentiable functions, then:
  (a) $[X,Y]=-[Y,X]$ \emph{(anticommutativity)};
  (b) $[aX+bY,Z]=a[X,Z]+b[Y,Z]$ \emph{(linearity)};
  (c) $[[X,Y],Z]+[[Y,Z],X]+[[Z,X],Y]=0$ \emph{(Jacobi identity)};
  (d) $[fX,gY]=fg[X,Y]+fX(g)Y-gY(f)X$.
\end{proposition}
\begin{proof}
  \sketch
  (a),(b) are immediate. For (c), expand
  $[[X,Y],Z]=XYZ-YXZ-ZXY+ZYX$ and
  $[X,[Y,Z]]+[Y,[Z,X]]=XYZ-XZY-YZX+ZYX+YZX-YXZ-ZXY+XZY$; the second members
  agree, and (c) follows using (a). For (d),
  $[fX,gY]=fX(gY)-gY(fX)=fgXY+fX(g)Y-gfYX-gY(f)X=fg[X,Y]+fX(g)Y-gY(f)X$.
\end{proof}

For a differentiable vector field $X$ and $p\in M$, there exist a neighborhood
$U$ of $p$, an interval $(-\delta,\delta)$, and a differentiable map
$\varphi:(-\delta,\delta)\times U\to M$ such that $t\mapsto\varphi(t,q)$ is the
unique curve with
$\frac{\partial\varphi}{\partial t}=X(\varphi(t,q))$ and $\varphi(0,q)=q$. A
curve $\alpha$ with $\alpha'(t)=X(\alpha(t))$, $\alpha(0)=q$, is a
\emph{trajectory} of $X$, and $\varphi_t(q)=\varphi(t,q)$ is the \emph{local
flow} of $X$.

\begin{proposition}[bracket as Lie derivative along the flow]
  \label{prop:dc-ch0-5-4}
  \lean{Riemannian.DCLieBracket_eq_flow_lieDerivative}
  \uses{lem:dc-ch0-5-5, lem:dc-ch0-5-2}
  \leanok
  \dcref{ch0:5.4}
  Let $X,Y$ be differentiable vector fields on $M$, $p\in M$, and $\varphi_t$ the
  local flow of $X$ in a neighborhood $U$ of $p$. Then

  $$
  [X,Y](p)=\lim_{t\to0}\frac1t\,[Y-d\varphi_t Y](\varphi_t(p)).
  $$
\end{proposition}
\begin{proof}
  \sketch
  \leanok
  For $f$ differentiable near $p$, put $h(t,q)=f(\varphi_t(q))-f(q)$; by
  \cref{lem:dc-ch0-5-5} there is differentiable $g(t,q)$ with
  $f\circ\varphi_t(q)=f(q)+tg(t,q)$ and $g(0,q)=Xf(q)$. Then
  $((d\varphi_t Y)f)(\varphi_t(p))=(Y(f\circ\varphi_t))(p)=Yf(p)+t(Yg(t,p))$, so

  $$
  \lim_{t\to0}\frac1t[Y-d\varphi_t Y]f(\varphi_t p)=\lim_{t\to0}\frac{(Yf)(\varphi_t p)-Yf(p)}{t}-(Yg(0,p))=(X(Yf))(p)-(Y(Xf))(p)=([X,Y]f)(p).
  $$
\end{proof}

\begin{lemma}[Hadamard-type lemma from calculus]
  \label{lem:dc-ch0-5-5}
  \lean{Riemannian.DCHadamardFactor, Riemannian.DCHadamardFactorDiff}
  \leanok
  \dcref{ch0:5.5}
  Let $h:(-\delta,\delta)\times U\to\mathbb{R}$ be differentiable with $h(0,q)=0$
  for all $q\in U$. Then there exists a differentiable
  $g:(-\delta,\delta)\times U\to\mathbb{R}$ with $h(t,q)=t\,g(t,q)$; in particular
  $g(0,q)=\frac{\partial h(t,q)}{\partial t}\big|_{t=0}$.
\end{lemma}
\begin{proof}
  \sketch
  \leanok
  Define $g(t,q)=\int_0^1\frac{\partial h(ts,q)}{\partial(ts)}\,ds$; after
  changing variables, $t\,g(t,q)=\int_0^t\frac{\partial h(ts,q)}{\partial(ts)}\,d(ts)=h(t,q)$.
  Differentiability of $g$ follows by differentiation under the integral sign
  (the pointwise Fréchet-derivative form of the dominated-convergence theorem, the
  domination bound coming from continuity of $\frac{\partial h}{\partial t}$ on the
  compact set swept by $s\in[0,1]$ and $q$ in a closed ball), and
  $g(0,q)=\int_0^1\frac{\partial h}{\partial t}(0,q)\,ds=\frac{\partial h}{\partial t}(0,q)$.
\end{proof}

A manifold need not be Hausdorff or have a countable coordinate basis. For a
connected manifold these two conditions are equivalent to existence of a
differentiable partition of unity (\cref{thm:dc-ch0-5-6}). Whitney's theorem
gives, for an $n$-manifold with both conditions, an immersion into
$\mathbb{R}^{2n}$ and an embedding into $\mathbb{R}^{2n+1}$; for $n>1$ the
bounds refine to $\mathbb{R}^{2n-1}$ and $\mathbb{R}^{2n}$, respectively.

A family of open sets $V_\alpha\subset M$ with $\bigcup_\alpha V_\alpha=M$ is
\emph{locally finite} if every $p\in M$ has a neighborhood $W$ meeting only
finitely many $V_\alpha$. The \emph{support} of $f:M\to\mathbb{R}$ is the
closure of $\{f\neq0\}$. A family $\{f_\alpha\}$ of differentiable functions is
a \emph{differentiable partition of unity} if $f_\alpha\ge0$, each support lies
in a coordinate neighborhood $V_\alpha$, the family $\{V_\alpha\}$ is locally
finite, and $\sum_\alpha f_\alpha(p)=1$ for every $p\in M$. Such a partition is
\emph{subordinate} to $\{V_\alpha\}$.

\begin{theorem}[existence of partitions of unity]
  \label{thm:dc-ch0-5-6}
  \notready
  \dcref{ch0:5.6}
  A differentiable manifold $M$ has a differentiable partition of unity if and only
  if every connected component of $M$ is Hausdorff and has a countable basis.
\end{theorem}
\begin{proof}
  \uses{lem:dc-ch0-5-6-exists, lem:dc-ch0-5-6-onlyif-complete, lem:dc-ch0-5-6-component}
  \sketch
  For the forward implication, Hausdorffness and a countable basis yield a subordinate
  partition of unity by \cref{lem:dc-ch0-5-6-exists}. Conversely, a partition forces
  Hausdorffness by \cref{lem:dc-ch0-5-6-hausdorff} and second countability by
  \cref{lem:dc-ch0-5-6-onlyif-pou}, using
  \cref{lem:dc-ch0-5-6-sigmacompact-core}. These conclusions combine in
  \cref{lem:dc-ch0-5-6-onlyif-complete}.
\end{proof}

\begin{lemma}[connected space with a locally finite $\sigma$-compact open cover]
  \label{lem:dc-ch0-5-6-sigmacompact-core}
  \lean{Riemannian.sigmaCompactSpace_of_locallyFinite_isSigmaCompact_cover}
  \uses{thm:dc-ch0-5-6}
  \leanok
  \dcref{ch0:5.6}
  Let $X$ be a connected topological space carrying a locally finite open cover
  $\{V_i\}_{i\in\iota}$ each of whose members is $\sigma$-compact. Then $X$ is
  $\sigma$-compact.
\end{lemma}
\begin{proof}
  \sketch
  \leanok
  If every $V_i$ is empty their union is empty, so $X=\varnothing$ is
  $\sigma$-compact; assume then $V_{i_0}\neq\varnothing$ for some $i_0$. Because each
  $V_i=\bigcup_n K_{i,n}$ is a countable union of compacts and the family
  $\{V_j\}$ is locally finite, each compact $K_{i,n}$ meets only finitely many
  $V_j$; hence the overlap set $\{j : V_i\cap V_j\neq\varnothing\}$ is countable for
  every $i$.
  Starting from $\{i_0\}$ and repeatedly adjoining all indices whose $V_j$ meets a
  set already chosen produces, after countably many layers, a countable index set
  $A$ that is closed under overlaps: if $i\in A$ and $V_i\cap V_j\neq\varnothing$
  then $j\in A$. Put $U=\bigcup_{i\in A}V_i$. Then $U$ is open (a union of opens) and
  closed: any $x\in\overline U$ lies in some $V_j$ (the $V$'s cover $X$), and this
  open $V_j$ meets $U$, so $V_j\cap V_i\neq\varnothing$ for some $i\in A$, whence
  $j\in A$ and $x\in V_j\subseteq U$. Thus $U$ is clopen and nonempty, so $U=X$ by
  connectedness. Finally $U=\bigcup_{i\in A}V_i$ is a countable union of
  $\sigma$-compact sets, so $X$ is $\sigma$-compact.
\end{proof}

\begin{lemma}[second countability from a partition of unity]
  \label{lem:dc-ch0-5-6-onlyif}
  \lean{Riemannian.DCSecondCountable_of_locallyFinite_isSigmaCompact_cover}
  \uses{lem:dc-ch0-5-6-sigmacompact-core, thm:dc-ch0-5-6}
  \leanok
  \dcref{ch0:5.6}
  Let $M$ be a connected manifold charted over a second-countable model $H$. If $M$
  carries a locally finite open cover by $\sigma$-compact sets — in particular, the
  coordinate neighborhoods $\{V_\alpha\}$ of a differentiable partition of unity,
  which are locally finite, cover $M$, and are each homeomorphic to an open subset
  of $\mathbb{R}^n$ (hence $\sigma$-compact) — then $M$ has a countable basis.
\end{lemma}
\begin{proof}
  \sketch
  \leanok
  By \cref{lem:dc-ch0-5-6-sigmacompact-core} the manifold $M$ is $\sigma$-compact.
  A charted space over a second-countable model that is $\sigma$-compact is itself
  second countable (cover each of the countably many compacts by finitely many chart
  domains and use the countable basis of $H$ inside each). Applying this argument to
  each connected component gives the second-countability
  conclusion in the theorem.
\end{proof}

\begin{lemma}[second countability from an honest partition of unity]
  \label{lem:dc-ch0-5-6-onlyif-pou}
  \lean{Riemannian.DCSecondCountable_of_partitionOfUnity}
  \uses{lem:dc-ch0-5-6-onlyif, thm:dc-ch0-5-6}
  \leanok
  \dcref{ch0:5.6}
  Let $M$ be a connected manifold charted over a second-countable model $H$, and let
  $\{f_\alpha\}$ be a differentiable partition of unity on $M$ (with $\sum_\alpha f_\alpha
  \equiv 1$) subordinate to a locally finite family $\{V_\alpha\}$ of $\sigma$-compact
  coordinate neighborhoods. Then $M$ has a countable basis.
\end{lemma}
\begin{proof}
  \sketch
  \leanok
  Unlike \cref{lem:dc-ch0-5-6-onlyif}, the covering hypothesis is not assumed but
  \emph{derived} from the partition. For each $p\in M$, condition $\sum_\alpha f_\alpha(p)=1$
  gives some index $\alpha$ with $f_\alpha(p)\neq0$, so $p\in\operatorname{supp}f_\alpha
  \subseteq\operatorname{tsupp}f_\alpha\subseteq V_\alpha$ by subordinacy; hence
  $\{V_\alpha\}$ covers $M$. As these neighborhoods are open, $\sigma$-compact, and locally
  finite, \cref{lem:dc-ch0-5-6-onlyif} (equivalently
  \cref{lem:dc-ch0-5-6-sigmacompact-core} followed by charted second-countability) yields a
  countable basis for $M$.
\end{proof}

\begin{lemma}[a partition of unity forces the Hausdorff axiom]
  \label{lem:dc-ch0-5-6-hausdorff}
  \lean{Riemannian.t2Space_of_partitionOfUnity}
  \uses{thm:dc-ch0-5-6}
  \leanok
  \dcref{ch0:5.6}
  Let $M$ be a manifold admitting a differentiable partition of unity $\{f_\alpha\}$
  subordinate to a family $\{V_\alpha\}$ of open, Hausdorff coordinate neighborhoods. Then
  $M$ is Hausdorff.
\end{lemma}
\begin{proof}
  \sketch
  \leanok
  Let $p\neq q$ have no disjoint neighborhoods; then the filter $\mathcal F=\mathcal N_p
  \sqcap\mathcal N_q$ is proper. Each $f_\alpha$ is continuous, so it tends to both
  $f_\alpha(p)$ and $f_\alpha(q)$ along $\mathcal F$; uniqueness of limits in $\mathbb R$
  forces $f_\alpha(p)=f_\alpha(q)$. Since $\sum_\alpha f_\alpha(p)=1$, some $f_\alpha(p)\neq
  0$, whence $f_\alpha(q)\neq0$ too, so both $p,q\in\operatorname{supp}f_\alpha\subseteq
  V_\alpha$. But $V_\alpha$ is open and Hausdorff, so $p,q$ are separated by sets open in
  $V_\alpha$ and hence open in $M$ — contradicting the choice of $p,q$. Thus $M$ is
  Hausdorff.
\end{proof}

\begin{lemma}[complete ``only if'' direction from a partition of unity]
  \label{lem:dc-ch0-5-6-onlyif-complete}
  \lean{Riemannian.DCManifoldAxioms_of_partitionOfUnity}
  \uses{lem:dc-ch0-5-6-hausdorff, lem:dc-ch0-5-6-onlyif-pou}
  \leanok
  \dcref{ch0:5.6}
  Let $M$ be a connected manifold charted over a second-countable model $H$ admitting a
  differentiable partition of unity subordinate to a locally finite family of open,
  Hausdorff, $\sigma$-compact coordinate neighborhoods. Then $M$ is Hausdorff and has a
  countable basis.
\end{lemma}
\begin{proof}
  \sketch
  \leanok
  Immediate from \cref{lem:dc-ch0-5-6-hausdorff} (Hausdorff) and
  \cref{lem:dc-ch0-5-6-onlyif-pou} (countable basis), component by component.
\end{proof}

\begin{lemma}[existence of a subordinate partition of unity]
  \label{lem:dc-ch0-5-6-exists}
  \lean{Riemannian.DCSmoothPartitionOfUnity.exists_isSubordinate}
  \uses{thm:dc-ch0-5-6}
  \leanok
  \dcref{ch0:5.6}
  Let $M$ be a Hausdorff, second-countable (equivalently $\sigma$-compact),
  finite-dimensional differentiable manifold. For every closed set $s\subseteq M$
  and every open covering $\{U_i\}$ of $s$ there is a differentiable partition of
  unity $\{f_i\}$ subordinate to $\{U_i\}$, i.e. with
  $\operatorname{supp} f_i\subseteq U_i$.
\end{lemma}
\begin{proof}
  \sketch
  \leanok
  Hausdorffness and second countability make $M$ paracompact and $\sigma$-compact.
  Local Euclidean bump functions supported in the $U_i$ form a locally finite family;
  normalizing their sum yields the required partition.
\end{proof}

\begin{lemma}[normalization of a partition of unity]
  \label{lem:dc-ch0-5-6-sum}
  \lean{Riemannian.DCSmoothPartitionOfUnity.sum_eq_one}
  \leanok
  \dcref{ch0:5.6}
  Let $\{f_\alpha\}$ be a differentiable partition of unity on $M$ subordinate to a
  covering of a set $s\subseteq M$. Then $\sum_\alpha f_\alpha(p)=1$ for every
  $p\in s$.
\end{lemma}
\begin{proof}
  \sketch
  \leanok
  This is condition (3) in the definition of a differentiable partition of unity:
  the family is locally finite, so the sum is finite at each point, and by
  definition it takes the value $1$ throughout the covered set $s$.
\end{proof}

\begin{remark}[bump functions and locality]
  \label{rem:dc-ch0-5-7}
  \lean{Riemannian.exists_ch0_bump_function, Riemannian.exists_ch0_manifold_bump_function}
  \leanok
  \dcref{ch0:5.7}
  Given $p\in\mathbb{R}^n$ and an open ball $B_r(p)\subset\mathbb{R}^n$, there exist
  a neighborhood $U$ of $p$ with $\overline U\subset B_r(p)$ and a differentiable
  $f:\mathbb{R}^n\to\mathbb{R}$ with $0\le f(q)\le1$ and
  $$
  f(q)=\begin{cases}1,&q\in\overline U,\\ 0,&q\notin B_r(p).\end{cases}
  $$
  For $r=2$ take $U=B_1(p)$ and $f(q)=\beta(-|p-q|)$, where
  $\beta(t)=\frac{\int_{-\infty}^t\alpha(s)\,ds}{\int_{-2}^{-1}\alpha(s)\,ds}$ and
  $\alpha:\mathbb{R}\to\mathbb{R}$ equals $\exp\big(-\frac{1}{(t+2)(-1-t)}\big)$ for
  $-2<t<-1$ and zero otherwise. The same holds on $M$: if $V\subset M$ is a
  neighborhood of $p$ contained in a coordinate neighborhood diffeomorphic to an
  open ball, there exist $U$ with $\overline U\subset V$ and differentiable
  $f:M\to\mathbb{R}$, $0\le f\le1$, $f=1$ on $\overline U$, $f=0$ off $V$.
  These bump functions localize globally defined constructions.
\end{remark}

\begin{example}[a non-differentiable curve]
  \label{ex:dc-ch0-3-2}
  \dcref{ch0:3.2}
  \lean{Riemannian.ch0AbsCurve, Riemannian.ch0AbsCurve_not_differentiableAt_zero}
  \leanok
  The curve $\alpha:\mathbb{R}\to\mathbb{R}^2$, $\alpha(t)=(t,|t|)$, is not
  differentiable at $t=0$ (Fig. 5).
\end{example}

\begin{example}[differentiable but not an immersion]
  \label{ex:dc-ch0-3-3}
  \dcref{ch0:3.3}
  \lean{Riemannian.ch0CuspCurve, Riemannian.ch0CuspCurve_differentiable,
    Riemannian.ch0CuspCurve_fderiv_zero, Riemannian.ch0CuspCurve_not_immersion}
  \leanok
  $\alpha:\mathbb{R}\to\mathbb{R}^2$, $\alpha(t)=(t^3,t^2)$, is differentiable but
  not an immersion: the immersion condition is $\alpha'(t)\neq0$, which fails at
  $t=0$ (Fig. 6).
\end{example}

\begin{example}[immersion with self-intersection]
  \label{ex:dc-ch0-3-4}
  \dcref{ch0:3.4}
  \lean{Riemannian.ch0SelfIntersectionCurve,
    Riemannian.ch0SelfIntersectionCurve_hasDerivAt,
    Riemannian.ch0SelfIntersectionCurve_differentiable,
    Riemannian.ch0SelfIntersectionCurve_deriv_ne_zero,
    Riemannian.ch0SelfIntersectionCurve_self_intersects}
  \leanok
  $\alpha(t)=(t^3-4t,\,t^2-4)$ (Fig. 7) is an immersion
  $\alpha:\mathbb{R}\to\mathbb{R}^2$ with a self-intersection at $t=2,t=-2$. Hence
  $\alpha$ is not an embedding.
\end{example}

\begin{example}[injective immersion that is not an embedding]
  \label{ex:dc-ch0-3-5}
  \dcref{ch0:3.5}
  \notready
  The curve $\alpha:(-3,0)\to\mathbb{R}^2$ (Fig. 8), defined piecewise as the
  segment $(0,-(t+2))$ for $t\in(-3,-1)$, a regular joining curve for
  $t\in(-1,-\frac1\pi)$, and $(-t,-\sin\frac1t)$ for $t\in(-\frac1\pi,0)$, is an
  immersion without self-intersections that is not an embedding: a neighborhood of
  a point $p$ in the vertical part has infinitely many connected components in the
  topology induced from $\mathbb{R}^2$, whereas in the topology induced from
  $\alpha$ (the line) it is an open interval, hence connected.
\end{example}

\begin{example}[the sphere $S^n$ is orientable]
  \label{ex:dc-ch0-4-6}
  \dcref{ch0:4.6}
  \notready
  \uses{ex:dc-ch0-4-5}
  $S^n=\{(x_1,\dots,x_{n+1})\in\mathbb{R}^{n+1};\sum_{i=1}^{n+1}x_i^2=1\}$. With
  north/south poles $N=(0,\dots,0,1)$, $S=(0,\dots,0,-1)$, the *stereographic
  projections* $\pi_1:S^n-\{N\}\to\mathbb{R}^n$,
  $\pi_1(x)=\big(\frac{x_1}{1-x_{n+1}},\dots,\frac{x_n}{1-x_{n+1}}\big)$, and
  $\pi_2:S^n-\{S\}\to\mathbb{R}^n$ (from the south pole) are differentiable
  injections onto $x_{n+1}=0$. The parametrizations
  $(\mathbb{R}^n,\pi_1^{-1}),(\mathbb{R}^n,\pi_2^{-1})$ cover $S^n$ with change of
  coordinates $y_j'=\frac{y_j}{\sum_{i=1}^n y_i^2}$, differentiable. For $n\geq2$,
  $\pi_1^{-1}(\mathbb{R}^n)\cap\pi_2^{-1}(\mathbb{R}^n)=S^n-\{N,S\}$ is connected,
  so by \cref{ex:dc-ch0-4-5}, $S^n$ is orientable. For $n=1$, the same atlas is
  orientable directly: on the two components of $S^1-\{N,S\}$, the change of
  coordinates is $y'=1/y$ with derivative $-1/y^2<0$, and changing the sign of one
  coordinate makes both transition derivatives positive. The antipodal map
  $A:S^n\to S^n$, $A(p)=-p$, is a diffeomorphism ($A^2=\mathrm{id}$); it reverses the
  orientation of $S^n$ when $n$ is even and preserves it when $n$ is odd.
\end{example}

\begin{example}[special cases of Example 4.8]
  \label{ex:dc-ch0-4-9}
  \dcref{ch0:4.9}
  \notready
  \uses{ex:dc-ch0-4-8}
  **4.9 (a) (the $k$-torus).** Let $G$ be the group of "integral" translations of
  $\mathbb{R}^k$: for $(n_1,\dots,n_k)\in\mathbb{Z}^k$, the corresponding
  translation sends $(x_1,\dots,x_k)$ to $(x_1+n_1,\dots,x_k+n_k)$.
  This is a properly discontinuous action; $\mathbb{R}^k/G$ (with the structure of
  \cref{ex:dc-ch0-4-8}) is the *$k$-torus $T^k$*. For $k=2$, $T^2$ is diffeomorphic to the
  torus of revolution in $\mathbb{R}^3$, the inverse image of $0$ of
  $f(x,y,z)=z^2+(\sqrt{x^2+y^2}-a)^2-r^2$.
  \textbf{4.9 (b) (Klein bottle; Möbius band).} Let $S\subset\mathbb{R}^3$ be a regular
  surface symmetric about $0$ (if $p\in S$ then $-p=A(p)\in S$). The group
  $\{A,\mathrm{Id}\}$ acts properly discontinuously on $S$; $S/G$ carries the
  structure of \cref{ex:dc-ch0-4-8}. When $S$ is the torus of revolution $T^2$, $S/G=K$ is
  \emph{the Klein bottle}; when $S=C=\{(x,y,z); x^2+y^2=1,\,-1<z<1\}$ (a right circular
  cylinder), $S/G$ is the \emph{Möbius band}. By Exercise 9 the Klein bottle and Möbius
  band are non-orientable.
\end{example}

\begin{lemma}[fixed-cover characterization on an open connected component]
  \label{lem:dc-ch0-5-6-component}
  \dcref{ch0:5.6}
  \uses{lem:dc-ch0-5-6-fixed-cover}
  \lean{Riemannian.DCPartitionOfUnity_iff_manifoldAxioms_on_connectedComponent}
  \leanok
  Let $p\in M$ and suppose its connected component $C_p$ is open in $M$. For a locally finite
  cover of $C_p$ by open Hausdorff $\sigma$-compact sets, a smooth partition of unity on $C_p$
  subordinate to this fixed cover exists if and only if $C_p$ is Hausdorff and has a countable
  basis. The cover and partition are genuinely component-local; no global countability or
  cross-component assembly is asserted.
\end{lemma}
\begin{proof}
  \leanok
  Regard $C_p$ as the open subspace of $M$. Connected components are preconnected, so this
  subspace carries the connectedness instance required by
  \cref{lem:dc-ch0-5-6-fixed-cover}; applying that characterization to the inherited manifold
  structure gives the result.
\end{proof}

\begin{lemma}[fixed-cover partition-of-unity characterization]
  \label{lem:dc-ch0-5-6-fixed-cover}
  \dcref{ch0:5.6}
  \lean{Riemannian.DCPartitionOfUnity_iff_manifoldAxioms}
  \leanok
  Let $M$ be a connected finite-dimensional manifold, charted over a second-countable model,
  and let $\{V_i\}_{i\in\iota}$ be a locally finite cover by open Hausdorff $\sigma$-compact
  sets. Then a smooth partition of unity subordinate to this fixed cover exists if and only if
  $M$ is Hausdorff and has a countable basis.
\end{lemma}
\begin{proof}
  \leanok
  This is the connected fixed-cover form of the theorem: the forward implication obtains the
  Hausdorff and second-countability axioms from the partition, while the reverse implication
  applies the smooth subordinate-partition construction to the closed set $M$ itself.
\end{proof}

\chapter{Riemannian Metrics}
\label{chap:dc-riemannian-metrics}

Throughout this chapter, manifolds are Hausdorff and have a countable basis;
``differentiable'' means $C^\infty$, and $M^n$ denotes an $n$-dimensional
manifold $M$.

\section{Riemannian Metrics}

\begin{definition}[Riemannian metric]
  \label{def:dc-ch1-2-1}
  \lean{Riemannian.RiemannianMetric}
  \leanok
  \dcref{ch1:2.1}
  A \emph{Riemannian metric} (or \emph{Riemannian structure}) on a differentiable
  manifold $M$ assigns an inner product $\langle\ ,\ \rangle_p$ on each $T_pM$,
  varying differentiably with $p$. In a coordinate system
  $\mathbf{x}:U\subset\mathbb{R}^n\to M$, for
  $q=\mathbf{x}(x_1,\ldots,x_n)$ and
  $\frac{\partial}{\partial x_i}(q)=d\mathbf{x}_{(x_1,\ldots,x_n)}e_i$,
  the coefficients
  $$
  g_{ij}(x_1,\ldots,x_n)=
  \left\langle\frac{\partial}{\partial x_i}(q),
  \frac{\partial}{\partial x_j}(q)\right\rangle_q
  $$
  are differentiable on $U$. This condition is independent of the chart;
  equivalently, $\langle X,Y\rangle$ is differentiable for differentiable vector
  fields $X,Y$. The functions $g_{ij}=g_{ji}$ are the \emph{local representation
  of the metric}; $M$ with such a metric is a \emph{Riemannian manifold}.
\end{definition}

\begin{definition}[isometry]
  \label{def:dc-ch1-2-2}
  \lean{Riemannian.DCIsometry, Riemannian.DCIsometry.preservesMetric}
  \uses{def:dc-ch1-2-1}
  \leanok
  \dcref{ch1:2.2}
  For Riemannian manifolds $M,N$, a diffeomorphism $f:M\to N$ is an
  \emph{isometry} when
  $$
  \langle u,v\rangle_p=\langle df_p(u),df_p(v)\rangle_{f(p)},\qquad\text{(1)}
  $$
  for every $p\in M$ and $u,v\in T_pM$.
\end{definition}

\begin{definition}[local isometry]
  \label{def:dc-ch1-2-3}
  \lean{Riemannian.DCIsLocalIsometryAtFull, Riemannian.DCIsLocalIsometryAtFull.isLocalDiffeomorphAt, Riemannian.DCIsLocalIsometryAtFull.preservesMetricNear}
  \uses{def:dc-ch1-2-1, def:dc-ch1-2-2}
  \leanok
  \dcref{ch1:2.3}
  A differentiable map $f:M\to N$ between Riemannian manifolds is a
  \emph{local isometry at} $p\in M$ if some neighborhood $U$ of $p$ is mapped
  diffeomorphically onto $f(U)$ and the restriction satisfies (1). The manifold
  $M$ is \emph{locally isometric} to $N$ if every $p\in M$ has such a neighborhood.
\end{definition}

\begin{example}[Euclidean space]
  \label{ex:dc-ch1-2-4}
  \lean{Riemannian.DCEuclideanMetric}
  \uses{def:dc-ch1-2-1}
  \leanok
  \dcref{ch1:2.4}
  On $\mathbb{R}^n$, identify $\frac{\partial}{\partial x_i}$ with the standard
  basis vector $e_i$ and set $\langle e_i,e_j\rangle=\delta_{ij}$. This is
  Euclidean space of dimension $n$ with its standard Riemannian metric.
\end{example}

\begin{example}[immersed manifolds]
  \label{ex:dc-ch1-2-5}
  \lean{Riemannian.DCInducedForm, Riemannian.DCInducedMetric}
  \uses{def:dc-ch1-2-1, ex:dc-ch1-2-4}
  \leanok
  \dcref{ch1:2.5}
  If $f:M^n\to N^{n+k}$ is an immersion into a Riemannian manifold, define
  $$
  \langle u,v\rangle_p=\langle df_p(u),df_p(v)\rangle_{f(p)}.
  $$
  Injectivity of $df_p$ gives positive definiteness, and differentiability follows
  in coordinates; this is the metric \emph{induced by $f$}, and $f$ is an
  \emph{isometric immersion}. If $h:M^{n+k}\to N^k$ has regular value $q$, then
  $h^{-1}(q)$ is an $n$-submanifold with the metric induced by inclusion. For
  $h:\mathbb{R}^n\to\mathbb{R}$, $h(x)=\sum_i x_i^2-1$, the regular level
  $h^{-1}(0)=S^{n-1}$ has its canonical metric.
\end{example}
\begin{proof}
  \leanok
  \uses{ex:dc-ch1-2-4}
  \sketch
  Pull back the target inner product along $df_p$. Bilinearity and symmetry are
  inherited. If $u\ne0$, injectivity gives $df_p(u)\ne0$, hence
  $\langle u,u\rangle_p>0$. In coordinates the resulting coefficients are obtained
  by composing the differentiable target coefficients with the coordinate
  differential of $f$, so they vary differentiably and satisfy
  \cref{def:dc-ch1-2-1}.
\end{proof}

\begin{example}[Lie groups]
  \label{ex:dc-ch1-2-6}
  \uses{prop:dc-ch0-5-4}
  \notready
  \dcref{ch1:2.6}
  \notready
  A \emph{Lie group} is a group $G$ with a differentiable structure for which
  $(x,y)\mapsto xy^{-1}$ is differentiable. Left and right translations
  $L_x(y)=xy$ and $R_x(y)=yx$ are diffeomorphisms. A metric is \emph{left invariant}
  if every $L_x$ is an isometry, \emph{right invariant} if every $R_x$ is an
  isometry, and \emph{bi-invariant} if both conditions hold.

  A vector field $X$ is left invariant when $dL_xX=X$ for every $x$; it is determined
  by $X_e$. The bracket of left-invariant fields is left invariant, so
  $[X_e,Y_e]=[X,Y]_e$ defines the Lie algebra $\mathcal{G}=T_eG$. Given an inner
  product on $\mathcal{G}$, the formula
  $$
  \langle u,v\rangle_x=
  \left\langle(dL_{x^{-1}})_x(u),(dL_{x^{-1}})_x(v)\right\rangle_e
  \qquad\text{(2)}
  $$
  defines a left-invariant metric; the right-invariant construction is analogous,
  and every compact Lie group has a bi-invariant metric. For a bi-invariant metric,
  $$
  \langle[U,X],V\rangle=-\langle U,[V,X]\rangle
  \qquad\text{for }U,V,X\in\mathcal{G}.\tag{3}
  $$
  Conversely, a positive bilinear form on $\mathcal{G}$ satisfying (3), extended
  by (2), gives a bi-invariant metric.
\end{example}
\begin{proof}
  \sketch
  For left-invariant $X,Y$,
  $$
  dL_x[X,Y]f=[X,Y](f\circ L_x)=X(dL_xY)f-Y(dL_xX)f=[X,Y]f,
  $$
  so $[X,Y]$ is left invariant. Consequently $[X_e,Y_e]=[X,Y]_e$ defines the Lie
  algebra $\mathcal{G}=T_eG$. Formula (2) is left invariant by construction.
  For $a\in G$, the inner automorphism $R_{a^{-1}}L_a$ fixes $e$ and has
  differential $\operatorname{Ad}(a)=dR_{a^{-1}}dL_a$. If $x_t$ is the flow of a
  left-invariant $X$, then $L_y\circ x_t=x_t\circ L_y$, hence
  $x_t(y)=R_{x_t(e)}(y)$ and $dx_t=dR_{x_t(e)}$. Using
  $[Y,X]=\lim_{t\to0}t^{-1}(dx_t(Y)-Y)$ from
  \cref{prop:dc-ch0-5-4} and
  $\operatorname{Ad}(a)Y=dR_{a^{-1}}Y$ gives
  $$
  [Y,X]=\lim_{t\to0}\frac{1}{t}
  \left(\operatorname{Ad}((x_t(e))^{-1})Y-Y\right).
  $$
  Bi-invariance gives
  $\langle U,V\rangle=\langle dR_{x_t(e)}U,dR_{x_t(e)}V\rangle$; differentiating at
  $t=0$ yields (3).
\end{proof}

\begin{proposition}[existence of a left-invariant metric]
  \label{prop:dc-ch1-2-6-leftinv}
  \lean{Riemannian.DCLeftInvariantForm, Riemannian.DCLeftInvariantMetric}
  \uses{def:dc-ch1-2-1}
  \leanok
  Let $G$ be a Lie group with Lie algebra $\mathcal{G}=T_eG$, and let
  $\langle\ ,\ \rangle_e$ be a positive-definite symmetric inner product on
  $\mathcal{G}$. Then (2) defines a left-invariant Riemannian metric on $G$.
\end{proposition}
\begin{proof}
  \leanok
  \uses{def:dc-ch1-2-1}
  \sketch
  Each $(dL_{x^{-1}})_x$ is a linear isomorphism, since $L_x$ and $L_{x^{-1}}$
  are smooth inverses. Pulling back the fixed inner product therefore gives a
  symmetric, bilinear, positive-definite form. Smoothness follows from the smooth
  map $(x,y)\mapsto x^{-1}y$ and its parametrized differential. Formula (2) is
  left invariant by construction.
\end{proof}

\begin{example}[the product metric]
  \label{ex:dc-ch1-2-7}
  \lean{Riemannian.DCProductForm, Riemannian.DCProductMetric}
  \uses{def:dc-ch1-2-1, ex:dc-ch1-2-5}
  \leanok
  \dcref{ch1:2.7}
  For Riemannian manifolds $M_1,M_2$ with projections $\pi_1,\pi_2$, define
  the \emph{product metric} on $M_1\times M_2$ by
  $$
  \langle u,v\rangle_{(p,q)}=
  \langle d\pi_1\cdot u,d\pi_1\cdot v\rangle_p+
  \langle d\pi_2\cdot u,d\pi_2\cdot v\rangle_q.
  $$
  This holds for $(p,q)\in M_1\times M_2$ and
  $u,v\in T_{(p,q)}(M_1\times M_2)$.
  The product of the induced metrics on $S^1\times\cdots\times S^1=T^n$ is the
  \emph{flat torus} metric.
\end{example}
\begin{proof}
  \leanok
  \uses{ex:dc-ch1-2-5}
  \sketch
  Each summand is the pullback of a factor metric along a projection and is
  symmetric, bilinear, and differentiable by \cref{ex:dc-ch1-2-5}; their sum has
  the same properties. If $u\ne0$, its two projected components are not both zero,
  so positive definiteness of one factor makes the sum positive. Thus
  \cref{def:dc-ch1-2-1} applies.
\end{proof}

\begin{definition}[parametrized curve]
  \label{def:dc-ch1-2-8}
  \lean{Riemannian.DCIsCurve}
  \dcref{ch1:2.8}
  A differentiable map $c:I\to M$ from an open interval $I\subset\mathbb{R}$ to
  a differentiable manifold is a \emph{(parametrized) curve}; it may have
  self-intersections or corners.
\end{definition}

\begin{definition}[vector field along a curve; length]
  \label{def:dc-ch1-2-9}
  \lean{Riemannian.DCIsVectorFieldAlong, Riemannian.DCVelocity, Riemannian.DCArcLength}
  \uses{def:dc-ch1-2-1, def:dc-ch1-2-8}
  \leanok
  \dcref{ch1:2.9}
  A \emph{vector field along a curve} $c:I\to M$ is a differentiable assignment
  $V(t)\in T_{c(t)}M$, where differentiability means that $t\mapsto V(t)f$ is
  differentiable for every differentiable $f$ on $M$. The velocity field is
  $dc(\frac{d}{dt})=\frac{dc}{dt}$. For $[a,b]\subset I$, the restriction is a
  \emph{segment} and its length is
  $$
  \ell_a^b(c)=\int_a^b\left\langle\frac{dc}{dt},\frac{dc}{dt}\right\rangle^{1/2}\,dt.
  $$
\end{definition}

\begin{lemma}[isometries preserve length]
  \label{lem:dc-ch1-2-9-isom}
  \lean{Riemannian.DCIsometry.dcArcLength, Riemannian.DCPreservesMetric.dcArcLength}
  \uses{def:dc-ch1-2-2, def:dc-ch1-2-9}
  \leanok
  If $f:M\to N$ is an isometry and $c:I\to M$ is differentiable, then
  $$
  \ell_a^b(f\circ c)=\ell_a^b(c)
  $$
  for every segment $[a,b]\subset I$.
\end{lemma}
\begin{proof}
  \leanok
  \sketch
  The chain rule gives $\frac{d(f\circ c)}{dt}=df_{c(t)}(\frac{dc}{dt})$. Applying
  (1) to the velocity vector makes the two length integrands equal pointwise, so
  their integrals agree.
\end{proof}

\begin{proposition}[existence of Riemannian metrics]
  \label{prop:dc-ch1-2-10}
  \lean{Riemannian.exists_riemannianMetric}
  \uses{def:dc-ch1-2-1}
  \leanok
  \dcref{ch1:2.10}
  Every differentiable Hausdorff manifold with a countable basis admits a
  Riemannian metric.
\end{proposition}
\begin{proof}
  \leanok
  \uses{thm:dc-ch0-5-6,ex:dc-ch1-2-4}
  \sketch
  By \cref{thm:dc-ch0-5-6}, choose a differentiable partition of unity
  $\{f_\alpha\}$ subordinate to a locally finite cover by coordinate neighborhoods
  $V_\alpha$, with $f_\alpha\ge0$, $\mathrm{supp}\,f_\alpha\subset V_\alpha$,
  and $\sum_\alpha f_\alpha=1$. In a chart
  $\mathbf{x}_\alpha:V_\alpha\to\mathbb{R}^n$, pull back the Euclidean metric to
  obtain
  $$
  \langle u,v\rangle_p^\alpha=
  \left\langle d(\mathbf{x}_\alpha)_p u,d(\mathbf{x}_\alpha)_p v\right\rangle,
  \qquad p\in V_\alpha.
  $$
  Each local form is symmetric, bilinear, positive definite, and differentiable.
  Define
  $$
  \langle u,v\rangle_p=\sum_\alpha f_\alpha(p)\langle u,v\rangle_p^\alpha,
  $$
  omitting terms with $f_\alpha(p)=0$. Local finiteness gives differentiability and
  a finite sum near each point. Nonnegativity of every summand and
  $\sum_\alpha f_\alpha(p)=1$ imply positivity for $u\ne0$ from any index with
  $f_\alpha(p)>0$. Thus this is a Riemannian metric.
\end{proof}

\subsection{Volume on an oriented manifold}

For a positive parametrization $\mathbf{x}:U\to M$ and a positive orthonormal
basis of $T_pM$, write
$X_i(p)=\frac{\partial}{\partial x_i}(p)=\sum_j a_{ij}e_j$. Then
$g_{ik}(p)=\sum_j a_{ij}a_{kj}$ and
$$
\mathrm{vol}(X_1(p),\ldots,X_n(p))=\det(a_{ij})=\sqrt{\det(g_{ij})}(p).
$$
For another positive parametrization $\mathbf{y}$ with coefficients $h_{ij}$,
$$
\sqrt{\det(g_{ij})}(p)=J\sqrt{\det(h_{ij})}(p),\qquad
J=\det\left(\frac{\partial y_i}{\partial x_j}\right)>0.\qquad\text{(4)}
$$
If $R\subset M$ is open, connected, has compact closure, lies in a positively
parametrized neighborhood $\mathbf{x}(U)$, and the boundary of
$\mathbf{x}^{-1}(R)$ has measure zero, define
$$
\mathrm{vol}(R)=\int_{\mathbf{x}^{-1}(R)}\sqrt{\det(g_{ij})}\,dx_1\cdots dx_n.\qquad\text{(5)}
$$
The change-of-variables theorem and (4) make this independent of the chosen
positive chart; orientability fixes the sign.

\begin{remark}[volume form]
  \label{rem:dc-ch1-2-11}
  \uses{lem:dc-ch1-2-11-pointwise}
  \dcref{ch1:2.11}
  Equation (4) identifies the integrand in (5) with a positive differential
  $n$-form, the \emph{volume form} (or volume element) $\nu$. For a compact region
  not contained in one chart, choose a partition of unity $\{\varphi_i\}$ subordinate
  to a finite chart cover and set
  $$
  \mathrm{vol}(R)=\sum_i\int_{\mathbf{x}_i^{-1}(R)}\varphi_i\nu;
  $$
  the result is independent of the partition.
\end{remark}

\begin{remark}[volume from a volume element]
  \label{rem:dc-ch1-2-12}
  \dcref{ch1:2.12}
  A globally defined positive differential $n$-form gives a notion of volume on a
  differentiable manifold; a Riemannian metric is only one way to obtain such a
  volume element.
\end{remark}

\chapter{Affine Connections; Riemannian Connections}
\label{chap:dc-affine-connections-riemannian-connections}

Following the conventions of \cite{doCarmo}, $\mathcal{X}(M)$ denotes the space
of smooth vector fields on $M$, and $\mathcal{D}(M)$ denotes the ring of smooth
real-valued functions on $M$.

\section{Affine connections}

\begin{definition}[affine connection]
  \label{def:dc-ch2-2-1}
  \lean{Riemannian.IsAffineConnectionMap, Riemannian.AffineConnection.ofIsAffineConnectionMap, Riemannian.AffineConnection.isAffineConnectionMap}
  \leanok
  \dcref{ch2:2.1}
  An \emph{affine connection} on a differentiable manifold $M$ is a map
  \[
    \nabla\colon \mathcal{X}(M)\times\mathcal{X}(M)\longrightarrow\mathcal{X}(M),
    \qquad (X,Y)\longmapsto\nabla_XY,
  \]
  such that, for $X,Y,Z\in\mathcal{X}(M)$ and $f,g\in\mathcal{D}(M)$,
  \begin{enumerate}
    \item[(i)] $\nabla_{fX+gY}Z=f\nabla_XZ+g\nabla_YZ$;
    \item[(ii)] $\nabla_X(Y+Z)=\nabla_XY+\nabla_XZ$;
    \item[(iii)] $\nabla_X(fY)=f\nabla_XY+X(f)Y$.
  \end{enumerate}
\end{definition}

\begin{proposition}[covariant derivative along a curve]
  \label{prop:dc-ch2-2-2}
  \lean{Riemannian.AffineConnection.covDerivAlong,
    Riemannian.AffineConnection.covDerivAlong_eq_chart,
    Riemannian.AffineConnection.covDerivAlong_add,
    Riemannian.AffineConnection.covDerivAlong_smul,
    Riemannian.AffineConnection.covDerivAlong_comp_smoothVectorField,
    Riemannian.AffineConnection.covDerivAlong_unique,
    Riemannian.AffineConnection.covDerivAlong_unique_global,
    Riemannian.RiemannianMetric.leviCivita_covDerivAlong_eq}
  \uses{def:dc-ch2-2-1, lem:dc-ch2-2-2-coord, lem:dc-ch2-2-2-c-coord, lem:dc-ch2-2-2-c-bridge}
  \leanok
  \dcref{ch2:2.2}
  Let $M$ carry an affine connection $\nabla$, and let $c\colon I\to M$ be a
  differentiable curve. There is a unique operation assigning to every vector
  field $V$ along $c$ a vector field $DV/dt$ along $c$ such that
  \begin{enumerate}
    \item[(a)] $\frac{D}{dt}(V+W)=\frac{DV}{dt}+\frac{DW}{dt}$;
    \item[(b)] $\frac{D}{dt}(fV)=\frac{df}{dt}V+f\frac{DV}{dt}$ for every
      differentiable $f\colon I\to\mathbb{R}$;
    \item[(c)] if $V(t)=Y(c(t))$ for $Y\in\mathcal{X}(M)$, then
      $\frac{DV}{dt}=\nabla_{dc/dt}Y$.
  \end{enumerate}
  This operation is the \emph{covariant derivative along $c$}.
\end{proposition}
\begin{proof}
  \leanok
  \sketch
  In coordinates, write $V=\sum_jv^jX_j$, where
  $X_j=\frac{\partial}{\partial x_j}(c(t))$. Properties (a)--(c) and
  \cref{def:dc-ch2-2-1} force
  \[
    \frac{DV}{dt}
    =\sum_j\frac{dv^j}{dt}X_j
      +\sum_{i,j}\frac{dx_i}{dt}v^j\nabla_{X_i}X_j.\tag{1}
  \]
  Thus the operation is unique. Conversely, formula (1) defines it locally and
  gives (a)--(c); uniqueness makes the local definitions agree on chart
  overlaps.
\end{proof}

\begin{lemma}[covariant derivative in coordinates; properties (a), (b)]
  \label{lem:dc-ch2-2-2-coord}
  \lean{Riemannian.covariantDerivCoord,
    Riemannian.covariantDerivCoord_add,
    Riemannian.covariantDerivCoord_smul,
    Riemannian.covariantDerivCoord_unique}
  \uses{def:dc-ch2-3-7-christoffel}
  \leanok
  \dcref{ch2:2.2}
  Fix a Riemannian metric and a chart with Christoffel symbols $\Gamma^k_{ij}$.
  Set
  \[
    \Gamma(v,w)(y)=\sum_k\left(\sum_{i,j}\Gamma^k_{ij}(y)v^iw^j\right)e_k.
  \]
  For a coordinate curve $u\colon I\to\mathbb{R}^n$ and a coordinate vector
  field $V$ along it, formula (1) becomes
  \[
    \frac{DV}{dt}=\dot V+\Gamma(\dot u,V)(u).
  \]
  This operator is additive and satisfies
  $\frac{D}{dt}(fV)=\dot f\,V+f\frac{DV}{dt}$.
\end{lemma}

\begin{lemma}[property (c) in coordinates: induced fields]
  \label{lem:dc-ch2-2-2-c-coord}
  \lean{Riemannian.chartCovariantDeriv, Riemannian.covariantDerivCoord_induced}
  \uses{lem:dc-ch2-2-2-coord, def:dc-ch2-3-7-christoffel}
  \leanok
  \dcref{ch2:2.2}
  For a coordinate vector field $Y$ and a tangent vector $v$ at $y$, define
  \[
    (\nabla_vY)(y)=DY(y)v+\Gamma(v,Y(y))(y).
  \]
  If $V=Y\circ u$ along a coordinate curve $u$, then
  \[
    \frac{D}{dt}(Y\circ u)=(\nabla_{\dot u}Y)(u).
  \]
\end{lemma}
\begin{proof}
  \leanok
  \sketch
  The chain rule gives $(Y\circ u)'=DY(u)\dot u$, and the Christoffel terms in
  the two expressions coincide.
\end{proof}

\begin{lemma}[formula (10): the Koszul reduction on a coordinate frame]
  \label{lem:dc-ch2-2-2-c-formula10}
  \lean{Riemannian.christoffel_bridge_inner, Riemannian.christoffel_bridge_vector}
  \uses{lem:dc-ch2-3-6-koszul, def:dc-ch2-3-7-christoffel, thm:dc-ch2-3-6}
  \leanok
  \dcref{ch2:2.2}
  Let $\nabla$ be Levi-Civita and let $X_1,\dots,X_n$ be a coordinate frame at
  $p$. Suppose $[X_a,X_b](p)=0$ and
  $X_r\langle X_a,X_b\rangle(p)=\partial_rG_{ab}$, where
  $G_{ab}=\langle X_a,X_b\rangle$. Then
  \[
    \langle X_k,\nabla_{X_i}X_j\rangle
      =\sum_mG_{km}\Gamma^m_{ij},
    \qquad
    \nabla_{X_i}X_j=\sum_m\Gamma^m_{ij}X_m.\tag{10}
  \]
  \end{lemma}
\begin{proof}
    \leanok
    \sketch
    Apply the Koszul formula (9) from \cref{lem:dc-ch2-3-6-koszul} with
    $Y=X_i$, $X=X_j$, and $Z=X_k$. The bracket terms vanish, while the derivative
    terms give
    $\partial_iG_{kj}+\partial_jG_{ki}-\partial_kG_{ij}
      =2\sum_mG_{km}\Gamma^m_{ij}$.
    This proves the inner-product identity. Since the $X_k(p)$ form a basis and
    the metric is nondegenerate, it also proves the vector identity.
  \end{proof}

\begin{lemma}[property (c): chart-Christoffel bridge to the abstract connection]
  \label{lem:dc-ch2-2-2-c-bridge}
  \lean{Riemannian.exists_chartFrame_leviCivita_christoffel}
  \uses{lem:dc-ch2-2-2-c-coord, lem:dc-ch2-2-2-c-formula10, thm:dc-ch2-3-6, def:dc-ch2-3-7-christoffel}
  \leanok
  \dcref{ch2:2.2}
  The coordinate derivative in \cref{lem:dc-ch2-2-2-c-coord}, formed using the
  metric Christoffel symbols, is the coordinate expression of the abstract
  Levi-Civita connection from \cref{thm:dc-ch2-3-6}. At every point of a chart
  there are global smooth fields $X_1,\dots,X_n$ agreeing locally with its
  coordinate frame and satisfying
  \[
    \nabla_{X_i}X_j=\sum_m\Gamma^m_{ij}X_m.
  \]
  \end{lemma}
\begin{proof}
    \leanok
    \sketch
    Extend the coordinate frame fields globally without changing their germs at
    the chosen point. Locally they are pullbacks of constant coordinate fields,
    so their brackets vanish, and the directional derivatives of their Gram
    entries are the corresponding partial derivatives. Formula (10) from
    \cref{lem:dc-ch2-2-2-c-formula10} then identifies the coordinate and abstract
    Christoffel symbols, establishing property (c) of
    \cref{prop:dc-ch2-2-2}.
  \end{proof}

\begin{remark}[Remark 2.3]
  \label{rem:dc-ch2-2-3}
  \lean{Riemannian.AffineConnection.cov_congr_apply_left, Riemannian.AffineConnection.cov_congr_apply_eventuallyEq_right, Riemannian.AffineConnection.cov_smulSum_apply}
  \uses{rem:dc-ch0-5-7, def:dc-ch2-2-1}
  \leanok
  \dcref{ch2:2.3}
  An affine connection is local. In coordinates, if
  $X=\sum_ix_iX_i$, $Y=\sum_jy_jX_j$, and
  $\nabla_{X_i}X_j=\sum_k\Gamma_{ij}^kX_k$, then
  \[
    \nabla_XY
      =\sum_{i,j}x_iy_j\nabla_{X_i}X_j+\sum_jX(y_j)X_j
      =\sum_k\left(\sum_{i,j}x_iy_j\Gamma_{ij}^k+X(y_k)\right)X_k.
  \]
  Hence $\nabla_XY(p)$ depends only on $X(p)$ and the germ of $Y$ at $p$.
\end{remark}

\begin{remark}[Remark 2.4]
  \label{rem:dc-ch2-2-4}
  \lean{Riemannian.AffineConnection.covDerivAlong,
    Riemannian.Geodesic.continuousAt_and_hasGeodesicEquationAt_iff_leviCivita_covDerivAlong_velocity_eq_zero,
    Riemannian.covariantDerivCoord,
    Riemannian.covariantDerivCoord_add,
    Riemannian.covariantDerivCoord_smul}
  \uses{prop:dc-ch2-2-2}
  \leanok
  \dcref{ch2:2.4}
  An affine connection supplies a derivative of vector fields along curves, and
  therefore defines the acceleration $D(dc/dt)/dt$ of a curve in $M$.
\end{remark}

\begin{definition}[parallel vector field]
  \label{def:dc-ch2-2-5}
  \lean{Riemannian.IsParallelAlongWithinOn,
    Riemannian.IsIntrinsicPiecewiseParallelAlongOn,
    Riemannian.IsParallelCoord,
    Riemannian.IsParallelAlong,
    Riemannian.IsParallelAlongOn,
    Riemannian.isParallelAlongWithinOn_iff_isParallelFieldAlongOn,
    Riemannian.RiemannianMetric.leviCivita_isParallelFieldAlongOn_iff_isParallelAlong}
  \uses{lem:dc-ch2-2-2-coord}
  \leanok
  \dcref{ch2:2.5}
  A vector field $V$ along $c\colon I\to M$ is \emph{parallel} if
  $DV/dt=0$ on $I$. In coordinates this is the linear system
  \[
    \dot V=-\Gamma(\dot u,V)(u).
  \]
\end{definition}

\begin{lemma}[the parallel-transport coefficient is a bounded linear operator]
  \label{lem:dc-ch2-2-6-coeff}
  \lean{Riemannian.chartChristoffelContractionRight, Riemannian.lipschitzWith_neg_chartChristoffelContraction}
  \uses{def:dc-ch2-2-5}
  \leanok
  \dcref{ch2:2.6}
  For a chart velocity $v$ and base point $y$, the map
  $w\mapsto\Gamma(v,w)(y)$ is a bounded linear operator $A(v,y)$ on the
  finite-dimensional model space. Thus
  $w\mapsto-\Gamma(\dot u(t),w)(u(t))$ is globally Lipschitz with constant
  $\lVert A(\dot u(t),u(t))\rVert$.
\end{lemma}

\begin{proposition}[parallel transport]
  \label{prop:dc-ch2-2-6}
  \lean{Riemannian.exists_unique_isParallelAlongWithinOn,
    Riemannian.exists_unique_isIntrinsicPiecewiseParallelAlongOn,
    Riemannian.intrinsicPiecewiseParallelTransportTangentEquiv,
    Riemannian.intrinsicPiecewiseParallelTransportTangentEquiv_eq_terminal,
    Riemannian.metricInner_intrinsicPiecewiseParallelTransportTangentEquiv}
  \uses{def:dc-ch2-2-5, lem:dc-ch2-2-6-exist, lem:dc-ch2-2-6-unique, lem:dc-ch2-2-6-transport}
  \leanok
  \dcref{ch2:2.6}
  Let $c\colon I\to M$ be differentiable, $t_0\in I$, and
  $V_0\in T_{c(t_0)}M$. There is a unique parallel vector field $V$ along $c$
  with $V(t_0)=V_0$. Its values are the \emph{parallel transports} of $V_0$
  along $c$.
\end{proposition}
\begin{proof}
  \uses{lem:dc-ch2-2-6-exist, lem:dc-ch2-2-6-unique, lem:dc-ch2-2-6-transport}
  \sketch
  In one chart, writing $V=\sum_jv^jX_j$ reduces parallelism to
  \[
    0=\frac{dv^k}{dt}
      +\sum_{i,j}\Gamma_{ij}^kv^j\frac{dx_i}{dt},
      \qquad k=1,\dots,n.\tag{2}
  \]
  This linear system has a unique solution for each initial value. On a compact
  subinterval, cover the curve by finitely many coordinate neighborhoods and
  glue the local solutions using uniqueness. These solutions extend throughout
  $I$.
\end{proof}

\begin{lemma}[existence of parallel transport]
  \label{lem:dc-ch2-2-6-exist}
  \lean{Riemannian.exists_isParallelCoord_Icc,
    Riemannian.exists_isParallelAlongWithinOn,
    Riemannian.exists_isParallelAlongPresentation}
  \uses{def:dc-ch2-2-5, lem:dc-ch2-2-6-coeff}
  \leanok
  \dcref{ch2:2.6}
  On a compact interval $[a,b]$, every initial coordinate vector $V_0$ extends
  to a solution of
  \[
    \dot V=-\Gamma(\dot u,V)(u),\qquad V(a)=V_0.
  \]
  It suffices that $u$ is $C^1$ and the Christoffel symbols are continuous.
\end{lemma}
\begin{proof}
  \sketch
  By \cref{lem:dc-ch2-2-6-coeff}, this is a linear ODE
  $\dot V=A(t)V$ with continuous bounded coefficients. Local existence and a
  finite subdivision of $[a,b]$ give a global solution.
\end{proof}

\begin{lemma}[uniqueness of parallel transport]
  \label{lem:dc-ch2-2-6-unique}
  \lean{Riemannian.isParallelCoord_eqOn_Icc,
    Riemannian.isParallelSol_eqOn_Icc,
    Riemannian.IsParallelAlongWithinOn.eqOn_of_eq_at,
    Riemannian.eqOn_of_two_intrinsicPresentations,
    Riemannian.IsIntrinsicPiecewiseParallelAlongOn.eqOn}
  \uses{def:dc-ch2-2-5, lem:dc-ch2-2-2-coord, lem:dc-ch2-2-6-coeff}
  \leanok
  \dcref{ch2:2.6}
  Two parallel coordinate vector fields along the same coordinate curve on
  $[a,b]$ that agree at $a$ agree throughout $[a,b]$.
\end{lemma}
\begin{proof}
  \sketch
  Their difference solves the homogeneous linear system (2) with zero initial
  value. The Lipschitz estimate from \cref{lem:dc-ch2-2-6-coeff} and
  Grönwall's inequality imply that it vanishes.
\end{proof}

\begin{lemma}[parallel transport is a linear isomorphism]
  \label{lem:dc-ch2-2-6-transport}
  \lean{Riemannian.parallelTransport,
    Riemannian.exists_parallelTransport_spec,
    Riemannian.intrinsicPiecewiseParallelTransportTangentEquiv,
    Riemannian.intrinsicPiecewiseParallelTransportTangentEquiv_eq_terminal,
    Riemannian.metricInner_intrinsicPiecewiseParallelTransportTangentEquiv}
  \uses{lem:dc-ch2-2-6-exist, lem:dc-ch2-2-6-unique}
  \leanok
  \dcref{ch2:2.6}
  On $[a,b]$, the endpoint map
  \[
    P_c\colon T_{c(a)}M\longrightarrow T_{c(b)}M,
    \qquad V_0\longmapsto V(b),
  \]
  where $V$ is the parallel field with $V(a)=V_0$, is a linear isomorphism.
  Its inverse is parallel transport along the reversed curve.
\end{lemma}
\begin{proof}
  \sketch
  Existence and uniqueness make the map well-defined. Superposition for the
  linear system (2) gives linearity, and backward uniqueness gives injectivity
  and the stated inverse.
\end{proof}

\section{Riemannian connections}

\begin{definition}[compatible with the metric]
  \label{def:dc-ch2-3-1}
  \lean{Riemannian.CompatibleGen,
    Riemannian.AffineConnection.PreservesMetricUnderParallelTransport,
    Riemannian.AffineConnection.PreservesMetricAlongOn}
  \leanok
  \dcref{ch2:3.1}
  An affine connection $\nabla$ on a Riemannian manifold is \emph{compatible
  with the metric} if the inner product of any two parallel vector fields along
  any smooth curve is constant.
\end{definition}

\begin{proposition}[product rule characterization]
  \label{prop:dc-ch2-3-2}
  \lean{Riemannian.compatibleGen_iff_productRuleGen,
    Riemannian.AffineConnection.HasMetricProductRuleAlong,
    Riemannian.AffineConnection.preservesMetricUnderParallelTransport_iff_hasMetricProductRuleAlong,
    Riemannian.hasDerivAt_metricInner_along_of_mdifferentiableAt,
    Riemannian.RiemannianMetric.leviCivitaConnection_hasMetricProductRuleAlong,
    Riemannian.metricInner_const_of_isParallelAlong,
    Riemannian.metricInner_eq_endpoints_of_isParallelAlongWithinOn}
  \uses{def:dc-ch2-3-1, prop:dc-ch2-2-6, lem:dc-ch2-3-2-coord}
  \leanok
  \dcref{ch2:3.2}
  A connection $\nabla$ is compatible with the metric if and only if, for all
  vector fields $V,W$ along every differentiable curve $c\colon I\to M$,
  \[
    \frac{d}{dt}\langle V,W\rangle
      =\left\langle\frac{DV}{dt},W\right\rangle
       +\left\langle V,\frac{DW}{dt}\right\rangle,
    \qquad t\in I.\tag{3}
  \]
\end{proposition}
\begin{proof}
  \leanok
  \sketch
  Formula (3) immediately makes inner products of parallel fields constant.
  Conversely, parallel-transport an orthonormal basis at one point along $c$;
  compatibility keeps it orthonormal. Expanding
  $V=\sum_iv^iP_i$ and $W=\sum_iw^iP_i$ in this parallel frame reduces (3) to
  the ordinary product rule for $\sum_iv^iw^i$.
\end{proof}

\begin{corollary}[compatibility on vector fields]
  \label{cor:dc-ch2-3-3}
  \lean{Riemannian.AffineConnection.IsMetricCompatible,
    Riemannian.AffineConnection.isMetricCompatible_iff_hasMetricProductRuleAlong,
    Riemannian.AffineConnection.preservesMetricUnderParallelTransport_iff_isMetricCompatible}
  \uses{prop:dc-ch2-3-2}
  \leanok
  \dcref{ch2:3.3}
  A connection $\nabla$ is compatible with the metric if and only if
  \[
    X\langle Y,Z\rangle
      =\langle\nabla_XY,Z\rangle+\langle Y,\nabla_XZ\rangle,
    \qquad X,Y,Z\in\mathcal{X}(M).\tag{4}
  \]
\end{corollary}
\begin{proof}
  \sketch
  Evaluate (3) at a point along a curve whose velocity there is $X$. This gives
  (4). Conversely, apply (4) pointwise to the velocity field of a curve to
  recover (3).
\end{proof}

\begin{lemma}[metric compatibility of $D/dt$, coordinate form]
  \label{lem:dc-ch2-3-2-coord}
  \lean{Riemannian.CompatibleGen,
    Riemannian.compatibleGen_iff_productRuleGen,
    Riemannian.hasDerivAt_chartMetricInner_along}
  \uses{lem:dc-ch2-2-2-coord, def:dc-ch2-3-7-christoffel}
  \leanok
  \dcref{ch2:3.2}
  In a chart, write
  $\langle a,b\rangle_y=\sum_{i,j}g_{ij}(y)a^ib^j$ and
  $DV/dt=\dot V+\Gamma(\dot u,V)(u)$. Then coordinate vector fields $V,W$
  along $u$ satisfy
  \[
    \frac{d}{dt}\langle V,W\rangle
      =\left\langle\frac{DV}{dt},W\right\rangle
       +\left\langle V,\frac{DW}{dt}\right\rangle.
  \]
\end{lemma}
\begin{proof}
  \leanok
  \uses{def:dc-ch2-3-7-christoffel}
  \sketch
  Differentiate $\sum_{i,j}g_{ij}(u)V^iW^j$ and use
  \[
    \partial_kg_{ij}
      =\sum_m\left(g_{mj}\Gamma^m_{ki}+g_{im}\Gamma^m_{kj}\right),
  \]
  which follows from \cref{def:dc-ch2-3-7-christoffel}. The two Christoffel
  sums combine respectively with $\dot V$ and $\dot W$ to give the two
  covariant derivatives.
\end{proof}

\begin{lemma}[parallel fields have constant inner product, coordinate form]
  \label{lem:dc-ch2-3-1-coord}
  \lean{Riemannian.CompatibleGen,
    Riemannian.hasDerivAt_chartMetricInner_eq_zero_of_isParallelCoord}
  \uses{def:dc-ch2-2-5, def:dc-ch2-3-1, lem:dc-ch2-3-2-coord}
  \leanok
  \dcref{ch2:3.1}
  If $V$ and $W$ are parallel coordinate vector fields along a coordinate
  curve, then $\langle V,W\rangle$ is constant.
\end{lemma}
\begin{proof}
  \leanok
  \uses{lem:dc-ch2-3-2-coord}
  \sketch
  In \cref{lem:dc-ch2-3-2-coord}, both covariant derivatives vanish, so
  $\frac{d}{dt}\langle V,W\rangle=0$.
\end{proof}

\begin{definition}[symmetric connection]
  \label{def:dc-ch2-3-4}
  \lean{Riemannian.AffineConnection.IsSymmetric}
  \uses{def:dc-ch2-2-1}
  \leanok
  \dcref{ch2:3.4}
  An affine connection $\nabla$ is \emph{symmetric} if
  \[
    \nabla_XY-\nabla_YX=[X,Y]
    \quad\text{for all }X,Y\in\mathcal{X}(M).\tag{5}
  \]
\end{definition}

\begin{remark}[Remark 3.5]
  \label{rem:dc-ch2-3-5}
  \lean{Riemannian.chartChristoffel_symm}
  \leanok
  \dcref{ch2:3.5}
  In a coordinate frame $X_i=\partial/\partial x_i$, symmetry is equivalent to
  \[
    \nabla_{X_i}X_j-\nabla_{X_j}X_i=[X_i,X_j]=0,
    \qquad
    \Gamma_{ij}^k=\Gamma_{ji}^k.\tag{5'}
  \]
\end{remark}

\begin{lemma}[Koszul formula]
  \label{lem:dc-ch2-3-6-koszul}
  \lean{Riemannian.AffineConnection.koszul_formula}
  \uses{def:dc-ch2-3-4, cor:dc-ch2-3-3}
  \leanok
  \dcref{ch2:3.6}
  If $\nabla$ is symmetric and metric-compatible, then, for all
  $X,Y,Z\in\mathcal{X}(M)$,
  \[
    \begin{aligned}
    2\langle Z,\nabla_YX\rangle
      ={}&X\langle Y,Z\rangle+Y\langle Z,X\rangle-Z\langle X,Y\rangle\\
       &-\langle[X,Z],Y\rangle-\langle[Y,Z],X\rangle
        -\langle[X,Y],Z\rangle.
    \end{aligned}\tag{9}
  \]
\end{lemma}
\begin{proof}
  \leanok
  \uses{def:dc-ch2-3-4, cor:dc-ch2-3-3}
  \sketch
  Apply (4) to the three cyclic orderings of $X,Y,Z$. Add the first two
  identities, subtract the third, and replace differences of covariant
  derivatives using symmetry (5). Symmetry of the metric gives (9).
\end{proof}

\begin{lemma}[tangent-vector extension]
  \label{lem:dc-ch2-tangent-extension}
  \lean{Riemannian.exists_smoothVectorField_eq}
  \leanok
  On a finite-dimensional, $\sigma$-compact, Hausdorff smooth manifold, every
  $v\in T_pM$ is the value $Z(p)$ of a global smooth vector field
  $Z\in\mathcal{X}(M)$.
\end{lemma}
\begin{proof}
  \leanok
  \sketch
  In a bundle trivialization near $p$, take the locally constant section with
  value $v$ at $p$; away from $p$, use the zero section. A partition of unity
  glues these local sections to the required global field.
\end{proof}

\begin{lemma}[uniqueness of the Levi-Civita connection]
  \label{lem:dc-ch2-3-6-unique}
  \lean{Riemannian.AffineConnection.leviCivita_unique'}
  \uses{lem:dc-ch2-3-6-koszul, lem:dc-ch2-tangent-extension}
  \leanok
  \dcref{ch2:3.6}
  On a finite-dimensional, $\sigma$-compact, Hausdorff Riemannian manifold,
  there is at most one affine connection that is symmetric and
  metric-compatible.
\end{lemma}
\begin{proof}
  \leanok
  \uses{lem:dc-ch2-3-6-koszul, lem:dc-ch2-tangent-extension}
  \sketch
  For two such connections, (9) gives equal inner products
  $\langle Z,(\nabla_1)_YX\rangle=\langle Z,(\nabla_2)_YX\rangle$ for every
  global field $Z$. By \cref{lem:dc-ch2-tangent-extension}, this tests against
  every tangent vector at each point; nondegeneracy therefore gives
  $\nabla_1=\nabla_2$.
\end{proof}

\begin{lemma}[well-definedness of the Koszul functional]
  \label{lem:dc-ch2-3-6-welldef}
  \lean{Riemannian.RiemannianMetric.koszulRHS, Riemannian.RiemannianMetric.koszulRHS_add_right, Riemannian.RiemannianMetric.koszulRHS_smul_right}
  \uses{lem:dc-ch2-3-6-koszul}
  \leanok
  \dcref{ch2:3.6}
  For fixed $X,Y\in\mathcal{X}(M)$, let
  \[
    F_{Y,X}(Z)
      =X\langle Y,Z\rangle+Y\langle Z,X\rangle-Z\langle X,Y\rangle
       -\langle[X,Z],Y\rangle-\langle[Y,Z],X\rangle
       -\langle[X,Y],Z\rangle.
  \]
  Then $F_{Y,X}$ is $\mathcal{D}(M)$-linear in $Z$. Hence
  $Z\mapsto\frac12F_{Y,X}(Z)$ is tensorial and has a metric Riesz dual, the
  candidate $\nabla_YX$.
\end{lemma}
\begin{proof}
  \leanok
  \uses{lem:dc-ch2-3-6-koszul}
  \sketch
  Additivity is termwise. For $f\in\mathcal{D}(M)$, expand
  $X(fh)=fX(h)+X(f)h$ and
  $[X,fZ]=f[X,Z]+X(f)Z$. The derivative-of-$f$ terms cancel by symmetry of the
  metric, leaving $F_{Y,X}(fZ)=fF_{Y,X}(Z)$.
\end{proof}

\begin{lemma}[the Koszul-dual connection is an affine connection]
  \label{lem:dc-ch2-3-6-affine}
  \lean{Riemannian.RiemannianMetric.koszulRHS_add_left, Riemannian.RiemannianMetric.koszulRHS_leibniz_left, Riemannian.RiemannianMetric.koszulRHS_add_middle, Riemannian.RiemannianMetric.koszulRHS_smul_middle}
  \uses{def:dc-ch2-2-1, lem:dc-ch2-3-6-welldef}
  \leanok
  \dcref{ch2:3.6}
  The candidate $\nabla_YX=\frac12(F_{Y,X})^\sharp$ satisfies
  \begin{align*}
    F_{Y,X_1+X_2}&=F_{Y,X_1}+F_{Y,X_2},\\
    F_{Y,fX}&=fF_{Y,X}+2Y(f)\langle X,Z\rangle,\\
    F_{Y_1+Y_2,X}&=F_{Y_1,X}+F_{Y_2,X},\\
    F_{fY,X}&=fF_{Y,X}.
  \end{align*}
  Therefore it is additive and $\mathcal{D}(M)$-linear in the direction field,
  and additive with the Leibniz rule in the differentiated field; hence it is an
  affine connection.
\end{lemma}
\begin{proof}
  \leanok
  \uses{lem:dc-ch2-3-6-koszul}
  \sketch
  Substitute the indicated sums and scalar multiples into the six terms of
  $F$. Additivity is termwise. The Leibniz rules for derivatives and brackets
  produce the displayed $2Y(f)\langle X,Z\rangle$ in the differentiated slot;
  in the direction slot all derivative-of-$f$ terms cancel in pairs. Together
  with \cref{lem:dc-ch2-3-6-welldef}, these identities give the connection
  axioms.
\end{proof}

\begin{lemma}[the Koszul-dual connection is Levi-Civita]
  \label{lem:dc-ch2-3-6-exist}
  \lean{Riemannian.AffineConnection.isLeviCivita_of_koszulDual, Riemannian.koszulDual_isSymmetric, Riemannian.koszulDual_isMetricCompatible}
  \uses{def:dc-ch2-2-1, def:dc-ch2-3-4, lem:dc-ch2-3-6-koszul, lem:dc-ch2-3-6-welldef, lem:dc-ch2-tangent-extension}
  \leanok
  \dcref{ch2:3.6}
  On a finite-dimensional, $\sigma$-compact, Hausdorff Riemannian manifold,
  suppose an affine connection satisfies
  \[
    2\langle\nabla_YX,Z\rangle=F_{Y,X}(Z)
  \]
  for every $X,Y,Z$. Then it is symmetric and metric-compatible.
\end{lemma}
\begin{proof}
  \leanok
  \uses{lem:dc-ch2-3-6-koszul, lem:dc-ch2-tangent-extension}
  \sketch
  The identities
  \[
    F_{Y,X}(Z)-F_{X,Y}(Z)=2\langle[X,Y],Z\rangle
  \]
  and
  \[
    F_{Y,X}(Z)+F_{Z,X}(Y)=2X\langle Y,Z\rangle
  \]
  follow by cancellation. The first, together with
  \cref{lem:dc-ch2-tangent-extension} and nondegeneracy, gives
  $\nabla_XY-\nabla_YX=[X,Y]$. The second gives the product rule (4), hence
  metric compatibility.
\end{proof}

\begin{theorem}[Levi-Civita]
  \label{thm:dc-ch2-3-6}
  \lean{Riemannian.RiemannianMetric.exists_unique_isLeviCivita}
  \uses{def:dc-ch2-3-4, cor:dc-ch2-3-3, lem:dc-ch2-3-6-koszul, lem:dc-ch2-3-6-unique, lem:dc-ch2-3-6-welldef, lem:dc-ch2-3-6-exist}
  \leanok
  \dcref{ch2:3.6}
  Every Riemannian manifold admits a unique affine connection $\nabla$ that is
  \begin{enumerate}
    \item[(a)] symmetric;
    \item[(b)] compatible with the Riemannian metric.
  \end{enumerate}
\end{theorem}
\begin{proof}
  \leanok
  \uses{lem:dc-ch2-3-6-unique, lem:dc-ch2-3-6-exist}
  \sketch
  Metric compatibility applied cyclically gives
  \[
    X\langle Y,Z\rangle
      =\langle\nabla_XY,Z\rangle+\langle Y,\nabla_XZ\rangle,\tag{6}
  \]
  \[
    Y\langle Z,X\rangle
      =\langle\nabla_YZ,X\rangle+\langle Z,\nabla_YX\rangle,\tag{7}
  \]
  \[
    Z\langle X,Y\rangle
      =\langle\nabla_ZX,Y\rangle+\langle X,\nabla_ZY\rangle.\tag{8}
  \]
  Adding (6) and (7), subtracting (8), and using symmetry gives
  \[
    \begin{aligned}
    \langle Z,\nabla_YX\rangle=\tfrac12\{&X\langle Y,Z\rangle
      +Y\langle Z,X\rangle-Z\langle X,Y\rangle\\
      &-\langle[X,Z],Y\rangle-\langle[Y,Z],X\rangle
       -\langle[X,Y],Z\rangle\}.\tag{9}
    \end{aligned}
  \]
  Thus uniqueness follows from \cref{lem:dc-ch2-3-6-unique}. For existence,
  take the metric Riesz dual of the right-hand side of (9); its
  well-definedness and Levi-Civita properties are
  \cref{lem:dc-ch2-3-6-welldef,lem:dc-ch2-3-6-exist}.
\end{proof}

\begin{definition}[Christoffel symbols]
  \label{def:dc-ch2-3-7-christoffel}
  \lean{Riemannian.chartChristoffel}
  \uses{def:dc-ch1-2-1}
  \leanok
  \dcref{ch2:3.7}
  In coordinates, let $g_{ij}=\langle X_i,X_j\rangle$ and let $(g^{km})$ be
  the inverse matrix. The \emph{Christoffel symbols of the second kind} are
  \[
    \Gamma_{ij}^m
      =\tfrac12\sum_k\left(
        \frac{\partial g_{jk}}{\partial x_i}
       +\frac{\partial g_{ki}}{\partial x_j}
       -\frac{\partial g_{ij}}{\partial x_k}\right)g^{km}.\tag{10}
  \]
\end{definition}

\begin{remark}[the Levi-Civita connection and its coefficients]
  \label{rem:dc-ch2-3-7}
  \lean{Riemannian.chartChristoffel, Riemannian.chartChristoffel_symm, Riemannian.chartGram_christoffel_contraction, Riemannian.christoffel_bridge_inner, Riemannian.christoffel_bridge_vector, Riemannian.exists_chartFrame_leviCivita_christoffel}
  \uses{thm:dc-ch2-3-6, def:dc-ch2-3-7-christoffel}
  \leanok
  \dcref{ch2:3.7}
  The connection of \cref{thm:dc-ch2-3-6} is the \emph{Levi-Civita} or
  \emph{Riemannian connection}. Its coordinate coefficients, defined by
  $\nabla_{X_i}X_j=\sum_k\Gamma_{ij}^kX_k$, satisfy
  \[
    \sum_\ell\Gamma_{ij}^\ell g_{\ell k}
      =\tfrac12\left(
        \frac{\partial g_{jk}}{\partial x_i}
       +\frac{\partial g_{ki}}{\partial x_j}
       -\frac{\partial g_{ij}}{\partial x_k}\right),
  \]
  and hence agree with \cref{def:dc-ch2-3-7-christoffel}. Moreover,
  \[
    \frac{DV}{dt}
      =\sum_k\left(\frac{dv^k}{dt}
        +\sum_{i,j}\Gamma_{ij}^kv^j\frac{dx_i}{dt}\right)X_k.
  \]
  In Euclidean coordinates, $\Gamma_{ij}^k=0$, so this is the ordinary
  derivative.
\end{remark}

\begin{proposition}[Riemannian parallel transport along a $C^1$ curve]
  \label{prop:dc-ch2-2-6-riemannian}
  \dcref{ch2:2.6}
  \uses{def:dc-ch2-2-5}
  \lean{Riemannian.Variation.exists_parallelCovariantFieldAlongOn_at}
  \leanok
  Let $M$ be a positive-dimensional boundaryless Riemannian manifold, let
  $\gamma:\mathbb{R}\to M$ be a $C^1$ curve, and let $a<b$. For every
  $t_0\in[a,b]$ and $V_0\in T_{\gamma(t_0)}M$, there exists a parallel vector
  field $V$ along $\gamma$ on $[a,b]$ such that $V(t_0)=V_0$. It is unique on
  $[a,b]$ among parallel fields having that value at $t_0$.
\end{proposition}
\begin{proof}
  \uses{lem:dc-ch2-2-6-exist, lem:dc-ch2-2-6-unique, lem:dc-ch2-2-6-transport}
  \leanok
  Starting at $a$, solve the parallelism system in a coordinate neighborhood.
  Compactness of $[a,b]$ gives finitely many such neighborhoods along the curve,
  and uniqueness makes the local solutions agree on every overlap; hence each
  $X\in T_{\gamma(a)}M$ determines a unique parallel field $V_X$ on $[a,b]$.
  Evaluation at $t_0$ defines a linear map
  $P_{a,t_0}:T_{\gamma(a)}M\to T_{\gamma(t_0)}M$. If
  $P_{a,t_0}(X)=P_{a,t_0}(Y)$, uniqueness applied from $t_0$ gives
  $V_X=V_Y$ on $[a,b]$, and therefore $X=Y$. Thus $P_{a,t_0}$ is injective;
  since its domain and codomain have the same finite dimension, it is an
  isomorphism. Set $X=P_{a,t_0}^{-1}(V_0)$ and $V=V_X$. Then
  $V(t_0)=V_0$, and a final application of uniqueness at $t_0$ proves the
  asserted uniqueness of $V$ on $[a,b]$.
\end{proof}

\chapter{Geodesics; Convex Neighborhoods}
\label{chap:dc-geodesics-convex-neighborhoods}

\section{The geodesic flow}

Throughout this chapter, $M$ is a Riemannian manifold with its Riemannian
connection.

\begin{definition}[geodesic]
  \label{def:dc-ch3-2-1}
  \lean{Riemannian.Geodesic.HasGeodesicEquationAt,
    Riemannian.Geodesic.hasGeodesicEquationAt_iff_covDerivAlong_velocity_eq_zero,
    Riemannian.Geodesic.continuousAt_and_hasGeodesicEquationAt_iff_leviCivita_covDerivAlong_velocity_eq_zero,
    Riemannian.Geodesic.IsGeodesicCurve,
    Riemannian.Geodesic.isGeodesicCurve_iff_covDerivAlong_velocity_eq_zero,
    Riemannian.Geodesic.isGeodesicCurve_iff_leviCivita_covDerivAlong_velocity_eq_zero,
    Riemannian.RiemannianMetric.leviCivita_covDerivAlong_eq}
  \leanok
  \dcref{ch3:2.1}
  A parametrized curve $\gamma:I\to M$ is a \emph{geodesic at} $t_0\in I$ if
  $\frac{D}{dt}\left(\frac{d\gamma}{dt}\right)=0$ at the point $t_0$; if $\gamma$ is
  a geodesic at $t$ for all $t\in I$, then $\gamma$ is a \emph{geodesic}. If
  $[a,b]\subset I$, the restriction of $\gamma$ to $[a,b]$ is a \emph{geodesic segment joining $\gamma(a)$ to $\gamma(b)$}. By abuse of language, the image $\gamma(I)$
  is also called a geodesic.
  If $\gamma$ is a geodesic then
  $$
  \frac{d}{dt}\Big\langle\frac{d\gamma}{dt},\frac{d\gamma}{dt}\Big\rangle=2\Big\langle\frac{D}{dt}\frac{d\gamma}{dt},\frac{d\gamma}{dt}\Big\rangle=0,
  $$
  so $\left|\frac{d\gamma}{dt}\right|=c\neq 0$ is constant (we exclude geodesics
  reducing to points). Arc length $s(t)=\int_{t_0}^t\left|\frac{d\gamma}{dt}\right|dt=c(t-t_0)$
  is then proportional to the parameter; when $c=1$ the geodesic is \emph{normalized}.
  In a coordinate system $(U,\mathbf{x})$ about $\gamma(t_0)$, a curve
  $\gamma(t)=(x_1(t),\dots,x_n(t))$ is a geodesic iff
  $$
  0=\frac{D}{dt}\Big(\frac{d\gamma}{dt}\Big)=\sum_k\Big(\frac{d^2x_k}{dt^2}+\sum_{i,j}\Gamma_{ij}^k\frac{dx_i}{dt}\frac{dx_j}{dt}\Big)\frac{\partial}{\partial x_k},
  $$
  i.e. iff the second-order system
  $$
  \frac{d^2x_k}{dt^2}+\sum_{i,j}\Gamma_{ij}^k\frac{dx_i}{dt}\frac{dx_j}{dt}=0,\qquad k=1,\dots,n,\qquad\text{(1)}
  $$
  holds.

\end{definition}

\begin{lemma}[geodesics have constant speed]
  \label{lem:dc-ch3-2-1-constspeed}
  \lean{Riemannian.hasDerivAt_chartMetricInner_geodesic_speed_zero, Riemannian.Geodesic.speedSq, Riemannian.Geodesic.IsGeodesicOn.hasDerivAt_speedSq_zero, Riemannian.Geodesic.IsGeodesicOn.speedSq_eq}
  \uses{def:dc-ch3-2-1, lem:dc-ch2-3-2-coord}
  \leanok
  \dcref{ch3:2.1}
  If $\gamma$ is a geodesic at $t$, then
  $\frac{d}{ds}\big\langle\gamma'(s),\gamma'(s)\big\rangle=0$ at $s=t$. Read in
  the chart at the foot $\gamma(t)$, with $u=\varphi_{\gamma(t)}\circ\gamma$ the
  chart representative and $\langle\cdot,\cdot\rangle$ the chart Gram inner
  product $g_{ij}(u)\,\cdot^i\,\cdot^j$, the squared speed
  $s\mapsto\langle u'(s),u'(s)\rangle$ has vanishing derivative at $s=t$.

  Intrinsically: the squared speed
  $t\mapsto\langle\gamma'(t),\gamma'(t)\rangle$ of the manifold velocity
  $\gamma'(t)\in T_{\gamma(t)}M$ has vanishing derivative at every time of an
  open set on which the (continuous) curve $\gamma$ is a geodesic, and is
  therefore constant on every such connected open set. Hence $|\gamma'|=c$ is
  constant along a geodesic (excluding geodesics reducing to points), and the
  arc length $s(t)=\int_{t_0}^t|\gamma'|=c(t-t_0)$ is proportional to the
  parameter.
\end{lemma}
\begin{proof}
  \sketch
  \leanok
  By metric compatibility in coordinate form
  (\cref{lem:dc-ch2-3-2-coord}),
  $$
  \frac{d}{ds}\big\langle u'(s),u'(s)\big\rangle
    =\Big\langle\tfrac{Du'}{ds},u'\Big\rangle+\Big\langle u',\tfrac{Du'}{ds}\Big\rangle,
  $$
  where $\frac{Du'}{ds}=u''+\Gamma(u',u')(u)$ is the covariant derivative of the
  velocity. The geodesic equation at $t$ is exactly $\frac{Du'}{ds}=0$ at $s=t$,
  so both inner products vanish and the derivative is $0$.

  For the intrinsic statement, fix a base time $t$ and read everything in the
  chart at $\gamma(t)$: for $\sigma$ near $t$ the chart Gram inner product of
  the chart velocity at $u(\sigma)$ equals
  $\langle\gamma'(\sigma),\gamma'(\sigma)\rangle$ at $\gamma(\sigma)$, because
  the Gram entries are by definition the inner products of the chart frame
  vectors, and the chart velocity taken at base time $\sigma$ transforms into
  the one taken at base time $t$ by the derivative of the smooth chart
  transition — exactly the change of frame between the two chart
  presentations of $T_{\gamma(\sigma)}M$. Hence the intrinsic squared speed
  agrees near $t$ with the chart expression above, so its derivative vanishes
  at $t$; a real function on a connected open set whose derivative vanishes
  everywhere is constant.
\end{proof}

\subsection{Tangent bundle}
$TM$ is the set of pairs $(q,v)$, $q\in M$, $v\in T_qM$. If $(U,\mathbf{x})$ is a
coordinate system on $M$, any $v\in T_qM$, $q\in\mathbf{x}(U)$, is
$\sum_i y_i\frac{\partial}{\partial x_i}$, and $(x_1,\dots,x_n,y_1,\dots,y_n)$ are
coordinates of $(q,v)$ on $TU$, giving $TM$ a differentiable structure (cf.
\cref{ex:dc-ch0-4-1} of Chap. 0). Locally $TU=U\times\mathbb{R}^n$, and the canonical
projection $\pi:TM\to M$, $\pi(q,v)=q$, is differentiable. A curve $t\mapsto\gamma(t)$
in $M$ determines $t\mapsto(\gamma(t),\frac{d\gamma}{dt}(t))$ in $TM$; if $\gamma$
is a geodesic, this curve satisfies the first-order system

$$
\begin{cases}\dfrac{dx_k}{dt}=y_k,\\[4pt]\dfrac{dy_k}{dt}=-\sum_{i,j}\Gamma_{ij}^k y_i y_j,\end{cases}\qquad k=1,\dots,n,\qquad\text{(1')}
$$

on $TU$. Thus the second-order system (1) on $U$ is equivalent to the first-order
system (1') on $TU$.

\begin{lemma}[trajectories of a vector field: existence and local uniqueness]
  \label{lem:dc-ch3-2-2-pointwise}
  \lean{exists_isMIntegralCurveAt_of_contMDiffAt_boundaryless, isMIntegralCurveAt_eventuallyEq_of_contMDiffAt_boundaryless}
  \dcref{ch3:2.2}
  \mathlibok
  Let $X$ be a $C^1$ vector field on a boundaryless manifold $M$ modelled on a
  complete normed space, and $p\in M$. Then there is a trajectory of $X$ through
  $p$ at $t=0$ (a curve $\gamma$ with $\gamma'=X(\gamma)$ defined near $0$), and
  any two trajectories of $X$ that agree at a time $t_0$ agree on a
  neighbourhood of $t_0$. (Picard--Lindel\"of and Gr\"onwall, transported to the
  manifold through a chart.)
\end{lemma}

\begin{lemma}[local flow with Lipschitz dependence on the initial condition]
  \label{lem:dc-ch3-2-2-flow}
  \lean{IsPicardLindelof.exists_forall_mem_closedBall_eq_hasDerivWithinAt_lipschitzOnWith, IsPicardLindelof.exists_forall_mem_closedBall_eq_hasDerivWithinAt_continuousOn, IsPicardLindelof.of_contDiffAt_one}
  \uses{lem:dc-ch3-2-2-pointwise}
  \dcref{ch3:2.2}
  \mathlibok
  Let $f$ be a $C^1$ vector field on a Banach space. Around any point there are
  $r > 0$, $\varepsilon > 0$ and a \emph{local flow}: a single family
  $Z(x, \cdot)$ of solutions of $z' = f(z)$, one for each initial condition $x$
  in the closed ball of radius $r$, all defined on the fixed time interval
  $[-\varepsilon, \varepsilon]$, and Lipschitz in $x$ uniformly in time (hence
  jointly continuous in $(x, t)$). This is the uniform-time and
  continuous-dependence content of \cref{thm:dc-ch3-2-2}; beyond it, that
  theorem asserts the joint $C^\infty$ regularity of the flow, whose $C^1$-in-%
  the-initial-condition part is \cref{lem:dc-ch3-2-2-c1dep}.
\end{lemma}

\begin{lemma}[$C^1$ dependence on the initial condition along a base solution]
  \label{lem:dc-ch3-2-2-c1dep}
  \lean{Riemannian.FlowDependence.hasStrictFDerivAt_of_picardResidual_curve, Riemannian.FlowDependence.exists_hasStrictFDerivAt_of_picardResidual_curve, Riemannian.FlowDependence.hasStrictFDerivAt_eval_of_picardResidual_curve, Riemannian.FlowDependence.hasStrictFDerivAt_superposition}
  \uses{lem:dc-ch3-2-2-flow}
  \leanok
  \dcref{ch3:2.2}
  Let $f$ be a vector field on a Banach space $E$, differentiable on an open
  set $u$ with derivative $f'$, and let $\alpha_0 : [0,T] \to u$ be a solution
  of $\alpha' = f(\alpha)$ with initial value $x_0$, with $f'$ continuous at
  every point of $\alpha_0$. Write $A_0(t) = f'(\alpha_0(t))$ for the operator
  curve along the solution and assume $T \sup_t \|A_0(t)\| < 1$. If $\sigma$
  assigns to every initial value $x$ near $x_0$ a curve
  $\sigma(x) \in C([0,T], E)$ solving the integral equation
  $\sigma(x)(t) = x + \int_0^t f(\sigma(x)(s))\,ds$, with
  $\sigma(x_0) = \alpha_0$ and $\sigma$ continuous at $x_0$, then $\sigma$ is
  strictly Fr\'echet differentiable at $x_0$ as a map $E \to C([0,T], E)$. Its
  derivative $D$ exists and is characterized by the linearized (variational)
  integral equation
  $$
  (Dv)(t) - \int_0^t A_0(s)\,\big((Dv)(s)\big)\,ds = v, \qquad v \in E,
  $$
  and for each fixed time $t$ the time-$t$ map $x \mapsto \sigma(x)(t)$ is
  strictly differentiable at $x_0$ with derivative $v \mapsto (Dv)(t)$. This is
  the $C^1$ dependence of the flow on its initial condition at an arbitrary
  (non-equilibrium) base solution, the pointwise-$C^1$ core of
  \cref{thm:dc-ch3-2-2}.
\end{lemma}
\begin{proof}
  \sketch
  \leanok
  A curve $\alpha$ solves the ODE with initial value $x$ iff the Picard
  residual $G(x, \alpha) = \alpha - x - \int_0^t f(\alpha)$ vanishes, viewed as
  a map $E \times C([0,T], E) \to C([0,T], E)$ of Banach spaces. The only
  nonlinearity of $G$ is the superposition operator
  $\alpha \mapsto f \circ \alpha$; its strict derivative at the base curve
  $\alpha_0$ is the postcomposition operator
  $\beta \mapsto (t \mapsto A_0(t)\,\beta(t))$. Indeed, for curves
  $\alpha, \beta$ uniformly within $\delta$ of $\alpha_0$, the pointwise error
  $f(\alpha(t)) - f(\beta(t)) - f'(\alpha_0(t))(\alpha(t) - \beta(t))$ is at
  most $\varepsilon\,\|\alpha(t) - \beta(t)\|$ by the mean value inequality
  applied to $z \mapsto f(z) - f'(\alpha_0(t))z$ on the ball
  $B_\delta(\alpha_0(t))$; the estimate holds simultaneously for all $t$
  because the range of $\alpha_0$ is compact, so a uniform $\delta$-thickening
  of the range stays inside $u$ and $f'$ is uniformly continuous at the range
  (Heine: continuity at each point of a compact set is uniform against nearby
  points, even ones outside the set). Hence $G$ is strictly differentiable at
  $(x_0, \alpha_0)$, with partial derivative in the curve variable
  $1 - J \circ M$, where $J$ is the Volterra primitive
  ($\|J\| \le T$) and $M$ the postcomposition by $A_0$
  ($\|M\| \le \sup_t \|A_0(t)\|$). Since $T \sup_t \|A_0(t)\| < 1$, the
  operator $1 - J \circ M$ is invertible by the Neumann series, and the Banach
  implicit function theorem applies to $G$ at $(x_0, \alpha_0)$: the local zero
  family is strictly differentiable at $x_0$, and it coincides with $\sigma$
  near $x_0$ because $\sigma$ is continuous at $x_0$ and solves
  $G(x, \sigma(x)) = 0$. The derivative equation
  $(1 - J \circ M) \circ D = \mathrm{const}$ is exactly the variational
  integral equation, solved by $D = (1 - J \circ M)^{-1} \circ \mathrm{const}$
  and determining $D$ uniquely; postcomposing with the bounded evaluation at a
  fixed time $t$ gives the time-$t$ statement.
\end{proof}

\begin{lemma}[the variational flow computes the derivative of the flow]
  \label{lem:dc-ch3-2-2-variational}
  \lean{Riemannian.FlowDependence.flowDeriv_eq_applyCurve_opFlow, Riemannian.FlowDependence.variationalField, Riemannian.FlowDependence.picardResidual_variationalField_compRight, Riemannian.FlowDependence.applyCurve_sub_intervalPrimitive_of_picardResidual_variationalField}
  \uses{lem:dc-ch3-2-2-c1dep}
  \leanok
  \dcref{ch3:2.2}
  Let $f$ be a vector field on a Banach space $E$, twice differentiable on an
  open set $u$, and let $\alpha_0 : [0,T] \to u$ solve $\alpha' = f(\alpha)$
  with $T \sup_t \|f'(\alpha_0(t))\| < 1$. Consider the \emph{extended
  (variational) field} on $E \times L(E,E)$,
  $$\hat f(z, W) = \big(f(z),\; f'(z) \circ W\big),$$
  whose trajectories pair a trajectory of $f$ with the solution of the
  linearized equation $W' = f'(\alpha(t)) \circ W$ along it. If $(\alpha_0, W)$
  solves the integral equation of $\hat f$ with initial condition $(x_0, 1)$,
  then the derivative $D$ of the solution family of $f$ at $x_0$
  (\cref{lem:dc-ch3-2-2-c1dep}) is computed by the operator flow:
  $Dv = (t \mapsto W(t)\,v)$ for every $v \in E$. Moreover right composition
  scales the operator slot: $(\alpha_0, W \circ W_0)$ solves the extended
  equation with initial condition $(x_0, W_0)$, for every $W_0 \in L(E,E)$.
\end{lemma}
\begin{proof}
  \sketch
  \leanok
  Applying the second component of the extended integral equation to a fixed
  vector $v$, and commuting the (bounded linear) evaluation with the Volterra
  primitive, shows that the curve $t \mapsto W(t)v$ satisfies
  $$W(t)v - \int_0^t f'(\alpha_0(s))\big(W(s)v\big)\,ds = v ,$$
  which is exactly the variational integral equation
  $(1 - J \circ M)\beta = \mathrm{const}\, v$ characterizing $Dv$
  (\cref{lem:dc-ch3-2-2-c1dep}). Since $T \sup_t \|f'(\alpha_0(t))\| < 1$, the
  operator $1 - J \circ M$ is invertible by the Neumann series, so the two
  solutions coincide: $Dv = (t \mapsto W(t)v)$. For the scaling clause, the map
  $\psi(z, W) = (z, W \circ W_0)$ is linear and bounded, commutes with the
  extended field ($\hat f \circ \psi = \psi \circ \hat f$, by associativity of
  composition), with the constant embedding ($\psi(x_0, 1) = (x_0, W_0)$), and
  with the primitive; hence it maps zeros of the extended residual with initial
  condition $(x_0, 1)$ to zeros with initial condition $(x_0, W_0)$.
\end{proof}

\begin{lemma}[second-order dependence on the initial condition]
  \label{lem:dc-ch3-2-2-c2dep}
  \lean{Riemannian.FlowDependence.exists_hasStrictFDerivAt_opFlow_of_picardResidual_curve, Riemannian.FlowDependence.hasFDerivAt_variationalField, Riemannian.FlowDependence.variationalFieldDeriv}
  \uses{lem:dc-ch3-2-2-c1dep, lem:dc-ch3-2-2-variational}
  \leanok
  \dcref{ch3:2.2}
  In the setting of \cref{lem:dc-ch3-2-2-variational}, suppose additionally
  that for every initial value $x$ near $x_0$ the operator flow
  $\tau(x) \in C([0,T], L(E,E))$ solves the extended integral equation along
  $\sigma(x)$ with initial condition $(x, 1)$, that $\tau$ is continuous at
  $x_0$, and that the linearization $\hat A_0$ of the extended field along the
  base pair $(\alpha_0, \tau(x_0))$ satisfies $T \sup_t \|\hat A_0(t)\| < 1$.
  Then the family of variational flows $x \mapsto \tau(x)$ is strictly
  Fr\'echet differentiable at $x_0$. Combined with
  \cref{lem:dc-ch3-2-2-variational} — the first derivative of the flow is
  computed by $\tau$ — this is second-order differentiability of the flow in
  its initial condition, the $C^2$ step of the joint regularity asserted by
  \cref{thm:dc-ch3-2-2}.
\end{lemma}
\begin{proof}
  \sketch
  \leanok
  The extended field $\hat f(z, W) = (f(z), f'(z) \circ W)$ is differentiable
  on $u \times L(E,E)$ with derivative
  $$D\hat f_{(z,W)}(u', V) = \big(f'(z)u',\; f'(z) \circ V
    + (f''(z)u') \circ W\big),$$
  continuous wherever $f'$ and $f''$ are continuous: the first slot is the
  chain rule, and the second is the derivative of the bounded bilinear
  composition. The pairs $\hat\sigma(x, W_0) = (\sigma(x), \tau(x) \circ W_0)$
  form a solution family of $\hat f$ for \emph{all} initial conditions
  $(x, W_0)$ near $(x_0, 1)$ — right composition scales the operator slot
  (\cref{lem:dc-ch3-2-2-variational}) — and $\hat\sigma$ is continuous at
  $(x_0, 1)$, since $(W_0, \gamma) \mapsto \gamma \circ W_0$ is bounded
  bilinear. The trajectories $\sigma(x)$ stay in $u$ for $x$ near $x_0$: the
  compact range of $\alpha_0$ has a uniform thickening inside $u$, and
  $\sigma$ is continuous at $x_0$ in the sup norm. Since
  $T \sup_t \|\hat A_0(t)\| < 1$, the $C^1$-dependence theorem
  (\cref{lem:dc-ch3-2-2-c1dep}, applied on the Banach space
  $E \times L(E,E)$) makes $\hat\sigma$ strictly differentiable at $(x_0, 1)$;
  composing with the affine embedding $x \mapsto (x, 1)$ and the bounded
  projection to the operator component yields the strict differentiability of
  $\tau$ at $x_0$.
\end{proof}

\begin{theorem}[flow of a vector field]
  \label{thm:dc-ch3-2-2}
  \uses{lem:dc-ch3-2-2-pointwise, lem:dc-ch3-2-2-flow, lem:dc-ch3-2-2-c1dep, lem:dc-ch3-2-2-variational, lem:dc-ch3-2-2-c2dep, lem:dc-ch3-2-2-cinfty}
  \notready
  \dcref{ch3:2.2}
  \notready
  If $X$ is a $C^\infty$ vector field on the open set $V$ in the manifold $M$ and
  $p\in V$, then there exist an open set $V_0\subset V$ with $p\in V_0$, a number
  $\delta>0$, and a $C^\infty$ mapping $\varphi:(-\delta,\delta)\times V_0\to V$ such
  that $t\mapsto\varphi(t,q)$, $t\in(-\delta,\delta)$, is the unique trajectory of
  $X$ passing through $q$ at $t=0$, for every $q\in V_0$. The map
  $\varphi_t:V_0\to V$, $\varphi_t(q)=\varphi(t,q)$, is the \emph{flow} of $X$ on $V$.
  (Beyond \cref{lem:dc-ch3-2-2-pointwise}, this asserts a uniform time of
  existence and the joint $C^\infty$ regularity of $\varphi$ in $(t,q)$ — the
  smooth dependence of solutions of an ODE on their initial conditions. The
  uniform time and Lipschitz dependence is \cref{lem:dc-ch3-2-2-flow}; the
  pointwise $C^1$ dependence on the initial condition, at an arbitrary base
  solution, is \cref{lem:dc-ch3-2-2-c1dep}; the second-order dependence, given
  the variational flow along the base solution, is \cref{lem:dc-ch3-2-2-c2dep}
  via \cref{lem:dc-ch3-2-2-variational}; the construction of the variational
  flow for the geodesic spray and the joint $C^k$ regularity for $k \ge 3$
  remain open.)
\end{theorem}

\begin{lemma}[the geodesic field]
  \label{lem:dc-ch3-2-3}
  \lean{Riemannian.Geodesic.existsUnique_geodesicField}
  \uses{def:dc-ch3-2-1}
  \leanok
  \dcref{ch3:2.3}
  \notready
  There exists a unique vector field $G$ on $TM$ whose trajectories are of the form
  $t\mapsto(\gamma(t),\gamma'(t))$, where $\gamma$ is a geodesic on $M$.
\end{lemma}
\begin{proof}
  \leanok
  \sketch
  Uniqueness: in a coordinate system $(U,\mathbf{x})$ the trajectories of
  $G$ on $TU$ are $t\mapsto(\gamma(t),\gamma'(t))$ with $\gamma$ a geodesic, hence
  solutions of (1'); by uniqueness of trajectories of such a system, $G$ is unique
  if it exists. Existence: define $G$ locally by (1'); uniqueness shows $G$ is
  well-defined on all of $TM$.
\end{proof}

\begin{definition}[geodesic field / geodesic flow]
  \label{def:dc-ch3-2-4}
  \lean{Riemannian.Geodesic.geodesicVectorField, Riemannian.Geodesic.geodesicVectorFieldChart, Riemannian.Geodesic.geodesicVectorFieldChart_contMDiffOn}
  \uses{thm:dc-ch3-2-2, lem:dc-ch3-2-3}
  \leanok
  \dcref{ch3:2.4}
  The vector field $G$ above is the \emph{geodesic field} on $TM$, and its flow is the
  \emph{geodesic flow} on $TM$. Applying \cref{thm:dc-ch3-2-2} to $G$ at $(p,0)\in TM$: for each
  $p\in M$ there exist an open set $\mathcal{U}\subset TU$ with $(p,0)\in\mathcal{U}$,
  a number $\delta>0$, and a $C^\infty$ map $\varphi:(-\delta,\delta)\times\mathcal{U}\to TU$
  such that $t\mapsto\varphi(t,q,v)$ is the unique trajectory of $G$ with
  $\varphi(0,q,v)=(q,v)$. One may take
  $$
  \mathcal{U}=\{(q,v)\in TU;\ q\in V,\ v\in T_qM,\ |v|<\varepsilon_1\},
  $$
  $V\subset U$ a neighborhood of $p$. Putting $\gamma=\pi\circ\varphi$, where
  $\pi:TM\to M$ is the canonical projection, we can describe the previous result as
  follows.
\end{definition}

\begin{lemma}[pointwise existence of a geodesic with prescribed initial velocity]
  \label{lem:dc-ch3-2-5-exist}
  \lean{Riemannian.Geodesic.exists_geodesic_with_initial_velocity_at}
  \uses{def:dc-ch3-2-1, lem:dc-ch3-2-2-pointwise}
  \leanok
  \dcref{ch3:2.5}
  For every $p\in M$ and every $v\in T_pM$ there is a curve $\gamma:\mathbb{R}\to M$
  with $\gamma(0)=p$ that is a local geodesic at $t=0$ and whose canonical tangent
  lift $t\mapsto(\gamma(t),\gamma'(t))\in TM$ takes the value $(p,v)$ at $t=0$; i.e.
  $\gamma$ is a geodesic through $p$ with initial velocity $v$.
\end{lemma}
\begin{proof}
  \sketch
  \leanok
  Read the geodesic equation on $TM$ as the first-order system (1') for the
  geodesic spray $G$. The chart-fixed spray is $C^\infty$ at $(p,v)$, so the
  integral-curve existence theorem (Picard--Lindel\"of,
  \cref{lem:dc-ch3-2-2-pointwise}) produces a local integral curve $f$ of $G$ with
  $f(0)=(p,v)$; its base projection $\gamma=\pi\circ f$ is the desired geodesic.
\end{proof}

\begin{lemma}[pointwise uniqueness of geodesics]
  \label{lem:dc-ch3-2-5-unique}
  \lean{Riemannian.Geodesic.isGeodesicAt_eventuallyEq_of_lift_eq}
  \uses{def:dc-ch3-2-1, lem:dc-ch3-2-2-pointwise}
  \leanok
  \dcref{ch3:2.5}
  If two curves $\gamma_1,\gamma_2$ are base projections of integral curves of the
  geodesic spray whose lifts agree at $t_0$ (equal initial point \emph{and} initial
  velocity), then $\gamma_1$ and $\gamma_2$ agree on a neighbourhood of $t_0$.
\end{lemma}
\begin{proof}
  \sketch
  \leanok
  The lifts are integral curves of the same $C^\infty$ vector field $G$ on $TM$
  with the same value at $t_0$, so by uniqueness of integral curves
  (\cref{lem:dc-ch3-2-2-pointwise}, Gr\"onwall) they agree near $t_0$; projecting to $M$
  gives $\gamma_1=\gamma_2$ near $t_0$.
\end{proof}

\begin{lemma}[a geodesic leaves its initial point with its initial velocity]
  \label{lem:dc-ch3-2-5-velocity}
  \lean{Riemannian.Geodesic.IsGeodesicOnWithInitial.hasDerivAt_extChartAt_zero, Riemannian.Geodesic.hasDerivAt_extChartAt_maximalGeodesic}
  \uses{lem:dc-ch3-2-5-exist, lem:dc-ch3-2-5-unique}
  \leanok
  \dcref{ch3:2.5}
  A geodesic with initial data $(p,v)$ really has velocity $v$ at time $0$: read
  in the chart at $p$, the curve $s\mapsto\varphi_p(\gamma(s))$ is differentiable
  at $s=0$ with derivative $v$. The same holds for the canonical maximal
  geodesic $\gamma(\cdot,p,v)$.
\end{lemma}
\begin{proof}
  \sketch
  \leanok
  The velocity lift $f$ of the geodesic is an integral curve of the geodesic
  field with $f(0)=(p,v)$. Reading the integral-curve equation in the chart of
  $TM$ at $(p,v)$ — which is based at the foot $p$ — the pair
  $(x,w)=(\varphi_p\circ\gamma,\ \text{fibre coordinate of }f)$ satisfies the
  first-order system $(x',w')=(w,\,-\Gamma_p(w,w)(x))$ near $0$; in particular
  $x'(0)=w(0)=v$, since the trivialisation at $p$ is the identity on the fibre
  over $p$. For the canonical maximal geodesic, local uniqueness
  (\cref{lem:dc-ch3-2-5-unique}) propagated along the connected witness interval
  identifies it with a single witness near $0$.
\end{proof}

\begin{lemma}[a geodesic satisfies the geodesic equation at the initial time]
  \label{lem:dc-ch3-2-5-equation}
  \lean{Riemannian.Geodesic.IsGeodesicOnWithInitial.hasGeodesicEquationAt_zero, Riemannian.Geodesic.hasGeodesicEquationAt_maximalGeodesic_zero}
  \uses{def:dc-ch3-2-1, lem:dc-ch3-2-5-exist, lem:dc-ch3-2-5-velocity}
  \leanok
  \dcref{ch3:2.5}
  A geodesic with initial data $(p,v)$ — the base projection of an integral
  curve of the geodesic field — satisfies the second-order geodesic equation of
  \cref{def:dc-ch3-2-1} at the initial time: in the chart at the foot
  $\gamma(0)=p$, the chart curve $x=\varphi_p\circ\gamma$ has a second
  derivative at $0$ and $x''(0)+\Gamma_p(v,v)(x(0))=0$. The same holds for the
  canonical maximal geodesic. This bridges the first-order (geodesic-field)
  formulation back to the second-order equation at the initial time.
\end{lemma}
\begin{proof}
  \sketch
  \leanok
  With $(x,w)$ as in \cref{lem:dc-ch3-2-5-velocity}, the first-order system
  gives $x'=w$ near $0$ and $w'(0)=-\Gamma_p(w(0),w(0))(x(0))$; hence $x'$ is
  differentiable at $0$ with
  $x''(0)=w'(0)=-\Gamma_p(v,v)(x(0))$, which is the geodesic equation at $0$.

\end{proof}

\begin{proposition}[local existence/uniqueness of geodesics]
  \label{prop:dc-ch3-2-5}
  \lean{Riemannian.GeodesicLocal.geodesic_local_existence_metricBall}
  \uses{lem:dc-ch3-2-5-exist, lem:dc-ch3-2-5-unique, lem:dc-ch3-2-5-chart-family, lem:dc-ch3-2-5-metric-family}
  \dcref{ch3:2.5}
  Given $p\in M$, there exist an open set $V\subset M$ with $p\in V$, numbers
  $\delta>0$, $\varepsilon_1>0$, and a $C^\infty$ map
  $$
  \gamma:(-\delta,\delta)\times\mathcal{U}\to M,\quad\mathcal{U}=\{(q,v);\ q\in V,\ v\in T_qM,\ |v|<\varepsilon_1\},
  $$
  such that $t\mapsto\gamma(t,q,v)$, $t\in(-\delta,\delta)$, is the unique geodesic
  passing through $q$ with velocity $v$ at $t=0$, for each $q\in V$ and each
  $v\in T_qM$ with $|v|<\varepsilon_1$.
\end{proposition}

\begin{lemma}[affine rescaling of a geodesic]
  \label{lem:dc-ch3-2-6-reparam}
  \lean{Riemannian.Geodesic.isGeodesic_comp_mul_left}
  \leanok
  \dcref{ch3:2.6}
  If $\gamma:\mathbb{R}\to M$ is a geodesic, then for every $a\in\mathbb{R}$ the
  affine reparametrisation $h(t)=\gamma(at)$ is again a geodesic, with velocity
  $h'(t)=a\,\gamma'(at)$.
\end{lemma}
\begin{proof}
  \sketch
  \leanok
  Read the geodesic equation in a chart at the moving foot $\gamma(at)$: with
  $u(s)=\mathbf{x}(\gamma(s))$ the chart representative, $h$ reads as $s\mapsto
  u(as)$, so $h'(t)=a\,u'(at)$ and $h''(t)=a^2\,u''(at)$. Since the Christoffel
  contraction $\Gamma$ scales as a $(2,0)$-tensor,
  $\Gamma(au'(at),au'(at))=a^2\,\Gamma(u'(at),u'(at))$, hence
  $$
  h''(t)+\Gamma\big(h'(t),h'(t)\big)
    =a^2\Big(u''(at)+\Gamma\big(u'(at),u'(at)\big)\Big)=a^2\cdot 0=0,
  $$
  so $h$ satisfies the geodesic equation at every $t$.
\end{proof}

\begin{lemma}[homogeneity of a geodesic]
  \label{lem:dc-ch3-2-6}
  \lean{Riemannian.Geodesic.IsGeodesicOnWithInitial.fiberScale, Riemannian.Geodesic.maximalGeodesicInterval_fiberScale, Riemannian.Geodesic.maximalGeodesic_fiberScale}
  \uses{lem:dc-ch3-2-6-reparam, lem:dc-ch3-2-5-unique}
  \leanok
  \dcref{ch3:2.6}
  If the geodesic $\gamma(t,q,v)$ is defined on $(-\delta,\delta)$, then for
  $a\in\mathbb{R}$, $a>0$, the geodesic $\gamma(t,q,av)$ is defined on
  $(-\frac{\delta}{a},\frac{\delta}{a})$ and

  $$
  \gamma(t,q,av)=\gamma(at,q,v).
  $$
\end{lemma}
\begin{proof}
  \sketch
  \leanok
  Let $h:(-\frac{\delta}{a},\frac{\delta}{a})\to M$, $h(t)=\gamma(at,q,v)$.
  Then $h(0)=q$, $\frac{dh}{dt}(0)=av$, and by \cref{lem:dc-ch3-2-6-reparam}
  $h$ is a geodesic (with $h'(t)=a\gamma'(at,q,v)$), so

  $$
  \frac{D}{dt}\Big(\frac{dh}{dt}\Big)=\nabla_{h'(t)}h'(t)=a^2\nabla_{\gamma'(at,q,v)}\gamma'(at,q,v)=0.
  $$

  So $h$ is a geodesic through $q$ with velocity $av$ at $t=0$; by uniqueness
  $h(t)=\gamma(at,q,v)=\gamma(t,q,av)$.
\end{proof}

\begin{lemma}[pointwise uniform interval of definition]
  \label{lem:dc-ch3-2-7-pointwise}
  \lean{Riemannian.Geodesic.exists_pos_smul_maximalGeodesicInterval}
  \uses{lem:dc-ch3-2-6}
  \leanok
  \dcref{ch3:2.7}
  For every $p\in M$, $v\in T_pM$ and every $T>0$ there is a $c>0$ such that
  the geodesic with initial data $(p,\,c\,v)$ is defined on $(-T,T)$. This is
  the pointwise-in-$(p,v)$ content of \cref{prop:dc-ch3-2-7}: shrinking the
  velocity extends the interval of definition. (The uniformity in the velocity
  over a ball $|v|\le r$ is now closed by \cref{lem:dc-ch3-2-7-vuniform}; the
  uniformity in the base point $q \in V$ together with the joint $C^\infty$
  regularity of the family — the smooth-family clause of
  \cref{prop:dc-ch3-2-7} — remains open, pending smooth dependence of ODE
  solutions on initial data.)
\end{lemma}
\begin{proof}
  \sketch
  \leanok
  The maximal interval $I_{\max}(p,v)$ is open and contains $0$
  (\cref{lem:dc-ch3-2-5-exist}), so it contains some $(-\delta,\delta)$,
  $\delta>0$. By the homogeneity of geodesics (\cref{lem:dc-ch3-2-6}),
  $I_{\max}(p,\,c\,v)=\{t:\ c\,t\in I_{\max}(p,v)\}$; taking
  $c=\frac{\delta}{2T}$ sends $(-T,T)$ inside $(-\delta/2,\delta/2)\subset
  I_{\max}(p,v)$, so $(-T,T)\subset I_{\max}(p,\,c\,v)$.
\end{proof}

\begin{lemma}[uniform interval in the velocity]
  \label{lem:dc-ch3-2-7-vuniform}
  \lean{Riemannian.Geodesic.exists_uniform_geodesic_flow, Riemannian.Geodesic.isGeodesicOnWithInitial_of_hasDerivAt_sprayCoord, Riemannian.Geodesic.maximalGeodesic_eq_witness_of_mem_chart, Riemannian.Geodesic.contDiffOn_geodesicSprayCoord_prod}
  \uses{lem:dc-ch3-2-2-flow, def:dc-ch3-2-4, lem:dc-ch3-2-5-exist, lem:dc-ch3-2-5-unique}
  \leanok
  \dcref{ch3:2.7}
  For each $p \in M$ there are $r > 0$, $\varepsilon > 0$ and a local flow $Z$
  of the coordinate geodesic spray at $p$ such that for \emph{every} velocity
  $v \in T_pM$ with $|v| \le r$: the geodesic $\gamma(\cdot, p, v)$ is defined
  on $(-\varepsilon, \varepsilon)$ — one interval for all such $v$ — and its
  chart reading is computed by the flow,
  $\varphi_p(\gamma(s, p, v)) = Z((\varphi_p(p), v), s)_1$; the flow is
  Lipschitz in $v$ uniformly in time, and the zero-velocity flow line is the
  constant $p$. This is the uniform-in-velocity clause of
  \cref{prop:dc-ch3-2-7} at a fixed base point.
\end{lemma}
\begin{proof}
  \sketch
  \leanok
  Apply the Picard--Lindelöf local flow (\cref{lem:dc-ch3-2-2-flow}) to the
  coordinate spray $F(x, w) = (w, -\Gamma_p(w, w)(x))$, which is $C^\infty$ on
  the chart target; trajectories are confined to the chart by Lipschitz
  dependence on the initial condition together with continuity of the centre
  trajectory (no compactness is needed). Each flow line, read through the
  tangent-bundle chart at $(p, 0)$, descends to a geodesic witness with foot in
  the chart at $p$ — the reverse of the chart reading of the integral-curve
  equation. Finally, the canonical maximal geodesic agrees with \emph{any}
  witness whose own foot stays in the chart: the clopen uniqueness propagation
  constrains only the given witness, so no chart-validity hypothesis about all
  witnesses is required.
\end{proof}

\begin{lemma}[the geodesic local flow is $C^1$ in the initial condition]
  \label{lem:dc-ch3-2-7-c1flow}
  \lean{Riemannian.Geodesic.exists_uniform_geodesic_flow_hasStrictFDerivAt}
  \uses{lem:dc-ch3-2-2-c1dep, lem:dc-ch3-2-7-vuniform, def:dc-ch3-2-4}
  \leanok
  \dcref{ch3:2.7}
  For each $p \in M$ there are $r, \varepsilon > 0$, a local flow $Z$ of the
  coordinate geodesic spray at $p$ with the properties of
  \cref{lem:dc-ch3-2-7-vuniform} (one time interval $[-\varepsilon,
  \varepsilon]$ for all initial conditions in the flow ball of radius $r$,
  Lipschitz dependence on the initial condition, chart confinement, and
  computation of the canonical maximal geodesic by the flow), and a time
  $0 < T < \varepsilon$, such that the flow — read as the family
  $\sigma : x \mapsto Z(x, \cdot)|_{[0,T]} \in C([0,T], \mathbb{R}^{2n})$ of
  its solution curves — is strictly Fr\'echet differentiable in the initial
  condition at \emph{every} $x_0$ in the open flow ball, equilibrium or not,
  with derivative characterized by the linearized (variational) integral
  equation of \cref{lem:dc-ch3-2-2-c1dep} along the base trajectory
  $\sigma(x_0)$. This is the pointwise-$C^1$ part, in the initial condition,
  of the $C^\infty$ family clause of \cref{prop:dc-ch3-2-7}.
\end{lemma}
\begin{proof}
  \sketch
  \leanok
  Trajectories starting in the flow ball are confined — by Lipschitz
  dependence on the initial condition together with constancy of the centre
  trajectory — to a ball whose closure is a compact subset of the chart
  region. The spray derivative $dF$ is continuous on the chart region, hence
  bounded by some $C \ge 0$ on the compact confinement ball; consequently the
  operator curve $A_0(t) = dF(Z(x_0, t))$ of \emph{every} base trajectory has
  sup norm at most $C$. Choose $T < \min(\varepsilon, 1/C)$; then
  $T \sup_t \|A_0(t)\| \le T C < 1$ simultaneously for all initial conditions
  in the flow ball. At each $x_0$ the family $\sigma$ satisfies the
  hypotheses of \cref{lem:dc-ch3-2-2-c1dep}: each flow line solves the Picard
  integral equation of the spray (its derivative is the spray value, and the
  values stay in the chart region where the spray is continuous), and
  $\sigma$ is Lipschitz — hence continuous — in the initial condition. The
  conclusion of that lemma is the claimed strict differentiability with the
  variational characterization.
\end{proof}

\begin{lemma}[the geodesic local flow is $C^2$ in the initial condition]
  \label{lem:dc-ch3-2-7-c2flow}
  \lean{Riemannian.Geodesic.exists_uniform_geodesic_flow_hasStrictFDerivAt_opFlow, Riemannian.FlowDependence.picardResidual_variationalField_eq_zero, Riemannian.FlowDependence.norm_variationalFieldDeriv_le}
  \uses{lem:dc-ch3-2-7-c1flow, lem:dc-ch3-2-2-variational, lem:dc-ch3-2-2-c2dep, def:dc-ch3-2-4}
  \leanok
  \dcref{ch3:2.7}
  For each $p \in M$ there are $r, \varepsilon > 0$, a local flow $Z$ of the
  coordinate geodesic spray at $p$ with the properties of
  \cref{lem:dc-ch3-2-7-vuniform}, a time $0 < T < \varepsilon$, the flow read
  as the family $\sigma : x \mapsto Z(x, \cdot)|_{[0,T]}$ of its solution
  curves, and the family $\tau : x \mapsto W_x \in C([0,T],
  L(\mathbb{R}^{2n}))$ of \emph{variational (operator) flows} along the
  trajectories — $W_x$ solves $W' = dF(\sigma(x)(t)) \circ W$, $W(0) = 1$ —
  such that at \emph{every} $x_0$ in the open flow ball:
  \begin{itemize}
  \item $\sigma$ is strictly Fr\'echet differentiable at $x_0$ and its
    derivative \emph{is computed by the variational flow}:
    $D\sigma(x_0)\,v = \big(t \mapsto \tau(x_0)(t)\,v\big)$
    (\cref{lem:dc-ch3-2-2-variational});
  \item the variational-flow family $\tau$ is itself strictly Fr\'echet
    differentiable at $x_0$ (\cref{lem:dc-ch3-2-2-c2dep}).
  \end{itemize}
  Together these say the local geodesic flow is $C^2$ in its initial
  condition; this is the $C^2$ part, in the initial condition, of the
  $C^\infty$ family clause of \cref{prop:dc-ch3-2-7}.
\end{lemma}
\begin{proof}
  \sketch
  \leanok
  As in \cref{lem:dc-ch3-2-7-c1flow}, trajectories from the flow ball are
  confined to a compact ball inside the chart region, on which the first and
  second derivatives of the spray are bounded, say $\|dF\| \le C$ and
  $\|d^2F\| \le C_2$. Choose the Picard time
  $T < \min\big(\varepsilon,\ \tfrac{1}{2(C + 2C_2 + 1)}\big)$, so that the
  contraction estimates below hold simultaneously for every initial condition
  in the flow ball.

  \emph{Construction of $\tau$.} For each $x$, the variational equation is
  the linear Volterra fixed-point equation $W = 1 + J_x W$ on $C([0,T],
  L(\mathbb{R}^{2n}))$, where $(J_x W)(t) = \int_0^t dF(\sigma(x)(s)) \circ
  W(s)\, ds$. Since $\|J_x\| \le T\,C \le \tfrac12 < 1$, the Neumann series
  inverts $1 - J_x$ and we set $\tau(x) = (1 - J_x)^{-1}(\mathbf{1})$. This
  single formula yields existence, the fixed-point identity
  $\tau(x) = \mathbf{1} + J_x \tau(x)$ (hence the uniform bound
  $\|\tau(x)\| \le (1 - \tfrac12)^{-1} = 2$), \emph{and} continuity of
  $x \mapsto \tau(x)$: the coefficient curve $t \mapsto dF(\sigma(x)(t))$
  depends continuously on $x$ (uniform continuity of $dF$ on the compact
  confinement ball together with Lipschitz dependence of the flow on its
  initial condition), $J_x$ depends continuously — indeed
  $\|J_x - J_{x'}\| \le T \sup_t \|dF(\sigma(x)(t)) - dF(\sigma(x')(t))\|$ —
  and inversion is continuous at units of a Banach algebra.

  \emph{The extended system.} The pair $(\sigma(x), \tau(x))$ is then a
  solution of the Picard integral equation of the extended (variational)
  field $\hat F(z, W) = (F(z),\ dF(z) \circ W)$ with initial condition
  $(x, 1)$: the residual splits into its two components, the first being the
  Picard residual of the base flow and the second the Volterra fixed-point
  identity above.

  \emph{Conclusion.} At each $x_0$ in the open flow ball, the derivative
  characterization of \cref{lem:dc-ch3-2-7-c1flow} and the uniqueness of the
  linearized fixed point identify $D\sigma(x_0)\,v$ with
  $t \mapsto \tau(x_0)(t)\,v$ (\cref{lem:dc-ch3-2-2-variational}). For the
  second derivative, apply the abstract second-order dependence theorem
  (\cref{lem:dc-ch3-2-2-c2dep}) to the extended system: its linearization
  along the base pair $(\sigma(x_0), \tau(x_0))$ is bounded by
  $C + C_2 \|\tau(x_0)\| \le C + 2C_2$ (the block estimate
  $\|D\hat F(z, W)\| \le \|dF(z)\| + \|d^2F(z)\|\,\|W\|$), so
  $T\,(C + 2C_2) < 1$ makes its hypotheses hold, and the family
  $x \mapsto \tau(x)$ is strictly differentiable at $x_0$.
\end{proof}

\begin{proposition}[uniform interval of definition]
  \label{prop:dc-ch3-2-7}
  \lean{Riemannian.GeodesicLocal.geodesic_local_existence_fixedInterval_metricBall}
  \uses{prop:dc-ch3-2-5, lem:dc-ch3-2-6, lem:dc-ch3-2-7-vuniform, lem:dc-ch3-2-7-c1flow, lem:dc-ch3-2-7-c2flow, lem:dc-ch3-2-5-unique, lem:dc-ch3-2-5-metric-ball}
  \leanok
  \dcref{ch3:2.7}
  \notready
  Given $p\in M$, there exist a neighborhood $V$ of $p$ in $M$, a number
  $\varepsilon>0$, and a $C^\infty$ map $\gamma:(-2,2)\times\mathcal{U}\to M$,
  $\mathcal{U}=\{(q,w)\in TM;\ q\in V,\ w\in T_qM,\ |w|<\varepsilon\}$, such that
  $t\mapsto\gamma(t,q,w)$, $t\in(-2,2)$, is the unique geodesic of $M$ passing
  through $q$ with velocity $w$ at $t=0$, for every $q\in V$ and $w\in T_qM$ with
  $|w|<\varepsilon$.
\end{proposition}
\begin{proof}
  \leanok
  \sketch
  The geodesic $\gamma(t,q,v)$ of \cref{prop:dc-ch3-2-5} is defined for
  $|t|<\delta$ and $|v|<\varepsilon_1$. By the lemma of homogeneity (2.6),
  $\gamma(t,q,\frac{\delta v}{2})$ is defined for $|t|<2$. Taking
  $\varepsilon<\frac{\delta\varepsilon_1}{2}$, the geodesic $\gamma(t,q,w)$ is
  defined for $|t|<2$ and $|w|<\varepsilon$.
\end{proof}

\begin{remark}[large velocity]
  \label{rem:dc-ch3-2-8}
  \lean{Riemannian.Geodesic.maximalGeodesicInterval_fiberScale}
  \uses{lem:dc-ch3-2-6}
  \leanok
  \dcref{ch3:2.8}
  By an analogous argument, one can make the velocity of a geodesic uniformly large
  in a neighborhood of $p$: the exact rescaling
  $I_{\max}(p,\,a\,v)=\{t:\ a\,t\in I_{\max}(p,v)\}$ of the maximal interval
  trades interval length against velocity in both directions.
\end{remark}

\subsection{Exponential map}
\cref{prop:dc-ch3-2-7} lets us define the \emph{exponential map}. Let $p\in M$ and let
$\mathcal{U}\subset TM$ be as in \cref{prop:dc-ch3-2-7}. The map $\exp:\mathcal{U}\to M$,

$$
\exp(q,v)=\gamma(1,q,v)=\gamma\Big(|v|,q,\frac{v}{|v|}\Big),\quad(q,v)\in\mathcal{U},
$$

is the \emph{exponential map} on $\mathcal{U}$; it is differentiable. In applications
we use the restriction to an open set of $T_qM$:

$$
\exp_q:B_\varepsilon(0)\subset T_qM\to M,\qquad\exp_q(v)=\exp(q,v),
$$

where $B_\varepsilon(0)$ is the open ball of center $0$ and radius $\varepsilon$
in $T_qM$. Then $\exp_q$ is differentiable and $\exp_q(0)=q$. Geometrically,
$\exp_q(v)$ is obtained by going a length $|v|$ from $q$ along the geodesic with
velocity $\frac{v}{|v|}$.

\begin{definition}[exponential map]
  \label{def:dc-ch3-expmap}
  \lean{Riemannian.Exponential.expMapIntrinsic, Riemannian.Exponential.expDomainIntrinsic, Riemannian.Exponential.zero_mem_expDomainIntrinsic, Riemannian.Exponential.expMapIntrinsic_zero}
  \uses{lem:dc-ch3-2-5-exist, lem:dc-ch3-2-5-unique}
  \leanok
  \dcref{ch3:after-2.8}
  The \emph{exponential map} at $p\in M$ is $\exp_p(v)=\gamma(1,p,v)$, the value at
  time $1$ of the (maximal) geodesic with initial data $(p,v)$; its natural
  domain is the set of $v\in T_pM$ whose maximal geodesic is defined at time
  $1$, which always contains $0$, and $\exp_p(0)=p$.

\end{definition}

\begin{lemma}[the exponential map traces geodesic rays]
  \label{lem:dc-ch3-expmap-ray}
  \lean{Riemannian.Exponential.expMapIntrinsic_smul_eq_intrinsicMaximalGeodesic_of_mem}
  \uses{def:dc-ch3-expmap, lem:dc-ch3-2-6}
  \leanok
  \dcref{ch3:after-2.8}
  For $v\in T_pM$ and $t$ in the maximal interval of the geodesic with initial
  data $(p,v)$,
  $$\exp_p(t\,v)=\gamma(t,p,v).$$
  In particular $t\mapsto\exp_p(t\,v)$ \emph{is} the geodesic through $p$ with
  initial velocity $v$; this is the computational heart of
  $d(\exp_p)_0=\mathrm{id}$ in \cref{prop:dc-ch3-2-9}.
\end{lemma}
\begin{proof}
  \sketch
  \leanok
  For $t=0$ both sides are $p$. For $t\neq0$ this is the homogeneity
  $\gamma(1,p,t\,v)=\gamma(t\cdot 1,p,v)$ of \cref{lem:dc-ch3-2-6} evaluated at
  time $1$.
\end{proof}

\begin{lemma}[$d(\exp_p)_0=\mathrm{id}$ along rays]
  \label{lem:dc-ch3-2-9-dexp}
  \lean{Riemannian.Exponential.hasDerivAt_extChartAt_expMap_smul}
  \uses{lem:dc-ch3-expmap-ray, lem:dc-ch3-2-5-velocity}
  \leanok
  \dcref{ch3:2.9}
  For every $v\in T_pM$, read in the chart at $p$, the curve
  $t\mapsto\exp_p(t\,v)$ is differentiable at $t=0$ with derivative $v$:
  $$
  d(\exp_p)_0(v)=\frac{d}{dt}\big(\exp_p(tv)\big)\Big|_{t=0}
    =\frac{d}{dt}\big(\gamma(t,p,v)\big)\Big|_{t=0}=v.
  $$
  Together with smooth dependence on initial conditions and the inverse function
  theorem, this derivative identity gives the local regularity and local
  invertibility of $\exp_p$ near $0$.
\end{lemma}
\begin{proof}
  \sketch
  \leanok
  By \cref{lem:dc-ch3-expmap-ray}, $\exp_p(t\,v)=\gamma(t,p,v)$ for $t$ in the
  maximal interval, hence for all $t$ near $0$; by
  \cref{lem:dc-ch3-2-5-velocity} the geodesic $\gamma(\cdot,p,v)$ has velocity
  $v$ at $0$.
\end{proof}

\begin{lemma}[$\exp_p$ is strictly differentiable at $0$ with derivative the identity]
  \label{lem:dc-ch3-2-9-strict}
  \lean{Riemannian.Exponential.exists_hasStrictFDerivAt_extChartAt_expMap, Riemannian.Exponential.map_expMap_nhds, Riemannian.Exponential.exists_injOn_expMap, Riemannian.FlowDependence.hasStrictFDerivAt_of_picardResidual}
  \uses{def:dc-ch3-expmap, lem:dc-ch3-2-6, lem:dc-ch3-2-7-vuniform, lem:dc-ch3-2-9-dexp}
  \leanok
  \dcref{ch3:2.9}
  There is $\rho > 0$ such that the ball $B_\rho(0) \subset T_pM$ lies in the
  domain of $\exp_p$, and, read in the chart at $p$, the map $\exp_p$ is
  \emph{strictly} Fréchet differentiable at $0$ with derivative the identity.
  Consequently (Banach inverse function theorem) $\exp_p$ maps the
  neighbourhood filter of $0$ exactly onto the neighbourhood filter of $p$ —
  the image of any neighbourhood of $0$ is a neighbourhood of $p$ — and
  $\exp_p$ is injective on a ball around $0$. This is
  \cref{prop:dc-ch3-2-9} minus the $C^\infty$ regularity of $\exp_p$ away from
  the strict-derivative point.
\end{lemma}
\begin{proof}
  \sketch
  \leanok
  $C^1$ dependence on initial conditions at an equilibrium, on the Banach
  space $C([0,T], E \times E)$: the Picard integral equation
  $G(x, \alpha) = \alpha - x - \int_0^t F(\alpha)$ is strictly differentiable
  at the zero-section equilibrium — the superposition operator
  $\alpha \mapsto F \circ \alpha$ has a strict derivative at a constant curve
  as soon as $DF$ is continuous at the point — and its curve-partial
  derivative $1 - J \circ M_{DF}$ is invertible by a Neumann series once
  $T\,\|DF\| < 1$; the Banach implicit function theorem then makes the
  solution family strictly differentiable in its initial condition. The
  linearization of the geodesic spray at the zero section is
  $A(u, w) = (w, 0)$, quadratically nilpotent ($A^2 = 0$, the Christoffel term
  being quadratic in the velocity), so the Neumann series terminates and the
  derivative is computed exactly; the homogeneity
  $\gamma(1, p, w) = \gamma(T, p, w/T)$ of \cref{lem:dc-ch3-2-6} trades the
  unit evaluation time for the short Picard time $T$. Evaluating at time $T$
  and projecting to the base gives the strict derivative
  $d(\exp_p)_0 = \mathrm{id}$ of the chart reading, and the inverse function
  theorem for strictly differentiable maps yields the neighbourhood and
  injectivity statements.
\end{proof}

\begin{lemma}[$\exp_p$ is strictly differentiable on a ball]
  \label{lem:dc-ch3-2-9-c1ball}
  \lean{Riemannian.Exponential.exists_hasStrictFDerivAt_extChartAt_expMap_ball}
  \uses{def:dc-ch3-expmap, lem:dc-ch3-2-6, lem:dc-ch3-2-7-c1flow}
  \leanok
  \dcref{ch3:2.9}
  There is $\rho > 0$ such that the ball $B_\rho(0) \subset T_pM$ lies in the
  domain of $\exp_p$, its image under $\exp_p$ stays in the chart at $p$, and
  at \emph{every} $v_0$ with $\|v_0\| < \rho$ — zero or not — the chart
  reading $w \mapsto \varphi_p(\exp_p(w))$ has a strict Fr\'echet derivative
  at $v_0$. In particular $\exp_p$ is differentiable at every point of a
  neighbourhood of the origin, which is the regularity clause of
  \cref{prop:dc-ch3-2-9} at the $C^1$ level; at $v_0 \ne 0$ the derivative is
  the solution of the variational equation along the geodesic
  $t \mapsto \exp_p(t\,v_0)$, not the identity.
\end{lemma}
\begin{proof}
  \sketch
  \leanok
  Let $Z$, $r$, $\varepsilon$, $T$ be as in \cref{lem:dc-ch3-2-7-c1flow} and
  put $\rho = rT$. For $\|w\| < \rho$, the homogeneity of geodesics
  (\cref{lem:dc-ch3-2-6}, realized at the level of the flow witness by fibre
  scaling) trades the fixed evaluation time $1$ for the short Picard time
  $T$: $\exp_p(w) = \gamma(1, p, w) = \gamma(T, p, w/T)$, and the chart
  reading of the right-hand side is computed by the flow,
  $\varphi_p(\exp_p(w)) = \big(Z((\varphi_p(p), w/T))(T)\big)_1$. In
  particular $B_\rho(0)$ lies in the exponential domain and its image stays
  in the chart. The map $w \mapsto (\varphi_p(p), w/T)$ is continuous affine
  with linear part $w \mapsto (0, w/T)$, and it sends $B_\rho(0)$ into the
  open flow ball; composing its strict derivative with the strict derivative
  in the initial condition of the flow family (\cref{lem:dc-ch3-2-7-c1flow})
  evaluated at the fixed time $T$, then with the (linear, bounded) projection
  to the base coordinate, yields the strict derivative of the chart reading
  of $\exp_p$ at every $v_0 \in B_\rho(0)$.
\end{proof}

\begin{lemma}[$\exp_p$ is $C^1$ near the origin]
  \label{lem:dc-ch3-2-9-c1}
  \lean{Riemannian.Exponential.exists_contDiffOn_extChartAt_expMap_ball, Riemannian.FlowDependence.contDiffOn_one_of_forall_hasStrictFDerivAt, Riemannian.FlowDependence.continuousOn_of_forall_hasStrictFDerivAt}
  \uses{lem:dc-ch3-2-9-c1ball}
  \leanok
  \dcref{ch3:2.9}
  There is $\rho > 0$ such that the ball $B_\rho(0) \subset T_pM$ lies in the
  domain of $\exp_p$, its image under $\exp_p$ stays in the chart at $p$, and
  the chart reading $w \mapsto \varphi_p(\exp_p(w))$ is $C^1$ on $B_\rho(0)$.
  This upgrades the pointwise strict differentiability of
  \cref{lem:dc-ch3-2-9-c1ball} to genuine $C^1$ regularity, the differentiable
  half of the local-diffeomorphism claim of \cref{prop:dc-ch3-2-9} at the
  $C^1$ level.
\end{lemma}
\begin{proof}
  \sketch
  \leanok
  A map that is \emph{strictly} differentiable at every point of an open set
  automatically has a continuous derivative there: if $x, x'$ are close, the
  strict estimates
  $\|f(y) - f(z) - df_x(y - z)\| \le \varepsilon \|y - z\|$ at $x$ and the
  corresponding estimate at $x'$ both hold for all pairs $y, z$ in a common
  small ball around $x'$; applying them to the pairs $(x' + w,\, x')$ for all
  small $w$ and subtracting gives
  $\|(df_{x'} - df_x)(w)\| \le 2\varepsilon \|w\|$, whence
  $\|df_{x'} - df_x\| \le 2\varepsilon$ by homogeneity. Thus the derivative
  map is continuous on the open set, and on open sets continuity of the
  derivative together with differentiability is exactly the $C^1$
  characterization. Applying this to the chart reading of $\exp_p$ on
  $B_\rho(0)$, which is strictly differentiable at every point by
  \cref{lem:dc-ch3-2-9-c1ball}, gives the claim.
\end{proof}

\begin{lemma}[the derivative of $\exp_p$ is invertible near the origin]
  \label{lem:dc-ch3-2-9-invertible}
  \lean{Riemannian.Exponential.exists_hasStrictFDerivAt_equiv_extChartAt_expMap_ball}
  \uses{lem:dc-ch3-2-9-c1ball, lem:dc-ch3-2-9-strict, lem:dc-ch3-2-9-c1}
  \leanok
  \dcref{ch3:2.9}
  There is $\rho > 0$ such that the ball $B_\rho(0) \subset T_pM$ lies in the
  domain of $\exp_p$, its image under $\exp_p$ stays in the chart at $p$, and
  at every $v_0$ with $\|v_0\| < \rho$ the strict Fr\'echet derivative of the
  chart reading $w \mapsto \varphi_p(\exp_p(w))$ is a continuous linear
  isomorphism of $T_pM$.
\end{lemma}
\begin{proof}
  \sketch
  \leanok
  By \cref{lem:dc-ch3-2-9-c1ball} the chart reading $f$ of $\exp_p$ is strictly
  differentiable at every point of a ball $B_{\rho_0}(0)$, so its derivative
  map $x \mapsto df_x$ is continuous there (\cref{lem:dc-ch3-2-9-c1}), and
  $df_0 = \mathrm{id}$ (\cref{lem:dc-ch3-2-9-strict}). Choose
  $\rho \le \rho_0$ such that $\|df_{v_0} - \mathrm{id}\| < 1$ for
  $\|v_0\| < \rho$. Writing $df_{v_0} = \mathrm{id} - t$ with
  $t = \mathrm{id} - df_{v_0}$, $\|t\| < 1$, the Neumann series
  $\sum_{n \ge 0} t^n$ converges in the Banach algebra of bounded endomorphisms
  of $T_pM$ and inverts $\mathrm{id} - t$; hence $df_{v_0}$ is a continuous
  linear isomorphism.
\end{proof}

\begin{lemma}[$\exp_p$ is a $C^1$ diffeomorphism near the origin]
  \label{lem:dc-ch3-2-9-c1diffeo}
  \lean{Riemannian.Exponential.exists_c1_local_diffeomorphism_expMap}
  \uses{def:dc-ch3-expmap, lem:dc-ch3-2-9-c1, lem:dc-ch3-2-9-invertible, lem:dc-ch3-2-9-strict}
  \leanok
  \dcref{ch3:2.9}
  There is $\varepsilon > 0$ such that the ball
  $B_\varepsilon(0) \subset T_pM$ lies in the domain of $\exp_p$ and
  $\exp_p : B_\varepsilon(0) \to M$ is a $C^1$ diffeomorphism onto an open
  subset of $M$: it is injective on $B_\varepsilon(0)$, its image is open in
  $M$, its chart reading is $C^1$ on $B_\varepsilon(0)$, and it admits a left
  inverse whose chart reading is $C^1$ on the image. This is
  \cref{prop:dc-ch3-2-9} with ``diffeomorphism'' read at the $C^1$ regularity
  level.
\end{lemma}
\begin{proof}
  \sketch
  \leanok
  Let $f = \varphi_p \circ \exp_p$ be the chart reading of $\exp_p$, which is
  $C^1$ on a ball around the origin by \cref{lem:dc-ch3-2-9-c1}, with
  $df_0 = \mathrm{id}$ (\cref{lem:dc-ch3-2-9-strict}). By the inverse function
  theorem for strictly differentiable maps, $f$ admits a local inverse
  $f^{-1}$ near $0$; since $f$ is $C^1$ at $0$ with invertible derivative,
  $f^{-1}$ is $C^1$ on a neighbourhood $V$ of $f(0)$. Moreover $\exp_p$ is
  injective on a small ball (\cref{lem:dc-ch3-2-9-strict}), and at every point
  of the ball of \cref{lem:dc-ch3-2-9-invertible} the strict derivative of $f$
  is invertible, so the inverse function theorem applied at each point makes
  $f$ an open map there: $f(B_\varepsilon(0))$ is open, and the image
  $\exp_p(B_\varepsilon(0))$ — the intersection of the chart domain with the
  $\varphi_p$-preimage of $f(B_\varepsilon(0))$ — is open in $M$. Shrinking
  $\varepsilon$ so that $f(B_\varepsilon(0)) \subset V$ and so that the
  left-inverse identity $f^{-1}(f(w)) = w$ holds on $B_\varepsilon(0)$ gives
  all the claims.
\end{proof}

\begin{lemma}[$\exp_p$ is $C^2$ on a ball]
  \label{lem:dc-ch3-2-9-c2ball}
  \lean{Riemannian.Exponential.exists_contDiffOn_two_extChartAt_expMap_ball}
  \uses{def:dc-ch3-expmap, lem:dc-ch3-2-6, lem:dc-ch3-2-7-c2flow, lem:dc-ch3-2-9-c1}
  \leanok
  \dcref{ch3:2.9}
  There is $\rho > 0$ such that the ball $B_\rho(0) \subset T_pM$ lies in the
  domain of $\exp_p$, its image under $\exp_p$ stays in the chart at $p$, and
  the chart reading $w \mapsto \varphi_p(\exp_p(w))$ is $C^2$ on $B_\rho(0)$.
  This upgrades the $C^1$ regularity of \cref{lem:dc-ch3-2-9-c1} one order:
  the derivative of the chart reading at $w$ is the explicit
  continuous-linear image $\Phi(\tau(\iota(w)))$ of the variational flow of
  \cref{lem:dc-ch3-2-7-c2flow}, and is itself $C^1$ in $w$. It is the
  regularity input for the Gauss lemma (\cref{lem:dc-ch3-3-5}): the mixed
  second derivatives of the chart reading of $\exp_p$ exist and are symmetric
  on the ball.
\end{lemma}
\begin{proof}
  \sketch
  \leanok
  Let $Z$, $r$, $\varepsilon$, $T$, $\sigma$, $\tau$ be as in
  \cref{lem:dc-ch3-2-7-c2flow} and put $\rho = rT$. As in
  \cref{lem:dc-ch3-2-9-c1ball}, the homogeneity of geodesics realized on the
  flow witness by fibre scaling (\cref{lem:dc-ch3-2-6}) identifies, for
  $\|w\| < \rho$,
  $\varphi_p(\exp_p(w)) = \big(\sigma(\iota(w))(T)\big)_1$ where
  $\iota(w) = (\varphi_p(p), w/T)$ is affine into the open flow ball. By the
  chain rule and the computation of $D\sigma$ by the variational flow
  (\cref{lem:dc-ch3-2-7-c2flow}), the strict derivative of the chart reading
  at $w$ is
  $$
  v \;\longmapsto\; \Big(\tau(\iota(w))(T)\,\big(d\iota(v)\big)\Big)_1
  \;=\; \Phi\big(\tau(\iota(w))\big)\,v,
  $$
  where $\Phi$ — evaluate the operator curve at time $T$, compose with the
  fixed linear map $d\iota$, project to the base component — is a fixed
  bounded linear map from $C([0,T], L(\mathbb{R}^{2n}))$ to
  $L(\mathbb{R}^n)$. Since $\tau$ is strictly differentiable at every point
  of the flow ball (\cref{lem:dc-ch3-2-7-c2flow}) and $\Phi$, $\iota$ are
  (affine) continuous linear, the derivative map
  $w \mapsto \Phi(\tau(\iota(w)))$ is strictly differentiable, hence $C^1$,
  on $B_\rho(0)$ (strict differentiability at every point of an open set
  self-improves to $C^1$, as in \cref{lem:dc-ch3-2-9-c1}). A differentiable
  map whose derivative map is $C^1$ on an open set is $C^2$ there.
\end{proof}

\begin{lemma}[$\exp_p$ is a $C^2$ diffeomorphism near the origin]
  \label{lem:dc-ch3-2-9-c2diffeo}
  \lean{Riemannian.Exponential.exists_c2_local_diffeomorphism_expMap}
  \uses{lem:dc-ch3-2-9-c2ball, lem:dc-ch3-2-9-invertible, lem:dc-ch3-2-9-c1diffeo}
  \leanok
  \dcref{ch3:2.9}
  There is $\varepsilon > 0$ such that $\exp_p$ is injective on
  $B_\varepsilon(0) \subset T_pM$ with open image in $M$, the chart reading
  $f = \varphi_p \circ \exp_p$ is $C^2$ on $B_\varepsilon(0)$ with open image
  in the chart, and the local inverse $f^{-1}$ is $C^2$ on the open chart
  image $f(B_\varepsilon(0))$. This upgrades the inverse-regularity of
  \cref{lem:dc-ch3-2-9-c1diffeo} one order; it is the regularity of
  $\exp_p^{-1}$ consumed by the second-derivative computation of
  \cref{lem:dc-ch3-4-1}.
\end{lemma}
\begin{proof}
  \sketch
  \leanok
  \uses{lem:dc-ch3-2-9-c2ball, lem:dc-ch3-2-9-invertible, lem:dc-ch3-2-9-c1diffeo}
  Take $\varepsilon$ below the radii of \cref{lem:dc-ch3-2-9-c2ball} ($f$ is
  $C^2$ on the ball), \cref{lem:dc-ch3-2-9-invertible} ($df$ is invertible at
  every point of the ball), and of the injectivity ball of
  \cref{lem:dc-ch3-2-9-c1diffeo}. Injectivity of $f$ and openness of both
  images follow as in \cref{lem:dc-ch3-2-9-c1diffeo} (an everywhere-invertible
  derivative makes $f$ an open map on the ball). Define $f^{-1}$ globally on
  the image of the injective $f$. At each image point $f(v_0)$ the inverse
  function theorem for $C^2$ maps — applicable since $f$ is $C^2$ at $v_0$
  with invertible derivative — provides a local inverse which is $C^2$ at
  $f(v_0)$; it agrees with $f^{-1}$ near $f(v_0)$, because both are left
  inverses of $f$ on a neighborhood of $v_0$ and $f$ maps neighborhoods of
  $v_0$ onto neighborhoods of $f(v_0)$. Hence $f^{-1}$ is $C^2$ at every
  point of the open image.
\end{proof}

\begin{proposition}[exp_q is a local diffeomorphism]
  \label{prop:dc-ch3-2-9}
  \lean{Riemannian.Exponential.exists_infty_local_diffeomorphism_expMap}
  \uses{def:dc-ch3-expmap, lem:dc-ch3-expmap-ray, lem:dc-ch3-2-9-dexp, lem:dc-ch3-2-9-strict, lem:dc-ch3-2-9-c1ball, lem:dc-ch3-2-9-c1, lem:dc-ch3-2-9-invertible, lem:dc-ch3-2-9-c1diffeo, lem:dc-ch3-2-9-c2ball, lem:dc-ch3-2-9-c2diffeo, lem:dc-ch3-2-9-cinfty}
  \leanok
  \dcref{ch3:2.9}
  \notready
  Given $q\in M$, there exists $\varepsilon>0$ such that
  $\exp_q:B_\varepsilon(0)\subset T_qM\to M$ is a diffeomorphism of
  $B_\varepsilon(0)$ onto an open subset of $M$.
\end{proposition}
\begin{proof}
  \sketch
  Compute $d(\exp_q)_0$:
  
  $$
  d(\exp_q)_0(v)=\frac{d}{dt}\big(\exp_q(tv)\big)\Big|_{t=0}=\frac{d}{dt}\big(\gamma(1,q,tv)\big)\Big|_{t=0}=\frac{d}{dt}\big(\gamma(t,q,v)\big)\Big|_{t=0}=v.
  $$
  
  Hence $d(\exp_q)_0=\mathrm{Id}$ on $T_qM$, and by the inverse function theorem
  $\exp_q$ is a local diffeomorphism on a neighborhood of $0$.
\end{proof}

\begin{example}[Euclidean space]
  \label{ex:dc-ch3-2-10}
  \lean{Riemannian.euclidean_chartBasisVecFiber, Riemannian.euclidean_chartGramMatrix, Riemannian.euclidean_chartChristoffel, Riemannian.euclidean_chartChristoffelContraction, Riemannian.euclideanGeodesic, Riemannian.isGeodesic_euclideanGeodesic, Riemannian.expMapIntrinsic_euclidean}
  \leanok
  \dcref{ch3:2.10}
  \notready
  Let $M=\mathbb{R}^n$. The covariant derivative is the usual derivative, so the
  geodesics are straight lines parametrized proportionally to arc length; the
  exponential is the identity (under the usual identification $T_p\mathbb{R}^n\cong\mathbb{R}^n$).
\end{example}

\begin{example}[the sphere]
  \label{ex:dc-ch3-2-11}
  \uses{prop:dc-ch3-2-5}
  \notready
  \dcref{ch3:2.11}
  \notready
  Let $M=S^n\subset\mathbb{R}^{n+1}$ be the unit sphere. Great circles
  parametrized by arc length are geodesics. All
  geodesics of $S^n$ are great circles parametrized proportionally to arc length:
  given $p\in S^n$ and a unit $v\in T_pS^n$, the intersection of $S^n$ with the
  plane through the origin, $p$, and $v$ is a great circle parametrizable as the
  geodesic through $p$ with velocity $v$; by uniqueness of \cref{prop:dc-ch3-2-5} the claim
  follows. Here $\exp_p$ is defined on all of $T_p S^n$: it maps $B_\pi(0)$
  injectively onto $S^n-\{q\}$ ($q$ the antipode of $p$), collapses
  $\partial B_\pi(0)$ to $q$, maps the annulus $B_{2\pi}(0)-\overline{B_\pi(0)}$
  injectively onto $S^n-\{p,q\}$, and collapses $\partial B_{2\pi}(0)$ to $p$.
  For the manifold $S^n-\{q\}$, however, $\exp_p$ is defined only on
  $B_\pi(0)\subset T_p(S^n-\{q\})$.
\end{example}

\section{Minimizing properties of geodesics}

\begin{definition}[piecewise differentiable curve]
  \label{def:dc-ch3-3-1}
  \lean{Riemannian.Geodesic.IsPiecewiseDifferentiableCurve, Riemannian.Geodesic.IsPiecewiseDifferentiableCurveOfOrder, Riemannian.Geodesic.IsPiecewiseSmoothCurve, Riemannian.Geodesic.SubdividedCurveOfOrder, Riemannian.Geodesic.SubdividedCurveOfOrder.vertex}
  \dcref{ch3:3.1}
  A \emph{piecewise differentiable curve} is a continuous map $c:[a,b]\to M$ for which
  there is a partition $a=t_0<t_1<\dots<t_k=b$ such that each restriction
  $c|_{[t_i,t_{i+1}]}$ is differentiable; $c$ \emph{joins} $c(a)$ and $c(b)$. $c(t_i)$ is
  a \emph{vertex} of $c$, and the \emph{vertex angle at $c(t_i)$} is the angle between
  $\lim_{t\to t_i^+}c'(t)$ and $\lim_{t\to t_i^-}c'(t)$. Parallel transport extends
  to such curves by transporting successively over each differentiable segment.
\end{definition}

\begin{definition}[minimizing geodesic]
  \label{def:dc-ch3-3-2}
  \lean{Riemannian.Geodesic.IsMinimizingGeodesicSegment}
  \uses{def:dc-ch3-3-1}
  \leanok
  \dcref{ch3:3.2}
  A segment of a geodesic $\gamma:[a,b]\to M$ is \emph{minimizing} if
  $\ell(\gamma)\le\ell(c)$ for every piecewise differentiable curve $c$ joining
  $\gamma(a)$ to $\gamma(b)$.
\end{definition}

\begin{definition}[parametrized surface]
  \label{def:dc-ch3-3-3}
  \lean{Riemannian.Geodesic.IntrinsicParametrizedSurface,
    Riemannian.Geodesic.IntrinsicParametrizedSurface.velocityU,
    Riemannian.Geodesic.IntrinsicParametrizedSurface.velocityV,
    Riemannian.Geodesic.IsSmoothParametrizedSurfaceWithBoundary, Riemannian.Geodesic.IsNonPiAngle, Riemannian.Geodesic.IsNonPiVertexAngle, Riemannian.Geodesic.isNonPiVertexAngle_iff_vertexAngle_ne_pi, Riemannian.Geodesic.vertexAngle_eq_pi_of_reversed, Riemannian.Geodesic.isNonPiAngle_iff_angle_ne_pi, Riemannian.Geodesic.PiecewiseC1BoundaryParametrization.interior_vertex_angle_ne_pi, Riemannian.Geodesic.PiecewiseC1BoundaryParametrization.closing_vertex_angle_ne_pi, Riemannian.Geodesic.IsParametrizedSurfaceWithBoundary.toExtended, Riemannian.Geodesic.IsParametrizedSurfaceWithBoundary.surfaceDomain, Riemannian.Geodesic.IsParametrizedSurfaceWithBoundary.regularOn, Riemannian.Geodesic.IsParametrizedSurfaceWithBoundary.isParametrizedSurface, Riemannian.Geodesic.PiecewiseC1BoundaryParametrization.interior_vertex_nonPi, Riemannian.Geodesic.PiecewiseC1BoundaryParametrization.closing_vertex_nonPi}
  \leanok
  \dcref{ch3:3.3}
  Let $A$ be a connected set in $\mathbb{R}^2$, $U\subset A\subset\overline{U}$, $U$
  open, such that $\partial A$ is a piecewise differentiable curve with vertex
  angles different from $\pi$. A \emph{parametrized surface} in $M$ is a differentiable
  map $s:A\subset\mathbb{R}^2\to M$ (differentiable on $A$ meaning extendable to an
  open $U\supset A$; the vertex-angle condition ensures $ds$ is independent of the
  extension). A vector field $V$ along $s$ assigns to each $q\in A$ a vector
  $V(q)\in T_{s(q)}M$; it is differentiable if $q\mapsto V(q)f$ is differentiable
  for every differentiable function $f$ on $M$. For cartesian $(u,v)$ on
  $\mathbb{R}^2$, $\frac{\partial s}{\partial u}(u_0,v_0)$ is the velocity of
  $u\mapsto s(u,v_0)$ at $u_0$, and $\frac{\partial s}{\partial v}$ is analogous.
  If $V$ is a vector field along $s$, $\frac{DV}{\partial u}(u_0,v_0)$ is the
  covariant derivative along $u\mapsto s(u,v_0)$ of the restriction of $V$, and
  $\frac{DV}{\partial v}$ is defined analogously.
\end{definition}

\begin{lemma}[symmetry]
  \label{lem:dc-ch3-3-4}
  \lean{Riemannian.Geodesic.IntrinsicParametrizedSurface.mixedCovariantDerivatives_commute,
    Riemannian.Geodesic.ParametrizedSurface.covariantDerivativeUOfVelocityV_eq_covariantDerivativeVOfVelocityU,
    Riemannian.Geodesic.covariant_sndFDeriv_symm, Riemannian.Geodesic.covariant_sndFDeriv_symm_of_eventually, Riemannian.Geodesic.covariant_sndFDeriv_symm_of_surface}
  \leanok
  \dcref{ch3:3.4}
  If $M$ has a symmetric connection and $s:A\to M$ is a parametrized surface, then
  $$
  \frac{D}{\partial v}\frac{\partial s}{\partial u}=\frac{D}{\partial u}\frac{\partial s}{\partial v}.
  $$
\end{lemma}
\begin{proof}
  \sketch
  \leanok
  In a coordinate system $\mathbf{x}$ with
  $\mathbf{x}^{-1}\circ s(u,v)=(x^1(u,v),\dots,x^n(u,v))$,
  
  $$
  \frac{D}{\partial v}\Big(\frac{\partial s}{\partial u}\Big)=\sum_i\frac{\partial^2 x^i}{\partial v\,\partial u}\frac{\partial}{\partial x^i}+\sum_{i,j}\frac{\partial x^i}{\partial u}\frac{\partial x^j}{\partial v}\nabla_{\partial/\partial x^j}\frac{\partial}{\partial x^i}.
  $$
  
  By symmetry of the connection,
  $\nabla_{\partial/\partial x^j}\frac{\partial}{\partial x^i}=\nabla_{\partial/\partial x^i}\frac{\partial}{\partial x^j}$,
  so computing $\frac{D}{\partial u}\frac{\partial s}{\partial v}$ gives the same
  expression.
\end{proof}

\begin{lemma}[rays of the exponential chart reading are coordinate geodesics]
  \label{lem:dc-ch3-3-5-rayode}
  \lean{Riemannian.Exponential.exists_expMap_ray_ode_ball}
  \uses{def:dc-ch3-expmap, lem:dc-ch3-2-6, lem:dc-ch3-2-7-vuniform, lem:dc-ch3-2-7-c2flow, lem:dc-ch3-2-9-c2ball}
  \leanok
  \dcref{ch3:3.5}
  There are $\rho > 0$ and $b > 1$ such that, writing
  $f = \varphi_p \circ \exp_p$ for the chart reading of the exponential map:
  (i) for $\|u\| < \rho$ and $|a| < b$ the vector $a\,u$ lies in the domain of
  $\exp_p$ and $\exp_p(a\,u)$ stays in the chart at $p$;
  (ii) $f$ is $C^2$ on $B_\rho(0)$;
  (iii) $(df)_0 = \mathrm{id}$;
  (iv) for $\|u\| < \rho$, $|t| < b$, $\|t\,u\| < \rho$, the ray velocity
  $V(t) = (df)_{tu}(u)$ satisfies the chart-coordinate geodesic equation
  $$V'(t) = -\,\Gamma\big(V(t), V(t)\big)\big(f(tu)\big),$$
  so each ray $t \mapsto f(t\,u)$ is a coordinate geodesic.
\end{lemma}
\begin{proof}
  \sketch
  \leanok
  Let $Z$, $r$, $\varepsilon$, $T$ be the uniform local flow of the coordinate
  spray (\cref{lem:dc-ch3-2-7-c2flow}); put
  $\rho = \min(\rho_2, rT)$, with $\rho_2$ the $C^2$ radius of
  \cref{lem:dc-ch3-2-9-c2ball}, and $b = \varepsilon/T > 1$. For
  $\|u\| \le rT$ and $|a| < b$, the fibre- and time-rescaled trajectory
  $$\zeta(s) = \big(Z_1(aTs),\; aT\,Z_2(aTs)\big),$$
  computed along the flow line of the initial condition
  $(\varphi_p(p), u/T)$, again solves the spray ODE — the first component
  because $\dot Z_1 = Z_2$, the second because the Christoffel contraction is
  quadratic in the velocity — now with initial data $(\varphi_p(p), a\,u)$,
  on the interval $\{s : |aTs| < \varepsilon\} \supseteq [0,1]$. By the
  witness identification (\cref{lem:dc-ch3-2-7-vuniform}) the canonical
  geodesic through $(p, a u)$ is defined at time $1$ with
  $\varphi_p(\exp_p(a u)) = Z_1(aT)$; this realizes the homogeneity
  $\exp_p(a u) = \gamma(a, p, u)$ of \cref{lem:dc-ch3-2-6} at the level of
  the flow, proving (i). Clause (ii) is \cref{lem:dc-ch3-2-9-c2ball}.
  Differentiating the identification in $a$ and comparing with the chain
  rule at points of the $C^2$ ball gives $(df)_{tu}(u) = T\,Z_2(tT)$; at
  $t = 0$ this is $u$, which proves (iii) after extending by homogeneity of
  both sides, and differentiating once more in $t$ through the flow ODE
  $\dot Z_2 = -\Gamma(Z_2, Z_2)(Z_1)$ gives (iv).
\end{proof}

\begin{lemma}[Gauss surface computation]
  \label{lem:dc-ch3-3-5-surface}
  \lean{Riemannian.Exponential.gauss_surface_computation}
  \uses{def:dc-ch3-3-3, lem:dc-ch3-3-4, lem:dc-ch2-3-2-coord}
  \leanok
  \dcref{ch3:3.5}
  Let $\rho > 0$, $b > 1$, and let $f$ be a map of $B_\rho(0)$ into the chart
  target at $p$ which is $C^2$ on $B_\rho(0)$ with $f(0) = \varphi_p(p)$ and
  $(df)_0 = \mathrm{id}$, and whose ray velocities satisfy the
  chart-coordinate geodesic equation as in \cref{lem:dc-ch3-3-5-rayode}(iv).
  Then for $\|v\| < \rho$ and every $w$,
  $$\big\langle (df)_v(v),\, (df)_v(w) \big\rangle_{f(v)}
    = \langle v, w\rangle_{\varphi_p(p)},$$
  the inner products being the chart Gram inner products at the indicated
  base points: $f$ is a radial isometry.
\end{lemma}
\begin{proof}
  \sketch
  \leanok
  Consider the parametrized surface (\cref{def:dc-ch3-3-3})
  $c(t,s) = f\big(t\,(v + s w)\big)$ on the open parameter set
  $\Omega = \{(t,s) : \|t(v+sw)\| < \rho,\ \|v+sw\| < \rho,\ |t| < b\}$,
  which contains $[0,1] \times \{0\}$; its partial velocities are
  $\partial_t c = (df)_{t(v+sw)}(v+sw)$ and $\partial_s c = (df)_{t(v+sw)}(tw)$.

  \emph{Constant speed.} For fixed $s$ the curve $t \mapsto c(t,s)$ is a
  coordinate geodesic, so by metric compatibility
  (\cref{lem:dc-ch2-3-2-coord})
  $\partial_t \langle \partial_t c, \partial_t c\rangle
  = 2\,\langle \tfrac{D}{\partial t}\partial_t c, \partial_t c\rangle = 0$:
  the squared speed $\psi(t,s)$ is constant in $t$ on the (connected)
  $t$-sections of $\Omega$, and $\psi(0,s) = \langle v+sw, v+sw\rangle$
  because $(df)_0 = \mathrm{id}$.

  \emph{The mixed term.} Let
  $h(t) = \langle \partial_s c, \partial_t c\rangle(t,0)$. By compatibility
  along the $t$-direction and the geodesic equation
  $\tfrac{D}{\partial t}\partial_t c = 0$,
  $h'(t) = \langle \tfrac{D}{\partial t}\partial_s c, \partial_t c\rangle$;
  by the symmetry lemma (\cref{lem:dc-ch3-3-4}),
  $\tfrac{D}{\partial t}\partial_s c = \tfrac{D}{\partial s}\partial_t c$,
  and by compatibility along the $s$-direction
  $$2\,\big\langle \tfrac{D}{\partial s}\partial_t c, \partial_t c\big\rangle
    = \partial_s\, \psi(t, \cdot)\big|_{s=0}
    = \tfrac{d}{ds}\,\langle v+sw, v+sw\rangle\big|_{s=0}
    = 2\,\langle v, w\rangle.$$
  Hence $h'(t) = \langle v, w\rangle$ for every $t \in [0,1]$.

  \emph{Endpoints.} $h(0) = 0$ since
  $\partial_s c(0, \cdot) = (df)_0(0 \cdot w) = 0$, so
  $h(1) = \langle v, w\rangle$. As
  $h(1) = \langle (df)_v(w), (df)_v(v)\rangle_{f(v)}$, symmetry of the
  metric gives the claim.
\end{proof}

\begin{lemma}[Gauss lemma on a ball]
  \label{lem:dc-ch3-3-5-ball}
  \lean{Riemannian.Exponential.exists_gauss_lemma_ball}
  \uses{def:dc-ch3-expmap, lem:dc-ch3-3-5-rayode, lem:dc-ch3-3-5-surface}
  \leanok
  \dcref{ch3:3.5}
  There is $\rho > 0$ such that the ball $B_\rho(0) \subset T_pM$ lies in the
  domain of $\exp_p$, its image under $\exp_p$ stays in the chart at $p$, and
  for every $v$ with $\|v\| < \rho$ and every
  $w \in T_pM \approx T_v(T_pM)$, read in the chart at $p$,
  $$\big\langle (d\exp_p)_v(v),\, (d\exp_p)_v(w) \big\rangle_{\exp_p(v)}
    = \langle v, w\rangle_p.$$
  In particular $\exp_p$ preserves the radial component of the metric: the
  geodesic spheres $\exp_p(\partial B_r(0))$, $r < \rho$, are orthogonal to
  the radial geodesics.
\end{lemma}
\begin{proof}
  \sketch
  \leanok
  Apply the surface computation (\cref{lem:dc-ch3-3-5-surface}) to the chart
  reading $f = \varphi_p \circ \exp_p$, whose ball domain, $C^2$ regularity,
  value $f(0)=\varphi_p(p)$, derivative $(df)_0 = \mathrm{id}$, and
  ray-geodesic property are supplied by \cref{lem:dc-ch3-3-5-rayode}.

\end{proof}

\begin{lemma}[radial lower bound]
  \label{lem:dc-ch3-3-5-radial}
  \lean{Riemannian.Exponential.exists_gauss_radial_lower_bound_ball}
  \uses{lem:dc-ch3-3-5-ball}
  \leanok
  \dcref{ch3:3.5}
  On the Gauss ball of \cref{lem:dc-ch3-3-5-ball}: for every $v$ with
  $\|v\| < \rho$ and every $\xi \in T_pM$, read in the chart at $p$,
  $$\langle v, \xi\rangle_p^2
    \;\le\; \langle v, v\rangle_p \cdot
      \big\langle (d\exp_p)_v(\xi),\, (d\exp_p)_v(\xi)\big\rangle_{\exp_p(v)}.$$
  The exponential map does not shrink the radial component of any vector;
  this is the pointwise inequality behind the minimizing property of radial
  geodesics (\cref{prop:dc-ch3-3-6}).
\end{lemma}
\begin{proof}
  \sketch
  \leanok
  For $v = 0$ both sides vanish (the right-hand side is nonnegative, the
  metric being positive semidefinite). For $v \ne 0$ put
  $\lambda = \langle v,\xi\rangle_p / \langle v,v\rangle_p$ and decompose
  $\xi = \lambda v + \xi_N$ with $\langle v, \xi_N\rangle_p = 0$. Expanding by
  bilinearity and applying the Gauss identity (\cref{lem:dc-ch3-3-5-ball})
  twice,
  $$\big|(d\exp_p)_v(\xi)\big|^2
    = \lambda^2 \langle v,v\rangle_p
      + \big|(d\exp_p)_v(\xi_N)\big|^2
    \;\ge\; \lambda^2 \langle v,v\rangle_p
    = \frac{\langle v,\xi\rangle_p^2}{\langle v,v\rangle_p},$$
  the cross term $2\lambda\,\langle (d\exp_p)_v(v), (d\exp_p)_v(\xi_N)\rangle
  = 2\lambda\,\langle v, \xi_N\rangle_p$ vanishing by orthogonality.
  Multiplying through by $\langle v,v\rangle_p > 0$ gives the claim.

\end{proof}

\begin{lemma}[Gauss]
  \label{lem:dc-ch3-3-5}
  \uses{def:dc-ch3-expmap, lem:dc-ch3-3-5-ball}
  \notready
  \dcref{ch3:3.5}
  \notready
  We identify $T_pM$ with $T_v(T_pM)$ at $v\in T_pM$, writing $T_pM\approx T_v(T_pM)$.
  Let $p\in M$ and $v\in T_pM$ with $\exp_p v$ defined, and $w\in T_pM\approx T_v(T_pM)$.
  Then
  $$
  \langle(d\exp_p)_v(v),\,(d\exp_p)_v(w)\rangle=\langle v,w\rangle.\qquad\text{(2)}
  $$
\end{lemma}
\begin{proof}
  \sketch
  Write $w=w_T+w_N$, $w_T$ parallel to $v$, $w_N$ normal to $v$. Since
  $d\exp_p$ is linear and
  $\langle(d\exp_p)_v(v),(d\exp_p)_v(w_T)\rangle=\langle v,w_T\rangle$, it suffices
  to prove (2) for $w=w_N$ ($w_N\ne 0$). Pick a curve $v(s)$ in $T_pM$ with
  $v(0)=v$, $v'(0)=w_N$, $|v(s)|=\text{const}$. Since $\exp_p v$ is defined,
  choose $\varepsilon>0$ so that $\exp_p(tv(s))$ is defined for $0\le t\le 1$,
  $|s|<\varepsilon$, and consider the parametrized surface
  $f(t,s)=\exp_p tv(s)$ whose curves $t\mapsto f(t,s_0)$ are geodesics.
  Then $\langle\frac{\partial f}{\partial s},\frac{\partial f}{\partial t}\rangle(1,0)=\langle(d\exp_p)_v(w_N),(d\exp_p)_v(v)\rangle$.
  Since $\frac{\partial f}{\partial t}$ is the tangent of a geodesic and by symmetry
  of the connection,
  
  $$
  \frac{\partial}{\partial t}\Big\langle\frac{\partial f}{\partial s},\frac{\partial f}{\partial t}\Big\rangle=\Big\langle\frac{D}{\partial t}\frac{\partial f}{\partial s},\frac{\partial f}{\partial t}\Big\rangle=\Big\langle\frac{D}{\partial s}\frac{\partial f}{\partial t},\frac{\partial f}{\partial t}\Big\rangle=\frac12\frac{\partial}{\partial s}\Big\langle\frac{\partial f}{\partial t},\frac{\partial f}{\partial t}\Big\rangle=0,
  $$
  
  so $\langle\frac{\partial f}{\partial s},\frac{\partial f}{\partial t}\rangle$ is
  independent of $t$. As $\lim_{t\to0}\frac{\partial f}{\partial s}(t,0)=\lim_{t\to0}(d\exp_p)_{tv}\,t w_N=0$,
  the inner product is $0$, proving (2).
\end{proof}

\subsection{Normal neighborhoods}
If $\exp_p$ is a diffeomorphism of a neighborhood $V\ni 0$ in $T_pM$ onto
$U=\exp_p(V)$, then $U$ is a \emph{normal neighborhood} of $p$. If
$\overline{B_\varepsilon(0)}\subset V$, then $\exp_p B_\varepsilon(0)=B_\varepsilon(p)$
is the \emph{normal ball} (or \emph{geodesic ball}) of center $p$ and radius $\varepsilon$.
By the Gauss lemma (3.5), its boundary is a hypersurface orthogonal to the
geodesics from $p$, the \emph{geodesic sphere} $S_\varepsilon(p)$; the geodesics in
$B_\varepsilon(p)$ that begin at $p$ are the \emph{radial geodesics}.

\begin{lemma}[smoothed-radius comparison]
  \label{lem:dc-ch3-3-6-radius}
  \lean{Riemannian.Exponential.gauss_radius_comparison, Riemannian.Exponential.exists_gauss_radius_comparison_ball, Riemannian.hasDerivAt_chartMetricInner_const_base, Riemannian.continuousOn_chartMetricInner_along, Riemannian.chartMetricInner_self_nonneg_of_mem_target}
  \uses{lem:dc-ch3-3-5-radial}
  \leanok
  \dcref{ch3:3.6}
  On the Gauss ball $B_\rho(0)\subset T_pM$ of \cref{lem:dc-ch3-3-5-radial}: for
  every differentiable curve $w:[a,b]\to B_\rho(0)$ with continuous derivative,
  write $c=\exp_p\circ w$ for the composed curve (read in the chart at $p$) and
  $r(t)=\sqrt{\langle w(t),w(t)\rangle_p}$ for the radius of the polar lift.
  Then
  $$r(b)-r(a)\;\le\;\int_a^b\big|\dot c(t)\big|_{c(t)}\,dt,$$
  where $\dot c(t)=(d\exp_p)_{w(t)}\,w'(t)$ by the chain rule and the norm is
  taken in the metric at the moving foot.
\end{lemma}
\begin{proof}
  \sketch
  \leanok
  For $\delta>0$ consider the smoothed radius
  $r_\delta(t)=\sqrt{\langle w(t),w(t)\rangle_p+\delta}$. The map
  $t\mapsto\langle w(t),w(t)\rangle_p$ is differentiable with derivative
  $2\langle w(t),w'(t)\rangle_p$ — the base point of the pairing is fixed, so
  this is the bilinear product rule with no Christoffel correction — hence
  $r_\delta$ is differentiable everywhere (the smoothing keeps the square root
  away from $0$) with
  $$r_\delta'(t)=\frac{\langle w(t),w'(t)\rangle_p}
    {\sqrt{\langle w(t),w(t)\rangle_p+\delta}}.$$
  The radial lower bound (\cref{lem:dc-ch3-3-5-radial}) at the pair
  $(w(t),w'(t))$ gives
  $$\langle w,w'\rangle_p^2\;\le\;\langle w,w\rangle_p\,\big|\dot c\big|^2
    \;\le\;\big(\langle w,w\rangle_p+\delta\big)\,\big|\dot c\big|^2,$$
  the second step because $|\dot c|^2\ge 0$ (the metric is positive
  semidefinite at every point of the chart image). Taking square roots and
  dividing, $r_\delta'(t)\le|\dot c(t)|$ for \emph{every} $t\in[a,b]$ —
  including the zeros of $w$, where the radius $r$ has a corner but
  $r_\delta$ does not. All the data are continuous ($w$, $w'$, the Gram
  entries along $c$, and $(d\exp_p)$ on the ball by $C^1$ regularity,
  \cref{lem:dc-ch3-2-9-c1ball}), so by the fundamental theorem of calculus
  and monotonicity of the integral
  $$r_\delta(b)-r_\delta(a)=\int_a^b r_\delta'(t)\,dt
    \le\int_a^b\big|\dot c(t)\big|\,dt.$$
  Letting $\delta\to0^+$ gives the claim.
\end{proof}

\begin{lemma}[reach estimate: length dominates the radius reached]
  \label{lem:dc-ch3-3-6-reach}
  \lean{Riemannian.Exponential.gauss_radius_reach, Riemannian.Exponential.exists_gauss_radius_reach_ball}
  \uses{lem:dc-ch3-3-6-radius}
  \leanok
  \dcref{ch3:3.6}
  On the Gauss ball: a curve $c=\exp_p\circ w$ starting at $p$ (i.e.\
  $w(a)=0$) has chart-read length at least the $g_p$-radius its polar lift
  reaches at \emph{any} time:
  $$\sqrt{\langle w(t_1),w(t_1)\rangle_p}\;\le\;\int_a^b|\dot c(t)|\,dt
    \qquad(t_1\in[a,b]).$$
  This is the inequality behind the escape case in
  \cref{prop:dc-ch3-3-6}: a competitor leaving the normal ball of
  $g_p$-radius $\rho'$ reaches radius $\rho'$ at its first exit, so its
  length is at least $\rho'>\ell(\gamma)$.
\end{lemma}
\begin{proof}
  \sketch
  \leanok
  Apply the smoothed-radius comparison (\cref{lem:dc-ch3-3-6-radius}) on
  $[a,t_1]$ and enlarge the integration interval: the speed integrand is
  pointwise nonnegative, so $\int_a^{t_1}\le\int_a^b$. With $w(a)=0$ the
  radius gain is the radius reached.
\end{proof}

\begin{lemma}[radial geodesics have unit chart speed]
  \label{lem:dc-ch3-3-6-rayspeed}
  \lean{Riemannian.Exponential.exists_expMap_ray_speed_ball}
  \uses{lem:dc-ch3-3-5-ball, lem:dc-ch3-3-5-rayode, def:dc-ch3-expmap}
  \leanok
  \dcref{ch3:3.6}
  There is $\rho>0$ such that for every $v$ with $\|v\|<\rho$ and every
  $t\in[0,1]$, read in the chart at $p$,
  $$\big\langle(d\exp_p)_{tv}(v),\,(d\exp_p)_{tv}(v)\big\rangle_{\exp_p(tv)}
    =\langle v,v\rangle_p.$$
  The radial geodesic $\gamma(t)=\exp_p(tv)$, $0\le t\le1$, therefore has
  constant speed $|v|_p$ and length
  $\ell(\gamma)=\int_0^1|v|_p\,dt=|v|_p$ (the normalized identity
  $|\partial f/\partial r|=1$).
\end{lemma}
\begin{proof}
  \sketch
  \leanok
  For $t\ne0$ apply the Gauss identity (\cref{lem:dc-ch3-3-5-ball}) at the
  pair $(tv,v)$:
  $$\big\langle(d\exp_p)_{tv}(tv),\,(d\exp_p)_{tv}(v)\big\rangle_{\exp_p(tv)}
    =\langle tv,v\rangle_p.$$
  Pulling the scalar $t$ out of both sides — by linearity of $(d\exp_p)_{tv}$
  on the left, by bilinearity of the metric on the right — and cancelling
  $t\ne0$ gives the claim. At $t=0$ the ray package
  (\cref{lem:dc-ch3-3-5-rayode}) supplies $(d\exp_p)_0=\mathrm{id}$ and
  $\exp_p(0)=p$, so the left-hand side is $\langle v,v\rangle_p$ directly.

\end{proof}

\begin{lemma}[radial geodesics minimize, polar form]
  \label{lem:dc-ch3-3-6-ball}
  \lean{Riemannian.Exponential.exists_minimizing_geodesic_ball}
  \uses{lem:dc-ch3-3-6-radius, lem:dc-ch3-3-6-rayspeed}
  \leanok
  \dcref{ch3:3.6}
  There is $\rho>0$ (a Gauss ball for $\exp_p$) such that for every $v$ with
  $\|v\|<\rho$ and every competing curve given in polar form through the
  exponential chart — a differentiable $w:[0,1]\to B_\rho(0)\subset T_pM$ with
  continuous derivative, $w(0)=0$, $w(1)=v$ — the chart-read length of the
  radial geodesic $\gamma(t)=\exp_p(tv)$ is at most that of
  $c(t)=\exp_p(w(t))$:
  $$\ell(\gamma)=\int_0^1\sqrt{\langle\dot\gamma,\dot\gamma\rangle}\,dt
    \;\le\;\int_0^1\sqrt{\langle\dot c,\dot c\rangle}\,dt=\ell(c).$$
\end{lemma}
\begin{proof}
  \sketch
  \leanok
  By \cref{lem:dc-ch3-3-6-rayspeed}, $\ell(\gamma)=|v|_p$. By the radius
  comparison (\cref{lem:dc-ch3-3-6-radius}),
  $$\ell(c)\;\ge\;r(1)-r(0)=|w(1)|_p-|w(0)|_p=|v|_p-0=\ell(\gamma).
  $$
\end{proof}

\begin{lemma}[radial geodesics minimize in a normal ball]
  \label{lem:dc-ch3-3-6-normal}
  \lean{Riemannian.Exponential.exists_minimizing_geodesic_normal_ball}
  \uses{lem:dc-ch3-3-6-ball, lem:dc-ch3-2-9-c1diffeo}
  \leanok
  \dcref{ch3:3.6}
  There is $\varepsilon>0$ such that $\exp_p$ is injective on
  $B_\varepsilon(0)\subset T_pM$ with open image — so
  $B=\exp_p(B_\varepsilon(0))$ is a normal ball — and for every $v$ with
  $\|v\|<\varepsilon$ and every curve $c:[0,1]\to B$ with $c(0)=p$,
  $c(1)=\exp_p(v)$, whose chart reading $t\mapsto\varphi_p(c(t))$ is
  differentiable with continuous derivative, the radial geodesic
  $\gamma(t)=\exp_p(tv)$ satisfies $\ell(\gamma)\le\ell(c)$ (chart-read
  lengths).
\end{lemma}
\begin{proof}
  \sketch
  \leanok
  Put the competing curve in polar form: $w=\exp_p^{-1}\circ c$ through the
  local $C^1$ inverse of $\exp_p$ (\cref{lem:dc-ch3-2-9-c1diffeo}), so
  $c=\exp_p\circ w$ with $w$ differentiable into $B_\varepsilon(0)$,
  $w(0)=0$, $w(1)=v$. The image $\varphi_p(\exp_p(B_\varepsilon(0)))$ is open
  — at every point of the ball the derivative of the chart reading of
  $\exp_p$ is a linear equivalence (\cref{lem:dc-ch3-2-9-invertible}), so the
  map is open — hence the chain rule applies to $w=\exp_p^{-1}\circ c$ and
  $w'$ is continuous. The polar minimizing property
  (\cref{lem:dc-ch3-3-6-ball}) gives $\ell(\gamma)\le\ell(\exp_p\circ w)$,
  and $\ell(\exp_p\circ w)=\ell(c)$: the two length integrands agree, the
  feet by $\exp_p(w(t))=c(t)$ and the velocities
  $(d\exp_p)_{w(t)}w'(t)=\frac{d}{dt}\varphi_p(c(t))$ by uniqueness of
  one-sided derivatives on $[0,1]$ (both compute the derivative of
  $u=\varphi_p\circ c=f\circ w$ within $[0,1]$).
\end{proof}

\begin{lemma}[radial geodesics minimize among piecewise competitors in a normal ball]
  \label{lem:dc-ch3-3-6-piecewise}
  \lean{Riemannian.Exponential.gauss_radius_comparison_piecewise, Riemannian.Exponential.gauss_radius_reach_piecewise, Riemannian.Exponential.exists_minimizing_geodesic_normal_ball_piecewise}
  \uses{def:dc-ch3-3-1, lem:dc-ch3-3-6-normal}
  \leanok
  \dcref{ch3:3.6}
  There is $\varepsilon>0$ such that $\exp_p$ is injective on
  $B_\varepsilon(0)\subset T_pM$ with open image — so
  $B=\exp_p(B_\varepsilon(0))$ is a normal ball — and for every $v$ with
  $\|v\|<\varepsilon$, every partition $0=\tau_0\le\dots\le\tau_n=1$, and
  every continuous curve $c:[0,1]\to B$ from $p$ to $\exp_p(v)$ whose chart
  reading is differentiable on each piece $[\tau_i,\tau_{i+1}]$ in the
  one-sided sense with continuous piece derivatives (a piecewise
  differentiable curve, \cref{def:dc-ch3-3-1}), the chart-read length of the
  radial geodesic $\gamma(t)=\exp_p(tv)$ is at most the total length of $c$:
  $$\ell(\gamma)\;\le\;\sum_{i<n}\int_{\tau_i}^{\tau_{i+1}}
    \big|\dot c(t)\big|_{c(t)}\,dt=\ell(c).$$
  Moreover the underlying radius comparison telescopes: for a continuous,
  piecewise differentiable polar lift $w$ with values in the Gauss ball, both
  the radius gain $|w(\tau_n)|_p-|w(\tau_0)|_p$ and — when $w(\tau_0)=0$ —
  the radius $|w(t_1)|_p$ reached at \emph{any} time $t_1$ are dominated by
  the total length of $\exp_p\circ w$.
\end{lemma}
\begin{proof}
  \sketch
  \leanok
  \uses{lem:dc-ch3-3-6-radius, lem:dc-ch3-3-6-rayspeed, lem:dc-ch3-2-9-c1diffeo, lem:dc-ch3-2-9-invertible}
  Each piece $[\tau_i,\tau_{i+1}]$ satisfies the single-piece smoothed-radius
  comparison (\cref{lem:dc-ch3-3-6-radius}): the smoothed radius
  $r_\delta=\sqrt{\langle w,w\rangle_p+\delta}$ is differentiable on the
  closed piece using only the one-sided derivatives of $w$ at its endpoints —
  exactly what a corner of a piecewise curve provides — and the fundamental
  theorem of calculus on the closed piece gives
  $r_\delta(\tau_{i+1})-r_\delta(\tau_i)\le\int_{\tau_i}^{\tau_{i+1}}|\dot c|$,
  hence the same bound for $r$ in the limit $\delta\to0^+$. Summing over the
  pieces telescopes the radius gains:
  $$|w(\tau_n)|_p-|w(\tau_0)|_p
    =\sum_{i<n}\big(|w(\tau_{i+1})|_p-|w(\tau_i)|_p\big)
    \le\sum_{i<n}\int_{\tau_i}^{\tau_{i+1}}|\dot c|\,dt.$$
  For the reach form at a time $t_1$ in the piece $[\tau_k,\tau_{k+1}]$, apply
  the telescoped comparison to the partition of $[\tau_0,t_1]$ obtained by
  cutting at $t_1$, then enlarge: the dropped terms are integrals of a
  nonnegative integrand. The minimizing statement follows as in
  \cref{lem:dc-ch3-3-6-normal}: the competitor is put in polar form
  $w=\exp_p^{-1}\circ c$ through the local $C^1$ inverse of $\exp_p$
  (\cref{lem:dc-ch3-2-9-c1diffeo}), which is differentiable piece by piece
  with continuous one-sided derivatives by the chain rule
  (\cref{lem:dc-ch3-2-9-invertible}), and the telescoped comparison with
  $w(0)=0$, $w(1)=v$ bounds the total length of $c$ below by
  $|v|_p=\ell(\gamma)$ (\cref{lem:dc-ch3-3-6-rayspeed}).
\end{proof}

\begin{lemma}[the competitor bound: piecewise curves, escape case included]
  \label{lem:dc-ch3-3-6-competitor}
  \lean{Riemannian.Exponential.pathELength_sum_partition, Riemannian.Exponential.exists_le_pathELength, Riemannian.Exponential.exists_le_pathELength_piecewise}
  \uses{def:dc-ch3-3-1, lem:dc-ch3-3-6-piecewise}
  \leanok
  \dcref{ch3:3.6}
  Let $p\in M$. There are $\varepsilon>0$ and a constant $C>0$ (comparing the
  coordinate norm at $p$ with the $g_p$-norm) such that $\exp_p$ is injective
  on $B_\varepsilon(0)\subset T_pM$, open on sub-balls, and — measuring the
  length of a piecewise differentiable curve $\sigma:[0,1]\to M$
  intrinsically, as the sum of the lengths of its differentiable pieces:
  \begin{enumerate}
  \item every piecewise differentiable curve $\sigma$ from $p$ to
    $\exp_p(v)$ with $\|v\|<\varepsilon$ satisfies
    $\ell(\sigma)\ge|v|_p=\ell(\gamma)$, where $\gamma(t)=\exp_p(tv)$ is the
    radial geodesic — the inequality, with competitors allowed to
    leave the normal ball;
  \item every piecewise differentiable curve from $p$ ending outside
    $\exp_p(B_r(0))$, $0<r\le\varepsilon$, has length at least $r/C$ (the
    escape estimate).
  \end{enumerate}
\end{lemma}
\begin{proof}
  \sketch
  \leanok
  \uses{lem:dc-ch3-3-6-piecewise, lem:dc-ch3-2-9-c1diffeo}
  Fix a radius $r'$ so small that on the coordinate ball $B_{r'}(0)$ the
  polar machinery of \cref{lem:dc-ch3-3-6-piecewise} applies and the
  coordinate norm is comparable to the $g_p$-norm, $\|z\|^2\le C^2\,
  \langle z,z\rangle_p$; write $U=\exp_p(B_{r'}(0))$, an open neighborhood of
  $p$. A competitor $\sigma$ either stays in $U$ or leaves it.

  If $\sigma$ stays in $U$, its polar lift $w=\exp_p^{-1}\circ\sigma$ exists
  on all of $[0,1]$, is continuous, and is piecewise differentiable with
  continuous one-sided piece derivatives (\cref{lem:dc-ch3-2-9-c1diffeo});
  the telescoped radius comparison (\cref{lem:dc-ch3-3-6-piecewise}) with
  $w(0)=0$ bounds $\ell(\sigma)$ below by the $g_p$-radius of the lifted
  endpoint — which is $|v|_p$ in case (1), and at least $r/C$ in case (2)
  since the endpoint has coordinate norm $\ge r$.

  If $\sigma$ leaves $U$, the image of the \emph{closed} coordinate ball
  $\overline{B_{r'}(0)}$ under $\exp_p$ is compact, hence closed; the set of
  times at which $\sigma$ lies outside $U$ is closed in $[0,1]$ and avoids
  $t=0$, so it has a smallest element $T>0$, and on $[0,T)$ the curve stays
  in $U$. By continuity $\sigma(T)$ lies in the compact image, and it is not
  in the open image $U$, so its polar coordinate has norm exactly $r'$: the
  curve first exits through the coordinate $r'$-sphere. Truncate the
  partition at $T$ — each piece $[\tau_i,\tau_{i+1}]$ becomes
  $[\min(\tau_i,T),\min(\tau_{i+1},T)]$, which either sits inside the
  original piece or degenerates to the point $T$ — and apply the reach form
  of the telescoped comparison on $[0,T]$: the length accumulated up to $T$
  is at least the $g_p$-radius of the exit point, which the norm comparison
  bounds below by $r'/C$. Choosing $\varepsilon\le r'$ so small that every
  $v$ with $\|v\|<\varepsilon$ has $|v|_p<r'/C$, both cases give (1); (2) is
  the same estimate with $r\le\varepsilon\le r'$.

  Finally, the intrinsic total length of $\sigma$ — the sum of the piece
  lengths — is the manifold path length of $\sigma$ over $[0,1]$, by
  additivity of the length integral over a monotone partition; on each piece
  the intrinsic speed agrees with the chart-read speed used by the polar
  comparison.
\end{proof}

\begin{lemma}[equality forces radial velocity]
  \label{lem:dc-ch3-3-6-equality-radial}
  \lean{Riemannian.Exponential.radial_equality_forcing, Riemannian.Exponential.gauss_radius_equality_ftc, Riemannian.Exponential.exists_gauss_equality_radial_ball}
  \uses{lem:dc-ch3-3-5-ball, lem:dc-ch3-3-6-radius}
  \leanok
  \dcref{ch3:3.6}
  There is $\rho>0$ (a Gauss ball for $\exp_p$) such that for every
  differentiable curve $w:[a,b]\to B_\rho(0)\subset T_pM$ with continuous
  derivative whose composed curve $c=\exp_p\circ w$ realizes the radius
  comparison with \emph{equality},
  $$|w(b)|_p-|w(a)|_p=\int_a^b|\dot c(t)|\,dt,$$
  the following hold at every $t\in(a,b)$ with $w(t)\neq0$: the radius
  $r=|w|_p$ is differentiable at $t$ with $r'(t)=|\dot c(t)|\ge0$, and the
  velocity is radial and outward,
  $$w'(t)=\frac{\langle w(t),w'(t)\rangle_p}{\langle w(t),w(t)\rangle_p}\,
    w(t),\qquad\langle w(t),w'(t)\rangle_p\ge0.$$
\end{lemma}
\begin{proof}
  \sketch
  \leanok
  \uses{lem:dc-ch3-3-5-ball, lem:dc-ch3-3-6-radius, lem:dc-ch3-2-9-invertible}
  The radius comparison holds on every subinterval $[c,d]\subseteq[a,b]$
  (\cref{lem:dc-ch3-3-6-radius}), so the deficit
  $\varphi(t)=\int_a^t|\dot c|-\big(r(t)-r(a)\big)$ is nondecreasing;
  vanishing at both ends, it vanishes identically, giving the FTC form
  $r(t)=r(a)+\int_a^t|\dot c|$ for \emph{all} $t$. The speed $|\dot c|$ is
  continuous, so $r$ is differentiable on $(a,b)$ with $r'=|\dot c|\ge0$.
  Where $w(t)\neq0$, the radius is also differentiable by the product rule
  with $r'=\langle w,w'\rangle_p/r$; identifying the two derivatives gives
  $\langle w,w'\rangle_p=r\,|\dot c|\ge0$, i.e.\ \emph{equality} in the
  radial lower bound
  $\langle w,w'\rangle_p^2\le\langle w,w\rangle_p\,|\dot c|^2$.

  Equality there forces the velocity to be radial. Decompose
  $w'=\lambda w+w_N'$ with
  $\lambda=\langle w,w'\rangle_p/\langle w,w\rangle_p$ and
  $\langle w,w_N'\rangle_p=0$. Expanding
  $|\dot c|^2=\big|(d\exp_p)_w(w')\big|^2$ by bilinearity and applying the
  Gauss identity (\cref{lem:dc-ch3-3-5-ball}) twice kills the cross terms:
  $$\big|(d\exp_p)_w(w')\big|^2
    =\lambda^2\langle w,w\rangle_p+\big|(d\exp_p)_w(w_N')\big|^2.$$
  Equality in the radial bound then forces
  $\big|(d\exp_p)_w(w_N')\big|^2=0$; the metric is positive definite at
  every point of a small chart ball around $p$ (the Gram comparison
  $\|z\|^2\le c\,\langle z,z\rangle_y$ holds for all $y$ near the pole), so
  $(d\exp_p)_w(w_N')=0$, and $(d\exp_p)_w$ is invertible on the ball
  (\cref{lem:dc-ch3-2-9-invertible}), so $w_N'=0$.
\end{proof}

\begin{lemma}[equality case in polar form: the competitor runs along the ray]
  \label{lem:dc-ch3-3-6-equality-ray}
  \lean{Riemannian.Exponential.exists_gauss_equality_ray_ball}
  \uses{lem:dc-ch3-3-6-equality-radial}
  \leanok
  \dcref{ch3:3.6}
  There is $\rho>0$ (a Gauss ball for $\exp_p$) such that every
  differentiable polar competitor $w:[0,1]\to B_\rho(0)\subset T_pM$ with
  continuous derivative and $w(0)=0$ that realizes the radius comparison with
  \emph{equality}, $|w(1)|_p=\int_0^1|\dot c|\,dt$ where $c=\exp_p\circ w$,
  satisfies, for every $t\in[0,1]$,
  $$w(t)=\frac{|w(t)|_p}{|w(1)|_p}\,w(1):$$
  the competitor traces the radial segment from $0$ to $w(1)$, with monotone
  radius $|w(t)|_p$.
\end{lemma}
\begin{proof}
  \sketch
  \leanok
  \uses{lem:dc-ch3-3-6-equality-radial}
  If $|w(1)|_p=0$ the monotone radius (its derivative is the nonnegative
  speed, by the FTC form of \cref{lem:dc-ch3-3-6-equality-radial}) is
  squeezed to $0$, so $w\equiv0$ by positive definiteness of
  $\langle\cdot,\cdot\rangle_p$ and both sides vanish. Otherwise let
  $t_*=\sup\{t:|w(t)|_p=0\}$, attained since the radius is continuous; then
  $t_*<1$ and $r=|w|_p$ is positive on $(t_*,1]$. On $(t_*,1)$ the
  normalized lift $\eta=w/r$ is differentiable with
  $$\eta'=\frac{w'}{r}-\frac{r'}{r^2}\,w=0,$$
  combining $r'=\langle w,w'\rangle_p/r$ (product rule) with the radial
  velocity $w'=(\langle w,w'\rangle_p/\langle w,w\rangle_p)\,w$
  (\cref{lem:dc-ch3-3-6-equality-radial}); a continuous function with
  vanishing right derivative on $[t,1]$ is constant, so $\eta\equiv\eta(1)$
  on $(t_*,1]$, i.e.\ $w(t)=r(t)\,w(1)/r(1)$ there. On $[0,t_*]$ the
  monotone radius satisfies $0\le r(t)\le r(t_*)=0$, so $w(t)=0$ and the
  identity holds with coefficient $0$.
\end{proof}

\begin{lemma}[confinement: a short curve from $p$ stays in the Gauss ball]
  \label{lem:dc-ch3-3-6-confinement}
  \lean{Riemannian.Exponential.exists_forall_mem_expMap_ball_of_pathELength_lt}
  \uses{lem:dc-ch3-3-6-competitor}
  \leanok
  \dcref{ch3:3.6}
  For every target radius $\rho'>0$ there is a length threshold $\delta>0$
  such that every $C^1$ curve $\sigma:[0,1]\to M$ starting at $p$ with
  $\ell(\sigma)<\delta$ satisfies $\sigma(t)\in\exp_p(B_{\rho'}(0))$ for all
  $t\in[0,1]$. This discharges the escape case in the equality analysis of
  \cref{prop:dc-ch3-3-6}: an equality competitor is automatically confined
  to the Gauss ball.
\end{lemma}
\begin{proof}
  \sketch
  \leanok
  \uses{lem:dc-ch3-3-6-competitor}
  Truncate $\sigma$ at time $t$ and reparametrize the truncation to $[0,1]$
  by the monotone affine map $\tau\mapsto t\tau$, which leaves the length
  unchanged. If the truncated endpoint $\sigma(t)$ escaped the sublevel set
  $\exp_p(\{z:\ \|z\|<\varepsilon,\ |z|_p<\delta\})$, the escape estimate of
  \cref{lem:dc-ch3-3-6-competitor} would force
  $\delta\le\ell(\sigma|_{[0,t]})\le\ell(\sigma)<\delta$, a contradiction.
  The Gram comparison $\|z\|^2\le c\,\langle z,z\rangle_p$ places the
  sublevel set inside the norm ball of radius $\sqrt{c}\,\delta$, which is
  inside $B_{\rho'}(0)$ for $\delta=\rho''/(\sqrt{c}+1)$,
  $\rho''=\min(\rho',\varepsilon)$.
\end{proof}

\begin{lemma}[the equality case in the manifold: monotone reparametrization of the radial geodesic]
  \label{lem:dc-ch3-3-6-equality-manifold}
  \lean{Riemannian.Exponential.exists_gauss_equality_manifold_ball, Riemannian.Exponential.exists_gauss_equality_manifold}
  \uses{lem:dc-ch3-3-6-equality-ray, lem:dc-ch3-3-6-equality-radial, lem:dc-ch3-2-9-c1diffeo, lem:dc-ch7-2-8-chart-length, lem:dc-ch3-3-6-confinement}
  \leanok
  \dcref{ch3:3.6}
  There is $\rho>0$ (a Gauss ball for $\exp_p$, with $\exp_p$ injective on
  $B_\rho(0)$) such that every $C^1$ competitor $\sigma:[0,1]\to M$ from $p$
  to $\exp_p(v)$, $|v|<\rho$, that stays in $\exp_p(B_\rho(0))$ and realizes
  the radial length $\ell(\sigma)=|v|_p$ is a \emph{monotone reparametrization
  of the radial geodesic}: there is a continuous monotone
  $s:[0,1]\to[0,1]$ with $s(0)=0$, $s(1)=1$,
  $$\sigma(t)=\exp_p\big(s(t)\,v\big)\quad\text{and}\quad
    \ell\big(\sigma|_{[0,t]}\big)=s(t)\,|v|_p\quad\text{for all }t\in[0,1];$$
  in particular $\sigma([0,1])=\exp_p([0,1]\cdot v)$. By the confinement lemma
  (\cref{lem:dc-ch3-3-6-confinement}) the stay-in-ball hypothesis is
  automatic: the statement holds for \emph{every} $C^1$ equality competitor
  from $p$ to $\exp_p(v)$, with no assumption on where it travels.
\end{lemma}
\begin{proof}
  \sketch
  \leanok
  \uses{lem:dc-ch3-3-6-equality-ray, lem:dc-ch3-3-6-equality-radial, lem:dc-ch3-2-9-c1diffeo, lem:dc-ch7-2-8-chart-length, lem:dc-ch3-3-6-confinement}
  Lift $\sigma$ through the $C^1$ local inverse of $\exp_p$
  (\cref{lem:dc-ch3-2-9-c1diffeo}): the polar lift
  $w=\exp_p^{-1}\circ\sigma$ is $C^1$ on $[0,1]$ with $w(0)=0$, $w(1)=v$, and
  values in $B_\rho(0)$. Reading the length in the chart at $p$
  (\cref{lem:dc-ch7-2-8-chart-length}) and computing the chart velocity of
  $\sigma$ by the chain rule through $w$, the hypothesis
  $\ell(\sigma)=|v|_p$ becomes the FTC-form equality of the polar lemmas.
  The ray lemma (\cref{lem:dc-ch3-3-6-equality-ray}) then forces
  $w(t)=\big(r(t)/r(1)\big)\,w(1)$ with $r=|w|_p$, and the radial lemma
  (\cref{lem:dc-ch3-3-6-equality-radial}) makes $r$ monotone — its derivative
  on $(0,1)$ is the nonnegative speed — and continuous, so
  $s=r/r(1)$ is the desired reparametrization (when $v=0$ the lift is
  identically $0$ and $s(t)=t$ works). The running-length identity
  $\ell(\sigma|_{[0,t]})=r(t)$ is the fundamental theorem of calculus for
  $r'=|\dot\sigma|$, and the image equality follows from the intermediate
  value theorem applied to the continuous monotone radius.
\end{proof}

\begin{lemma}[an arclength-proportional minimizer is the radial geodesic]
  \label{lem:dc-ch3-3-6-equality-geodesic}
  \lean{Riemannian.Exponential.exists_gauss_equality_geodesic_ball, Riemannian.Exponential.exists_gauss_equality_geodesic}
  \uses{lem:dc-ch3-3-6-equality-manifold, lem:dc-ch7-2-8-ray-geodesic, lem:dc-ch3-3-6-confinement}
  \leanok
  \dcref{ch3:3.6}
  There is $\rho>0$ such that every $C^1$ curve $\sigma:[0,1]\to M$ from $p$
  to $\exp_p(v)$, $|v|<\rho$, staying in $\exp_p(B_\rho(0))$, that is
  \emph{parametrized proportionally to arc length while realizing the radial
  length} — $\ell(\sigma|_{[0,t]})=t\,|v|_p$ for all $t\in[0,1]$ — coincides
  with the radial geodesic with matching parametrization:
  $\sigma(t)=\exp_p(t\,v)$ for all $t\in[0,1]$, and $\sigma$ satisfies the
  intrinsic geodesic equation on $(0,1)$. This is the normal-ball core of
  \cref{cor:dc-ch3-3-9}: the equality case of the minimizing property turns
  an arclength-proportional minimizer into a geodesic. By the confinement
  lemma (\cref{lem:dc-ch3-3-6-confinement}) the stay-in-ball hypothesis is
  again automatic.
\end{lemma}
\begin{proof}
  \sketch
  \leanok
  \uses{lem:dc-ch3-3-6-equality-manifold, lem:dc-ch7-2-8-ray-geodesic, lem:dc-ch3-3-6-confinement}
  The running-length hypothesis at $t=1$ realizes the radial length, so
  \cref{lem:dc-ch3-3-6-equality-manifold} yields the monotone
  reparametrization $s$ with $\sigma(t)=\exp_p(s(t)v)$ and
  $\ell(\sigma|_{[0,t]})=s(t)|v|_p$. Comparing with the hypothesis
  $\ell(\sigma|_{[0,t]})=t\,|v|_p$ and cancelling $|v|_p>0$ (for $v\ne0$;
  the case $v=0$ is the constant curve) pins $s(t)=t$. The geodesic equation
  transfers from the radial geodesic (\cref{lem:dc-ch7-2-8-ray-geodesic}) at
  every interior time by locality of the moving-foot equation, since
  $\sigma$ agrees with the ray on a neighborhood of each $t\in(0,1)$.

\end{proof}

\begin{lemma}[equality case in polar form, piecewise: the competitor runs along the ray]
  \label{lem:dc-ch3-3-6-equality-ray-piecewise}
  \lean{Riemannian.Exponential.exists_gauss_equality_ray_interval, Riemannian.Exponential.exists_gauss_equality_ray_piecewise, Riemannian.Exponential.monotoneOn_of_partition}
  \uses{lem:dc-ch3-3-6-equality-radial, lem:dc-ch3-3-6-equality-ray, lem:dc-ch3-3-6-radius}
  \leanok
  \dcref{ch3:3.6}
  There is $\rho>0$ (a Gauss ball for $\exp_p$) such that:
  \begin{itemize}
  \item (\emph{general interval, no anchoring}) every differentiable polar
    competitor $w:[a,b]\to B_\rho(0)\subset T_pM$ with continuous derivative
    that realizes the radius comparison with \emph{equality},
    $|w(b)|_p-|w(a)|_p=\int_a^b|\dot c|\,dt$ where $c=\exp_p\circ w$, has
    monotone radius $r=|w|_p$ and satisfies
    $w(t)=\bigl(r(t)/r(b)\bigr)\,w(b)$ for every $t\in[a,b]$ — the lift is
    \emph{not} required to start at the origin;
  \item (\emph{piecewise}) every polar competitor on a partition
    $\tau_0\le\dots\le\tau_n$ — continuous on the span, $C^1$ on each piece —
    realizing the total radius comparison with equality has monotone radius
    on the span, satisfies the equality in FTC form on every piece,
    $r(t)=r(\tau_i)+\int_{\tau_i}^t|\dot c|$, and traces the ray of its
    final point: $w(t)=\bigl(r(t)/r(\tau_n)\bigr)\,w(\tau_n)$ for every $t$
    in the span.
  \end{itemize}
\end{lemma}
\begin{proof}
  \sketch
  \leanok
  \uses{lem:dc-ch3-3-6-equality-radial, lem:dc-ch3-3-6-equality-ray, lem:dc-ch3-3-6-radius}
  For the general interval: the argument of
  \cref{lem:dc-ch3-3-6-equality-ray} never used the anchoring $w(a)=0$
  beyond producing a zero of the radius. If the radius never vanishes on
  $[a,b]$, the normalized lift $\eta=w/r$ has vanishing derivative on the
  interior — combining $r'=\langle w,w'\rangle_p/r$ with the radial velocity
  of \cref{lem:dc-ch3-3-6-equality-radial} — hence is constant on $(a,b]$,
  and takes the same value in the limit at $a$ by continuity of $w$ and $r$;
  so $w(t)=r(t)\,\eta(b)=\bigl(r(t)/r(b)\bigr)w(b)$ throughout. If it
  vanishes somewhere, let $t_*$ be its last zero: on $[a,t_*]$ the monotone
  radius is squeezed to $0$, so $w=0$ and the identity holds with
  coefficient $0$; on $(t_*,b]$ the previous argument applies verbatim.

  For the piecewise form: the radius comparison of
  \cref{lem:dc-ch3-3-6-radius} holds on every piece, and the per-piece
  radius gains telescope to $r(\tau_n)-r(\tau_0)$; total equality therefore
  forces \emph{equality on every piece} (a sum of nonnegative deficits
  vanishes only if each does). Each piece then has monotone radius and
  traces the ray of its right endpoint by the general-interval case, the
  per-piece monotone radii glue along the partition, and a backward
  induction over the pieces composes the scalar factors:
  $$\frac{r(t)}{r(\tau_{i+1})}\cdot\frac{r(\tau_{i+1})}{r(\tau_n)}
    =\frac{r(t)}{r(\tau_n)},$$
  the degenerate case $r(\tau_{i+1})=0$ being absorbed by monotonicity
  ($r(t)=0$ forces $w(t)=0$). No gluing of directions across the corners is
  needed.
\end{proof}

\begin{lemma}[confinement, piecewise: a short piecewise curve from $p$ stays in the Gauss ball]
  \label{lem:dc-ch3-3-6-confinement-piecewise}
  \lean{Riemannian.Exponential.exists_forall_mem_expMap_ball_of_pathELength_lt_piecewise}
  \uses{lem:dc-ch3-3-6-competitor, lem:dc-ch3-3-6-confinement}
  \leanok
  \dcref{ch3:3.6}
  For every target radius $\rho'>0$ there is a length threshold $\delta>0$
  such that every piecewise-$C^1$ curve $\sigma:[0,1]\to M$ starting at $p$
  with $\ell(\sigma)<\delta$ satisfies $\sigma(t)\in\exp_p(B_{\rho'}(0))$
  for all $t\in[0,1]$: the piecewise analogue of
  \cref{lem:dc-ch3-3-6-confinement}, discharging the escape case of the
  piecewise equality analysis.
\end{lemma}
\begin{proof}
  \sketch
  \leanok
  \uses{lem:dc-ch3-3-6-competitor}
  Truncate $\sigma$ at time $t$ along the truncated partition
  $\min(\tau_j,t)$ and reparametrize each truncated piece affinely to a
  piecewise-$C^1$ curve on $[0,1]$ (collapsed pieces are constants); by
  additivity of the length over the partition and invariance of each
  piece's length under its monotone affine reparametrization, the
  reparametrized truncation has length
  $\ell(\sigma|_{[0,t]})\le\ell(\sigma)$. If the endpoint $\sigma(t)$
  escaped $\exp_p(B_{\rho''}(0))$ with $\rho''=\min(\rho',\varepsilon)$, the
  piecewise escape estimate of \cref{lem:dc-ch3-3-6-competitor} would force
  $\rho''/\sqrt{c}\le\ell(\sigma|_{[0,t]})<\delta=\rho''/\sqrt{c}$, a
  contradiction.
\end{proof}

\begin{lemma}[the equality case in the manifold, piecewise competitors]
  \label{lem:dc-ch3-3-6-equality-manifold-piecewise}
  \lean{Riemannian.Exponential.exists_gauss_equality_manifold_piecewise_ball, Riemannian.Exponential.exists_gauss_equality_manifold_piecewise}
  \uses{lem:dc-ch3-3-6-equality-ray-piecewise, lem:dc-ch3-3-6-confinement-piecewise, lem:dc-ch3-2-9-c1diffeo, lem:dc-ch7-2-8-chart-length, lem:dc-ch3-3-6-equality-manifold}
  \leanok
  \dcref{ch3:3.6}
  There is $\rho>0$ (a Gauss ball for $\exp_p$, with $\exp_p$ injective on
  $B_\rho(0)$) such that every \emph{piecewise-$C^1$} competitor
  $\sigma:[0,1]\to M$ from $p$ to $\exp_p(v)$, $|v|<\rho$ — continuous on
  $[0,1]$, $C^1$ on each piece of a partition $0=\tau_0\le\dots\le\tau_n=1$ —
  that realizes the radial length $\ell(\sigma)=|v|_p$ is a \emph{monotone
  reparametrization of the radial geodesic}: there is a continuous monotone
  $s:[0,1]\to[0,1]$ with $s(0)=0$, $s(1)=1$,
  $$\sigma(t)=\exp_p\big(s(t)\,v\big)\quad\text{and}\quad
    \ell\big(\sigma|_{[0,t]}\big)=s(t)\,|v|_p\quad\text{for all }t\in[0,1];$$
  in particular $\sigma([0,1])=\exp_p([0,1]\cdot v)$ in its literal regularity
  class. By the piecewise
  confinement lemma (\cref{lem:dc-ch3-3-6-confinement-piecewise}) no
  stay-in-ball hypothesis is needed.
\end{lemma}
\begin{proof}
  \sketch
  \leanok
  \uses{lem:dc-ch3-3-6-equality-ray-piecewise, lem:dc-ch3-3-6-confinement-piecewise, lem:dc-ch3-2-9-c1diffeo, lem:dc-ch7-2-8-chart-length}
  As in the $C^1$ case (\cref{lem:dc-ch3-3-6-equality-manifold}), lift
  $\sigma$ through the $C^1$ local inverse of $\exp_p$
  (\cref{lem:dc-ch3-2-9-c1diffeo}): the polar lift $w=\exp_p^{-1}\circ\sigma$
  is continuous on $[0,1]$ and $C^1$ on each piece, with $w(0)=0$,
  $w(1)=v$. Reading the length in the chart at $p$
  (\cref{lem:dc-ch7-2-8-chart-length}) piece by piece and summing over the
  partition, the hypothesis $\ell(\sigma)=|v|_p$ becomes the total FTC-form
  equality of the piecewise polar lemma
  (\cref{lem:dc-ch3-3-6-equality-ray-piecewise}), which forces
  $w(t)=\bigl(r(t)/r(1)\bigr)\,v$ with $r=|w|_p$ monotone; so $s=r/r(1)$ is
  the desired reparametrization (for $v=0$ the lift is identically $0$ and
  $s(t)=t$ works). The running-length identity
  $\ell(\sigma|_{[0,t]})=r(t)$ telescopes over the truncated partition
  $\min(\tau_j,t)$, evaluating each truncated piece by the per-piece FTC
  clause; the image equality is the intermediate value theorem for the
  continuous monotone radius.
\end{proof}

\begin{proposition}[geodesics locally minimize]
  \label{prop:dc-ch3-3-6}
  \lean{Riemannian.Exponential.exists_le_pathELength_piecewise, Riemannian.Exponential.exists_gauss_equality_manifold_piecewise}
  \uses{lem:dc-ch3-3-6-normal, lem:dc-ch3-3-6-piecewise, lem:dc-ch3-3-6-competitor, lem:dc-ch3-3-6-equality-radial, lem:dc-ch3-3-6-equality-ray, lem:dc-ch3-3-6-equality-manifold, lem:dc-ch3-3-6-equality-geodesic, lem:dc-ch3-3-6-confinement, lem:dc-ch3-3-6-equality-ray-piecewise, lem:dc-ch3-3-6-confinement-piecewise, lem:dc-ch3-3-6-equality-manifold-piecewise, prop:dc-ch3-2-9}
  \leanok
  \dcref{ch3:3.6}
  Let $p\in M$, $U$ a normal neighborhood of $p$, and $B\subset U$ a normal ball of
  center $p$. Let $\gamma:[0,1]\to B$ be a geodesic segment with $\gamma(0)=p$. If
  $c:[0,1]\to M$ is any piecewise differentiable curve joining $\gamma(0)$ to
  $\gamma(1)$, then $\ell(\gamma)\le\ell(c)$, and if equality holds then
  $\gamma([0,1])=c([0,1])$.
\end{proposition}
\begin{proof}
  \sketch
  \leanok
  \uses{lem:dc-ch3-3-6-competitor, lem:dc-ch3-3-6-equality-manifold-piecewise, lem:dc-ch3-3-6-piecewise, lem:dc-ch3-3-6-radius, lem:dc-ch3-3-6-rayspeed, lem:dc-ch3-3-6-equality-radial, lem:dc-ch3-3-6-equality-ray, lem:dc-ch3-3-6-equality-manifold, lem:dc-ch3-3-6-equality-geodesic, lem:dc-ch3-3-6-confinement, lem:dc-ch3-3-6-equality-ray-piecewise, lem:dc-ch3-3-6-confinement-piecewise}
  Suppose first $c([0,1])\subset B$. For $t\ne 0$ write uniquely
  $c(t)=\exp_p(r(t)\,v(t))=f(r(t),t)$, $|v(t)|=1$, $r:(0,1]\to\mathbb{R}$ positive
  piecewise differentiable. Then except at finitely many points
  $\frac{dc}{dt}=\frac{\partial f}{\partial r}r'(t)+\frac{\partial f}{\partial t}$.
  By the Gauss lemma $\langle\frac{\partial f}{\partial r},\frac{\partial f}{\partial t}\rangle=0$
  and $|\frac{\partial f}{\partial r}|=1$, so
  
  $$
  \Big|\frac{dc}{dt}\Big|^2=|r'(t)|^2+\Big|\frac{\partial f}{\partial t}\Big|^2\ge|r'(t)|^2,\qquad\text{(1)}
  $$
  
  $$
  \int_\varepsilon^1\Big|\frac{dc}{dt}\Big|dt\ge\int_\varepsilon^1|r'(t)|dt\ge r(1)-r(\varepsilon).\qquad\text{(2)}
  $$
  
  Letting $\varepsilon\to0$ gives $\ell(c)\ge r(1)=\ell(\gamma)$. If equality holds
  then $|\frac{\partial f}{\partial t}|=0$ and $r'(t)>0$, so $c$ is a monotone
  reparametrization of $\gamma$ and $c([0,1])=\gamma([0,1])$. If $c([0,1])\not\subset B$,
  let $t_1$ be the first time $c(t_1)\in\partial B$; with $\rho$ the radius of $B$,
  $\ell(c)\ge\ell_{[0,t_1]}(c)\ge\rho>\ell(\gamma)$.
\end{proof}

\begin{lemma}[geodesic segments of the chart flow, from any base point]
  \label{lem:dc-ch3-3-7-descent}
  \lean{Riemannian.Exponential.isGeodesicOn_uniform_flow_segment}
  \uses{def:dc-ch3-2-1, lem:dc-ch3-2-7-vuniform, lem:dc-ch7-2-8-chart-independence}
  \leanok
  \dcref{ch3:3.7}
  Let $p\in M$, let $\mathbf{x}$ be the chart at $p$, and let $Z$ be a local flow
  of the coordinate geodesic equation as in \cref{lem:dc-ch3-2-7-vuniform}: for
  every initial condition $\zeta$ in a ball around $(\mathbf{x}(p),0)$, the curve
  $Z(\zeta)$ solves $z'=(z_2,\,-\Gamma_p(z_2,z_2)(z_1))$ on
  $[-\varepsilon,\varepsilon]$ with $Z(\zeta)(0)=\zeta$ and position component in
  the chart image. Let $0<T<\varepsilon$ and let $(y,w/T)$ be an initial
  condition in the ball. Then the curve
  $\gamma(s)=\mathbf{x}^{-1}\bigl(Z(y,w/T)(sT)_1\bigr)$, $s\in[0,1]$,

  \begin{itemize}
  \item starts at $\gamma(0)=\mathbf{x}^{-1}(y)$ and is continuous;
  \item stays in the chart domain at $p$, where its $\mathbf{x}$-reading is the
    position component of the flow line;
  \item has $\mathbf{x}$-velocity $w$ at $s=0$;
  \item is a geodesic on $[0,1]$: it satisfies the geodesic equation of
    \cref{def:dc-ch3-2-1} at every $s$, read in the chart at its own foot
    $\gamma(s)$;
  \item has $\mathbf{x}$-acceleration
    $x''(0)=-\Gamma_p(w,w)(y)=\bigl(\mathrm{spray}(y,w)\bigr)_2$ at $s=0$ — the
    fixed-chart-at-$p$ reading of the geodesic equation $x''+\Gamma(x',x')=0$,
    read algebraically off the spray with no time integration.
  \end{itemize}

  The base point $\mathbf{x}^{-1}(y)$ ranges over the whole chart-image ball —
  the initial condition is not required to lie over $p$ — so a single flow
  provides geodesic segments with a uniform velocity ball at every base point
  near $p$.
\end{lemma}
\begin{proof}
  \sketch
  \leanok
  Write $u(s)=Z(y,w/T)(sT)_1$ and $v(s)=Z(y,w/T)(sT)_2$ for the position and
  velocity components of the reparametrized flow line. By the chain rule
  $u'(s)=T\,v(s)$ and $v'(s)=-T\,\Gamma_p(v(s),v(s))(u(s))$ on the open window
  $|s|<\varepsilon/T\supset[0,1]$, so, the Christoffel contraction being
  quadratic in its velocity slot,
  $$
  u''(s)=T\,v'(s)\cdot T/T=-\Gamma_p\bigl(T v(s),T v(s)\bigr)(u(s))
    =-\Gamma_p\bigl(u'(s),u'(s)\bigr)(u(s)):
  $$
  the reading $u$ solves the second-order geodesic equation in the fixed chart
  at $p$. The initial conditions give $u(0)=y$ and $u'(0)=T\cdot w/T=w$; hence
  $u''(0)=-\Gamma_p(w,w)(y)=\bigl(\mathrm{spray}(y,w)\bigr)_2$. Every
  foot $\gamma(s)=\mathbf{x}^{-1}(u(s))$ lies in the chart domain at $p$, so by
  the chart-independence of the geodesic equation
  (\cref{lem:dc-ch7-2-8-chart-independence}) $\gamma$ satisfies the geodesic
  equation read in the chart at $\gamma(s)$ as well. Continuity of $\gamma$
  follows since $u$ is differentiable and $\mathbf{x}^{-1}$ is continuous on the
  chart image.
\end{proof}

\begin{lemma}[interior chart velocity of the chart-flow segment]
  \label{lem:dc-ch3-3-7-descent-velocity}
  \lean{Riemannian.Geodesic.isGeodesicOn_uniform_flow_segment_Ioo}
  \uses{lem:dc-ch3-3-7-descent}
  \leanok
  \dcref{ch3:3.7}
  In the setting of \cref{lem:dc-ch3-3-7-descent}, the flow-geodesic segment
  $\gamma(s)=\mathbf{x}^{-1}\bigl(Z(y,w/T)(sT)_1\bigr)$ has a well-defined
  $\mathbf{x}$-velocity \emph{at every} interior time, not only at $s=0$: on the
  open window $|s|<\varepsilon/T\supset[0,1]$ one has
  $$
  (\mathbf{x}\circ\gamma)'(s)=T\,Z(y,w/T)(sT)_2,
  $$
  the rescaled velocity component of the flow line. In particular the joining
  geodesic of \cref{thm:dc-ch3-3-7} carries at each interior time $t_0\in(0,1)$ the
  chart-velocity datum $\mathbf{x}$-$\gamma'(t_0)=T\,Z(\cdot)(t_0T)_2$ that the
  convex-neighborhood admissibility \cref{lem:dc-ch3-4-2-interior} consumes —
  closing the interface gap between \cref{thm:dc-ch3-3-7} (whose $s=0$ velocity
  clause is \cref{lem:dc-ch3-3-7-descent}) and the interior deduction.
\end{lemma}
\begin{proof}
  \sketch
  \leanok
  Differentiate the position component $u(s)=Z(y,w/T)(sT)_1$ by the chain rule:
  the flow solves $z'(t)=\bigl(z_2(t),-\Gamma_p(z_2,z_2)(z_1)\bigr)$
  (\cref{lem:dc-ch3-3-7-descent}), so the time-derivative of its first component
  is its second component; the reparametrisation $t=sT$ multiplies the derivative
  by $T$, giving $u'(s)=T\,Z(y,w/T)(sT)_2$ on the whole open window where the flow
  derivative is interior to $[-\varepsilon,\varepsilon]$. This is the chart velocity
  of $\gamma$ since $\mathbf{x}\circ\gamma=u$ on the chart-image window.
\end{proof}

\begin{lemma}[differential of the pair map]
  \label{lem:dc-ch3-3-7-pairmap}
  \lean{Riemannian.Exponential.exists_pairMap_hasStrictFDerivAt}
  \uses{lem:dc-ch3-2-7-vuniform, lem:dc-ch3-2-2-c1dep, lem:dc-ch3-2-9-strict, lem:dc-ch3-3-7-descent}
  \leanok
  \dcref{ch3:3.7}
  In the setting of \cref{lem:dc-ch3-3-7-descent}, the pair map
  $$
  G(y,w)=\bigl(y,\;Z(y,w/T)(T)_1\bigr)
  $$
  — the $\mathbf{x}$-coordinate expression of $F(q,v)=(q,\exp_q v)$, since by
  \cref{lem:dc-ch3-3-7-descent} the point $Z(y,w/T)(T)_1$ is the time-$1$ value
  of the geodesic through $q=\mathbf{x}^{-1}(y)$ with $\mathbf{x}$-velocity $w$ —
  is strictly differentiable at $(\mathbf{x}(p),0)$ with differential
  $$
  dG_{(\mathbf{x}(p),0)}(a,b)=(a,\;a+b),
  $$
  the block matrix $\begin{pmatrix}I & 0\\ I & I\end{pmatrix}$. In
  particular the differential is an invertible linear map.
\end{lemma}
\begin{proof}
  \sketch
  \leanok
  The point $(\mathbf{x}(p),0)$ is an equilibrium of the geodesic spray
  $F(x,u)=(u,\,-\Gamma_p(u,u)(x))$, and $F$ is $C^1$ near it. By the
  $C^1$-dependence theorem for solution families
  (\cref{lem:dc-ch3-2-2-c1dep}, applied at the equilibrium exactly as in
  \cref{lem:dc-ch3-2-9-strict}), the map sending an initial condition $\zeta$
  near the equilibrium to the solution curve $Z(\zeta)|_{[0,T]}$ is strictly
  differentiable there, with differential $D$ the solution of the variational
  equation $Dv(t)-\int_0^t A\,(Dv)(s)\,ds=v$ along the constant operator curve
  $A=dF_{(\mathbf{x}(p),0)}$. Since $\Gamma_p$ is quadratic in the velocity,
  $A(a,b)=(b,0)$; hence $A^2=0$ and $Dv(t)=v+tAv$ solves the variational
  equation. The pair map is the composition of $\zeta\mapsto Z(\zeta)|_{[0,T]}$
  with the linear rescaling $(a,b)\mapsto(a,b/T)$, the evaluation at time $T$
  and the projection to the position component, so
  $$
  dG_{(\mathbf{x}(p),0)}(a,b)
    =\Bigl(a,\;\bigl[(a,b/T)+T\,A(a,b/T)\bigr]_1\Bigr)
    =\bigl(a,\;a+T\cdot b/T\bigr)=(a,\;a+b).
  $$
  The inverse of the differential is $(a,c)\mapsto(a,c-a)$.
\end{proof}

\begin{lemma}[the pair map is $C^1$ on a ball]
  \label{lem:dc-ch3-3-7-c1pair}
  \lean{Riemannian.Exponential.exists_pairMap_contDiffOn}
  \uses{lem:dc-ch3-2-7-c1flow, lem:dc-ch3-3-7-pairmap, lem:dc-ch3-2-9-c1}
  \leanok
  \dcref{ch3:3.7}
  In the setting of \cref{lem:dc-ch3-3-7-pairmap}, the pair map
  $G(y,w)=(y,\,Z(y,w/T)(T)_1)$ is $C^1$ on the open set of admissible initial
  conditions $\{(y,w):\ (y,w/T)\in B((\mathbf{x}(p),0),r)\}$. In particular,
  slicing at each base point: for every $y$ with $|y-\mathbf{x}(p)|<r$, the
  chart exponential $w\mapsto Z(y,w/T)(T)_1$ at the base point
  $\mathbf{x}^{-1}(y)$ is $C^1$ on the uniform velocity ball $|w|<Tr$.
\end{lemma}
\begin{proof}
  \sketch
  \leanok
  By the $C^1$-dependence of the flow on its full initial condition, at an
  arbitrary admissible base solution (\cref{lem:dc-ch3-2-7-c1flow}), the map
  $\zeta\mapsto Z(\zeta)|_{[0,T]}$ is strictly differentiable at every $\zeta$
  in the ball around $(\mathbf{x}(p),0)$ — the differential being the solution
  of the variational equation along the flow line through $\zeta$, no longer
  the explicit shear of \cref{lem:dc-ch3-3-7-pairmap} away from the
  equilibrium. Composing with the linear rescaling, the evaluation at $T$ and
  the projection, $G$ is strictly differentiable at every point of the stated
  open set, and a map strictly differentiable at every point of an open set is
  $C^1$ there (as in \cref{lem:dc-ch3-2-9-c1}). The slice statement follows by
  composing with the affine embedding $w\mapsto(y,w)$, which maps the ball
  $|w|<Tr$ into the admissible set.
\end{proof}

\begin{lemma}[the geodesic segment determines the pair map]
  \label{lem:dc-ch3-3-7-flow-agree}
  \lean{Riemannian.Exponential.uniform_flow_pairMap_agree}
  \uses{lem:dc-ch3-3-7-descent}
  \leanok
  \dcref{ch3:3.7}
  Let $Z$ and $Z'$ be two local flows of the chart-$p$ geodesic system as in
  \cref{lem:dc-ch3-3-7-descent}, with time scales $0<T<\varepsilon$ and
  $0<T'<\varepsilon'$ and admissibility balls of radii $r$, $r'$ around
  $(\mathbf{x}(p),0)$. For any initial data $(y,w)$ with $(y,w/T)$ admissible
  for $Z$ and $(y,w/T')$ admissible for $Z'$,
  $$Z(y,w/T)(T)_1=Z'(y,w/T')(T')_1:$$
  the chart reading of the time-$1$ endpoint of the geodesic from
  $\mathbf{x}^{-1}(y)$ with chart velocity $w$ does not depend on the flow
  used to compute it. In particular the pair maps
  $G(y,w)=(y,\,Z(y,w/T)(T)_1)$ built from any two such flows agree on the
  common admissible set.
\end{lemma}
\begin{proof}
  \sketch
  \leanok
  \uses{lem:dc-ch3-3-7-descent, lem:dc-ch7-2-8-uniqueness}
  By \cref{lem:dc-ch3-3-7-descent} — whose proof gives the geodesic property
  on the open window $|s|<\varepsilon/T\supset[0,1]$, and likewise
  $|s|<\varepsilon'/T'$ for $Z'$ — the two rescaled segments
  $s\mapsto\mathbf{x}^{-1}(Z(y,w/T)(sT)_1)$ and
  $s\mapsto\mathbf{x}^{-1}(Z'(y,w/T')(sT')_1)$ are continuous geodesics
  passing through $\mathbf{x}^{-1}(y)$ at $s=0$ with the same chart velocity
  $w$. By uniqueness of geodesics with given initial data
  (\cref{lem:dc-ch7-2-8-uniqueness}) they agree on the common open window,
  which contains $s=1$ since $T<\varepsilon$ and $T'<\varepsilon'$. Both
  endpoint readings lie in the chart image, where $\mathbf{x}^{-1}$ is
  injective, so evaluating the identity at $s=1$ equates them.
\end{proof}

\begin{lemma}[the derivative of the pair map is invertible near the zero section]
  \label{lem:dc-ch3-3-7-invertible}
  \lean{Riemannian.Exponential.exists_pairMap_hasStrictFDerivAt_equiv_ball}
  \uses{lem:dc-ch3-3-7-pairmap, lem:dc-ch3-3-7-c1pair}
  \leanok
  \dcref{ch3:3.7}
  There are a flow package $(r,\varepsilon,T,Z)$ as in
  \cref{lem:dc-ch3-3-7-descent} and a radius $\rho>0$ such that the pair map
  $G(y,w)=(y,\,Z(y,w/T)(T)_1)$ fixes the center
  ($G(\mathbf{x}(p),0)=(\mathbf{x}(p),\mathbf{x}(p))$), is $C^1$ on the open
  set of admissible initial conditions, is differentiable at the center with
  derivative the shear $(a,b)\mapsto(a,a+b)$ of
  \cref{lem:dc-ch3-3-7-pairmap}, and has an \emph{invertible} differential at
  every point of the ball $B((\mathbf{x}(p),0),\rho)$.
\end{lemma}
\begin{proof}
  \sketch
  \leanok
  \uses{lem:dc-ch3-3-7-pairmap, lem:dc-ch3-3-7-c1pair, lem:dc-ch3-3-7-flow-agree}
  Take the flow package of \cref{lem:dc-ch3-3-7-c1pair}, so that $G$ is $C^1$
  on the open admissible set and its differential $dG$ is continuous there.
  The derivative at the center is the shear: \cref{lem:dc-ch3-3-7-pairmap}
  computes it for \emph{its} flow package, and by
  \cref{lem:dc-ch3-3-7-flow-agree} the two pair maps agree on a neighborhood
  of $(\mathbf{x}(p),0)$, so the derivative transports; the center is fixed
  because the zero-velocity flow line is constant. The shear $S(a,b)=(a,a+b)$
  is invertible with inverse $(a,c)\mapsto(a,c-a)$, and invertibility is
  stable under small perturbation: choose $\rho>0$ so that
  $\|dG_x-S\|<(\|S^{-1}\|+1)^{-1}$ for $|x-(\mathbf{x}(p),0)|<\rho$; then
  $t:=S^{-1}\circ(S-dG_x)$ has $\|t\|\le\|S^{-1}\|\,\|S-dG_x\|<1$, the
  Neumann series inverts $1-t$, and $dG_x=S\circ(1-t)$ is invertible — as in
  \cref{lem:dc-ch3-2-9-invertible}.
\end{proof}

\begin{theorem}[totally normal neighborhood]
  \label{thm:dc-ch3-3-7}
  \lean{Riemannian.Exponential.exists_totallyNormal_c1_diffeo, Riemannian.Exponential.exists_totallyNormal_neighborhood}
  \uses{prop:dc-ch3-2-7, lem:dc-ch3-3-7-descent, lem:dc-ch3-3-7-descent-velocity, lem:dc-ch3-3-7-pairmap, lem:dc-ch3-3-7-c1pair}
  \leanok
  \dcref{ch3:3.7}
  For any $p\in M$ there exist a neighborhood $W$ of $p$ and a number $\delta>0$
  such that for every $q\in W$, $\exp_q$ is a diffeomorphism on
  $B_\delta(0)\subset T_qM$ and $\exp_q(B_\delta(0))\supset W$; that is, $W$ is a
  normal neighborhood of each of its points.
\end{theorem}
\begin{proof}
  \sketch
  \leanok
  \uses{lem:dc-ch3-3-7-descent, lem:dc-ch3-3-7-descent-velocity, lem:dc-ch3-3-7-pairmap, lem:dc-ch3-3-7-c1pair, lem:dc-ch3-3-7-invertible}
  Let $\varepsilon$, $V$, $\mathcal{U}$ be as in \cref{prop:dc-ch3-2-7}, with
  $\mathcal{U}\subset TU$, $V\subset\mathbf{x}(U)$. Define $F:\mathcal{U}\to M\times M$
  by $F(q,v)=(q,\exp_q v)$. Around $F(p,0)=(p,p)$ take coordinates
  $(U\times U;\mathbf{x},\mathbf{x})$; then since $(d\exp_p)_0=I$,

  $$
  dF_{(p,0)}=\begin{pmatrix}I & I\\ 0 & I\end{pmatrix},
  $$

  so $F$ is a local diffeomorphism near $(p,0)$, mapping a neighborhood
  $\mathcal{U}'=\{(q,v);\ q\in V',\ v\in T_qM,\ |v|<\delta\}$ diffeomorphically onto
  a neighborhood $W'$ of $(p,p)$. Choose $W\ni p$ with $W\times W\subset W'$. For
  $q\in W$ and $B_\delta(0)\subset T_qM$, $F(\{q\}\times B_\delta(0))\supset\{q\}\times W$,
  so by definition of $F$, $\exp_q(B_\delta(0))\supset W$.
\end{proof}

\begin{remark}[uniqueness and differentiable dependence]
  \label{rem:dc-ch3-3-8}
  \uses{thm:dc-ch3-3-7}
  \dcref{ch3:3.8}
  From \cref{thm:dc-ch3-3-7} and the minimizing property of geodesics, any two points
  $q_1,q_2\in W$ are joined by a unique minimizing geodesic $\gamma$ of length
  $<\delta$; moreover $\gamma$ depends differentiably on $(q_1,q_2)$: there is a
  unique $v\in T_{q_1}M$ with $F^{-1}(q_1,q_2)=(q_1,v)$, depending differentiably on
  $(q_1,q_2)$, with $\gamma'(0)=v$. Such a $W$ is a \emph{totally normal neighborhood} of
  $p\in M$.
\end{remark}

\begin{lemma}[a minimizing broken pair of radial geodesics has no corner]
  \label{lem:dc-ch3-3-9-corner}
  \lean{Riemannian.Exponential.eq_neg_of_forall_edist_expMap_eq}
  \uses{def:dc-ch3-expmap, def:dc-ch7-2-4, lem:dc-ch3-2-9-c1ball, lem:dc-ch7-2-8-gram-bound, lem:dc-ch7-2-8-chart-length}
  \leanok
  \dcref{ch3:3.9}
  Let $M$ be a Riemannian manifold whose distance $d$ is the Riemannian
  distance, $x\in M$, and let $u_1,u_2\in T_xM$ be $g$-unit vectors. If the
  two radial legs through $x$ jointly realize the distance between their
  endpoints for every small radius — there is $\eta_0>0$ with
  $d\big(\exp_x(\eta u_1),\exp_x(\eta u_2)\big)=2\eta$ for all
  $0<\eta<\eta_0$ — then $u_2=-u_1$: the broken curve
  $\exp_x(\eta u_1)\rightsquigarrow x\rightsquigarrow\exp_x(\eta u_2)$ has no
  corner at $x$.
\end{lemma}
\begin{proof}
  \sketch
  \leanok
  The engine is the \emph{chord upper bound}: for every $\theta>1$ there
  is $\rho>0$ with $d(\exp_x v,\exp_x w)\le\theta\,|w-v|_x$ for
  $\|v\|,\|w\|<\rho$. Its witness is the exponential image
  $s\mapsto\exp_x\big(v+s\,(w-v)\big)$ of the straight segment: since
  $\exp_x$ is $C^1$ near $0$ with $d(\exp_x)_0=\operatorname{id}$
  (\cref{lem:dc-ch3-2-9-c1ball}) and the Gram form is jointly continuous
  and positive definite at the pole (\cref{lem:dc-ch7-2-8-gram-bound}), the
  $g$-speed of this curve is at most $\theta\,|w-v|_x$ on a small enough
  ball, uniformly in the direction; integrating
  (\cref{lem:dc-ch7-2-8-chart-length}) bounds its length, and the distance
  is at most the length of any such competitor (\cref{def:dc-ch7-2-4}).
  Applying the bound to $v=\eta u_1$, $w=\eta u_2$ for small $\eta$ gives
  $2\eta\le\theta\,\eta\,|u_2-u_1|_x$, i.e. $2\le\theta\,|u_2-u_1|_x$ for
  every $\theta>1$, hence $|u_2-u_1|_x\ge2$. Expanding with
  $|u_1|_x=|u_2|_x=1$:
  $4\le|u_2-u_1|_x^2=2-2\langle u_1,u_2\rangle_x$, so
  $\langle u_1,u_2\rangle_x\le-1$, whence
  $|u_1+u_2|_x^2=2+2\langle u_1,u_2\rangle_x\le0$; positive definiteness of
  $g_x$ gives $u_1+u_2=0$.
\end{proof}

\begin{lemma}[anchored radial identification of a distance-realizing segment]
  \label{lem:dc-ch3-3-9-anchored}
  \lean{Riemannian.Exponential.exists_radial_of_anchored}
  \uses{lem:dc-ch3-3-6-equality-geodesic, lem:dc-ch7-2-8-normal-ball-dist, lem:dc-ch7-2-8-gram-bound}
  \leanok
  \dcref{ch3:3.9}
  Let $\gamma$ be $C^1$ on $[t_0,B]$, $t_0<B$, with speed $\ell>0$ in the
  sense that $\ell(\gamma|_{[t_0,t]})=\ell\,(t-t_0)$ and
  $d\big(\gamma(t_0),\gamma(t)\big)=\ell\,(t-t_0)$ for all $t\in[t_0,B]$.
  Then, with $q=\gamma(t_0)$, there are $t_1\in(t_0,B]$ and $v\in T_qM$ with
  $|v|_q=\ell\,(t_1-t_0)$ such that
  $$\gamma(r)=\exp_q\Big(\tfrac{r-t_0}{t_1-t_0}\,v\Big)
    \qquad\text{for all } r\in[t_0,t_1]:$$
  some initial subsegment of $\gamma$ is exactly the radial geodesic of the
  normal ball at its own starting point.
\end{lemma}
\begin{proof}
  \sketch
  \leanok
  Choose $t_1\in(t_0,B]$ with $\ell\,(t_1-t_0)$ smaller than both the
  metric-ball threshold $\delta$ of the normal ball at $q$
  (\cref{lem:dc-ch7-2-8-normal-ball-dist}) and, after conversion by the Gram
  comparison $\|w\|^2\le c\,\langle w,w\rangle_q$
  (\cref{lem:dc-ch7-2-8-gram-bound}, as packaged in
  \cref{lem:dc-ch7-2-8-normal-ball-dist}), the Gauss radius of the equality
  case at $q$. Since $d\big(q,\gamma(t_1)\big)=\ell\,(t_1-t_0)<\delta$, the point
  $\gamma(t_1)$ lies in the normal ball: $\gamma(t_1)=\exp_q v$ with
  $|v|_q=d\big(q,\gamma(t_1)\big)=\ell\,(t_1-t_0)$ by the distance formula
  of \cref{lem:dc-ch7-2-8-normal-ball-dist}, and the Gram bound puts the
  coordinate norm of $v$ inside the Gauss radius. The affine
  reparametrization $\sigma(s)=\gamma\big(t_0+s\,(t_1-t_0)\big)$ is $C^1$ on
  $[0,1]$, runs from $q$ to $\exp_q v$, and has running length
  $\ell(\sigma|_{[0,s]})=\ell\cdot s\,(t_1-t_0)=s\,|v|_q$ — exactly the
  arclength-proportional hypothesis of the equality case
  \cref{lem:dc-ch3-3-6-equality-geodesic}, which yields
  $\sigma(s)=\exp_q(s\,v)$ on $[0,1]$; undoing the reparametrization gives
  the claim.
\end{proof}

\begin{lemma}[localization at moving centers: the two-sided corner argument]
  \label{lem:dc-ch3-3-9-localization}
  \lean{Riemannian.Exponential.hasGeodesicEquationAt_of_arclength_edist}
  \uses{lem:dc-ch3-3-9-anchored, lem:dc-ch3-3-9-corner, lem:dc-ch7-2-8-ray-geodesic, lem:dc-ch3-2-6-reparam, def:dc-ch3-2-1}
  \leanok
  \dcref{ch3:3.9}
  Let $\gamma$ be parametrized proportionally to arc length with speed
  $\ell>0$ and distance-realizing on $[A,B]$ — i.e.
  $\ell(\gamma|_{[s,t]})=d\big(\gamma(s),\gamma(t)\big)=\ell\,(t-s)$ for all
  $A\le s\le t\le B$ — and let $A<u<B$ be an interior time such that
  $\gamma$ is $C^1$ on $[A,u]$ and on $[u,B]$ separately ($u$ may be a corner
  of a piecewise curve). Then $\gamma$ satisfies the geodesic equation of
  \cref{def:dc-ch3-2-1} at $u$.
\end{lemma}
\begin{proof}
  \sketch
  \leanok
  Anchor \cref{lem:dc-ch3-3-9-anchored} at $q=\gamma(u)$ twice: applied to
  $\gamma|_{[u,B]}$ it exhibits a forward subsegment
  $\gamma(r)=\exp_q\big(\tfrac{r-u}{t_1-u}v_2\big)$ on $[u,t_1]$; applied to
  the time reversal $s\mapsto\gamma(2u-s)$ — which is again
  arclength-proportional and distance-realizing, with the same speed — it
  exhibits a backward subsegment
  $\gamma(r)=\exp_q\big(\tfrac{u-r}{u-t_0'}v_1\big)$ on $[t_0',u]$. Let
  $u_i=v_i/|v_i|_q$. For $0<\eta<\ell\min(u-t_0',t_1-u)$ the points
  $\exp_q(\eta u_1)=\gamma(u-\eta/\ell)$ and
  $\exp_q(\eta u_2)=\gamma(u+\eta/\ell)$ realize the distance
  $d=\ell\cdot(2\eta/\ell)=2\eta$, so \cref{lem:dc-ch3-3-9-corner} gives
  $u_2=-u_1$. Substituting back, both subsegments trace the \emph{single}
  ray $\gamma(r)=\exp_q\big(\ell\,(r-u)\,u_2\big)$ for
  $r\in[t_0',t_1]$. Exponential rays are geodesics on an open interval
  around time $0$ (\cref{lem:dc-ch7-2-8-ray-geodesic}), affine
  reparametrization preserves the geodesic equation
  (\cref{lem:dc-ch3-2-6-reparam}), and the geodesic equation at $u$ only
  depends on the curve near $u$, where $\gamma$ coincides with this ray.

\end{proof}

\begin{corollary}[minimizing curve is a geodesic]
  \label{cor:dc-ch3-3-9}
  \lean{Riemannian.Exponential.isGeodesicOn_piecewise_of_arclength_forall_le, Riemannian.Exponential.isGeodesicOn_piecewise_of_arclength_edist, Riemannian.Exponential.isGeodesicOn_of_arclength_edist, Riemannian.Exponential.edist_le_pathELength_piecewise_partition}
  \uses{prop:dc-ch3-3-6, lem:dc-ch3-3-6-equality-geodesic, lem:dc-ch3-3-9-localization, lem:dc-ch3-3-9-corner, lem:dc-ch3-3-9-anchored, lem:dc-ch7-2-8-chart-length, def:dc-ch3-3-1, def:dc-ch3-2-1}
  \leanok
  \dcref{ch3:3.9}
  If a piecewise differentiable curve $\gamma:[a,b]\to M$, parametrized
  proportionally to arc length, has length $\le$ that of any other piecewise
  differentiable curve joining $\gamma(a)$ to $\gamma(b)$, then $\gamma$ is a
  geodesic; in particular it is regular.
\end{corollary}
\begin{proof}
  \sketch
  \leanok
  For $t\in[a,b]$ let $W$ be a totally normal neighborhood of $\gamma(t)$
  and $I\subset[a,b]$ a closed interval with $t\in I$ and $\gamma(I)\subset W$.
  Then $\gamma_I$ joins two points of a normal ball; by the hypothesis and
  \cref{prop:dc-ch3-3-6}, $\ell(\gamma_I)$ equals the length of the radial geodesic joining
  them, so $\gamma_I$ coincides with that geodesic. As $\gamma_I$ is parametrized
  proportionally to arc length it is a geodesic on $I$, hence at $t$.
\end{proof}

\begin{example}[the Lobatchevski plane]
  \label{ex:dc-ch3-3-10}
  \uses{cor:dc-ch3-3-9}
  \notready
  \dcref{ch3:3.10}
  \notready
  The isometries of a Riemannian manifold take geodesics into geodesics.
  Let $G=\{(x,y)\in\mathbb{R}^2;\ y>0\}$ with metric $g_{11}=g_{22}=\frac{1}{y^2}$,
  $g_{12}=g_{21}=0$. The $y$-axis segment $\gamma(t)=(0,t)$, $a\le t\le b$ ($a>0$),
  is the image of a geodesic: for any arc $c(t)=(x(t),y(t))$ with $c(a)=(0,a)$,
  $c(b)=(0,b)$,
  $$
  \ell(c)=\int_a^b\sqrt{(x')^2+(y')^2}\,\frac{dt}{y}\ge\int_a^b\frac{|y'|}{y}\,dt\ge\int_a^b\frac{dy}{y}=\ell(\gamma),
  $$
  so $\gamma$ minimizes, and by \cref{cor:dc-ch3-3-9} its image is a geodesic. The
  isometries $z\mapsto\frac{az+b}{cz+d}$, $z=x+iy$, $ad-bc=1$ carry the $y$-axis
  into upper semicircles or rays $x=x_0$, $y>0$, which
  are therefore geodesics; these are all the geodesics, since through each $p\in G$
  and each direction there passes such a circle centered on the $x$-axis (a normal
  direction giving a vertical ray).
\end{example}

\section{Convex neighborhoods}

By \cref{thm:dc-ch3-3-7} (cf. \cref{rem:dc-ch3-3-8}) every $p\in M$ has a totally normal neighborhood
$W$ with $\delta>0$ so that any two $q_1,q_2\in W$ are joined by a minimizing
geodesic of length $<\delta$; but such a geodesic need not lie entirely in $W$. A
subset $S\subset M$ is \emph{strongly convex} if for any two points $q_1,q_2$ in the
closure $\overline{S}$ there is a unique minimizing geodesic $\gamma$ joining them
whose interior lies in $S$. We show the radius of a totally normal ball can be
chosen so the ball is strongly convex.

The analytic heart of \cref{lem:dc-ch3-4-1} is the behaviour of the squared radial
distance $F(t)=|u(t)|_p^2$, $u(t)=\exp_p^{-1}(\gamma(t))$, along a geodesic $\gamma$
read in normal coordinates at $p$. We isolate its fibrewise (fixed base point $q$
and velocity $v$) calculus in the following lemmas; the metric
$\langle\cdot,\cdot\rangle_p$ is the fixed chart Gram inner product at $p$.

\begin{lemma}[Leibniz rule for the fixed metric at $p$]
  \label{lem:dc-ch3-4-1-leibniz}
  \lean{Riemannian.hasDerivAt_chartMetricInner}
  \leanok
  \dcref{ch3:4.1}
  For curves $f,h:\mathbb{R}\to T_pM$ differentiable at $t$, read in the chart at
  $p$,
  $$\frac{d}{dt}\langle f(t),h(t)\rangle_p
    = \langle f'(t),h(t)\rangle_p + \langle f(t),h'(t)\rangle_p.$$
\end{lemma}
\begin{proof}
  \sketch
  \leanok
  The form is the finite sum $\sum_{i,j}G_{ij}\,f^i(t)\,h^j(t)$ with constant
  coefficients $G_{ij}$ (the base point is fixed); each coordinate $f^i,h^j$ is a
  continuous linear functional of the curve, so the product rule applies termwise.
\end{proof}

\begin{lemma}[first time-derivative of the squared radial distance]
  \label{lem:dc-ch3-4-1-dF}
  \lean{Riemannian.hasDerivAt_chartMetricInner_diag}
  \uses{lem:dc-ch3-4-1-leibniz}
  \leanok
  \dcref{ch3:4.1}
  For $u:\mathbb{R}\to T_pM$ differentiable at $t$, the squared radial distance
  $F(t)=\langle u(t),u(t)\rangle_p$ satisfies
  $$\frac{\partial F}{\partial t}=2\Big\langle\frac{\partial u}{\partial t},u\Big\rangle_p.$$
\end{lemma}
\begin{proof}
  \sketch
  \leanok
  Apply \cref{lem:dc-ch3-4-1-leibniz} with $f=h=u$ and use the symmetry of
  $\langle\cdot,\cdot\rangle_p$.
\end{proof}

\begin{lemma}[second time-derivative of the squared radial distance]
  \label{lem:dc-ch3-4-1-ddF}
  \lean{Riemannian.hasDerivAt_deriv_chartMetricInner_diag}
  \uses{lem:dc-ch3-4-1-leibniz}
  \leanok
  \dcref{ch3:4.1}
  If $u$ is twice differentiable at $t$, then the first-derivative field
  $s\mapsto 2\langle u'(s),u(s)\rangle_p$ is differentiable at $t$ with
  $$\frac{\partial^2 F}{\partial t^2}
    = 2\Big\langle\frac{\partial^2 u}{\partial t^2},u\Big\rangle_p
      + 2\Big|\frac{\partial u}{\partial t}\Big|_p^2,$$
  the displayed second time-derivative of $F=|u|_p^2$.
\end{lemma}
\begin{proof}
  \sketch
  \leanok
  The field $s\mapsto 2\langle u'(s),u(s)\rangle_p$ is twice the bilinear pairing
  of $u'$ and $u$; differentiate it by \cref{lem:dc-ch3-4-1-leibniz} with
  $f=u'$, $h=u$.
\end{proof}

\begin{lemma}[normal coordinate along the radial geodesic]
  \label{lem:dc-ch3-4-1-radialcoord}
  \lean{Riemannian.Exponential.exists_expMap_inv_smul_eq}
  \uses{def:dc-ch3-expmap, lem:dc-ch3-2-9-c1diffeo}
  \leanok
  \dcref{ch3:4.1}
  There is $\varepsilon>0$ and a local inverse $\exp_p^{-1}$ of $\exp_p$ on
  $B_\varepsilon(0)$ such that for every $v$ and every $t$ with
  $\|tv\|<\varepsilon$, $\exp_p^{-1}(\exp_p(tv))=tv$. Hence along the radial
  geodesic $t\mapsto\exp_p(tv)$ the normal coordinate $u(t)=\exp_p^{-1}(\gamma(t,p,v))$
  is $tv$ itself, and $F(t)=|u(t)|_p^2=t^2|v|_p^2$.
\end{lemma}
\begin{proof}
  \sketch
  \leanok
  This is the left-inverse clause of the $C^1$ diffeomorphism
  \cref{lem:dc-ch3-2-9-c1diffeo}, instantiated at $w=tv$.
\end{proof}

\begin{lemma}[center Hessian: $\partial^2 F/\partial t^2(0,p,v)=2|v|^2$]
  \label{lem:dc-ch3-4-1-hessian-center}
  \lean{Riemannian.Exponential.hasDerivAt_secondDeriv_expMap_inv_sqNorm_radial}
  \uses{lem:dc-ch3-4-1-ddF, lem:dc-ch3-4-1-radialcoord}
  \leanok
  \dcref{ch3:4.1}
  For $q=p$ the geodesic through $p$ with velocity $v$ is the radial ray, so
  $u(t)=tv$ (\cref{lem:dc-ch3-4-1-radialcoord}), $u'=v$, $u''=0$, and by
  \cref{lem:dc-ch3-4-1-ddF}
  $$\frac{\partial^2 F}{\partial t^2}(0,p,v)=2|v|_p^2.$$
  In particular for a unit $v$ this equals $2>0$, the strict-minimum seed.
\end{lemma}
\begin{proof}
  \sketch
  \leanok
  With $u(t)=tv$ the first-derivative field is
  $s\mapsto 2\langle v,sv\rangle_p=2s\,|v|_p^2$, whose derivative at $0$ is
  $2|v|_p^2$.
\end{proof}

\begin{lemma}[the metric at $p$ is positive definite]
  \label{lem:dc-ch3-4-1-posdef}
  \lean{Riemannian.Exponential.chartMetricInner_extChartAt_self_pos}
  \leanok
  \dcref{ch3:4.1}
  For $v\neq0$, $|v|_p^2=\langle v,v\rangle_p>0$. Hence the center Hessian
  $\frac{\partial^2 F}{\partial t^2}(0,p,v)=2|v|_p^2$
  (\cref{lem:dc-ch3-4-1-hessian-center}) is strictly positive, giving the
  strict minimum of $F$ at $q=p$.
\end{lemma}
\begin{proof}
  \sketch
  \leanok
  Positive-definiteness of the Riemannian metric $g$ at $p$, read in the chart
  at $p$.
\end{proof}

\begin{lemma}[a tangent geodesic with positive Hessian stays outside]
  \label{lem:dc-ch3-4-1-outside}
  \lean{Riemannian.eventually_ge_of_deriv_deriv_pos}
  \leanok
  \dcref{ch3:4.1}
  If $F:\mathbb{R}\to\mathbb{R}$ is continuous at $0$ with $F'(0)=0$ and
  $F''(0)>0$, then $F(t)\ge F(0)$ for all $t$ near $0$. With $F=|u|_p^2$ and
  $F(0)=r^2$, a geodesic tangent to $S_r(p)$ at $\gamma(0)$ keeps
  $|\exp_p^{-1}\gamma(t)|^2\ge r^2$, i.e. does not re-enter the open ball
  $B_r(p)$.
\end{lemma}
\begin{proof}
  \sketch
  \leanok
  The second-derivative test: $F'(0)=0$ and $F''(0)>0$ give a local minimum of
  $F$ at $0$, so $F(t)\ge F(0)$ near $0$.
\end{proof}

\begin{lemma}[a tangent geodesic with positive Hessian stays \emph{strictly} outside]
  \label{lem:dc-ch3-4-1-strict-outside}
  \lean{Riemannian.eventually_lt_of_deriv_deriv_pos}
  \leanok
  \dcref{ch3:4.1}
  If $F:\mathbb{R}\to\mathbb{R}$ is continuous at $0$ with $F'(0)=0$ and
  $F''(0)>0$, then $F(0)<F(t)$ for \emph{every} $t\neq0$ near $0$ (a strict
  punctured minimum). With $F=|u|_p^2$ and $F(0)=r^2$, a geodesic tangent to
  $S_r(p)$ at $\gamma(0)$ keeps $|\exp_p^{-1}\gamma(t)|^2>r^2$ for $t\neq0$, i.e.
  stays \emph{strictly} outside the closed ball $\overline{B_r(p)}$ except at the
  point of tangency. This strengthens \cref{lem:dc-ch3-4-1-outside} and is the
  separation the convex-neighborhood argument \cref{prop:dc-ch3-4-2}
  consumes: it rules out the interior distance-maximum whose existence the
  proposition assumes for contradiction.
\end{lemma}
\begin{proof}
  \sketch
  \leanok
  By the sign form of the second-derivative test, $\operatorname{sign} F'(t)$
  agrees with $\operatorname{sign}(t)$ near $0$ (from $F'(0)=0$, $F''(0)>0$): $F'>0$
  on a right interval $(0,c)$ and $F'<0$ on a left interval $(a,0)$. Hence $F$ is
  strictly increasing on $[0,c)$ and strictly decreasing on $(a,0]$, so $F(0)<F(t)$
  on both punctured sides.
\end{proof}

\begin{lemma}[joint continuity of the second time-derivative in the base point]
  \label{lem:dc-ch3-4-1-continuity}
  \lean{Riemannian.Exponential.continuousOn_secondDerivChartForm}
  \uses{lem:dc-ch3-4-1-ddF, lem:dc-ch3-2-9-c2diffeo}
  \leanok
  \dcref{ch3:4.1}
  Write $G(q,v)=\frac{\partial^2 F}{\partial t^2}(0,q,v)$. Read in normal coordinates
  at $p$, with $a=\varphi_p(q)$ the chart position and $w$ the chart velocity, this
  value is the algebraic expression
  $$G(q,v)=2\big\langle D^2(\exp_p^{-1})_a(w,w)+D(\exp_p^{-1})_a\!\big(-\Gamma_a(w,w)\big),\ \exp_p^{-1}(q)\big\rangle_p
    +2\big|D(\exp_p^{-1})_a(w)\big|_p^2,$$
  where $-\Gamma_a(w,w)$ is the geodesic acceleration $x''(0)$ read in the chart at
  $p$ (from the geodesic equation $x''+\Gamma(x',x')=0$). Since $\exp_p^{-1}$ is $C^2$
  and the Christoffel symbols are smooth, $G$ is jointly continuous in $(q,v)$.
\end{lemma}
\begin{proof}
  \sketch
  \leanok
  The value $G(q,v)$ needs no time integration: the geodesic acceleration at $t=0$
  is $x''(0)=-\Gamma_a(w,w)$ directly from the geodesic ODE, so by the second-order
  chain rule $u''(0)=D^2(\exp_p^{-1})_a(w,w)+D(\exp_p^{-1})_a(x''(0))$. The three maps
  $\exp_p^{-1}$, $D(\exp_p^{-1})$, $D^2(\exp_p^{-1})$ are continuous on the normal chart
  image ($\exp_p^{-1}$ is $C^2$, \cref{lem:dc-ch3-2-9-c2diffeo}), $x''(0)=-\Gamma_a(w,w)$
  is continuous in $(a,w)$ (the geodesic spray is smooth), and the chart Gram form
  $\langle\cdot,\cdot\rangle_p$ is a continuous bilinear form; $G$ is a finite algebraic
  combination of these, hence continuous.
\end{proof}

\begin{lemma}[strict positivity of the second time-derivative at $q=p$]
  \label{lem:dc-ch3-4-1-poscenter}
  \lean{Riemannian.Exponential.secondDerivChartForm_extChartAt_self_pos}
  \uses{lem:dc-ch3-4-1-posdef, lem:dc-ch3-2-9-c2diffeo}
  \leanok
  \dcref{ch3:4.1}
  At $q=p$ one has $\exp_p^{-1}(p)=0$, so the $\langle u'',u\rangle_p$ term of $G(p,v)$
  vanishes and $G(p,v)=2\big|D(\exp_p^{-1})_0(w)\big|_p^2$. Since $D(\exp_p^{-1})_0$ is a
  linear isomorphism and $\langle\cdot,\cdot\rangle_p$ is positive definite, $G(p,v)>0$
  for $v\neq0$.
\end{lemma}
\begin{proof}
  \sketch
  \leanok
  $u(0,p,v)=\exp_p^{-1}(p)=0$ kills the first term. The derivative $D(\exp_p^{-1})_0$ is
  the inverse of $D(\varphi_p\circ\exp_p)_0$ (inverse function theorem), hence injective,
  so $D(\exp_p^{-1})_0(w)\neq0$ for $w\neq0$; then
  $\big|D(\exp_p^{-1})_0(w)\big|_p^2>0$ by positive definiteness of the metric at $p$
  (\cref{lem:dc-ch3-4-1-posdef}).
\end{proof}

\begin{lemma}[a neighborhood of $p$ with strictly positive Hessian]
  \label{lem:dc-ch3-4-1-posnbhd}
  \lean{Riemannian.Exponential.exists_secondDerivChartForm_pos_nhds}
  \uses{lem:dc-ch3-4-1-continuity, lem:dc-ch3-4-1-poscenter}
  \leanok
  \dcref{ch3:4.1}
  There is a neighborhood $V$ of $p$ (inside a normal chart) such that
  $\frac{\partial^2 F}{\partial t^2}(0,q,v)>0$ for every $q\in V$ and every unit chart
  velocity $v$.
\end{lemma}
\begin{proof}
  \sketch
  \leanok
  $G(q,v)$ is jointly continuous (\cref{lem:dc-ch3-4-1-continuity}) and $G(p,v)>0$ for
  $v\neq0$ (\cref{lem:dc-ch3-4-1-poscenter}). Since the unit sphere of chart velocities
  is compact, the tube lemma produces an open neighborhood $V$ of $p$ on which
  $G(q,v)>0$ for all $q\in V$ and all unit $v$.
\end{proof}

\begin{lemma}[degree-2 homogeneity: strict positivity for every nonzero velocity]
  \label{lem:dc-ch3-4-1-homog}
  \lean{Riemannian.Exponential.exists_secondDerivChartForm_pos_nhds_ne}
  \uses{lem:dc-ch3-4-1-posnbhd, lem:dc-ch3-2-6}
  \leanok
  \dcref{ch3:4.1}
  The value $G(q,v)=\frac{\partial^2 F}{\partial t^2}(0,q,v)$ is degree-2 homogeneous in
  the chart velocity: $G(q,cv)=c^2\,G(q,v)$. Consequently the neighborhood $V$ of
  \cref{lem:dc-ch3-4-1-posnbhd} satisfies $G(q,v)>0$ for every $q\in V$ and every
  \emph{nonzero} chart velocity $v$, removing the unit-speed restriction.
\end{lemma}
\begin{proof}
  \sketch
  \leanok
  Under $v\mapsto cv$ every building block of $G$ scales:
  $D(\exp_p^{-1})_a(cv)=c\,D(\exp_p^{-1})_a(v)$,
  $D^2(\exp_p^{-1})_a(cv,cv)=c^2\,D^2(\exp_p^{-1})_a(v,v)$, and the geodesic acceleration
  $-\Gamma_a(cv,cv)=c^2\,\big(-\Gamma_a(v,v)\big)$ by the degree-2 homogeneity of the
  Christoffel contraction (the coordinate form of the homogeneity \cref{lem:dc-ch3-2-6});
  since $\langle\cdot,\cdot\rangle_p$ is bilinear both summands acquire the factor $c^2$.
  Writing a nonzero $v=|v|\,(v/|v|)$ with $v/|v|$ of unit chart norm,
  $G(q,v)=|v|^2\,G(q,v/|v|)>0$ by \cref{lem:dc-ch3-4-1-posnbhd}.
\end{proof}

\begin{lemma}[the F-identity: second-order chain rule for the squared radial distance]
  \label{lem:dc-ch3-4-1-fidentity}
  \lean{Riemannian.Exponential.hasDerivAt_and_deriv_deriv_sqNormComp}
  \uses{lem:dc-ch3-4-1-dF, lem:dc-ch3-4-1-ddF}
  \leanok
  \dcref{ch3:4.1}
  Let $\mathrm{finv}$ be $C^2$ on an open set $S\ni a$, and let $x:\mathbb R\to T_pM$ be a
  curve differentiable near $0$ with values eventually in $S$, and with $x(0)=a$, $x'(0)=w$,
  second derivative $x''(0)=\xi$. Then, writing $u=\mathrm{finv}\circ x$, the squared radial
  distance $F(s)=\langle u(s),u(s)\rangle_p$ satisfies
  $$\frac{\partial F}{\partial t}(0)=2\langle D\mathrm{finv}(a)(w),\mathrm{finv}(a)\rangle_p,$$
  $$\frac{\partial^2 F}{\partial t^2}(0)
    =2\big\langle D^2\mathrm{finv}(a)(w,w)+D\mathrm{finv}(a)(\xi),\ \mathrm{finv}(a)\big\rangle_p
      +2\,\big|D\mathrm{finv}(a)(w)\big|_p^2,$$
  the displayed derivatives of $F=|u|_p^2$ follow from the second-order chain rule
  $u''(0)=D^2\mathrm{finv}(a)(w,w)+D\mathrm{finv}(a)(x''(0))$ made explicit.
\end{lemma}
\begin{proof}
  \sketch
  \leanok
  \uses{lem:dc-ch3-4-1-leibniz, lem:dc-ch3-4-1-dF, lem:dc-ch3-4-1-ddF}
  Near $0$, $u$ is differentiable with $u'(s)=D\mathrm{finv}(x(s))\big(x'(s)\big)$ (chain rule:
  $\mathrm{finv}$ is differentiable on the open $S$ and $x(s)\in S$ for $s$ near $0$), so
  $\partial F/\partial t(s)=2\langle u'(s),u(s)\rangle_p$ (\cref{lem:dc-ch3-4-1-dF}); evaluating
  at $0$ gives the first identity. Differentiating the field $s\mapsto D\mathrm{finv}(x(s))$
  against $x'$ at $0$ by the composition and bilinear-application rules gives
  $u''(0)=D^2\mathrm{finv}(a)(w,w)+D\mathrm{finv}(a)(x''(0))$, and \cref{lem:dc-ch3-4-1-ddF}
  turns $\partial^2 F/\partial t^2(0)=2\langle u''(0),u(0)\rangle_p+2|u'(0)|_p^2$ into the
  second identity.
\end{proof}

\begin{lemma}[the F-identity along the moving-base geodesic flow]
  \label{lem:dc-ch3-4-1-flowreading}
  \lean{Riemannian.Exponential.hasDerivAt_and_deriv_deriv_sqNorm_flowReading}
  \uses{lem:dc-ch3-4-1-fidentity, lem:dc-ch3-3-7-descent, lem:dc-ch3-4-1-continuity}
  \leanok
  \dcref{ch3:4.1}
  Let $Z$ be a local flow of the chart-$p$ geodesic spray, $\mathrm{finv}$ a $C^2$ exponential
  inverse on an open $S\ni y$, and $(y,T^{-1}w)$ an admissible initial condition. Along the
  geodesic reading $x(s)=\varphi_p(\gamma(s))$ of the geodesic through $q=\varphi_p^{-1}(y)$
  with chart velocity $w$, the squared radial distance $F(s)=|\mathrm{finv}(x(s))|_p^2$ has
  $\partial F/\partial t(0)=2\langle D\mathrm{finv}(y)(w),\mathrm{finv}(y)\rangle_p$ and
  $$\frac{\partial^2 F}{\partial t^2}(0,q,v)=G(q,v),$$
  the algebraic second-derivative form of \cref{lem:dc-ch3-4-1-continuity}. This is the
  identification of the abstract form $G$ with the genuine second time-derivative of $F$ along
  the geodesic family — the assembly step of the Lemma 4.1.
\end{lemma}
\begin{proof}
  \sketch
  \leanok
  \uses{lem:dc-ch3-4-1-fidentity, lem:dc-ch3-3-7-descent}
  The base descent (\cref{lem:dc-ch3-3-7-descent}) gives the chart reading with $x(0)=y$,
  $x'(0)=w$ and the fixed-chart acceleration $x''(0)=(\mathrm{spray}(y,w))_2$; its open-window
  form gives differentiability of $x$ on a neighborhood of $0$, and continuity of $x$ places
  $x(s)$ in the $C^2$ domain $S$ for $s$ near $0$. Applying the F-identity
  (\cref{lem:dc-ch3-4-1-fidentity}) with $a=y$ and $\xi=(\mathrm{spray}(y,w))_2$ yields exactly
  the algebraic form $G(q,v)$ of \cref{lem:dc-ch3-4-1-continuity}.
\end{proof}

\begin{lemma}[geodesic tangent to a sphere stays outside]
  \label{lem:dc-ch3-4-1}
  \lean{Riemannian.Exponential.exists_forall_geodesic_tangent_stays_outside_ball}
  \uses{lem:dc-ch3-4-1-flowreading, lem:dc-ch3-4-1-homog, lem:dc-ch3-4-1-outside, lem:dc-ch3-4-1-posnbhd, lem:dc-ch3-3-7-descent}
  \leanok
  \dcref{ch3:4.1}
  For any $p\in M$ there exists $c>0$ such that any geodesic of $M$ tangent at
  $q\in M$ to the geodesic sphere $S_r(p)$ of radius $r<c$ stays outside the
  geodesic ball $B_r(p)$ in some neighborhood of $q$.
\end{lemma}
\begin{proof}
  \sketch
  \leanok
  \uses{lem:dc-ch3-4-1-dF, lem:dc-ch3-4-1-ddF, lem:dc-ch3-4-1-hessian-center,
    lem:dc-ch3-4-1-radialcoord, lem:dc-ch3-4-1-posdef, lem:dc-ch3-4-1-outside,
    lem:dc-ch3-4-1-continuity, lem:dc-ch3-4-1-poscenter, lem:dc-ch3-4-1-posnbhd,
    lem:dc-ch3-3-7-descent, lem:dc-ch3-4-1-homog, lem:dc-ch3-3-5-ball,
    lem:dc-ch3-4-1-fidentity, lem:dc-ch3-4-1-flowreading}
  Let $W$ be a totally normal neighborhood of $p$; by the lemma of
  homogeneity restrict to unit-speed geodesics, so to the unit bundle
  $T_1W=\{(q,v);\ q\in W,\ v\in T_qM,\ |v|=1\}$. After shrinking $W$, assume that
  $\overline{W}$ is compactly contained in a normal neighborhood $U$ of $p$; since
  $T_1\overline{W}$ is compact, decrease $I$ so that
  $\gamma(I\times T_1W)\subset U$. Let
  $\gamma:I\times T_1W\to M$, $I=(-\varepsilon,\varepsilon)$, be the differentiable
  map with $t\mapsto\gamma(t,q,v)$ the unit-speed geodesic through $q$ with velocity
  $v$ at $t=0$. Set $u(t,q,v)=\exp_p^{-1}(\gamma(t,q,v))$ and
  $F:I\times T_1W\to\mathbb{R}$, $F(t,q,v)=|u(t,q,v)|^2$, the squared "distance"
  from $p$ to a point moving along $\gamma$. Then
  
  $$
  \frac{\partial F}{\partial t}=2\Big\langle\frac{\partial u}{\partial t},u\Big\rangle,\qquad\frac{\partial^2 F}{\partial t^2}=2\Big\langle\frac{\partial^2 u}{\partial t^2},u\Big\rangle+2\Big|\frac{\partial u}{\partial t}\Big|^2.
  $$
  
  Choose $r>0$ with $\exp_p B_r(0)=B_r(p)\subset W$. If $\gamma$ is tangent to
  $S_r(p)$ at $q=\gamma(0,q,v)$, then by the Gauss lemma
  $\langle\frac{\partial u}{\partial t}(0,q,v),u(0,q,v)\rangle=0$, i.e.
  $\frac{\partial F}{\partial t}(0,q,v)=0$. It suffices to show that for $r$ small
  the critical point $(0,q,v)$ is a strict minimum. For $q=p$, $u(t,p,v)=tv$, so
  $\frac{\partial^2 F}{\partial t^2}(0,p,v)=2|v|^2=2>0$. Hence there is a
  neighborhood $V\subset W$ of $p$ with $\frac{\partial^2 F}{\partial t^2}(0,q,v)>0$
  for all $q\in V$, $|v|=1$. Take $c>0$ with $\exp_p B_c(0)\subset V$: any geodesic
  in $B_c(p)$ tangent to $S_r(p)$, $r<c$, has a strict local minimum of $F$ at
  $(0,q,v)$, so in a neighborhood of $q$ its points stay outside $B_r(p)$.
\end{proof}

\begin{lemma}[strict form of Lemma 4.1: tangent geodesic stays \emph{strictly} outside]
  \label{lem:dc-ch3-4-1-strict}
  \lean{Riemannian.Exponential.exists_forall_geodesic_tangent_stays_strictly_outside_ball}
  \uses{lem:dc-ch3-4-1-flowreading, lem:dc-ch3-4-1-homog, lem:dc-ch3-4-1-strict-outside, lem:dc-ch3-4-1-posnbhd, lem:dc-ch3-3-7-descent}
  \leanok
  \dcref{ch3:4.1}
  Same data as \cref{lem:dc-ch3-4-1}, but with the \emph{strict} conclusion. For any
  $p\in M$ there are a $C^2$ exponential inverse $\mathrm{finv}$, an open neighborhood
  $V$ of $\varphi_p(p)$, and a local geodesic flow package $(r,\varepsilon,T,Z)$ such
  that: for every base position $y\in V$ and every nonzero chart velocity $w$ with
  $(y,T^{-1}w)$ admissible, if the geodesic reading $F(s)=|\exp_p^{-1}(\gamma(s))|_p^2$
  is \emph{tangent} to the geodesic sphere at $s=0$ (i.e. $\partial F/\partial t(0)=0$,
  the Gauss lemma), then $F(0)<F(s)$ for \emph{every} $s\neq0$ near $0$: the tangent
  geodesic stays \emph{strictly} outside the ball in a punctured neighborhood of its base
  point. This is the form the convex-neighborhood proof \cref{prop:dc-ch3-4-2} consumes.
\end{lemma}
\begin{proof}
  \sketch
  \leanok
  \uses{lem:dc-ch3-4-1-flowreading, lem:dc-ch3-4-1-homog, lem:dc-ch3-4-1-strict-outside,
    lem:dc-ch3-4-1-posnbhd, lem:dc-ch3-3-7-descent}
  Verbatim the assembly of \cref{lem:dc-ch3-4-1} — the F-identity
  $\partial^2 F/\partial t^2(0,q,v)=G(q,v)$ (\cref{lem:dc-ch3-4-1-flowreading}), the strict
  positivity of $G$ on $V$ for nonzero velocities (\cref{lem:dc-ch3-4-1-posnbhd},
  \cref{lem:dc-ch3-4-1-homog}), and the tangency $\partial F/\partial t(0)=0$ — but the
  final second-derivative test is the \emph{strict} one (\cref{lem:dc-ch3-4-1-strict-outside},
  in place of the non-strict
  \cref{lem:dc-ch3-4-1-outside}, giving $F(0)<F(s)$ on the punctured neighborhood.
\end{proof}

\begin{lemma}[intrinsic strict separation: an arbitrary geodesic tangent to a sphere stays strictly outside]
  \label{lem:dc-ch3-4-2-rebase}
  \lean{Riemannian.Exponential.exists_forall_intrinsic_geodesic_tangent_strictly_outside_ball}
  \uses{lem:dc-ch3-4-1-strict, lem:dc-ch3-2-5-unique}
  \leanok
  \dcref{ch3:4.1}
  The strict Lemma 4.1 (\cref{lem:dc-ch3-4-1-strict}) is phrased for the geodesic
  \emph{flow reading} $s\mapsto\varphi_p^{-1}((Z(y,T^{-1}w)(sT))_1)$ of its own flow $Z$;
  the convex-neighborhood proof must apply it to an \emph{arbitrary} intrinsic geodesic — the
  joining geodesic of \cref{thm:dc-ch3-3-7}, re-based at its interior distance-maximum. For any
  $p\in M$ there are a $C^2$ exponential inverse $\mathrm{finv}$, an open neighborhood $V$ of
  $\varphi_p(p)$, and radii $0<r$, $0<T<\varepsilon$ such that for \emph{every} continuous
  intrinsic geodesic $\sigma$ on an open interval $(-a,a)$, $a>0$, whose base point $\sigma(0)$
  reads into $V$ (with $\sigma(0)$ in the chart at $p$) and whose chart velocity
  $w=(\varphi_p\circ\sigma)'(0)\neq0$ is admissible ($(\varphi_p(\sigma 0),T^{-1}w)$ in the flow's
  ball of initial conditions), the radial functional $F(s)=|\mathrm{finv}(\varphi_p(\sigma s))|_p^2$
  satisfies: if $\sigma$ is tangent to the geodesic sphere at $0$ ($\partial F/\partial t(0)=0$),
  then $F(0)<F(s)$ for every $s\neq0$ near $0$.
\end{lemma}
\begin{proof}
  \sketch
  \leanok
  \uses{lem:dc-ch3-4-1-strict, lem:dc-ch3-2-5-unique}
  Feed $\sigma$ into the \emph{readback} of \cref{lem:dc-ch3-2-5-unique} for the strict
  lemma's own flow $Z$: since
  the readback consumes only the generic flow clauses of $Z$ (not a particular construction), it
  applies verbatim, and on the overlap window $\sigma(s)=\varphi_p^{-1}((Z(\varphi_p(\sigma 0),
  T^{-1}w)(sT))_1)$. Hence $F$ coincides near $0$ with the strict lemma's radial functional, so
  the tangency hypothesis and the strict separation $F(0)<F(s)$ transfer across the equality of
  functions (using $\varphi_p(\sigma s)$ equals the flow reading's chart position on the window).

\end{proof}

\begin{lemma}[the radial distance has no interior local maximum along a geodesic]
  \label{lem:dc-ch3-4-2-nomax}
  \lean{Riemannian.Exponential.exists_forall_intrinsic_geodesic_not_isLocalMax_radial}
  \uses{lem:dc-ch3-4-2-rebase}
  \leanok
  \dcref{ch3:4.2}
  With the package of \cref{lem:dc-ch3-4-2-rebase}, for every admissible continuous intrinsic
  geodesic $\sigma$ the radial functional $F(s)=|\mathrm{finv}(\varphi_p(\sigma s))|_p^2$ has
  \emph{no} local maximum at the base point $0$. This is the exclusion the
  convex-neighborhood contradiction rests on: at an interior point where the distance from $p$ to
  a joining geodesic attains a maximum, $F$ would have a local maximum, which this rules out.
\end{lemma}
\begin{proof}
  \sketch
  \leanok
  \uses{lem:dc-ch3-4-2-rebase}
  A local maximum of $F$ at $0$ forces $\partial F/\partial t(0)=0$ (Fermat's theorem,
  ), whence \cref{lem:dc-ch3-4-2-rebase} gives the strict
  separation $F(0)<F(s)$ for every $s\neq0$ near $0$. But a local maximum also gives $F(s)\le F(0)$
  nearby; the two contradict on the (nonempty) punctured neighborhood $\mathbb R\setminus\{0\}$.

\end{proof}

\begin{lemma}[metric--radial bridge: the radial functional is the squared distance]
  \label{lem:dc-ch3-4-2-bridge}
  \lean{Riemannian.Exponential.sq_dist_eq_chartMetricInner_expMapInv}
  \uses{lem:dc-ch3-4-2-nomax, prop:dc-ch3-3-6}
  \leanok
  \dcref{ch3:4.2}
  Let $\mathrm{finv}$ be the $C^2$ exponential inverse of \cref{lem:dc-ch3-4-2-nomax} (which now
  exposes its left-inverse clause $\mathrm{finv}(\varphi_p(\exp_p w))=w$ for $\|w\|<\varepsilon_L$),
  and let the ambient distance be the Riemannian distance of $g$. Then there is $\rho>0$ such that
  for every chart velocity $v$ with $\|v\|<\rho$,
  $$ \big(d(p,\exp_p v)\big)^2 = \big\langle \mathrm{finv}(\varphi_p(\exp_p v)),\,
     \mathrm{finv}(\varphi_p(\exp_p v))\big\rangle_p . $$
  This is the identification $F=d^2$ that the convex-neighborhood contradiction consumes:
  the chart radial functional $F(x)=\langle\exp_p^{-1}(x),\exp_p^{-1}(x)\rangle_p$ of
  \cref{lem:dc-ch3-4-2-nomax} equals the squared metric distance from $p$, so a maximum of
  $d(p,\cdot)$ along a geodesic is a maximum of $F$.
\end{lemma}
\begin{proof}
  \sketch
  \leanok
  \uses{lem:dc-ch3-4-2-nomax, prop:dc-ch3-3-6}
  Rewrite $\mathrm{finv}(\varphi_p(\exp_p v))=v$ by the left inverse (valid for $\|v\|<\varepsilon_L$),
  reducing the claim to $\big(d(p,\exp_p v)\big)^2=\langle v,v\rangle_p$. The radial distance
  realization $d(p,\exp_p v)=\sqrt{\langle v,v\rangle_p}$ (\cref{prop:dc-ch3-3-6}, the Gauss-lemma
  consequence for the normal ball) together with $\mathrm{edist}=\mathrm{ofReal}\circ\mathrm{dist}$ in
  the ambient metric space and the positive-semidefiniteness $\langle v,v\rangle_p\ge0$ of the chart
  Gram form gives the identity on the overlap $\|v\|<\rho=\min(\varepsilon_e,\varepsilon_L)$.
\end{proof}

\begin{lemma}[the radial functional is the squared distance on a geodesic ball]
  \label{lem:dc-ch3-4-2-bridge-ball}
  \lean{Riemannian.Exponential.exists_ball_sq_dist_eq_chartMetricInner}
  \uses{lem:dc-ch3-4-2-bridge, prop:dc-ch3-3-6}
  \leanok
  \dcref{ch3:4.2}
  With the same data as \cref{lem:dc-ch3-4-2-bridge}, there is $\delta'>0$ such that on the whole
  metric ball $B_{\delta'}(p)$ the chart radial functional is the squared distance:
  $\big(d(p,x)\big)^2 = \langle \mathrm{finv}(\varphi_p x), \mathrm{finv}(\varphi_p x)\rangle_p$ for
  every $x$ with $d(p,x)<\delta'$. This is the form the max-distance contradiction consumes
  directly: on a small geodesic ball around $p$, a local maximum of $d(p,\cdot)$ along a geodesic is
  a local maximum of $F$, which \cref{lem:dc-ch3-4-2-nomax} forbids.
\end{lemma}
\begin{proof}
  \sketch
  \leanok
  \uses{lem:dc-ch3-4-2-bridge, prop:dc-ch3-3-6}
  Every $x$ with $d(p,x)<\delta'$ lies in the normal $\varepsilon_e$-ball (the escape clause of the
  normal-ball realization \cref{prop:dc-ch3-3-6}: points outside $\exp_p(B_{\varepsilon_e}(0))$ are at
  distance $\ge\delta_e$), so $x=\exp_p v$ with $d(p,x)=\sqrt{\langle v,v\rangle_p}$. The local
  coordinate-norm bound $\|v\|^2\le c\,\langle v,v\rangle_p$ (positive-definiteness of the chart Gram
  form at $p$) forces $\|v\|<\min(\varepsilon_e,\varepsilon_L)$ once
  $\delta'\le\min(\delta_e,\min(\varepsilon_e,\varepsilon_L)/\sqrt c)$, so the left inverse gives
  $\mathrm{finv}(\varphi_p x)=v$ and \cref{lem:dc-ch3-4-2-bridge} applies.
\end{proof}

\begin{lemma}[the max-distance interior deduction: convexity from no interior maximum]
  \label{lem:dc-ch3-4-2-maxdeduction}
  \lean{Riemannian.lt_of_forall_not_isLocalMax_of_le}
  \leanok
  \dcref{ch3:4.2}
  A continuous function $h$ on $[0,1]$ with $h(0)\le\beta$, $h(1)\le\beta$ and \emph{no} interior
  local maximum on $(0,1)$ stays \emph{strictly} below $\beta$ on the open interval: $h(s)<\beta$ for
  every $s\in(0,1)$. This is the convex-neighborhood contradiction stripped of the geometry:
  with $h(s)=d(p,\gamma(s))$ the distance from $p$ to a joining geodesic, the maximum of $h$ over
  $[0,1]$ sits at an endpoint (an interior maximizer would be a local maximum, forbidden by
  \cref{lem:dc-ch3-4-2-nomax} through the bridge $F=d^2$ of \cref{lem:dc-ch3-4-2-bridge-ball}), hence
  is $\le\beta$, and no interior point can equal it.
\end{lemma}
\begin{proof}
  \sketch
  \leanok
  The maximum of $h$ over the compact $[0,1]$ is attained at some $s_0$; if $s_0$ were interior it
  would be a local maximum, contradicting the hypothesis, so $s_0\in\{0,1\}$ and $h(s_0)\le\beta$.
  If some interior $s$ had $h(s)=\beta$, then $\beta\le h(s)\le h(s_0)\le\beta$ forces $h(s)=h(s_0)$,
  so $s$ is itself an interior maximizer — again a contradiction. Hence $h(s)<\beta$.
\end{proof}

\begin{definition}[strongly convex set]
  \label{def:dc-ch3-4-2-stronglyconvex}
  \lean{Riemannian.Exponential.StronglyConvex}
  \leanok
  \dcref{ch3:4.2}
  A set $S\subset M$ is \emph{strongly convex} if for any two points $q_1,q_2\in\overline{S}$
  there is a minimizing geodesic $\gamma:[0,1]\to M$ (constant-speed, realizing the distance
  $d(q_1,q_2)$ proportionally to arc length: $d(\gamma(s),\gamma(t))=|s-t|\,d(q_1,q_2)$) whose
  open-arc interior $\gamma((0,1))$ lies in $S$, and this joining geodesic is unique among such
  competitors (pointwise equality on $[0,1]$).
\end{definition}

\begin{lemma}[intrinsic speed in a fixed chart]
  \label{lem:dc-ch3-4-2-speedchart}
  \lean{Riemannian.Exponential.speedSq_eq_chartMetricInner_extChartAt}
  \uses{lem:dc-ch3-2-1-constspeed}
  \leanok
  \dcref{ch3:4.2}
  For a curve $\gamma$ continuous at $t$ with $\gamma(t)$ in the source of the chart at $\alpha$ and
  coordinate velocity $\xi=(\varphi_\alpha\circ\gamma)'(t)$, the intrinsic squared speed
  $\langle\gamma'(t),\gamma'(t)\rangle_g$ equals the chart-Gram value
  $\langle\xi,\xi\rangle$ read in the \emph{fixed} chart at $\alpha$ (not the moving chart at
  $\gamma(t)$), at the chart position $\varphi_\alpha(\gamma(t))$.
\end{lemma}
\begin{proof}
  \sketch
  \leanok
  The intrinsic velocity is the tangent-coordinate-change readback of $\xi$, and the chart-Gram form
  is the intrinsic inner product of the readbacks; composing the two identities gives the claim.
\end{proof}

\begin{lemma}[chart velocity bounded by the conserved speed]
  \label{lem:dc-ch3-4-2-velbound}
  \lean{Riemannian.Exponential.exists_sq_norm_deriv_le_speedSq}
  \uses{lem:dc-ch3-4-2-speedchart}
  \leanok
  \dcref{ch3:4.2}
  For every $p$ there are a constant $c>0$ and a neighborhood $V$ of $\varphi_p(p)$ such that, whenever
  a curve $\gamma$ reads into $V$ at $t$ with coordinate velocity $\xi=(\varphi_p\circ\gamma)'(t)$, one
  has $\|\xi\|^2\le c\,\langle\gamma'(t),\gamma'(t)\rangle_g$. Along a geodesic the right-hand side is
  constant, so this bounds the chart velocity \emph{uniformly along the geodesic} by its initial speed.
\end{lemma}
\begin{proof}
  \sketch
  \leanok
  The uniform coordinate-norm bound $\|\xi\|^2\le c\,\langle\xi,\xi\rangle_{p,y}$ (positive-definiteness
  of the Gram form, spread over a neighborhood of $\varphi_p(p)$ by the tube lemma) combined with the
  fixed-chart speed reading \cref{lem:dc-ch3-4-2-speedchart}.
\end{proof}

\begin{lemma}[the inverse pair map fixes the center diagonal]
  \label{lem:dc-ch3-4-2-invdiagonal}
  \lean{Riemannian.Exponential.exists_totallyNormal_c1_diffeo}
  \uses{thm:dc-ch3-3-7}
  \leanok
  \dcref{ch3:4.2}
  The inverse pair map $\mathrm{Ginv}$ of \cref{thm:dc-ch3-3-7} fixes the center
  diagonal: $\mathrm{Ginv}(\varphi_p(p),\varphi_p(p))=(\varphi_p(p),0)$. This is
  the $\exp_p^{-1}(p)=0$ read in the chart at $p$ — the joining velocity
  from $p$ to itself is zero.
\end{lemma}
\begin{proof}
  \sketch
  \leanok
  The pair map $G(y,w)=(y,(Z(y,w/T)T)_1)$ sends $(\varphi_p(p),0)$ to
  $(\varphi_p(p),\varphi_p(p))$, since the zero-velocity flow segment is the
  constant geodesic at $p$: $(Z(\varphi_p(p),0)T)_1=\varphi_p(p)$ (the center
  fixed point of the spray flow). As $(\varphi_p(p),0)$ lies in the product ball
  $B$ on which $\mathrm{Ginv}\circ G=\mathrm{id}$, applying the two-sided inverse
  gives $\mathrm{Ginv}(\varphi_p(p),\varphi_p(p))=(\varphi_p(p),0)$.
\end{proof}

\begin{lemma}[uniform smallness of a chart-pair functional near the center]
  \label{lem:dc-ch3-4-2-pairsmall}
  \lean{Riemannian.Exponential.exists_closedBall_forall_pair_norm_lt}
  \leanok
  \dcref{ch3:4.2}
  Let $\Psi$ be a map of chart-coordinate pairs $E\times E\to F$, continuous at the
  diagonal center $(\varphi_p(p),\varphi_p(p))$ and vanishing there. Then for every
  tolerance $r>0$ there is a radius $\beta>0$ such that the closed geodesic ball
  $\overline{B_\beta(p)}$ lies inside the chart source at $p$ and every pair of
  points $q_1,q_2\in\overline{B_\beta(p)}$ reads to a small value:
  $\|\Psi(\varphi_p(q_1),\varphi_p(q_2))\|<r$.
\end{lemma}
\begin{proof}
  \sketch
  \leanok
  By continuity of $\Psi$ at the diagonal center and $\Psi=0$ there, the set
  $\{x:\|\Psi(x)\|<r\}$ is a neighborhood of $(\varphi_p(p),\varphi_p(p))$, hence
  contains a product ball of some radius $\rho$. The chart $\varphi_p$ is
  continuous at $p$ with $\varphi_p(p)$ the center, so
  $\varphi_p^{-1}(B_\rho(\varphi_p(p)))\cap(\text{chart source})$ is a neighborhood
  of $p$ and contains $\overline{B_\beta(p)}$ for some $\beta>0$. For
  $q_1,q_2\in\overline{B_\beta(p)}$ both chart images lie in $B_\rho(\varphi_p(p))$,
  so the pair lies in the product ball and $\|\Psi(\varphi_p(q_1),\varphi_p(q_2))\|<r$.

\end{proof}

\begin{lemma}[the joining velocity is uniformly small on a small ball]
  \label{lem:dc-ch3-4-2-velsmall}
  \lean{Riemannian.Exponential.exists_closedBall_forall_ginvSnd_norm_lt}
  \uses{lem:dc-ch3-4-2-pairsmall, lem:dc-ch3-4-2-invdiagonal, thm:dc-ch3-3-7}
  \leanok
  \dcref{ch3:4.2}
  With the inverse pair map $\mathrm{Ginv}$ of \cref{thm:dc-ch3-3-7}, $C^1$ (hence
  continuous) on its open diffeomorphism image $U\ni(\varphi_p(p),\varphi_p(p))$ and
  fixing the center diagonal (\cref{lem:dc-ch3-4-2-invdiagonal}): for every $r>0$
  there is $\beta>0$ with $\overline{B_\beta(p)}$ inside the chart source and, for all
  $q_1,q_2\in\overline{B_\beta(p)}$, the joining chart velocity
  $v=(\mathrm{Ginv}(\varphi_p(q_1),\varphi_p(q_2)))_2$ obeys $\|v\|<r$.
\end{lemma}
\begin{proof}
  \sketch
  \leanok
  Apply \cref{lem:dc-ch3-4-2-pairsmall} to $\Psi(x)=(\mathrm{Ginv}(x))_2$. This
  $\Psi$ is continuous at the center (restriction of the continuous $\mathrm{Ginv}$
  to the open $U$, followed by the second projection) and vanishes there by
  \cref{lem:dc-ch3-4-2-invdiagonal}.
\end{proof}

\begin{lemma}[the interior of a joining geodesic stays inside the ball]
  \label{lem:dc-ch3-4-2-interior}
  \lean{Riemannian.Exponential.exists_forall_geodesic_dist_lt_of_admissible}
  \uses{lem:dc-ch3-4-2-nomax, lem:dc-ch3-4-2-bridge-ball, lem:dc-ch3-4-2-maxdeduction}
  \leanok
  \dcref{ch3:4.2}
  With the exponential-inverse package of \cref{lem:dc-ch3-4-2-nomax} and the bridge radius $\delta'$ of
  \cref{lem:dc-ch3-4-2-bridge-ball}: for every continuous intrinsic geodesic $\gamma$ on an open interval
  $(l,h)\supsetneq[0,1]$ whose endpoints lie within $\beta\le\delta'$ of $p$, which stays inside
  $B_{\delta'}(p)$ over $[0,1]$, and which is \cref{lem:dc-ch3-4-2-nomax}-admissible at every interior
  time (base reads into $V$, nonzero chart velocity in the flow's ball of initial conditions), the whole
  open arc stays strictly inside $B_\beta(p)$: $d(p,\gamma(t))<\beta$ for all $t\in(0,1)$.
\end{lemma}
\begin{proof}
  \sketch
  \leanok
  The distance $h(t)=d(p,\gamma(t))$ is continuous on $[0,1]$ with $h(0),h(1)\le\beta$. If $h$ had an
  interior local maximum at $t_0$, re-base $\sigma(s)=\gamma(t_0+s)$ (an affine reparametrisation of a
  geodesic is a geodesic): the metric$\leftrightarrow$radial bridge $F=d^2$ of
  \cref{lem:dc-ch3-4-2-bridge-ball} turns the local maximum of $h$ into a local maximum of the chart
  radial functional at $0$, which \cref{lem:dc-ch3-4-2-nomax} forbids. So $h$ has no interior local
  maximum, and \cref{lem:dc-ch3-4-2-maxdeduction} yields $h<\beta$ on $(0,1)$.
\end{proof}

\begin{lemma}[the chart Gram quadratic form is uniformly small near the center]
  \label{lem:dc-ch3-4-2-gramsmall}
  \lean{Riemannian.Exponential.exists_ball_forall_chartMetricInner_lt}
  \leanok
  \dcref{ch3:4.2}
  For every tolerance $\eta>0$ there is a radius $\rho>0$ such that the open chart ball
  $B_\rho(\varphi_p(p))$ lies in the chart target and, for every base position $y$ in it and every
  chart vector $w$ with $\|w\|<\rho$, the chart-Gram value $\langle w,w\rangle_{p,y}<\eta$.
\end{lemma}
\begin{proof}
  \sketch
  \leanok
  The functional $\Psi(y,w)=\langle w,w\rangle_{p,y}$ is jointly continuous (a finite sum of the
  continuous Gram components composed with the base, times the continuous chart coordinates of $w$)
  and vanishes at $(\varphi_p(p),0)$ (both vector slots are $0$), so $\{\Psi<\eta\}$ is a
  neighborhood of the center and contains a product ball.
\end{proof}

\begin{lemma}[interior arc in the ball from the minimizing property and conserved speed]
  \label{lem:dc-ch3-4-2-interior-minimizing}
  \lean{Riemannian.Exponential.exists_forall_minimizing_geodesic_interior_ball}
  \uses{lem:dc-ch3-4-2-interior, lem:dc-ch3-4-2-velbound, lem:dc-ch3-4-2-speedchart, lem:dc-ch3-4-2-gramsmall}
  \leanok
  \dcref{ch3:4.2}
  For every $p$ there are radii $\beta_0,\rho>0$ such that any minimizing geodesic $\gamma$
  (with $d(\gamma(s),\gamma(t))=|s-t|\,d(q_1,q_2)$) joining $q_1,q_2\in\overline{B_\beta(p)}$
  ($\beta\le\beta_0$), extended as a continuous intrinsic geodesic on an open window
  $(l,h)\supsetneq[0,1]$ with nonzero small initial chart velocity $w$ ($w\neq0$, $\|w\|<\rho$) and a
  chart-velocity derivative at every interior time, has its open arc strictly inside $B_\beta(p)$:
  $d(p,\gamma(t))<\beta$ for all $t\in(0,1)$.
\end{lemma}
\begin{proof}
  \sketch
  \leanok
  This discharges the admissibility of \cref{lem:dc-ch3-4-2-interior} without any flow-continuity
  (tube) plumbing. \emph{Base confinement} $d(p,\gamma(t))\le3\beta$ follows from the triangle
  inequality on the minimizing length ($d(\gamma(0),\gamma(t))=t\,d(q_1,q_2)\le2\beta$), which puts
  the base into every prescribed chart neighborhood of $p$ and inside $B_{\delta'}(p)$. The squared
  speed is conserved along the geodesic, and at $t=0$ equals the small chart-Gram value
  $\langle w,w\rangle$ (\cref{lem:dc-ch3-4-2-speedchart}); the coordinate-velocity bound
  \cref{lem:dc-ch3-4-2-velbound} together with the chart-Gram smallness \cref{lem:dc-ch3-4-2-gramsmall}
  then bounds the interior chart velocity $T^{-1}w_0$ inside the flow's ball of initial conditions,
  while the Gram lower bound keeps the conserved speed positive so that $w_0\neq0$.
  \cref{lem:dc-ch3-4-2-interior} then yields the strict interior confinement.
\end{proof}

\begin{lemma}[a geodesic segment has length equal to speed times elapsed parameter]
  \label{lem:dc-ch3-geodesic-pathlength}
  \lean{Riemannian.Geodesic.IsGeodesicOn.pathELength_eq}
  \uses{lem:dc-ch3-2-1-constspeed}
  \leanok
  \dcref{ch3:3.6}
  Let $\gamma$ be a continuous intrinsic geodesic on a connected open set of times containing
  $a$ and $b$. Then the length of the arc over $[a,b]$ is the constant speed times the elapsed
  parameter,
  \[
    \ell\big(\gamma|_{[a,b]}\big)=\sqrt{\langle\gamma'(a),\gamma'(a)\rangle_g}\,(b-a).
  \]
\end{lemma}
\begin{proof}
  \sketch
  \leanok
  The length is the integral of the intrinsic speed $\sqrt{\langle\gamma',\gamma'\rangle_g}$ over
  $[a,b]$. Along a geodesic the squared speed is constant (\cref{lem:dc-ch3-2-1-constspeed}), equal
  to its value at $a$, so the integrand is the constant $\sqrt{\langle\gamma'(a),\gamma'(a)\rangle_g}$
  and the integral is that constant times $b-a$.
\end{proof}

\begin{lemma}[the joining geodesic segment has length $\sqrt{\langle w,w\rangle}$, uniformly in the base]
  \label{lem:dc-ch3-4-2-segment-length}
  \lean{Riemannian.Geodesic.pathELength_uniform_flow_segment_Ioo}
  \uses{lem:dc-ch3-geodesic-pathlength, lem:dc-ch3-3-7-descent, lem:dc-ch3-3-7-descent-velocity, lem:dc-ch3-4-2-speedchart}
  \leanok
  \dcref{ch3:4.2}
  For the rescaled chart-$p$ spray segment $\gamma(s)=\varphi_p^{-1}\big((Z(y,T^{-1}w)(sT))_1\big)$
  joining $q=\varphi_p^{-1}(y)$ to a nearby point (the geodesic of \cref{thm:dc-ch3-3-7}), the length
  over $[0,1]$ equals the chart-Gram norm of the initial velocity read at the base $y$:
  \[
    \ell\big(\gamma|_{[0,1]}\big)=\sqrt{\langle w,w\rangle_{p,y}}.
  \]
  This is uniform in the base point $y$: the same flow $Z$ and time scale $T$ serve every $q$ in the
  neighborhood. It is the metric payload that \cref{thm:dc-ch3-3-7} otherwise omits.
\end{lemma}
\begin{proof}
  \sketch
  \leanok
  The segment is a continuous intrinsic geodesic on the open time window, so
  \cref{lem:dc-ch3-geodesic-pathlength} gives its length as
  $\sqrt{\langle\gamma'(0),\gamma'(0)\rangle_g}$. Its squared speed at $s=0$ is the chart-Gram value
  $\langle w,w\rangle$ read in the fixed chart at $p$ (\cref{lem:dc-ch3-4-2-speedchart}), at the base
  position $\varphi_p(\gamma(0))=y$ (the flow starts at its initial condition).
\end{proof}

\begin{lemma}[uniform radial upper bound]
  \label{lem:dc-ch3-4-2-segment-edist}
  \lean{Riemannian.Geodesic.edist_uniform_flow_segment_le}
  \uses{lem:dc-ch3-4-2-segment-length}
  \leanok
  \dcref{ch3:4.2}
  For the joining geodesic segment of \cref{thm:dc-ch3-3-7}, the distance from the base
  $q=\varphi_p^{-1}(y)$ to its endpoint is at most the chart-Gram norm of the initial velocity,
  \[
    d\big(q,\gamma(1)\big)\le\sqrt{\langle w,w\rangle_{p,y}},
  \]
  uniformly in the base point $y$.
\end{lemma}
\begin{proof}
  \sketch
  \leanok
  The distance between the endpoints of a $C^1$ curve is at most its length; combine this with the
  length identity \cref{lem:dc-ch3-4-2-segment-length}. This is the $\le$ half of the normal-ball
  minimizing property (\cref{prop:dc-ch3-3-6}), now available \emph{uniformly} in the base point $y$
  --- the direction that does not require a base-uniform Gauss estimate.
\end{proof}

\begin{lemma}[the joining geodesic carries the base-uniform metric payload]
  \label{lem:dc-ch3-4-2-join-metric}
  \lean{Riemannian.Exponential.exists_convex_join_ball}
  \uses{thm:dc-ch3-3-7, lem:dc-ch3-3-7-descent, lem:dc-ch3-4-2-segment-length, lem:dc-ch3-4-2-segment-edist, lem:dc-ch3-4-2-velsmall, lem:dc-ch3-geodesic-pathlength}
  \leanok
  \dcref{ch3:4.2}
  For every $p$ and every velocity tolerance $\rho>0$ there is $\beta>0$ such that any two points
  $q_1,q_2\in\overline{B_\beta(p)}$ are joined by a geodesic
  $\gamma(s)=\varphi_p^{-1}\big((Z(\varphi_p q_1,T^{-1}w)(sT))_1\big)$ on an open time window
  $(l,h)\supsetneq[0,1]$, with joining chart velocity $\|w\|<\rho$ (and $w\neq0$ unless $q_1=q_2$),
  whose length is arclength-proportional,
  $\ell(\gamma|_{[0,t]})=\sqrt{\langle w,w\rangle_{p,\varphi_p q_1}}\,t$, and whose endpoints obey the
  radial upper bound $d(q_1,q_2)\le\sqrt{\langle w,w\rangle_{p,\varphi_p q_1}}$. The chart velocity at
  every interior time and the geodesic equation are also available. This is the metric-laden
  totally-normal joining geodesic, uniform over the ball of base points.
\end{lemma}
\begin{proof}
  \sketch
  \leanok
  The totally-normal diffeomorphism \cref{thm:dc-ch3-3-7} joins $q_1,q_2$ by the flow segment and now
  re-exposes the raw spray-flow predicate with the same $(T,Z)$, so the descent
  \cref{lem:dc-ch3-3-7-descent} gives the geodesic on the open window and
  \cref{lem:dc-ch3-4-2-segment-length} its length $\sqrt{\langle w,w\rangle}$ (the constant-speed
  identity \cref{lem:dc-ch3-geodesic-pathlength}); the radial upper bound is
  \cref{lem:dc-ch3-4-2-segment-edist}. The joining velocity is uniformly small by
  \cref{lem:dc-ch3-4-2-velsmall}, and $w=0$ forces zero length hence $q_1=q_2$.
\end{proof}

\begin{lemma}[a short geodesic reaches the radial flow endpoint]
  \label{lem:dc-ch3-4-2-endpoint-id}
  \lean{Riemannian.Exponential.movingBase_geodesic_endpoint_eq_flow_reading}
  \uses{thm:dc-ch3-3-7, lem:dc-ch3-3-7-descent, def:dc-ch3-2-1}
  \leanok
  \dcref{ch3:3.6}
  Let $Z$ be a local flow of the coordinate geodesic spray at $p$ on a ball around the zero section
  (as in \cref{thm:dc-ch3-3-7}). Let $\gamma$ be a continuous intrinsic geodesic on an open window
  $(l,h)\supsetneq[0,1]$ with $\gamma(0)=q_1$ near $p$ and initial chart velocity $w$ small enough that
  $(\varphi_p q_1,T^{-1}w)$ lies in the flow ball. Then $\gamma$ coincides with the rescaled flow line
  $s\mapsto\varphi_p^{-1}\big((Z(\varphi_p q_1,T^{-1}w)(sT))_1\big)$ on the overlap window; in
  particular $\gamma(1)=\varphi_p^{-1}\big((Z(\varphi_p q_1,T^{-1}w)T)_1\big)$.
\end{lemma}
\begin{proof}
  \sketch
  \leanok
  The flow line is a continuous intrinsic geodesic on $(-(\varepsilon/T),\varepsilon/T)$
  (\cref{lem:dc-ch3-3-7-descent}) starting at $\varphi_p^{-1}(\varphi_p q_1)=q_1$ with chart velocity
  $w$ at time $0$ --- the same position and velocity as $\gamma$. On the asymmetric open overlap
  $(l,h)\cap(-(\varepsilon/T),\varepsilon/T)$, which contains $[0,1]$ because $h>1$ and
  $\varepsilon/T>1$, intrinsic geodesic uniqueness identifies the two curves; evaluating at $s=1$ gives
  the endpoint identity. Using the asymmetric overlap avoids the symmetric window $(-a,a)$, so the raw
  $l<0<1<h$ data suffices.
\end{proof}

\begin{lemma}[base-uniform competitor bound with escape]
  \label{lem:dc-ch3-4-2-competitor-lb}
  \lean{Riemannian.Exponential.exists_movingBase_competitor_le_pathELength}
  \uses{thm:dc-ch3-3-7, lem:dc-ch3-3-6-radius, lem:dc-ch3-3-6-competitor, lem:dc-ch3-4-2-gramsmall}
  \leanok
  \dcref{ch3:3.6}
  For every $p$ there are $\beta,\rho'>0$, a time scale $T$ and a spray flow $Z$ such that for every
  base point $q\in\overline{B_\beta(p)}$ and every chart velocity $\|w\|<\rho'$, every $C^1$ curve
  $\sigma$ from $q$ to the radial flow endpoint $\varphi_p^{-1}\big((Z(\varphi_p q,T^{-1}w)T)_1\big)$
  has length at least $\sqrt{\langle w,w\rangle_{p,\varphi_p q}}$. This is the Proposition 3.6
  competitor inequality (\cref{lem:dc-ch3-3-6-competitor}), stated over the whole ball of centres $q$
  at once.
\end{lemma}
\begin{proof}
  \sketch
  \leanok
  Fix $q$ and read the totally-normal diffeomorphism (\cref{thm:dc-ch3-3-7}) at $y=\varphi_p q$: the
  slice $f_y(z)=(Z(y,T^{-1}z)T)_1$ is a $C^1$ diffeomorphism of a ball onto a chart region, with
  inverse $\mathrm{finv}_y=(\mathrm{Ginv}(y,\cdot))_2$. If $\sigma$ stays in the closed $r'$-region
  $\{\varphi_p^{-1}(f_y z):\|z\|\le r'\}$, the (moving-base) Gauss radius comparison
  (\cref{lem:dc-ch3-3-6-radius}) applied to the polar lift $\mathrm{finv}_y\circ\varphi_p\circ\sigma$
  gives $\sqrt{\langle\mathrm{finv}_y(\varphi_p\sigma(1)),\cdot\rangle_y}\le\ell(\sigma)$; since the
  endpoint is $\varphi_p^{-1}(f_y w)$ and $\mathrm{finv}_y(f_y w)=w$, this is
  $\sqrt{\langle w,w\rangle_y}$. If instead $\sigma$ leaves the region, at the first exit time it
  crosses the coordinate $r'$-sphere at some $z_0$ with $\|z_0\|=r'$; the confined comparison up to
  that time plus the coordinate-norm bound $\|z_0\|^2\le c\,\langle z_0,z_0\rangle_y$ gives
  $\ell(\sigma)\ge r'/\sqrt c>\sqrt{\langle w,w\rangle_y}$, the last inequality by the base-uniform
  smallness (\cref{lem:dc-ch3-4-2-gramsmall}) for $\|w\|<\rho'$. The region is open (the slice inverse
  is a homeomorphism) and its closure is compact, so the first exit time is well defined.
\end{proof}

\begin{lemma}[base-uniform lower bound: short geodesics realize the distance]
  \label{lem:dc-ch3-4-2-lowerbound}
  \lean{Riemannian.Exponential.exists_movingBase_prop36_lower_bound}
  \uses{lem:dc-ch3-4-2-competitor-lb, lem:dc-ch3-4-2-endpoint-id, lem:dc-ch3-geodesic-pathlength, lem:dc-ch3-4-2-speedchart, prop:dc-ch3-3-6}
  \leanok
  \dcref{ch3:3.6}
  For every $p$ there are $\beta_H,\rho_H>0$ such that every geodesic $\gamma$ on an open window
  $(l,h)\supsetneq[0,1]$ joining $q_1,q_2$ within $\beta_H$ of $p$, with small initial chart velocity
  $\|w\|<\rho_H$ and constant-speed length rate $\ell$ (i.e.\ $\ell(\gamma|_{[0,t]})=\ell t$), realizes
  the distance between its endpoints: $\ell\le d(q_1,q_2)$. This is the
  base-uniform form of the distance-realization result in
  \cref{prop:dc-ch3-3-6}.
\end{lemma}
\begin{proof}
  \sketch
  \leanok
  Constant speed (\cref{lem:dc-ch3-geodesic-pathlength}), read in the fixed chart at $p$
  (\cref{lem:dc-ch3-4-2-speedchart}), forces $\ell=\sqrt{\langle w,w\rangle_{p,\varphi_p q_1}}$. The
  flow-reading endpoint identification (\cref{lem:dc-ch3-4-2-endpoint-id}) pins $q_2$ as the radial
  endpoint of $w$, so the base-uniform competitor bound (\cref{lem:dc-ch3-4-2-competitor-lb}) applies:
  every $C^1$ curve from $q_1$ to $q_2$ has length at least $\sqrt{\langle w,w\rangle}=\ell$. Since the
  distance is the infimum of $C^1$-curve lengths (the metric is the Riemannian distance),
  $d(q_1,q_2)\ge\ell$.
\end{proof}

\begin{lemma}[minimizing geodesics with interior in the ball, modulo the lower bound]
  \label{lem:dc-ch3-4-2-minimizing-interior}
  \lean{Riemannian.Exponential.exists_minimizing_interior_ball_of_lower_bound}
  \uses{lem:dc-ch3-4-2-join-metric, lem:dc-ch3-4-2-interior-minimizing, prop:dc-ch3-3-6, cor:dc-ch3-3-9, def:dc-ch3-4-2-stronglyconvex, lem:dc-ch3-4-2-lowerbound}
  \leanok
  \dcref{ch3:4.2}
  Assume the following base-uniform lower-bound property: there are $\beta_H,\rho_H>0$ such that every
  \emph{geodesic} $\gamma$ (on an open window $\supsetneq[0,1]$) joining $q_1,q_2$ within $\beta_H$ of
  $p$, with \emph{small} initial chart velocity ($\|w\|<\rho_H$) and constant-speed length
  $\ell(\gamma|_{[0,t]})=\ell t$, satisfies $\ell\le d(q_1,q_2)$ (a \emph{short} geodesic near $p$
  \emph{realizes} the distance between its endpoints --- \cref{prop:dc-ch3-3-6}, phrased
  base-uniformly). The geodesic and smallness hypotheses are essential: for an \emph{arbitrary}
  constant-speed curve one always has $d(q_1,q_2)\le\ell$, so a bare $\ell\le d(q_1,q_2)$ would force
  every constant-speed curve minimizing --- false on any manifold of dimension $\ge2$ (a constant-speed
  loop through $q_1=q_2$ has $\ell>0=d$). Then there is
  $\beta>0$ such that any $q_1,q_2\in\overline{B_\beta(p)}$ are joined by a \emph{minimizing} geodesic
  ($d(\gamma(s),\gamma(t))=|s-t|\,d(q_1,q_2)$) whose open arc lies in $\overline{B_\beta(p)}$. This is
  the strong convexity (\cref{def:dc-ch3-4-2-stronglyconvex}) \emph{minus} the uniqueness
  clause, under this lower-bound hypothesis.
\end{lemma}
\begin{proof}
  \sketch
  \leanok
  The joining geodesic and its length payload $\ell=\sqrt{\langle w,w\rangle}$ come from
  \cref{lem:dc-ch3-4-2-join-metric}; the lower-bound hypothesis supplies the matching lower bound, so
  $d(q_1,q_2)=\ell$,
  and the arclength distance-realization of \cref{cor:dc-ch3-3-9}'s segment machinery propagates
  $d(\gamma(s),\gamma(t))=|s-t|\,d(q_1,q_2)$ from the endpoints. The interior-arc clause is
  \cref{lem:dc-ch3-4-2-interior-minimizing}; the degenerate $q_1=q_2$ case is the constant geodesic.

\end{proof}

\begin{lemma}[strong convexity of the closed ball, modulo the lower bound and uniqueness]
  \label{lem:dc-ch3-4-2-stronglyconvex-reduction}
  \lean{Riemannian.Exponential.exists_stronglyConvex_closedBall_of_lower_bound}
  \uses{lem:dc-ch3-4-2-minimizing-interior, def:dc-ch3-4-2-stronglyconvex, thm:dc-ch3-3-7, prop:dc-ch3-3-6}
  \leanok
  \dcref{ch3:4.2}
  Assume the base-uniform lower-bound property from
  \cref{lem:dc-ch3-4-2-minimizing-interior} and local uniqueness of minimizing
  geodesics near $p$: there is $\beta_U>0$ such that
  any two constant-speed distance-realizing geodesics joining the same $q_1,q_2$ within $\beta_U$ of
  $p$ coincide on $[0,1]$ (the reading of the injectivity of $\exp_{q_1}$). Then there is
  $\beta>0$ such that the \emph{closed} ball $\overline{B_\beta(p)}$ is strongly convex
  (\cref{def:dc-ch3-4-2-stronglyconvex}). This is the whole of the Proposition 4.2 for the
  closed ball.
\end{lemma}
\begin{proof}
  \sketch
  \leanok
  The existence, minimizing and interior clauses are \cref{lem:dc-ch3-4-2-minimizing-interior}
  (shrinkable in the radius, so the strong-convexity radius $\beta=\min(B_0,\beta_U)$ couples
  correctly); the closure of a closed ball is itself, so the quantifier over $\overline{S}$ reads on
  $\overline{B_\beta(p)}$; local uniqueness supplies the uniqueness clause. The closed ball is used
  because the literal open-ball statement fails at a boundary diagonal $q_1=q_2\in\partial B_\beta(p)$
  (constant geodesic, interior $\{q_1\}\not\subset$ open ball).
\end{proof}

\begin{lemma}[minimizing geodesics with interior in the ball --- unconditional]
  \label{lem:dc-ch3-4-2-minimizing-interior-uncond}
  \lean{Riemannian.Exponential.exists_minimizing_interior_ball}
  \uses{lem:dc-ch3-4-2-minimizing-interior, lem:dc-ch3-4-2-lowerbound}
  \leanok
  \dcref{ch3:4.2}
  There is $\beta>0$ such that any $q_1,q_2\in\overline{B_\beta(p)}$ are joined by a \emph{minimizing}
  geodesic ($d(\gamma(s),\gamma(t))=|s-t|\,d(q_1,q_2)$) whose open arc lies in $\overline{B_\beta(p)}$.
  The lower-bound hypothesis in \cref{lem:dc-ch3-4-2-minimizing-interior} is supplied by
  \cref{lem:dc-ch3-4-2-lowerbound}; hence this gives the existence, minimizing, and interior
  clauses of strong convexity (\cref{def:dc-ch3-4-2-stronglyconvex}).
\end{lemma}
\begin{proof}
  \sketch
  \leanok
  Apply the lower-bound lemma \cref{lem:dc-ch3-4-2-lowerbound} to the reduction
  \cref{lem:dc-ch3-4-2-minimizing-interior}.
\end{proof}

\begin{lemma}[strong convexity of the closed ball, reduced to uniqueness alone]
  \label{lem:dc-ch3-4-2-stronglyconvex-uniq}
  \lean{Riemannian.Exponential.exists_stronglyConvex_closedBall_of_uniq}
  \uses{lem:dc-ch3-4-2-stronglyconvex-reduction, lem:dc-ch3-4-2-lowerbound, lem:dc-ch3-4-2-minimizing-interior-uncond}
  \leanok
  \dcref{ch3:4.2}
  Assume only local uniqueness of minimizing geodesics near $p$ (any two
  constant-speed distance-realizing geodesics joining the same $q_1,q_2$ within $\beta_U$ of $p$
  coincide on $[0,1]$). Then there is $\beta>0$ such that the \emph{closed} ball
  $\overline{B_\beta(p)}$ is strongly convex (\cref{def:dc-ch3-4-2-stronglyconvex}). This is
  \cref{lem:dc-ch3-4-2-stronglyconvex-reduction} together with
  \cref{lem:dc-ch3-4-2-lowerbound}.
\end{lemma}
\begin{proof}
  \sketch
  \leanok
  Apply \cref{lem:dc-ch3-4-2-lowerbound} to
  \cref{lem:dc-ch3-4-2-stronglyconvex-reduction}.
\end{proof}

\begin{lemma}[moving-base geodesic coincides with its flow reading on the overlap window]
  \label{lem:dc-ch3-4-2-eqon-flow-reading}
  \lean{Riemannian.Exponential.movingBase_geodesic_eqOn_flow_reading}
  \uses{lem:dc-ch3-4-2-endpoint-id, thm:dc-ch3-3-7, lem:dc-ch3-3-4}
  \leanok
  \dcref{ch3:4.2}
  Let $Z$ be the geodesic-spray flow of \cref{thm:dc-ch3-3-7}. A continuous intrinsic geodesic
  $\gamma$ on an open window $(l,h)\ni0$ whose foot $\gamma(0)=q_1$ lies in the chart at $p$ and whose
  chart-$p$ coordinate velocity at $0$ is $w$ \emph{coincides} with the flow-reading geodesic
  $s\mapsto\varphi_p^{-1}\bigl((Z(\varphi_p q_1,T^{-1}w)(sT))_1\bigr)$ on the overlap window
  $(l,h)\cap(-(\varepsilon/T),\varepsilon/T)$.
\end{lemma}
\begin{proof}
  \sketch
  \leanok
  Both curves are geodesics on the overlap window and share the position $q_1$ and the chart-$p$
  velocity $w$ at $0$; intrinsic uniqueness (\cref{lem:dc-ch3-3-4}) identifies them on the whole
  window. This is the \emph{EqOn} strengthening of the single-endpoint identification
  \cref{lem:dc-ch3-4-2-endpoint-id}, freed of the $1<h$ restriction (only $l<0<h$ is used).
\end{proof}

\begin{lemma}[base-uniform geodesic-uniqueness engine]
  \label{lem:dc-ch3-4-2-uniqueness-engine}
  \lean{Riemannian.Exponential.exists_movingBase_geodesic_uniqueness}
  \uses{lem:dc-ch3-4-2-eqon-flow-reading, thm:dc-ch3-3-7}
  \leanok
  \dcref{ch3:4.2}
  For any $p$ there are radii $\beta,\rho_w>0$ such that any two continuous intrinsic geodesics
  $\gamma_1,\gamma_2$ on open windows $\supsetneq[0,1]$, both joining $q_1\in\overline{B_\beta(p)}$ to
  a common $q_2\in\overline{B_\beta(p)}$, with initial chart-$p$ coordinate velocities $w_1,w_2$ of
  norm $<\rho_w$, coincide on $[0,1]$. Thus sufficiently short geodesics with
  the same endpoints and initial chart velocity are unique in this neighborhood.
\end{lemma}
\begin{proof}
  \sketch
  \leanok
  Each $\gamma_i$ coincides with the flow-reading geodesic $g_{w_i}$ on a window containing $[0,1]$
  (\cref{lem:dc-ch3-4-2-eqon-flow-reading}), so $g_{w_i}(1)=\gamma_i(1)=q_2$. The unique-velocity
  clause of the totally-normal $C^1$ diffeomorphism (\cref{thm:dc-ch3-3-7}) --- the injectivity
  of $\exp_{q_1}$ --- forces $w_1=w_2$, whence $g_{w_1}=g_{w_2}$ and therefore $\gamma_1=\gamma_2$ on
  $[0,1]$.
\end{proof}

\begin{lemma}[a closed-interval geodesic extends to an open time window]
  \label{lem:dc-ch3-4-2-window}
  \lean{Riemannian.Exponential.exists_geodesic_window_of_isGeodesicOn_Icc}
  \uses{prop:dc-ch3-2-5, lem:dc-ch3-3-4}
  \leanok
  \dcref{ch3:4.2}
  A continuous intrinsic geodesic $\alpha$ on the \emph{closed} interval $[0,1]$ extends to a
  continuous intrinsic geodesic on an \emph{open} window $(l,h)$ with $l<0<1<h$, agreeing with
  $\alpha$ on $[0,1]$.
\end{lemma}
\begin{proof}
  \sketch
  \leanok
  The endpoint values $\alpha(0),\alpha(1)$ are known points, not Cauchy limits, so a
  completeness-free forward prolongation of a geodesic on $(a,b)$ --- given the (known) endpoint
  limit --- applies at each end: once forward past $1$ with limit $\alpha(1)$, and once, after the
  time reversal $t\mapsto\alpha(-t)$, forward past $0$ with limit $\alpha(0)$. Each prolongation
  solves the geodesic ODE at the known endpoint by local existence (\cref{prop:dc-ch3-2-5}). The two
  one-sided extensions are glued at the interior time $1/2$, where they both equal $\alpha$; they
  agree there because a geodesic is locally determined by its position and velocity
  (\cref{lem:dc-ch3-3-4}).
\end{proof}

\begin{lemma}[a minimizing geodesic has speed equal to the distance it realizes]
  \label{lem:dc-ch3-4-2-speed-dist}
  \lean{Riemannian.Exponential.sqrt_speedSq_eq_dist_of_minimizing}
  \uses{prop:dc-ch3-3-6, lem:dc-ch3-3-4}
  \leanok
  \dcref{ch3:4.2}
  Let $\gamma$ be a continuous geodesic on an open window $(l,h)$ with $l<0<1<h$, joining
  $q_1=\gamma(0)$ to $q_2=\gamma(1)$ and realizing the distance on $[0,1]$
  (i.e. $d(\gamma(s),\gamma(t))=|s-t|\,d(q_1,q_2)$). Then its constant speed equals that distance:
  $\sqrt{\langle\gamma'(0),\gamma'(0)\rangle_g}=d(q_1,q_2)$.
\end{lemma}
\begin{proof}
  \sketch
  \leanok
  The inequality $d(q_1,q_2)\le\sqrt{S}$ (with $S$ the squared speed) is the length bound: a geodesic
  is at least as long as the chord between its endpoints. For the reverse, rescale the velocity by a
  small $\kappa>0$ so that the exponential ray $\eta(s)=\exp_{q_1}(s\,\kappa w)$ is defined; it is a
  geodesic through $q_1$ with the same chart velocity at $0$ as the rescaled curve
  $s\mapsto\gamma(\kappa s)$, so the two coincide near $0$ by intrinsic uniqueness
  (\cref{lem:dc-ch3-3-4}). The radial isometry of the normal ball (\cref{prop:dc-ch3-3-6}) then reads
  $d(q_1,\gamma(\kappa s))=s\,\kappa\sqrt{S}$ for small $s>0$, while the distance-realizing hypothesis
  reads the same distance as $\kappa s\,d(q_1,q_2)$; cancelling $\kappa s>0$ gives
  $\sqrt{S}=d(q_1,q_2)$.
\end{proof}

\begin{proposition}[convex neighborhoods]
  \label{prop:dc-ch3-4-2}
  \lean{Riemannian.Exponential.exists_stronglyConvex_closedBall}
  \uses{thm:dc-ch3-3-7, lem:dc-ch3-4-1, lem:dc-ch3-4-1-strict, lem:dc-ch3-4-2-rebase, lem:dc-ch3-4-2-window, lem:dc-ch3-4-2-speed-dist, lem:dc-ch3-4-2-nomax, lem:dc-ch3-4-2-bridge, lem:dc-ch3-4-2-bridge-ball, lem:dc-ch3-4-2-maxdeduction, lem:dc-ch3-4-2-speedchart, lem:dc-ch3-4-2-velbound, lem:dc-ch3-4-2-interior, lem:dc-ch3-4-2-interior-minimizing, lem:dc-ch3-4-2-gramsmall, lem:dc-ch3-4-2-invdiagonal, lem:dc-ch3-4-2-pairsmall, lem:dc-ch3-4-2-velsmall, lem:dc-ch3-4-2-segment-length, lem:dc-ch3-4-2-segment-edist, lem:dc-ch3-4-2-join-metric, lem:dc-ch3-4-2-minimizing-interior, lem:dc-ch3-4-2-stronglyconvex-reduction, lem:dc-ch3-4-2-lowerbound, lem:dc-ch3-4-2-minimizing-interior-uncond, lem:dc-ch3-4-2-stronglyconvex-uniq, lem:dc-ch3-4-2-endpoint-id, lem:dc-ch3-4-2-competitor-lb, lem:dc-ch3-4-2-eqon-flow-reading, lem:dc-ch3-4-2-uniqueness-engine, def:dc-ch3-4-2-stronglyconvex}
  \leanok
  \dcref{ch3:4.2}
  For any $p\in M$ there exists $\beta>0$ such that the closed geodesic ball $\overline{B_\beta(p)}$ is
  strongly convex (\cref{def:dc-ch3-4-2-stronglyconvex}). (The closed ball is used deliberately: the
  literal open-ball form is unsatisfiable at a boundary diagonal, see below.)
\end{proposition}
\begin{proof}
  \sketch
  \leanok
  Let $c$ be as in \cref{lem:dc-ch3-4-1}. Choose $\delta>0$ and $W$ as in \cref{thm:dc-ch3-3-7}
  with $\delta<\frac{c}{2}$, and take $\beta<\delta$ with $\overline{B_\beta(p)}\subset W$. Let
  $q_1,q_2\in\overline{B_\beta(p)}$ and let $\gamma$ be the unique geodesic of length
  $<2\delta<c$ joining them; then $\gamma\subset B_c(p)$. If the interior of
  $\gamma$ is not contained in $B_\beta(p)$, there is an interior point $m$ where the
  maximum distance $r$ from $p$ to $\gamma$ is attained; near $m$ the points of
  $\gamma$ remain in $\overline{B_r(p)}$. Since $m\in B_c(p)$ this contradicts
  \cref{lem:dc-ch3-4-1-strict}: at $m$ the tangency $\partial F/\partial t(0)=0$ holds
  (Gauss lemma) and the strict separation $F(0)<F(s)$ for $s\neq0$ forces points of
  $\gamma$ near $m$ \emph{strictly} outside $\overline{B_r(p)}$, contradicting the
  maximum. Hence the interior of $\gamma$ lies in $B_\beta(p)$, proving the proposition.
\end{proof}

\begin{definition}[piecewise parallel transport]
  \label{def:dc-ch3-3-1-piecewise-transport}
  \dcref{ch3:3.1}
  \uses{def:dc-ch3-3-1}
  \lean{Riemannian.Geodesic.IsPiecewiseParallelAlong, Riemannian.Geodesic.piecewiseTransport,
    Riemannian.Geodesic.Segmentation.ofLocal,
    Riemannian.Geodesic.local_of_segmentation,
    Riemannian.Geodesic.exists_isPiecewiseParallelAlong_of_local}
  \leanok
  Parallel transport along a piecewise differentiable curve is obtained by
  transporting a single own-foot field in succession on each differentiable
  segment of a finite partition, with the values agreeing at every vertex.  The
  endpoint transport is the relation between the initial and terminal values of
  such a field.
\end{definition}

\begin{definition}[vertex angle]
  \label{def:dc-ch3-3-1-vertex-angle}
  \dcref{ch3:3.1}
  \uses{def:dc-ch3-3-1}
  \lean{Riemannian.Geodesic.leftVelocity, Riemannian.Geodesic.rightVelocity, Riemannian.Geodesic.vertexAngle,
    Riemannian.Geodesic.HasLeftVelocity, Riemannian.Geodesic.HasRightVelocity,
    Riemannian.Geodesic.hasLeftVelocity_leftVelocity,
    Riemannian.Geodesic.hasRightVelocity_rightVelocity,
    Riemannian.Geodesic.hasLeftVelocity_of_segmentation,
    Riemannian.Geodesic.hasRightVelocity_of_segmentation,
    Riemannian.Geodesic.hasVelocities_of_segmentation,
    Riemannian.RiemannianMetric.angle_eq_pi_iff}
  \leanok
  At an interior vertex $c(t_i)$, the \emph{vertex angle} is the angle between the outgoing
  one-sided velocity $\lim_{t\to t_i^+}c'(t)$ and the incoming one-sided velocity
  $\lim_{t\to t_i^-}c'(t)$, measured by the Riemannian metric.
\end{definition}

\begin{definition}[formal parametrized-surface interface]
  \label{def:dc-ch3-3-3-formal}
  \dcref{ch3:3.3}
  \lean{Riemannian.Geodesic.IsParametrizedSurface,
    Riemannian.Geodesic.IsParametrizedSurfaceOfOrder,
    Riemannian.Geodesic.IsSmoothParametrizedSurface,
    Riemannian.Geodesic.IsSmoothExtendedParametrizedSurface,
    Riemannian.Geodesic.IsExtendedParametrizedSurfaceOfOrder.regularOn,
    Riemannian.Geodesic.IsExtendedParametrizedSurfaceOfOrder.mdifferentiableAt,
    Riemannian.Geodesic.ParametrizedSurface.IsVectorFieldAlong,
    Riemannian.Geodesic.ParametrizedSurface.partialU,
    Riemannian.Geodesic.ParametrizedSurface.partialV,
    Riemannian.Geodesic.IsParametrizedSurfaceOfOrder.surfaceSliceU_contMDiffOn,
    Riemannian.Geodesic.IsParametrizedSurfaceOfOrder.surfaceSliceV_contMDiffOn,
    Riemannian.Geodesic.IsExtendedParametrizedSurfaceOfOrder.partialU_eq_slice_mfderiv,
    Riemannian.Geodesic.IsExtendedParametrizedSurfaceOfOrder.partialV_eq_slice_mfderiv,
    Riemannian.Geodesic.ParametrizedSurface.partialU_eq_slice_mfderiv,
    Riemannian.Geodesic.ParametrizedSurface.partialV_eq_slice_mfderiv,
    Riemannian.Geodesic.IsSmoothParametrizedSurfaceWithBoundary,
    Riemannian.Geodesic.IsParametrizedSurfaceWithBoundary.surfaceDomain,
    Riemannian.Geodesic.PiecewiseC1BoundaryParametrization.interior_vertex_nonPi,
    Riemannian.Geodesic.PiecewiseC1BoundaryParametrization.closing_vertex_nonPi,
    Riemannian.Geodesic.IsNonPiAngle}
  \leanok
  The formal interface consists of the order-graded map predicate on $A$,
  the connected/open-core domain record, and the open-neighborhood extension
  record.  It also supplies the $C^r$ slice restrictions and identifies the
  ambient coordinate partials with the corresponding slice derivatives.  The
  The full boundary-data refinement is supplied by
  `IsParametrizedSurfaceWithBoundary`; its `Nonempty` boundary component keeps
  this predicate proof-irrelevant while exposing the finite parametrization as
  a separate structure.
\end{definition}

\begin{lemma}[smooth dependence of a coordinate flow]
  \label{lem:dc-ch3-2-2-cinfty}
  \dcref{ch3:2.2}
  \uses{lem:dc-ch3-2-2-flow,lem:dc-ch3-2-2-c1dep}
  \lean{Riemannian.GenericFlow.exists_local_flow_contDiffAt,
    Riemannian.GenericFlow.exists_local_flow_joint_contDiff}
  \leanok
  Let $F$ be a finite-dimensional Banach space and let $f:F\to F$ be a globally
  $C^\infty$ autonomous vector field.  At every $z_0\in F$ there are a uniform
  initial-value ball, a fixed time interval, and a family of solutions of
  $z'=f(z)$ which is Lipschitz in the initial value.  The curve-space family and,
  by the autonomous time-state augmentation, its endpoint map are $C^\infty$;
  in particular the solution map is jointly smooth in initial value and time on
  a smaller window.  This is the coordinate-space smooth-dependence core of the
  local-flow theorem; transporting it through a manifold chart remains a
  separate step.
\end{lemma}
\begin{proof}
  \leanok
  Apply Picard--Lindelöf on a ball around $z_0$ to obtain the uniform solution
  family and its Lipschitz estimate.  A compact-ball bound on $\|f'\|$ permits
  choosing $T$ with $T\sup\|f'\|<1$.  The Picard residual is then $C^\infty$ and
  its curve-variable derivative is invertible by the Neumann series, so the
  smooth implicit-function theorem gives the stated $C^\infty$ dependence.
\end{proof}

\begin{lemma}[chart-local uniform geodesic family]
  \label{lem:dc-ch3-2-5-chart-family}
  \dcref{ch3:2.5}
  \lean{Riemannian.GeodesicLocal.geodesic_local_existence,
    Riemannian.GeodesicLocal.geodesic_local_existence_tangentBundle,
    Riemannian.GeodesicLocal.geodesic_local_uniqueness}
  \leanok
  Fix $p\in M$ and a model-space representative $v\in E$ of an initial
  tangent vector in the chart at $p$.  There are $\varepsilon>0$, neighborhoods
  $V_1\in\mathcal N(p)$ and $V_2\in\mathcal N(v)$, and a family
  $c:M\to E\to\mathbb R\to M$ such that, for $q\in V_1$ and $w\in V_2$,
  the curve $c(q,w,\cdot)$ stays in the fixed chart at $p$, solves the
  first-order chart geodesic system on $(-\varepsilon,\varepsilon)$, and has
  $c(q,w,0)=q$ and chart velocity $w$.  Its chart reading is jointly
  $C^\infty$ in $(q,w,t)$.  Two such chart geodesics on open
  order-connected time sets with the same position and chart velocity at one
  common time agree on their overlap.  At the zero tangent vector over $p$,
  the foot and fiber restrictions assemble into a single neighborhood
  $\mathcal U\subset TM$; every $z\in\mathcal U$ seeds one of these curves at
  its foot with the fiber coordinate of $z$ as initial velocity.

  The further refinement of $\mathcal U$ to the explicit base-uniform metric
  disc in $\cref{prop:dc-ch3-2-5}$ is supplied by
  $\cref{lem:dc-ch3-2-5-metric-ball,lem:dc-ch3-2-5-metric-family}$.
\end{lemma}

\begin{lemma}[uniform metric ball in the chart fiber]
  \label{lem:dc-ch3-2-5-metric-ball}
  \dcref{ch3:2.5}
  \uses{lem:dc-ch7-2-8-gram-bound}
  \lean{Riemannian.GeodesicLocal.exists_chartFiberCoord_mem_of_metricInner_lt}
  \leanok
  Let $g$ be a Riemannian metric, let $p\in M$, and let $W$ be a neighborhood of
  $0$ in the model fiber of a chart at $p$.  There are a neighborhood $V$ of $p$
  and a number $\rho>0$ such that, for every $q\in V$ and
  $v\in T_qM$ with $\langle v,v\rangle_q<\rho$, the chart-$p$ coordinate of $v$
  belongs to $W$.
\end{lemma}
\begin{proof}
  \leanok
  Choose $\varepsilon>0$ with the model ball $B_\varepsilon(0)\subset W$.
  The positive-definite Gram form of $g$ in the chart at $p$ dominates a fixed
  positive multiple of the squared model norm on a neighborhood $Y$ of the
  chart image of $p$: there is $c>0$ with
  $\lVert w\rVert^2\le c\,G_y(w,w)$ for $y\in Y$.  Shrink the base neighborhood
  to $V$ so that every $q\in V$ has chart image in $Y$, and put
  $\rho=\varepsilon^2/c$.  If $v\in T_qM$ has squared norm below $\rho$, its
  chart coordinate $w$ satisfies
  $\lVert w\rVert^2\le c\langle v,v\rangle_q<\varepsilon^2$, hence
  $w\in B_\varepsilon(0)\subset W$.
\end{proof}

\begin{lemma}[uniform local geodesic family on a metric velocity ball]
  \label{lem:dc-ch3-2-5-metric-family}
  \dcref{ch3:2.5}
  \uses{lem:dc-ch3-2-5-chart-family,lem:dc-ch3-2-5-metric-ball}
  \lean{Riemannian.GeodesicLocal.geodesic_local_existence_metricBall}
  \leanok
  For every $p\in M$ there are $\delta>0$, a neighborhood $V$ of $p$,
  a model-fiber neighborhood $W$ of $0$, a number $\rho>0$, and a family
  $c:M\times E\times\mathbb R\to M$ such that, whenever $q\in V$ and
  $v\in T_qM$ has $\langle v,v\rangle_q<\rho$, the curve obtained from
  $c(q,\operatorname{coord}_p(v),t)$ is a geodesic on $(-\delta,\delta)$,
  starts at $q$ with the prescribed chart velocity, and is the unique chart
  geodesic with those initial data.  In chart coordinates the family is
  $C^\infty$ on $\bigl(\varphi_p(V)\times W\bigr)\times(-\delta,\delta)$.
\end{lemma}
\begin{proof}
  \leanok
  Apply the chart-local uniform family of
  \cref{lem:dc-ch3-2-5-chart-family} at the zero velocity.  Its model velocity
  neighborhood is a neighborhood of $0$, so
  \cref{lem:dc-ch3-2-5-metric-ball} supplies a smaller base neighborhood and a
  fixed $\rho$ whose metric ball is carried into that model neighborhood.
  Restricting the family and its smoothness statement to the smaller product
  gives the asserted existence and regularity.  If another chart geodesic has
  the same initial position and chart velocity, the phase lifts solve the same
  first-order geodesic system; uniqueness on the connected interval
  $(-\delta,\delta)$ makes the two curves equal.
\end{proof}

\begin{lemma}[$\exp_p$ is a $C^\infty$ diffeomorphism near the origin]
  \label{lem:dc-ch3-2-9-cinfty}
  \dcref{ch3:2.9}
  \uses{lem:dc-ch3-2-9-c2diffeo, lem:dc-ch3-2-9-invertible,
    lem:dc-ch3-2-9-c1diffeo}
  \lean{Riemannian.Exponential.exists_infty_local_diffeomorphism_expMap}
  \leanok
  There is $\varepsilon>0$ such that $\exp_p$ is injective on
  $B_\varepsilon(0)\subset T_pM$, has open image, and its chart reading and
  inverse are $C^\infty$ on the corresponding open images. In particular the
  exponential is a diffeomorphism of a neighbourhood of the origin onto an
  open neighbourhood of $p$.
\end{lemma}
\begin{proof}
  \leanok
  The $C^\infty$ flow-dependence theorem gives a $C^\infty$ chart reading of
  $\exp_p$ on a ball. Shrink this ball inside the strict-derivative and
  injectivity balls. The derivative is then an equivalence at every point, so
  the inverse function theorem makes the chart image open. Define the inverse
  by the unique preimage; at each image point it agrees locally with the
  $C^\infty$ inverse supplied by the inverse function theorem, hence is
  $C^\infty$ on the whole image.
\end{proof}

\begin{lemma}[piecewise curves have length at least distance]
  \label{lem:dc-ch3-3-1-edist}
  \dcref{ch3:3.1}
  \uses{def:dc-ch3-3-1}
  \lean{Riemannian.Geodesic.IsPiecewiseDifferentiableCurve.edist_le_pathELength}
  \leanok
  If $c:[a,b]\to M$ is piecewise differentiable and the ambient metric is the
  Riemannian distance, then
  $$
    d(c(a),c(b))\leq \ell(c|_{[a,b]}).
  $$
  The estimate is applied on each $C^1$ piece and combined across the vertices
  using the triangle inequality and additivity of path length.
\end{lemma}

\chapter{Curvature}
\label{chap:dc-curvature}

\section{Curvature}

\begin{definition}[curvature]
  \label{def:dc-ch4-2-1}
  \lean{Riemannian.AffineConnection.curvature}
  \uses{def:dc-ch2-2-1}
  \leanok
  \dcref{ch4:2.1}
  The \emph{curvature} $R$ of a Riemannian manifold $M$ is a correspondence that
  associates to every pair $X,Y\in\mathcal{X}(M)$ a mapping
  $R(X,Y):\mathcal{X}(M)\to\mathcal{X}(M)$ given by
  $$
  R(X,Y)Z=\nabla_Y\nabla_X Z-\nabla_X\nabla_Y Z+\nabla_{[X,Y]}Z,\quad Z\in\mathcal{X}(M),
  $$
  where $\nabla$ is the Riemannian connection of $M$.
  For $M=\mathbb{R}^n$, the coordinate expression
  $\nabla_X Z=(Xz_1,\dots,Xz_n)$ gives $R(X,Y)Z=0$; thus $R$ measures the
  deviation from the Euclidean connection.
  In a coordinate system $\{x_i\}$ around $p\in M$, since
  $[\frac{\partial}{\partial x_i},\frac{\partial}{\partial x_j}]=0$,
  $$
  R\Big(\frac{\partial}{\partial x_i},\frac{\partial}{\partial x_j}\Big)\frac{\partial}{\partial x_k}=\Big(\nabla_{\partial/\partial x_j}\nabla_{\partial/\partial x_i}-\nabla_{\partial/\partial x_i}\nabla_{\partial/\partial x_j}\Big)\frac{\partial}{\partial x_k},
  $$
  Hence curvature measures the non-commutativity of covariant derivatives. We use
  this sign convention; the opposite sign also occurs in the literature.
\end{definition}

\begin{proposition}[linearity properties of $R$]
  \label{prop:dc-ch4-2-2}
  \lean{Riemannian.AffineConnection.curvature_add_left, Riemannian.AffineConnection.curvature_add_middle, Riemannian.AffineConnection.curvature_add_right, Riemannian.AffineConnection.curvature_smul_left, Riemannian.AffineConnection.curvature_smul_middle, Riemannian.AffineConnection.curvature_smul_right}
  \uses{def:dc-ch4-2-1}
  \leanok
  \dcref{ch4:2.2}
  The curvature $R$ has the following properties:
  \begin{enumerate}
  \item[(i)] $R$ is bilinear in $\mathcal{X}(M)\times\mathcal{X}(M)$:
  $$
  R(fX_1+gX_2,Y_1)=fR(X_1,Y_1)+gR(X_2,Y_1),
  $$
  $$
  R(X_1,fY_1+gY_2)=fR(X_1,Y_1)+gR(X_1,Y_2),\quad f,g\in\mathcal{D}(M),\ X_1,X_2,Y_1,Y_2\in\mathcal{X}(M).
  $$
  \item[(ii)] For $X,Y\in\mathcal{X}(M)$, the operator
  $R(X,Y):\mathcal{X}(M)\to\mathcal{X}(M)$ is linear, that is,
  $$
  R(X,Y)(Z+W)=R(X,Y)Z+R(X,Y)W,\qquad R(X,Y)fZ=fR(X,Y)Z,
  \quad f\in\mathcal{D}(M),\ Z,W\in\mathcal{X}(M).
  $$
  \end{enumerate}
\end{proposition}
\begin{proof}
  \leanok
  \sketch
  Bilinearity in (i) follows by expanding the definition and applying the
  connection axioms in the first two arguments. Additivity in (ii) is immediate.
  For homogeneity,
  
  $$
  \nabla_Y\nabla_X(fZ)=\nabla_Y(f\nabla_X Z+(Xf)Z)=f\nabla_Y\nabla_X Z+(Yf)(\nabla_X Z)+(Xf)(\nabla_Y Z)+(Y(Xf))Z.
  $$
  
  Therefore
  
  $$
  \nabla_Y\nabla_X(fZ)-\nabla_X\nabla_Y(fZ)=f(\nabla_Y\nabla_X-\nabla_X\nabla_Y)Z+((YX-XY)f)Z,
  $$
  
  hence
  
  $$
  R(X,Y)fZ=f\nabla_Y\nabla_X Z-f\nabla_X\nabla_Y Z+([Y,X]f)Z+f\nabla_{[X,Y]}Z+([X,Y]f)Z=fR(X,Y)Z.\ \blacksquare
  $$
\end{proof}

\begin{remark}[Remark 2.3]
  \label{rem:dc-ch4-2-3}
  \uses{rem:dc-ch4-2-6}
  \dcref{ch4:2.3}
  The term $\nabla_{[X,Y]}Z$ is exactly what makes $R(X,Y)$ linear in its
  third argument; the coordinate form is recorded in \cref{rem:dc-ch4-2-6}.
\end{remark}

\begin{proposition}[Bianchi Identity]
  \label{prop:dc-ch4-2-4}
  \lean{Riemannian.AffineConnection.curvature_bianchi}
  \uses{def:dc-ch4-2-1, def:dc-ch2-3-4, lem:dc-ch0-5-3-jacobi}
  \leanok
  \dcref{ch4:2.4}
  $$
  R(X,Y)Z+R(Y,Z)X+R(Z,X)Y=0.
  $$
\end{proposition}
\begin{proof}
  \leanok
  \sketch
  Expanding the definition and using torsion-freeness gives
  
  $$
  \begin{aligned}R(X,Y)Z+R(Y,Z)X+R(Z,X)Y&=\nabla_Y\nabla_X Z-\nabla_X\nabla_Y Z+\nabla_{[X,Y]}Z\\&\quad+\nabla_Z\nabla_Y X-\nabla_Y\nabla_Z X+\nabla_{[Y,Z]}X\\&\quad+\nabla_X\nabla_Z Y-\nabla_Z\nabla_X Y+\nabla_{[Z,X]}Y\\&=\nabla_Y[X,Z]+\nabla_Z[Y,X]+\nabla_X[Z,Y]-\nabla_{[X,Z]}Y-\nabla_{[Y,X]}Z-\nabla_{[Z,Y]}X\\&=[Y,[X,Z]]+[Z,[Y,X]]+[X,[Z,Y]]=0,\end{aligned}
  $$
  
  where the last equality follows from the Jacobi identity for vector fields.
  $\blacksquare$
\end{proof}

\begin{proposition}[symmetries of the curvature]
  \label{prop:dc-ch4-2-5}
  \lean{Riemannian.AffineConnection.curvatureForm_antisymm_left, Riemannian.AffineConnection.curvatureForm_antisymm_right, Riemannian.AffineConnection.curvatureForm_bianchi, Riemannian.AffineConnection.curvatureForm_pairSwap}
  \uses{def:dc-ch4-2-1, prop:dc-ch4-2-4, def:dc-ch2-3-4, cor:dc-ch2-3-3}
  \leanok
  \dcref{ch4:2.5}
  % (a) = curvatureForm_bianchi (= prop:dc-ch4-2-4); (b) = curvatureForm_antisymm_left;
  % (c) = curvatureForm_antisymm_right (metric compatibility); (d) = curvatureForm_pairSwap.
  Write $\langle R(X,Y)Z,T\rangle=(X,Y,Z,T)$. Then
  \begin{enumerate}
  \item[(a)] $(X,Y,Z,T)+(Y,Z,X,T)+(Z,X,Y,T)=0$;
  \item[(b)] $(X,Y,Z,T)=-(Y,X,Z,T)$;
  \item[(c)] $(X,Y,Z,T)=-(X,Y,T,Z)$;
  \item[(d)] $(X,Y,Z,T)=(Z,T,X,Y)$.
  \end{enumerate}
\end{proposition}
\begin{proof}
  \leanok
  \sketch
  (a) is just the Bianchi identity again; (b) follows directly from \cref{def:dc-ch4-2-1};
  (c) is equivalent to $(X,Y,Z,Z)=0$, whose proof follows:
  
  $$
  (X,Y,Z,Z)=\langle\nabla_Y\nabla_X Z-\nabla_X\nabla_Y Z+\nabla_{[X,Y]}Z,Z\rangle.
  $$
  
  But
  
  $$
  \langle\nabla_Y\nabla_X Z,Z\rangle=Y\langle\nabla_X Z,Z\rangle-\langle\nabla_X Z,\nabla_Y Z\rangle,\qquad\langle\nabla_{[X,Y]}Z,Z\rangle=\tfrac{1}{2}[X,Y]\langle Z,Z\rangle.
  $$
  
  Hence
  
  $$
  \begin{aligned}(X,Y,Z,Z)&=Y\langle\nabla_X Z,Z\rangle-X\langle\nabla_Y Z,Z\rangle+\tfrac{1}{2}[X,Y]\langle Z,Z\rangle\\&=\tfrac{1}{2}Y(X\langle Z,Z\rangle)-\tfrac{1}{2}X(Y\langle Z,Z\rangle)+\tfrac{1}{2}[X,Y]\langle Z,Z\rangle\\&=-\tfrac{1}{2}[X,Y]\langle Z,Z\rangle+\tfrac{1}{2}[X,Y]\langle Z,Z\rangle=0,\end{aligned}
  $$
  
  which proves (c). To prove (d), we use (a), and write:
  
  $$
  \begin{aligned}(X,Y,Z,T)+(Y,Z,X,T)+(Z,X,Y,T)&=0,\\(Y,Z,T,X)+(Z,T,Y,X)+(T,Y,Z,X)&=0,\\(Z,T,X,Y)+(T,X,Z,Y)+(X,Z,T,Y)&=0,\\(T,X,Y,Z)+(X,Y,T,Z)+(Y,T,X,Z)&=0.\end{aligned}
  $$
  
  Summing the equations above we obtain $2(Z,X,Y,T)+2(T,Y,Z,X)=0$ and therefore
  $(Z,X,Y,T)=(Y,T,Z,X)$. $\blacksquare$
\end{proof}

\begin{remark}[Remark 2.6]
  \label{rem:dc-ch4-2-6}
  \lean{Riemannian.leviCivita_curvature_chartFrame_expansion, Riemannian.AffineConnection.curvature_apply_congr, Riemannian.AffineConnection.curvatureFormAt, Riemannian.AffineConnection.isAlgCurvatureForm_curvatureFormAt, Riemannian.leviCivita_curvatureFormAt_chartFrame}
  \uses{rem:dc-ch2-2-3, def:dc-ch2-3-7-christoffel, prop:dc-ch4-2-2, prop:dc-ch4-2-5}
  \leanok
  \dcref{ch4:2.6}
  It is convenient to express what was seen above in a coordinate system $(U,\mathbf{x})$
  based at $p\in M$. Indicating $\frac{\partial}{\partial x_i}=X_i$, we put
  $$
  R(X_i,X_j)X_k=\sum_\ell R_{ijk}^\ell X_\ell,
  $$
  so the $R_{ijk}^\ell$ are the components of the curvature $R$ in $(U,\mathbf{x})$.
  If $X=\sum_i u^i X_i$, $Y=\sum_j v^j X_j$, $Z=\sum_k w^k X_k$, then by linearity of
  $R$,
  $$
  R(X,Y)Z=\sum_{i,j,k,\ell}R_{ijk}^\ell u^i v^j w^k X_\ell.\qquad(1)
  $$
  To express $R_{ijk}^\ell$ in terms of the coefficients $\Gamma_{ij}^k$ of the
  Riemannian connection, we write
  $$
  R(X_i,X_j)X_k=\nabla_{X_j}\nabla_{X_i}X_k-\nabla_{X_i}\nabla_{X_j}X_k=\nabla_{X_j}\Big(\sum_\ell\Gamma_{ik}^\ell X_\ell\Big)-\nabla_{X_i}\Big(\sum_\ell\Gamma_{jk}^\ell X_\ell\Big),
  $$
  which by a direct calculation yields
  $$
  R_{ijk}^s=\sum_\ell\Gamma_{ik}^\ell\Gamma_{j\ell}^s-\sum_\ell\Gamma_{jk}^\ell\Gamma_{i\ell}^s+\frac{\partial}{\partial x_j}\Gamma_{ik}^s-\frac{\partial}{\partial x_i}\Gamma_{jk}^s.\qquad(2)
  $$
  Putting $\langle R(X_i,X_j)X_k,X_s\rangle=\sum_\ell R_{ijk}^\ell g_{\ell s}=R_{ijks}$,
  the identities of Proposition 2.5 can be written as
  $$
  R_{ijks}+R_{jkis}+R_{kijs}=0,\qquad R_{ijks}=-R_{jiks},\qquad R_{ijks}=-R_{ijsk},\qquad R_{ijks}=R_{ksij}.
  $$
  Equation (1), which depends on the linearity of the operator $R$, shows that the
  value of $R(X,Y)Z$ at the point $p$ depends uniquely on the values of $X,Y,Z$ at
  $p$ and the values of the functions $R_{ijk}^\ell$ at $p$. This contrasts with the
  behavior of the covariant derivative (see \cref{rem:dc-ch2-2-3}), the reason
  being that the covariant derivative is not linear in all of its arguments. In
  general, entities such as the curvature that are linear are called \emph{tensors} on
  $M$.
\end{remark}

\section{Sectional curvature}

Given a vector space $V$, we denote by $|x\wedge y|$ the expression

$$
\sqrt{|x|^2|y|^2-\langle x,y\rangle^2},
$$

which represents the area of a two-dimensional parallelogram determined by the
pair of vectors $x,y\in V$.

\begin{proposition}[sectional curvature independent of basis]
  \label{prop:dc-ch4-3-1}
  \lean{Riemannian.IsAlgCurvatureForm.sectionalCurvature_changeBasis}
  \uses{prop:dc-ch4-2-5, def:dc-ch4-3-2}
  \leanok
  \dcref{ch4:3.1}
  Let $\sigma\subset T_pM$ be a two-dimensional subspace of the tangent space
  $T_pM$ and let $x,y\in\sigma$ be two linearly independent vectors. Then
  $$
  K(x,y)=\frac{(x,y,x,y)}{|x\wedge y|^2}
  $$
  does not depend on the choice of the vectors $x,y\in\sigma$.
\end{proposition}
\begin{proof}
  \leanok
  \sketch
  Any two bases of $\sigma$ are related by swaps, nonzero rescalings, and shears
  $\{x,y\}\mapsto\{x+\lambda y,y\}$. The curvature symmetries and the identity
  $|x\wedge y|^2=|x|^2|y|^2-\langle x,y\rangle^2$ show that the numerator and
  denominator transform by the same square factor, so $K(x,y)$ is invariant.
\end{proof}

\begin{definition}[sectional curvature]
  \label{def:dc-ch4-3-2}
  \lean{Riemannian.sectionalCurvature}
  \uses{def:dc-ch4-2-1}
  \leanok
  \dcref{ch4:3.2}
  Given a point $p\in M$ and a two-dimensional subspace $\sigma\subset T_pM$, the
  real number $K(x,y)=K(\sigma)$, where $\{x,y\}$ is any basis of $\sigma$, is
  called the \emph{sectional curvature} of $\sigma$ at $p$.
  The importance of the sectional curvature comes from the fact that knowledge of
  $K(\sigma)$ for all $\sigma$ determines the curvature $R$ completely. This is a
  purely algebraic fact:
\end{definition}

\begin{lemma}[sectional curvature determines $R$]
  \label{lem:dc-ch4-3-3}
  \lean{Riemannian.IsAlgCurvatureForm.ext}
  \uses{prop:dc-ch4-2-5}
  \leanok
  \dcref{ch4:3.3}
  Let $V$ be a vector space of dimension $\geq 2$, provided with an inner product
  $\langle\,,\,\rangle$. Let $R:V\times V\times V\to V$ and $R':V\times V\times V\to V$
  be tri-linear mappings such that conditions (a), (b), (c) and (d) of \cref{prop:dc-ch4-2-5} are satisfied by
  $$
  (x,y,z,t)=\langle R(x,y)z,t\rangle,\qquad (x,y,z,t)'=\langle R'(x,y)z,t\rangle.
  $$
  If $x,y$ are two linearly independent vectors, we may write
  $$
  K(\sigma)=\frac{(x,y,x,y)}{|x\wedge y|^2},\qquad K'(\sigma)=\frac{(x,y,x,y)'}{|x\wedge y|^2},
  $$
  where $\sigma$ is the bi-dimensional subspace generated by $x$ and $y$. If for all
  $\sigma\subset V$, $K(\sigma)=K'(\sigma)$, then $R=R'$.
\end{lemma}
\begin{proof}
  \leanok
  \sketch
  It suffices to prove that $(x,y,z,t)=(x,y,z,t)'$ for any $x,y,z,t\in V$.
  Observe first that, by hypothesis, $(x,y,x,y)=(x,y,x,y)'$ for all $x,y\in V$. Then
  
  $$
  (x+z,y,x+z,y)=(x+z,y,x+z,y)',
  $$
  
  hence
  
  $$
  (x,y,x,y)+2(x,y,z,y)+(z,y,z,y)=(x,y,x,y)'+2(x,y,z,y)'+(z,y,z,y)',
  $$
  
  and therefore $(x,y,z,y)=(x,y,z,y)'$ for all $x,y,z\in V$. Using what we have just
  proved,
  
  $$
  (x,y+t,z,y+t)=(x,y+t,z,y+t)',
  $$
  
  hence
  
  $$
  (x,y,z,t)+(x,t,z,y)=(x,y,z,t)'+(x,t,z,y)',
  $$
  
  which can be written as
  
  $$
  (x,y,z,t)-(x,y,z,t)'=(y,z,x,t)-(y,z,x,t)'.
  $$
  
  It follows that the expression $(x,y,z,t)-(x,y,z,t)'$ is invariant by cyclic
  permutations of the first three elements. Therefore, by (a) of \cref{prop:dc-ch4-2-5},
  
  $$
  3[(x,y,z,t)-(x,y,z,t)']=0,
  $$
  
  hence $(x,y,z,t)=(x,y,z,t)'$ for all $x,y,z,t\in V$. $\blacksquare$
\end{proof}

\begin{lemma}[constant curvature characterization]
  \label{lem:dc-ch4-3-4}
  \lean{Riemannian.IsAlgCurvatureForm.eq_smul_stdCurvForm_of_const, Riemannian.IsAlgCurvatureForm.const_of_eq_smul_stdCurvForm}
  \uses{lem:dc-ch4-3-3, prop:dc-ch4-2-5}
  \leanok
  \dcref{ch4:3.4}
  Let $M$ be a Riemannian manifold and $p$ a point of $M$. Define a tri-linear
  mapping $R':T_pM\times T_pM\times T_pM\to T_pM$ by
  $$
  \langle R'(X,Y,W),Z\rangle=\langle X,W\rangle\langle Y,Z\rangle-\langle Y,W\rangle\langle X,Z\rangle,
  $$
  for all $X,Y,W,Z\in T_pM$. Then $M$ has constant sectional curvature equal to
  $K_o$ if and only if $R=K_o R'$, where $R$ is the curvature of $M$.
\end{lemma}
\begin{proof}
  \leanok
  \sketch
  Assume that $K(p,\sigma)=K_o$ for all $\sigma\subset T_pM$, and set
  $\langle R'(X,Y,W),Z\rangle=(X,Y,W,Z)'$. Observe that $R'$ satisfies the properties
  (a), (b), (c) and (d) of \cref{prop:dc-ch4-2-5}. Since
  
  $$
  (X,Y,X,Y)'=\langle X,X\rangle\langle Y,Y\rangle-\langle X,Y\rangle^2,
  $$
  
  we have that, for all pairs of vectors $X,Y\in T_pM$,
  
  $$
  R(X,Y,X,Y)=K_o(|X|^2|Y|^2-\langle X,Y\rangle^2)=K_o R'(X,Y,X,Y).
  $$
  
  \cref{lem:dc-ch4-3-3} implies that, for all $X,Y,W,Z$,
  
  $$
  R(X,Y,W,Z)=K_o R'(X,Y,W,Z),
  $$
  
  hence $R=K_o R'$. The converse is immediate. $\blacksquare$
\end{proof}

\begin{corollary}[constant curvature in an orthonormal basis]
  \label{cor:dc-ch4-3-5}
  \lean{Riemannian.IsAlgCurvatureForm.eq_kronecker_iff_const}
  \uses{lem:dc-ch4-3-4, prop:dc-ch4-2-5}
  \leanok
  \dcref{ch4:3.5}
  Let $M$ be a Riemannian manifold, $p$ a point of $M$ and $\{e_1,\dots,e_n\}$,
  $n=\dim M$, an orthonormal basis of $T_pM$. Define
  $R_{ijk\ell}=\langle R(e_i,e_j)e_k,e_\ell\rangle$, $i,j,k,\ell=1,\dots,n$. Then
  $K(p,\sigma)=K_o$ for all $\sigma\subset T_pM$ if and only if
  $$
  R_{ijk\ell}=K_o(\delta_{ik}\delta_{j\ell}-\delta_{i\ell}\delta_{jk}),
  $$
  where $\delta_{ij}=1$ if $i=j$ and $\delta_{ij}=0$ if $i\neq j$; equivalently,
  $K(p,\sigma)=K_o$ for all $\sigma\subset T_pM$ if and only if
  $R_{ijij}=-R_{ijji}=K_o$ for all $i\neq j$, and $R_{ijk\ell}=0$ in the other
  cases.
\end{corollary}
\begin{proof}
  \leanok
  \uses{lem:dc-ch4-3-4, lem:dc-ch4-3-3}
  \sketch
  By \cref{lem:dc-ch4-3-4}, $K(p,\sigma)=K_o$ for all $\sigma$ is equivalent to
  $R=K_oR'$, where $R'(X,Y,Z,T)=\langle X,Z\rangle\langle Y,T\rangle-\langle
  Y,Z\rangle\langle X,T\rangle$. Evaluating $R'$ on the orthonormal basis gives
  $R'(e_i,e_j,e_k,e_\ell)=\delta_{ik}\delta_{j\ell}-\delta_{i\ell}\delta_{jk}$,
  so $R=K_oR'$ yields
  $R_{ijk\ell}=K_o(\delta_{ik}\delta_{j\ell}-\delta_{i\ell}\delta_{jk})$.
  Conversely, since an algebraic curvature form is determined by its components
  in a basis (multilinearity, \cref{lem:dc-ch4-3-3}), the Kronecker identity on
  all basis $4$-tuples forces $R=K_oR'$, hence constant sectional curvature
  $K_o$. $\blacksquare$
\end{proof}

\section{Ricci curvature and scalar curvature}

For a unit vector $x=z_n\in T_pM$ and an orthonormal basis
$\{z_1,\dots,z_{n-1}\}$ of $x^\perp$, set

$$
\mathrm{Ric}_p(x)=\frac{1}{n-1}\sum_i\langle R(x,z_i)x,z_i\rangle,\quad i=1,2,\dots,n-1,
$$

$$
K(p)=\frac{1}{n}\sum_j\mathrm{Ric}_p(z_j)=\frac{1}{n(n-1)}\sum_{ij}\langle R(z_i,z_j)z_i,z_j\rangle,\quad j=1,\dots,n.
$$

These are the Ricci curvature in direction $x$ and the normalized scalar
curvature. Define the Ricci form by

$$
Q(x,y)=\text{trace of the mapping }z\mapsto R(x,z)y.
$$

For $x$ unit and $\{z_1,\dots,z_{n-1},z_n=x\}$ orthonormal,

$$
Q(x,y)=\sum_i\langle R(x,z_i)y,z_i\rangle=\sum_i\langle R(y,z_i)x,z_i\rangle=Q(y,x),
$$

Thus $Q$ is symmetric and $Q(x,x)=(n-1)\mathrm{Ric}_p(x)$, so Ricci curvature is
independent of the chosen orthonormal basis. If $K$ is defined by
$\langle K(x),y\rangle=Q(x,y)$, then for any orthonormal basis
$\{z_1,\dots,z_n\}$,

$$
\text{Trace of }K=\sum_j\langle K(z_j),z_j\rangle=\sum_j Q(z_j,z_j)=(n-1)\sum_j\mathrm{Ric}_p(z_j)=n(n-1)K(p),
$$

Consequently the normalized scalar curvature is basis-independent. The bilinear
form $\frac{1}{n-1}Q$ is the Ricci tensor. In coordinates $(x_i)$ with
$X_i=\frac{\partial}{\partial x_i}$, $g_{ij}=\langle X_i,X_j\rangle$ and $g^{ij}$
the inverse matrix of $g_{ij}$, the coefficients of the bilinear form
$\frac{1}{n-1}Q$ in the basis $\{X_i\}$ are

$$
\frac{1}{n-1}R_{ik}=\frac{1}{n-1}\sum_j R_{ijk}^j=\frac{1}{n-1}\sum_{sj}R_{ijks}g^{sj},
$$

and

$$
K=\frac{1}{n(n-1)}\sum_{ik}R_{ik}g^{ik}.
$$

\begin{definition}[Ricci form]
  \label{def:dc-ch4-4-ricci-form}
  \lean{Riemannian.ricciForm, Riemannian.ricciForm_eq_sum}
  \uses{prop:dc-ch4-2-5}
  \leanok
  \dcref{ch4:4}
  Let $p\in M$ and let $B(x,y,z,t)=\langle R(x,y)z,t\rangle$ be the algebraic
  curvature form on $T_pM$. The \emph{Ricci form} is the bilinear form
  $$
  Q(x,y)=\text{trace of the mapping }z\mapsto R(x,z)y,
  $$
  the trace of the endomorphism of $T_pM$ whose associated bilinear form is
  $(z,w)\mapsto B(x,z,y,w)$. In any orthonormal basis $\{e_i\}$ of $T_pM$,
  $$
  Q(x,y)=\sum_i\langle R(x,e_i)y,e_i\rangle=\sum_i B(x,e_i,y,e_i),
  $$
  and this sum is independent of the orthonormal basis, since it is a trace. The
  \emph{Ricci curvature} in the direction of a unit vector $x$ is
  $\mathrm{Ric}_p(x)=\frac{1}{n-1}Q(x,x)$.
\end{definition}

\begin{proposition}[the Ricci form is symmetric]
  \label{prop:dc-ch4-4-ricci-symm}
  \lean{Riemannian.ricciForm_symm}
  \uses{def:dc-ch4-4-ricci-form, prop:dc-ch4-2-5}
  \leanok
  \dcref{ch4:4}
  The Ricci form is symmetric, $Q(x,y)=Q(y,x)$ for all $x,y\in T_pM$. In
  particular $\mathrm{Ric}_p(x)$ is intrinsically defined.
\end{proposition}
\begin{proof}
  \leanok
  \sketch
  In an orthonormal basis $\{e_i\}$ we have $Q(x,y)=\sum_i B(x,e_i,y,e_i)$. By the
  pair-swap symmetry (d) of \cref{prop:dc-ch4-2-5}, $B(x,e_i,y,e_i)=B(y,e_i,x,e_i)$,
  hence $Q(x,y)=\sum_i B(y,e_i,x,e_i)=Q(y,x)$. $\blacksquare$
\end{proof}

\begin{definition}[scalar curvature]
  \label{def:dc-ch4-4-scalar}
  \lean{Riemannian.scalarCurvature, Riemannian.scalarCurvature_eq_sum}
  \uses{def:dc-ch4-4-ricci-form}
  \leanok
  \dcref{ch4:4}
  The (unnormalized) \emph{scalar curvature} at $p$ is the trace of the Ricci
  form $Q$,
  $$
  \mathrm{trace}(Q)=\sum_j Q(e_j,e_j)=\sum_{ij}\langle R(e_i,e_j)e_i,e_j\rangle
    =\sum_{ij}B(e_i,e_j,e_i,e_j),
  $$
  again independent of the orthonormal basis $\{e_i\}$ (a trace of a trace). do
  Carmo's normalized scalar curvature is $K(p)=\frac{1}{n(n-1)}\mathrm{trace}(Q)$.
\end{definition}

\begin{lemma}[curvature on a parametrized surface]
  \label{lem:dc-ch4-4-1}
  \lean{Riemannian.Jacobi.surface_covariant_commutator, Riemannian.Jacobi.surface_covariant_commutator_of_eventually}
  \uses{def:dc-ch4-2-1, prop:dc-ch2-2-2}
  \leanok
  \dcref{ch4:4.1}
  Let $f:A\subset\mathbb{R}^2\to M$ be a parametrized surface with coordinates
  $(s,t)$, and let $V=V(s,t)$ be a vector field along $f$. Then
  $$
  \frac{D}{dt}\frac{D}{ds}V-\frac{D}{ds}\frac{D}{dt}V=R\Big(\frac{\partial f}{\partial s},\frac{\partial f}{\partial t}\Big)V.\qquad(3)
  $$
\end{lemma}
\begin{proof}
  \sketch
  In local coordinates
  $(U,\mathbf{x})$ based at $p\in M$. Let $V=\sum_i v^i X_i$, where $v^i=v^i(s,t)$ and
  $X_i=\frac{\partial}{\partial x_i}$. Then
  
  $$
  \frac{D}{ds}V=\sum_i v^i\frac{D}{ds}X_i+\sum_i\frac{\partial v^i}{\partial s}X_i,
  $$
  
  $$
  \frac{D}{dt}\Big(\frac{D}{ds}V\Big)=\sum_i v^i\frac{D}{dt}\frac{D}{ds}X_i+\sum_i\frac{\partial v^i}{\partial t}\frac{D}{ds}X_i+\sum_i\frac{\partial v^i}{\partial s}\frac{D}{dt}X_i+\sum_i\frac{\partial^2 v^i}{\partial t\partial s}X_i.
  $$
  
  Interchanging the roles of $s$ and $t$ and subtracting,
  
  $$
  \frac{D}{dt}\frac{D}{ds}V-\frac{D}{ds}\frac{D}{dt}V=\sum_i v^i\Big(\frac{D}{dt}\frac{D}{ds}X_i-\frac{D}{ds}\frac{D}{dt}X_i\Big).
  $$
  
  Putting $f(s,t)=(x_1(s,t),\dots,x_n(s,t))$, so that
  $\frac{\partial f}{\partial s}=\sum_j\frac{\partial x_j}{\partial s}X_j$ and
  $\frac{\partial f}{\partial t}=\sum_k\frac{\partial x_k}{\partial t}X_k$, we have
  $\frac{D}{ds}X_i=\sum_j\frac{\partial x_j}{\partial s}\nabla_{X_j}X_i$ and
  
  $$
  \frac{D}{dt}\frac{D}{ds}X_i=\sum_j\frac{\partial^2 x_j}{\partial t\partial s}\nabla_{X_j}X_i+\sum_{jk}\frac{\partial x_j}{\partial s}\frac{\partial x_k}{\partial t}\nabla_{X_k}\nabla_{X_j}X_i,
  $$
  
  or
  
  $$
  \Big(\frac{D}{dt}\frac{D}{ds}-\frac{D}{ds}\frac{D}{dt}\Big)X_i=\sum_{jk}\frac{\partial x_j}{\partial s}\frac{\partial x_k}{\partial t}(\nabla_{X_k}\nabla_{X_j}X_i-\nabla_{X_j}\nabla_{X_k}X_i).
  $$
  
  Joining everything together, we finally get
  
  $$
  \Big(\frac{D}{dt}\frac{D}{ds}-\frac{D}{ds}\frac{D}{dt}\Big)V=\sum_{ijk}v^i\frac{\partial x_j}{\partial s}\frac{\partial x_k}{\partial t}R(X_j,X_k)X_i=R\Big(\frac{\partial f}{\partial s},\frac{\partial f}{\partial t}\Big)V.\ \blacksquare
  $$
\end{proof}

\section{Tensors on Riemannian manifolds}

The space $\mathcal{X}(M)$ is a module over $\mathcal{D}(M)$.

\begin{definition}[tensor of order $r$]
  \label{def:dc-ch4-5-1}
  \lean{Riemannian.CovariantTensor2Predicate, Riemannian.CovariantTensor4Predicate}
  \leanok
  \dcref{ch4:5.1}
  A \emph{tensor} $T$ of order $r$ on a Riemannian manifold is a multilinear mapping
  $$
  T:\underbrace{\mathcal{X}(M)\times\cdots\times\mathcal{X}(M)}_{r\text{ factors}}\to\mathcal{D}(M).
  $$
  This means that given $Y_1,\dots,Y_r\in\mathcal{X}(M)$, $T(Y_1,\dots,Y_r)$ is a
  differentiable function on $M$, and that $T$ is linear in each argument, that is,
  $$
  T(Y_1,\dots,fX+gY,\dots,Y_r)=fT(Y_1,\dots,X,\dots,Y_r)+gT(Y_1,\dots,Y,\dots,Y_r),
  $$
  for all $X,Y\in\mathcal{X}(M)$, $f,g\in\mathcal{D}(M)$.
  A tensor is a pointwise object. Fix $p\in M$ and let $U$ be a neighborhood of $p$
  on which it is possible to define vector fields $E_1,\dots,E_n\in\mathcal{X}(U)$
  such that at each $q\in U$ the vectors $\{E_i(q)\}$, $i=1,\dots,n$, form a basis of
  $T_qM$; we say $\{E_i\}$ is a \emph{moving frame} on $U$. Writing
  $Y_1=\sum_{i_1}y_{i_1}E_{i_1},\dots,Y_r=\sum_{i_r}y_{i_r}E_{i_r}$, by linearity
  $$
  T(Y_1,\dots,Y_r)=\sum_{i_1,\dots,i_r}y_{i_1}\cdots y_{i_r}T(E_{i_1},\dots,E_{i_r}).
  $$
  The functions $T(E_{i_1},\dots,E_{i_r})=T_{i_1\dots i_r}$ on $U$ are called the
  \emph{components} of $T$ in the frame $\{E_i\}$. The value of $T(Y_1,\dots,Y_r)$ at $p$
  depends only on the values at $p$ of the components of $T$ and the values of
  $Y_1,\dots,Y_r$ at $p$.
\end{definition}

\begin{example}[curvature tensor]
  \label{ex:dc-ch4-5-2}
  \lean{Riemannian.AffineConnection.curvatureForm, Riemannian.AffineConnection.curvatureForm_isCovariantTensor4}
  \uses{def:dc-ch4-5-1, prop:dc-ch4-2-2, prop:dc-ch4-2-5}
  \leanok
  \dcref{ch4:5.2}
  The \emph{curvature tensor}

  $$
  R:\mathcal{X}(M)\times\mathcal{X}(M)\times\mathcal{X}(M)\times\mathcal{X}(M)\to\mathcal{D}(M)
  $$

  is defined by

  $$
  R(X,Y,Z,W)=\langle R(X,Y)Z,W\rangle,\quad X,Y,Z,W\in\mathcal{X}(M).
  $$

  The components in the coordinate frame $\{X_i=\frac{\partial}{\partial x_i}\}$
  are $R(X_i,X_j,X_k,X_\ell)=R_{ijk\ell}$.
\end{example}
\begin{proof}
  \leanok
  \sketch
  Multilinearity of $R(X,Y,Z,W)=\langle R(X,Y)Z,W\rangle$ over $\mathcal{D}(M)$
  in the first three arguments is the $\mathcal{D}(M)$-linearity of the curvature
  operator $R(X,Y)Z$ (\cref{prop:dc-ch4-2-2}: additivity and homogeneity
  $R(fX,Y)Z=R(X,fY)Z=R(X,Y)fZ=fR(X,Y)Z$); linearity in the fourth argument $W$
  is the bilinearity of the metric $\langle\cdot,\cdot\rangle$. $\blacksquare$
\end{proof}

\begin{example}[metric tensor]
  \label{ex:dc-ch4-5-3}
  \lean{Riemannian.metricTensor, Riemannian.metricTensor_isCovariantTensor2}
  \uses{def:dc-ch4-5-1}
  \leanok
  \dcref{ch4:5.3}
  The "\emph{metric tensor}" $G:\mathcal{X}(M)\times\mathcal{X}(M)\to\mathcal{D}(M)$ is
  defined by $G(X,Y)=\langle X,Y\rangle$, $X,Y\in\mathcal{X}(M)$. $G$ is a tensor of
  order 2 and its components in the frame $\{X_i\}$ are the coefficients $g_{ij}$ of
  the Riemannian metric in the given system of coordinates.
\end{example}
\begin{proof}
  \leanok
  \sketch
  $\mathcal{D}(M)$-bilinearity of $G(X,Y)=\langle X,Y\rangle$ is inherited slotwise
  from the bilinearity of the inner product $g_p$ on each tangent space:
  $G(fX+gY,Z)=fG(X,Z)+gG(Y,Z)$ and symmetrically in the second argument.
  $\blacksquare$
\end{proof}

\begin{example}[the connection is not a tensor]
  \label{ex:dc-ch4-5-4}
  \lean{Riemannian.AffineConnection.connectionForm, Riemannian.AffineConnection.connectionForm_not_isCovariantTensor}
  \uses{def:dc-ch4-5-1, def:dc-ch2-2-1}
  \leanok
  \dcref{ch4:5.4}
  The Riemannian connection $\nabla$ defined by

  $$
  \nabla:\mathcal{X}(M)\times\mathcal{X}(M)\times\mathcal{X}(M)\to\mathcal{D}(M),\qquad\nabla(X,Y,Z)=\langle\nabla_X Y,Z\rangle,\quad X,Y,Z\in\mathcal{X}(M),
  $$

  is \emph{not} a tensor, because $\nabla$ is not linear with respect to the argument
  $Y$.
\end{example}
\begin{proof}
  \leanok
  \sketch
  The form is $\mathcal{D}(M)$-linear in the direction slot $X$, by the
  $\mathcal{D}(M)$-linearity of $\nabla$ in its first argument, and in the metric
  slot $Z$, by bilinearity of $\langle\cdot,\cdot\rangle$. In the differentiated
  slot $Y$, however, the Leibniz rule of the connection (\cref{def:dc-ch2-2-1},
  property (iii)) gives
  $$
  \nabla(X, fY, Z)=f\,\nabla(X, Y, Z)+(Xf)\,\langle Y, Z\rangle .
  $$
  The correction term $(Xf)\langle Y,Z\rangle$ is not of the tensorial form
  $f\,\nabla(X,Y,Z)$: whenever it is nonzero at some point — which occurs on any
  manifold of positive dimension — the middle-slot homogeneity required of an
  order-$3$ tensor (\cref{def:dc-ch4-5-1}) fails, so $\nabla$ is not a tensor.
  $\blacksquare$
\end{proof}

\begin{remark}[Remark 5.5]
  \label{rem:dc-ch4-5-5}
  \lean{Riemannian.RiemannianMetric.metricToDualEquiv, Riemannian.RiemannianMetric.metricToDual_apply, Riemannian.RiemannianMetric.metricRiesz_inner}
  \uses{def:dc-ch1-2-1, def:dc-ch4-5-1}
  \leanok
  \dcref{ch4:5.5}
  On a differentiable manifold without a metric, covariant tensors and
  contravariant tensors, defined using the dual $\mathcal{X}^*(M)$, are distinct.
  A Riemannian metric associates to each $X\in\mathcal{X}(M)$ a unique
  $\omega\in\mathcal{X}^*(M)$ given by
  $$
  \omega(Y)=\langle X,Y\rangle,\quad\text{for all }Y\in\mathcal{X}(M).
  $$
  This correspondence identifies covariant and contravariant tensors.
\end{remark}

\begin{remark}[Remark 5.6]
  \label{rem:dc-ch4-5-6}
  \lean{Riemannian.AffineConnection.covectorTensor, Riemannian.AffineConnection.covectorTensor_apply}
  \uses{rem:dc-ch4-5-5, ex:dc-ch4-5-3}
  \leanok
  \dcref{ch4:5.6}
  Identify the field $X\in\mathcal{X}(M)$ with the tensor
  $X:\mathcal{X}(M)\to\mathcal{D}(M)$ given by $X(Y)=\langle X,Y\rangle$,
  for all $Y\in\mathcal{X}(M)$.
\end{remark}

\begin{definition}[covariant differential of a tensor]
  \label{def:dc-ch4-5-7}
  \lean{Riemannian.AffineConnection.covariantDifferential1, Riemannian.AffineConnection.covariantDifferential2}
  \uses{def:dc-ch4-5-1, def:dc-ch2-2-1}
  \leanok
  \dcref{ch4:5.7}
  Let $T$ be a tensor of order $r$. The \emph{covariant differential} $\nabla T$ of $T$
  is a tensor of order $(r+1)$ given by
  $$
  \nabla T(Y_1,\dots,Y_r,Z)=Z(T(Y_1,\dots,Y_r))-T(\nabla_Z Y_1,\dots,Y_r)-\cdots-T(Y_1,\dots,Y_{r-1},\nabla_Z Y_r).
  $$
  For each $Z\in\mathcal{X}(M)$, the \emph{covariant derivative} $\nabla_Z T$ of $T$
  relative to $Z$ is a tensor of order $r$ given by
  $$
  \nabla_Z T(Y_1,\dots,Y_r)=\nabla T(Y_1,\dots,Y_r,Z).
  $$
  In a convenient frame this becomes natural. Let $p\in M$ and let
  $\alpha:(-\varepsilon,\varepsilon)\to M$ be a differentiable curve with $\alpha(0)=p$,
  $\alpha'(t)=Z(\alpha(t))$. Let $\{e_1,\dots,e_n\}$ be a basis of $T_pM$ and let
  $e_i(t)$ be the parallel transport of $e_i$ along $\alpha$. Let $T_{i_1\dots i_r}(t)$
  be the components, in the basis $\{e_i(t)\}$, of the restriction $T(\alpha(t))$ of
  $T$ to the curve $\alpha$. Then, by the definition of $\nabla_Z T$,
  $$
  (\nabla_Z T)(e_{i_1}(t),\dots,e_{i_r}(t))=\frac{d}{dt}T_{i_1\dots i_r}(t)-T(\nabla_Z e_{i_1},\dots,e_{i_r}(t))-\cdots-T(e_{i_1}(t),\dots,\nabla_Z e_{i_r}(t)).
  $$
  Since $\nabla_Z e_i(t)=0$, we have by linearity
  $$
  (\nabla_Z T)_{i_1\dots i_r}=(\nabla_Z T)(e_{i_1}(t),\dots,e_{i_r}(t))=\frac{d}{dt}T_{i_1\dots i_r}.
  $$
  Thus, in this frame, the components of $\nabla_ZT$ are the ordinary derivatives
  of the components of $T$.
\end{definition}

\begin{example}[covariant differential of the metric tensor]
  \label{ex:dc-ch4-5-8}
  \lean{Riemannian.AffineConnection.covariantDifferential2_metricTensor_eq_zero}
  \uses{def:dc-ch4-5-7, ex:dc-ch4-5-3, def:dc-ch2-3-1}
  \leanok
  \dcref{ch4:5.8}
  The covariant differential of the metric tensor is the zero tensor. Indeed, for
  all $X,Y,Z\in\mathcal{X}(M)$,

  $$
  \nabla G(X,Y,Z)=Z\langle X,Y\rangle-\langle\nabla_Z X,Y\rangle-\langle X,\nabla_Z Y\rangle=0,
  $$

  because $\nabla$ is the Riemannian connection.
\end{example}
\begin{proof}
  \leanok
  \sketch
  By \cref{def:dc-ch4-5-7}, $\nabla G(X,Y,Z)=Z\langle X,Y\rangle-\langle\nabla_Z X,Y\rangle
  -\langle X,\nabla_Z Y\rangle$. Compatibility of the Levi-Civita connection with the
  metric (\cref{def:dc-ch2-3-1}) gives $Z\langle X,Y\rangle=\langle\nabla_Z X,Y\rangle
  +\langle X,\nabla_Z Y\rangle$, so the three terms cancel. $\blacksquare$
\end{proof}

\begin{example}[covariant derivative of a vector field as a tensor]
  \label{ex:dc-ch4-5-9}
  \lean{Riemannian.AffineConnection.covectorTensor, Riemannian.AffineConnection.covariantDifferential1_covectorTensor_eq}
  \uses{rem:dc-ch4-5-6, def:dc-ch4-5-7, def:dc-ch2-3-1}
  \leanok
  \dcref{ch4:5.9}
  Let $X\in\mathcal{X}(M)$. Identify $X$ with the tensor that associates to the
  vector field $Y\in\mathcal{X}(M)$ the function $\langle X,Y\rangle$ (see \cref{rem:dc-ch4-5-6}). The covariant derivative of the tensor $X$ relative to the vector field
  $Z\in\mathcal{X}(M)$ is such that, for all $Y\in\mathcal{X}(M)$,

  $$
  \nabla_Z X(Y)=\nabla X(Y,Z)=Z(X(Y))-X(\nabla_Z Y)=Z\langle X,Y\rangle-\langle X,\nabla_Z Y\rangle=\langle\nabla_Z X,Y\rangle.
  $$

  We conclude that the tensor $\nabla_Z X$ can be identified with the vector field
  $\nabla_Z X$. This justifies the notation adopted, and shows that the covariant
  derivative of tensors is a generalization of the covariant derivative of vector
  fields.
\end{example}
\begin{proof}
  \leanok
  \sketch
  By \cref{def:dc-ch4-5-7}, $\nabla_Z X(Y)=Z\langle X,Y\rangle-\langle X,\nabla_Z Y\rangle$.
  Metric compatibility (\cref{def:dc-ch2-3-1}) gives $Z\langle X,Y\rangle
  =\langle\nabla_Z X,Y\rangle+\langle X,\nabla_Z Y\rangle$, whence
  $\nabla_Z X(Y)=\langle\nabla_Z X,Y\rangle$; i.e. the tensor $\nabla_Z X$ is the
  covector tensor of the vector field $\nabla_Z X$. $\blacksquare$
\end{proof}

\chapter{Jacobi Fields}
\label{chap:dc-jacobi-fields}

\section{The Jacobi equation}

\begin{definition}[Jacobi field]
  \label{def:dc-ch5-2-1}
  \lean{Riemannian.Jacobi.exists_jacobiField, Riemannian.Jacobi.jacobiField_eqOn}
  \uses{lem:dc-ch5-2-1-parallel-frame, lem:dc-ch5-2-1-ode, lem:dc-ch5-2-1-frame-reduction, lem:dc-ch5-2-1-frame-existence, lem:dc-ch5-2-1-real-curvature}
  \dcref{ch5:2.1}
  Let $\gamma:[0,a]\to M$ be a geodesic in $M$. A vector field $J$ along $\gamma$ is
  a \emph{Jacobi field} if it satisfies
  $$
  \frac{D^2 J}{dt^2}+R(\gamma'(t),J(t))\gamma'(t)=0.\tag{1}
  $$
  on $[0,a]$.

  If $p\in M$, $v\in T_pM$ lies in the domain of $\exp_p$, and we identify
  $T_v(T_pM)$ with $T_pM$, let $v(s)$ satisfy $v(0)=v$, $v'(0)=w$, and set
  $f(t,s)=\exp_p(tv(s))$. Then the variational field
  $J(t)=\frac{\partial f}{\partial s}(t,0)$ is Jacobi and
  $$
  (d\exp_p)_v w=\frac{\partial f}{\partial s}(1,0),\qquad
  (d\exp_p)_{tv}(tw)=\frac{\partial f}{\partial s}(t,0).
  $$

  A Jacobi field is uniquely determined by $J(0)$ and $\frac{DJ}{dt}(0)$, and
  every pair of initial values occurs. Consequently, if $n=\dim M$, the space
  of Jacobi fields along $\gamma$ has dimension $2n$.
\end{definition}
\begin{proof}
  \sketch
  For each fixed $s$, the curve $t\mapsto f(t,s)$ is geodesic. Commuting the
  covariant derivatives in the variation gives
  $$
  0=\frac{D}{\partial s}\Big(\frac{D}{\partial t}\frac{\partial f}{\partial t}\Big)
  =\frac{D}{\partial t}\frac{D}{\partial t}\frac{\partial f}{\partial s}
  +R\Big(\frac{\partial f}{\partial t},\frac{\partial f}{\partial s}\Big)
    \frac{\partial f}{\partial t},
  $$
  which is (1) for $J=\partial f/\partial s$.

  In a parallel orthonormal frame $e_1(t),\dots,e_n(t)$, write
  
  $$
  J(t)=\sum_i f_i(t)e_i(t),\quad a_{ij}=\langle R(\gamma'(t),e_i(t))\gamma'(t),e_j(t)\rangle,\quad i,j=1,\dots,n=\dim M,
  $$
  
  one has $\frac{D^2 J}{dt^2}=\sum_i f_i''(t)e_i(t)$ and
  $R(\gamma',J)\gamma'=\sum_{ij}f_i a_{ij}e_j$. Therefore equation (1) is equivalent
  to the system
  
  $$
  f_j''(t)+\sum_i a_{ij}(t)f_i(t)=0,\quad j=1,\dots,n,
  $$
  
  Standard existence and uniqueness for this second-order linear system gives
  the claim.
\end{proof}

\begin{lemma}[parallel orthonormal frame along a geodesic, coordinate form]
  \label{lem:dc-ch5-2-1-parallel-frame}
  \lean{Riemannian.Jacobi.exists_parallelOrthoFrame_self, Riemannian.Jacobi.exists_parallelOrthoFrame, Riemannian.Jacobi.exists_chartOrthonormalBasis_self, Riemannian.Jacobi.chartMetricInner_const_of_parallelSol}
  \uses{def:dc-ch2-2-5, lem:dc-ch2-3-1-coord}
  \dcref{ch5:2.1}
  Along a geodesic $\gamma:[0,a]\to M$, every orthonormal basis of
  $T_{\gamma(0)}M$ extends uniquely to a parallel orthonormal frame
  $e_1(t),\dots,e_n(t)$.
\end{lemma}
\begin{proof}

  \sketch
  Extend the initial basis by parallel transport. The inner products
  $\langle e_i,e_j\rangle$ are constant by
  \cref{lem:dc-ch2-3-1-coord}, hence remain equal to $\delta_{ij}$.
\end{proof}

\begin{lemma}[Jacobi system in a parallel frame: existence and uniqueness]
  \label{lem:dc-ch5-2-1-ode}
  \lean{Riemannian.Jacobi.exists_isJacobiPairOn, Riemannian.Jacobi.IsJacobiPairOn.eqOn, Riemannian.Jacobi.IsJacobiPairOn.add, Riemannian.Jacobi.IsJacobiPairOn.smul}
  \leanok
  \dcref{ch5:2.1}
  Let $A:[0,a]\to\operatorname{End}(\mathbb R^n)$ be continuous. For every
  $(f_0,f_1)\in\mathbb R^n\times\mathbb R^n$, the system
  $f''+A(t)f=0$ has a unique solution with $f(0)=f_0$ and $f'(0)=f_1$.
  Its solutions form a $2n$-dimensional vector space.
\end{lemma}

\begin{lemma}[Jacobi reduction in a parallel orthonormal frame]
  \label{lem:dc-ch5-2-1-frame-reduction}
  \lean{Riemannian.Jacobi.covariantDerivCoord_frameCombination, Riemannian.Jacobi.covariantDerivCoord2_frameCombination, Riemannian.Jacobi.chartMetricInner_covariantDerivCoord2_frame, Riemannian.Jacobi.chartMetricInner_map_frameCombination_left, Riemannian.Jacobi.frameJacobiComponent, Riemannian.Jacobi.frameExpansion, Riemannian.Jacobi.covariantDerivCoord2_add_map_frameCombination_expand}
  \uses{lem:dc-ch5-2-1-parallel-frame, prop:dc-ch2-2-2}
  \leanok
  \dcref{ch5:2.1}
  In a parallel orthonormal frame $e_1(t),\dots,e_n(t)$, a field
  $J=\sum_i f_i e_i$ is Jacobi if and only if
  $f_j''+\sum_i a_{ij}f_i=0$ for every $j$, where
  $a_{ij}=\langle R(\gamma',e_i)\gamma',e_j\rangle$.
\end{lemma}
\begin{proof}

  \sketch
  Since $\frac{De_i}{dt}=0$, the Leibniz rule of
  \cref{prop:dc-ch2-2-2} gives
  $$
  \frac{D}{dt}\sum_i f_i e_i=\sum_i f_i'\,e_i,\qquad
  \frac{D^2}{dt^2}\sum_i f_i e_i=\sum_i f_i''\,e_i.
  $$
  Pairing the Jacobi equation with $e_j$ extracts
  $f_j''+\sum_i a_{ij}f_i$. Conversely, expansion in the frame gives
  $\frac{D^2J}{dt^2}+R(\gamma',J)\gamma'=\sum_j\big(f_j''+\sum_i a_{ij}f_i\big)e_j$.
\end{proof}

\begin{lemma}[existence and uniqueness of a Jacobi field by its initial data, frame form]
  \label{lem:dc-ch5-2-1-frame-existence}
  \lean{Riemannian.Jacobi.jacobiCoefOp, Riemannian.Jacobi.exists_jacobiField_frame, Riemannian.Jacobi.jacobiField_frame_eqOn}
  \uses{lem:dc-ch5-2-1-parallel-frame, lem:dc-ch5-2-1-ode, lem:dc-ch5-2-1-frame-reduction}
  \leanok
  \dcref{ch5:2.1}
  For every pair $(J_0,w_0)\in T_{\gamma(0)}M\times T_{\gamma(0)}M$, there is a
  unique Jacobi field $J$ along $\gamma$ such that
  $J(0)=J_0$ and $\frac{DJ}{dt}(0)=w_0$.
\end{lemma}
\begin{proof}
  \leanok
  \uses{lem:dc-ch5-2-1-parallel-frame, lem:dc-ch5-2-1-ode, lem:dc-ch5-2-1-frame-reduction}
  \sketch
  Package the coefficients $a_{ij}(t)$ as a continuous operator field $A(t)$ acting on the
  coefficient vector $(f_i)$ by $(A(t)f)_j=\sum_i a_{ij}(t)f_i$. By \cref{lem:dc-ch5-2-1-ode}
  the second-order linear system $f''+A(t)f=0$ has a solution with initial data
  $f_i(0)=\langle J_0,e_i(0)\rangle$, $f_i'(0)=\langle w_0,e_i(0)\rangle$, unique given that
  data. Put $J=\sum_i f_i e_i$. Pairing the Jacobi operator with the frame
  (\cref{lem:dc-ch5-2-1-frame-reduction}) expands
  $\frac{D^2J}{dt^2}+R(\gamma',J)\gamma'=\sum_j\big(f_j''+\sum_i a_{ij}f_i\big)e_j=0$, so $J$
  is a Jacobi field; the orthonormal expansion gives
  $J(0)=\sum_i\langle J_0,e_i(0)\rangle e_i(0)=J_0$, and the frame being parallel gives the
  initial velocity $w_0$. Uniqueness follows from uniqueness of the scalar system.
\end{proof}

\begin{lemma}[the Jacobi coefficient is the curvature contraction read in the chart]
  \label{lem:dc-ch5-2-1-real-curvature}
  \lean{Riemannian.Jacobi.chartCurvatureCoef, Riemannian.Jacobi.chartCurvatureCoef_contDiffOn, Riemannian.Jacobi.chartCurvatureOp, Riemannian.Jacobi.continuousOn_chartCurvatureOp, Riemannian.Jacobi.exists_jacobiField, Riemannian.Jacobi.jacobiField_eqOn}
  \uses{lem:dc-ch5-2-1-frame-existence, rem:dc-ch4-2-6}
  \dcref{ch5:2.1}
  In coordinates $u$ along $\gamma$, the curvature contraction
  $w\mapsto R(\gamma',w)\gamma'$ has components
  $$
  \big(R(t)\,w\big)^l=\sum_{i,j,k}R^l{}_{ijk}\big(u(t)\big)\,\dot u^i(t)\,w^j\,\dot u^k(t),
  \qquad
  R^l{}_{ijk}=\partial_j\Gamma^l_{ik}-\partial_i\Gamma^l_{jk}
    +\sum_s\big(\Gamma^s_{ik}\Gamma^l_{js}-\Gamma^s_{jk}\Gamma^l_{is}\big),
  $$
  with the coefficients of \cref{rem:dc-ch4-2-6}.
\end{lemma}
\begin{proof}

  \sketch
  These coefficients are smooth, so the operator field is continuous along $\gamma$. Hence
  \cref{lem:dc-ch5-2-1-frame-existence} applies to the intrinsic Jacobi equation.
\end{proof}

\begin{remark}[standard tangential Jacobi fields]
  \label{rem:dc-ch5-2-2}
  \lean{Riemannian.Jacobi.isJacobiFieldAlongOn_velocity, Riemannian.Jacobi.isJacobiFieldAlongOn_smul_velocity}
  \uses{def:dc-ch5-2-1}
  \leanok
  \dcref{ch5:2.2}
  Both $\gamma'$ and $t\gamma'$ are Jacobi fields along $\gamma$. The former is
  parallel and nowhere zero for a nonconstant geodesic, while the latter
  vanishes exactly at $t=0$. Jacobi fields orthogonal to $\gamma'$ will be called
  \emph{normal Jacobi fields}.
\end{remark}
\begin{proof}

  \sketch
  The assertions follow from $\nabla_{\gamma'}\gamma'=0$ and
  $R(\gamma',\gamma')\gamma'=0$.
\end{proof}

\begin{example}[Jacobi fields on manifolds of constant curvature]
  \label{ex:dc-ch5-2-3}
  \lean{Riemannian.Jacobi.isJacobiFieldOn_of_constantCurvature, Riemannian.Jacobi.isJacobiFieldOn_constCurvatureSol_pos, Riemannian.Jacobi.isJacobiFieldOn_constCurvatureSol_zero, Riemannian.Jacobi.isJacobiFieldOn_constCurvatureSol_neg}
  \uses{lem:dc-ch4-3-4, lem:dc-ch5-2-3-const-form, lem:dc-ch5-2-3-jacobi-op}
  \leanok
  \dcref{ch5:2.3}
  Let $M$ have constant sectional curvature $K$, let
  $\gamma:[0,\ell]\to M$ be unit speed, and let $J$ be a normal Jacobi field.
  Then
  $$
  R(\gamma',J)\gamma'=KJ.
  $$
  and the Jacobi equation becomes

  $$
  \frac{D^2 J}{dt^2}+KJ=0.\qquad(2)
  $$

  If $w$ is a parallel unit field orthogonal to $\gamma'$, the solution of (2)
  with $J(0)=0$ and $J'(0)=w(0)$ is
  $$
  J(t)=\begin{cases}\dfrac{\sin(t\sqrt{K})}{\sqrt{K}}\,w(t),&\text{if }K>0,\\[1ex] t\,w(t),&\text{if }K=0,\\[1ex]\dfrac{\sinh(t\sqrt{-K})}{\sqrt{-K}}\,w(t),&\text{if }K<0,\end{cases}
  $$
\end{example}
\begin{proof}
  

  \sketch
  By \cref{lem:dc-ch4-3-4}, for every field $T$ along $\gamma$,
  $$
  \langle R(\gamma',J)\gamma',T\rangle
  =K\big(\langle\gamma',\gamma'\rangle\langle J,T\rangle
  -\langle\gamma',T\rangle\langle J,\gamma'\rangle\big)
  =K\langle J,T\rangle.
  $$
  Thus $R(\gamma',J)\gamma'=KJ$. Since $w$ is parallel, substituting
  $J=h w$ reduces (2) to $h''+Kh=0$, whose normalized solutions are the three
  displayed functions.
\end{proof}

\begin{lemma}[constant curvature: the pointwise $(0,4)$ form, coordinate form]
  \label{lem:dc-ch5-2-3-const-form}
  \lean{Riemannian.Jacobi.curvatureFormAt_isConstantCurvature}
  \uses{lem:dc-ch4-3-4}
  \leanok
  \dcref{ch5:2.3}
  At each $q\in M$, the constant-curvature identity
  $\langle R(X,Y)Z,W\rangle=K_0(\langle X,Z\rangle\langle Y,W\rangle-\langle Y,Z\rangle\langle X,W\rangle)$
  from \cref{lem:dc-ch4-3-4} gives the pointwise $(0,4)$-form
  $\langle R(x,y)z,t\rangle_g=K_0(\langle x,z\rangle\langle y,t\rangle-\langle y,z\rangle\langle x,t\rangle)$
  for all $x,y,z,t\in T_qM$.
\end{lemma}

\begin{lemma}[constant curvature: $R(\gamma',J)\gamma'=K J$, coordinate form]
  \label{lem:dc-ch5-2-3-jacobi-op}
  \lean{Riemannian.Jacobi.chartCurvatureOp_isConstantCurvature}
  \uses{lem:dc-ch5-2-3-const-form, cor:dc-ch5-2-9}
  \leanok
  \dcref{ch5:2.3}
  On a manifold of constant sectional curvature $K_0$, if $J$ is normal to a
  unit-speed geodesic $\gamma$, then
  $R(\gamma',J)\gamma'=K_0J$. Consequently the Jacobi equation is
  $D^2J/dt^2+K_0J=0$.
\end{lemma}
\begin{proof}

  \sketch
  Pair $R(\gamma',J)\gamma'$ with an arbitrary vector and apply
  \cref{lem:dc-ch5-2-3-const-form}; the unit-speed and orthogonality hypotheses
  leave $K_0\langle J,\cdot\rangle$.
\end{proof}

\begin{proposition}[Jacobi fields from exponential-map variations]
  \label{prop:dc-ch5-2-4}
  \lean{Riemannian.Jacobi.expDifferential_eq_jacobiField, Riemannian.Jacobi.IsJacobiFieldAlongOn.eqOn_of_initial}
  \uses{def:dc-ch5-2-1}
  \leanok
  \dcref{ch5:2.4}
  Let $\gamma:[0,a]\to M$ be a geodesic and let $J$ be a Jacobi field along
  $\gamma$ with $J(0)=0$. Put $\frac{DJ}{dt}(0)=w$ and $\gamma'(0)=v$. Using the
  natural identification $T_{av}(T_{\gamma(0)}M)\simeq T_{\gamma(0)}M$, construct a
  curve $v(s)$ in $T_{\gamma(0)}M$ with $v(0)=av$, $v'(0)=aw$. Put
  $f(t,s)=\exp_p(\frac{t}{a}v(s))$,
  $p=\gamma(0)$, and define a Jacobi field $\bar J$ by
  $\bar J(t)=\frac{\partial f}{\partial s}(t,0)$. Then $\bar J=J$ on $[0,a]$.
\end{proposition}
\begin{proof}
  \leanok
  \uses{def:dc-ch5-2-1}
  \sketch
  For $s=0$,
  
  $$
  \frac{D}{dt}\frac{\partial f}{\partial s}=\frac{D}{\partial t}\big((d\exp_p)_{tv}(tw)\big)
  =\frac{D}{\partial t}\big(t(d\exp_p)_{tv}(w)\big)
  =(d\exp_p)_{tv}(w)+t\frac{D}{\partial t}\big((d\exp_p)_{tv}(w)\big).
  $$
  
  Therefore, for $t=0$,
  
  $$
  \frac{D\bar J}{dt}(0)=\frac{D}{\partial t}\frac{\partial f}{\partial s}(0,0)=(d\exp_p)_o(w)=w.
  $$
  
  Since $J(0)=\bar J(0)=0$ and $\frac{DJ}{dt}(0)=\frac{D\bar J}{dt}(0)=w$, we
  conclude from uniqueness that $J=\bar J$.
\end{proof}

\begin{corollary}[exponential differential as a Jacobi field]
  \label{cor:dc-ch5-2-5}
  \lean{Riemannian.Jacobi.expDifferential_eq_jacobiField}
  \uses{prop:dc-ch5-2-4, def:dc-ch5-2-1}
  \leanok
  \dcref{ch5:2.5}
  Let $\gamma:[0,a]\to M$ be a geodesic and put $p=\gamma(0)$. Then a Jacobi field
  $J$ along $\gamma$ with $J(0)=0$ is given by

  $$
  J(t)=(d\exp_p)_{t\gamma'(0)}(tJ'(0)),\quad t\in[0,a].
  $$

\end{corollary}

\begin{remark}[arbitrary initial values and derivative notation]
  \label{rem:dc-ch5-2-6}
  \uses{prop:dc-ch5-2-4}
  \dcref{ch5:2.6}
  The construction of \cref{prop:dc-ch5-2-4} extends to Jacobi fields with
  arbitrary initial value. We henceforth write $\frac{DJ}{dt}=J'$ and
  $\frac{D^2 J}{dt^2}=J''$, etc.
\end{remark}

\begin{lemma}[Taylor expansion of $|J|^2$ in a parallel frame (analytic core)]
  \label{lem:dc-ch5-2-7-taylor-ode}
  \lean{Riemannian.Jacobi.norm_sq_jacobi_isLittleO, Riemannian.Jacobi.norm_sq_jacobi_isLittleO_local}
  \uses{def:dc-ch5-2-1}
  \leanok
  \dcref{ch5:2.7}
  Let $J=\sum_i f_i e_i$ be a Jacobi field in a parallel orthonormal frame,
  and put
  $A(t)=\big(\langle R(\gamma',e_i)\gamma',e_j\rangle\big)_{ij}$.
  If $f(0)=0$ and $w=f'(0)$, then
  $$
  |J(t)|^2=\langle w,w\rangle t^2
  -\tfrac13\langle w,A(0)w\rangle t^4+o(t^4),\qquad t\to0.
  $$
\end{lemma}
\begin{proof}

  \sketch
  The frame equation is $f''=-A(t)f$. For
  $g(t)=\langle f(t),f(t)\rangle$, differentiation at $0$ gives
  $$
  g(0)=g'(0)=g'''(0)=0,\quad
  g''(0)=2\langle w,w\rangle,\quad
  g''''(0)=-8\langle w,A(0)w\rangle.
  $$
  Taylor's theorem with Peano remainder yields the expansion. Moreover,
  $\langle w,A(0)w\rangle
  =\langle R(\gamma'(0),w)\gamma'(0),w\rangle$.
\end{proof}

\begin{lemma}[$C^\infty$ smoothness of the parallel frame and the Jacobi coefficient]
  \label{lem:dc-ch5-2-7-frame-smooth}
  \lean{Riemannian.Jacobi.exists_contDiffOn_parallelOrthoFrame, Riemannian.Jacobi.contDiffOn_infty_jacobiCoefOp}
  \uses{lem:dc-ch5-2-1-parallel-frame, lem:dc-ch5-2-1-real-curvature}
  \leanok
  \dcref{ch5:2.7}
  Along a smooth geodesic in a coordinate chart, the parallel orthonormal frame
  of \cref{lem:dc-ch5-2-1-parallel-frame} may be chosen smooth. Consequently,
  $$
  A(t)=\big(\langle R(\gamma',e_i)\gamma',e_j\rangle\big)_{ij}
  $$
  is smooth in $t$.
\end{lemma}
\begin{proof}

  \sketch
  In coordinates the parallel-transport equation is the linear ODE
  $\dot e_i=-\Gamma(\dot u,e_i)(u)$. Its coefficients are smooth, so its
  solutions are smooth. The assertion for $A$ follows from
  \cref{lem:dc-ch5-2-1-real-curvature}.
\end{proof}

\begin{proposition}[fourth-order norm expansion for exponential variations]
  \label{prop:dc-ch5-2-7}
  \lean{Riemannian.Jacobi.norm_sq_jacobi_frame_isLittleO}
  \uses{lem:dc-ch5-2-7-taylor-ode, lem:dc-ch5-2-7-frame-smooth, def:dc-ch5-2-1}
  \leanok
  \dcref{ch5:2.7}
  Let $p\in M$ and $\gamma:[0,a]\to M$ be a geodesic with $\gamma(0)=p$,
  $\gamma'(0)=v$. Let $w\in T_v(T_p M)$ with $|w|=1$ and let $J$ be a Jacobi field
  along $\gamma$ given by
  $$
  J(t)=(d\exp_p)_{tv}(tw),\quad 0\le t\le a.
  $$
  Then the Taylor expansion of $|J(t)|^2$ about $t=0$ is given by
  $$
  |J(t)|^2=t^2-\tfrac{1}{3}\langle R(v,w)v,w\rangle t^4+R(t),\qquad(3)
  $$
  where $\lim_{t\to 0}\frac{R(t)}{t^4}=0$.
\end{proposition}
\begin{proof}
  \leanok
  \uses{lem:dc-ch5-2-7-taylor-ode, lem:dc-ch5-2-7-frame-smooth}
  \sketch
  Since $J(0)=0$ and $J'(0)=w$, differentiation gives
  
  $$
  \langle J,J\rangle(0)=0,\quad\langle J,J\rangle'(0)=2\langle J,J'\rangle(0)=0,\quad\langle J,J\rangle''(0)=2\langle J',J'\rangle(0)+2\langle J'',J\rangle(0)=2.
  $$
  
  Since $J''(0)=-R(\gamma',J)\gamma'(0)=0$, we have
  $\langle J,J\rangle'''(0)=6\langle J',J''\rangle(0)+2\langle J''',J\rangle(0)=0$. Now
  we need the fact
  
  $$
  \nabla_{\gamma'}(R(\gamma',J)\gamma')(0)=R(\gamma',J')\gamma'(0).\qquad(4)
  $$
  
  For any field $W$, at $t=0$,
  
  $$
  \Big\langle\frac{D}{dt}(R(\gamma',J)\gamma'),W\Big\rangle
  =\frac{d}{dt}\langle R(\gamma',W)\gamma',J\rangle-\langle R(\gamma',J)\gamma',W'\rangle
  =\Big\langle\frac{D}{dt}(R(\gamma',W)\gamma'),J\Big\rangle+\langle R(\gamma',W)\gamma',J'\rangle
  =\langle R(\gamma',J')\gamma',W\rangle,
  $$
  
  which proves (4). The Jacobi equation then gives
  $J'''(0)=-R(\gamma',J')\gamma'(0)$. Therefore,
  
  $$
  \langle J,J\rangle''''(0)=8\langle J',J'''\rangle(0)+6\langle J'',J''\rangle(0)+2\langle J'''',J\rangle(0)=-8\langle J',R(\gamma',J')\gamma'\rangle(0)=-8\langle R(v,w)v,w\rangle.
  $$
  
  Taylor's theorem now gives (3).
\end{proof}

\begin{remark}[curvature derivative identity at the initial point]
  \label{rem:dc-ch5-2-8}
  \dcref{ch5:2.8}
  Identity (4) also follows by covariantly differentiating the curvature
  tensor along $\gamma$ and using $J(0)=0$.
\end{remark}

\begin{corollary}[sectional-curvature form of the norm expansion]
  \label{cor:dc-ch5-2-9}
  \lean{Riemannian.Jacobi.chartMetricInner_chartCurvatureOp_eq_sectionalCurvature, Riemannian.Jacobi.chartMetricInner_chartCurvatureOp_eq_curvatureFormAt}
  \uses{def:dc-ch4-3-2, prop:dc-ch5-2-7}
  \leanok
  \dcref{ch5:2.9}
  If $\gamma:[0,\ell]\to M$ is parametrized by arc length (i.e. $|v|=1$) and
  $\langle w,v\rangle=0$, the expression $\langle R(v,w)v,w\rangle$ is the sectional
  curvature at $p$ with respect to the plane $\sigma$ generated by $v$ and $w$.
  Therefore, in this situation,

  $$
  |J(t)|^2=t^2-\tfrac{1}{3}K(p,\sigma)t^4+R(t).\qquad(5)
  $$

\end{corollary}

\begin{lemma}[square root of the $|J|^2$ expansion (analytic core)]
  \label{lem:dc-ch5-2-10-sqrt}
  \lean{Riemannian.Jacobi.sqrt_isLittleO_of_sq_isLittleO}
  \uses{cor:dc-ch5-2-9}
  \leanok
  \dcref{ch5:2.10}
  Let $g\ge0$ satisfy $g(t)=t^2-ct^4+o(t^4)$ as $t\to0$. Then
  $$
  \sqrt{g(t)}=t-\tfrac{c}{2}\,t^3+o(t^3),\qquad t\to 0^+.
  $$
  The one-sided limit is necessary because $\sqrt{t^2}=|t|$.
\end{lemma}
\begin{proof}

  \sketch
  Put $p(t)=t-\tfrac c2t^3$. Then $g-p^2=o(t^4)$ and
  $\sqrt g-p=(g-p^2)/(\sqrt g+p)$. For sufficiently small $t>0$, the denominator
  is at least $t/2$, so the quotient is $o(t^3)$.
\end{proof}

\begin{corollary}[norm expansion and sectional curvature]
  \label{cor:dc-ch5-2-10}
  \lean{Riemannian.Jacobi.norm_jacobi_frame_isLittleO}
  \uses{prop:dc-ch5-2-7, lem:dc-ch5-2-10-sqrt, cor:dc-ch5-2-9}
  \leanok
  \dcref{ch5:2.10}
  With the same conditions as in the previous corollary,

  $$
  |J(t)|=t-\tfrac{1}{6}K(p,\sigma)t^3+\tilde R(t),\quad\lim_{t\to 0}\frac{\tilde R(t)}{t^3}=0.\qquad(6)
  $$

\end{corollary}

\begin{remark}[curvature controls radial spreading]
  \label{rem:dc-ch5-2-11}
  \dcref{ch5:2.11}
  Let
  $$
  f(t,s)=\exp_p tv(s),\quad t\in[0,\delta],\quad s\in(-\varepsilon,\varepsilon),
  $$
  where $|v(s)|=1$, $v(0)=v$, $v'(0)=w$, and $\delta>0$ is chosen so that
  $\exp_p(tv(s))$ is defined. The variation of the radial lines in $T_pM$ has
  norm
  $$
  \Big|\Big(\frac{\partial}{\partial s}tv(s)\Big)(0)\Big|=|tw|=t.
  $$
  By (6), the norm of the corresponding geodesic variation is
  $t-\frac16K(p,\sigma)t^3+o(t^3)$. Thus positive sectional curvature decreases,
  and negative sectional curvature increases, the infinitesimal radial
  separation relative to $T_pM$.
\end{remark}

\section{Conjugate points}

\begin{definition}[conjugate point, multiplicity]
  \label{def:dc-ch5-3-1}
  \lean{Riemannian.Jacobi.IsConjugatePointAt, Riemannian.Jacobi.IsGeodesicConjugatePointAt}
  \uses{def:dc-ch5-2-1}
  \leanok
  \dcref{ch5:3.1}
  Let $\gamma:[0,a]\to M$ be a geodesic and $t_0\in(0,a]$. The point
  $\gamma(t_0)$ is \emph{conjugate} to $\gamma(0)$ along $\gamma$ if some
  nonzero Jacobi field $J$ satisfies $J(0)=J(t_0)=0$. Its \emph{multiplicity}
  is the dimension of the space of such fields. Conjugacy along a geodesic
  segment is symmetric under reversal of the segment.
\end{definition}

\begin{remark}[dimension bound for conjugate multiplicity]
  \label{rem:dc-ch5-3-2}
  \lean{Riemannian.Jacobi.conjugateMultiplicity, Riemannian.Jacobi.conjugateMultiplicity_le_finrank, Riemannian.Jacobi.conjugateMultiplicity_le_finrank_sub_one}
  \uses{rem:dc-ch5-2-2}
  \dcref{ch5:3.2}
  If $\dim M=n$, the Jacobi fields along $\gamma$ that vanish at $0$ form an
  $n$-dimensional space. A conjugate point along a nonconstant geodesic has
  multiplicity at most $n-1$.
\end{remark}
\begin{proof}

  \sketch
  The initial-derivative map $J\mapsto J'(0)$ identifies this space with
  $T_{\gamma(0)}M$; in particular, fields $J_i$ vanishing at $0$ are linearly
  independent exactly when the vectors $J_i'(0)$ are. At a positive endpoint,
  the field $t\gamma'(t)$ of \cref{rem:dc-ch5-2-2} does not vanish, so the kernel
  of endpoint evaluation is a proper subspace and has dimension at most $n-1$.
\end{proof}

\begin{example}[conjugate points in positive constant curvature]
  \label{ex:dc-ch5-3-3}
  \lean{Riemannian.Jacobi.isConjugatePointAt_constCurvature_pos, Riemannian.Jacobi.conjugateMultiplicity_constCurvature_pos_eq}
  \uses{ex:dc-ch5-2-3, def:dc-ch5-3-1, rem:dc-ch5-3-2, prop:dc-ch5-3-5}
  \leanok
  \dcref{ch5:3.3}
  Let $M^n$ have constant sectional curvature $K_0>0$, and let
  $\gamma:[0,\pi/\sqrt{K_0}]\to M$ be unit speed. Then
  $\gamma(\pi/\sqrt{K_0})$ is conjugate to $\gamma(0)$ with multiplicity $n-1$.
  In particular, on the unit sphere $S^n$, the antipode $\gamma(\pi)$ has
  multiplicity $n-1$ along every unit-speed geodesic from $\gamma(0)$.
\end{example}
\begin{proof}

  \sketch
  For each $w_0\perp\gamma'(0)$, let $w$ be its parallel translate along
  $\gamma$. By \cref{ex:dc-ch5-2-3},
  $$
  J(t)=\frac{\sin(t\sqrt{K_0})}{\sqrt{K_0}}\,w(t)
  $$
  is Jacobi and vanishes at both endpoints. These fields form an
  $(n-1)$-dimensional family, while \cref{rem:dc-ch5-3-2} gives the reverse
  bound on multiplicity.
\end{proof}

\begin{definition}[conjugate locus]
  \label{def:dc-ch5-3-4}
  \lean{Riemannian.Jacobi.IsFirstConjugatePointAt, Riemannian.Jacobi.conjugateLocus}
  \uses{def:dc-ch5-3-1}
  \dcref{ch5:3.4}
  The \emph{conjugate locus} $C(p)$ of $p\in M$ is the set of first conjugate
  points along all geodesics starting at $p$. For the unit sphere,
  $C(p)=\{-p\}$.
\end{definition}

\begin{proposition}[conjugate points and critical points of the exponential map]
  \label{prop:dc-ch5-3-5}
  \lean{Riemannian.Jacobi.expDifferential_injective_iff_not_conjugate, Riemannian.Jacobi.injective_jacobiEndpointOfVel_iff_not_conjugate, Riemannian.Jacobi.conjugateMultiplicity, Riemannian.Jacobi.isConjugatePointAt_iff_conjugateMultiplicity_pos, Riemannian.Jacobi.conjugateMultiplicity_eq_finrank_ker_expDifferential}
  \uses{cor:dc-ch5-2-5, def:dc-ch5-3-1}
  \leanok
  \dcref{ch5:3.5}
  Let $\gamma:[0,a]\to M$ be a geodesic with $\gamma(0)=p$, and let
  $t_0\in(0,a]$. Then $q=\gamma(t_0)$ is conjugate to $p$ along $\gamma$ if and
  only if $v_0=t_0\gamma'(0)$ is a critical point of $\exp_p$. Moreover, the
  multiplicity of $q$ equals $\dim\ker(d\exp_p)_{v_0}$.
\end{proposition}
\begin{proof}
  \sketch
  Put $v=\gamma'(0)$. By \cref{cor:dc-ch5-2-5}, every Jacobi field with
  $J(0)=0$ is determined by $w=J'(0)$ and satisfies
  $$
  J(t_0)=(d\exp_p)_{t_0v}(t_0w).
  $$
  Since $t_0\ne0$, nonzero initial derivatives in the kernel correspond exactly
  to nonzero Jacobi fields vanishing at both endpoints. The map
  $w\mapsto t_0w$ is an isomorphism, so it also identifies their space with
  $\ker(d\exp_p)_{t_0v}$ and preserves dimension.
\end{proof}

\begin{proposition}[affine tangential component of a Jacobi field]
  \label{prop:dc-ch5-3-6}
  \lean{Riemannian.Jacobi.chartMetricInner_jacobi_velocity_affine}
  \uses{def:dc-ch5-2-1}
  \leanok
  \dcref{ch5:3.6}
  Let $J$ be a Jacobi field along the geodesic $\gamma:[0,a]\to M$. Then

  $$
  \langle J(t),\gamma'(t)\rangle=\langle J'(0),\gamma'(0)\rangle t+\langle J(0),\gamma'(0)\rangle,\quad t\in[0,a].
  $$

\end{proposition}
\begin{proof}
  \leanok
  \uses{def:dc-ch5-2-1}
  \sketch
  Metric compatibility, the geodesic equation, the Jacobi equation, and the
  curvature symmetries give
  $$
  \langle J',\gamma'\rangle'=\langle J'',\gamma'\rangle=-\langle R(\gamma',J)\gamma',\gamma'\rangle=0.
  $$
  Hence $\langle J',\gamma'\rangle$ is constant, and
  $$
  \langle J,\gamma'\rangle'=\langle J',\gamma'\rangle=\langle J'(0),\gamma'(0)\rangle.
  $$
  Integration yields the stated affine formula.
\end{proof}

\begin{corollary}[constancy of the tangential component]
  \label{cor:dc-ch5-3-7}
  \lean{Riemannian.Jacobi.chartMetricInner_jacobi_velocity_const_of_eq}
  \uses{prop:dc-ch5-3-6}
  \leanok
  \dcref{ch5:3.7}
  If $\langle J,\gamma'\rangle(t_1)=\langle J,\gamma'\rangle(t_2)$, $t_1,t_2\in[0,a]$,
  $t_1\neq t_2$, then $\langle J,\gamma'\rangle$ does not depend on $t$; in
  particular, if $J(0)=J(a)=0$, then $\langle J,\gamma'\rangle(t)\equiv 0$.
\end{corollary}
\begin{proof}

  \sketch
  An affine function taking the same value at two distinct times has zero
  slope. If $J$ vanishes at both endpoints, that constant value is zero.
\end{proof}

\begin{corollary}[normal Jacobi fields with zero initial value]
  \label{cor:dc-ch5-3-8}
  \lean{Riemannian.Jacobi.chartMetricInner_jacobi_velocity_eq_zero_iff, Riemannian.Jacobi.velocityFunctional, Riemannian.Jacobi.finrank_velocityPerp_eq}
  \uses{prop:dc-ch5-3-6, rem:dc-ch5-3-2}
  \leanok
  \dcref{ch5:3.8}
  Suppose $J(0)=0$. Then
  $\langle J'(0),\gamma'(0)\rangle=0$ if and only if
  $\langle J(t),\gamma'(t)\rangle=0$ for every $t$. If $\dim M=n$ and $\gamma$
  is nonconstant, the space of such Jacobi fields has dimension $n-1$.
\end{corollary}
\begin{proof}

  \sketch
  By \cref{prop:dc-ch5-3-6}, the pairing is
  $t\langle J'(0),\gamma'(0)\rangle$. Under the identification
  $J\mapsto J'(0)$, the normal fields form the hyperplane
  $\gamma'(0)^\perp\subset T_{\gamma(0)}M$.
\end{proof}

\begin{proposition}[Jacobi field with prescribed nonconjugate endpoints]
  \label{prop:dc-ch5-3-9}
  \lean{Riemannian.Jacobi.exists_unique_jacobi_of_endpoints}
  \uses{cor:dc-ch5-2-5, def:dc-ch5-3-1}
  \leanok
  \dcref{ch5:3.9}
  Let $\gamma:[0,a]\to M$ be a geodesic, let
  $V_1\in T_{\gamma(0)}M$, and let $V_2\in T_{\gamma(a)}M$. If $\gamma(a)$ is
  not conjugate to $\gamma(0)$ along $\gamma$, there is a unique Jacobi field
  $J$ such that $J(0)=V_1$ and $J(a)=V_2$.
\end{proposition}
\begin{proof}
  \leanok
  \uses{cor:dc-ch5-2-5, def:dc-ch5-3-1}
  \sketch
  For $w\in T_{\gamma(0)}M$, let $J_w$ be the Jacobi field with
  $J_w(0)=0$ and $J_w'(0)=w$. The endpoint map
  $$
  \Theta:T_{\gamma(0)}M\longrightarrow T_{\gamma(a)}M,\qquad
  \Theta(w)=J_w(a),
  $$
  is injective by nonconjugacy, hence is an isomorphism because both spaces
  have dimension $\dim M$. Let $J_0$ be the Jacobi field with initial data
  $(V_1,0)$, choose $w$ with $\Theta(w)=V_2-J_0(a)$, and set $J=J_0+J_w$.
  The difference of two solutions lies in $\ker\Theta$, proving uniqueness.
\end{proof}

\begin{corollary}[normal Jacobi fields form an endpoint basis]
  \label{cor:dc-ch5-3-10}
  \lean{Riemannian.Jacobi.jacobiConjugateBasis, Riemannian.Jacobi.jacobiConjugateEquiv, Riemannian.Jacobi.jacobiEndpointOfVel_map_velocityPerp_eq, Riemannian.Jacobi.metricInner_jacobiJ_velocity_eq_zero}
  \uses{cor:dc-ch5-3-8, prop:dc-ch5-3-5, prop:dc-ch5-3-6, def:dc-ch5-3-1}
  \leanok
  \dcref{ch5:3.10}
  Let $\gamma:[0,a]\to M$ be a geodesic in $M$, $\dim M=n$, and let
  $\mathcal{J}^\perp$ be the space of Jacobi fields with $J(0)=0$,
  $J'(0)\perp\gamma'(0)$. Let $\{J_1,\dots,J_{n-1}\}$ be a basis of
  $\mathcal{J}^\perp$. If $\gamma(t)$, $t\in(0,a]$, is not conjugate to $\gamma(0)$,
  then $\{J_1(t),\dots,J_{n-1}(t)\}$ is a basis for the orthogonal complement
  $\{\gamma'(t)\}^\perp\subset T_{\gamma(t)}M$ of $\gamma'(t)$.
\end{corollary}
\begin{proof}

  \sketch
  By \cref{cor:dc-ch5-3-8}, evaluation at $t$ sends
  $\mathcal J^\perp$ into $\gamma'(t)^\perp$. Nonconjugacy makes this endpoint
  map injective by \cref{prop:dc-ch5-3-5}. Both spaces have dimension $n-1$,
  so it is an isomorphism and carries any basis of $\mathcal J^\perp$ to the
  displayed basis.
\end{proof}

\begin{example}[conjugate loci of the sphere and ellipsoid]
  \label{ex:dc-ch5-3-4}
  \dcref{ch5:3.4}
  \uses{def:dc-ch5-3-4}
  \notready
  On $S^n$, $C(p)=\{-p\}$, for all $p$. The case of $S^n$, however, is not typical.
  A more typical example is given by the ellipsoid, where $C(p)$ is, in general, a
  curve with four singular points (see Fig. 4 of Chap. 13; cf. Braunmühl, A.,
  "Geodätische Linien auf dreiachsigen Flächen 2-Grades", Math. Ann. 20 (1882),
  557-586).
\end{example}

\chapter{Isometric Immersions}
\label{chap:dc-isometric-immersions}
An isometric immersion $f:M^n\to\overline M^{n+m}$ induces the metric
$\langle v,w\rangle=\langle df(v),df(w)\rangle$ on $M$.
\section{The second fundamental form}

Locally identify an immersed patch with its image in $\overline M$. Write $T_p\overline M=T_pM\oplus(T_pM)^\perp$, and decompose fields as $X=X^T+X^N$.

\begin{definition}[the identified setup: tangent distribution and orthogonal splitting]
  \label{def:dc-ch6-2-patch}
  \lean{Riemannian.DCImmersedPatch.normalSpace, Riemannian.DCImmersedPatch.mem_normalSpace_iff, Riemannian.DCImmersedPatch.normalProj, Riemannian.DCImmersedPatch.IsTangentField, Riemannian.DCImmersedPatch.IsNormalField, Riemannian.DCImmersedPatch.isTangentField_bracketField}
  \leanok
  Let $(\overline{M},\langle\,,\rangle)$ be a Riemannian manifold. An
  \emph{immersed patch} of dimension $n$ in $\overline{M}$ — the data furnished
  by an isometric immersion $f:M^n\to\overline{M}$ after the identification of
  $U\subset M$ with $f(U)\subset\overline{M}$ — consists of a family of
  subspaces $T_pM\subseteq T_p\overline{M}$, $p\in\overline{M}$, of constant
  dimension $n$, such that:

  (i) the orthogonal splitting $T_p\overline{M}=T_pM\oplus(T_pM)^\perp$ is
  \emph{differentiable}: every differentiable vector field $X$ on
  $\overline{M}$ admits a differentiable field $X^T$ with $X^T(p)\in T_pM$ and
  $\langle v,\,X(p)-X^T(p)\rangle=0$ for every $v\in T_pM$ and every $p$;

  (ii) fields tangent to $M$ are closed under the Lie bracket: if
  $X(p),Y(p)\in T_pM$ for all $p$, then $[X,Y](p)\in T_pM$ for all $p$ (on $M$,
  $[\overline{X},\overline{Y}]=[X,Y]$).

  The \emph{normal space} at $p$ is
  $(T_pM)^\perp=\{v\in T_p\overline{M};\ \langle w,v\rangle=0\ \text{for all}\
  w\in T_pM\}$, and the \emph{normal part} of a field is $X^N=X-X^T$, so that
  $X=X^T+X^N$. A field is \emph{tangent} when $X(p)\in T_pM$ for every $p$ and
  \emph{normal} when $X(p)\in(T_pM)^\perp$ for every $p$; under the
  identification, the tangent fields are the local fields $\mathcal{X}(U)$ of
  $M$ and the normal fields are $\mathcal{X}(U)^\perp$.
\end{definition}

\begin{lemma}[uniqueness and tensoriality of the tangential projection]
  \label{lem:dc-ch6-2-proj}
  \lean{Riemannian.DCImmersedPatch.eq_zero_of_mem_tang_of_mem_normalSpace, Riemannian.DCImmersedPatch.tangentProj_eq_at, Riemannian.DCImmersedPatch.tangentProj_congr_apply, Riemannian.DCImmersedPatch.tangentProj_add, Riemannian.DCImmersedPatch.tangentProj_smul, Riemannian.DCImmersedPatch.tangentProj_apply_of_mem, Riemannian.DCImmersedPatch.tangentProj_apply_of_mem_normalSpace, Riemannian.DCImmersedPatch.normalProj_add, Riemannian.DCImmersedPatch.normalProj_smul}
  \uses{def:dc-ch6-2-patch}
  \leanok
  The splitting $X(p)=X^T(p)+X^N(p)$ at a point is the unique decomposition of
  $X(p)$ into a tangent and a normal vector: if $u\in T_pM$ and
  $\langle v,\,X(p)-u\rangle=0$ for all $v\in T_pM$, then $u=X^T(p)$. In
  particular $X^T(p)$ depends only on the value $X(p)$; the projections
  $X\mapsto X^T$ and $X\mapsto X^N$ are additive and
  $\mathcal{D}(\overline{M})$-homogeneous; $X^T=X$ and $X^N=0$ for tangent
  fields; and $X^T=0$, $X^N=X$ for normal fields.
\end{lemma}
\begin{proof}
  \leanok
  \sketch
  If $d\in T_pM$ is also normal at $p$ then $\langle d,d\rangle=0$, so $d=0$ by
  positive-definiteness: the splitting is direct. Let $u$ be as in the
  statement. Then $X^T(p)-u\in T_pM$, and for every tangent $v$,
  $\langle v,\,X^T(p)-u\rangle
  =\langle v,\,X(p)-u\rangle-\langle v,\,X(p)-X^T(p)\rangle=0$, so $X^T(p)-u$
  is also normal, hence zero. Every remaining assertion is an instance of this
  uniqueness: $X^T(p)+Y^T(p)$ is tangent and
  $(X+Y)(p)-(X^T(p)+Y^T(p))=(X(p)-X^T(p))+(Y(p)-Y^T(p))$ is orthogonal to
  $T_pM$, so $(X+Y)^T=X^T+Y^T$; likewise $f(p)X^T(p)$ is tangent with
  $f(p)X(p)-f(p)X^T(p)=f(p)(X(p)-X^T(p))$ orthogonal to $T_pM$, so
  $(fX)^T=fX^T$; if $X(p)\in T_pM$ then $u=X(p)$ satisfies the
  characterization, and if $X(p)$ is normal then $u=0$ does. The assertions
  for $X^N=X-X^T$ follow by subtraction.
\end{proof}

\begin{definition}[the induced connection]
  \label{def:dc-ch6-2-induced-connection}
  \lean{Riemannian.DCImmersedPatch.inducedCov}
  \uses{def:dc-ch6-2-patch, thm:dc-ch2-3-6}
  \leanok
  Let $\overline{\nabla}$ be the Riemannian connection of $\overline{M}$. The
  \emph{induced connection} of the patch is
  $\nabla_X Y=(\overline{\nabla}_X Y)^T$.
\end{definition}

\begin{proposition}[the induced connection is the Riemannian connection of the induced metric]
  \label{prop:dc-ch6-2-induced-connection}
  \lean{Riemannian.DCImmersedPatch.isTangentField_inducedCov, Riemannian.DCImmersedPatch.inducedCov_add_left, Riemannian.DCImmersedPatch.inducedCov_smul_left, Riemannian.DCImmersedPatch.inducedCov_add_right, Riemannian.DCImmersedPatch.inducedCov_smul_right, Riemannian.DCImmersedPatch.inducedCov_sub_swap, Riemannian.DCImmersedPatch.inducedCov_compat}
  \uses{def:dc-ch6-2-induced-connection, lem:dc-ch6-2-proj}
  \leanok
  The induced connection is tangent-valued and, on tangent fields, satisfies
  the axioms of an affine connection: it is additive in both arguments,
  $\mathcal{D}$-homogeneous in the direction ($\nabla_{fX}Y=f\nabla_XY$), and
  satisfies the Leibniz rule $\nabla_X(fY)=f\nabla_XY+X(f)\,Y$ for $Y$ tangent.
  Moreover it is symmetric, $\nabla_XY-\nabla_YX=[X,Y]$ for tangent $X,Y$, and
  compatible with the metric on tangent fields,
  $X\langle Y,Z\rangle=\langle\nabla_XY,Z\rangle+\langle Y,\nabla_XZ\rangle$.
  By the uniqueness of the Riemannian connection (\cref{thm:dc-ch2-3-6}) it is
  the Riemannian connection of the induced metric.
\end{proposition}
\begin{proof}
  \leanok
  \sketch
  Additivity and $\mathcal{D}$-homogeneity in the direction slot follow from
  the corresponding laws of $\overline{\nabla}$ and of the projection
  (\cref{lem:dc-ch6-2-proj}). For the Leibniz rule,
  $\overline{\nabla}_X(fY)=f\overline{\nabla}_XY+X(f)\,Y$; taking tangential
  parts, the term $X(f)\,Y$ is already tangent, giving
  $\nabla_X(fY)=f\nabla_XY+X(f)\,Y$. For symmetry, the difference
  $\overline{\nabla}_XY-\overline{\nabla}_YX=[X,Y]$ is tangent for tangent
  $X,Y$ (\cref{def:dc-ch6-2-patch}(ii)), so applying the projection gives
  $\nabla_XY-\nabla_YX=[X,Y]$. For compatibility, write
  $\overline{\nabla}_XY=(\overline{\nabla}_XY)^T+(\overline{\nabla}_XY)^N$ in
  $X\langle Y,Z\rangle
  =\langle\overline{\nabla}_XY,Z\rangle+\langle Y,\overline{\nabla}_XZ\rangle$;
  the normal parts pair to zero against the tangent fields $Z$ and $Y$.
\end{proof}

\begin{definition}[the second fundamental form field]
  \label{def:dc-ch6-2-B}
  \lean{Riemannian.DCImmersedPatch.secondFundForm, Riemannian.DCImmersedPatch.cov_eq_inducedCov_add_secondFundForm, Riemannian.DCImmersedPatch.secondFundForm_mem}
  \uses{def:dc-ch6-2-induced-connection}
  \leanok
  The \emph{second fundamental form} of the patch is
  $B(X,Y)=\overline{\nabla}_XY-\nabla_XY=(\overline{\nabla}_XY)^N$, a
  normal-valued form on fields, giving the \emph{Gauss decomposition}
  $\overline{\nabla}_XY=\nabla_XY+B(X,Y)$.
\end{definition}

\begin{proposition}[the second fundamental form mapping $B$]
  \label{prop:dc-ch6-2-1}
  \lean{Riemannian.DCImmersedPatch.secondFundForm_add_left, Riemannian.DCImmersedPatch.secondFundForm_add_right, Riemannian.DCImmersedPatch.secondFundForm_sub_left, Riemannian.DCImmersedPatch.secondFundForm_smul_left, Riemannian.DCImmersedPatch.secondFundForm_smul_right, Riemannian.DCImmersedPatch.secondFundForm_symm, Riemannian.DCImmersedPatch.isNormalField_secondFundForm}
  \uses{def:dc-ch6-2-B, lem:dc-ch6-2-proj, prop:dc-ch6-2-induced-connection}
  \leanok
  \dcref{ch6:2.1}
  If $X,Y\in\mathcal{X}(U)$, the mapping
  $B:\mathcal{X}(U)\times\mathcal{X}(U)\to\mathcal{X}(U)^\perp$ given by

  $$
  B(X,Y)=\overline{\nabla}_{\overline{X}}\overline{Y}-\nabla_X Y
  $$

  is bilinear and symmetric.
\end{proposition}
\begin{proof}
  \leanok
  \sketch
  From linearity of a connection, $B$ is additive in $X$ and $Y$ and
  $B(fX,Y)=fB(X,Y)$, $f\in\mathcal{D}(U)$. To show $B(X,fY)=fB(X,Y)$, denote the
  extension of $f$ to $\overline{U}$ by $\overline{f}$; then

  $$
  B(X,fY)=\overline{\nabla}_{\overline{X}}(\overline{f}\,\overline{Y})-\nabla_X(fY)=\overline{f}\,\overline{\nabla}_{\overline{X}}\overline{Y}-f\nabla_X Y+\overline{X}(\overline{f})\overline{Y}-X(f)Y.
  $$

  Since $f=\overline{f}$, $\overline{X}(\overline{f})=X(f)$, and $Y=\overline{Y}$ on
  $M$, the last two terms cancel, giving $B(X,fY)=fB(X,Y)$. To show $B$ is
  symmetric, by symmetry of the Riemannian connection,

  $$
  B(X,Y)=\overline{\nabla}_{\overline{X}}\overline{Y}-\nabla_X Y=\overline{\nabla}_{\overline{Y}}\overline{X}+[\overline{X},\overline{Y}]-\nabla_Y X-[X,Y].
  $$

  Since $[\overline{X},\overline{Y}]=[X,Y]$ on $M$, we conclude $B(X,Y)=B(Y,X)$.
\end{proof}

\begin{lemma}[pointwise locality of $B$ and of the covariant derivative]
  \label{lem:dc-ch6-2-locality}
  \lean{Riemannian.AffineConnection.cov_congr_apply_left, Riemannian.exists_decomposition_of_apply_eq_zero, Riemannian.DCImmersedPatch.exists_tangent_decomposition_of_apply_eq_zero, Riemannian.DCImmersedPatch.secondFundForm_congr_apply, Riemannian.DCImmersedPatch.exists_isTangentField_eq, Riemannian.DCImmersedPatch.exists_isNormalField_eq}
  \uses{prop:dc-ch6-2-1, def:dc-ch6-2-patch}
  \leanok
  Because $B$ is bilinear over the differentiable functions, the value
  $B(X,Y)(p)$ depends only on the values $X(p)$ and $Y(p)$: if $X(p)=X'(p)$ and
  the tangent fields $Y,Y'$ satisfy $Y(p)=Y'(p)$, then
  $B(X,Y)(p)=B(X',Y')(p)$. Likewise $\overline{\nabla}_XZ(p)$ depends on the
  direction only through $X(p)$. Moreover every tangent vector $v\in T_pM$
  (resp.\ normal vector $\eta\in(T_pM)^\perp$) is the value at $p$ of a tangent
  (resp.\ normal) field, so the pointwise objects are
  well-defined.
\end{lemma}
\begin{proof}
  \leanok
  \sketch
  A field $\sigma$ with $\sigma(p)=0$ can be written, near $p$, as a finite sum
  $\sum_i f_i W_i$ with differentiable functions $f_i$ vanishing at $p$: expand
  $\sigma$ in a local frame of the tangent bundle of $\overline{M}$ around $p$
  and cut the coefficients and the frame off by a bump function equal to $1$
  near $p$; the discrepancy is a field vanishing on a neighbourhood of $p$.
  For a mapping $\Phi$ that is additive and $\mathcal{D}$-homogeneous in
  $\sigma$, each summand contributes $f_i(p)\,\Phi(W_i)(p)=0$, and a field
  vanishing near $p$ contributes $0$ because it equals $(1-\chi)\tau$ for a
  bump $\chi$ with $\chi(p)=1$ supported in the vanishing neighbourhood.
  Applying this to $\Phi=\overline{\nabla}_{(\cdot)}Z$ and, in the second slot,
  to $\Phi=B(X,\cdot)$ — where for tangent $\sigma$ the frame decomposition
  may be projected tangentially, so the $W_i$ can be taken tangent — gives the
  locality claims. For the extensions: any vector extends to a global
  differentiable field, and composing with the tangential (resp.\ normal)
  projection fixes the value at $p$ when it is tangent (resp.\ normal).
\end{proof}

By \cref{lem:dc-ch6-2-locality}, the value $B(X,Y)(p)$ depends only on the
values $X(p)$ and $Y(p)$. Now let $p\in M$ and $\eta\in(T_pM)^\perp$. The
mapping $H_\eta:T_pM\times T_pM\to\mathbb{R}$ given by

$$
H_\eta(x,y)=\langle B(x,y),\eta\rangle,\quad x,y\in T_pM,
$$

is, by Proposition 2.1, a symmetric bilinear form.

\begin{definition}[second fundamental form $\mathit{II}_\eta$]
  \label{def:dc-ch6-2-2}
  \lean{Riemannian.DCImmersedPatch.secondFundFormAt, Riemannian.DCImmersedPatch.secondFundScalarAt, Riemannian.DCImmersedPatch.secondFundQuadAt, Riemannian.DCImmersedPatch.shapeOperatorAt, Riemannian.DCImmersedPatch.inner_shapeOperatorAt, Riemannian.DCImmersedPatch.inner_shapeOperatorAt_symm, Riemannian.DCImmersedPatch.secondFundScalarAt_symm, Riemannian.DCImmersedPatch.secondFundFormAt_symm}
  \uses{lem:dc-ch6-2-locality, prop:dc-ch6-2-1, def:dc-ch6-3-normal-connection, lem:dc-ch6-3-weingarten}
  \leanok
  \dcref{ch6:2.2}
  The quadratic form $\mathit{II}_\eta$ defined on $T_pM$ by
  $$
  \mathit{II}_\eta(x)=H_\eta(x,x)
  $$
  is called \emph{the second fundamental form} of $f$ at $p$ along the normal vector
  $\eta$. Sometimes the expression \emph{second fundamental form} also designates the
  mapping $B$. The bilinear mapping $H_\eta$ is associated to a linear self-adjoint
  operator $S_\eta:T_pM\to T_pM$ by
  $$
  \langle S_\eta(x),y\rangle=H_\eta(x,y)=\langle B(x,y),\eta\rangle.
  $$
\end{definition}

\begin{proposition}[the operator $S_\eta$ via the covariant derivative]
  \label{prop:dc-ch6-2-3}
  \lean{Riemannian.DCImmersedPatch.shapeOperatorAt_eq_shapeOperator}
  \uses{def:dc-ch6-2-2, def:dc-ch6-3-normal-connection, lem:dc-ch6-3-weingarten, lem:dc-ch6-2-locality}
  \leanok
  \dcref{ch6:2.3}
  Let $p\in M$, $x\in T_pM$ and $\eta\in(T_pM)^\perp$. Let $N$ be a local extension
  of $\eta$ normal to $M$. Then

  $$
  S_\eta(x)=-(\overline{\nabla}_x N)^T.
  $$
\end{proposition}
\begin{proof}
  \leanok
  \sketch
  Let $y\in T_pM$ and let $X,Y$ be local extensions of $x,y$, tangent to
  $M$. Then $\langle N,Y\rangle=0$, and therefore
  
  $$
  \langle S_\eta(x),y\rangle=\langle B(X,Y)(p),N\rangle=\langle\overline{\nabla}_X Y-\nabla_X Y,N\rangle(p)=\langle\overline{\nabla}_X Y,N\rangle(p)=-\langle Y,\overline{\nabla}_X N\rangle(p)=\langle-\overline{\nabla}_x N,y\rangle,
  $$
  
  for all $y\in T_pM$.
\end{proof}

\begin{example}[hypersurface; principal curvatures; Gauss spherical mapping]
  \label{ex:dc-ch6-2-4}
  \lean{Riemannian.gaussMap, Riemannian.gaussMap_norm, Riemannian.normalSpace_eq_span_of_unit_normal, Riemannian.mem_tang_of_inner_unit_normal_eq_zero, Riemannian.fderiv_gaussMap_mem_tang, Riemannian.fderiv_gaussMap_eq_neg_shapeOperatorAt, Riemannian.fderiv_gaussMap_mem_sphereTang}
  \uses{def:dc-ch6-2-2, prop:dc-ch6-2-3, def:dc-ch6-2-4-principal, lem:dc-ch6-2-8-euclidean-opens, ex:dc-ch6-2-8}
  \leanok
  \dcref{ch6:2.4}
  An immersion $f:M^n\to\overline{M}^{n+1}$ has image a \emph{hypersurface}.
  For $p\in M$ and a unit normal $\eta\in(T_pM)^\perp$, the symmetric operator
  $S_\eta:T_pM\to T_pM$ has an orthonormal basis of
  eigenvectors $\{e_1,\dots,e_n\}$ of $T_pM$ with real eigenvalues
  $\lambda_1,\dots,\lambda_n$, $S_\eta(e_i)=\lambda_i e_i$. If $M$ and $\overline{M}$
  are both orientable and oriented, then $\eta$ is uniquely determined by requiring
  $\{e_1,\dots,e_n\}$ to be in the orientation of $M$ and $\{e_1,\dots,e_n,\eta\}$ in
  the orientation of $\overline{M}$. The $e_i$ are \emph{principal directions} and the
  $\lambda_i=k_i$ are \emph{principal curvatures} of $f$. The symmetric functions of the
  $\lambda_i$ are invariants of the immersion: $\det(S_\eta)=\lambda_1\cdots\lambda_n$
  is the \emph{Gauss-Kronecker curvature} and
  $\frac{1}{n}(\lambda_1+\cdots+\lambda_n)$ the \emph{mean curvature} of $f$.
  When $\overline{M}=\mathbb{R}^{n+1}$, a local unit normal $N$ defines the
  \emph{Gauss spherical map} $g:M^n\to S_1^n$ by $g(q)=N(q)$. Under the natural
  identification $T_qM=T_{g(q)}S_1^n$, its differential is
  $$
  dg_q(x)=\overline{\nabla}_xN=-S_{N(q)}(x).
  $$
\end{example}

\begin{remark}[compact hypersurfaces with non-vanishing Gauss-Kronecker curvature]
  \label{rem:dc-ch6-2-4-covering}
  \uses{ex:dc-ch6-2-4}
  \notready
  \notready
  If $M^n$, $n\geq2$, is connected, compact, and orientable and admits an
  immersion into $\mathbb R^{n+1}$ with nonzero Gauss--Kronecker curvature,
  then its Gauss map is a local diffeomorphism. Compactness makes it a covering
  of the simply connected sphere, hence $M\cong S^n$.
\end{remark}

\begin{definition}[principal curvatures and directions; Gauss-Kronecker and mean curvature]
  \label{def:dc-ch6-2-4-principal}
  \lean{Riemannian.DCImmersedPatch.shapeOperatorHom, Riemannian.DCImmersedPatch.shapeOperatorHom_isSymmetric, Riemannian.DCImmersedPatch.principalCurvatures, Riemannian.DCImmersedPatch.principalDirections, Riemannian.DCImmersedPatch.principalDirection, Riemannian.DCImmersedPatch.metricInner_principalDirection, Riemannian.DCImmersedPatch.shapeOperatorAt_principalDirection, Riemannian.DCImmersedPatch.secondFundScalarAt_principalDirection, Riemannian.DCImmersedPatch.gaussKroneckerCurvature, Riemannian.DCImmersedPatch.meanCurvature, Riemannian.DCImmersedPatch.trace_shapeOperatorHom, Riemannian.DCImmersedPatch.meanCurvature_eq_sum, Riemannian.DCImmersedPatch.gaussKroneckerCurvature_eq_prod}
  \uses{def:dc-ch6-2-2}
  \leanok
  Let $p\in M$ and $\eta\in(T_pM)^\perp$. Since $S_\eta:T_pM\to T_pM$ is linear
  and self-adjoint with respect to the metric, the spectral theorem yields an
  orthonormal basis of eigenvectors $\{e_1,\dots,e_n\}$ of $T_pM$ with real
  eigenvalues $\lambda_1,\dots,\lambda_n$, $S_\eta(e_i)=\lambda_ie_i$, and
  $H_\eta(e_i,e_j)=\lambda_i\delta_{ij}$. The $e_i$ are the \emph{principal
  directions} and the $\lambda_i=k_i$ the \emph{principal curvatures} of $f$ at
  $p$ along $\eta$. Their symmetric functions are invariants of the immersion:
  $\det(S_\eta)=\lambda_1\cdots\lambda_n$ is the \emph{Gauss-Kronecker
  curvature} and $\frac1n\operatorname{trace}(S_\eta)
  =\frac1n(\lambda_1+\cdots+\lambda_n)$ the \emph{mean curvature} of $f$ at $p$
  along $\eta$.
\end{definition}

\begin{theorem}[Gauss]
  \label{thm:dc-ch6-2-5}
  \lean{Riemannian.DCImmersedPatch.gauss_theorem, Riemannian.DCImmersedPatch.inducedCurvature_inner_sub_curvature_inner}
  \uses{prop:dc-ch6-2-1, def:dc-ch6-3-curvatures, def:dc-ch4-3-2, prop:dc-ch6-3-1}
  \leanok
  \dcref{ch6:2.5}
  Let $p\in M$ and let $x,y$ be orthonormal vectors in $T_pM$. Then

  $$
  K(x,y)-\overline{K}(x,y)=\langle B(x,x),B(y,y)\rangle-|B(x,y)|^2,\qquad(1)
  $$

  where $K(x,y)$ and $\overline{K}(x,y)$ denote the sectional curvatures of $M$ and
  $\overline{M}$, respectively, in the plane generated by $x$ and $y$.
\end{theorem}
\begin{proof}
  \leanok
  \sketch
  Extend $x,y$ to tangent fields $X,Y$ and choose an orthonormal normal frame
  $E_i$. Insert the Gauss decomposition in the curvature definition. The normal
  component of the bracket term pairs trivially with $Y$; pairing the two
  second-covariant-derivative terms with $Y$ gives respectively
  $-\sum_iH_i(X,X)H_i(Y,Y)$ and $-\sum_iH_i(X,Y)^2$, while the tangential
  terms reproduce the curvature of $M$. Their difference is (1). Equivalently,
  (1) is the $Z=X,T=Y$ specialization of the Gauss equation
  (\cref{prop:dc-ch6-3-1}) for an orthonormal pair.
\end{proof}

\begin{remark}[Gauss formula for a hypersurface]
  \label{rem:dc-ch6-2-6}
  \lean{Riemannian.DCImmersedPatch.secondFundFormAt_eq_secondFundScalarAt_smul, Riemannian.DCImmersedPatch.hypersurface_gauss}
  \uses{thm:dc-ch6-2-5, def:dc-ch6-2-4-principal}
  \leanok
  \dcref{ch6:2.6}
  In the case of a hypersurface $f:M^n\to\overline{M}^{n+1}$, let $p\in M$,
  $\eta\in(T_pM)^\perp$, $|\eta|=1$, and $\{e_1,\dots,e_n\}$ an orthonormal basis of $T_pM$ in
  which $S_\eta=S$ is diagonal, $S(e_i)=\lambda_ie_i$. Then $H(e_i,e_i)=\lambda_i$
  and $H(e_i,e_j)=0$ if $i\neq j$, so (1) becomes
  $$
  K(e_i,e_j)-\overline{K}(e_i,e_j)=\lambda_i\lambda_j.\qquad(4)
  $$
\end{remark}

\begin{lemma}[Koszul formula for the induced connection]
  \label{lem:dc-ch6-2-7-koszul}
  \lean{Riemannian.DCImmersedPatch.inducedCov_koszul}
  \uses{prop:dc-ch6-2-induced-connection}
  \leanok
  For fields $X,Y,Z$ tangent to the patch,
  $$
  2\langle Z,\nabla_YX\rangle
    = X\langle Y,Z\rangle + Y\langle Z,X\rangle - Z\langle X,Y\rangle
      - \langle[X,Z],Y\rangle - \langle[Y,Z],X\rangle - \langle[X,Y],Z\rangle.
  $$
  Every quantity on the right-hand side — inner products of tangent fields,
  their derivatives along tangent fields, and brackets of tangent fields — is
  intrinsic to the immersed manifold: the induced connection is determined by
  the induced metric.
\end{lemma}
\begin{proof}
  \leanok
  \sketch
  Write the compatibility identity of \cref{prop:dc-ch6-2-induced-connection}
  for the three cyclic orderings,
  $X\langle Y,Z\rangle=\langle\nabla_XY,Z\rangle+\langle Y,\nabla_XZ\rangle$,
  $Y\langle Z,X\rangle=\langle\nabla_YZ,X\rangle+\langle Z,\nabla_YX\rangle$,
  $Z\langle X,Y\rangle=\langle\nabla_ZX,Y\rangle+\langle X,\nabla_ZY\rangle$,
  and substitute the symmetry identities
  $\nabla_XY-\nabla_YX=[X,Y]$, $\nabla_XZ-\nabla_ZX=[X,Z]$,
  $\nabla_YZ-\nabla_ZY=[Y,Z]$ of the same proposition. Adding the first two
  identities and subtracting the third, all covariant-derivative terms cancel
  against bracket terms except $2\langle Z,\nabla_YX\rangle$, which gives the
  formula.
\end{proof}

\begin{lemma}[orthonormal pairs spanning a plane give equal curvature values]
  \label{lem:dc-ch6-2-7-plane}
  \lean{Riemannian.linearIndependent_pair_of_orthonormal, Riemannian.IsAlgCurvatureForm.apply_eq_of_orthonormal_of_mem_span, Riemannian.IsAlgCurvatureForm.apply_eq_of_orthonormal_of_finrank_eq_two}
  \uses{prop:dc-ch4-2-5, prop:dc-ch4-3-1}
  \leanok
  Let $B$ be an algebraic curvature form (\cref{prop:dc-ch4-2-5}) on an inner
  product space $V$, and let $\{u_1,u_2\}$ and $\{v_1,v_2\}$ be orthonormal
  pairs with $v_1,v_2\in\operatorname{span}\{u_1,u_2\}$ — in particular any two
  orthonormal pairs when $\dim V=2$. Then
  $$
  B(v_1,v_2,v_1,v_2)=B(u_1,u_2,u_1,u_2).
  $$
\end{lemma}
\begin{proof}
  \leanok
  \sketch
  Write $v_1=au_1+bu_2$ and $v_2=cu_1+du_2$. Orthonormality of the two pairs
  gives $a^2+b^2=1$, $c^2+d^2=1$ and $ac+bd=0$, so by the Lagrange identity
  $$
  (ad-bc)^2=(a^2+b^2)(c^2+d^2)-(ac+bd)^2=1,
  $$
  and the change of basis is invertible. An orthonormal pair is linearly
  independent (pair a vanishing combination against $u_1$ and $u_2$), so by
  the basis independence of the sectional curvature (\cref{prop:dc-ch4-3-1})
  $$
  \frac{B(v_1,v_2,v_1,v_2)}{|v_1\wedge v_2|^2}
    =\frac{B(u_1,u_2,u_1,u_2)}{|u_1\wedge u_2|^2},
  $$
  and both denominators equal $1$ by orthonormality. When $\dim V=2$ the pair
  $\{u_1,u_2\}$ is a basis, so the span hypothesis is automatic.
\end{proof}

\begin{definition}[pushforward of fields along a diffeomorphism of Euclidean opens]
  \label{def:dc-ch6-2-7-pushforward}
  \lean{Riemannian.opensMapExt, Riemannian.opensFDeriv, Riemannian.opensPushforward, Riemannian.fderiv_extendZero_comp, Riemannian.opensFDeriv_symm_comp}
  \uses{lem:dc-ch6-2-8-euclidean-opens}
  \leanok
  Let $S\subseteq F$ and $T\subseteq G$ be open subsets of Euclidean spaces and
  $\varphi:S\to T$ a smooth map. Its \emph{derivative} $d\varphi_p:F\to G$ at
  $p\in S$ is the derivative of any local extension of $\varphi$ to the
  ambient space; it obeys the chain rule
  $d(g\circ\varphi)_p=dg_{\varphi(p)}\circ d\varphi_p$, and if $\varphi$ is a
  diffeomorphism with inverse $\psi$ then
  $d\psi_{\varphi(p)}\circ d\varphi_p=\mathrm{id}$. For a diffeomorphism
  $\varphi$ and a smooth vector field $X$ on $S$, the \emph{pushforward}
  $$
  (\varphi_*X)(q)=d\varphi_{\varphi^{-1}(q)}\bigl(X(\varphi^{-1}(q))\bigr)
  $$
  is a smooth vector field on $T$ (the derivative of a smooth map depends
  smoothly on the base point).
\end{definition}

\begin{lemma}[naturality of derivations and brackets under pushforward]
  \label{lem:dc-ch6-2-7-pushforward-naturality}
  \lean{Riemannian.dir_opensPushforward, Riemannian.DCLieBracket_opensPushforward, Riemannian.bracketField_opensPushforward, Riemannian.metricInner_field_contMDiff_opens}
  \uses{def:dc-ch6-2-7-pushforward, lem:dc-ch0-5-2}
  \leanok
  Let $\varphi$ be a diffeomorphism of Euclidean opens. For smooth vector
  fields $X,Y$ and a smooth scalar $h$,
  $$
  (\varphi_*X)(h\circ\varphi^{-1})=(Xh)\circ\varphi^{-1}
  \qquad\text{and}\qquad
  [\varphi_*X,\varphi_*Y]=\varphi_*[X,Y].
  $$
\end{lemma}
\begin{proof}
  \leanok
  \sketch
  The first identity is the chain rule of \cref{def:dc-ch6-2-7-pushforward}
  together with $d\psi_{\varphi(p)}\circ d\varphi_p=\mathrm{id}$, where
  $\psi=\varphi^{-1}$. For the bracket, compute the derivative of the
  pushforward field $\varphi_*Y=(d\varphi\circ Y)\circ\psi$ at
  $q=\varphi(p)$ by the chain and product rules:
  $$
  d(\varphi_*Y)_q\bigl(d\varphi_p(X_p)\bigr)
    = d^2\varphi_p(X_p,Y_p)+d\varphi_p\bigl(dY_p(X_p)\bigr).
  $$
  Subtracting the same expression with $X$ and $Y$ interchanged, the
  second-derivative terms cancel by symmetry of $d^2\varphi_p$,
  leaving
  $[\varphi_*X,\varphi_*Y](q)
    =d\varphi_p\bigl(dY_p(X_p)-dX_p(Y_p)\bigr)=d\varphi_p([X,Y](p))$,
  the bracket of fields on a Euclidean open being the commutator of directional
  derivatives.
\end{proof}

\begin{definition}[isometry of immersed patches]
  \label{def:dc-ch6-2-7-isometry}
  \lean{Riemannian.DCPatchIsometry.isTangentField_opensPushforward, Riemannian.DCPatchIsometry.metricInner_opensPushforward}
  \uses{def:dc-ch6-2-patch, def:dc-ch1-2-2, def:dc-ch6-2-7-pushforward}
  \leanok
  Let $D$ and $D'$ be immersed patches in open subsets of Euclidean spaces. An
  \emph{isometry} $\varphi:D\to D'$ is a diffeomorphism of the ambient opens
  whose derivative carries each tangent space $T_pM$ into
  $T_{\varphi(p)}M'$ (and, for the inverse, $T_qM'$ into
  $T_{\varphi^{-1}(q)}M$) and preserves the inner product of tangent vectors:
  $$
  \langle d\varphi_p(v),d\varphi_p(w)\rangle=\langle v,w\rangle
  \qquad\text{for } v,w\in T_pM.
  $$
  No condition is imposed on $d\varphi$ transverse to the tangent distributions;
  tangent fields push forward to tangent fields and tangent inner products are
  preserved pointwise.
\end{definition}

\begin{lemma}[isometries preserve the induced connection]
  \label{lem:dc-ch6-2-7-connection}
  \lean{Riemannian.DCPatchIsometry.inducedCov_opensPushforward}
  \uses{def:dc-ch6-2-7-isometry, lem:dc-ch6-2-7-koszul, lem:dc-ch6-2-7-pushforward-naturality}
  \leanok
  Let $\varphi:D\to D'$ be an isometry of immersed patches in flat Euclidean
  ambients, and let $\nabla,\nabla'$ be the induced connections. Then for
  tangent fields $X,Y$,
  $$
  \nabla'_{\varphi_*X}(\varphi_*Y)=\varphi_*(\nabla_XY).
  $$
\end{lemma}
\begin{proof}
  \leanok
  \sketch
  Both sides are tangent fields, so by positive-definiteness it suffices that
  their inner products against every tangent vector at every $q=\varphi(p)$
  agree. A tangent vector $w\in T_qM'$ is $d\varphi_p(z)$ for the tangent
  vector $z=d\varphi^{-1}_q(w)\in T_pM$; extend $z$ to a tangent field $Z$.
  By the Koszul formula (\cref{lem:dc-ch6-2-7-koszul}) applied upstairs to
  the tangent fields $\varphi_*X,\varphi_*Y,\varphi_*Z$, each term of
  $2\langle\varphi_*Z,\nabla'_{\varphi_*X}\varphi_*Y\rangle$ is a
  derivative of an inner product of tangent fields along a tangent field, or
  an inner product with a bracket of tangent fields; by
  \cref{lem:dc-ch6-2-7-pushforward-naturality} and the metric preservation of
  \cref{def:dc-ch6-2-7-isometry}, each equals the corresponding term of the
  Koszul formula downstairs, so
  $\langle\varphi_*Z,\nabla'_{\varphi_*X}\varphi_*Y\rangle(q)
    =\langle Z,\nabla_XY\rangle(p)
    =\langle\varphi_*Z,\varphi_*(\nabla_XY)\rangle(q)$,
  the last step again by metric preservation, since $\nabla_XY$ is tangent.
\end{proof}

\begin{lemma}[isometries preserve the induced curvature]
  \label{lem:dc-ch6-2-7-curvature}
  \lean{Riemannian.DCPatchIsometry.inducedCurvature_opensPushforward, Riemannian.DCPatchIsometry.inducedCurvature_inner_opensPushforward}
  \uses{lem:dc-ch6-2-7-connection, def:dc-ch6-3-curvatures}
  \leanok
  Under the hypotheses of \cref{lem:dc-ch6-2-7-connection}, for tangent fields
  $X,Y,Z$,
  $$
  R'(\varphi_*X,\varphi_*Y)(\varphi_*Z)=\varphi_*\bigl(R(X,Y)Z\bigr),
  $$
  and consequently the sectional numerator is invariant:
  $\langle R'(\varphi_*X,\varphi_*Y)\varphi_*X,\varphi_*Y\rangle(\varphi(p))
    =\langle R(X,Y)X,Y\rangle(p)$.
\end{lemma}
\begin{proof}
  \leanok
  \sketch
  Expand $R(X,Y)Z=\nabla_Y\nabla_XZ-\nabla_X\nabla_YZ+\nabla_{[X,Y]}Z$
  (\cref{def:dc-ch6-3-curvatures}) and apply
  \cref{lem:dc-ch6-2-7-connection} to each covariant derivative — the
  intermediate fields $\nabla_XZ$, $\nabla_YZ$ and the bracket $[X,Y]$ are
  again tangent, the bracket pushes forward by
  \cref{lem:dc-ch6-2-7-pushforward-naturality}, and the pushforward is
  pointwise linear. The numerator identity follows by pairing with
  $\varphi_*Y$ and metric preservation, $R(X,Y)X$ being tangent.
\end{proof}

\begin{lemma}[the pointwise Gauss form of a patch in a flat ambient]
  \label{lem:dc-ch6-2-7-gauss-form}
  \lean{Riemannian.DCImmersedPatch.gaussFormAt, Riemannian.DCImmersedPatch.isAlgCurvatureForm_gaussFormAt, Riemannian.inducedCurvature_inner_eq_gaussFormAt}
  \uses{def:dc-ch6-2-2, prop:dc-ch6-3-1, prop:dc-ch4-2-5, lem:dc-ch6-2-8-euclidean-opens}
  \leanok
  For an immersed patch and $p\in M$, the \emph{Gauss form}
  $$
  A(x,y,z,t)=\langle B(y,t),B(x,z)\rangle-\langle B(x,t),B(y,z)\rangle,
  \qquad x,y,z,t\in T_pM,
  $$
  built from the pointwise second fundamental form of \cref{def:dc-ch6-2-2},
  is an algebraic curvature form on $T_pM$ (\cref{prop:dc-ch4-2-5}). If the
  ambient space is a flat Euclidean open, then for tangent fields
  $$
  \langle R(X,Y)Z,T\rangle(p)=A\bigl(X_p,Y_p,Z_p,T_p\bigr):
  $$
  the curvature pairing depends only on the values of the fields at $p$.
\end{lemma}
\begin{proof}
  \uses{prop:dc-ch6-2-1}
  \leanok
  \sketch
  Multilinearity of $A$ follows from bilinearity of the pointwise $B$ and of
  the inner product. Antisymmetry in $(x,y)$ is immediate from the shape of
  $A$; antisymmetry in $(z,t)$ follows using the symmetry of the inner
  product; the first Bianchi identity follows from the symmetry
  $B(x,y)=B(y,x)$ on tangent vectors (\cref{prop:dc-ch6-2-1} read at a
  point): the six terms of the cyclic sum cancel in pairs. The identity with
  the curvature pairing is the Gauss equation (\cref{prop:dc-ch6-3-1}) with
  flat ambient curvature $\overline{R}\equiv 0$
  (\cref{lem:dc-ch6-2-8-euclidean-opens}), with $B$ evaluated pointwise as in
  \cref{def:dc-ch6-2-2}.
\end{proof}

\begin{remark}[isometry invariance of Gaussian curvature]
  \label{rem:dc-ch6-2-7}
  \lean{Riemannian.DCPatchIsometry.theorema_egregium}
  \uses{rem:dc-ch6-2-6, def:dc-ch6-2-7-isometry, def:dc-ch6-2-4-principal}
  \leanok
  \dcref{ch6:2.7}
  In the case $M=M^2\subset\overline{M}=\mathbb{R}^3$, the product $\lambda_1\lambda_2$
  of the principal curvatures coincides with the Gaussian curvature of the surface.
  \cref{rem:dc-ch6-2-6} then shows that the Gaussian curvature coincides with the sectional
  curvature of the surface, and gives the \emph{Theorema Egregium}:
  the Gaussian curvature of $M^2\subset\mathbb{R}^3$ is an invariant under
  isometries. Precisely: let $D,D'$ be two-dimensional hypersurface patches in
  flat Euclidean ambients, with unit normals $\eta$ at $p$ and $\eta'$ at
  $\varphi(p)$ spanning the normal lines, and let $\varphi:D\to D'$ be an
  isometry (\cref{def:dc-ch6-2-7-isometry}). Then
  $$
  \lambda_1(p)\,\lambda_2(p)=\lambda'_1(\varphi(p))\,\lambda'_2(\varphi(p)).
  $$
\end{remark}
\begin{proof}
  \uses{lem:dc-ch6-2-7-curvature, lem:dc-ch6-2-7-gauss-form,
    lem:dc-ch6-2-7-plane}
  \leanok
  \sketch
  Let $e_1,e_2$ be principal directions at $p$, extended to tangent fields
  $E_1,E_2$, and $e'_1,e'_2$ principal directions at $\varphi(p)$, extended to
  $E'_1,E'_2$. Since the ambients are flat, the Gauss formula for a
  hypersurface (\cref{rem:dc-ch6-2-6}) reads
  $$
  \lambda_1\lambda_2=\langle R(E_1,E_2)E_1,E_2\rangle(p),
  \qquad
  \lambda'_1\lambda'_2=\langle R'(E'_1,E'_2)E'_1,E'_2\rangle(\varphi(p)).
  $$
  By the curvature invariance (\cref{lem:dc-ch6-2-7-curvature}),
  $\langle R(E_1,E_2)E_1,E_2\rangle(p)
    =\langle R'(\varphi_*E_1,\varphi_*E_2)\varphi_*E_1,\varphi_*E_2\rangle(\varphi(p))$.
  By \cref{lem:dc-ch6-2-7-gauss-form} both curvature pairings at
  $\varphi(p)$ are values of the Gauss form $A$ of $D'$: on the pair
  $u_i=d\varphi_p(e_i)$, and on the pair $e'_i$, respectively. Both pairs are
  orthonormal in the two-dimensional plane $T_{\varphi(p)}M'$ — the first
  because $\varphi$ is an isometry — so the two values of $A$ agree by
  \cref{lem:dc-ch6-2-7-plane}, and the chain
  $\lambda_1\lambda_2
    =A(u_1,u_2,u_1,u_2)=A(e'_1,e'_2,e'_1,e'_2)=\lambda'_1\lambda'_2$
  closes.
\end{proof}

\begin{lemma}[the Euclidean ambient: flat connection, vanishing curvature]
  \label{lem:dc-ch6-2-8-euclidean}
  \lean{Riemannian.SmoothVectorField.contDiff, Riemannian.DCLieBracket_eq_fderiv, Riemannian.euclideanCov, Riemannian.euclideanConnection, Riemannian.euclideanConnection_isLeviCivita, Riemannian.euclideanConnection_curvature, Riemannian.euclideanConnection_isConstantCurvature_zero}
  \uses{def:dc-ch1-2-1, thm:dc-ch2-3-6, def:dc-ch4-2-1, def:dc-ch6-3-constant-curvature}
  \leanok
  On Euclidean space — a real inner-product space $F$ with the metric of
  \cref{def:dc-ch1-2-1} given by the ambient inner product — the ordinary
  directional derivative $\nabla_XY=dY(X)$ is an affine connection, it is the
  Levi-Civita connection of the Euclidean metric, and its curvature tensor
  vanishes identically; in particular Euclidean space has constant sectional
  curvature $0$ in the field form of \cref{def:dc-ch6-3-constant-curvature}.
\end{lemma}
\begin{proof}
  \leanok
  \sketch
  On the vector-space model the tangent bundle is canonically trivial, so a
  smooth vector field \emph{is} a smooth map $F\to F$, and the manifold Lie
  bracket is the classical commutator $[X,Y]=dY(X)-dX(Y)$. Additivity and
  $\mathcal{D}$-homogeneity of $\nabla_XY=dY(X)$ in $X$ are linearity of
  $dY_p$; additivity and the Leibniz rule in $Y$ are the sum and product rules
  for the derivative. Symmetry is $dY(X)-dX(Y)=[X,Y]$; metric compatibility is
  the product rule for the inner product. For the curvature,
  $$
  R(X,Y)Z=\nabla_Y\nabla_XZ-\nabla_X\nabla_YZ+\nabla_{[X,Y]}Z ,
  $$
  the second-order terms cancel by equality of mixed derivatives, and the remaining first-order terms cancel exactly against
  $\nabla_{[X,Y]}Z=dZ([X,Y])$. Constant curvature $0$ is then immediate.
\end{proof}

\begin{lemma}[open subsets of Euclidean space as ambient manifolds]
  \label{lem:dc-ch6-2-8-euclidean-opens}
  \lean{Riemannian.opensEuclideanMetric, Riemannian.opensEuclideanConnection, Riemannian.opensEuclideanConnection_isLeviCivita, Riemannian.opensEuclideanConnection_curvature, Riemannian.opensEuclideanConnection_isConstantCurvature_zero, Riemannian.DCLieBracket_opens_eq_fderiv}
  \uses{def:dc-ch1-2-1, thm:dc-ch2-3-6, def:dc-ch4-2-1, def:dc-ch6-3-constant-curvature, lem:dc-ch6-2-8-euclidean}
  \leanok
  Let $A\subseteq F$ be an open subset of a Euclidean space, regarded as an
  open submanifold. A vector field on $A$ is canonically a smooth map
  $A\to F$, and under this identification $A$ carries the restricted Euclidean
  metric $\langle v,w\rangle_q=\langle v,w\rangle$, the Lie bracket is the
  commutator $[X,Y]=dY(X)-dX(Y)$, and the ordinary directional derivative
  $\nabla_XY=dY(X)$ is the Levi-Civita connection of this metric, with
  identically vanishing curvature. Thus $A$ is a flat ambient manifold of
  constant sectional curvature $0$ in the field form of
  \cref{def:dc-ch6-3-constant-curvature}.
\end{lemma}
\begin{proof}
  \leanok
  \sketch
  Smoothness of a field on $A$ and its derivative at $q\in A$ depend only on
  the germ at $q$, and $A$ is open in $F$, so both may be computed after
  extending the field to $F$. The pointwise arguments of
  \cref{lem:dc-ch6-2-8-euclidean} then apply verbatim at each $q\in A$:
  linearity and the Leibniz rule are the sum and product rules of the
  derivative, symmetry of the connection is the commutator identity, metric
  compatibility is the product rule for the inner product, and the curvature
  vanishes by equality of mixed derivatives.
\end{proof}

\begin{example}[curvature of $S^n$]
  \label{ex:dc-ch6-2-8}
  \lean{Riemannian.punctured, Riemannian.sphereTang, Riemannian.spherePatch, Riemannian.spherePatch_secondFundForm_apply, Riemannian.spherePatch_secondFundScalarAt, Riemannian.spherePatch_shapeOperatorAt, Riemannian.spherePatch_inducedCurvature_inner, Riemannian.sphere_sectionalCurvature_one}
  \uses{def:dc-ch6-2-patch, def:dc-ch6-2-B, def:dc-ch6-2-2, thm:dc-ch6-2-5, rem:dc-ch6-2-6, lem:dc-ch6-2-8-euclidean-opens}
  \leanok
  \dcref{ch6:2.8}
  The sphere of radius $r$ in Euclidean space has constant sectional curvature
  $1/r^2$; in particular, the unit sphere has constant sectional curvature $1$.
\end{example}
\begin{proof}
  \sketch
  With the inward unit normal $\eta=-p/r$, the second fundamental form is
  $H_\eta(x,y)=\langle x,y\rangle/r$, so $S_\eta=r^{-1}\mathrm{id}$. Gauss'
  equation against the flat ambient gives sectional curvature $1/r^2$.
\end{proof}

\begin{definition}[geodesic at a point; totally geodesic immersion]
  \label{def:dc-ch6-2-8-geodesic}
  \lean{Riemannian.DCImmersedPatch.IsGeodesicAt, Riemannian.DCImmersedPatch.IsTotallyGeodesic, Riemannian.DCImmersedPatch.isGeodesicAt_iff_secondFundFormAt_eq_zero}
  \source{docarmo:page-0147}
  \uses{def:dc-ch6-2-2}
  \leanok
  An immersion $f:M\to\overline{M}$ is \emph{geodesic at} $p\in M$ if for every
  $\eta\in(T_pM)^\perp$ the second fundamental form $\mathit{II}_\eta$ is
  identically zero at $p$ — equivalently, by polarization, $B$ vanishes at $p$.
  An immersion is \emph{totally geodesic} if it is geodesic at every
  $p\in M$.
\end{definition}

\begin{definition}[geodesic field at a point]
  \label{def:dc-ch6-2-9-geodesic-field}
  \lean{Riemannian.DCImmersedPatch.IsGeodesicFieldAt}
  \uses{def:dc-ch6-2-induced-connection, def:dc-ch6-2-8-geodesic}
  \leanok
  A tangent field $X\in\mathcal{X}(U)$ is a \emph{geodesic field at} $p$ when
  $(\nabla_X X)(p)=0$. The velocity field of a geodesic satisfies this condition
  along the geodesic.
\end{definition}

\begin{lemma}[existence of geodesic fields; first-order scalar data]
  \label{lem:dc-ch6-2-9-existence}
  \lean{Riemannian.exists_contMDiff_dirTangent_eq_one, Riemannian.AffineConnection.cov_apply_eq_zero_of_apply_eq_zero_left, Riemannian.DCImmersedPatch.exists_isGeodesicFieldAt}
  \uses{def:dc-ch6-2-9-geodesic-field, def:dc-ch6-2-induced-connection, lem:dc-ch6-2-proj}
  \leanok
  (a) For every $p\in\overline{M}$ and $v\in T_p\overline{M}$, $v\neq0$, there is a
  globally smooth $f:\overline{M}\to\mathbb{R}$ with $f(p)=0$ and $df_p(v)=1$.
  (b) Through every $x\in T_pM$ there passes a geodesic field at $p$: a tangent
  field $X$ with $X(p)=x$ and $(\nabla_X X)(p)=0$.
\end{lemma}
\begin{proof}
  \leanok
  \sketch
  (a) Choose by Hahn--Banach a continuous linear functional $\ell$ on the model
  space with $\ell(v)=1$, and let $\varphi$ be the coordinate reading of $\ell$
  in the chart at $p$, shifted so that $\varphi(p)=0$. Cutting $\varphi$ off by
  a smooth bump function equal to $1$ near $p$ and supported in the chart source
  yields a global smooth $f$ with $f(p)=0$; since the cut-off is constant near
  $p$ it does not affect the differential, and $df_p(v)=\ell(d(\mathrm{chart})_p
  v)=\ell(v)=1$ because the differential of the chart at its base point is the
  identity.

  (b) If $x=0$ the zero field works, by tensoriality of the covariant derivative
  in its direction slot. Otherwise let $X_0$ be a tangent extension of $x$
  (\cref{lem:dc-ch6-2-proj}), set $a:=(\nabla_{X_0}X_0)(p)\in T_pM$, let $W$ be a
  tangent extension of $-a$, and take $f$ as in (a) with $df_p(x)=1$, $f(p)=0$.
  Put $X:=X_0+fW$; then $X$ is tangent and $X(p)=x$. Expanding
  $\nabla_{X}X$ by additivity in both slots, $\mathcal{D}$-homogeneity in the
  direction slot and the Leibniz rule in the derived slot,
  $$
  \nabla_X X=\nabla_{X_0}X_0+f\,\nabla_{X_0}W+X_0(f)\,W+f\,\nabla_{W}X ,
  $$
  and at $p$ all terms carrying the factor $f(p)=0$ vanish, leaving
  $(\nabla_XX)(p)=a+X_0(f)(p)\,W(p)=a+1\cdot(-a)=0$.
\end{proof}

\begin{proposition}[totally geodesic criterion]
  \label{prop:dc-ch6-2-9}
  \source{docarmo:page-0147}
  \uses{def:dc-ch6-2-8-geodesic, def:dc-ch6-2-9-geodesic-field, lem:dc-ch6-2-9-existence, def:dc-ch6-2-B, def:dc-ch6-2-2, lem:dc-ch6-2-locality, prop:dc-ch6-2-1, lem:dc-ch6-2-9-field}
  \notready
  \dcref{ch6:2.9}
  An immersion $f:M\to\overline{M}$ is geodesic at $p\in M$ if and only if every
  geodesic of $M$ at $p$ is a geodesic of $\overline{M}$ at $p$; precisely: if
  and only if every geodesic field $X$ at $p$
  (\cref{def:dc-ch6-2-9-geodesic-field}) satisfies
  $(\overline{\nabla}_X X)(p)=0$. The two readings agree through the displayed
  identity of the proof: for any tangent field $X$ with $X(p)=x$ and any
  $\eta\in(T_pM)^\perp$,
  $$
  H_\eta(x,x)=\langle\overline{\nabla}_X X,\eta\rangle(p).
  $$
\end{proposition}
\begin{proof}
  \sketch
  For any tangent field $X$ with $X(p)=x$ and $\eta\in(T_pM)^\perp$, the Gauss
  decomposition $\overline{\nabla}_X X=\nabla_X X+B(X,X)$ paired against $\eta$
  gives
  $$
  H_\eta(x,x)=\langle B(x,x),\eta\rangle=\langle\overline{\nabla}_X X,\eta\rangle(p),
  $$
  since the tangential part $\nabla_X X$ pairs to zero against $\eta$.

  ($\Rightarrow$) If $f$ is geodesic at $p$ then $B$ vanishes at $p$, so for a
  geodesic field $X$ at $p$ the ambient acceleration
  $(\overline{\nabla}_X X)(p)=(\nabla_X X)(p)+B(x,x)=0+0=0$.

  ($\Leftarrow$) Let $\eta\in(T_pM)^\perp$ and $x\in T_pM$. By
  \cref{lem:dc-ch6-2-9-existence} there is a geodesic field $X$ at $p$ with
  $X(p)=x$; by hypothesis $(\overline{\nabla}_X X)(p)=0$, hence
  $H_\eta(x,x)=\langle\overline{\nabla}_X X,\eta\rangle(p)=0$. Thus every
  $\mathit{II}_\eta$ vanishes at $p$, i.e.\ $f$ is geodesic at $p$.
\end{proof}

\begin{remark}[geometric interpretation of sectional curvature]
  \label{rem:dc-ch6-2-9-interpretation}
  \lean{Riemannian.DCImmersedPatch.secondFundForm_apply_eq_zero_of_isGeodesicAt, Riemannian.DCImmersedPatch.inducedCurvature_inner_eq_curvature_inner_of_isGeodesicAt, Riemannian.DCImmersedPatch.sectionalCurvature_eq_of_isGeodesicAt, Riemannian.DCImmersedPatch.inducedCurvatureAt_inner_eq_of_isGeodesicAt}
  \source{docarmo:page-0147,docarmo:page-0148}
  \uses{prop:dc-ch6-2-9, thm:dc-ch6-2-5, def:dc-ch6-2-8-geodesic, def:dc-ch4-3-2}
  \leanok
  If $B\subset T_pM$ is a normal neighborhood and $\sigma\subset T_pM$ is
  two-dimensional, then $S=\exp_p(\sigma\cap B)$ is geodesic at $p$ and
  $$
  K_S(p)=K(p,\sigma).
  $$
  More generally, any submanifold geodesic at $p$ has the same sectional
  curvature as the ambient manifold on each tangent plane at $p$.
\end{remark}
\begin{proof}
  \sketch
  The exponential construction makes the second fundamental form of $S$ vanish
  at $p$. The equality then follows from Gauss' equation; the same argument
  applies to any submanifold geodesic at $p$.
\end{proof}

\begin{remark}[geodesic surfaces and constant sectional curvature]
  \label{rem:dc-ch6-2-9-exp-surface}
  \source{docarmo:page-0148}
  \uses{prop:dc-ch6-2-9}
  \notready
  \notready
  When $\exp_p$ is a smooth local diffeomorphism and
  $\exp_p|_{\sigma\cap B}$ is an immersion, the surface
  $S=\exp_p(\sigma\cap B)$ is geodesic at $p$. Linear subspaces and translates
  in Euclidean space, and intersections of linear subspaces with $S^n$, are
  totally geodesic. Conversely, if every two-plane in every tangent space is
  tangent to a totally geodesic submanifold, then sectional curvature is
  constant.
\end{remark}

\begin{definition}[minimal immersion; mean curvature vector]
  \label{def:dc-ch6-2-10}
  \lean{Riemannian.DCImmersedPatch.IsMinimal, Riemannian.DCImmersedPatch.meanCurvatureVector, Riemannian.DCImmersedPatch.meanCurvatureVector_mem, Riemannian.DCImmersedPatch.inner_meanCurvatureVector, Riemannian.DCImmersedPatch.meanCurvatureVector_unique, Riemannian.DCImmersedPatch.existsUnique_meanCurvatureVector, Riemannian.DCImmersedPatch.isMinimal_iff_meanCurvatureVector_eq_zero, Riemannian.DCImmersedPatch.shapeOperatorAt_add_normal, Riemannian.DCImmersedPatch.shapeOperatorAt_smul_normal}
  \source{docarmo:page-0148}
  \uses{def:dc-ch6-2-4-principal, def:dc-ch6-2-2}
  \leanok
  \dcref{ch6:2.10}
  An immersion $f:M\to\overline{M}$ is called \emph{minimal} if for every $p\in M$ and
  every $\eta\in(T_pM)^\perp$ the trace of $S_\eta$ is zero. Choosing an orthonormal frame
  $E_1,\dots,E_m$ of vectors in $\mathcal{X}(U)^\perp$, where $U$ is a neighborhood
  of $p$ in which $f$ is an embedding, we write at $p$
  $$
  B(x,y)=\sum_i H_i(x,y)E_i,\quad x,y\in T_pM,\ i=1,\dots,m,
  $$
  with $H_i=H_{E_i}$. The normal vector
  $$
  H=\frac{1}{n}\sum_i(\operatorname{trace}S_i)E_i,
  $$
  where $S_i=S_{E_i}$, does not depend on the chosen frame; it is the \emph{mean
  curvature vector} of $f$. Thus $f$ is minimal iff $H(p)=0$ for all $p\in M$.
\end{definition}

\section{The fundamental equations}

Use Latin letters $X,Y,Z$ for tangent fields and Greek letters $\xi,\eta,\zeta$
for normal fields. The tangential and normal components of
$\overline{\nabla}_X\eta$ are respectively $-S_\eta(X)$ and $\nabla_X^\perp\eta$:

$$
\nabla_X^\perp\eta=(\overline{\nabla}_X\eta)^N=\overline{\nabla}_X\eta-(\overline{\nabla}_X\eta)^T=\overline{\nabla}_X\eta+S_\eta(X).\qquad(5)
$$

\begin{definition}[shape operator and normal connection: the Weingarten decomposition]
  \label{def:dc-ch6-3-normal-connection}
  \lean{Riemannian.DCImmersedPatch.shapeOperator, Riemannian.DCImmersedPatch.normalCov, Riemannian.DCImmersedPatch.cov_eq_neg_shapeOperator_add_normalCov, Riemannian.DCImmersedPatch.shapeOperator_mem, Riemannian.DCImmersedPatch.normalCov_mem}
  \uses{def:dc-ch6-2-patch, def:dc-ch6-2-B}
  \leanok
  For fields $X$, $\eta$ on the patch set
  $S_\eta(X)=-(\overline{\nabla}_X\eta)^T$ (the \emph{shape operator}, tangent
  valued) and $\nabla^\perp_X\eta=(\overline{\nabla}_X\eta)^N$ (the \emph{normal
  connection}, normal valued), so that equation (5) reads

  $$
  \overline{\nabla}_X\eta=-S_\eta(X)+\nabla^\perp_X\eta.
  $$
\end{definition}

\begin{proposition}[connection laws of the normal connection]
  \label{prop:dc-ch6-3-normal-connection-laws}
  \lean{Riemannian.DCImmersedPatch.normalCov_add_left, Riemannian.DCImmersedPatch.normalCov_smul_left, Riemannian.DCImmersedPatch.normalCov_add_right, Riemannian.DCImmersedPatch.normalCov_smul_right}
  \uses{def:dc-ch6-3-normal-connection, lem:dc-ch6-2-proj}
  \leanok
  The normal connection has the usual properties of a connection: it is
  additive in both arguments, $\mathcal{D}$-homogeneous in the direction
  ($\nabla^\perp_{fX}\eta=f\nabla^\perp_X\eta$), and satisfies the Leibniz rule
  $\nabla^\perp_X(f\eta)=f\nabla^\perp_X\eta+X(f)\,\eta$ for $\eta$ normal.
\end{proposition}
\begin{proof}
  \leanok
  \sketch
  Immediate from the corresponding laws of $\overline{\nabla}$ and the
  additivity and $\mathcal{D}$-homogeneity of the normal projection
  (\cref{lem:dc-ch6-2-proj}); in the Leibniz rule the first-order term
  $X(f)\,\eta$ is already normal because $\eta$ is.
\end{proof}

\begin{lemma}[Weingarten pairing and self-adjointness]
  \label{lem:dc-ch6-3-weingarten}
  \lean{Riemannian.DCImmersedPatch.inner_shapeOperator_apply, Riemannian.DCImmersedPatch.inner_shapeOperator_symm}
  \uses{def:dc-ch6-3-normal-connection, prop:dc-ch6-2-1, def:dc-ch6-2-B}
  \leanok
  For $Y$ tangent and $\eta$ normal,

  $$
  \langle S_\eta(X),Y\rangle=\langle B(X,Y),\eta\rangle
  $$

  at every point. Consequently, for tangent $X,Y$,
  $\langle S_\eta(X),Y\rangle=\langle X,S_\eta(Y)\rangle$: the shape operator
  is self-adjoint with respect to the metric.
\end{lemma}
\begin{proof}
  \leanok
  \sketch
  Since $\langle\eta,Y\rangle=0$ identically, compatibility of
  $\overline{\nabla}$ with the metric gives
  $0=X\langle\eta,Y\rangle
  =\langle\overline{\nabla}_X\eta,Y\rangle+\langle\eta,\overline{\nabla}_XY\rangle$.
  In the first term only the tangential part of $\overline{\nabla}_X\eta$
  survives against the tangent field $Y$, giving $-\langle S_\eta(X),Y\rangle$;
  in the second only the normal part of $\overline{\nabla}_XY$ survives against
  the normal field $\eta$, giving $\langle B(X,Y),\eta\rangle$. The
  self-adjointness follows from the symmetry of $B$
  (\cref{prop:dc-ch6-2-1}):
  $\langle S_\eta(X),Y\rangle=\langle B(X,Y),\eta\rangle
  =\langle B(Y,X),\eta\rangle=\langle S_\eta(Y),X\rangle$.
\end{proof}

\begin{definition}[curvatures of the induced and normal connections]
  \label{def:dc-ch6-3-curvatures}
  \lean{Riemannian.DCImmersedPatch.inducedCurvature, Riemannian.DCImmersedPatch.inducedCurvature_mem, Riemannian.DCImmersedPatch.normalCurvature, Riemannian.DCImmersedPatch.normalCurvature_mem}
  \uses{def:dc-ch6-2-induced-connection, def:dc-ch6-3-normal-connection, def:dc-ch4-2-1}
  \leanok
  The \emph{curvature of the induced connection} is
  $R(X,Y)Z=\nabla_Y\nabla_XZ-\nabla_X\nabla_YZ+\nabla_{[X,Y]}Z$ (tangent
  valued), and the \emph{normal curvature} is
  $R^\perp(X,Y)\eta=\nabla^\perp_Y\nabla^\perp_X\eta
  -\nabla^\perp_X\nabla^\perp_Y\eta+\nabla^\perp_{[X,Y]}\eta$ (normal valued);
  on tangent fields $R$ is the curvature of the immersed manifold.
\end{definition}

\begin{lemma}[decomposition of the ambient curvature]
  \label{lem:dc-ch6-3-decomposition}
  \lean{Riemannian.DCImmersedPatch.curvature_decomposition, Riemannian.DCImmersedPatch.curvature_normal_decomposition}
  \uses{def:dc-ch6-3-curvatures, lem:dc-ch6-2-proj, def:dc-ch6-2-B, def:dc-ch6-3-normal-connection}
  \leanok
  For any fields $X,Y,Z,\eta$ on the patch,

  $$
  \overline{R}(X,Y)Z=R(X,Y)Z+B(Y,\nabla_X Z)-B(X,\nabla_Y Z)+B([X,Y],Z)-S_{B(X,Z)}Y+S_{B(Y,Z)}X+\nabla_Y^\perp B(X,Z)-\nabla_X^\perp B(Y,Z)\qquad(6)
  $$

  and

  $$
  \overline{R}(X,Y)\eta=R^\perp(X,Y)\eta-S_{\nabla_X^\perp\eta}Y+S_{\nabla_Y^\perp\eta}X-\nabla_Y(S_\eta X)+\nabla_X(S_\eta Y)-B(Y,S_\eta X)+B(X,S_\eta Y)-S_\eta([X,Y]).
  $$
\end{lemma}
\begin{proof}
  \leanok
  \sketch
  Substitute the Gauss decomposition
  $\overline{\nabla}_XW=\nabla_XW+B(X,W)$ (\cref{def:dc-ch6-2-B}) and the
  Weingarten decomposition
  $\overline{\nabla}_X\eta=-S_\eta(X)+\nabla^\perp_X\eta$
  (\cref{def:dc-ch6-3-normal-connection}) into
  $\overline{R}(X,Y)W=\overline{\nabla}_Y\overline{\nabla}_XW
  -\overline{\nabla}_X\overline{\nabla}_YW+\overline{\nabla}_{[X,Y]}W$ and
  collect terms: for the first identity split
  $\overline{\nabla}_XZ=\nabla_XZ+B(X,Z)$, then split
  $\overline{\nabla}_Y(\nabla_XZ)=\nabla_Y\nabla_XZ+B(Y,\nabla_XZ)$ and
  $\overline{\nabla}_Y(B(X,Z))=-S_{B(X,Z)}Y+\nabla^\perp_YB(X,Z)$, and likewise
  for the terms with $X$ and $Y$ interchanged and for
  $\overline{\nabla}_{[X,Y]}Z$; for the second identity split
  $\overline{\nabla}_X\eta=-S_\eta(X)+\nabla^\perp_X\eta$, then
  $\overline{\nabla}_Y(S_\eta X)=\nabla_Y(S_\eta X)+B(Y,S_\eta X)$ and
  $\overline{\nabla}_Y(\nabla^\perp_X\eta)
  =-S_{\nabla^\perp_X\eta}Y+\nabla^\perp_Y\nabla^\perp_X\eta$, and likewise for
  the remaining terms.
\end{proof}

\begin{proposition}[Gauss equation and Ricci equation]
  \label{prop:dc-ch6-3-1}
  \lean{Riemannian.DCImmersedPatch.gauss_equation, Riemannian.DCImmersedPatch.ricci_equation}
  \uses{def:dc-ch6-3-curvatures, lem:dc-ch6-3-decomposition, lem:dc-ch6-3-weingarten}
  \leanok
  \dcref{ch6:3.1}
  The equations are:
  (a) \emph{Gauss equation}
  $$
  \langle\overline{R}(X,Y)Z,T\rangle=\langle R(X,Y)Z,T\rangle-\langle B(Y,T),B(X,Z)\rangle+\langle B(X,T),B(Y,Z)\rangle.
  $$
  (b) \emph{Ricci equation}
  $$
  \langle\overline{R}(X,Y)\eta,\zeta\rangle-\langle R^\perp(X,Y)\eta,\zeta\rangle=\langle[S_\eta,S_\zeta]X,Y\rangle,
  $$
  where $[S_\eta,S_\zeta]$ denotes the operator $S_\eta\circ S_\zeta-S_\zeta\circ S_\eta$.
\end{proposition}
\begin{proof}
  \leanok
  \sketch
  Substitute the Gauss and Weingarten decompositions into the ambient curvature
  operator and collect tangent and normal parts, as in
  \cref{lem:dc-ch6-3-decomposition}. Pairing the first resulting identity with
  $T$ removes normal terms; the Weingarten pairing converts the remaining shape
  terms to $B$-pairings and yields Gauss' equation. Pairing the second identity
  with $\zeta$ removes tangent terms and gives
  $\langle R^\perp(X,Y)\eta,\zeta\rangle$ plus the commutator term
  $\langle[S_\eta,S_\zeta]X,Y\rangle$, which is Ricci's equation.
\end{proof}

\begin{remark}[Gauss equation as a sectional-curvature identity]
  \label{rem:dc-ch6-3-2}
  \lean{Riemannian.DCImmersedPatch.inducedCurvature_inner_sub_curvature_inner}
  \uses{thm:dc-ch6-2-5, prop:dc-ch6-3-1}
  \leanok
  \dcref{ch6:3.2}
  \cref{thm:dc-ch6-2-5} is a special case of Gauss' equation:
  taking $Z=X$ and $T=Y$ in the Gauss equation and using the symmetry of $B$
  yields $\langle R(X,Y)X,Y\rangle-\langle\overline{R}(X,Y)X,Y\rangle
  =\langle B(X,X),B(Y,Y)\rangle-|B(X,Y)|^2$.
\end{remark}

\begin{definition}[constant sectional curvature, field form]
  \label{def:dc-ch6-3-constant-curvature}
  \lean{Riemannian.AffineConnection.IsConstantCurvature}
  \uses{def:dc-ch4-2-1, lem:dc-ch4-3-4}
  \leanok
  A Riemannian manifold with Riemannian connection $\overline{\nabla}$ has
  \emph{constant sectional curvature} $K_0$ when

  $$
  \langle\overline{R}(X,Y)Z,W\rangle=K_0\big(\langle X,Z\rangle\langle Y,W\rangle-\langle Y,Z\rangle\langle X,W\rangle\big)
  $$

  for all fields $X,Y,Z,W$ — the characterization of constant sectional
  curvature of \cref{lem:dc-ch4-3-4}.
\end{definition}

\begin{remark}[flat normal bundle]
  \label{rem:dc-ch6-3-3}
  \lean{Riemannian.DCImmersedPatch.ricci_equation_of_isConstantCurvature, Riemannian.DCImmersedPatch.HasFlatNormalBundle, Riemannian.DCImmersedPatch.ShapeOperatorsCommuteAt, Riemannian.DCImmersedPatch.hasFlatNormalBundle_iff_shapeOperatorsCommute, Riemannian.DCImmersedPatch.eq_of_inner_eq_of_mem_normalSpace}
  \uses{prop:dc-ch6-3-1, def:dc-ch6-3-covderiv-B, def:dc-ch6-3-constant-curvature, def:dc-ch6-2-2}
  \leanok
  \dcref{ch6:3.3}
  We say the normal bundle of an immersion is \emph{flat} if $R^\perp=0$. Assume the
  ambient space $\overline{M}$ has constant sectional curvature. Then the Ricci
  equation can be written
  $$
  \langle R^\perp(X,Y)\eta,\zeta\rangle=-\langle[S_\eta,S_\zeta]X,Y\rangle.
  $$
  Thus $R^\perp=0$ if and only if $[S_\eta,S_\zeta]=0$ for all
  $\eta,\zeta$, that is, if and only if for all $p\in M$ there exists a basis of
  $T_pM$ which diagonalizes all of the $S_\eta$ simultaneously.
\end{remark}

\begin{definition}[the covariant derivative of the second fundamental form]
  \label{def:dc-ch6-3-covderiv-B}
  \lean{Riemannian.DCImmersedPatch.secondFundTensor, Riemannian.DCImmersedPatch.secondFundTensorCovDeriv}
  \uses{def:dc-ch6-2-B, def:dc-ch6-2-induced-connection, def:dc-ch6-3-normal-connection}
  \leanok
  The second fundamental form regarded as a tensor
  $B:\mathcal{X}(M)\times\mathcal{X}(M)\times\mathcal{X}(M)^\perp\to\mathbb{R}$
  is $B(X,Y,\eta)=\langle B(X,Y),\eta\rangle$, with covariant derivative

  $$
  (\overline{\nabla}_X B)(Y,Z,\eta)=X(B(Y,Z,\eta))-B(\nabla_X Y,Z,\eta)-B(Y,\nabla_X Z,\eta)-B(Y,Z,\nabla_X^\perp\eta).
  $$
\end{definition}

\begin{lemma}[reduction of the covariant derivative of $B$]
  \label{lem:dc-ch6-3-covderiv-reduce}
  \lean{Riemannian.DCImmersedPatch.secondFundTensorCovDeriv_eq}
  \uses{def:dc-ch6-3-covderiv-B, lem:dc-ch6-3-weingarten}
  \leanok
  For $\eta$ normal,

  $$
  (\overline{\nabla}_X B)(Y,Z,\eta)=\langle\nabla_X^\perp(B(Y,Z)),\eta\rangle-\langle B(\nabla_X Y,Z),\eta\rangle-\langle B(Y,\nabla_X Z),\eta\rangle.
  $$
\end{lemma}
\begin{proof}
  \leanok
  \sketch
  By compatibility of $\overline{\nabla}$ with the metric,
  $X\langle B(Y,Z),\eta\rangle
  =\langle\overline{\nabla}_X(B(Y,Z)),\eta\rangle
  +\langle B(Y,Z),\overline{\nabla}_X\eta\rangle$. In the first term only the
  normal part $\nabla^\perp_X(B(Y,Z))$ of $\overline{\nabla}_X(B(Y,Z))$
  survives against the normal field $\eta$; in the second only the normal part
  $\nabla^\perp_X\eta$ of $\overline{\nabla}_X\eta$ survives against the normal
  field $B(Y,Z)$. Substituting into the definition of
  $(\overline{\nabla}_X B)(Y,Z,\eta)$ cancels the term
  $\langle B(Y,Z),\nabla^\perp_X\eta\rangle$ and leaves the stated formula.
\end{proof}

\begin{proposition}[Codazzi's equation]
  \label{prop:dc-ch6-3-4}
  \lean{Riemannian.DCImmersedPatch.codazzi_equation}
  \uses{def:dc-ch6-3-covderiv-B, lem:dc-ch6-3-covderiv-reduce, lem:dc-ch6-3-decomposition, prop:dc-ch6-2-induced-connection}
  \leanok
  \dcref{ch6:3.4}
  With the notation above,

  $$
  \langle\overline{R}(X,Y)Z,\eta\rangle=(\overline{\nabla}_Y B)(X,Z,\eta)-(\overline{\nabla}_X B)(Y,Z,\eta).
  $$
\end{proposition}
\begin{proof}
  \leanok
  \sketch
  Observe first that, by \cref{lem:dc-ch6-3-covderiv-reduce},

  $$
  (\overline{\nabla}_X B)(Y,Z,\eta)=X\langle B(Y,Z),\eta\rangle-\langle B(\nabla_X Y,Z),\eta\rangle-\langle B(Y,\nabla_X Z),\eta\rangle-\langle B(Y,Z),\nabla_X^\perp\eta\rangle=\langle\nabla_X^\perp(B(Y,Z)),\eta\rangle-\langle B(\nabla_X Y,Z),\eta\rangle-\langle B(Y,\nabla_X Z),\eta\rangle.
  $$

  Consider expression (6) of \cref{lem:dc-ch6-3-decomposition} and multiply both sides by
  $\eta$. Taking the remark above into account,
  
  $$
  \langle\overline{R}(X,Y)Z,\eta\rangle=\langle B(Y,\nabla_X Z),\eta\rangle+\langle\nabla_Y^\perp B(X,Z),\eta\rangle-\langle B(X,\nabla_Y Z),\eta\rangle-\langle\nabla_X^\perp B(Y,Z),\eta\rangle+\langle B(\nabla_X Y,Z),\eta\rangle-\langle B(\nabla_Y X,Z),\eta\rangle=-(\overline{\nabla}_X B)(Y,Z,\eta)+(\overline{\nabla}_Y B)(X,Z,\eta),
  $$
  
  which is Codazzi's equation.
\end{proof}

\begin{remark}[Codazzi for constant curvature ambient and hypersurfaces]
  \label{rem:dc-ch6-3-5}
  \lean{Riemannian.DCImmersedPatch.codazzi_symm_of_isConstantCurvature, Riemannian.DCImmersedPatch.normalCov_eq_zero_of_unit_spanning, Riemannian.DCImmersedPatch.codazzi_hypersurface}
  \uses{prop:dc-ch6-3-4, def:dc-ch6-3-constant-curvature, def:dc-ch6-3-normal-connection}
  \leanok
  \dcref{ch6:3.5}
  If the ambient space $\overline{M}$ has constant sectional curvature, the Codazzi
  equation reduces to
  $$
  (\overline{\nabla}_X B)(Y,Z,\eta)=(\overline{\nabla}_Y B)(X,Z,\eta).
  $$
  If, in addition, the codimension of the immersion is $1$ and $\eta$ is chosen as a
  local unit normal field, then
  $\nabla_X^\perp\eta=0$, hence
  $$
  (\overline{\nabla}_X B)(Y,Z,\eta)=X\langle S_\eta(Y),Z\rangle-\langle S_\eta(\nabla_X Y),Z\rangle-\langle S_\eta(Y),\nabla_X Z\rangle=\langle\nabla_X(S_\eta(Y)),Z\rangle-\langle S_\eta(\nabla_X Y),Z\rangle.
  $$
  Therefore, in this case, the Codazzi equation can be written
  $$
  \nabla_X(S_\eta(Y))-\nabla_Y(S_\eta(X))=S_\eta([X,Y]).
  $$
\end{remark}

\begin{lemma}[pointwise geodesic-field criterion]
  \label{lem:dc-ch6-2-9-field}
  \source{docarmo:page-0147}
  \dcref{ch6:2.9}
  \uses{def:dc-ch6-2-8-geodesic, def:dc-ch6-2-9-geodesic-field,
    lem:dc-ch6-2-9-existence, def:dc-ch6-2-B, def:dc-ch6-2-2,
    lem:dc-ch6-2-locality, prop:dc-ch6-2-1}
  \lean{Riemannian.DCImmersedPatch.isGeodesicAt_iff_forall_isGeodesicFieldAt,
    Riemannian.DCImmersedPatch.secondFundQuadAt_eq_metricInner_cov,
    Riemannian.DCImmersedPatch.IsTotallyGeodesic.cov_apply_eq_zero}
  \leanok
  An immersion $f:M\to\overline{M}$ is geodesic at $p\in M$ if and only if every
  geodesic field $X$ at $p$
  (\cref{def:dc-ch6-2-9-geodesic-field}) satisfies
  $(\overline{\nabla}_X X)(p)=0$. The two readings agree through the displayed
  identity of the proof: for any tangent field $X$ with $X(p)=x$ and any
  $\eta\in(T_pM)^\perp$,
  $$
  H_\eta(x,x)=\langle\overline{\nabla}_X X,\eta\rangle(p).
  $$
\end{lemma}
\begin{proof}
  \leanok
  For any tangent field $X$ with $X(p)=x$ and $\eta\in(T_pM)^\perp$, the Gauss
  decomposition $\overline{\nabla}_X X=\nabla_X X+B(X,X)$ paired against $\eta$
  gives (do Carmo computes the same identity through the Weingarten pairing
  $\langle N,\overline{\nabla}_X X\rangle=-\langle\overline{\nabla}_X N,X\rangle$
  for a normal extension $N$ of $\eta$)
  $$
  H_\eta(x,x)=\langle B(x,x),\eta\rangle=\langle\overline{\nabla}_X X,\eta\rangle(p),
  $$
  since the tangential part $\nabla_X X$ pairs to zero against $\eta$.

  ($\Rightarrow$) If $f$ is geodesic at $p$ then $B$ vanishes at $p$, so for a
  geodesic field $X$ at $p$ the ambient acceleration
  $(\overline{\nabla}_X X)(p)=(\nabla_X X)(p)+B(x,x)=0+0=0$.

  ($\Leftarrow$) Let $\eta\in(T_pM)^\perp$ and $x\in T_pM$. By
  \cref{lem:dc-ch6-2-9-existence} there is a geodesic field $X$ at $p$ with
  $X(p)=x$; by hypothesis $(\overline{\nabla}_X X)(p)=0$, hence
  $H_\eta(x,x)=\langle\overline{\nabla}_X X,\eta\rangle(p)=0$. Thus every
  $\mathit{II}_\eta$ vanishes at $p$, i.e.\ $f$ is geodesic at $p$. $\square$
\end{proof}

\chapter{Complete Manifolds; Hopf-Rinow and Hadamard Theorems}
\label{chap:dc-complete-manifolds-hopf-rinow-hadamard}

\section{Complete manifolds; Hopf-Rinow Theorem}

All manifolds in this chapter are assumed connected.

\begin{definition}[extendible]

  \label{def:dc-ch7-2-1}
  \lean{Riemannian.IsExtendible}
  \leanok
  \dcref{ch7:2.1}
A Riemannian manifold $M$ is \emph{extendible} if there exists a Riemannian manifold
  $M'$ such that $M$ is isometric to a proper open subset of $M'$. In the opposite
  case, $M$ is \emph{non-extendible}.
\end{definition}

\begin{definition}[geodesically complete]

  \label{def:dc-ch7-2-2}
  \lean{Riemannian.Exponential.expDomainIntrinsic, Riemannian.Geodesic.expDomainIntrinsic_eq_univ_iff_isGeodesicallyCompleteAt, Riemannian.Geodesic.IsGeodesicallyComplete}
  \uses{def:dc-ch3-expmap}
  \leanok
  \dcref{ch7:2.2}
A Riemannian manifold $M$ is \emph{(geodesically) complete} if for all $p\in M$ the
  exponential map $\exp_p$ is defined for all $v\in T_pM$; i.e. if any geodesic
  $\gamma(t)$ starting from $p$ is defined for all values of the parameter
  $t\in\mathbb{R}$.
\end{definition}

\begin{lemma}[a local isometry maps geodesics to geodesics]

  \label{lem:dc-ch7-2-3-geodesic-push}
  \lean{Riemannian.RiemannianMetric.eq_of_metricInner_eq, Riemannian.Geodesic.solvesGeodesicODEAt_comp_of, Riemannian.Geodesic.isGeodesic_comp_of_isLocalDiffeomorph}
  \uses{def:dc-ch1-2-2, lem:dc-ch7-3-4-rays-are-geodesics}
  \leanok
  \dcref{ch7:2.3}
Let $f:M\to M'$ be a smooth local diffeomorphism that preserves the metric
  ($|df\,v|=|v|$ for all $v$; equivalently $f^*g'=g$). Then $f$ carries geodesics of $g$ to
  geodesics of $g'$: for every continuous $g$-geodesic $\gamma$, the image $f\circ\gamma$ is a
  $g'$-geodesic.
\end{lemma}
\begin{proof}

  \leanok
  \uses{lem:dc-ch7-3-4-rays-are-geodesics}
\sketch Metric preservation identifies $g=f^*g'$. The Christoffel transformation law for a local diffeomorphism sends the geodesic operator for $f^*g'$ to the corresponding operator for $g'$, so $f\circ\gamma$ satisfies the geodesic equation whenever $\gamma$ does.
\end{proof}

\begin{proposition}[Completeness prevents extension]

  \label{prop:dc-ch7-2-3}
  \lean{Riemannian.not_isExtendible_of_isGeodesicallyComplete}
  \uses{def:dc-ch7-2-1, def:dc-ch7-2-2, lem:dc-ch7-2-3-geodesic-push, thm:dc-ch3-3-7}
  \leanok
  \dcref{ch7:2.3}
If $M$ is complete then $M$ is non-extendible.
\end{proposition}
\begin{proof}

  \leanok
  \uses{def:dc-ch7-2-1, def:dc-ch7-2-2, lem:dc-ch7-2-3-geodesic-push, thm:dc-ch3-3-7}
\sketch If $M$ were isometric to a proper open subset of connected $M'$, choose a boundary point and a normal neighborhood there. A geodesic from the boundary point to an interior point restricts, after reversing parameter, to a geodesic in $M$ that cannot be extended to the boundary time, contradicting completeness.
\end{proof}

\begin{definition}[distance]

  \label{def:dc-ch7-2-4}
  \lean{Riemannian.piecewiseRiemannianEDist, Riemannian.piecewiseRiemannianEDist_eq_riemannianEDist, Riemannian.riemannianEDist_eq_iInf_isPiecewiseSmoothCurve, Riemannian.RiemannianMetric.IsRiemannianDist}
  \uses{def:dc-ch1-2-9, def:dc-ch3-3-1}
  \leanok
  \dcref{ch7:2.4}
The \emph{distance} $d(p,q)$ is defined by $d(p,q)=$ infimum of the lengths of all
  curves $f_{p,q}$, where $f_{p,q}$ is a piecewise differentiable curve joining $p$
  to $q$. Equivalently, one may take the infimum over curves that are smooth on
  the subintervals of a finite partition.
\end{definition}
\begin{proof}
\sketch Additivity of length and the triangle inequality give one inequality.
  Conversely, subdivide a $C^1$ path into small chart windows and replace each
  piece by the pullback of a straight segment. Continuity of the chart Gram
  forms makes the polygonal lengths converge to the original length.
\end{proof}

\begin{lemma}[finiteness of the distance]

  \label{lem:dc-ch7-2-4-finite}
  \lean{Manifold.riemannianEDist_ne_top}
  \uses{def:dc-ch7-2-4}
  \leanok
  \dcref{ch7:2.4}
If $M$ is connected, then $d(p,q)<\infty$ for all $p,q\in M$.
\end{lemma}
\begin{proof}

  \leanok
\sketch In a chart, nearby points are joined to a fixed point by finite-length coordinate segments, so the set of points at finite distance from $p$ is open. The triangle inequality makes its complement open as well; connectedness implies it is all of $M$.
\end{proof}

\begin{lemma}[short curves stay near their origin]

  \label{lem:dc-ch7-2-5-chart}
  \lean{setOf_riemannianEDist_lt_subset_nhds}
  \uses{def:dc-ch7-2-4}
  \dcref{ch7:2.5}
  \mathlibok
Let $p\in M$ and let $S$ be a neighborhood of $p$. Then there exists $c>0$
  such that every $q\in M$ with $d(p,q)<c$ lies in $S$.
\end{lemma}

\begin{lemma}[positivity of the distance]

  \label{lem:dc-ch7-2-5-pos}
  \lean{Manifold.riemannianEDist_pos}
  \uses{def:dc-ch7-2-4}
  \leanok
  \dcref{ch7:2.5}
If $p\ne q$ then $d(p,q)>0$.
\end{lemma}
\begin{proof}

  \leanok
  \uses{lem:dc-ch7-2-5-chart}
\sketch Apply the short-curve neighborhood lemma to $S=M\smallsetminus\{q\}$, obtaining $c>0$ such that $d(p,x)<c$ implies $x\in S$. Since $q\notin S$, $d(p,q)\ge c>0$.
\end{proof}

\begin{proposition}[$d$ is a metric]

  \label{prop:dc-ch7-2-5}
  \lean{Manifold.riemannianEDist_ne_top, Manifold.riemannianEDist_pos, _root_.MetricSpace.ofRiemannianMetric}
  \uses{def:dc-ch7-2-4}
  \leanok
  \dcref{ch7:2.5}
With the distance $d$, $M$ is a metric space, that is:
  1. $d(p,r)\le d(p,q)+d(q,r)$,
  2. $d(p,q)=d(q,p)$,
  3. $d(p,q)\ge 0$, and $d(p,q)=0\iff p=q$.
\end{proposition}
\begin{proof}

  \leanok
  \uses{lem:dc-ch7-2-4-finite, lem:dc-ch7-2-5-pos}
\sketch Reversal and concatenation of curves give symmetry and the triangle inequality, while constant curves give $d(p,p)=0$ and finiteness follows from \cref{lem:dc-ch7-2-4-finite}. The separation axiom is \cref{lem:dc-ch7-2-5-pos}.
\end{proof}

\begin{proposition}[topology coincides]

  \label{prop:dc-ch7-2-6}
  \lean{eventually_riemannianEDist_lt, setOf_riemannianEDist_lt_subset_nhds, PseudoEMetricSpace.ofRiemannianMetric}
  \uses{def:dc-ch7-2-4}
  \dcref{ch7:2.6}
  \mathlibok
The topology induced by $d$ on $M$ coincides with the original topology on $M$.
\end{proposition}
\begin{proof}

  \uses{lem:dc-ch7-2-4-finite, lem:dc-ch7-2-5-chart}
\sketch Small coordinate segments have arbitrarily small Riemannian length, so metric balls are manifold neighborhoods. The short-curve lemma gives a metric ball inside every manifold neighborhood; hence the two topologies coincide.
\end{proof}

\begin{corollary}[distance is continuous]

  \label{cor:dc-ch7-2-7}
  \lean{Manifold.continuous_riemannianEDist_right}
  \uses{def:dc-ch7-2-4}
  \leanok
  \dcref{ch7:2.7}
If $p_0\in M$, the function $f:M\to\mathbb{R}$ given by $f(p)=d(p,p_0)$ is
  continuous.
\end{corollary}
\begin{proof}

  \leanok
  \uses{prop:dc-ch7-2-5, prop:dc-ch7-2-6}
\sketch The triangle inequality gives $|d(p,p_0)-d(q,p_0)|\le d(p,q)$, so the distance-to-$p_0$ function is $1$-Lipschitz and therefore continuous.
\end{proof}

\begin{theorem}[Hopf--Rinow equivalences]
  \label{thm:dc-ch7-2-8}
  \lean{Riemannian.Geodesic.hopfRinow}
  \uses{def:dc-ch7-2-2, def:dc-ch7-2-4, prop:dc-ch7-2-5}
  \dcref{ch7:2.8}
  Let $M$ be a Riemannian manifold and let $p\in M$. The following assertions are
  equivalent:
a) The exponential map $\exp_p$ is defined on all of $T_p(M)$.
b) The closed and bounded sets of $M$ are compact.
c) $M$ is complete as a metric space.
d) $M$ is geodesically complete.
e) There exists a sequence of compact subsets $K_n\subset M$, $K_n\subset K_{n+1}$
     and $\bigcup_n K_n=M$, such that if $q_n\notin K_n$ then $d(p,q_n)\to\infty$.

In addition, any of the statements above implies that
f) For any $q\in M$ there exists a geodesic $\gamma$ joining $p$ to $q$ with
     $\ell(\gamma)=d(p,q)$.
\end{theorem}
\begin{proof}

  \leanok
  \uses{cor:dc-ch3-3-9, lem:dc-ch7-2-8-b-to-c, lem:dc-ch7-2-8-a-to-b, lem:dc-ch7-2-8-b-iff-e, lem:dc-ch7-2-8-lipschitz, lem:dc-ch7-2-8-c-to-d}
\sketch The implication (a)$\Rightarrow$(f) follows by minimizing distance on a small geodesic sphere and extending the minimizing radial segment step by step. Then (a)$\Rightarrow$(b) follows because bounded sets lie in the compact image of a closed tangent ball. Properness implies metric completeness, while a Cauchy sequence along an incomplete geodesic converges and the uniform local geodesic flow extends it, proving (c)$\Rightarrow$(d); (d)$\Rightarrow$(a) is immediate. The compact-exhaustion criterion is equivalent to properness.
\end{proof}

\begin{lemma}[proper $\Rightarrow$ complete]

  \label{lem:dc-ch7-2-8-b-to-c}
  \lean{complete_of_proper}
  \dcref{ch7:2.8}
  \mathlibok
A metric space in which closed bounded sets are compact is complete.
\end{lemma}
\begin{proof}
\sketch A Cauchy sequence is bounded, so its closure is compact and it has a
  convergent subsequence. The Cauchy property then forces the whole sequence to
  converge.
\end{proof}

\begin{lemma}[properness via divergent compact exhaustions]

  \label{lem:dc-ch7-2-8-b-iff-e}
  \lean{OpenGA.properSpace_iff_exists_compact_exhaustion}
  \leanok
  \dcref{ch7:2.8}
Let $X$ be a metric space and $p\in X$. Closed bounded subsets of $X$ are
  compact if and only if there exists a sequence of compact subsets
  $K_n\subset X$ with $K_n\subset K_{n+1}$ and $\bigcup_n K_n=X$, such that
  $q_n\notin K_n$ implies $d(p,q_n)\to\infty$.
\end{lemma}
\begin{proof}

  \leanok
\sketch For a proper space, take $K_n$ to be closed balls about $p$ with radii tending to infinity. Conversely, the stated exhaustion makes every bounded set lie in some $K_n$; a closed bounded set is then closed in a compact $K_n$ and is compact.
\end{proof}

\begin{lemma}[geodesics are locally Lipschitz]

  \label{lem:dc-ch7-2-8-lipschitz}
  \lean{Riemannian.Geodesic.IsGeodesicOn.contMDiffOn, Riemannian.Geodesic.IsGeodesicOn.edist_le, Riemannian.Geodesic.IsGeodesicOn.dist_le}
  \uses{def:dc-ch3-2-1, lem:dc-ch3-2-1-constspeed, def:dc-ch7-2-4}
  \leanok
  \dcref{ch7:2.8}
Let $M$ be a Riemannian manifold whose distance $d$ is the Riemannian
  distance of \cref{def:dc-ch7-2-4}, and let $\gamma$ be a (continuous)
  geodesic defined on a connected open set $I$ of times. Then $\gamma$ is
  continuously differentiable, and for $a\le b$ in $I$,
  $$
  d(\gamma(a),\gamma(b))\ \le\ |\gamma'|\,(b-a),
  $$
  where $|\gamma'|=\sqrt{\langle\gamma',\gamma'\rangle}$ is the (constant)
  speed of $\gamma$ (\cref{lem:dc-ch3-2-1-constspeed}). In particular a
  normalized geodesic satisfies $d(\gamma(s),\gamma(t))\le|s-t|$: geodesics
  are locally Lipschitz.
\end{lemma}
\begin{proof}

  \leanok
  \uses{lem:dc-ch3-2-1-constspeed}
\sketch Geodesics have constant speed and the length of their restriction to $[a,b]$ is $|\gamma'|(b-a)$. Since distance is bounded by curve length, $d(\gamma(a),\gamma(b))\le |\gamma'|(b-a)$.
\end{proof}

\begin{lemma}[transformation law for the Christoffel symbols]

  \label{lem:dc-ch7-2-8-christoffel-transform}
  \lean{Riemannian.chartChristoffelContraction_change}
  \uses{def:dc-ch2-3-1, thm:dc-ch2-3-6}
  \leanok
  \dcref{ch7:2.8}
Let $\varphi_\alpha,\varphi_\beta$ be two charts of $M$ whose domains contain a
  point $x$, let $\tau=\varphi_\beta\circ\varphi_\alpha^{-1}$ be the transition map,
  $A=D\tau(\varphi_\alpha(x))$ its derivative and $D^2\tau$ its second derivative at
  $\varphi_\alpha(x)$. Then the Christoffel contractions
  $\Gamma^\alpha,\Gamma^\beta$ of the Levi–Civita connection of $g$ in the two
  charts satisfy, for all $v,w\in\mathbb R^n$,
  $$
  A\,\Gamma^\alpha(v,w)\;=\;\Gamma^\beta(Av,Aw)\;+\;D^2\tau(v,w).
  $$
\end{lemma}
\begin{proof}

  \leanok
\sketch Apply the second-order chain rule to the transition map $\tau$ and substitute the coordinate geodesic equation in chart $\alpha$. Rearranging yields the displayed transformation law for the Christoffel contractions.
\end{proof}

\begin{lemma}[the geodesic equation is chart-independent]

  \label{lem:dc-ch7-2-8-chart-independence}
  \lean{Riemannian.Geodesic.SolvesGeodesicODEAt, Riemannian.Geodesic.SolvesGeodesicODEAt.transfer, Riemannian.Geodesic.HasGeodesicEquationAt.solvesGeodesicODEAt, Riemannian.Geodesic.SolvesGeodesicODEAt.hasGeodesicEquationAt}
  \uses{def:dc-ch3-2-1, lem:dc-ch7-2-8-christoffel-transform}
  \leanok
  \dcref{ch7:2.8}
Let $\gamma$ be a curve in $M$ continuous at $t_0$ whose foot $\gamma(t_0)$ lies
  in the domains of two charts $\varphi_\alpha,\varphi_\beta$, and write
  $u_\alpha=\varphi_\alpha\circ\gamma$, $u_\beta=\varphi_\beta\circ\gamma$. If
  $u_\alpha$ is differentiable near $t_0$, its derivative is differentiable at
  $t_0$, and the second-order geodesic equation
  $u_\alpha''+\Gamma^\alpha(u_\alpha',u_\alpha')(u_\alpha)=0$ holds at $t_0$, then
  the same three statements hold for $u_\beta$. In particular the geodesic
  equation of \cref{def:dc-ch3-2-1}, read in the chart at the moving foot
  $\gamma(t_0)$, holds if and only if it holds in any fixed chart whose domain
  contains the foot.
\end{lemma}
\begin{proof}

  \leanok
  \uses{lem:dc-ch7-2-8-christoffel-transform}
\sketch Write $u_\beta=\tau\circ u_\alpha$ and use the first and second derivative chain rules together with \cref{lem:dc-ch7-2-8-christoffel-transform}. The differentiability and geodesic-equation conditions therefore transfer between the two chart readings.
\end{proof}

\begin{lemma}[flow solutions are geodesics; uniform local existence]

  \label{lem:dc-ch7-2-8-flow-geodesic}
  \lean{Riemannian.Geodesic.sprayBase, Riemannian.Geodesic.isGeodesicOn_sprayBase, Riemannian.Geodesic.hasDerivAt_chartReading_sprayBase, Riemannian.Geodesic.continuousOn_sprayBase, Riemannian.Geodesic.exists_seed_geodesic}
  \uses{def:dc-ch3-2-1, lem:dc-ch7-2-8-chart-independence, lem:dc-ch3-2-2-pointwise, lem:dc-ch3-2-6-reparam}
  \leanok
  \dcref{ch7:2.8}
Let $q\in M$ and let $\zeta=(x,w):J\to\mathbb R^n\times\mathbb R^n$ solve the
  first-order geodesic system $x'=w$, $w'=-\Gamma^q(w,w)(x)$ of the chart at $q$
  on an open set $J$ of times, with $x$ staying in the chart image. Then the base
  curve $\sigma=\varphi_q^{-1}\circ x$ is continuous and satisfies the (intrinsic)
  geodesic equation at every time of $J$, with chart velocity reading $w$.
  Consequently -- by Picard–Lindelöf applied to this system, whose right-hand side
  is $C^1$ on the chart image -- every initial datum $(p,v)$ admits, for some
  $b>0$, a continuous geodesic on $(-b,b)$ through $p$ with chart-$p$ velocity
  $v$ at time $0$.
\end{lemma}
\begin{proof}

  \leanok
  \uses{lem:dc-ch7-2-8-chart-independence, lem:dc-ch3-2-6-reparam}
\sketch The first-order system $x'=w$, $w'=-\Gamma^q(w,w)(x)$ is equivalent to the coordinate geodesic equation, and chart independence gives the intrinsic equation. Picard--Lindelöf supplies a uniform local solution for every initial state.
\end{proof}

\begin{lemma}[uniqueness of geodesics with given initial data]

  \label{lem:dc-ch7-2-8-uniqueness}
  \lean{Riemannian.Geodesic.IsGeodesicOn.eventuallyEq_of_deriv_chartReading_eq, Riemannian.Geodesic.IsGeodesicOn.eqOn_of_deriv_chartReading_eq}
  \uses{def:dc-ch3-2-1, lem:dc-ch7-2-8-chart-independence}
  \leanok
  \dcref{ch7:2.8}
Let $M$ be Hausdorff and let $\gamma_1,\gamma_2$ be continuous geodesics on an
  open interval $I$ which at some $t_0\in I$ have the same value and the same
  velocity reading in some chart around the common foot. Then
  $\gamma_1=\gamma_2$ on $I$. (Hausdorffness is necessary: on the line with two
  origins, two geodesics can agree except at a single time, where they pass
  through the two origins.)
\end{lemma}
\begin{proof}

  \leanok
  \uses{lem:dc-ch7-2-8-chart-independence}
\sketch In a common chart, both geodesics solve the same smooth first-order geodesic system with identical initial data, so ODE uniqueness gives equality near $t_0$. The coincidence set is open and closed in the connected interval, hence the curves agree on all of $I$.
\end{proof}

\begin{lemma}[local comparison of coordinate and Riemannian norms]

  \label{lem:dc-ch7-2-8-gram-bound}
  \lean{Riemannian.Geodesic.exists_sq_norm_le_chartMetricInner, Riemannian.Geodesic.HasGeodesicEquationAt.speedSq_eq_chartMetricInner_of_mem_source}
  \uses{def:dc-ch1-2-1, lem:dc-ch3-2-1-constspeed}
  \leanok
  \dcref{ch7:2.8}
Let $p_0\in M$ and let $G$ be the Gram matrix of $g$ in the chart at $p_0$.
  There are $c>0$ and a neighbourhood $V$ of $\varphi_{p_0}(p_0)$ inside the chart
  image such that $\|w\|^2\le c\,\langle w,w\rangle_{G(y)}$ for all $y\in V$ and
  all $w\in\mathbb R^n$. Moreover, along any geodesic whose foot lies in the chart
  at $p_0$, the Gram reading $\langle u',u'\rangle_{G(u)}$ of the chart velocity
  equals the intrinsic squared speed $\langle\gamma',\gamma'\rangle$.
\end{lemma}
\begin{proof}

  \leanok
\sketch Positive-definiteness and continuity of the Gram matrix on a compact chart neighborhood give a uniform comparison $\|w\|^2\le c\langle w,w\rangle_{G(y)}$. Along a geodesic, the coordinate Gram expression equals the intrinsic constant speed by the geodesic equation.
\end{proof}

\begin{lemma}[c $\Rightarrow$ d: complete metrics have global geodesics]

  \label{lem:dc-ch7-2-8-c-to-d}
  \lean{Riemannian.Geodesic.exists_tendsto_of_isGeodesicOn, Riemannian.Geodesic.IsGeodesicOn.exists_forward_extension, Riemannian.Geodesic.exists_global_geodesic, Riemannian.Geodesic.isGeodesicallyComplete_of_complete}
  \uses{def:dc-ch7-2-2, def:dc-ch7-2-4, lem:dc-ch7-2-8-lipschitz, lem:dc-ch7-2-8-flow-geodesic, lem:dc-ch7-2-8-uniqueness, lem:dc-ch7-2-8-gram-bound, lem:dc-ch3-2-6-reparam}
  \leanok
  \dcref{ch7:2.8}
Let $M$ be a Riemannian manifold whose distance is the Riemannian distance of
  $g$ and which is complete as a metric space. Then for every $p\in M$ and
  $v\in T_pM$ there is a geodesic $\gamma:\mathbb R\to M$ with $\gamma(0)=p$ and
  $\gamma'(0)=v$: $M$ is geodesically complete.
\end{lemma}
\begin{proof}

  \leanok
  \uses{lem:dc-ch7-2-8-lipschitz, lem:dc-ch7-2-8-flow-geodesic, lem:dc-ch7-2-8-uniqueness, lem:dc-ch7-2-8-gram-bound, lem:dc-ch3-2-6-reparam}
\sketch If a geodesic stopped at a finite endpoint, the local Lipschitz estimate would make its values a Cauchy family. Metric completeness gives a limit point, and the uniform local flow near that point extends the geodesic past the endpoint. Rescaling handles arbitrary initial velocities.
\end{proof}

\begin{lemma}[a $\Rightarrow$ b: geodesic completeness at a point makes closed balls compact]

  \label{lem:dc-ch7-2-8-a-to-b}
  \lean{Riemannian.Geodesic.properSpace_of_expDomainIntrinsic_eq_univ}
  \uses{def:dc-ch7-2-2, def:dc-ch7-2-4, lem:dc-ch7-2-8-step}
  \leanok
  \dcref{ch7:2.8}
Let $M$ be a connected Riemannian manifold whose distance is the Riemannian distance of
  $g$, and let $p\in M$ be a point at which $M$ is geodesically complete: every $v\in T_pM$
  generates a geodesic $\gamma:\mathbb R\to M$, defined for all time, with $\gamma(0)=p$ and
  $\gamma'(0)=v$. Then the closed and bounded subsets of $M$ are compact, i.e. $M$ is a
  \emph{proper} metric space.
\end{lemma}
\begin{proof}

  \leanok
  \uses{def:dc-ch7-2-4, lem:dc-ch7-2-8-step}
\sketch Geodesic completeness at $p$ makes every bounded set lie in the image under $\exp_p$ of a closed tangent ball. The image is compact, and closed subsets of it are compact; the geodesic-sphere step supplies the required radial continuation to reach any point in the bounded set.
\end{proof}

\begin{lemma}[length of a curve read in one chart]

  \label{lem:dc-ch7-2-8-chart-length}
  \lean{Riemannian.Geodesic.hasMFDerivAt_of_hasDerivAt_extChartAt, Riemannian.Geodesic.mfderiv_eq_of_hasDerivAt_extChartAt, Riemannian.Geodesic.enorm_mfderiv_eq_of_hasDerivAt_extChartAt, Riemannian.Geodesic.contDiffOn_extChartAt_comp, Riemannian.Geodesic.pathELength_eq_ofReal_integral_chartMetricInner}
  \uses{def:dc-ch7-2-4, lem:dc-ch7-2-8-gram-bound}
  \leanok
  \dcref{ch7:2.8}
Let $\gamma:[a,b]\to M$ be a $C^1$ curve whose image lies in the domain of a
  single chart $\varphi_\alpha$, and let $u=\varphi_\alpha\circ\gamma$ be its
  chart reading. Then at every $t$ the intrinsic velocity $\gamma'(t)\in
  T_{\gamma(t)}M$ is the readback of the coordinate velocity $\dot u(t)$ under
  the chart frame at $\gamma(t)$, its Riemannian norm is the Gram norm
  $|\gamma'(t)| = \sqrt{\langle\dot u(t),\dot u(t)\rangle_{G(u(t))}}$ of the
  coordinate velocity, and the length of $\gamma$ (\cref{def:dc-ch7-2-4}) is
  the coordinate integral
  $$
  \ell(\gamma)\;=\;\int_a^b \sqrt{\big\langle\dot u(t),\dot u(t)\big\rangle_{G(u(t))}}\,dt ,
  $$
  where $G$ is the Gram matrix of $g$ in the chart at $\alpha$.
\end{lemma}
\begin{proof}

  \leanok
  \uses{lem:dc-ch7-2-8-gram-bound}
\sketch The chart differential identifies $\gamma'(t)$ with $\dot u(t)$, so its norm is the Gram norm $\sqrt{\langle\dot u,\dot u\rangle_{G(u)}}$. Integrating this norm gives the stated length formula.
\end{proof}

\begin{lemma}[the metric normal ball]

  \label{lem:dc-ch7-2-8-normal-ball-dist}
  \lean{Riemannian.Exponential.exists_isOpen_expMap_image, Riemannian.Exponential.exists_pathELength_expMap_ray, Riemannian.Exponential.exists_le_pathELength, Riemannian.Exponential.exists_edist_expMap_ball}
  \uses{def:dc-ch7-2-4, def:dc-ch3-expmap, lem:dc-ch3-2-9-c1diffeo, lem:dc-ch3-3-6-radius, lem:dc-ch3-3-6-reach, lem:dc-ch3-3-6-rayspeed, lem:dc-ch7-2-8-gram-bound, lem:dc-ch7-2-8-chart-length}
  \leanok
  \dcref{ch7:2.8}
Let $M$ be a Riemannian manifold whose distance $d$ is the Riemannian
  distance of \cref{def:dc-ch7-2-4} and let $p\in M$. There are
  $\varepsilon,\delta>0$ such that $\exp_p$ is injective and open on
  $B_\varepsilon(0)\subset T_pM$ and:
  \begin{enumerate}
  \item (\emph{radial geodesics realize the distance})
    $d\big(p,\exp_p v\big)=|v|_p$ for every $v$ with $\|v\|<\varepsilon$;
  \item (\emph{escape estimate}) $d(p,q)\ge\delta$ for every
    $q\notin\exp_p\big(B_\varepsilon(0)\big)$ -- the normal ball contains the
    metric ball $B_\delta(p)$.
  \end{enumerate}
\end{lemma}
\begin{proof}

  \leanok
  \uses{lem:dc-ch3-2-9-c1diffeo, lem:dc-ch3-3-6-radius, lem:dc-ch3-3-6-reach, lem:dc-ch3-3-6-rayspeed, lem:dc-ch7-2-8-gram-bound, lem:dc-ch7-2-8-chart-length}
\sketch The Gauss lemma and the local diffeomorphism properties of $\exp_p$ give radial minimizing geodesics inside a sufficiently small normal ball. Compactness of the unit sphere and the chart length comparison give a positive lower bound for curves that leave this ball.
\end{proof}

\begin{lemma}[exponential rays are geodesics]

  \label{lem:dc-ch7-2-8-ray-geodesic}
  \lean{Riemannian.Exponential.exists_isGeodesicOn_expMap_ray}
  \uses{def:dc-ch3-expmap, def:dc-ch3-2-1, lem:dc-ch7-2-8-chart-independence}
  \leanok
  \dcref{ch7:2.8}
Let $p\in M$. There are $\rho>0$ and $b>1$ such that for every $u\in T_pM$
  with $\|u\|<\rho$ the ray
  $$
  \gamma_u : t \longmapsto \exp_p(t\,u)
  $$
  is defined and stays in the chart at $p$ for $|t|<b$, starts at
  $\gamma_u(0)=p$ with velocity $u$, is continuous, and satisfies the geodesic
  equation (\cref{def:dc-ch3-2-1}) at every $t\in(-b,b)$. In particular the
  segment $t\in[0,1]$, joining $p$ to $\exp_p u$, is a geodesic segment.
\end{lemma}
\begin{proof}

  \leanok
  \uses{lem:dc-ch7-2-8-chart-independence}
\sketch The radial identity $\exp_p(tu)$ is the geodesic with initial velocity $u$ on a uniform small interval, by the exponential-map construction and chart-independent geodesic equation.
\end{proof}

\begin{lemma}[sphere-minimum decomposition of the distance]

  \label{lem:dc-ch7-2-8-sphere-min}
  \lean{Riemannian.Exponential.exists_normalSphere_min_edist}
  \uses{def:dc-ch7-2-4, lem:dc-ch7-2-8-gram-bound, lem:dc-ch7-2-8-chart-length, lem:dc-ch7-2-8-normal-ball-dist}
  \leanok
  \dcref{ch7:2.8}
Let $M$ be a Riemannian manifold whose distance $d$ is the Riemannian
  distance of \cref{def:dc-ch7-2-4} and let $p\in M$. There are
  $\varepsilon, c>0$, depending only on $p$, with the following property: for
  every $\delta>0$ with $\sqrt c\,\delta<\varepsilon$, writing
  $S_\delta(p)=\exp_p\{\,|z|_p=\delta\,\}$ for the geodesic sphere of radius
  $\delta$, and for every $q\in M$ with $d(p,q)\ge\delta$, there is a point
  $x_0=\exp_p z_0\in S_\delta(p)$ such that
  $$
  d(p,q)\;=\;\delta+d(x_0,q),
  \qquad
  d(x_0,q)\;=\;\min_{x\in S_\delta(p)} d(x,q).
  $$
\end{lemma}
\begin{proof}

  \leanok
  \uses{lem:dc-ch7-2-8-gram-bound, lem:dc-ch7-2-8-chart-length, lem:dc-ch7-2-8-normal-ball-dist}
\sketch The geodesic sphere $S_\delta(p)$ is compact for small $\delta$, so $x\mapsto d(x,q)$ attains a minimum. Any curve from $p$ to $q$ crosses this sphere; the normal-ball distance formula shows that the first crossing contributes exactly $\delta$, yielding the displayed decomposition.
\end{proof}

\begin{lemma}[the geodesic-sphere step]

  \label{lem:dc-ch7-2-8-step}
  \lean{Riemannian.Exponential.exists_minimizing_step}
  \uses{def:dc-ch7-2-4, lem:dc-ch7-2-8-ray-geodesic, lem:dc-ch7-2-8-sphere-min}
  \leanok
  \dcref{ch7:2.8}
Let $M$ be a Riemannian manifold whose distance $d$ is the Riemannian
  distance of \cref{def:dc-ch7-2-4} and let $p\in M$. There is $\delta_0>0$
  such that for every $0<\delta<\delta_0$ and every $q\in M$ with
  $d(p,q)\ge\delta$ there is a continuous geodesic segment
  $\gamma:[0,1]\to M$ with
  $$
  \gamma(0)=p,\qquad d\big(p,\gamma(1)\big)=\delta,\qquad
  d(p,q)\;=\;\delta+d\big(\gamma(1),q\big).
  $$
\end{lemma}
\begin{proof}

  \leanok
  \uses{lem:dc-ch7-2-8-ray-geodesic, lem:dc-ch7-2-8-sphere-min}
\sketch Choose a minimizing point $x_0\in S_\delta(p)$ from \cref{lem:dc-ch7-2-8-sphere-min}. The radial geodesic from $p$ to $x_0$ and the distance decomposition give the stated geodesic segment and equality.
\end{proof}

\begin{corollary}[compact $\Rightarrow$ complete]

  \label{cor:dc-ch7-2-9}
  \lean{Riemannian.Geodesic.isGeodesicallyComplete_of_compactSpace}
  \uses{def:dc-ch7-2-2}
  \leanok
  \dcref{ch7:2.9}
If $M$ is compact then $M$ is complete.
\end{corollary}
\begin{proof}

  \leanok
  \uses{lem:dc-ch7-2-8-c-to-d}
\sketch A compact metric space is complete. The metric induced by a Riemannian metric agrees with the manifold topology, so compactness of $M$ gives metric completeness and then geodesic completeness by \cref{thm:dc-ch7-2-8}.
\end{proof}

\begin{lemma}[a metric-preserving inclusion is $1$-Lipschitz for the Riemannian distances]

  \label{lem:dc-ch7-2-10-lipschitz}
  \lean{Riemannian.DCPreservesMetric.enorm_mfderiv_eq, Riemannian.DCPreservesMetric.pathELength_comp_eq, Riemannian.DCPreservesMetric.riemannianEDist_le}
  \uses{def:dc-ch7-2-4, def:dc-ch1-2-2}
  \leanok
  \dcref{ch7:2.10}
Let $\iota:N\to M$ be a smooth map which \emph{preserves the metric} -- a local isometry
  for the induced metric $g_N=\iota^*g_M$ (\cref{def:dc-ch1-2-2}), i.e.
  $\langle d\iota_p u,d\iota_p v\rangle_M=\langle u,v\rangle_N$ for all $p\in N$ and all
  $u,v\in T_pN$. Then $\iota$ does not increase the Riemannian distance:
  $d_M(\iota a,\iota b)\le d_N(a,b)$ for all $a,b\in N$.
\end{lemma}
\begin{proof}

  \leanok
  \uses{def:dc-ch7-2-4}
\sketch Metric preservation makes the length of every piecewise differentiable curve in $N$ equal to the length of its image in $M$. Taking infima over curves gives $d_M(\iota a,\iota b)\le d_N(a,b)$.
\end{proof}

\begin{corollary}[closed submanifold complete]

  \label{cor:dc-ch7-2-10}
  \lean{Riemannian.completeSpace_of_isClosedEmbedding_dcPreservesMetric}
  \uses{lem:dc-ch7-2-10-lipschitz, prop:dc-ch7-2-6, thm:dc-ch7-2-8}
  \leanok
  \dcref{ch7:2.10}
A closed submanifold of a complete Riemannian manifold is complete in the induced
  metric; in particular, the closed submanifolds of Euclidean space are complete.
\end{corollary}
\begin{proof}

  \leanok
  \uses{lem:dc-ch7-2-10-lipschitz, prop:dc-ch7-2-6}
\sketch A Cauchy sequence in the closed submanifold is Cauchy in the complete ambient manifold by the preceding Lipschitz estimate, hence converges there. Closedness puts the limit in the submanifold, proving completeness.
\end{proof}

\section{The Theorem of Hadamard}

\begin{theorem}[Hadamard theorem]

  \label{thm:dc-ch7-3-1}
  \lean{Riemannian.Jacobi.hadamardDiffeomorphOfNonpos_complete, Riemannian.Jacobi.hadamardDiffeomorphOfNonpos_complete_coe}
  \uses{lem:dc-ch7-3-2, lem:dc-ch7-3-3, lem:dc-ch7-3-1-covering-diffeo, lem:dc-ch7-3-1-assembly, lem:dc-ch7-3-4-pullback-metric, lem:dc-ch7-3-4-pole-diffeo, lem:dc-ch7-3-expmap-global, lem:dc-ch7-3-expmap-smooth, thm:dc-ch7-2-8}
  \leanok
  \dcref{ch7:3.1}
Let $M$ be a complete Riemannian manifold, simply connected, with sectional
  curvature $K(p,\sigma)\le 0$ for all $p\in M$ and for all $\sigma\subset T_p(M)$.
  Then $M$ is diffeomorphic to $\mathbb{R}^n$, $n=\dim M$; more precisely,
  $\exp_p:T_pM\to M$ is a diffeomorphism.
\end{theorem}
\begin{proof}
\sketch Nonpositive curvature makes the Jacobi energy strictly positive away from its initial zero, so $d\exp_p$ is everywhere invertible and $\exp_p$ is a local diffeomorphism. The pullback metric on $T_pM$ makes $\exp_p$ metric-preserving; completeness and path lifting make it a covering map. Simple connectedness of $M$ makes the covering bijective, hence $\exp_p$ is a diffeomorphism.
\end{proof}

\begin{lemma}[the Jacobi energy is convex, and $J$ has no interior zero, when $K\le 0$]

  \label{lem:dc-ch7-3-2-no-conjugate}
  \lean{Riemannian.Jacobi.jacobiCoefOp_quadratic_nonpos, Riemannian.Jacobi.frameJacobi_ne_zero_of_nonpos}
  \uses{def:dc-ch5-2-1, lem:dc-ch5-2-1-real-curvature}
  \leanok
  \dcref{ch7:3.2}
Let $J=\sum_i f_i e_i$ be a Jacobi field in a parallel orthonormal frame along the
  geodesic $\gamma$, with $J(0)=0$ and nonzero initial velocity $J'(0)\ne 0$. Under
  nonpositive sectional curvature, read in the chart as $\langle R(\gamma',w)\gamma',w\rangle\le 0$
  for every $w$, the energy $q=\langle J,J\rangle$ satisfies
  $$
  \langle J,J\rangle''=2|J'|^2-2\langle R(\gamma',J)\gamma',J\rangle\ge 0,
  $$
  because the curvature term $\langle R(\gamma',J)\gamma',J\rangle=\sum_{ij}a_{ij}f_if_j\le 0$
  is the (negative semidefinite) quadratic form of the Jacobi coefficient
  $a_{ij}=\langle R(\gamma',e_i)\gamma',e_j\rangle$. Since $q(0)=0$ and $q'(0)=0$, convexity
  gives $q'\ge 0$, so $q$ is nondecreasing; nontriviality ($J'(0)\ne 0$) then forces
  $q(t)>0$ for every $t>0$, i.e. $J(t)\ne 0$. Hence no point $\gamma(t)$, $t>0$, is
  conjugate to $\gamma(0)$.
\end{lemma}

\begin{lemma}[no conjugate points when $K\le 0$]

  \label{lem:dc-ch7-3-2}
  \lean{Riemannian.Jacobi.not_isConjugatePointAt_one_of_nonpos, Riemannian.Jacobi.expDifferential_isEquiv_of_nonpos, Riemannian.Jacobi.expMapGlobal_locallyInjective_of_nonpos, Riemannian.Jacobi.isLocalDiffeomorph_expMapGlobal_hadamard_of_nonpos_complete}
  \uses{def:dc-ch5-2-1, def:dc-ch5-3-1, def:dc-ch5-3-4, lem:dc-ch7-3-2-no-conjugate, lem:dc-ch7-3-expmap-smooth}
  \leanok
  \dcref{ch7:3.2}
Let $M$ be a complete Riemannian manifold with $K(p,\sigma)\le 0$ for all $p\in M$
  and for all $\sigma\subset T_pM$. Then for all $p\in M$ the conjugate locus
  $C(p)=\varnothing$; in particular the exponential map $\exp_p:T_pM\to M$ is a local
  diffeomorphism.
\end{lemma}
\begin{proof}

  \uses{lem:dc-ch7-3-2-no-conjugate, prop:dc-ch5-3-5}
\sketch Apply \cref{lem:dc-ch7-3-2-no-conjugate} to every geodesic from $p$. The absence of conjugate points makes the differential of $\exp_p$ an isomorphism at every vector; the inverse function theorem then gives a smooth local diffeomorphism.
\end{proof}

\begin{lemma}[length expansion under a metric-expanding map]

  \label{lem:dc-ch7-3-3-length-expansion}
  \lean{Riemannian.DCExpandsMetric, Riemannian.DCExpandsMetric.pathELength_le, Riemannian.DCExpandsMetric.riemannianEDist_le_pathELength_comp}
  \uses{def:dc-ch7-2-4, def:dc-ch1-2-9}
  \leanok
  \dcref{ch7:3.3}
Let $f:M\to N$ be a differentiable map between Riemannian manifolds which
  \emph{expands the metric}: $|df_p(v)|\ge|v|$ for every $p\in M$ and every
  $v\in T_pM$. Then for every piecewise differentiable curve
  $\bar c:[a,b]\to M$,
  $$
  \ell(f\circ\bar c)\ \ge\ \ell(\bar c)\ \ge\ d\big(\bar c(a),\bar c(b)\big).
  $$
  In particular $d\big(\bar c(a),\bar c(b)\big)\le\ell(f\circ\bar c)$: the
  source distance between the endpoints of a curve is bounded by the length of
  its image.
\end{lemma}
\begin{proof}

  \leanok
  \uses{def:dc-ch7-2-4}
\sketch Along a piecewise differentiable curve, $|d(f\circ\bar c)/dt|\ge|d\bar c/dt|$. Integrating yields $\ell(f\circ\bar c)\ge\ell(\bar c)$, and the definition of distance gives the second inequality.
\end{proof}

\begin{lemma}[lifts of finite-length curves are relatively compact]

  \label{lem:dc-ch7-3-3-relatively-compact-lift}
  \lean{Riemannian.DCExpandsMetric.dist_le_pathELength_comp, Riemannian.DCExpandsMetric.lift_mem_closedBall, Riemannian.DCExpandsMetric.lift_image_subset_isCompact}
  \uses{lem:dc-ch7-3-3-length-expansion, thm:dc-ch7-2-8, def:dc-ch7-2-4}
  \leanok
  \dcref{ch7:3.3}
Let $f:M\to N$ expand the metric ($|df_p(v)|\ge|v|$ for all $p,v$), and let $M$
  be a complete Riemannian manifold. Then for every piecewise differentiable
  curve $\bar c:[a,b]\to M$ whose image $f\circ\bar c$ has finite length, and for
  every $t\in[a,b]$,
  $$
  d\big(\bar c(a),\bar c(t)\big)\ \le\ \ell(f\circ\bar c),
  $$
  so $\bar c(t)$ lies in the closed metric ball $\bar B\big(\bar c(a),\ell(f\circ\bar c)\big)$.
  Consequently the whole image $\bar c([a,b])$ is contained in a single
  \emph{compact} subset $K\subset M$.
\end{lemma}
\begin{proof}

  \leanok
  \uses{lem:dc-ch7-3-3-length-expansion, thm:dc-ch7-2-8}
\sketch Apply the length estimate to each restriction of the lifted curve. Every lift point lies in the closed ball of radius $\ell(f\circ\bar c)$ about its initial point, which is compact by Hopf--Rinow.
\end{proof}

\begin{lemma}[local liftability of a curve through a local homeomorphism]

  \label{lem:dc-ch7-3-3-local-lift}
  \lean{Riemannian.IsLocalHomeomorph.exists_rightward_lift, Riemannian.IsLocalHomeomorph.exists_extend_lift}
  \leanok
  \dcref{ch7:3.3}
Let $f:M\to N$ be a local homeomorphism and $c:[0,1]\to N$ a continuous curve.
  \begin{enumerate}
    \item For every $s$ and every point $q\in M$ over $c(s)$ (that is,
      $f(q)=c(s)$) the curve $c$ lifts through $q$ on a nondegenerate interval to
      the right: there is a continuous $\bar c:[s,s+\varepsilon]\to M$ with
      $\bar c(s)=q$ and $f\circ\bar c=c$.
    \item Consequently, if $c$ lifts to a continuous $\bar c:[a,t]\to M$ with
      $f\circ\bar c=c$, then the lift extends to $[a,t+\delta]$ for some
      $\delta>0$: the set of parameters over which $c$ lifts starting from a fixed
      initial point is open to the right.
  \end{enumerate}
\end{lemma}
\begin{proof}

  \leanok
\sketch A local inverse of $f$ near any point over $c(s)$ lifts $c$ on a small interval. Applying this at the endpoint of a partial lift extends the lift to the right.
\end{proof}

\begin{lemma}[path lifting for a metric-expanding local diffeomorphism]

  \label{lem:dc-ch7-3-3-path-lifting}
  \lean{Riemannian.DCExpandsMetric.exists_pathLift, Riemannian.IsLocalDiffeomorph.contMDiffOn_lift, Riemannian.DCExpandsMetric.dist_lift_le}
  \uses{thm:dc-ch7-2-8, def:dc-ch7-2-4, lem:dc-ch7-3-3-length-expansion, lem:dc-ch7-3-3-relatively-compact-lift, lem:dc-ch7-3-3-local-lift}
  \leanok
  \dcref{ch7:3.3}
Let $M$ be a complete Riemannian manifold and let $f:M\to N$ be a local
  diffeomorphism which expands the metric ($|df_p(v)|\ge|v|$ for all $p\in M$ and
  all $v\in T_pM$). Then $f$ has the \emph{path-lifting property}: for every
  $C^1$ curve $c:[0,1]\to N$ of finite length and every $q\in M$ with $f(q)=c(0)$,
  there exists a continuous curve $\bar c:[0,1]\to M$ with $\bar c(0)=q$ and
  $f\circ\bar c=c$.

Moreover any continuous lift of a $C^1$ curve through a local diffeomorphism is
  itself $C^1$ (locally it equals $(f|_V)^{-1}\circ c$ for a local inverse), and
  the endpoints $\bar c(t)$ of a partial lift satisfy
  $d\big(\bar c(0),\bar c(t)\big)\le\ell(c)$ (the estimate that traps the lift in a
  compact ball).
\end{lemma}
\begin{proof}

  \leanok
  \uses{thm:dc-ch7-2-8, lem:dc-ch7-3-3-length-expansion, lem:dc-ch7-3-3-relatively-compact-lift, lem:dc-ch7-3-3-local-lift}
\sketch Start with the local lift and extend it maximally. The length estimate confines its image to a compact ball, so an accumulation point at a finite endpoint exists; local liftability then extends the lift, proving existence on all of $[0,1]$. Local inverses also show any lift of a $C^1$ curve is $C^1$ and satisfy the stated distance bound.
\end{proof}

\begin{lemma}[uniqueness of continuous lifts]

  \label{lem:dc-ch7-3-3-lift-unique}
  \lean{Riemannian.IsLocalHomeomorph.eqOn_lift, Riemannian.IsLocalDiffeomorph.eqOn_pathLift, Riemannian.DCExpandsMetric.existsUnique_pathLift_eqOn}
  \uses{lem:dc-ch7-3-3-path-lifting}
  \leanok
  \dcref{ch7:3.3}
Let $f:M\to N$ be a local diffeomorphism out of a Riemannian manifold $M$ (in
  particular a Hausdorff space). If $\bar c_1,\bar c_2:S\to M$ are two continuous
  lifts of the same curve $c$ (i.e.\ $f\circ\bar c_i=c$) over a connected parameter
  set $S$, and $\bar c_1(t_0)=\bar c_2(t_0)$ for some $t_0\in S$, then
  $\bar c_1=\bar c_2$ on all of $S$. In particular the continuous lift of a $C^1$
  curve $c:[0,1]\to N$ through a chosen point $q$ over $c(0)$, produced by
  \cref{lem:dc-ch7-3-3-path-lifting}, is \emph{unique}.
\end{lemma}
\begin{proof}

  \leanok
  \uses{lem:dc-ch7-3-3-path-lifting}
\sketch Two lifts agreeing at one parameter agree on a neighborhood by local injectivity. Their coincidence set is open and closed in the connected parameter domain, so the lifts coincide everywhere.
\end{proof}

\begin{lemma}[finite length of a chart-contained $C^1$ curve]

  \label{lem:dc-ch7-3-3-chart-length-finite}
  \lean{Riemannian.Geodesic.pathELength_ne_top_of_mapsTo_chartSource}
  \uses{def:dc-ch1-2-9}
  \leanok
  \dcref{ch7:3.3}
If $\gamma:[a,b]\to N$ is a $C^1$ curve whose image lies in the source of a single
  chart of $N$, then its length $\ell(\gamma)$ is finite. Reading the length through
  the chart presents it as $\mathrm{ofReal}\big(\int_a^b\sqrt{\langle\gamma',\gamma'\rangle}\big)$,
  a real number.
\end{lemma}
\begin{proof}

  \leanok
  \uses{def:dc-ch1-2-9}
\sketch In a single chart the metric coefficients and the $C^1$ velocity are bounded on the compact parameter interval, so the coordinate length integral is finite.
\end{proof}

\begin{lemma}[closed image; surjectivity onto a connected target]

  \label{lem:dc-ch7-3-3-closed-image}
  \lean{Riemannian.DCExpandsMetric.isClosed_range, Riemannian.DCExpandsMetric.surjective}
  \uses{lem:dc-ch7-3-3-path-lifting, lem:dc-ch7-3-3-chart-length-finite}
  \leanok
  \dcref{ch7:3.3}
Let $M$ be a complete Riemannian manifold and $f:M\to N$ a local diffeomorphism
  which expands the metric. Then the image $f(M)$ is \emph{closed} in $N$; if $N$ is
  connected and $M$ nonempty, $f$ is \emph{surjective}.
\end{lemma}
\begin{proof}

  \leanok
  \uses{lem:dc-ch7-3-3-path-lifting, lem:dc-ch7-3-3-chart-length-finite}
\sketch If $f(x_n)$ converges, join a fixed image point to the limit by a finite-length chart curve and lift its initial segments. The lift remains in a compact ball, hence has an accumulation point whose image is the limit; therefore $f(M)$ is closed. A nonempty closed image in connected $N$ is all of $N$.
\end{proof}

\begin{lemma}[expansion $\Rightarrow$ covering map]

  \label{lem:dc-ch7-3-3}
  \lean{Riemannian.DCExpandsMetric.isCoveringMap, Riemannian.DCExpandsMetric.exists_trivialization_at, Riemannian.IsLocalHomeomorph.discreteTopology_fiber}
  \uses{thm:dc-ch7-2-8, lem:dc-ch7-3-3-length-expansion, lem:dc-ch7-3-3-relatively-compact-lift, lem:dc-ch7-3-3-local-lift, lem:dc-ch7-3-3-path-lifting, lem:dc-ch7-3-3-lift-unique, lem:dc-ch7-3-3-chart-length-finite, lem:dc-ch7-3-3-closed-image}
  \leanok
  \dcref{ch7:3.3}
Let $M$ be a complete Riemannian manifold and let $f:M\to N$ be a local
  diffeomorphism onto a Riemannian manifold $N$ which has the following property:
  for all $p\in M$ and for all $v\in T_pM$, we have $|df_p(v)|\ge|v|$. Then $f$ is a
  covering map.
\end{lemma}
\begin{proof}

  \leanok
\sketch Path lifting and uniqueness give local trivializations over points in the image; the closed-image result supplies a neighborhood disjoint from the image at points outside it. Thus every point has an evenly covered neighborhood and $f$ is a covering map.
\end{proof}

\begin{lemma}[a covering map onto a simply connected base is a homeomorphism]

  \label{lem:dc-ch7-3-1-covering-simplyconnected}
  \lean{IsCoveringMap.bijective_of_simplyConnected, IsCoveringMap.homeomorphOfSimplyConnected}
  \leanok
  \dcref{ch7:3.1}
Let $f:E\to X$ be a covering map whose total space $E$ is connected and nonempty,
  and whose base $X$ is simply connected and locally path connected. Then $f$ is a
  homeomorphism; in particular it is bijective.
\end{lemma}
\begin{proof}

  \leanok
\sketch Lift the identity of the simply connected base to a global section. It proves surjectivity. The section composed with $f$ and the identity on the connected total space are lifts of the same map agreeing at one point, so uniqueness gives injectivity; an open continuous bijection is a homeomorphism.
\end{proof}

\begin{lemma}[metric-expanding local diffeomorphism onto a simply connected manifold]

  \label{lem:dc-ch7-3-1-covering-diffeo}
  \lean{Riemannian.DCExpandsMetric.diffeomorphOfSimplyConnected}
  \uses{lem:dc-ch7-3-3, lem:dc-ch7-3-1-covering-simplyconnected}
  \leanok
  \dcref{ch7:3.1}
Let $f:M\to N$ be a smooth local diffeomorphism between Riemannian manifolds which
  expands the metric ($|df_p(v)|\ge|v|$ for all $p\in M$ and all $v\in T_pM$), let $M$
  be complete, and let $N$ be simply connected. Then $f$ is a diffeomorphism.
\end{lemma}
\begin{proof}

  \leanok
  \uses{lem:dc-ch7-3-3, lem:dc-ch7-3-1-covering-simplyconnected}
\sketch The expansion and completeness hypotheses give a covering map by \cref{lem:dc-ch7-3-3}. The simply connected-base lemma makes it bijective, and a bijective local diffeomorphism is a diffeomorphism.
\end{proof}

\begin{lemma}[the global exponential map of a complete manifold]

  \label{lem:dc-ch7-3-expmap-global}
  \lean{Riemannian.Geodesic.intrinsicMaximalGeodesic_eq_globalGeodesic, Riemannian.Geodesic.globalGeodesic, Riemannian.Geodesic.globalGeodesic_eq, Riemannian.Geodesic.globalGeodesic_smul, Riemannian.Exponential.expMapIntrinsic_eq_expMapGlobal, Riemannian.Exponential.expMapGlobal, Riemannian.Exponential.expMapGlobal_eq_of_isGeodesic, Riemannian.Exponential.expMapGlobal_smul, Riemannian.Exponential.expMapIntrinsic_surjective}
  \uses{def:dc-ch7-2-2, thm:dc-ch7-2-8}
  \leanok
  \dcref{ch7:2.2}
Let $M$ be a complete Riemannian manifold and $p\in M$. For every $v\in T_pM$ there is a
  geodesic $\gamma_v:\mathbb{R}\to M$, defined on all of $\mathbb{R}$, with $\gamma_v(0)=p$
  and $\gamma_v'(0)=v$, and it is the unique such geodesic. The exponential map
  $\exp_p:T_pM\to M$, $\exp_p(v)=\gamma_v(1)$, is therefore defined on all of $T_pM$, and it
  is traced radially: $\exp_p(cv)=\gamma_v(c)$ for all $c\in\mathbb{R}$, so the curves
  $c\mapsto\exp_p(cv)$ are precisely the geodesics of $M$ issuing from $p$. If moreover $M$
  is connected, then $\exp_p$ is surjective.
\end{lemma}
\begin{proof}

  \leanok
  \uses{def:dc-ch7-2-2, thm:dc-ch7-2-8}
\sketch Hopf--Rinow gives a global geodesic for each initial vector, and uniqueness makes $\gamma_v$ and $\exp_p$ well defined. Affine reparametrization gives $\exp_p(cv)=\gamma_v(c)$; minimizing geodesics from $p$ show surjectivity when $M$ is connected.
\end{proof}

\begin{lemma}[the exponential map of a complete manifold is $C^\infty$]

  \label{lem:dc-ch7-3-expmap-smooth}
  \lean{Riemannian.Exponential.contMDiff_expMapIntrinsic, Riemannian.Exponential.contMDiff_expMapGlobal, Riemannian.Jacobi.exists_geodesic_flow_step_contDiff, Riemannian.Jacobi.exists_geodesic_flow_step_contDiff_manifold_ball, Riemannian.Jacobi.exists_geodesic_flowstep_partition_contDiff, Riemannian.Jacobi.exists_geodesic_contDiff_chain_nbhd, Riemannian.Jacobi.contDiffAt_stateTransition}
  \uses{lem:dc-ch7-3-expmap-global}
  \leanok
  \dcref{ch7:3.1}
Let $M$ be a complete Riemannian manifold and $p\in M$. Then the exponential map
  $\exp_p:T_pM\to M$ is $C^\infty$. This is the smooth dependence of the geodesic flow on its
  initial conditions, specialised to the exponential map.
\end{lemma}
\begin{proof}

  \leanok
  \uses{lem:dc-ch7-3-expmap-global}
\sketch Cover the compact image of a geodesic segment by finitely many charts. Smooth local geodesic flows and smooth chart transitions compose to a smooth endpoint map, and the initial velocity enters affinely; this proves smoothness of $\exp_p$ near each vector.
\end{proof}

\begin{lemma}[Hadamard assembly from geodesic completeness at a point]

  \label{lem:dc-ch7-3-1-assembly}
  \lean{Riemannian.DCExpandsMetric.diffeomorphOfSimplyConnectedOfGeodesicCompleteAt}
  \uses{lem:dc-ch7-3-1-covering-diffeo, lem:dc-ch7-2-8-a-to-b, lem:dc-ch7-3-3}
  \leanok
  \dcref{ch7:3.1}
Let $f:M\to N$ be a smooth local diffeomorphism between Riemannian manifolds which
  \emph{expands the metric} ($|df_p(v)|\ge|v|$ for all $p\in M$ and all $v\in T_pM$), let
  $N$ be simply connected, and suppose $M$ is connected and \emph{geodesically complete at
  some point $o$}: every $v\in T_oM$ generates a geodesic $\gamma:\mathbb R\to M$, defined
  for all time, with $\gamma(0)=o$ and $\gamma'(0)=v$. Then $f$ is a diffeomorphism.
\end{lemma}
\begin{proof}

  \leanok
  \uses{lem:dc-ch7-3-1-covering-diffeo, lem:dc-ch7-2-8-a-to-b}
\sketch Completeness at the base point implies properness by \cref{lem:dc-ch7-2-8-a-to-b}, hence metric completeness. The expansion lemma gives a covering map, and the simply connected-base result upgrades it to a diffeomorphism.
\end{proof}

\subsection{Proof of the Hadamard theorem}

\begin{lemma}[the tangent space with the pulled-back metric]

  \label{lem:dc-ch7-3-4-pullback-metric}
  \lean{Riemannian.HadamardModel, Riemannian.HadamardModel.flatMetric, Riemannian.HadamardModel.pullbackMetric, Riemannian.HadamardModel.dcExpandsMetric_pullbackMetric}
  \uses{def:dc-ch7-2-4, def:dc-ch1-2-2, lem:dc-ch7-3-3-length-expansion}
  \leanok
  \dcref{ch7:3.4}
Let $M$ be a Riemannian manifold, $p\in M$, and $f:T_pM\to M$ a smooth local
  diffeomorphism (for instance $f=\exp_p$ at a pole). Regard the tangent space $T_pM$ as
  a smooth manifold in its own right, modelled on $\mathbb R^n=T_pM$ via the single global
  identity chart. Then $T_pM$ carries the \emph{pulled-back Riemannian metric} $f^*g$,
  $\langle u,v\rangle^{f^*g}_x=\langle df_x(u),df_x(v)\rangle_{f(x)}$, for which $f$ is a
  local isometry, hence a metric-expander: $|df_x(v)|=|v|$ for all $x,v$.
\end{lemma}
\begin{proof}

  \leanok
  \uses{lem:dc-ch7-3-3-length-expansion}
\sketch Define $\langle u,v\rangle^{f^*g}_x=\langle df_xu,df_xv\rangle_{f(x)}$. Since $df_x$ is an isomorphism, this is a smooth positive-definite metric and $f$ is a local isometry.
\end{proof}

\begin{lemma}[pole $\Rightarrow$ diffeomorphism]

  \label{lem:dc-ch7-3-4-pole-diffeo}
  \lean{Riemannian.HadamardModel.diffeomorphOfPole, Riemannian.HadamardModel.diffeomorphOfPole_coe, Riemannian.HadamardModel.expDiffeomorphOfPole}
  \uses{lem:dc-ch7-3-1-assembly, lem:dc-ch7-3-4-pullback-metric, lem:dc-ch7-3-expmap-global}
  \leanok
  \dcref{ch7:3.4}
Let $M$ be a complete, simply connected Riemannian manifold, $p\in M$, and $f:T_pM\to M$
  a smooth local diffeomorphism which is metric-expanding for its own pulled-back metric
  $f^*g$. If the radial lines $s\mapsto sv$ of $T_pM$ are geodesics of $f^*g$, then $f$ is a
  diffeomorphism $T_pM\to M$.
\end{lemma}
\begin{proof}

  \leanok
  \uses{lem:dc-ch7-3-1-assembly, lem:dc-ch7-3-4-pullback-metric}
\sketch Equip $T_pM$ with $f^*g$. The radial lines are complete geodesics through the origin, so the source is complete at that point; the assembly lemma then makes the metric-preserving local diffeomorphism $f$ a diffeomorphism.
\end{proof}

\begin{lemma}[the poles theorem reduces to the rays being pullback-geodesics]

  \label{lem:dc-ch7-3-4-ray-reduction}
  \lean{Riemannian.HadamardModel.diffeomorphOfPole_of_rayGeodesic, Riemannian.HadamardModel.diffeomorphOfPole_of_rayGeodesic_coe, Riemannian.HadamardModel.hrays_of_rayGeodesic, Riemannian.HadamardModel.isGeodesic_rayCurve_of_christoffel, Riemannian.HadamardModel.rayCurve}
  \uses{lem:dc-ch7-3-4-pole-diffeo, lem:dc-ch7-3-4-pullback-metric}
  \leanok
  \dcref{ch7:3.4}
Let $M$ be a complete, simply connected Riemannian manifold, $p\in M$, and
  $f:T_pM\to M$ a smooth local diffeomorphism which is metric-expanding for its own
  pulled-back metric $f^*g$. If for every $v\in T_pM$ the radial line
  $\gamma_v(s)=s\,v$ is a geodesic of $f^*g$, then $f$ is a diffeomorphism
  $T_pM\to M$.
\end{lemma}
\begin{proof}

  \leanok
  \uses{lem:dc-ch7-3-4-pole-diffeo}
\sketch In the identity chart on $T_pM$, $\gamma_v(s)=sv$ has zero second derivative. Thus the geodesic condition is exactly $\Gamma^{f^*g}_{sv}(v,v)=0$; the remaining initial-data conditions are immediate, so the pole lemma applies.
\end{proof}

\begin{lemma}[$\exp_p$ is a local isometry, so the radial lines are $\exp_p^*g$-geodesics]

  \label{lem:dc-ch7-3-4-rays-are-geodesics}
  \lean{Riemannian.HadamardModel.expMap_rays_are_geodesics}
  \uses{lem:dc-ch7-3-4-pullback-metric, def:dc-ch7-2-2, def:dc-ch1-2-2, lem:dc-ch7-3-expmap-global}
  \leanok
  \dcref{ch7:3.4}
For $f=\exp_p$ and the pulled-back metric $\exp_p^*g$ on $T_pM$, every radial line
  $\gamma_v(s)=s\,v$ is a geodesic of $\exp_p^*g$.
\end{lemma}
\begin{proof}

  \uses{lem:dc-ch7-3-4-pullback-metric, lem:dc-ch7-3-expmap-global}
  \leanok
\sketch The pullback construction makes $\exp_p$ a local isometry. The image of $s\mapsto sv$ is the geodesic $s\mapsto\exp_p(sv)$, and local isometries reflect geodesics, so the radial line is a geodesic for $\exp_p^*g$.
\end{proof}

\begin{remark}[poles]

  \label{rem:dc-ch7-3-4}
  \lean{Riemannian.HadamardModel.expDiffeomorphOfPole_of_pole}
  \uses{lem:dc-ch7-3-4-pole-diffeo, lem:dc-ch7-3-4-ray-reduction, lem:dc-ch7-3-4-rays-are-geodesics, lem:dc-ch7-3-expmap-global}
  \leanok
  \dcref{ch7:3.4}
A point $p$ of a complete Riemannian manifold is a \emph{pole} if no point along any geodesic from $p$ is conjugate to $p$. Nonpositive sectional curvature implies that every point is a pole. If a complete simply connected manifold has a pole, then it is diffeomorphic to $\mathbb{R}^n$.
\end{remark}

\begin{lemma}[minimizing geodesics realize distance]
  \label{lem:dc-ch7-2-5-minimizing}
  \dcref{ch7:2.5}
  \uses{def:dc-ch3-3-2, def:dc-ch7-2-4}
  \lean{Riemannian.Geodesic.IsMinimizingGeodesicSegment.riemannianEDist_eq_pathELength,
    Riemannian.Geodesic.IsMinimizingGeodesicSegment.piecewiseRiemannianEDist_eq_pathELength}
  \leanok
  A minimizing geodesic segment $\gamma:[0,1]\to M$ satisfies
  $d(\gamma(0),\gamma(1))=\ell(\gamma)$.
\end{lemma}
\begin{proof}
  \leanok
  The geodesic itself is a competitor in the defining infimum, so the distance is at most
  its length.  Conversely, minimality says that its length is at most every piecewise
  differentiable competitor with the same endpoints; hence it is at most their infimum.
\end{proof}

\chapter{Spaces of Constant Curvature}
\label{chap:dc-spaces-of-constant-curvature}

Multiplying a Riemannian metric by $c>0$ divides its sectional curvature by
$c$. Thus constant-curvature metrics may be normalized to curvature $1$, $0$,
or $-1$. The corresponding simply connected models are $S^n$, $\mathbb R^n$,
and $H^n$.

\section{Theorem of Cartan on the determination of the metric by means of the curvature}

Let $M$ and $\tilde{M}$ be two Riemannian manifolds of dimension $n$ and let
$p\in M$ and $\tilde{p}\in\tilde{M}$. Choose a linear isometry
$i:T_p(M)\to T_{\tilde{p}}(\tilde{M})$. Let $V\subset M$ be a normal neighborhood
of $p$ such that $\exp_{\tilde{p}}$ is defined at $i\circ\exp_p^{-1}(V)$. Define a
mapping $f:V\to\tilde{M}$ by

$$
f(q)=\exp_{\tilde{p}}\circ\,i\circ\exp_p^{-1}(q),\qquad q\in V.
$$

For $q\in V$, let $\gamma:[0,t]\to M$
with $\gamma(0)=p$, $\gamma(t)=q$. Denote by $P_t$ the parallel transport along
$\gamma$ from $\gamma(0)$ to $\gamma(t)$. Define $\phi_t:T_q(M)\to T_{f(q)}(\tilde{M})$
by

$$
\phi_t(v)=\tilde{P}_t\circ i\circ P_t^{-1}(v),\qquad v\in T_q(M),
$$

where $\tilde{P}_t$ is the parallel transport along the normalized geodesic
$\tilde{\gamma}:[0,t]\to\tilde{M}$ given by $\tilde{\gamma}(0)=\tilde{p}$,
$\tilde{\gamma}'(0)=i(\gamma'(0))$. Denote by $R$ and $\tilde{R}$ the curvatures of
$M$ and $\tilde{M}$, respectively.

\begin{lemma}[Jacobi transfer in parallel frames]
  \label{lem:dc-ch8-2-1-jacobi-transfer}
  \lean{Riemannian.Jacobi.jacobiFrameTransfer}
  \uses{def:dc-ch5-2-1}
  \leanok
  Let $\gamma$ and $\tilde\gamma$ be geodesics with parallel orthonormal frames
  $(e_i)$ and $(\tilde e_i)$, where $e_n=\gamma'$ and
  $\tilde e_n=\tilde\gamma'$. Suppose that $J=\sum_i y_i e_i$ satisfies the Jacobi
  equation $\tfrac{D^2 J}{dt^2}+R(u',J)u'=0$ along $\gamma$, and that at the parameter
  $t$ the two frames have matching curvature coefficients,

  $$
  \langle R(e_n,e_i)e_n,e_j\rangle=\langle\tilde R(\tilde e_n,\tilde e_i)\tilde e_n,\tilde e_j\rangle,
  \qquad i,j=1,\dots,n.
  $$

  Then $\tilde J=\sum_i y_i\tilde e_i$ is a Jacobi field along $\tilde\gamma$.
\end{lemma}
\begin{proof}
  \leanok
\sketch
  Parallelism reduces the Jacobi equation to
  $y_j''+\sum_i a_{ij}y_i=0$, where
  $a_{ij}=\langle R(e_n,e_i)e_n,e_j\rangle$. The matching hypothesis gives the
  identical system for $(\tilde e_i)$, and reassembling its solution yields
  $D^2\tilde J/dt^2+\tilde R(\tilde\gamma',\tilde J)\tilde\gamma'=0$.
\end{proof}

\begin{lemma}[Jacobi transfer in constant curvature]
  \label{lem:dc-ch8-2-1-jacobi-transfer-const}
  \lean{Riemannian.Jacobi.jacobiFrameTransfer_isConstantCurvature}
  \uses{lem:dc-ch8-2-1-jacobi-transfer}
  \leanok
  Suppose $M$ and $\tilde M$ both have constant sectional curvature $K_0$ (the same value),
  and let $e_1,\dots,e_n$ resp.\ $\tilde e_1,\dots,\tilde e_n$ be parallel orthonormal frames
  along geodesics $\gamma$ resp.\ $\tilde\gamma$, with the velocities $u'=e_{n_0}$,
  $\tilde u'=\tilde e_{n_0}$ among the frame vectors. Then the curvature-matching hypothesis
  of \cref{lem:dc-ch8-2-1-jacobi-transfer},

  $$
  \langle R(e_i,u')u',e_j\rangle=\langle\tilde R(\tilde e_i,\tilde u')\tilde u',\tilde e_j\rangle,
  $$

  holds automatically: the constant-curvature identity and orthonormality make both sides
  $K_0(\delta_{ij}-\delta_{i n_0}\delta_{n_0 j})$. Hence a Jacobi field $J=\sum_i y_i e_i$ along
  $\gamma$ carries over, with the \emph{same} scalar coefficients, to a Jacobi field
  $\tilde J=\sum_i y_i\tilde e_i$ along $\tilde\gamma$.
\end{lemma}
\begin{proof}
  \leanok
\sketch
  The constant-curvature identity gives
  $\langle R(a,u')u',b\rangle=K_0(\langle u',u'\rangle\langle a,b\rangle-\langle a,u'\rangle\langle u',b\rangle)$
  Substituting
  $a=e_i$, $b=e_j$, $u'=e_{n_0}$ and using orthonormality
  ($\langle e_i,e_j\rangle=\delta_{ij}$, $\langle e_i,e_{n_0}\rangle=\delta_{i n_0}$,
  $\langle e_{n_0},e_{n_0}\rangle=1$) collapses each entry to
  $K_0(\delta_{ij}-\delta_{i n_0}\delta_{n_0 j})$. The two frame-coefficient
  matrices therefore coincide, discharging the hypothesis of
  \cref{lem:dc-ch8-2-1-jacobi-transfer}, whose conclusion is the claimed carry-over.
\end{proof}

\begin{lemma}[velocity-seeded parallel orthonormal frame]
  \label{lem:dc-ch8-2-1-velocity-frame}
  \lean{Riemannian.Jacobi.exists_velocitySeededParallelOrthoFrame}
  \uses{def:dc-ch5-2-1}
  \leanok
  Let $\gamma:[a,b]\to M$ be a geodesic whose initial velocity is a unit vector. Then there is a parallel
  orthonormal frame $e_1(t),\dots,e_n(t)$ along $\gamma$ (parallel: $De_i/dt=0$; orthonormal:
  $\langle e_i(t),e_j(t)\rangle=\delta_{ij}$ for all $t$) whose distinguished member is the
  velocity, $e_{n_0}(t)=\gamma'(t)$ for all $t\in[a,b]$.
\end{lemma}
\begin{proof}
  \leanok
\sketch
  Extend $\gamma'(a)$ to an orthonormal basis and parallel-transport that basis.
  Parallel transport preserves the metric, while uniqueness identifies the transported
  copy of $\gamma'(a)$ with the parallel field $\gamma'$.
\end{proof}

\begin{lemma}[the parallel orthonormal frame along a geodesic, intrinsically]
  \label{lem:dc-ch8-2-1-parallel-ortho-frame}
  \lean{Riemannian.Jacobi.exists_parallelOrthoFrameAlongOn}
  \uses{def:dc-ch5-2-1}
  \leanok
  Let $\gamma:[a,b]\to M$ be a geodesic and let $\{e^0_i\}$ be an orthonormal family in
  $T_{\gamma(a)}M$. Then there is a family $\{e_i(t)\}$ of vector fields along $\gamma$ with
  $e_i(a)=e^0_i$, each $e_i$ parallel along $\gamma$, and
  $\langle e_i(t),e_j(t)\rangle_{\gamma(t)}=\delta_{ij}$ for every $t\in[a,b]$.
\end{lemma}
\begin{proof}
  \leanok
\sketch
  Parallel-transport each $e^0_i$ along $\gamma$; this exists with any prescribed initial
  value. The pairing of two parallel fields is constant, so the Gram matrix remains
  equal to $(\delta_{ij})$.
\end{proof}

\begin{lemma}[Parallel frame transported by a linear isometry]
  \label{lem:dc-ch8-2-1-transported-frame}
  \lean{Riemannian.Jacobi.exists_transportedParallelOrthoFrame, Riemannian.Jacobi.exists_transportedParallelOrthoFrame_pair}
  \uses{lem:dc-ch8-2-1-parallel-ortho-frame}
  Let $i:T_pM\to T_{\tilde p}\tilde M$ be a linear isometry, let $\{e^0_j\}$ be an orthonormal
  family at $p=\gamma(a)$, and let $\tilde\gamma:[a,b]\to\tilde M$ be a geodesic with
  $\tilde\gamma(a)=\tilde p$. Then there is a family $\{\tilde e_j(t)\}$ along $\tilde\gamma$
  with $\tilde e_j(a)=i(e^0_j)$, each $\tilde e_j$ parallel along $\tilde\gamma$, and
  $\langle\tilde e_j(t),\tilde e_k(t)\rangle_{\tilde\gamma(t)}=\delta_{jk}$ for every
  $t\in[a,b]$. If $e_j$ is the frame from the corresponding initial data along
  $\gamma$, then $\tilde e_j(t)=\phi_t(e_j(t))$ for
  $\phi_t=\tilde P_t\circ i\circ P_t^{-1}$.
\end{lemma}
\begin{proof}
  \leanok
\sketch
  The family $(i(e_j^0))$ is orthonormal. Apply
  \cref{lem:dc-ch8-2-1-parallel-ortho-frame} and use uniqueness of parallel
  transport together with \cref{lem:dc-ch8-2-1-phi}.
\end{proof}

\begin{lemma}[uniqueness of parallel transport, and the transport map $P_t$]
  \label{lem:dc-ch8-2-1-parallel-transport-map}
  \lean{Riemannian.Jacobi.IsParallelFieldAlongOn.eqOn_of_initial, Riemannian.Jacobi.parallelTransportAlong, Riemannian.Jacobi.metricInner_parallelTransportAlong, Riemannian.Jacobi.parallelTransportAlongEquiv}
  \uses{def:dc-ch5-2-1}
  \leanok
  Let $\gamma:[a,b]\to M$ be a geodesic. Two fields parallel along $\gamma$ that agree at $a$
  agree on all of $[a,b]$. Consequently the assignment carrying $w_0\in T_{\gamma(a)}M$ to the
  value at $t$ of the parallel field along $\gamma$ seeded by $w_0$ is a well-defined map

  $$
  P_t:T_{\gamma(a)}M\longrightarrow T_{\gamma(t)}M,\qquad t\in[a,b],
  $$

  which is linear, satisfies $P_a=\mathrm{id}$, and is an isometry:
  $\langle P_tu,P_tw\rangle_{\gamma(t)}=\langle u,w\rangle_{\gamma(a)}$. In particular $P_t$ is
  injective, hence a linear isomorphism onto $T_{\gamma(t)}M$.
\end{lemma}
\begin{proof}
  \leanok
  \uses{def:dc-ch5-2-1}
\sketch
  Let $v,w$ be parallel along $\gamma$ with $v(a)=w(a)$, and set $d=v-w$. The parallel
  transport equation $\frac{Dd}{dt}=0$ is linear in $d$, so $d$ is parallel along $\gamma$,
  and $d(a)=0$. The pairing of two parallel fields along $\gamma$ is constant, so

  $$
  \langle d(t),d(t)\rangle_{\gamma(t)}=\langle d(a),d(a)\rangle_{\gamma(a)}=0
  \qquad(t\in[a,b]),
  $$

  and since $g_{\gamma(t)}$ is positive definite, $d(t)=0$; that is, $v(t)=w(t)$.

  Uniqueness makes $P_t$ well defined, and existence of a parallel field with prescribed value
  at $a$ makes it total. For linearity, let $v,w$ be the parallel fields seeded by $w_0,w_1$.
  Then $v+w$ is parallel and $(v+w)(a)=w_0+w_1$, so by uniqueness $v+w$ is the parallel field
  seeded by $w_0+w_1$; evaluating at $t$ gives $P_t(w_0+w_1)=P_tw_0+P_tw_1$, and homogeneity is
  the same argument applied to $rv$. $P_a=\mathrm{id}$ is the seeding condition. The isometry
  statement is once more the constancy of the pairing of two parallel fields, applied to the
  fields seeded by $u$ and $w$. Finally, if $P_tw_0=0$ then
  $\langle w_0,w_0\rangle_{\gamma(a)}=\langle P_tw_0,P_tw_0\rangle_{\gamma(t)}=0$, so $w_0=0$;
  an injective endomorphism of a finite-dimensional space is an isomorphism.
\end{proof}

\begin{lemma}[Parallel-transport conjugation]
  \label{lem:dc-ch8-2-1-phi}
  \lean{Riemannian.Jacobi.parallelTransportConjugate, Riemannian.Jacobi.parallelTransportConjugate_left, Riemannian.Jacobi.metricInner_parallelTransportConjugate, Riemannian.Jacobi.eq_parallelTransportConjugate_of_isParallelFieldAlongOn}
  \uses{lem:dc-ch8-2-1-parallel-transport-map}
  \leanok
  Let $\gamma:[a,b]\to M$ and $\tilde\gamma:[a,b]\to\tilde M$ be geodesics and let
  $i:T_{\gamma(a)}M\to T_{\tilde\gamma(a)}\tilde M$ be linear. For $t\in[a,b]$ put

  $$
  \phi_t=\tilde P_t\circ i\circ P_t^{-1}:
  T_{\gamma(t)}M\longrightarrow T_{\tilde\gamma(t)}\tilde M,
  $$

  with $P_t$, $\tilde P_t$ the parallel transports of
  \cref{lem:dc-ch8-2-1-parallel-transport-map}. Then $\phi_a=i$; if $i$ is a linear isometry
  then so is every $\phi_t$; and if $v$ is parallel along $\gamma$ and $\tilde v$ is parallel
  along $\tilde\gamma$ with $\tilde v(a)=i(v(a))$, then

  $$
  \tilde v(t)=\phi_t(v(t))\qquad\text{for every }t\in[a,b].
  $$

\end{lemma}
\begin{proof}
  \leanok
  \uses{lem:dc-ch8-2-1-parallel-transport-map}
\sketch
  $\phi_a=i$ because $P_a$ and $\tilde P_a$ are the identity. If $i$ is a linear isometry then
  $\phi_t$ is a composite of three isometries, $P_t^{-1}$ and $\tilde P_t$ being isometries by
  \cref{lem:dc-ch8-2-1-parallel-transport-map}; hence $\phi_t$ is one.

  For the last claim, $v$ is parallel with value $v(a)$ at $a$, so by the uniqueness clause of
  \cref{lem:dc-ch8-2-1-parallel-transport-map} it is the parallel field seeded by $v(a)$, i.e.
  $P_t(v(a))=v(t)$, whence $P_t^{-1}(v(t))=v(a)$. Therefore

  $$
  \phi_t(v(t))=\tilde P_t\big(i(v(a))\big)=\tilde P_t\big(\tilde v(a)\big)=\tilde v(t),
  $$

  the last equality because $\tilde v$ is parallel along $\tilde\gamma$ with value $\tilde v(a)$
  at $a$, hence is the parallel field seeded by $\tilde v(a)$.
\end{proof}

\begin{lemma}[the chart frame coefficient is the intrinsic curvature form]
  \label{lem:dc-ch8-2-1-curvature-bridge}
  \lean{Riemannian.Jacobi.chartMetricInner_chartCurvatureEndo_eq_curvatureFormAt, Riemannian.Jacobi.chartFrameRealize_tangentCoordChange, Riemannian.Jacobi.chartMetricInner_chartCurvatureEndo_chartVectorRep_eq_curvatureFormAt}
  \leanok
  For $v,a,b\in T_qM$, the frame coefficient $\langle\mathcal R(a,v)v,b\rangle$ computed in a
  chart --- the form in which the Jacobi equation is actually run --- equals the intrinsic
  curvature form $R(v,a,v,b)$ evaluated on $v,a,v,b$ themselves. No curvature hypothesis is
  assumed.
\end{lemma}
\begin{proof}
  \leanok
\sketch
  Pairing the chart Jacobi operator against a chart vector gives
  $\langle\mathcal R(a,v)v,b\rangle$. The manifold--chart curvature bridge identifies
  $R(\hat a,\hat v,\hat v,\hat b)$ --- the intrinsic form on the chart-frame realizations
  $\hat\cdot$ of the chart vectors --- with the \emph{negative} of that quantity, and
  antisymmetry of $R$ in its first pair rewrites $R(\hat a,\hat v,\hat v,\hat b)$ as
  $-R(\hat v,\hat a,\hat v,\hat b)$; the two signs cancel. Finally, realizing the chart
  reading of a tangent vector in the chart frame at its own foot returns the vector
  (the readback inverse to the chart reading), so for intrinsic $v,a,b$ no realizations
  remain.
\end{proof}

\begin{lemma}[curvature intertwining in frame coefficients]
  \label{lem:dc-ch8-2-1-hmatch-transfer}
  \lean{Riemannian.Jacobi.chartMetricInner_chartCurvatureEndo_transfer_of_curvatureFormAt}
  \uses{lem:dc-ch8-2-1-curvature-bridge}
  \leanok
  Let $\varphi:T_qM\to T_{\tilde q}\tilde M$ be any map under which the curvature forms
  correspond,

  $$
  \langle R(x,y)z,w\rangle=\big\langle\tilde R(\varphi x,\varphi y)\varphi z,\varphi w\big\rangle
  \qquad\text{for all }x,y,z,w\in T_qM ,
  $$

  which is the hypothesis of \cref{thm:dc-ch8-2-1} with $\varphi=\phi_t$. Then the chart frame
  coefficients of $M$ and $\tilde M$ agree: for all $v,a,b\in T_qM$,

  $$
  \langle\mathcal R(a,v)v,b\rangle=\big\langle\tilde{\mathcal R}(\varphi a,\varphi v)\varphi v,
  \varphi b\big\rangle .
  $$

  Taking $v=\gamma'$ and $a,b$ frame vectors, this is the intrinsic content of the matching
  hypothesis of \cref{lem:dc-ch8-2-1-jacobi-transfer} and of
  \cref{lem:dc-ch8-2-1-jacobi-norm-transfer-general} --- the single point at which the
  curvature hypothesis of \cref{thm:dc-ch8-2-1} is consumed.

  It is not yet that matching hypothesis itself. Those lemmas read
  their frames in a fixed chart along a coordinate curve, so identifying the two demands three
  further things, none of them about curvature: that $\varphi(\gamma'(t))=\tilde\gamma'(t)$
  (which needs $\tilde\gamma'(a)=i(\gamma'(a))$); that its frame vectors are the chart
  readings of the present intrinsic ones under the present $\varphi$; and that its
  parallelism hypothesis, stated in one fixed chart, follows from parallelism along
  $\gamma$ in the sense of \cref{lem:dc-ch8-2-1-parallel-ortho-frame}, which is chart-local.
  That interface is the unbuilt residual of \cref{thm:dc-ch8-2-1}.

  Note $\varphi$ is arbitrary: no linearity, continuity or isometry is used, and no
  parallel-transport theory enters --- $\varphi$ is data, supplied by \cref{thm:dc-ch8-2-1}'s
  own hypothesis as $\phi_t$. The additional structure of $\phi_t$ is what makes the
  transported frame orthonormal (\cref{lem:dc-ch8-2-1-transported-frame}), which is a separate
  requirement.
\end{lemma}
\begin{proof}
  \leanok
\sketch
  Rewrite each side by \cref{lem:dc-ch8-2-1-curvature-bridge} into the intrinsic curvature
  forms $R(v,a,v,b)$ and $\tilde R(\varphi v,\varphi a,\varphi v,\varphi b)$, and apply the
  hypothesis at $(v,a,v,b)$.
\end{proof}

\begin{lemma}[the Jacobi field read in a parallel orthonormal frame]
  \label{lem:dc-ch8-2-1-frame-coeff-bridge}
  \lean{Riemannian.Jacobi.isJacobiPairOn_of_isJacobiFieldOn}
  \uses{def:dc-ch5-2-1, lem:dc-ch8-2-1-parallel-ortho-frame}
  \leanok
  Let $\gamma$ be a geodesic of $M$, read in a fixed chart as the coordinate curve $u$
  staying in the chart domain, and let the pair $(J,DJ)$ satisfy the Jacobi pair system
  along $u$ on a parameter window $[a',b']$,

  $$
  \tfrac{D}{dt}J=DJ,\qquad \tfrac{D}{dt}DJ=-\mathcal R(J,u')u' ,
  $$

  where $\langle\cdot,\cdot\rangle$ and $\mathcal R$ denote the metric and curvature read in
  that chart. Let $e_1(t),\dots,e_n(t)$ ($n=\dim M$) be a parallel orthonormal frame along
  $u$ in the same chart. For a parameter interval $[a,b]\subset(a',b')$ set

  $$
  F_i(t)=\langle J(t),e_i(t)\rangle,\qquad V_i(t)=\langle DJ(t),e_i(t)\rangle .
  $$

  Then on $[a,b]$ the pair $(F,V)$ solves the first-order system

  $$
  F'=V,\qquad V'=-A(t)F,\qquad A(t)_{ij}=a_{ij}(t)=\langle\mathcal R(e_i,u')u',e_j\rangle ,
  $$

  i.e.\ $F_j''+\sum_i a_{ij}F_i=0$.
\end{lemma}
\begin{proof}
  \leanok
\sketch
  Metric compatibility and $De_j/dt=0$ give
  $\frac d{dt}\langle X,e_j\rangle=\langle DX/dt,e_j\rangle$.
  Apply this first to $J$ and then to $DJ$. Expanding $J=\sum_iF_ie_i$ in the
  second identity gives $F'=V$ and $V'=-AF$.
\end{proof}

\begin{lemma}[transfer of the Jacobi norm under matching curvature coefficients]
  \label{lem:dc-ch8-2-1-jacobi-norm-transfer-general}
  \lean{Riemannian.Jacobi.chartMetricInner_transfer_of_jacobiCoef_match}
  \uses{lem:dc-ch8-2-1-frame-coeff-bridge, lem:dc-ch8-2-1-hmatch-transfer, def:dc-ch5-2-1}
  \leanok
  Let $\gamma$ in $M$ and $\tilde\gamma$ in $\tilde M$ ($\dim M=\dim\tilde M=n$) be geodesics
  read in fixed charts as the coordinate curves $u$ and $\tilde u$, each staying in its chart
  domain, let $(J,DJ)$ and $(\tilde J,D\tilde J)$ satisfy the Jacobi pair systems along $u$
  and $\tilde u$ on a window $[a',b']$, and let $e$ and $\tilde e$ be parallel orthonormal
  frames along $u$ and $\tilde u$ in those charts; as in
  \cref{lem:dc-ch8-2-1-frame-coeff-bridge}, $\langle\cdot,\cdot\rangle$ and $\mathcal R$
  denote the metric and curvature read in those charts. Let $[a,b]\subset(a',b')$, and
  suppose that on $[a,b]$ the chart curvature operator $w\mapsto\mathcal R(w,u')u'$ along
  $u$, the frame vectors $e_i$, and the chart metric coefficients along $u$ are all
  continuous, and that the curvature coefficients agree there,

  $$
  \langle\mathcal R(e_i,u')u',e_j\rangle
  =\big\langle\tilde{\mathcal R}(\tilde e_i,\tilde u')\tilde u',\tilde e_j\big\rangle
  \qquad\text{for all }i,j\text{ and all }t\in[a,b],
  $$

  and that the two fields have the same frame data at $a$,
  $\langle J(a),e_i(a)\rangle=\langle\tilde J(a),\tilde e_i(a)\rangle$ and
  $\langle DJ(a),e_i(a)\rangle=\langle D\tilde J(a),\tilde e_i(a)\rangle$ for all $i$.
  Then

  $$
  |\tilde J(t)|^2=|J(t)|^2\qquad\text{for all }t\in[a,b] .
  $$

\end{lemma}
\begin{proof}
  \leanok
\sketch
  By \cref{lem:dc-ch8-2-1-frame-coeff-bridge}, both coefficient pairs solve
  $F'=V$, $V'=-A(t)F$ with the same continuous matrix $A$. Their initial data agree,
  so uniqueness gives $F=\tilde F$.

  Finally, orthonormality of each frame reads the norm off the coefficients: expanding
  $J=\sum_i F_ie_i$ and $\tilde J=\sum_i\tilde F_i\tilde e_i$,

  $$
  |J(t)|^2=\sum_i F_i(t)^2,\qquad |\tilde J(t)|^2=\sum_i\tilde F_i(t)^2 ,
  $$

  and these agree because $F=\tilde F$ on $[a,b]$.
\end{proof}

\begin{lemma}[the transported frame and velocity, read in a fixed chart]
  \label{lem:dc-ch8-2-1-frame-chart-interface}
  \lean{Riemannian.Jacobi.jacobiCoef_match_of_curvatureFormAt, Riemannian.Jacobi.jacobiCoef_match_of_curvatureFormAt_parallelTransportConjugate}
  \uses{lem:dc-ch8-2-1-hmatch-transfer, lem:dc-ch8-2-1-transported-frame, lem:dc-ch8-2-1-parallel-transport-map, lem:dc-ch8-2-1-phi}
  \leanok
  Let $\gamma:[a,b]\to M$ and $\tilde\gamma:[a,b]\to\tilde M$ be geodesics, each of which
  \emph{stays in the source of one fixed chart}, read there as the coordinate curves $u$ and
  $\tilde u$. Let $e_1,\dots,e_n$ be parallel fields along $\gamma$ and $\tilde e_1,\dots,
  \tilde e_n$ parallel fields along $\tilde\gamma$ with matched seeds
  $\tilde e_j(a)=i(e_j(a))$, and suppose $\tilde\gamma'(a)=i(\gamma'(a))$; read all of them
  in those same charts. Then the curvature hypothesis of \cref{thm:dc-ch8-2-1} implies the
  matching hypothesis of \cref{lem:dc-ch8-2-1-jacobi-norm-transfer-general} in the form that
  lemma requires it, namely between the \emph{chart} frame coefficients:

  $$
  \langle\mathcal R(e_i,u')u',e_j\rangle
  =\big\langle\tilde{\mathcal R}(\tilde e_i,\tilde u')\tilde u',\tilde e_j\big\rangle
  \qquad\text{for all }i,j\text{ and all }t\in[a,b] .
  $$

  More generally, the same conclusion holds with $\phi_t$ replaced by any family of
  maps $T_{\gamma(t)}M\to T_{\tilde\gamma(t)}\tilde M$ under which the curvature forms
  correspond and which carries $\gamma'(t)$ to $\tilde\gamma'(t)$ and $e_j(t)$ to
  $\tilde e_j(t)$: no linearity, continuity or isometry of the family is used, and no
  parallel-transport theory enters. The frames need not contain the velocity or be
  orthonormal for this coefficient identity.
\end{lemma}
\begin{proof}
  \leanok
\sketch
  Since the velocity fields are parallel,

  $$
  \phi_t(\gamma'(t))=\tilde P_t\big(i(P_t^{-1}\gamma'(t))\big)=\tilde P_t\big(i(\gamma'(a))\big)
  =\tilde P_t(\tilde\gamma'(a))=\tilde\gamma'(t) ,
  $$

  using $\tilde\gamma'(a)=i(\gamma'(a))$. Also
  $\tilde e_j(t)=\phi_t(e_j(t))$ by \cref{lem:dc-ch8-2-1-transported-frame}.
  Apply \cref{lem:dc-ch8-2-1-hmatch-transfer} with $\varphi=\phi_t$ at
  $(v,a,b)=(\gamma'(t),e_i(t),e_j(t))$ gives the matching of the chart frame coefficients of
  the \emph{chart readings} of $\gamma'(t),e_i(t),e_j(t)$ and of their images under $\phi_t$,
  which gives the required chart coefficient identity.
\end{proof}

\begin{lemma}[the Cartan Jacobi norm transfer in variable curvature, intrinsically]
  \label{lem:dc-ch8-2-1-jacobi-norm-transfer-manifold}
  \lean{Riemannian.Jacobi.metricInner_jacobiField_transfer_of_curvatureFormAt}
  \uses{lem:dc-ch8-2-1-frame-chart-interface, lem:dc-ch8-2-1-jacobi-norm-transfer-general, lem:dc-ch8-2-1-parallel-ortho-frame, def:dc-ch5-2-1}
  \leanok
  Let $\gamma,\tilde\gamma:[a',b']\to M,\tilde M$ be geodesics, each staying in the source of
  one fixed chart; let $e_1,\dots,e_n$ and $\tilde e_1,\dots,\tilde e_n$ be parallel
  \emph{orthonormal} frames along them; and let $\phi_t$ carry $\gamma'(t)$ to
  $\tilde\gamma'(t)$, $e_j(t)$ to $\tilde e_j(t)$, and the curvature form of $M$ at
  $\gamma(t)$ to that of $\tilde M$ at $\tilde\gamma(t)$. Let $J$ be
  a Jacobi field along $\gamma$ and $\tilde J$ one along $\tilde\gamma$, and let
  $[a,b]\subset(a',b')$. If the two have the same frame data at $a$, then

  $$
  |\tilde J(t)|=|J(t)|\qquad\text{for all }t\in[a,b] .
  $$

  No constancy or sign condition on the curvature is required.
\end{lemma}
\begin{proof}
  \leanok
\sketch
  Read the fields and frames in the fixed charts. Their chart readings preserve norms and
  satisfy the hypotheses of \cref{lem:dc-ch8-2-1-jacobi-norm-transfer-general}; coefficient
  matching is \cref{lem:dc-ch8-2-1-frame-chart-interface}. Translate the resulting norm
  equality back to the tangent bundles.
\end{proof}

\begin{lemma}[time reversal of the Jacobi system]
  \label{lem:dc-ch8-2-1-jacobi-reversal}
  \lean{Riemannian.Jacobi.IsJacobiFieldAlongOn.comp_neg}
  \uses{def:dc-ch5-2-1}
  \leanok
  Let $\gamma$ be a geodesic on $[a,b]$ and let $(J,DJ)$ be a Jacobi pair along it. Then
  $t\mapsto\gamma(-t)$ is a geodesic on $[-b,-a]$, and

  $$
  J^-(t)=J(-t),\qquad DJ^-(t)=-DJ(-t)
  $$

  is a Jacobi pair along it.
\end{lemma}
\begin{proof}
  \leanok
\sketch
  Substitute $t\mapsto-t$ into the chart Jacobi system. Bilinearity supplies the signs in
  the connection terms, while the two reversed velocity slots leave the curvature term
  unchanged.
\end{proof}

\begin{lemma}[backward uniqueness of Jacobi fields]
  \label{lem:dc-ch8-2-1-jacobi-backward-uniqueness}
  \lean{Riemannian.Jacobi.IsJacobiFieldAlongOn.eqOn_zero_of_right}
  \uses{lem:dc-ch8-2-1-jacobi-reversal, def:dc-ch5-2-1}
  \leanok
  A Jacobi field along a geodesic on $[a,b]$ with $J(b)=0$ and $DJ(b)=0$ vanishes identically
  on $[a,b]$.
\end{lemma}
\begin{proof}
  \leanok
\sketch
  Reverse time using \cref{lem:dc-ch8-2-1-jacobi-reversal}. The reversed pair has zero
  initial data at $-b$, so forward uniqueness makes it, and hence the original pair, zero.
\end{proof}

\begin{lemma}[a Jacobi field with prescribed data at an interior time]
  \label{lem:dc-ch8-2-1-jacobi-interior-data}
  \lean{Riemannian.Jacobi.exists_isJacobiFieldAlongOn_at}
  \uses{lem:dc-ch8-2-1-jacobi-backward-uniqueness, def:dc-ch5-2-1}
  \leanok
  Let $\gamma$ be a geodesic on $[a,b]$, let $t_0\in[a,b]$, and let $J_0,DJ_0\in
  T_{\gamma(t_0)}M$. Then there is a Jacobi field $(J,DJ)$ on $[a,b]$ with
  $J(t_0)=J_0$ and $DJ(t_0)=DJ_0$.
\end{lemma}
\begin{proof}
  \leanok
\sketch
  Send initial data at $a$ to the resulting data at $t_0$. This linear map is injective by
  \cref{lem:dc-ch8-2-1-jacobi-backward-uniqueness}, and its domain and codomain both have
  dimension $2n$. It is therefore surjective.
\end{proof}

\begin{lemma}[from the norm transfer to the local isometry, in variable curvature]
  \label{lem:dc-ch8-2-1-exp-norm-transfer-general}
  \lean{Riemannian.Jacobi.metricInner_mfderiv_eq_of_semiconjugacy_of_curvatureFormAt}
  \uses{lem:dc-ch8-2-1-jacobi-norm-transfer-manifold, cor:dc-ch5-2-5, lem:dc-ch8-2-1-jacobi-interior-data, lem:dc-ch8-2-1-exp-norm-transfer, lem:dc-ch8-2-1-polarization, lem:dc-ch8-2-1-exp-equiv, lem:dc-ch8-2-1-mfderiv-bridge, lem:dc-ch8-2-2-semiconjugacy, lem:dc-ch8-2-1-no-conjugate}
  \leanok
  Let $M,\tilde M$ be complete, $p\in M$, $\tilde p\in\tilde M$, let $i:T_pM\to
  T_{\tilde p}\tilde M$ be a linear isometry, and let $v\in T_pM$. Write $\gamma=\gamma_v$
  and $\tilde\gamma=\gamma_{i(v)}$ for the two geodesics, and $q=\exp_p(v)$. Suppose:
  \begin{itemize}
    \item \emph{(single chart)} for some window $[a',b']$ with $a'<0<1<b'$, each of $\gamma$,
      $\tilde\gamma$ stays in the source of \emph{one} chart on $[a',b']$;
    \item \emph{(frames)} $e_1,\dots,e_n$ and $\tilde e_1,\dots,\tilde e_n$ are parallel
      \emph{orthonormal} frames along $\gamma$, $\tilde\gamma$ on $[a',b']$;
    \item \emph{(curvature intertwining)} $\phi_t$ carries $\gamma'(t)$ to
      $\tilde\gamma'(t)$, $e_j(t)$ to $\tilde e_j(t)$, and the curvature form of $M$ at
      $\gamma(t)$ to that of $\tilde M$ at $\tilde\gamma(t)$, for $t\in[a',b']$; and
      $\phi_0=i$;
    \item \emph{(no conjugacy)} $1$ is not a conjugate parameter along $\gamma$ or along
      $\tilde\gamma$.
  \end{itemize}
  Let $f$ be \emph{any} map differentiable at $q$ that satisfies the semiconjugacy
  $f\circ\exp_p=\exp_{\tilde p}\circ\,i$ near $v$. Then $f$ preserves the metric at $q$:

  $$
  \langle u,u'\rangle_q=\langle df_q(u),df_q(u')\rangle_{f(q)}\qquad\text{for all }u,u'\in T_qM .
  $$

  This is the variable-curvature counterpart of \cref{lem:dc-ch8-2-2-semiconjugacy}.
\end{lemma}
\begin{proof}
  \leanok
\sketch
  Since $1$ is not conjugate along $\gamma$, $d(\exp_p)_v$ is onto
  (\cref{lem:dc-ch8-2-1-exp-equiv}), so it suffices to prove the claim on vectors
  $d(\exp_p)_v(Z)$. Realize each side as the endpoint of a Jacobi field vanishing at $0$:
  by \cref{cor:dc-ch5-2-5}, $d(\exp_p)_v(Z)=J(1)$ where $J(0)=0$, $DJ(0)=Z$, and likewise
  $d(\exp_{\tilde p})_{i(v)}(i(Z))=\tilde J(1)$ with
  $\tilde J(0)=0$, $D\tilde J(0)=i(Z)$. The fields on the full window are supplied by
  \cref{lem:dc-ch8-2-1-jacobi-interior-data}.

  The two fields have matching frame data at $0$. For the value this is trivial, both being
  $0$. For the covariant derivative, $\tilde e_k(0)=\phi_0(e_k(0))=i(e_k(0))$, so

  $$
  \langle i(Z),\tilde e_k(0)\rangle_{\tilde p}=\langle i(Z),i(e_k(0))\rangle_{\tilde p}
    =\langle Z,e_k(0)\rangle_p
  $$

  because $i$ is a linear isometry.

  So \cref{lem:dc-ch8-2-1-jacobi-norm-transfer-manifold} applies and gives
  $|\tilde J(1)|=|J(1)|$, i.e. $|d(\exp_{\tilde p})_{i(v)}(i(Z))|=|d(\exp_p)_v(Z)|$. The
  chain rule along the semiconjugacy identifies $df_q(d(\exp_p)_v(Z))$ with
  $d(\exp_{\tilde p})_{i(v)}(i(Z))$ (\cref{lem:dc-ch8-2-1-mfderiv-bridge}), which turns that
  into $|df_q(w)|=|w|$ for every $w\in T_qM$; polarization
  (\cref{lem:dc-ch8-2-1-polarization}) upgrades the norm identity to the full inner product.
\end{proof}

\begin{lemma}[an orthonormal basis for the metric at a point]
  \label{lem:dc-ch8-2-1-metric-ortho-basis}
  \lean{Riemannian.Jacobi.exists_metricOrthonormalBasis}
  \leanok
  Let $g$ be a Riemannian metric on $M$ and $q\in M$. Then $T_qM$ has an orthonormal basis for
  $\langle\cdot,\cdot\rangle_q$.
\end{lemma}
\begin{proof}
  \leanok
\sketch
  $\langle\cdot,\cdot\rangle_q$ is a symmetric bilinear form on the finite-dimensional space
  $T_qM$, positive definite by the definition of a Riemannian metric. Diagonalize it and
  normalize.
\end{proof}

\begin{lemma}[Cartan transport and matched frames]
  \label{lem:dc-ch8-2-1-cartan-frame-data}
  \lean{Riemannian.Jacobi.cartanTransportZero, Riemannian.Jacobi.cartanSeed, Riemannian.Jacobi.metricInner_cartanSeed, Riemannian.Jacobi.cartanPhi, Riemannian.Jacobi.cartanPhi_zero, Riemannian.Jacobi.cartanPhi_velocity, Riemannian.Jacobi.cartanPhi_isLinear, Riemannian.Jacobi.metricInner_cartanPhi, Riemannian.Jacobi.exists_cartanFrameData}
  \uses{lem:dc-ch8-2-1-phi, lem:dc-ch8-2-1-transported-frame, lem:dc-ch8-2-1-parallel-transport-map, lem:dc-ch8-2-1-metric-ortho-basis, lem:dc-ch8-2-1-velocity-frame}
  \leanok
  Let $M,\tilde M$ be complete, $p\in M$, $\tilde p\in\tilde M$, let $i:T_pM\to T_{\tilde p}\tilde
  M$ be a linear isometry, let $v\in T_pM$, and let $[a',b']$ be a window with $a'<0<1<b'$. Write
  $\gamma=\gamma_v$ and $\tilde\gamma=\gamma_{i(v)}$.

  Define $j=\tilde P_0^{-1}\circ i\circ P_0$ and
  $\phi^{\mathrm{Cartan}}_t=\tilde P_t\circ j\circ P_t^{-1}$. Then for
  $t\in[a',b']$:
  \begin{itemize}
    \item $\phi^{\mathrm{Cartan}}_0=i$;
    \item $\phi^{\mathrm{Cartan}}_t(\gamma'(t))=\tilde\gamma'(t)$;
    \item $\phi^{\mathrm{Cartan}}_t$ is linear, and preserves the metric:
      $\langle\phi^{\mathrm{Cartan}}_tu,\phi^{\mathrm{Cartan}}_tw\rangle_{\tilde\gamma(t)}
      =\langle u,w\rangle_{\gamma(t)}$.
  \end{itemize}
  Moreover there exist parallel \emph{orthonormal} frames $e_1,\dots,e_n$ along $\gamma$ and
  $\tilde e_1,\dots,\tilde e_n$ along $\tilde\gamma$ on $[a',b']$ with
  $\phi^{\mathrm{Cartan}}_t(e_j(t))=\tilde e_j(t)$ for every $t\in[a',b']$.
\end{lemma}
\begin{proof}
  \leanok
  \uses{lem:dc-ch8-2-1-phi,lem:dc-ch8-2-1-transported-frame,
    lem:dc-ch8-2-1-metric-ortho-basis,lem:dc-ch8-2-1-velocity-frame}
\sketch
  Use the conjugated seed

  $$
  j=\tilde P_0^{-1}\circ i\circ P_0 ,
  $$

  Then $\phi_0=i$, and $j$ is a linear isometry because parallel transport and $i$
  preserve the metric.

  With $j$ in hand, seed the frames by an orthonormal basis at $\gamma(a')$
  (\cref{lem:dc-ch8-2-1-metric-ortho-basis}) and apply
  \cref{lem:dc-ch8-2-1-transported-frame}. The frame and velocity identities follow from
  uniqueness of parallel transport via \cref{lem:dc-ch8-2-1-phi}.
\end{proof}

\begin{lemma}[Cartan transfer with nonconjugacy discharged]
  \label{lem:dc-ch8-2-1-exp-norm-transfer-general-discharged}
  \lean{Riemannian.Jacobi.metricInner_mfderiv_eq_of_semiconjugacy_of_curvatureFormAt_of_isLocalDiffeomorphAt}
  \uses{lem:dc-ch8-2-1-exp-norm-transfer-general, lem:dc-ch8-2-1-cartan-frame-data, lem:dc-ch8-2-1-no-conjugate-normal}
  \leanok
  Same setting and same conclusion as \cref{lem:dc-ch8-2-1-exp-norm-transfer-general} --- $f$
  preserves the metric at $q=\exp_p(v)$ --- but with the frames, $\phi$, and \emph{both}
  no-conjugacy clauses discharged rather than assumed. In place of them it assumes only that
  $\exp_p$ and $\exp_{\tilde p}$ are local diffeomorphisms at $v$ and $i(v)$: the
  normal-neighborhood hypothesis. The remaining assumptions are the single-chart clauses
  and the curvature match for the Cartan transport of
  \cref{lem:dc-ch8-2-1-cartan-frame-data}.
\end{lemma}
\begin{proof}
  \leanok
  \uses{lem:dc-ch8-2-1-exp-norm-transfer-general,lem:dc-ch8-2-1-cartan-frame-data,
    lem:dc-ch8-2-1-no-conjugate-normal}
\sketch
  Obtain the frames and $\phi$ from \cref{lem:dc-ch8-2-1-cartan-frame-data} and feed the curvature
  hypothesis the produced $\phi$, which satisfies its three antecedents by construction. Obtain
  the two no-conjugacy clauses from \cref{lem:dc-ch8-2-1-no-conjugate-normal} applied on each
  side. Then apply \cref{lem:dc-ch8-2-1-exp-norm-transfer-general}.
\end{proof}

\begin{lemma}[from pointwise metric preservation to a local isometry on $V$]
  \label{lem:dc-ch8-2-1-local-isometry-packaging}
  \lean{Riemannian.DCIsLocalIsometryAt}
  \uses{lem:dc-ch8-2-1-exp-norm-transfer-general-discharged, lem:dc-ch8-2-1-single-chart, lem:dc-ch8-2-1-exp-diff-zero, lem:dc-ch8-2-1-no-conjugate-zero, lem:dc-ch8-2-1-window-transfer}
  \notready
  \notready
  The preceding result is pointwise at $q=\exp_p(v)$; this lemma packages it into a local
  isometry on the whole normal neighborhood. The remaining requirements are a single chart
  containing each relevant geodesic and the assembly of the pointwise differential identities.
  The base point satisfies $df_p=i$ by \cref{lem:dc-ch8-2-1-exp-diff-zero}.
\end{lemma}

\begin{lemma}[widening a single-chart time window]
  \label{lem:dc-ch8-2-1-window}
  \lean{Riemannian.Jacobi.exists_window_of_mem_chartAt_source}
  \leanok
  Let $\gamma:\mathbb R\to M$ be continuous and suppose $\gamma(t)$ lies in the source of a
  \emph{single} chart for every $t\in[0,1]$. Then it still does on some $[a',b']$ with $a'<0<1<b'$.

  This only widens a given single-chart hypothesis.
\end{lemma}
\begin{proof}
  \leanok
\sketch
  The preimage of the chart source is open and contains compact $[0,1]$, hence contains
  $(-\delta,1+\delta)$ for some $\delta>0$.
\end{proof}

\begin{lemma}[a geodesic of a normal neighborhood, confined to one chart]
  \label{lem:dc-ch8-2-1-single-chart}
  \uses{lem:dc-ch8-2-1-frame-chart-interface, lem:dc-ch8-2-1-jacobi-norm-transfer-general, lem:dc-ch8-2-1-jacobi-norm-transfer-manifold, lem:dc-ch8-2-1-chart-partition, lem:dc-ch8-2-1-radial-chart}
  \notready
  \notready
  The conclusion needed by the transfer is the single-chart condition on each geodesic over
  $[0,1]$ (or a finite chart-chain with compatible ODE data). The existence of such a chart
  uniformly for the normal neighborhood is the remaining obligation.
\end{lemma}

\begin{lemma}[the constant-curvature normal Jacobi field, any sign of $K_0$]
  \label{lem:dc-ch8-2-1-jacobi-const-solution}
  \lean{Riemannian.Jacobi.isJacobiFieldAlongOn_of_constantCurvature}
  \uses{def:dc-ch5-2-1}
  \leanok
  Let $M$ have constant sectional curvature $K_0$ (any sign), let $\gamma$ be a
  unit-speed geodesic, let $w$ be a \emph{parallel} field along $\gamma$ \emph{normal}
  to $\gamma'$, and let $h$ be any scalar solution of $h''+K_0 h=0$ (given by its
  derivative data $h'=Dh$, $Dh'=-K_0 h$). Then $J(t)=h(t)\,w(t)$ is a Jacobi field
  along $\gamma$, with covariant derivative $DJ(t)=Dh(t)\,w(t)$. This covers
  all curvature signs, with
  $h=\sin(t\sqrt{K_0})/\sqrt{K_0}$, $h=t$, $h=\sinh(t\sqrt{-K_0})/\sqrt{-K_0}$).
\end{lemma}
\begin{proof}
  \leanok
\sketch
  Since $w$ is parallel, $\tfrac{D}{dt}(h w)=h'\,w$ and
  $\tfrac{D}{dt}(Dh\,w)=Dh'\,w=-K_0 h\,w$; and by the constant-curvature model form
  $R(\gamma',w)\gamma'=K_0 w$ (using $|\gamma'|^2=1$ and
  $\langle w,\gamma'\rangle=0$) together with homogeneity in the middle slot,
  $R(\gamma',h w)\gamma'=h K_0\,w$. Hence
  $\tfrac{D^2}{dt^2}(h w)+R(\gamma',h w)\gamma'=-K_0 h\,w+K_0 h\,w=0$: the pair
  $(h w, Dh\,w)$ solves the Jacobi system. This generalizes the special case
  $K_0>0$, $h=\sin(t\sqrt{K_0})/\sqrt{K_0}$ used in
  \cref{ex:dc-ch5-3-3} to an arbitrary scalar solution.
\end{proof}

\begin{lemma}[the norm of a constant-curvature Jacobi field is manifold-independent]
  \label{lem:dc-ch8-2-1-jacobi-norm-transfer}
  \lean{Riemannian.Jacobi.metricInner_jacobiField_eq_of_constantCurvature_of_speedSq, Riemannian.Jacobi.metricInner_jacobiField_transfer_of_constantCurvature_of_speedSq}
  \uses{lem:dc-ch8-2-1-jacobi-const-solution, cor:dc-ch5-2-5, def:dc-ch5-2-1}
  \leanok
  Let $\gamma:[0,\ell]\to M$ be a geodesic of squared speed $c=|\gamma'|^2\neq0$ on a
  manifold of constant sectional curvature $K_0$, let $J$ be the Jacobi field with $J(0)=0$
  and covariant derivative $DJ$, and let $h$ solve $h''+K_0c\,h=0$, $h(0)=0$, $h'(0)=1$.
  Writing $Z=DJ(0)$ and $a=\langle Z,\gamma'(0)\rangle$,

  $$
  |J(t)|^2 = \frac{a^2}{c}\,t^2 + h(t)^2\Big(|Z|^2-\frac{a^2}{c}\Big).
  $$

  The right-hand side involves \emph{no} manifold data beyond $K_0$ and $c$ (through $h$),
  the time $t$, and the two scalar invariants $a=\langle Z,\gamma'(0)\rangle$ and $|Z|^2$.
  Consequently, Jacobi fields in any two spaces of curvature $K_0$ and equal speed $c$
  have equal norms when these two invariants agree. For $c=1$ this is
  $|J(t)|^2=t^2a^2+h(t)^2(|Z|^2-a^2)$.
\end{lemma}
\begin{proof}
  \leanok
\sketch
  Decompose $Z=\frac{a}{c}\,\gamma'(0)+Z_\perp$ with $Z_\perp=Z-\frac{a}{c}\gamma'(0)$; then
  $\langle Z_\perp,\gamma'(0)\rangle=a-\frac{a}{c}\cdot c=0$ and
  $|Z_\perp|^2=|Z|^2-\frac{a^2}{c}$. On a space of constant curvature $K_0$ a field normal to
  $\gamma'$ satisfies $R(\gamma',w)\gamma'=K_0|\gamma'|^2w=K_0c\,w$, so the normal Jacobi
  equation is the scalar equation $h''+K_0c\,h=0$. The tangential Jacobi field
  $\frac{a}{c}\,t\,\gamma'(t)$ and the normal Jacobi field
  $h(t)\,w(t)$, with $w$ the parallel transport of $Z_\perp$
  (\cref{lem:dc-ch8-2-1-jacobi-const-solution}), both vanish at $t=0$ and have covariant
  derivatives $\frac{a}{c}\gamma'(0)$ resp. $Z_\perp$ there; their sum has initial data
  $(0,Z)$, so by uniqueness of Jacobi fields with prescribed initial data it equals $J$.
  Parallel transport is an isometry preserving orthogonality, so
  $\langle\gamma'(t),w(t)\rangle=\langle\gamma'(0),Z_\perp\rangle=0$, $|\gamma'(t)|^2=c$
  (geodesics have constant speed) and $|w(t)|^2=|Z_\perp|^2=|Z|^2-\frac{a^2}{c}$; expanding
  the squared norm of the orthogonal combination
  $\big|\frac{a}{c}\,t\,\gamma'(t)+h(t)\,w(t)\big|^2$ gives
  $\frac{a^2}{c}t^2+h(t)^2(|Z|^2-\frac{a^2}{c})$, which proves the formula.
\end{proof}

\begin{lemma}[the differential of $\exp$ is norm-preserving between spaces of the same
  constant curvature]
  \label{lem:dc-ch8-2-1-exp-norm-transfer}
  \lean{Riemannian.Jacobi.chartMetricInner_expDifferential_eq_metricInner_jacobiField, Riemannian.Jacobi.chartMetricInner_expDifferential_transfer_of_constantCurvature_of_speedSq}
  \uses{lem:dc-ch8-2-1-jacobi-norm-transfer, cor:dc-ch5-2-5, def:dc-ch5-2-1}
  \leanok
  Let $M$ and $\tilde M$ be complete manifolds of the same constant sectional curvature
  $K_0$, let $p\in M$, $\tilde p\in\tilde M$, and let $v\in T_p M$, $\tilde v\in T_{\tilde p}\tilde M$
  be \emph{nonzero} vectors of the same length, $|v|_p^2=|\tilde v|_{\tilde p}^2=c$. Let
  $Z\in T_p M$ and $\tilde Z\in T_{\tilde p}\tilde M$ have matching scalar invariants,

  $$
  \langle \tilde Z,\tilde v\rangle_{\tilde p}=\langle Z,v\rangle_p,
  \qquad |\tilde Z|^2_{\tilde p}=|Z|^2_p .
  $$

  Then the two exponential differentials have equal norms,

  $$
  \big|(d\exp_{\tilde p})_{\tilde v}(\tilde Z)\big| = \big|(d\exp_p)_v(Z)\big| .
  $$

  If $i$ is a linear isometry with $\tilde v=i(v)$ and $\tilde Z=i(Z)$, these
  hypotheses hold automatically.
\end{lemma}
\begin{proof}
  \leanok
\sketch
  By \cref{cor:dc-ch5-2-5}, each differential is the value at $t=1$ of the
  Jacobi field with initial data $(0,Z)$ or $(0,\tilde Z)$. Constant speed gives the
  common squared speed $c$, so \cref{lem:dc-ch8-2-1-jacobi-norm-transfer} applies.
\end{proof}

\begin{lemma}[in constant curvature, no conjugate point where the scalar solution is nonzero]
  \label{lem:dc-ch8-2-1-no-conjugate}
  \lean{Riemannian.Jacobi.not_isConjugatePointAt_of_constantCurvature_of_speedSq, Riemannian.Jacobi.not_isConjugatePointAt_of_constantCurvature_of_lt_pi, Riemannian.Jacobi.exists_constCurvatureSol_ne_zero}
  \uses{lem:dc-ch8-2-1-jacobi-norm-transfer, lem:dc-ch8-2-1-jacobi-const-solution, def:dc-ch5-3-1}
  \leanok
  Let $\gamma:[0,\ell]\to M$ be a geodesic of squared speed $c\neq0$ on a manifold of constant
  sectional curvature $K_0$, and let $h$ solve $h''+K_0c\,h=0$, $h(0)=0$, $h'(0)=1$. If
  $h(\ell)\neq0$, then $\gamma(\ell)$ is \emph{not} conjugate to $\gamma(0)$ along $\gamma$.

  Only this direction is asserted. The converse --- that a zero of $h$ \emph{is} a conjugate
  point --- is not proved here; for $K_0>0$ it is witnessed at the first zero by
  \cref{ex:dc-ch5-3-3}, which exhibits $\gamma(\pi/\sqrt{K_0})$ as conjugate to $\gamma(0)$ in
  the unit-speed case (and additionally needs $n\geq2$).

  The criterion is sign-uniform. For $K_0\leq0$ the solution is $h(t)=t$ resp.
  $h(t)=\sinh(t\sqrt{-K_0c})/\sqrt{-K_0c}$, which never vanishes for $t>0$: there are no
  conjugate points at all, recovering the nonpositive-curvature statement of
  \cref{lem:dc-ch7-3-2}. For $K_0>0$ the solution is $h(t)=\sin(t\sqrt{K_0c})/\sqrt{K_0c}$,
  whose first positive zero is $\pi/\sqrt{K_0c}$, so the hypothesis holds throughout
  $0<\ell<\pi/\sqrt{K_0c}$ --- the interval that matters here, though $h$ is of course nonzero
  on further intervals beyond the first zero as well.

  Equivalently, in the $h$-free form: if

  $$
  K_0\,c\,\ell^2 < \pi^2,
  $$

  then $\gamma(\ell)$ is not conjugate to $\gamma(0)$. This single numerical condition covers
  all three signs at once --- it is vacuous for $K_0\leq0$, and for $K_0>0$ it reads
  $\ell\sqrt{K_0c}<\pi$.

  At parameter $1$ and $K_0>0$ this gives $|v|<\pi/\sqrt{K_0}$.
\end{lemma}
\begin{proof}
  \leanok
\sketch
  Let $J$ be a Jacobi field along $\gamma$ with $J(0)=0$ and $J(\ell)=0$; we show $J\equiv0$,
  which contradicts the nontriviality required of a conjugacy witness. Write $Z=DJ(0)$ and
  $a=\langle Z,\gamma'(0)\rangle$. By \cref{lem:dc-ch8-2-1-jacobi-norm-transfer},

  $$
  0=|J(\ell)|^2=\frac{a^2}{c}\,\ell^2+h(\ell)^2\Big(|Z|^2-\frac{a^2}{c}\Big).
  $$

  Both summands are nonnegative: $c>0$ because $c=|\gamma'(0)|^2$ and the metric is positive
  definite with $c\neq0$; and $|Z|^2-\frac{a^2}{c}\geq0$ by Cauchy--Schwarz, being the squared
  norm of the component of $Z$ orthogonal to $\gamma'(0)$. A sum of two nonnegative reals
  vanishes only if both do. The first gives $\frac{a^2}{c}\ell^2=0$, hence $a=0$ since
  $\ell>0$ and $c>0$. The second then reads $h(\ell)^2|Z|^2=0$, and $h(\ell)\neq0$ forces
  $|Z|^2=0$, so $Z=0$. Uniqueness then gives $J\equiv0$.

  For the $h$-free form, it remains to see that $K_0c\,\ell^2<\pi^2$ gives $h(\ell)\neq0$.
  Write $\kappa=K_0c$. If $\kappa<0$ then $h(t)=\sinh(t\sqrt{-\kappa})/\sqrt{-\kappa}$ and
  $\sinh$ is strictly increasing with $\sinh 0=0$, so $h(\ell)>0$; if $\kappa=0$ then
  $h(t)=t$ and $h(\ell)=\ell>0$; if $\kappa>0$ then $h(t)=\sin(t\sqrt\kappa)/\sqrt\kappa$, and
  $\kappa\ell^2<\pi^2$ with $\ell,\sqrt\kappa>0$ gives $0<\ell\sqrt\kappa<\pi$, whence
  $\sin(\ell\sqrt\kappa)>0$.
\end{proof}

\begin{lemma}[$\exp_p$ is a local diffeomorphism away from conjugate points]
  \label{lem:dc-ch8-2-1-exp-local-diffeo}
  \lean{Riemannian.Jacobi.isLocalDiffeomorphAt_expMapGlobal_of_not_conjugate}
  \uses{lem:dc-ch8-2-1-exp-equiv, prop:dc-ch5-3-5, lem:dc-ch7-3-2}
  \leanok
  Let $M$ be complete and $p\in M$. If the geodesic $\gamma_v$ from $p$ has no conjugate point
  of $p$ at parameter $1$, then $\exp_p$ is a $C^\infty$ local diffeomorphism at $v$; in
  particular it has a $C^\infty$ local inverse $\exp_p^{-1}$ defined near $\exp_p(v)$.

  Consequently $\exp_p^{-1}$ is smooth near $\exp_p(v)$.
\end{lemma}
\begin{proof}
  \leanok
\sketch
  By \cref{lem:dc-ch8-2-1-exp-equiv} the differential $d(\exp_p)_v$ is a continuous linear
  isomorphism: some chart $\zeta$ around $\exp_p(v)$ reads $\exp_p$ as a map with strict
  Fréchet derivative an isomorphism at $v$. That chart reading is $C^\infty$ on an open
  neighborhood of $v$, since $\exp_p$ is globally $C^\infty$ on a complete manifold, so the
  inverse function theorem makes the reading a local diffeomorphism at $v$. Composing with the
  chart inverse --- itself a local diffeomorphism --- recovers $\exp_p$ near $v$, which is
  therefore a local diffeomorphism at $v$ as well.
\end{proof}

\begin{lemma}[non-conjugacy from the normal neighborhood, in variable curvature]
  \label{lem:dc-ch8-2-1-no-conjugate-normal}
  \lean{Riemannian.Jacobi.not_isConjugatePointAt_globalGeodesic_of_injective_mfderiv, Riemannian.Jacobi.not_isConjugatePointAt_globalGeodesic_of_isLocalDiffeomorphAt, Riemannian.Jacobi.isLocalDiffeomorphAt_expMapGlobal_iff_not_isConjugatePointAt}
  \uses{lem:dc-ch8-2-1-exp-local-diffeo, lem:dc-ch8-2-1-mfderiv-bridge, prop:dc-ch5-3-5}
  \leanok
  Let $M$ be complete and $p\in M$, and let $v\in T_pM$. If $d(\exp_p)_v$ is injective --- in
  particular, if $\exp_p$ is a local diffeomorphism at $v$ --- then $1$ is \emph{not} a conjugate
  parameter along $\gamma_v$. No curvature hypothesis is needed.

  Together with \cref{lem:dc-ch8-2-1-exp-local-diffeo} this makes the two conditions
  \emph{equivalent}: $\exp_p$ is a $C^\infty$ local diffeomorphism at $v$ if and only if $1$ is
  not conjugate along $\gamma_v$.

  Thus the criterion applies in variable curvature whenever $\exp_p$ is locally invertible.
\end{lemma}
\begin{proof}
  \leanok
  \uses{lem:dc-ch8-2-1-mfderiv-bridge,prop:dc-ch5-3-5}
\sketch
  By \cref{prop:dc-ch5-3-5} it suffices to show that the chart reading $D$ of $d(\exp_p)_v$ is
  injective. \cref{lem:dc-ch8-2-1-mfderiv-bridge} factors it as

  $$
  D = C_{\exp_p(v)\to\zeta}\circ d(\exp_p)_v ,
  $$

  and the coordinate change $C_{\exp_p(v)\to\zeta}$ is injective, its cocycle inverse being a left
  inverse; so $D$ is a composite of two injective maps.

  A local diffeomorphism has an injective differential, giving the second assertion.
\end{proof}

\begin{lemma}[polarization: a linear map preserving lengths preserves the metric]
  \label{lem:dc-ch8-2-1-polarization}
  \lean{Riemannian.Jacobi.metricInner_polarization, Riemannian.Jacobi.metricInner_transfer_of_norm_transfer}
  \leanok
  Let $q\in M$, $\tilde q\in\tilde M$, and let $L:T_qM\to T_{\tilde q}\tilde M$ be additive. If
  $L$ preserves squared lengths, $|L(w)|^2_{\tilde q}=|w|^2_q$ for every $w\in T_qM$, then it
  preserves the metric itself:

  $$
  \langle L(u),L(v)\rangle_{\tilde q} = \langle u,v\rangle_q
  \qquad\text{for all } u,v\in T_qM .
  $$

  This is the step from \cref{lem:dc-ch8-2-1-exp-norm-transfer}, which delivers only the
  equality of \emph{lengths} $|df_q(u)|=|u|$, to the statement that $f$ is a local isometry,
  which is phrased with the full inner product. Additivity cannot be dropped: a
  length-preserving map that is not additive need not preserve the metric. In the application
  $L=df_q$ is linear, so the hypothesis is automatic.
\end{lemma}
\begin{proof}
  \leanok
\sketch
  A symmetric bilinear form over $\mathbb R$ is determined by its diagonal through the
  polarization identity: expanding $\langle u+v,u+v\rangle=\langle u,u\rangle+2\langle
  u,v\rangle+\langle v,v\rangle$ by bilinearity and symmetry gives

  $$
  \langle u,v\rangle=\tfrac12\big(\langle u+v,u+v\rangle-\langle u,u\rangle-\langle
  v,v\rangle\big).
  $$

  Apply this to $L(u),L(v)$ on $\tilde M$, use additivity to rewrite $L(u)+L(v)=L(u+v)$, then
  the length hypothesis at $u+v$, $u$ and $v$ to replace each of the three terms by its
  counterpart on $M$, and finally the same identity on $M$ read backwards:

  \langle L(u),L(v)\rangle=\tfrac12\big(|L(u+v)|^2-|L(u)|^2-|L(v)|^2\big)
  =\tfrac12\big(|u+v|^2-|u|^2-|v|^2\big)=\langle u,v\rangle. \qquad

\end{proof}

\begin{lemma}[$d(\exp_p)_v$ is an isomorphism, with the Jacobi identification]
  \label{lem:dc-ch8-2-1-exp-equiv}
  \lean{Riemannian.Jacobi.expDifferential_isEquiv_jacobi_of_not_conjugate}
  \uses{cor:dc-ch5-2-5, prop:dc-ch5-3-5}
  \leanok
  If the geodesic $\gamma_v$ from $p$ has no conjugate point of $p$ at parameter $1$, then
  $d(\exp_p)_v$ is a continuous linear isomorphism $T_pM\to T_{\exp_p(v)}M$, \emph{and} it still
  satisfies the identification of \cref{cor:dc-ch5-2-5}: it carries $\nabla J(0)$ to $J(1)$ for
  every Jacobi field $J$ along $\gamma_v$ with $J(0)=0$. Both facts come from the same strictly
  differentiable chart reading of $\exp_p$; the assembly of \cref{cor:dc-ch8-2-2} needs them
  simultaneously --- the isomorphism to invert $d\exp_p$ (so that every tangent vector at
  $q=\exp_p(v)$ is the endpoint value of a Jacobi field), and the identification to compute
  norms through \cref{lem:dc-ch8-2-1-exp-norm-transfer}.
\end{lemma}
\begin{proof}
  \leanok
\sketch
  \cref{prop:dc-ch5-3-5} makes the differential injective from the no-conjugate-point
  hypothesis; finite-dimensionality makes it an isomorphism. The Jacobi identification is
  \cref{cor:dc-ch5-2-5}.
\end{proof}

\begin{lemma}[the chart reading of a differential is the coordinate change of the intrinsic one]
  \label{lem:dc-ch8-2-1-mfderiv-bridge}
  \lean{Riemannian.Jacobi.chartReading_fderiv_apply_eq, Riemannian.Jacobi.chartMetricInner_chartReading_fderiv_eq_metricInner_mfderiv, Riemannian.Jacobi.chartMetricInner_expDifferential_eq_metricInner_mfderiv}
  \leanok
  Let $F:T_pM\to M$ be $C^\infty$, let $\zeta\in M$ be a chart base point with $F(v)$ in the
  chart source, and suppose the chart-$\zeta$ reading $w\mapsto\varphi_\zeta(F(w))$ has
  differential $D$ at $v$. Then $D$ is the intrinsic differential of $F$ followed by the
  coordinate change of its own target fibre:

  $$
  D(Z)=C_{F(v)\to\zeta}\big(dF_v(Z)\big)\qquad\text{for all }Z ,
  $$

  and consequently, for the chart Gram form,

  $$
  \langle D(Z),D(Z')\rangle_{\zeta,\varphi_\zeta(F(v))}=\langle dF_v(Z),dF_v(Z')\rangle_{F(v)} .
  $$

  This is the bridge between the chart-Gram form in which \cref{lem:dc-ch8-2-1-exp-norm-transfer}
  is stated and the intrinsic form that \cref{cor:dc-ch8-2-2} requires; applied to $F=\exp_p$ it
  turns every chart-level statement about $d(\exp_p)_v$ into an intrinsic one.
\end{lemma}
\begin{proof}
  \leanok
\sketch
  Apply the chain rule to $\varphi_\zeta\circ F$. Its differential is
  $C_{F(v)\to\zeta}\circ dF_v$, and uniqueness of the derivative gives the first
  display. Norm-faithfulness of chart readings gives the second.
\end{proof}

\begin{lemma}[$d(\exp_p)_0$ is the identity]
  \label{lem:dc-ch8-2-1-exp-diff-zero}
  \lean{Riemannian.Jacobi.mfderiv_expMapGlobal_zero_apply}
  \leanok
  For $M$ complete and $p\in M$, the differential of $\exp_p$ at the origin of $T_pM$ is the
  identity: $d(\exp_p)_0(v)=v$ for every $v\in T_pM$.
\end{lemma}
\begin{proof}
  \leanok
\sketch
  Fix $v$ and consider the ray $c:\mathbb R\to T_pM$, $c(t)=t\,v$, a linear map, so its own
  differential at $0$ is $c$ itself, and $c(1)=v$. Radial homogeneity of the exponential map gives
  $\exp_p(c(t))=\exp_p(t\,v)=\gamma_v(t)$, i.e. $\exp_p\circ\,c=\gamma_v$. Differentiating this
  identity at $t=0$ by the chain rule:
  $d(\exp_p)_0(c(1))=d(\gamma_v)_0(1)=\gamma_v'(0)=v$.
  Since $c(1)=v$, this is $d(\exp_p)_0(v)=v$.
\end{proof}

\begin{lemma}[the base point is never conjugate to itself at parameter $1$]
  \label{lem:dc-ch8-2-1-no-conjugate-zero}
  \lean{Riemannian.Jacobi.not_isConjugatePointAt_globalGeodesic_zero}
  \uses{lem:dc-ch8-2-1-mfderiv-bridge, lem:dc-ch8-2-1-exp-diff-zero, prop:dc-ch5-3-5}
  \leanok
  For $M$ complete and $p\in M$, the parameter $1$ is not conjugate along the constant geodesic
  $\gamma_0$. No curvature hypothesis is needed. Consequently, by
  \cref{lem:dc-ch8-2-1-exp-local-diffeo}, $\exp_p$ is a $C^\infty$ local diffeomorphism at
  $0\in T_pM$ on \emph{any} complete Riemannian manifold --- the normal-neighborhood fact.

\end{lemma}
\begin{proof}
  \leanok
\sketch
  By \cref{prop:dc-ch5-3-5} the parameter $1$ fails to be conjugate along $\gamma_v$ exactly when
  the chart reading $D$ of $d(\exp_p)_v$ is injective, so it suffices to prove that for $v=0$. By
  \cref{lem:dc-ch8-2-1-mfderiv-bridge}, $D(Z)=C_{\exp_p(0)\to\zeta}\big(d(\exp_p)_0(Z)\big)$, and
  $\exp_p(0)=p$ while $d(\exp_p)_0$ is the identity by \cref{lem:dc-ch8-2-1-exp-diff-zero}. Hence
  $D=C_{p\to\zeta}$, which is injective, since the cocycle identity
  $C_{\zeta\to p}\circ C_{p\to\zeta}=C_{p\to p}=\mathrm{id}$ exhibits a left inverse.
\end{proof}

\begin{lemma}[non-conjugacy along $\gamma_v$, and surjectivity of $d(\exp_p)_v$]
  \label{lem:dc-ch8-2-1-no-conjugate-ball}
  \lean{Riemannian.Jacobi.not_isConjugatePointAt_globalGeodesic_of_constantCurvature_of_lt_pi, Riemannian.Jacobi.surjective_mfderiv_expMapGlobal_of_not_conjugate}
  \uses{lem:dc-ch8-2-1-no-conjugate, lem:dc-ch8-2-1-exp-local-diffeo}
  \leanok
  Let $M$ have constant curvature $K_0$ and let $v\in T_pM$ with $v\neq 0$. If

  $$
  K_0\,\langle v,v\rangle_p<\pi^2 ,
  $$

  then $1$ is not conjugate along $\gamma_v$. The condition is vacuous for $K_0\leq 0$ and reads
  $|v|<\pi/\sqrt{K_0}$ for $K_0>0$; it is sharp there, by \cref{ex:dc-ch5-3-3}. Moreover, whenever
  $1$ is not conjugate along $\gamma_v$, the differential $d(\exp_p)_v$ is surjective onto
  $T_{\exp_p(v)}M$.

  Thus every tangent vector at $\exp_p(v)$ is $d(\exp_p)_v(Z)$ for some $Z$.
\end{lemma}
\begin{proof}
  \leanok
\sketch
  For the first claim, apply \cref{lem:dc-ch8-2-1-no-conjugate} along $\gamma_v$ at $\ell=1$: the
  geodesic $\gamma_v$ has constant speed $c=\langle v,v\rangle_p$, which is nonzero because the
  metric is positive definite and $v\neq 0$, and the numerical hypothesis
  $K_0\,c\,\ell^2<\pi^2$ is the assumption. For the second, \cref{lem:dc-ch8-2-1-exp-local-diffeo}
  makes $\exp_p$ a local diffeomorphism at $v$, and the differential of a local diffeomorphism is
  a linear isomorphism, in particular surjective.
\end{proof}

\begin{lemma}[the norm transfer, intrinsically]
  \label{lem:dc-ch8-2-1-exp-norm-transfer-mfderiv}
  \lean{Riemannian.Jacobi.metricInner_mfderiv_expMapGlobal_transfer_of_constantCurvature_of_speedSq}
  \uses{lem:dc-ch8-2-1-exp-norm-transfer, lem:dc-ch8-2-1-mfderiv-bridge, lem:dc-ch8-2-1-exp-equiv, lem:dc-ch8-2-1-jacobi-const-solution}
  \leanok
  Let $M$ and $\tilde M$ have the same constant curvature $K_0$, let $p\in M$, $\tilde p\in\tilde M$,
  and let $v,Z\in T_pM$, $\tilde v,\tilde Z\in T_{\tilde p}\tilde M$ with $v\neq 0$ and

  $$
  \langle v,v\rangle_p=\langle\tilde v,\tilde v\rangle_{\tilde p}=c,\qquad
  \langle \tilde Z,\tilde v\rangle_{\tilde p}=\langle Z,v\rangle_p,\qquad
  \langle \tilde Z,\tilde Z\rangle_{\tilde p}=\langle Z,Z\rangle_p .
  $$

  If $1$ is conjugate along neither $\gamma_v$ nor $\gamma_{\tilde v}$, then the two exponential
  differentials agree in length:

  $$
  \big|d(\exp_{\tilde p})_{\tilde v}(\tilde Z)\big|_{\exp_{\tilde p}(\tilde v)}
  =\big|d(\exp_p)_v(Z)\big|_{\exp_p(v)} .
  $$

  This is \cref{lem:dc-ch8-2-1-exp-norm-transfer} freed of its chart, and of the auxiliary scalar
  solution $h$.
\end{lemma}
\begin{proof}
  \leanok
\sketch
  \cref{lem:dc-ch8-2-1-exp-equiv} supplies, on each side, a chart around the endpoint in which the
  reading of the exponential map is differentiable with derivative $D$ (resp.\ $\tilde D$) still
  carrying $\nabla J(0)$ to $J(1)$. \cref{lem:dc-ch8-2-1-exp-norm-transfer} then equates the two
  chart-Gram norms of $D(Z)$ and $\tilde D(\tilde Z)$, its scalar hypothesis being discharged by
  \cref{lem:dc-ch8-2-1-jacobi-const-solution}, which produces a solution of $h''+(K_0c)h=0$ with
  $h(0)=0$, $h'(0)=1$ for every $K_0$ and $c$. Finally
  \cref{lem:dc-ch8-2-1-mfderiv-bridge} rewrites each chart-Gram norm as the intrinsic norm of the
  corresponding intrinsic differential.
\end{proof}

\begin{lemma}[metric preservation along a semiconjugacy]
  \label{lem:dc-ch8-2-2-semiconjugacy}
  \lean{Riemannian.Jacobi.mfderiv_apply_eq_of_semiconjugacy_zero, Riemannian.Jacobi.metricInner_mfderiv_eq_of_semiconjugacy, Riemannian.Jacobi.dcIsLocalIsometryAt_of_semiconjugacy, Riemannian.Jacobi.dcIsLocalIsometryAt_of_semiconjugacy_self}
  \uses{lem:dc-ch8-2-1-exp-norm-transfer-mfderiv, lem:dc-ch8-2-1-no-conjugate-ball, lem:dc-ch8-2-1-no-conjugate-zero, lem:dc-ch8-2-1-exp-diff-zero, lem:dc-ch8-2-1-polarization, lem:dc-ch8-2-1-exp-local-diffeo}
  \leanok
  Let $M$ and $\tilde M$ have the same constant curvature $K_0$, let $i:T_pM\to T_{\tilde p}\tilde M$
  be a linear isometry, and let $W\subset T_pM$ be an open neighborhood of $0$ on which
  $K_0\langle w,w\rangle_p<\pi^2$. Suppose $f:M\to\tilde M$ is differentiable at each $\exp_p(w)$,
  $w\in W$, and \emph{semiconjugates} $\exp_p$ to $\exp_{\tilde p}$ through $i$:

  $$
  f\big(\exp_p(w)\big)=\exp_{\tilde p}\big(i(w)\big)\qquad\text{for all }w\in W .
  $$

  Then $df_p=i$, and $f$ is a local isometry at $p$: there is a neighborhood $V$ of $p$ with
  $\langle u,u'\rangle_q=\langle df_q(u),df_q(u')\rangle_{f(q)}$ for all $q\in V$ and
  $u,u'\in T_qM$.

  Taking the semiconjugacy as the hypothesis, rather than $f=\exp_{\tilde p}\circ\,i\circ\exp_p^{-1}$,
  is what makes the differential of $\exp_p^{-1}$ unnecessary: $d(\exp_p)_v$ is surjective, so $df_q$
  is pinned down on all of $T_qM$ by its values on the image of $d(\exp_p)_v$.
\end{lemma}
\begin{proof}
  \leanok
\sketch
  Take $V=\exp_p(W')$, where $W'\subset W$ is the part of $W$ on which $\exp_p$ is inverted by the
  local inverse supplied by \cref{lem:dc-ch8-2-1-exp-local-diffeo} at $0$, available by
  \cref{lem:dc-ch8-2-1-no-conjugate-zero}. This $V$ is a neighborhood of $p$, and each $q\in V$ is
  $\exp_p(w)$ for a unique $w\in W'$.

  Suppose first $w\neq0$. Since $i$ is a linear isometry, $\langle i(w),i(w)\rangle_{\tilde p}
  =\langle w,w\rangle_p$, so the numerical hypothesis holds on both sides and by
  \cref{lem:dc-ch8-2-1-no-conjugate-ball} the parameter $1$ is conjugate along neither $\gamma_w$
  nor $\gamma_{i(w)}$; in particular $d(\exp_p)_w$ is surjective. Differentiating the semiconjugacy
  at $w$ by the chain rule --- both sides are differentiable, and the derivative of a map is unique
  --- gives

  $$
  df_q\big(d(\exp_p)_w(Z)\big)=d(\exp_{\tilde p})_{i(w)}\big(i(Z)\big)\qquad\text{for all }Z .
  $$

  The invariants $\langle i(Z),i(w)\rangle=\langle Z,w\rangle$ and $|i(Z)|^2=|Z|^2$ are preserved
  because $i$ is a linear isometry, so \cref{lem:dc-ch8-2-1-exp-norm-transfer-mfderiv} equates the
  lengths of the two sides. As $Z$ ranges over $T_pM$ the vector $d(\exp_p)_w(Z)$ ranges over all of
  $T_qM$, so $df_q$ preserves lengths; being linear, it preserves the metric, by
  \cref{lem:dc-ch8-2-1-polarization}.

  If instead $w=0$ then $q=\exp_p(0)=p$. The same chain rule at $0$, together with
  $d(\exp_p)_0=d(\exp_{\tilde p})_0=\mathrm{id}$ (\cref{lem:dc-ch8-2-1-exp-diff-zero}) and
  $i(0)=0$, gives $df_p=i$; and $f(p)=\exp_{\tilde p}(i(0))=\tilde p$. The claim at $q=p$ is then
  exactly the hypothesis that $i$ is a linear isometry.
\end{proof}

\begin{theorem}[Cartan local isometry theorem]
  \label{thm:dc-ch8-2-1}
  \uses{cor:dc-ch5-2-5, lem:dc-ch8-2-1-jacobi-transfer, lem:dc-ch8-2-1-parallel-ortho-frame, lem:dc-ch8-2-1-transported-frame, lem:dc-ch8-2-1-hmatch-transfer, lem:dc-ch8-2-1-curvature-bridge, lem:dc-ch8-2-1-frame-coeff-bridge, lem:dc-ch8-2-1-jacobi-norm-transfer-general, lem:dc-ch8-2-1-frame-chart-interface, lem:dc-ch8-2-1-jacobi-norm-transfer-manifold, lem:dc-ch8-2-1-single-chart, lem:dc-ch8-2-1-exp-norm-transfer-general, lem:dc-ch8-2-1-phi, lem:dc-ch8-2-1-parallel-transport-map, lem:dc-ch8-2-1-exp-norm-transfer, lem:dc-ch8-2-1-exp-equiv, lem:dc-ch8-2-1-exp-local-diffeo, lem:dc-ch8-2-1-polarization, lem:dc-ch8-2-1-no-conjugate, lem:dc-ch8-2-1-mfderiv-bridge, lem:dc-ch8-2-1-exp-diff-zero, lem:dc-ch8-2-1-no-conjugate-zero, lem:dc-ch8-2-1-no-conjugate-ball, lem:dc-ch8-2-1-exp-norm-transfer-mfderiv, lem:dc-ch8-2-1-jacobi-reversal, lem:dc-ch8-2-1-jacobi-backward-uniqueness, lem:dc-ch8-2-1-jacobi-interior-data, lem:dc-ch8-2-1-local-isometry-packaging}
  \notready
  \dcref{ch8:2.1}
  \notready
  With the notation above, if for all $q\in V$ and all $x,y,u,v\in T_q(M)$ we have


  $$
  \langle R(x,y)u,v\rangle=\big\langle\tilde{R}(\phi_t(x),\phi_t(y))\phi_t(u),\phi_t(v)\big\rangle,
  $$


  then $f:V\to f(V)\subset\tilde{M}$ is a local isometry and $df_p=i$.
\end{theorem}
\begin{proof}
\sketch
  Let $q\in V$ and let $\gamma:[0,\ell]\to M$ be a normalized geodesic with
  $\gamma(0)=p$, $\gamma(\ell)=q$. Let $v\in T_q(M)$ and let $J$ be a Jacobi field
  along $\gamma$ with $J(0)=0$ and $J(\ell)=v$. Let $e_1,\dots,e_n=\gamma'(0)$ be an
  orthonormal basis of $T_p(M)$ and let $e_i(t)$ be the parallel transport of $e_i$
  along $\gamma$. Writing $J(t)=\sum_i y_i(t)e_i(t)$, the Jacobi equation gives


  $$
  y_j''+\sum_i\langle R(e_n,e_i)e_n,e_j\rangle y_i=0,\qquad j=1,\dots,n.
  $$


  Let $\tilde{\gamma}:[0,\ell]\to\tilde{M}$ be the normalized geodesic with
  $\tilde{\gamma}(0)=\tilde{p}$, $\tilde{\gamma}'(0)=i(\gamma'(0))$, and set
  $\tilde{J}(t)=\phi_t(J(t))$, $\tilde{e}_j(t)=\phi_t(e_j(t))$. By linearity of
  $\phi_t$, $\tilde{J}(t)=\sum_i y_i(t)\tilde{e}_i(t)$. Since by hypothesis
  $\langle R(e_n,e_i)e_n,e_j\rangle=\langle\tilde{R}(\tilde{e}_n,\tilde{e}_i)\tilde{e}_n,\tilde{e}_j\rangle$,
  $\tilde{J}$ satisfies the same equation, hence $\tilde{J}$ is a Jacobi field along
  $\tilde{\gamma}$ with $\tilde{J}(0)=0$. Since parallel transport is an isometry,
  $|\tilde{J}(\ell)|=|J(\ell)|$. It remains to show that $\tilde{J}(\ell)=df_q(v)=df_q(J(\ell))$.
  Since $\tilde{J}(t)=\phi_t(J(t))$, $\tilde{J}'(0)=i(J'(0))$. Since $J$ and
  $\tilde{J}$ are Jacobi fields vanishing at $t=0$, by \cref{cor:dc-ch5-2-5},


  $$
  J(t)=(d\exp_p)_{t\gamma'(0)}(tJ'(0)),\qquad\tilde{J}(t)=(d\exp_{\tilde{p}})_{t\tilde{\gamma}'(0)}(t\tilde{J}'(0)).
  $$


  Therefore


  $$
  \tilde{J}(\ell)=(d\exp_{\tilde{p}})_{\ell\tilde{\gamma}'(0)}\circ i\circ\big((d\exp_p)_{\ell\gamma'(0)}\big)^{-1}(J(\ell))=df_q(J(\ell)),
  $$


  which proves the assertion.

  If both exponential maps are global diffeomorphisms, the same argument makes $f$
  a global isometry.
\end{proof}

\begin{corollary}[Local equivalence of equal-curvature spaces]
  \label{cor:dc-ch8-2-2}
  \lean{Riemannian.Jacobi.exists_dcIsLocalIsometryAt_of_constantCurvature_of_orthonormal_bases}
  \uses{lem:dc-ch8-2-1-jacobi-transfer-const, lem:dc-ch8-2-1-velocity-frame, lem:dc-ch8-2-1-jacobi-norm-transfer, lem:dc-ch8-2-1-jacobi-const-solution, lem:dc-ch8-2-1-exp-norm-transfer, lem:dc-ch8-2-1-exp-equiv, lem:dc-ch8-2-1-exp-local-diffeo, lem:dc-ch8-2-1-polarization, lem:dc-ch8-2-1-no-conjugate, lem:dc-ch8-2-1-mfderiv-bridge, lem:dc-ch8-2-1-exp-diff-zero, lem:dc-ch8-2-1-no-conjugate-zero, lem:dc-ch8-2-1-no-conjugate-ball, lem:dc-ch8-2-1-exp-norm-transfer-mfderiv, lem:dc-ch8-2-2-semiconjugacy}
  \leanok
  \dcref{ch8:2.2}
  Let $M$ and $\tilde{M}$ be spaces with the same constant curvature and the same
  dimension $n$. Let $p\in M$ and $\tilde{p}\in\tilde{M}$. Choose arbitrary
  orthonormal bases $\{e_j\}$ of $T_p(M)$ and $\{\tilde{e}_j\}$ of $T_{\tilde{p}}(\tilde{M})$,
  $j=1,\dots,n$. Then there exist a neighborhood $V\subset M$ of $p$, a neighborhood
  $\tilde{V}\subset\tilde{M}$ of $\tilde{p}$, and an isometry $f:V\to\tilde{V}$ such
  that $df_p(e_j)=\tilde{e}_j$.
\end{corollary}
\begin{proof}
  \leanok
\sketch
  Let $i(e_j)=\tilde e_j$ and define
  $f=\exp_{\tilde p}\circ i\circ\exp_p^{-1}$ on a sufficiently small normal
  neighborhood. The constant-curvature identity gives the curvature match; equivalently,
  \cref{lem:dc-ch8-2-2-semiconjugacy} applies on
  $K_0\langle w,w\rangle_p<\pi^2$. It yields $df_p=i$ and local metric preservation.
  Restricting the resulting local diffeomorphism gives $f:V\to\tilde V$.
\end{proof}

\begin{corollary}[Local frame transitivity]
  \label{cor:dc-ch8-2-3}
  \lean{Riemannian.Jacobi.exists_dcIsLocalIsometryAt_of_constantCurvature_self_of_orthonormal_bases}
  \uses{lem:dc-ch8-2-1-jacobi-transfer-const, lem:dc-ch8-2-1-velocity-frame, lem:dc-ch8-2-1-jacobi-norm-transfer, lem:dc-ch8-2-1-jacobi-const-solution, lem:dc-ch8-2-1-exp-norm-transfer, lem:dc-ch8-2-1-exp-equiv, lem:dc-ch8-2-1-exp-local-diffeo, lem:dc-ch8-2-1-polarization, lem:dc-ch8-2-1-no-conjugate, lem:dc-ch8-2-1-mfderiv-bridge, lem:dc-ch8-2-1-exp-diff-zero, lem:dc-ch8-2-1-no-conjugate-zero, lem:dc-ch8-2-1-no-conjugate-ball, lem:dc-ch8-2-1-exp-norm-transfer-mfderiv, lem:dc-ch8-2-2-semiconjugacy}
  \leanok
  \dcref{ch8:2.3}
  Let $M$ be a space of constant curvature and let $p$ and $q$ be any two points of
  $M$. Let $\{e_j\}$ and $\{f_j\}$ be arbitrary orthonormal bases of $T_p(M)$ and
  $T_q(M)$, respectively. Then there exist neighborhoods $U$ of $p$ and $V$ of $q$,
  and an isometry $g:U\to V$ such that $dg_p(e_j)=f_j$.

\end{corollary}
\begin{proof}
  \leanok
  \uses{cor:dc-ch8-2-2}
\sketch
  Apply \cref{cor:dc-ch8-2-2} with $\tilde M=M$ and $\tilde p=q$, taking for $i$ the linear
  isometry $T_pM\to T_qM$ carrying $e_j$ to $f_j$.
\end{proof}

\section{Hyperbolic space}

\begin{definition}[hyperbolic space $H^n$ as a Riemannian manifold]
  \label{def:dc-ch8-3-hyperbolic-space}
  \lean{Riemannian.opensConformalMetric, Riemannian.Hyperbolic.hyperbolicMetric, Riemannian.Hyperbolic.hyperbolicMetric_apply, Riemannian.Hyperbolic.hyperbolicMetric_gram}
  The \emph{hyperbolic space} $H^n$ is the open half-space
  $H^n=\{x\in\mathbb R^n : x_n>0\}$ carrying the metric $g_{ij}=\delta_{ij}/x_n^2$
  of $(1)$, that is, the conformal rescaling of the ambient Euclidean metric by
  the positive smooth factor $1/x_n^2$: for $u,v\in T_pH^n$,
  $\langle u,v\rangle_g=\tfrac{1}{x_n^2}\langle u,v\rangle$.
\end{definition}

\begin{lemma}[Christoffel symbols and coordinate curvature of a conformal metric]
  \label{lem:dc-ch8-3-conformal-curvature}
  \lean{Riemannian.Hyperbolic.christoffelFromMetric_eq, Riemannian.Hyperbolic.Rcoeff_diag, Riemannian.Hyperbolic.hyperbolic_sectionalCurvature}
  \leanok
  For the conformal metric $g_{ij}=\delta_{ij}/F^2$ ($f=\log F$), the Christoffel
  symbols computed from the Koszul formula
  $\Gamma^k_{ij}=\tfrac12\sum_\ell g^{k\ell}(\partial_i g_{j\ell}+\partial_j g_{i\ell}-\partial_\ell g_{ij})$
  equal $\Gamma^k_{ij}=-\delta_{jk}f_i-\delta_{ki}f_j+\delta_{ij}f_k$, and the
  associated coordinate curvature coefficients satisfy, for $i\neq j$,
  $$F^2R_{ijij}=-\sum_\ell f_\ell^2+f_i^2+f_j^2+f_{ii}+f_{jj}.$$
  Specializing to $F=x_n$ (so $f=\log x_n$, $f_a=\delta_{an}/x_n$,
  $f_{ab}=-\delta_{an}\delta_{bn}/x_n^2$) gives, for every \emph{coordinate}
  $2$-plane spanned by $\partial/\partial x_i,\partial/\partial x_j$ with $i\neq j$,
  the coordinate-frame sectional curvature $K_{ij}=R_{ijij}F^4=-1$. (This is the
  coordinate identity; extending it to all $2$-planes and to the intrinsic
  curvature of the manifold is \cref{prop:dc-ch8-3-const-curv}.)
\end{lemma}
\begin{proof}
  \leanok
\sketch
  The Koszul identity is a direct computation from
  $g^{k\ell}=\delta_{k\ell}F^2$ and $\partial_m g_{ab}=-2\delta_{ab}f_m/F^2$.
  The curvature coefficient is
  $R^s_{ijk}=\partial_j\Gamma^s_{ik}-\partial_i\Gamma^s_{jk}+\sum_\ell(\Gamma^\ell_{ik}\Gamma^s_{j\ell}-\Gamma^\ell_{jk}\Gamma^s_{i\ell})$;
  substituting the closed form for $\Gamma$ and contracting the Kronecker deltas
  yields $R^j_{iji}=-\sum_\ell f_\ell^2+f_i^2+f_j^2+f_{ii}+f_{jj}$, which is
  $F^2R_{ijij}$ since $g_{jj}=1/F^2$. For $F=x_n$ one has
  $(f_a)^2+f_{aa}=0$ at every index, so the bracketed sum collapses to
  $-\sum_\ell f_\ell^2=-1/x_n^2$, and $K_{ij}=R_{ijij}F^4=(F^2R_{ijij})F^2=-1$.
\end{proof}

\begin{lemma}[the Riemann curvature tensor in chart coordinates]
  \label{lem:dc-ch8-3-chart-curvature}
  \lean{Riemannian.leviCivita_curvature_chartFrame_expansion}
  \leanok
  Let $g$ be a Riemannian metric on $M$ and $\nabla$ its Levi-Civita connection.
  In the coordinate frame $\partial_i$ of a chart at a point $p$, the intrinsic
  curvature operator is computed from the chart Christoffel symbols
  $\Gamma^k_{ij}$ by the classical coordinate formula

  $$
  R(\partial_i,\partial_j)\partial_k
    = \sum_\ell\Big(\partial_j\Gamma^\ell_{ik}-\partial_i\Gamma^\ell_{jk}
      +\sum_s(\Gamma^s_{ik}\Gamma^\ell_{js}-\Gamma^s_{jk}\Gamma^\ell_{is})\Big)\partial_\ell,
  $$

  where $\partial_j\Gamma^\ell_{ik}$ is the chart partial derivative of the
  coordinate Christoffel function. With the convention
  $R(X,Y)Z=\nabla_Y\nabla_X Z-\nabla_X\nabla_Y Z+\nabla_{[X,Y]}Z$ the bracketed
  coefficient is exactly the $R^\ell_{ijk}$ of $(2)$.
\end{lemma}
\begin{proof}
  \leanok
\sketch
  In a coordinate chart the frame vectors have vanishing pairwise Lie brackets,
  $[\partial_i,\partial_j]=0$, so the bracket term
  $\nabla_{[\partial_i,\partial_j]}\partial_k$ drops out. The field
  $\nabla_{\partial_j}\partial_k$ is, throughout the chart, the frame combination
  $\sum_m\Gamma^m_{jk}\partial_m$ (formula $(10)$ of Chap.~2, the Koszul collapse
  on the coordinate frame, valid at every chart point). Differentiating this
  along $\partial_i$ by the Leibniz rule for $\nabla$ gives
  $\nabla_{\partial_i}\nabla_{\partial_j}\partial_k
   =\sum_\ell\big(\partial_i\Gamma^\ell_{jk}
   +\sum_s\Gamma^s_{jk}\Gamma^\ell_{is}\big)\partial_\ell$, using
  $\nabla_{\partial_i}\partial_m=\sum_\ell\Gamma^\ell_{im}\partial_\ell$ for the
  inner derivative. Subtracting the term with $i$ and $j$ interchanged yields the
  stated coefficients.
\end{proof}

\begin{lemma}[hyperbolic chart Gram and inverse-Gram in the abstract chart frame]
  \label{lem:dc-ch8-3-hyperbolic-chart-gram}
  \lean{Riemannian.Hyperbolic.hyperbolic_chartGramMatrix, Riemannian.Hyperbolic.hyperbolic_chartGramMatrix_eq, Riemannian.Hyperbolic.hyperbolic_chartInvGramMatrix, Riemannian.Hyperbolic.finBasisGramMatrix_det_isUnit}
  \uses{def:dc-ch8-3-hyperbolic-space}
  \leanok
  Because $H^n=\{x_n>0\}$ is an open subset of $\mathbb R^n$, its tangent bundle is
  canonically trivial and the abstract chart frame vector
  $\partial_i=$ \emph{finBasis}$_i$ (the frame underlying \cref{lem:dc-ch8-3-chart-curvature})
  is literally the constant basis vector. Hence the chart Gram matrix of the
  hyperbolic metric is the conformal rescaling
  $G_{ij}(x)=x_n^{-2}\,B_{ij}$ of the constant Gram matrix
  $B_{ij}=\langle\text{finBasis}_i,\text{finBasis}_j\rangle$ of the abstract basis,
  and its inverse is $G^{ij}(x)=x_n^{2}\,B^{ij}$ (with $B$ invertible, being the
  Gram matrix of a basis). Since the abstract basis need not be
  $\langle\,,\rangle$-orthonormal, $B$ is a \emph{general} constant symmetric
  positive-definite matrix, not $\delta$: $\delta_{ij}$ is
  replaced by $B_{ij}$ in the abstract chart frame, and the conformal
  computation goes through with $\delta\rightsquigarrow B$, the extra factors
  collapsing by the inverse-Gram identity $\sum_{ij}B^{ij}c_ic_j=1$ (for
  $c_i=\langle e_n,\text{finBasis}_i\rangle$).
\end{lemma}
\begin{proof}
  \leanok
\sketch
  The trivialisation of the tangent bundle of an open subset of $\mathbb R^n$ is
  the identity on fibres, so $\partial_i(x)=$ finBasis$_i$ and
  $G_{ij}(x)=\langle\partial_i,\partial_j\rangle_g=x_n^{-2}\langle\text{finBasis}_i,\text{finBasis}_j\rangle=x_n^{-2}B_{ij}$.
  The Gram matrix at the basepoint is positive definite, so $\det B$ is a unit and
  the inverse of the scalar multiple $x_n^{-2}B$ is $x_n^{2}B^{-1}$.
\end{proof}

\begin{lemma}[hyperbolic chart Christoffel symbols and their derivatives, closed form]
  \label{lem:dc-ch8-3-hyperbolic-christoffel}
  \lean{Riemannian.Hyperbolic.hyperbolic_chartChristoffel, Riemannian.Hyperbolic.hyperbolic_partialDeriv_chartChristoffel, Riemannian.Hyperbolic.hyperbolic_partialDeriv_chartGramOnE}
  \uses{lem:dc-ch8-3-hyperbolic-chart-gram}
  \leanok
  In the abstract chart frame the Christoffel symbols of the hyperbolic metric are,
  in closed form,
  $$\Gamma^k_{ij}(y)=-f_i\delta^k_j-f_j\delta^k_i+d^k B_{ij},\qquad
    f_i=\frac{c_i}{y_n},\ c_i=(\text{finBasis}_i)_n,\ d^k=\frac{\sum_\ell B^{k\ell}c_\ell}{y_n},$$
  the constant abstract-basis Gram $B$ (\cref{lem:dc-ch8-3-hyperbolic-chart-gram})
  replacing $\delta$ in the abstract frame.
  Because every term carries the single factor $1/y_n$, one may write
  $\Gamma^k_{ij}(y)=K^k_{ij}\,(y_n)^{-1}$ with $K^k_{ij}$ the ($y$-independent)
  numerator, so the abstract-chart-frame partial derivatives are
  $$\partial_m\Gamma^k_{ij}(y)=K^k_{ij}\cdot\bigl(-c_m\,(y_n)^{-2}\bigr).$$
  This is exactly the $\Gamma$ and $\partial\Gamma$ data consumed by the chart-frame
  curvature expansion \cref{lem:dc-ch8-3-chart-curvature}.
\end{lemma}
\begin{proof}
  \leanok
\sketch
  The chart Gram is $G_{ij}(y)=(y_n^2)^{-1}B_{ij}$ and its inverse
  $G^{k\ell}(y)=y_n^2 B^{k\ell}$ (\cref{lem:dc-ch8-3-hyperbolic-chart-gram}). The
  distinguished coordinate $y_n=y_e$ has derivative
  $\partial_m y_n=(\text{finBasis}_m)_n=c_m$ along $\text{finBasis}_m$, so
  $\partial_m\bigl((y_n^2)^{-1}\bigr)=-2\,y_n^{-3}c_m$ and hence
  $\partial_m G_{ij}=B_{ij}\cdot(-2\,y_n^{-3}c_m)$. Substituting into the Koszul
  formula
  $\Gamma^k_{ij}=\tfrac12\sum_\ell G^{k\ell}(\partial_iG_{j\ell}+\partial_jG_{i\ell}-\partial_\ell G_{ij})$
  and contracting $\sum_\ell B^{k\ell}B_{\ell j}=\delta^k_j$,
  $\sum_\ell B^{k\ell}B_{\ell i}=\delta^k_i$ (from $B^{-1}B=1$) collapses the factors
  $y_n^2\cdot y_n^{-3}=y_n^{-1}$, giving the closed form. Since only the $1/y_n$ factor
  depends on $y$, $\partial_m\Gamma^k_{ij}=K^k_{ij}\,\partial_m\bigl(y_n^{-1}\bigr)
  =K^k_{ij}\,(-c_m y_n^{-2})$.
\end{proof}

\begin{lemma}[inverse-Gram/coordinate identity $\sum_{ij}B^{ij}c_ic_j=1$]
  \label{lem:dc-ch8-3-inverse-gram-coord}
  \lean{Riemannian.Hyperbolic.finBasis_inv_gram_coord}
  \uses{lem:dc-ch8-3-hyperbolic-chart-gram}
  \leanok
  For the constant abstract-basis Gram matrix
  $B_{ij}=\langle\text{finBasis}_i,\text{finBasis}_j\rangle$ with inverse
  $B^{ij}=(B^{-1})_{ij}$, and $c_i=(\text{finBasis}_i)_n$ the $n$-th coordinate of
  the $i$-th basis vector,
  $$\sum_{i,j} B^{ij}\,c_i\,c_j = 1.$$
  Equivalently $c_i=\langle\text{finBasis}_i, w_n\rangle$ for the standard unit
  vector $w_n=\text{basisFun}_n$, and the $B^{-1}$-contraction reconstructs
  $\langle w_n, w_n\rangle=1$.
\end{lemma}
\begin{proof}
  \leanok
\sketch
  Write $w_n=\sum_i a_i\,\text{finBasis}_i$ in the abstract basis. Then
  $c_i=\langle\text{finBasis}_i,w_n\rangle=\sum_j
  B_{ij}a_j$, so $\sum_j B^{ij}c_j=(B^{-1}B a)_i=a_i$ and
  $\sum_{ij}B^{ij}c_ic_j=\sum_i a_i c_i=\langle w_n,w_n\rangle=1$.
\end{proof}

\begin{lemma}[the pointwise $(0,4)$ curvature form and its chart-frame value]
  \label{lem:dc-ch8-3-pointwise-curvature-form}
  \lean{Riemannian.AffineConnection.curvatureFormAt, Riemannian.AffineConnection.isAlgCurvatureForm_curvatureFormAt, Riemannian.leviCivita_curvatureFormAt_chartFrame}
  \uses{lem:dc-ch8-3-chart-curvature}
  \leanok
  The curvature operator $R(X,Y)Z$ is a tensor: its value at $p$ depends only on
  $X(p),Y(p),Z(p)$. Lowering by the
  metric gives the pointwise $(0,4)$ form $\langle R(x,y)z,t\rangle_g$ on
  $T_pM$, which for the Levi-Civita connection is an \emph{algebraic curvature
  form} (multilinear, antisymmetric in each pair, and satisfying first Bianchi).
  In the chart frame $\partial_a=\text{finBasis}_a$ its value is the
  coordinate curvature coefficient contracted with the chart Gram,
  $$\langle R(\partial_i,\partial_j)\partial_k,\partial_\ell\rangle_g
    = \sum_m R^m_{ijk}\,G_{m\ell},\qquad
    R^m_{ijk}=\partial_j\Gamma^m_{ik}-\partial_i\Gamma^m_{jk}
      +\sum_s(\Gamma^s_{ik}\Gamma^m_{js}-\Gamma^s_{jk}\Gamma^m_{is}),$$
  the metric-lowered avatar of \cref{lem:dc-ch8-3-chart-curvature}. This bridges
  the \emph{intrinsic} manifold curvature to the coordinate Christoffel data, and
  is the input to the all-planes identification \cref{cor:dc-ch4-3-5}: it makes the
  chart coefficients the components of an algebraic curvature form on $T_pM$.
\end{lemma}
\begin{proof}
  \leanok
\sketch
  Tensoriality is the pointwise reading of the $\mathcal D(M)$-multilinearity of
  $R$: a field vanishing at $p$ decomposes as $\sum_i f_i W_i+\tau$ with
  $f_i(p)=0$ and $\tau$ vanishing near $p$, and homogeneity of $R$ in that slot
  kills both parts, so $R(\cdot)$ against a field vanishing at $p$ vanishes at
  $p$; hence $R(X,Y)Z|_p$ depends only on the values at $p$. The
  algebraic-curvature-form axioms transport slotwise from the field-level
  symmetries (Prop.~2.5, via the Levi-Civita property). The chart-frame value
  pairs the frame expansion of $R(\partial_i,\partial_j)\partial_k$
  (\cref{lem:dc-ch8-3-chart-curvature}) against $\partial_\ell$, distributing the
  metric over the finite sum.
\end{proof}

\begin{proposition}[$H^n$ has constant sectional curvature $-1$]
  \label{prop:dc-ch8-3-const-curv}
  \lean{Riemannian.Hyperbolic.hyperbolic_isConstantCurvature, Riemannian.Hyperbolic.hyperbolicMetric_sectionalCurvature_eq_neg_one, Riemannian.Hyperbolic.hyperbolic_curvatureFormAt_eq, Riemannian.Hyperbolic.hyperbolic_curvatureFormAt_frame}
  \uses{lem:dc-ch8-3-conformal-curvature, lem:dc-ch8-3-chart-curvature, lem:dc-ch8-3-pointwise-curvature-form, lem:dc-ch8-3-hyperbolic-chart-gram, lem:dc-ch8-3-hyperbolic-christoffel, lem:dc-ch8-3-inverse-gram-coord, cor:dc-ch4-3-5, def:dc-ch8-3-hyperbolic-space}
  \leanok
  The hyperbolic space $H^n$, with the metric $g_{ij}=\delta_{ij}/x_n^2$, has
  constant sectional curvature equal to $-1$: for every point and every
  $2$-plane $\sigma\subset T_pH^n$, $K(p,\sigma)=-1$.
\end{proposition}
\begin{proof}
  \leanok
\sketch
  Work in the abstract chart frame $\partial_a=\text{finBasis}_a$ of $T_pH^n$. By
  \cref{lem:dc-ch8-3-pointwise-curvature-form} the pointwise $(0,4)$ curvature form
  of the Levi-Civita connection is, in this frame,
  $\langle R(\partial_i,\partial_j)\partial_k,\partial_\ell\rangle_g=\sum_m
  R^m_{ijk}G_{m\ell}$, with $R^m_{ijk}=\partial_j\Gamma^m_{ik}-\partial_i\Gamma^m_{jk}
  +\sum_s(\Gamma^s_{ik}\Gamma^m_{js}-\Gamma^s_{jk}\Gamma^m_{is})$ and $G_{m\ell}$
  the chart Gram matrix. Insert the closed forms
  \cref{lem:dc-ch8-3-hyperbolic-christoffel}, $\Gamma^k_{ij}(y)=y_n^{-1}K^k_{ij}$
  and $\partial_m\Gamma^k_{ij}(y)=-c_m\,y_n^{-2}K^k_{ij}$ with the \emph{constant}
  coefficient $K^k_{ij}=-c_i\delta^k_j-c_j\delta^k_i+D^kB_{ij}$
  ($c_i=(\text{finBasis}_i)_n$, $D^k=\sum_\ell B^{k\ell}c_\ell$), and the chart Gram
  \cref{lem:dc-ch8-3-hyperbolic-chart-gram}, $G_{m\ell}=y_n^{-2}B_{m\ell}$ with $B$
  the constant abstract-basis Gram matrix. Every term carries a factor $y_n^{-4}$,
  which cancels against the right-hand side, and the identity reduces to the purely
  finite contraction $\sum_m R^m_{ijk}B_{m\ell}=-(B_{ik}B_{j\ell}-B_{jk}B_{i\ell})$
  in the constants $c,B,D$. Expanding the Christoffel products and contracting each
  index sum with the ``lowered'' identity $\sum_m K^m_{ab}B_{m\ell}=-c_aB_{b\ell}
  -c_bB_{a\ell}+c_\ell B_{ab}$ (which uses $\sum_m D^mB_{m\ell}=c_\ell$), every
  monomial cancels except $-(B_{ik}B_{j\ell}-B_{jk}B_{i\ell})$, the last step using
  the single nontrivial input $\sum_{ij}B^{ij}c_ic_j=1$
  (\cref{lem:dc-ch8-3-inverse-gram-coord}). Hence
  $\langle R(\partial_i,\partial_j)\partial_k,\partial_\ell\rangle_g
  =-(G_{ik}G_{j\ell}-G_{jk}G_{i\ell})$ on the chart frame. Both the pointwise
  curvature form and $-1$ times the standard curvature form
  $\langle x,z\rangle\langle y,t\rangle-\langle y,z\rangle\langle x,t\rangle$ are
  algebraic curvature forms (\cref{lem:dc-ch8-3-pointwise-curvature-form}) that
  agree on this basis frame, so by the basis-determination step of
  \cref{cor:dc-ch4-3-5} they agree on all of $T_pH^n$:
  $\langle R(x,y)z,t\rangle_g=-(\langle x,z\rangle_g\langle y,t\rangle_g
  -\langle y,z\rangle_g\langle x,t\rangle_g)$. This is constant curvature $-1$ at the
  field level; dividing the diagonal $\langle R(x,y)x,y\rangle_g=-|x\wedge y|^2$ by
  $|x\wedge y|^2$ gives $K(p,\sigma)=-1$ for every $2$-plane
  $\sigma=\mathrm{span}\{x,y\}$.
\end{proof}

\begin{lemma}[Christoffel contraction of the hyperbolic metric, closed form]
  \label{lem:dc-ch8-3-christoffel-contraction}
  \lean{Riemannian.Hyperbolic.hyperbolic_chartChristoffelContraction_eq}
  \uses{lem:dc-ch8-3-hyperbolic-christoffel, lem:dc-ch8-3-inverse-gram-coord}
  \leanok
  Contracting the closed-form chart Christoffel symbols
  $\Gamma^k_{ij}=-\delta_{jk}f_i-\delta_{ki}f_j+\delta_{ij}f_k$
  ($f_i=(\partial_i)_n/x_n$) of \cref{lem:dc-ch8-3-hyperbolic-christoffel} against
  two tangent vectors $v,w\in T_pH^n$ collapses, via the inverse-Gram/coordinate
  identity \cref{lem:dc-ch8-3-inverse-gram-coord}, to the basis-independent
  Euclidean-valued expression

  $$
  \Gamma(v,w)(y)=\nabla_v w
    =\frac{1}{y_n}\Big(\langle v,w\rangle\, e_n-v_n\,w-w_n\,v\Big),
  $$

  where $y_n$ is the distinguished coordinate of the base point, $v_n=\langle v,e_n\rangle$,
  $e_n$ is the distinguished unit vector, and $\langle\,,\rangle$ is the ambient
  Euclidean inner product.
\end{lemma}
\begin{proof}
  \leanok
\sketch
  In the abstract chart frame the contraction is
  $\Gamma(v,w)(y)=\sum_k\big(\sum_{ij}\Gamma^k_{ij}(y)\,v^i w^j\big)\partial_k$ with
  $v^i,w^j$ the frame coordinates. Substituting the closed form
  \cref{lem:dc-ch8-3-hyperbolic-christoffel} and summing the two Kronecker terms
  gives $-\big(\sum_i (\partial_i)_n v^i\big)w-\big(\sum_j(\partial_j)_n w^j\big)v
  =-v_n w-w_n v$ (using $\sum_i(\partial_i)_n v^i=v_n$), while the $\delta_{ij}f_k$
  term contracts to $\langle v,w\rangle\sum_k(\sum_\ell B^{k\ell}c_\ell)\partial_k
  =\langle v,w\rangle\,e_n$ by \cref{lem:dc-ch8-3-inverse-gram-coord}; the common
  factor $1/y_n$ is carried throughout.
\end{proof}

\begin{lemma}[geodesic equation of $H^n$ in Euclidean form]
  \label{lem:dc-ch8-3-geodesic-equation}
  \lean{Riemannian.Hyperbolic.hyperbolic_isGeodesic_of}
  \uses{lem:dc-ch8-3-christoffel-contraction}
  \leanok
  A twice-differentiable curve $u:\mathbb R\to\mathbb R^n$ staying in the
  half-space ($u_n(t)>0$) whose acceleration satisfies the ordinary differential
  equation

  $$
  u''(t)+\frac{1}{u_n(t)}\Big(\langle u'(t),u'(t)\rangle\, e_n-2\,u'_n(t)\,u'(t)\Big)=0
  $$

  is a geodesic of $H^n$. This is the geodesic equation $u''+\Gamma(u',u')=0$ read
  through \cref{lem:dc-ch8-3-christoffel-contraction}; because the chart of the open
  half-space is the inclusion, the moving-foot chart curve of the intrinsic geodesic
  predicate is literally $u$.
\end{lemma}
\begin{proof}
  \leanok
\sketch
  For an open subset of $\mathbb R^n$ the extended chart at any base point is the
  inclusion, so the chart-local curve $s\mapsto\varphi_{\gamma t}(\gamma s)$ equals
  $u$ for every $t$; its first and second derivatives are $u'$ and $u''$. The
  intrinsic moving-foot geodesic condition at time $t$ is
  $u''(t)+\Gamma(u'(t),u'(t))(u(t))=0$, and substituting the closed form
  \cref{lem:dc-ch8-3-christoffel-contraction} (with $v=w=u'(t)$, so the two mixed
  terms coincide) gives exactly the displayed equation.
\end{proof}

\begin{lemma}[vertical lines are geodesics of $H^n$]
  \label{lem:dc-ch8-3-vertical-geodesic}
  \lean{Riemannian.Hyperbolic.hyperbolic_vertical_isGeodesic}
  \uses{lem:dc-ch8-3-geodesic-equation}
  \leanok
  For a horizontal base point $c$ (with $c_n=0$) and a constant $a>0$, the affinely
  parametrised vertical line

  $$
  t\longmapsto c+(a\,e^{t})\,e_n,\qquad t\in\mathbb R,
  $$

  perpendicular to the boundary hyperplane $\{x_n=0\}$, is a geodesic of $H^n$. Its
  distinguished coordinate $u_n(t)=a\,e^{t}$ satisfies $u_n''\,u_n=(u_n')^2$.
\end{lemma}
\begin{proof}
  \leanok
\sketch
  Here $u(t)=c+(a e^{t})e_n$, $u'(t)=u''(t)=(a e^{t})e_n$, $u_n(t)=u'_n(t)=a e^{t}$
  and $\langle u',u'\rangle=(a e^{t})^2$. The geodesic equation
  \cref{lem:dc-ch8-3-geodesic-equation} reads
  $(a e^{t})e_n+(a e^{t})^{-1}\big((a e^{t})^2 e_n-2(a e^{t})^2 e_n\big)
  =(a e^{t})e_n-(a e^{t})e_n=0$, which holds identically.
\end{proof}

\begin{lemma}[semicircles perpendicular to $\partial H^n$ are geodesics of $H^n$]
  \label{lem:dc-ch8-3-semicircle-geodesic}
  \lean{Riemannian.Hyperbolic.hyperbolic_semicircle_isGeodesic}
  \uses{lem:dc-ch8-3-geodesic-equation}
  \leanok
  For a center $m$ on the boundary hyperplane ($m_n=0$), a unit horizontal vector
  $\hat u$ ($\hat u_n=0$, $\langle\hat u,\hat u\rangle=1$) and a radius $r>0$, the
  affinely parametrised semicircle

  $$
  t\longmapsto m+(r\tanh t)\,\hat u+(r\,\mathrm{sech}\,t)\,e_n,\qquad t\in\mathbb R,
  $$

  (with $\tanh t=\sinh t/\cosh t$, $\mathrm{sech}\,t=1/\cosh t$), perpendicular to
  the boundary $\{x_n=0\}$ and centered on it, is a geodesic of $H^n$.
\end{lemma}
\begin{proof}
  \leanok
\sketch
  With $u(t)=m+(r\tanh t)\hat u+(r\,\mathrm{sech}\,t)e_n$ one has
  $u'(t)=(r\,\mathrm{sech}^2 t)\hat u-(r\tanh t\,\mathrm{sech}\,t)e_n$; using
  $\mathrm{sech}'=-\tanh\,\mathrm{sech}$, $\tanh'=\mathrm{sech}^2$, and
  $\cosh^2-\sinh^2=1$ one computes $u_n=r\,\mathrm{sech}\,t$,
  $\langle u',u'\rangle=r^2\mathrm{sech}^2 t$, and both the $\hat u$- and
  $e_n$-components of $u''+u_n^{-1}(\langle u',u'\rangle e_n-2u'_n u')$ vanish
  identically. Hence \cref{lem:dc-ch8-3-geodesic-equation} applies.
\end{proof}

\begin{proposition}[geodesics of $H^n$]
  \label{prop:dc-ch8-3-1}
  \lean{Riemannian.Hyperbolic.hyperbolic_geodesic_classification}
  \uses{ex:dc-ch3-3-10, lem:dc-ch8-3-geodesic-equation, lem:dc-ch8-3-vertical-geodesic, lem:dc-ch8-3-semicircle-geodesic}
  \leanok
  \dcref{ch8:3.1}
  The straight lines perpendicular to the hyperplane $x_n=0$, and the circles of
  $H^n$ whose planes are perpendicular to the hyperplane $x_n=0$ and whose centers
  are in this hyperplane, are the geodesics of $H^n$.
\end{proposition}
\begin{proof}
  \leanok
\sketch
  That each stated curve is a geodesic is the forward direction, already established:
  the vertical lines by \cref{lem:dc-ch8-3-vertical-geodesic} and the semicircles
  perpendicular to $\{x_n=0\}$ with centre on it by
  \cref{lem:dc-ch8-3-semicircle-geodesic}, both through the geodesic equation
  \cref{lem:dc-ch8-3-geodesic-equation}. It remains to prove the converse inclusion:
  that these are \emph{all} the geodesics.

  Let $\gamma$ be a geodesic of $H^n$, and read it through the inclusion chart of the
  open half-space, so that $\gamma$ is an ambient curve with $\gamma(0)=p$ and chart
  velocity $\gamma'(0)=v\in\mathbb R^n$. Decompose $v=v_h+v_n\,e_n$, where
  $e_n=\partial/\partial x_n$, $v_n=\langle v,e_n\rangle$, and $v_h$ is the horizontal
  part (so $\langle v_h,e_n\rangle=0$). Through $(p,v)$ passes exactly one member of
  the stated family:
  \begin{itemize}
    \item if $v_h=0$, the reparametrised vertical line
      $s\mapsto c+p_n\,e^{Bs}\,e_n$, with $c=p-p_n e_n$ (so $c_n=0$), $p_n=\langle
      p,e_n\rangle>0$ and $B=v_n/p_n$;
    \item if $v_h\neq0$, the semicircle
      $s\mapsto m+r\tanh(\sigma s+s_0)\,\hat u+r\,\operatorname{sech}(\sigma s+s_0)\,e_n$
      in the plane $\{e_n,\hat u\}$, with $\hat u=v_h/|v_h|$, and centre, radius,
      speed and phase determined by $(p,v)$: writing $s_0=\operatorname{arsinh}(-v_n/|v_h|)$,
      $r=p_n\cosh s_0$, $\sigma=|v_h|\cosh s_0/p_n$ and
      $m=(p-p_n e_n)+\bigl(p_n v_n/|v_h|\bigr)\hat u$, so that $m_n=0$, $\hat u_n=0$,
      $|\hat u|=1$, $r>0$.
  \end{itemize}
  A direct computation (using $\sinh s_0=-v_n/|v_h|$ and
  $\cosh^2-\sinh^2=1$) shows that in either case this family member is a geodesic
  passing through $p$ with chart velocity $v$ at $s=0$. Uniqueness of geodesics with
  prescribed initial data makes $\gamma$ that vertical line or semicircle.
\end{proof}

\section{Space forms}

\begin{theorem}[universal covering of a space of constant curvature]
  \label{thm:dc-ch8-4-1}
  \notready
  \dcref{ch8:4.1}
  \notready
  Let $M^n$ be a complete Riemannian manifold with constant sectional curvature $K$.
  Then the universal covering $\tilde{M}$ of $M$, with the covering metric, is
  isometric to:
  \begin{enumerate}
    \item[(a)] $H^n$, if $K=-1$;
    \item[(b)] $\mathbb{R}^n$, if $K=0$;
    \item[(c)] $S^n$, if $K=1$.
  \end{enumerate}
\end{theorem}
\begin{proof}
  \uses{thm:dc-ch8-2-1,lem:dc-ch8-4-2,prop:dc-ch8-3-const-curv,
    lem:dc-ch7-3-1-covering-diffeo,lem:dc-ch7-3-1-covering-simplyconnected}
\sketch
  $\tilde{M}$ is a simply connected, complete Riemannian
  manifold with constant sectional curvature $K$. Consider first the cases (a) and
  (b) and denote, for convenience, both $H^n$ and $\mathbb{R}^n$ by $\triangle$. Fix
  points $p\in\triangle$, $\tilde{p}\in\tilde{M}$ and a linear isometry
  $i:T_p(\triangle)\to T_{\tilde{p}}(\tilde{M})$. Consider the map
  $f=\exp_{\tilde{p}}\circ\,i\circ\exp_p^{-1}:\triangle\to\tilde{M}$. Since $\triangle$
  and $\tilde{M}$ are complete with non-positive sectional curvature, $f$ is
  well-defined. By \cref{thm:dc-ch8-2-1} it is a local isometry, and
  \cref{lem:dc-ch7-3-1-covering-diffeo} makes it a diffeomorphism.

  For case (c), fix $p\in S^n$, $\tilde{p}\in\tilde{M}$ and a linear isometry
  $i:T_p(S^n)\to T_{\tilde{p}}(\tilde{M})$. Let $q\in S^n$ be the antipodal point of
  $p$ and define $f=\exp_{\tilde{p}}\circ\,i\circ\exp_p^{-1}:S^n-\{q\}\to\tilde{M}$.
  From the Theorem of Cartan, $f$ is a local isometry. Choose $p'\in S^n$,
  $p'\neq p$, $p'\neq q$, set $\tilde{p}'=f(p')$, $i'=df_{p'}$ and define
  $f'=\exp_{\tilde{p}'}\circ\,i'\circ\exp_{p'}^{-1}:S^n-\{q'\}\to\tilde{M}$, where
  $q'$ is the antipodal point of $p'$. On the connected set
  $W=S^n-(\{q\}\cup\{q'\})$ we have $p'\in W$ and $f(p')=\tilde{p}'=f'(p')$,
  $df_{p'}=i'=df'_{p'}$, so by \cref{lem:dc-ch8-4-2}, $f=f'$ on $W$. Define
  $g:S^n\to\tilde{M}$ by $g(r)=f(r)$ if $r\in S^n-\{q\}$ and $g(r)=f'(r)$ if
  $r\in S^n-\{q'\}$. Then $g$ is a local isometry, hence a local diffeomorphism. By
  the compactness of $S^n$, $g$ is a covering map, and since $\tilde{M}$ is simply
  connected, $g$ is a diffeomorphism. Therefore $g$ is an isometry.
\end{proof}

\begin{lemma}[a local isometry pushes geodesics forward]
  \label{lem:dc-ch8-4-2-geodesic-push}
  \lean{Riemannian.Geodesic.isGeodesicOn_comp_of_isLocalDiffeomorph}
  \uses{def:dc-ch1-2-3, def:dc-ch3-2-1}
  Let $f:M\to N$ be a local isometry and let $\gamma$ be a geodesic of $M$ defined on
  an open set $s$ of times. Then $f\circ\gamma$ is a geodesic of $N$ on $s$.
\end{lemma}
\begin{proof}
  \leanok
\sketch
  Being a local isometry, $f$ identifies the metric of $M$ with the pullback $f^*g_N$,
  so $\gamma$ satisfies the $f^*g_N$-geodesic equation. The geodesic equation is a
  pointwise condition, and under $f$ the geodesic operator transforms by the
  differential $df$ alone: from $df\big(\gamma''+\Gamma^{f^*g_N}(\gamma',\gamma')\big)
  =(f\circ\gamma)''+\Gamma^{g_N}\big((f\circ\gamma)',(f\circ\gamma)'\big)$ and the
  vanishing of the left-hand argument one reads off that $f\circ\gamma$ satisfies the
  $g_N$-geodesic equation at each time of $s$.
\end{proof}

\begin{lemma}[the agreement locus is open and closed]
  \label{lem:dc-ch8-4-2-agree-locus}
  \lean{Riemannian.AgreeLocus, Riemannian.mem_agreeLocus_iff, Riemannian.isClosed_agreeLocus, Riemannian.eq_of_pointDatum_of_preconnected}
  Let $f_1,f_2:M\to N$ be differentiable, $M$ connected, and let

  $$
  A=\{x\in M;\ f_1(x)=f_2(x)\ \text{and}\ (df_1)_x=(df_2)_x\}
  $$

  be their \emph{agreement locus}. Then $A$ is closed. If moreover $A$ is open, for
  instance if agreement to first order at a point always propagates to a neighborhood
  of that point, and $A\neq\emptyset$, then $f_1=f_2$.
\end{lemma}
\begin{proof}
  \leanok
\sketch
  A point $x$ lies in $A$ exactly when the tangent maps $df_1$ and $df_2$ agree over
  $x$, since the tangent map of $f$ sends $(x,v)$ to $(f(x),(df)_x v)$: the base
  components give $f_1(x)=f_2(x)$ and the fibre components $(df_1)_x=(df_2)_x$. So $A$
  is the complement of the image, under the bundle projection $TM\to M$, of the locus
  where the two tangent maps differ. The tangent maps are continuous and $TN$ is
  Hausdorff, so that locus is open; the bundle projection is an open map, so its image
  is open and $A$ is closed. If $A$ is also open and nonempty then, $M$ being
  connected, $A=M$, i.e. $f_1=f_2$.
\end{proof}

\begin{lemma}[first-order agreement propagates]
  \label{lem:dc-ch8-4-2-propagation}
  \lean{Riemannian.eventuallyEq_of_pointDatum}
  \uses{lem:dc-ch8-4-2-geodesic-push, prop:dc-ch3-2-9, def:dc-ch3-expmap, lem:dc-ch3-expmap-ray, thm:dc-ch3-2-2}
  \leanok
  Let $f_1,f_2:M\to N$ be local isometries and let $q\in M$ satisfy $f_1(q)=f_2(q)$ and
  $(df_1)_q=(df_2)_q$. Then $f_1=f_2$ on a neighborhood of $q$.
\end{lemma}
\begin{proof}
  \leanok
  \uses{lem:dc-ch8-4-2-geodesic-push,prop:dc-ch3-2-9,def:dc-ch3-expmap,
    lem:dc-ch3-expmap-ray,thm:dc-ch3-2-2}
\sketch
  By \cref{prop:dc-ch3-2-9}, $\exp_q$ carries a ball $B_\varepsilon(0)\subset T_qM$
  diffeomorphically onto an open neighborhood of $q$; so it suffices to prove
  $f_1(\exp_q u)=f_2(\exp_q u)$ for all sufficiently small $u$. Fix such a $u$ and let
  $\gamma_u(t)=\exp_q(tu)$, a geodesic of $M$ on an interval $(-b,b)$ with $b>1$, with
  $\gamma_u(0)=q$ and $\gamma_u'(0)=u$ (\cref{lem:dc-ch3-expmap-ray}). By
  \cref{lem:dc-ch8-4-2-geodesic-push} both $f_1\circ\gamma_u$ and $f_2\circ\gamma_u$ are
  geodesics of $N$ on $(-b,b)$. At $t=0$ they have the same initial point, since
  $f_1(q)=f_2(q)$, and the same initial velocity, since
  $(df_1)_q u=(df_2)_q u$. By uniqueness of geodesics with given initial conditions
  (\cref{thm:dc-ch3-2-2}) they coincide on the connected interval $(-b,b)$; evaluating
  at $t=1$ gives $f_1(\exp_q u)=f_2(\exp_q u)$.
\end{proof}

\begin{lemma}[uniqueness of a local isometry from a point datum]
  \label{lem:dc-ch8-4-2}
  \lean{Riemannian.eq_of_pointDatum_of_localIsometry}
  \uses{def:dc-ch1-2-3}
  \leanok
  \dcref{ch8:4.2}
  Let $f_i:M\to N$, $i=1,2$, be two local isometries of the (connected) Riemannian
  manifold $M$ to the Riemannian manifold $N$, both of dimension $n$ and modelled on a
  common Euclidean space. Suppose that there exists a point $p\in M$ such that
  $f_1(p)=f_2(p)$ and $(df_1)_p=(df_2)_p$. Then $f_1=f_2$.
\end{lemma}
\begin{proof}
  \leanok
  \uses{lem:dc-ch8-4-2-agree-locus,lem:dc-ch8-4-2-propagation}
\sketch
  Consider the agreement locus
  $A=\{x\in M;\ f_1(x)=f_2(x)\ \text{and}\ (df_1)_x=(df_2)_x\}$. It contains $p$, so it
  is nonempty; it is closed by \cref{lem:dc-ch8-4-2-agree-locus}; and it is open,
  because by \cref{lem:dc-ch8-4-2-propagation} each of its points has a neighborhood on
  which $f_1$ and $f_2$ agree, and on such a neighborhood the differentials agree too.
  Since $M$ is connected, $A=M$, that is, $f_1=f_2$.
\end{proof}

The complete constant-curvature manifolds are called \emph{space forms}. An action
of $G$ on $M$ is free and totally discontinuous when $gx=x$ implies $g=e$ and every
$x$ has a neighborhood $U$ with $g(U)\cap U=\varnothing$ for $g\ne e$. Then
$M\to M/G$ is a regular covering.

If $M$ is Riemannian and $\Gamma$ is a totally discontinuous subgroup of its
isometry group, $M/\Gamma$ has a Riemannian structure for which $\pi$ is a local
isometry. For $\tilde p\in\pi^{-1}(p)$ define

$$
\langle u,v\rangle=\langle d\pi^{-1}(u),d\pi^{-1}(v)\rangle_{\tilde{p}},
$$

The definition is independent of the lift because $\Gamma$ acts by isometries.
Completeness and constant curvature descend from $M$ to $M/\Gamma$.

\begin{proposition}[every space form is a quotient of a model]
  \label{prop:dc-ch8-4-3}
  \notready
  \dcref{ch8:4.3}
  \notready
  Let $M$ be a complete Riemannian manifold with constant sectional curvature $K$
  ($1,0,-1$). Then $M$ is isometric to $\tilde{M}/\Gamma$, where $\tilde{M}$ is $S^n$
  (if $K=1$), $\mathbb{R}^n$ (if $K=0$) or $H^n$ (if $K=-1$), $\Gamma$ is a subgroup
  of the group of isometries of $\tilde{M}$ which acts in a totally discontinuous
  manner on $\tilde{M}$, and the metric on $\tilde{M}/\Gamma$ is induced from the
  covering $\pi:\tilde{M}\to\tilde{M}/\Gamma$.
\end{proposition}
\begin{proof}
  \uses{thm:dc-ch8-4-1,lem:dc-ch7-3-1-covering-simplyconnected}
\sketch
  Consider the universal covering $p:\tilde{M}\to M$ with the covering
  metric. Let $\Gamma$ be the group of covering transformations of $p$. Then
  $\Gamma$ is a subgroup of the isometries of $\tilde{M}$ and acts in a totally
  discontinuous manner. Introduce on $\tilde{M}/\Gamma$ the metric induced by
  $\pi:\tilde{M}\to\tilde{M}/\Gamma$. Since the covering $p$ is regular, given
  $\tilde{x},\tilde{y}\in\tilde{M}$, $p(\tilde{x})=p(\tilde{y})$ iff
  $\Gamma\tilde{x}=\Gamma\tilde{y}$, equivalently $\pi(\tilde{x})=\pi(\tilde{y})$.
  The equivalence classes given by $p$ and $\pi$ are the same, which induces a
  bijection $\xi:M\to\tilde{M}/\Gamma$ such that $\pi=\xi\circ p$. Since $\pi$ and
  $p$ are local isometries, $\xi$ is a local isometry, and being a bijection, is an
  isometry of $M$ onto $\tilde{M}/\Gamma$.
\end{proof}

\begin{proposition}[even-dimensional space forms of $K=1$]
  \label{prop:dc-ch8-4-4}
  \notready
  \dcref{ch8:4.4}
  \notready
  Let $M^n$ be a complete Riemannian manifold of even dimension $n=2m$, with constant
  sectional curvature $K=1$. Then $M^n$ is isometric to the sphere $S^n$ or the real
  projective space $P^n$ of the same dimension.
\end{proposition}
\begin{proof}
  \uses{prop:dc-ch8-4-3}
\sketch
  The orthogonal group $O(2m+1)$ is the (transitive) group of isometries of
  $S^n$. Let $\Gamma$ be a subgroup that acts in a totally discontinuous manner on
  $S^n$. If $\gamma\in\Gamma$ has determinant $+1$, there exists an eigenvalue of
  $\gamma$ equal to $1$; then $\gamma$ has a fixed point, which implies $\gamma=e$.
  If $\tilde{\gamma}\in\Gamma$ has determinant $-1$, then $\tilde{\gamma}^2$ has
  determinant $1$, therefore $\tilde{\gamma}^2=e$. The eigenvalues of
  $\tilde{\gamma}$ are $1$ or $-1$. If $1$ is an eigenvalue then $\tilde{\gamma}=e$,
  contradicting $\det\tilde{\gamma}=-1$. As a result $\tilde{\gamma}=-e$, hence
  $\Gamma=\{e,-e\}$. Using \cref{prop:dc-ch8-4-3} we obtain that $M^n$ is either $S^n$, if
  $\Gamma=\{e\}$, or $P^n$, if $\Gamma=\{e,-e\}$.
\end{proof}

\begin{proposition}[metric of constant negative curvature on surfaces of genus $>1$]
  \label{prop:dc-ch8-4-5}
  \notready
  \dcref{ch8:4.5}
  \notready
  Every compact orientable surface of genus $p>1$ can be provided with a metric of
  constant negative curvature.
\end{proposition}
\begin{proof}
\sketch
  In $H^2$ take a closed geodesic polygon $P$ with $4p$ equal sides and identify
  paired sides to form a surface $M^2$. Let $\Gamma$ be
  the subgroup of isometries of $H^2$ generated by the isometries that identify the
  sides of $P$. The translates of $P$ tile $H^2$ iff
  the sum $\alpha$ of the interior angles at the vertices of $P$ equals $2\pi$; in
  this case $M^2=H^2/\Gamma$ has a metric of constant curvature $-1$. We assert that
  it is possible to find a polygon $P$ satisfying $\alpha=2\pi$. By the Gauss–Bonnet
  Theorem, $\alpha=-A+2\pi(2p-1)$, where $A$ is the area of $P$ in the hyperbolic
  metric. If $P$ is arbitrarily small, $\alpha$ is arbitrarily close to $2\pi(2p-1)$;
  if we enlarge the area $A$, $\alpha$ becomes arbitrarily small. By continuously
  deforming $P$, there exists a geodesic polygon such that $\alpha=2\pi$.

\end{proof}

\section{Isometries of the hyperbolic space; Theorem of Liouville}

For a differentiable map $f:U\subset\mathbb R^n\to\mathbb R^n$, conformality at
$p$ is equivalent to the existence of $\lambda(p)>0$ such that

$$
\langle df_p(v_1),df_p(v_2)\rangle=\lambda^2(p)\langle v_1,v_2\rangle,\qquad\lambda^2\neq 0.\qquad(3)
$$

Indeed, if $(3)$ holds then $|df_p(v)|^2=\lambda^2|v|^2$, hence

$$
\cos\angle(df_p(v_1),df_p(v_2))=\cos\angle(v_1,v_2),\qquad(4)
$$

so angles are preserved.

\begin{definition}[conformal map with a coefficient of conformality]
  \label{def:dc-ch8-5-conformal}
  \lean{Riemannian.IsConformalWithCoeff, Riemannian.IsConformalWithCoeff.isConformalMap, Riemannian.IsConformalWithCoeff.conformalAt}
  \leanok
  A map $f:U\subset\mathbb{R}^n\to\mathbb{R}^n$ is \emph{conformal at $p$ with
  coefficient of conformality $\lambda(p)>0$} if $f$ is differentiable at $p$ and
  its differential rescales every inner product by $\lambda^2$, i.e. condition
  $(3)$ holds:
  $\langle df_p(v_1),df_p(v_2)\rangle=\lambda^2(p)\langle v_1,v_2\rangle$ for all
  $v_1,v_2$. Equivalently $df_p$ is a conformal linear map, so $f$ is conformal at
  $p$ in the sense of preserving non-oriented angles.
\end{definition}

\begin{lemma}[a conformal map preserves norms up to $\lambda$ and preserves angles: $(3)\Rightarrow(4)$]
  \label{lem:dc-ch8-5-conformal-angle}
  \lean{Riemannian.IsConformalWithCoeff.norm_fderiv_eq, Riemannian.IsConformalWithCoeff.angle_fderiv_eq}
  \uses{def:dc-ch8-5-conformal}
  \leanok
  If $f$ is conformal at $p$ with coefficient $\lambda$, then $|df_p(v)|=\lambda|v|$
  for every $v$, and $df_p$ preserves the non-oriented angle between any two
  vectors: $\angle(df_p(v_1),df_p(v_2))=\angle(v_1,v_2)$.
\end{lemma}
\begin{proof}
  \leanok
\sketch
  From $(3)$ with $v_1=v_2=v$, $|df_p(v)|^2=\langle df_p(v),df_p(v)\rangle
  =\lambda^2\langle v,v\rangle=\lambda^2|v|^2$, so $|df_p(v)|=\lambda|v|$ since
  $\lambda>0$. Hence the cosine of the angle
  $\langle df_p(v_1),df_p(v_2)\rangle/(|df_p(v_1)|\,|df_p(v_2)|)
  =\lambda^2\langle v_1,v_2\rangle/(\lambda^2|v_1|\,|v_2|)
  =\langle v_1,v_2\rangle/(|v_1|\,|v_2|)$
  is unchanged, giving $(4)$.
\end{proof}

\begin{example}[Basic conformal transformations]
  \label{ex:dc-ch8-5-1}
  \lean{Riemannian.isConformalWithCoeff_isometry, Riemannian.isConformalWithCoeff_dilatation, Riemannian.isConformalWithCoeff_inversion, Riemannian.inversion_unit_involutive, Riemannian.inversion_unit_fixed_of_mem_sphere, Riemannian.inversion_unit_eq_lineMap, Riemannian.dist_inversion_unit_center}
  \uses{def:dc-ch8-5-conformal}
  \dcref{ch8:5.1}
  Euclidean isometries are conformal with coefficient $1$, and dilatations
  $f(p)=\lambda p$, $\lambda>0$, are conformal with coefficient $\lambda$. The
  inversion in the unit sphere centered at $p_0$,
  $$
  f(p)=\frac{p-p_0}{|p-p_0|^2}+p_0,\qquad p\in\mathbb{R}^n-\{p_0\},\qquad(5)
  $$
  is a conformal involution with coefficient $1/|p-p_0|^2$.
\end{example}
\begin{proof}
  \sketch
  The isometry and dilatation cases are immediate. For inversion,
  $$
  df_p(v)=\frac{v|p-p_0|^2-2\langle v,p-p_0\rangle(p-p_0)}{|p-p_0|^4},
  $$
  and direct expansion gives
  $|df_p(v)|^2=|v|^2/|p-p_0|^4$.
\end{proof}

\begin{lemma}[conformal maps are closed under composition]
  \label{lem:dc-ch8-5-conformal-comp}
  \lean{Riemannian.IsConformalWithCoeff.comp, Riemannian.isConformalWithCoeff_id}
  \uses{def:dc-ch8-5-conformal}
  \leanok
  If $f$ is conformal at $p$ with coefficient $\lambda$ and $g$ is conformal at
  $f(p)$ with coefficient $\mu$, then $g\circ f$ is conformal at $p$ with
  coefficient $\mu\lambda$. In particular the identity is conformal with
  coefficient $1$, and any composition of isometries, dilatations, and inversions,
  the class of maps in the Theorem of Liouville, is conformal.
\end{lemma}
\begin{proof}
  \leanok
\sketch
  By the chain rule $d(g\circ f)_p=dg_{f(p)}\circ df_p$. For all $v_1,v_2$,
  $\langle d(g\circ f)_p(v_1),d(g\circ f)_p(v_2)\rangle
  =\mu^2\langle df_p(v_1),df_p(v_2)\rangle
  =\mu^2\lambda^2\langle v_1,v_2\rangle=(\mu\lambda)^2\langle v_1,v_2\rangle$,
  and $\mu\lambda>0$.
\end{proof}

\begin{theorem}[Liouville]
  \label{thm:dc-ch8-5-2}
  \uses{lem:dc-ch8-5-differential-core}
  \notready
  \dcref{ch8:5.2}
  \notready
  Let $f:U\to\mathbb{R}^n$, $n\geq 3$, be a conformal transformation of an open set
  $U\subset\mathbb{R}^n$. Then $f$ is the restriction to $U$ of a composition of
  isometries, dilatations or inversions, at most one of each.
\end{theorem}
\begin{proof}
\sketch
  Let $a_1=(1,0,\dots,0),\dots,a_n=(0,\dots,0,1)$ be the canonical basis and
  $(x_1,\dots,x_n)$ the cartesian coordinates. Let $e_1,\dots,e_n$ be parallel
  differentiable vector fields on $U$ with $\langle e_i,e_j\rangle=\delta_{ij}$. If
  $\lambda$ is the coefficient of conformality of $f$,


  $$
  \langle df(e_i),df(e_k)\rangle=\lambda^2\delta_{ik},\qquad i,k=1,\dots,n.\qquad(6)
  $$


  Let $d^2 f$ be the second differential of $f$, a symmetric bilinear map with
  $d^2 f(a_i,a_j)=\partial^2 f/\partial x_i\partial x_j$. Differentiating $(6)$ with
  $i,j,k$ distinct and combining the three cyclic equations yields
  $\langle d^2 f(e_k,e_j),df(e_i)\rangle=0$ for $i,j,k$ distinct. Letting $i$ vary,
  $d^2 f(e_k,e_j)$ belongs to the plane generated by $df(e_j)$ and $df(e_k)$, so
  $d^2 f(e_k,e_j)=\mu\,df(e_k)+\nu\,df(e_j)$ with $\mu=d\lambda(e_j)/\lambda$,
  $\nu=d\lambda(e_k)/\lambda$, that is,


  $$
  d^2 f(e_k,e_j)=\frac{1}{\lambda}(df(e_k)\,d\lambda(e_j)+df(e_j)\,d\lambda(e_k)).\qquad(7)
  $$


  Put $\rho=1/\lambda$. Using $d(\rho f)=d\rho\,f+\rho\,df$ and $(7)$,


  $$
  d^2(\rho f)(e_k,e_j)=d^2\rho(e_k,e_j)f.\qquad(8)
  $$


  We claim $d^2\rho(e_k,e_j)=0$ for $k\neq j$. Computing the third differential
  $d^3(\rho f)$ via $(8)$ and using the symmetry of the first member in $i,j,k$, one
  obtains $d^2\rho(e_k,e_j)df(e_i)=d^2\rho(e_k,e_i)df(e_j)$; since $df(e_i)$ and
  $df(e_j)$ are linearly independent and $i,j,k$ are distinct, $d^2\rho(e_k,e_j)=0$
  for all $j\neq k$. Choosing $e_1,\dots,e_n$ to form a prescribed orthonormal basis
  at any $p$, this holds for every orthonormal basis; since
  $0=d^2\rho\big(\tfrac{e_j+e_k}{\sqrt 2},\tfrac{e_j-e_k}{\sqrt 2}\big)=\tfrac12\{d^2\rho(e_j,e_j)-d^2\rho(e_k,e_k)\}$,
  we get $d^2\rho(e_j,e_j)=d^2\rho(e_k,e_k)$ for $j\neq k$. In summary, for every
  $p\in U$ and any orthonormal basis, $d^2\rho(e_j,e_k)=\sigma\delta_{jk}$; in the
  canonical basis,


  $$
  \frac{\partial^2\rho}{\partial x_i\partial x_j}=\sigma\delta_{ij}.\qquad(9)
  $$


  Differentiating $(9)$ gives $\partial\sigma/\partial x_i=0$ ($i\neq j$), so
  $\sigma=\text{const.}$

  \emph{Case $\sigma\neq 0$.} Equation $(9)$ implies


  $$
  \rho=\frac{\sigma}{2}\sum x_i^2+\sigma\sum b_i x_i+c,\quad b_i,c\text{ constants},\qquad(10)
  $$


  which we write as $1/\lambda=\rho=a_1|p-p_0|^2+k_1$ with $a_1=\sigma/2$,
  $k_1=\text{const.}$, $p_0\in\mathbb{R}^n$. The proof in this case will be complete
  once we show $k_1=0$: if $k_1=0$, taking the inversion
  $g(p)=\frac{p-p_0}{|p-p_0|^2}+p_0$ and $h=g\circ f^{-1}$, $h$ is conformal with
  coefficient $a_1$, hence an isometry followed by a dilatation; thus
  $f=h^{-1}\circ g$ is an inversion followed by a dilatation, followed by an
  isometry. To show $k_1=0$, apply the argument to $f^{-1}$:
  $\lambda=a_2|f(p)-q_0|^2+k_2$, hence


  $$
  (a_1|p-p_0|^2+k_1)(a_2|f(p)-q_0|^2+k_2)=1.\qquad(11)
  $$


  Equation $(11)$ shows that a sphere of center $p_0$ is mapped by $f$ into a sphere
  of center $q_0$, and since $f$ preserves angles, radii map to radii. For a radial
  segment $p(s)$, $0\leq s\leq s_0$, of arc length $s$,


  $$
  \int_0^{s_0}\Big|df\Big(\frac{dp}{ds}\Big)\Big|\,ds=\int_0^{s_0}\frac{ds}{a_1|p(s)-p_0|^2+k_1}=|f(p(s_0))-f(p(0))|.
  $$


  If $k_1\neq 0$, the left side is a transcendental function of $|p(s_0)-p_0|$, while
  $(11)$ implies it is algebraic. This contradiction shows $k_1=0$, completing the
  case $\sigma\neq0$.

  \emph{Case $\sigma=0$.} Here $\rho=1/\lambda=\sum a_i x_i+c_1=A_1(x)+c_1$. Applying the
  initial argument to $f^{-1}$,


  $$
  (A_1(x)+c_1)(A_2(f(x))+c_2)=1.\qquad(11')
  $$


  Equation $(11')$ shows a hyperplane parallel to $A_1=0$ maps into a hyperplane
  parallel to $A_2=0$; a line perpendicular to $A_1=0$ maps into a line. For a
  segment $p(s)$ of such a line,


  $$
  |f(p(s_0))-f(p(0))|=\int_0^{s_0}\frac{ds}{A_1(p(s))+c_1},
  $$


  which contradicts $(11')$ unless $A_1(p(s))\equiv 0$. We conclude that if
  $\sigma=0$, $\lambda=\text{const.}$; then lengths of tangent vectors are multiplied
  by a constant $\lambda$ and $f$ is an isometry followed by a dilatation.
\end{proof}

\begin{theorem}[isometries of $H^n$ are conformal transformations preserving $H^n$]
  \label{thm:dc-ch8-5-3}
  \uses{def:dc-ch8-5-conformal, def:dc-ch8-3-hyperbolic-space}
  \notready
  \dcref{ch8:5.3}
  \notready
  The isometries of $H^n$ are the restrictions to $H^n\subset\mathbb{R}^n$ of the
  conformal transformations of $\mathbb{R}^n$ that take $H^n$ onto itself.
\end{theorem}
\begin{proof}
  \uses{cor:dc-ch8-2-3,thm:dc-ch8-5-2}
\sketch
  Suppose first $n\geq 3$. Let $f:H^n\to H^n$ be an isometry in the metric
  $g_{ij}$. By Liouville's Theorem, $f$ extends to a conformal map of $\mathbb{R}^n$
  with the metric $\delta_{ij}$. Since $H^n$ is complete in $g_{ij}$, $f$ maps $H^n$
  onto $H^n$. Conversely, let $f:H^n\to H^n$ be a conformal transformation of $H^n$
  onto $H^n$ and let $e_1,\dots,e_n$ be an orthonormal basis (in $g_{ij}$) at
  $p\in H^n$. Since $g_{ij}=\delta_{ij}/x_n^2$ and $f$ is conformal, there exists
  $\lambda^2>0$ with $\langle df_p(e_i),df_p(e_j)\rangle=\lambda^2\delta_{ij}$, so
  $\{df_p(e_i)/\lambda\}$ is orthonormal at $f(p)$. By
  \cref{cor:dc-ch8-2-3} and the global exponential coordinates of $H^n$, there is an isometry $g$ taking
  $p$ to $f(p)$ with $dg(e_i)=df(e_i)/\lambda$. By the first part, $g$ is conformal. Hence
  $h=g^{-1}\circ f$ is the restriction to $H^n$ of a conformal map of $\mathbb{R}^n$
  taking $H^n$ onto $H^n$, leaving $p$ fixed and satisfying $dh_p=\lambda I$. It
  remains to show $h=\text{identity}$.

  Let $P$ be a hyperplane through $p$. By Liouville's theorem, $h(P)$ is a hyperplane
  or a sphere through $p$; since $dh_p$ is a multiple of the identity, $P$ and $h(P)$
  are tangent at $p$. Since $h$ takes $\partial H^n$ into itself and is conformal,
  the angle of $P$ with $\partial H^n$ equals the angle of $h(P)$ with $\partial H^n$.
  We claim $h(P)=P$. Consider a line $r_1$ through $p$ perpendicular to
  $\partial H^n$, $q_1=r_1\cap\partial H^n$. Since $h(r_1)$ is a circle or a line
  making the same angle with $\partial H^n$, $h(r_1)=r_1$, hence $h(q_1)=q_1$. So if
  $P$ is perpendicular to $\partial H^n$, then $h(P)=P$. If $P$ is not perpendicular,
  let $r$ be a line in $P$; for $h(r)$ to be a circle making the same angle
  $\alpha$ it would be necessary that $q_1\in h(r)$, contradicting $h(q_1)=q_1$;
  hence $h(r)$ is a line, $h(r)=r$, and $h(P)=P$. It follows that $h$ cannot be an
  inversion (the image of some plane would be a sphere) or an isometry distinct from
  the identity. From Liouville's theorem, $h$ is a dilatation; since $h$ takes $H^n$
  onto $H^n$, $h$ is the identity.

  For $n=2$, a simple calculation shows the conformal transformations of the form


  $$
  f(z)=\frac{az+b}{cz+d},\quad z\in H^2\subset\mathbb{C},\ a,b,c,d\in\mathbb{R},\ ad-bc>0\qquad(12)
  $$


  (which map $H^2$ onto $H^2$) are isometries of $H^2$ with the metric $g_{ij}$.
  Moreover, for a fixed point $p_0\in H^2$ and unit vector $v_0$ at $p_0$, there
  exists a transformation $(12)$ taking an arbitrary $(p,v)$ to $(p_0,v_0)$ ($p$ and
  $v$ are determined by three parameters, the number of parameters of $f$). Since
  there is a unique orientation-preserving isometry of $H^2$ taking $(p,v)$ to
  $(p_0,v_0)$, all orientation-preserving isometries of $H^2$ are of the form (12).
  Composing these with $z\mapsto-\bar z$ gives the orientation-reversing
  anti-holomorphic isometries, so together they give all isometries of $H^2$.

\end{proof}

\begin{remark}[the ball model; horospheres and hyperspheres]
  \label{rem:dc-ch8-5-4}
  \uses{def:dc-ch8-5-conformal, def:dc-ch8-3-hyperbolic-space, thm:dc-ch8-5-3}
  \notready
  \dcref{ch8:5.4}
  \notready
  Let $B^n=\{p\in\mathbb R^n:|p|<2\}$ carry the metric
  $$
  h_{ij}(p)=\frac{\delta_{ij}}{(1-\tfrac14|p|^2)^2}.
  $$
  It is isometric to $H^n$ via
  $$
  f(p)=4\frac{p-p_0}{|p-p_0|^2}-(0,\dots,1),\qquad p_0=(0,\dots,-2),\qquad(13)
  $$
  A map $g:B^n\to B^n$ is an isometry for $h_{ij}$ if and only if it is the
  restriction of a conformal transformation of $\mathbb{R}^n$ taking $B^n$ onto
  $B^n$.
  Euclidean spheres contained in $H^n$ are geodesic spheres. Spheres tangent to
  $\partial H^n$ are intrinsically flat \emph{horospheres}. Spheres meeting
  $\partial H^n$ are \emph{equidistant hypersurfaces}: they lie at fixed normal
  distance from a totally geodesic hyperplane.
\end{remark}
\begin{proof}
  \sketch
  The identities
  $|df_p(v)|^2=16|v|^2/|p-p_0|^4$ and
  $f_n(p)=(4-|p|^2)/|p-p_0|^2$ show that (13) pulls back the half-space metric to
  $h$. Conjugating by (13) gives the isometry classification. Inversions send
  the three types of Euclidean spheres to centered spheres, boundary-parallel
  hyperplanes, and hyperplanes at fixed distance from a totally geodesic one,
  respectively.
\end{proof}

\begin{lemma}[a uniform radial chart neighborhood]
  \label{lem:dc-ch8-2-1-radial-chart}
  \dcref{ch8:2.1}
  \lean{Riemannian.Jacobi.exists_radial_chart_neighborhood,
    Riemannian.Jacobi.exists_pair_radial_chart_neighborhood}
  \leanok
  Let $p$ be a point of a complete Riemannian manifold and let $\exp_p$ be defined on the
  whole tangent space. There is an open neighborhood $W$ of $0\in T_pM$ such that every
  radial geodesic $t\mapsto\exp_p(tw)$ with $w\in W$ stays in the source of one fixed chart
  at $p$ for $0\leq t\leq1$. For two complete manifolds modelled on the same space $E$ and a
  continuous linear equivalence $i:E\to E$ between their tangent coordinates, one may choose
  the same $W$ for the two radial families, using the intersection with the inverse image of the
  second neighborhood under $i$.
\end{lemma}
\begin{proof}
  \leanok
  The global exponential map is continuous and sends $0$ to $p$, while the chart source is
  open and contains $p$. Its inverse image therefore contains a metric ball about $0$; radial
  homogeneity and $0\leq t\leq1$ keep every scaled vector in that ball. For a pair, intersect
  the first ball with the pullback of the second ball under $i$.
\end{proof}

\chapter{Variations of Energy}
\label{chap:dc-variations-of-energy}
\section{Formulas for the first and second variations of energy}

\begin{definition}[variation]
  \label{def:dc-ch9-2-1}
  \lean{Riemannian.Variation.IsVariation, Riemannian.Variation.IsVariationOfOrder, Riemannian.Variation.IsSmoothVariation, Riemannian.Variation.IsProperVariation, Riemannian.Variation.IsDifferentiableVariation, Riemannian.Variation.IsDifferentiableVariationOfOrder, Riemannian.Variation.IsSmoothDifferentiableVariation, Riemannian.Variation.curve, Riemannian.Variation.transversal, Riemannian.Variation.variationalField, Riemannian.Variation.variationalFieldBundle, Riemannian.Variation.IsVariationOfOrder.variationalFieldBundle_proj_eq, Riemannian.Variation.IsVariationOfOrder.variationalFieldBundle_piecewise_contMDiff, Riemannian.Variation.IsVariation.curve_zero, Riemannian.Variation.IsVariation.isVariationOfOrderOne, Riemannian.Variation.IsVariation.curve_isPiecewiseDifferentiableCurve, Riemannian.Variation.IsVariationOfOrder.isVariation, Riemannian.Variation.IsVariationOfOrder.curve_piecewise_contMDiff, Riemannian.Variation.IsVariationOfOrder.curve_isPiecewiseDifferentiableCurve, Riemannian.Variation.isProperVariation_iff, Riemannian.Variation.IsDifferentiableVariation.isVariation, Riemannian.Variation.IsDifferentiableVariationOfOrder.isVariationOfOrder}
  \leanok
  \dcref{ch9:2.1}
  Let $c:[0,a]\to M$ be a piecewise differentiable curve in a manifold $M$. A
  \emph{variation} of $c$ is a continuous mapping $f:(-\varepsilon,\varepsilon)\times[0,a]\to M$
  such that:
  \emph{(a)} $f(0,t)=c(t)$, $t\in[0,a]$,
  \emph{(b)} there exists a subdivision of $[0,a]$ by points
  $0=t_0<t_1<\cdots<t_{k+1}=a$, such that the restriction of $f$ to each
  $(-\varepsilon,\varepsilon)\times[t_i,t_{i+1}]$, $i=0,1,\dots,k$, is differentiable.

  A variation is said to be \emph{proper} if $f(s,0)=c(0)$ and $f(s,a)=c(a)$, for all
  $s\in(-\varepsilon,\varepsilon)$. If $f$ is differentiable, the variation is said to be
  \emph{differentiable}. For each $s\in(-\varepsilon,\varepsilon)$, the parametrized curve
  $f_s:[0,a]\to M$ given by $f_s(t)=f(s,t)$ is called a \emph{curve in the variation}, so
  a variation determines a family $f_s(t)$ of neighboring curves of $f_0(t)=c(t)$.
  The variation is proper exactly when the curves in this family have the same
  initial point $c(0)$ and endpoint $c(a)$.
  The parametrized differentiable curve given by $f_t(s)=f(s,t)$, $t$ fixed, is a
  \emph{transversal curve of the variation}. The velocity vector of a transversal curve
  at $s=0$, defined by $V(t)=\frac{\partial f}{\partial s}(0,t)$, is a (piecewise
  differentiable) vector field along $c(t)$ and is called the \emph{variational field} of
  $f$.
\end{definition}
\begin{proposition}[existence of a variation with given variational field]
  \label{prop:dc-ch9-2-2}
  \uses{def:dc-ch9-2-1}
  \lean{Riemannian.Variation.exists_properVariation_with_variationalField}
  \dcref{ch9:2.2}
  Given a piecewise differentiable field $V(t)$, along the piecewise differentiable
  curve $c:[0,a]\to M$, there exists a variation $f:(-\varepsilon,\varepsilon)\times[0,a]\to M$
  of $c$, such that $V(t)$ is the variational field of $f$; in addition, if
  $V(0)=V(a)=0$, it is possible to choose $f$ as a proper variation.
\end{proposition}
\begin{proof}
  \uses{thm:dc-ch3-3-7, lem:dc-ch9-2-2-exponential}
  \sketch
  Since $c([0,a])\subset M$ is compact it is possible to find a $\delta>0$
  such that $\exp_{c(t)}$ is well-defined for all $v\in T_{c(t)}M$ with $|v|<\delta$:
  for each $c(t)$ consider a totally normal neighborhood $W_t$ of $c(t)$ and the
  number $\delta_t>0$ associated to it; a finite subcover
  $W_1,\dots,W_n$ of $c([0,a])$ gives $\delta=\min(\delta_1,\dots,\delta_n)$.
  Let $N=\max_{t\in[0,a]}|V(t)|$, choose $\varepsilon>0$ with
  $\varepsilon N<\delta$, and define
  $f(s,t)=\exp_{c(t)}sV(t)$. Since $\exp_{c(t)}sV(t)=\gamma(1,c(t),sV(t))$ and the
  geodesic depends differentiably on the initial conditions, $f$ is piecewise
  differentiable, and $f(0,t)=c(t)$. The variational field of $f$ is

  $$
  \frac{\partial f}{\partial s}(0,t)=\frac{d}{ds}(\exp_{c(t)}sV(t))\Big|_{s=0}=(d\exp_{c(t)})_0 V(t)=V(t),
  $$

  and if $V(0)=V(a)=0$ then $f$ is proper.
\end{proof}
\begin{definition}[energy of a curve]
  \label{def:dc-ch9-2-2-energy}
  \lean{Riemannian.DCEnergy, Riemannian.dcSpeed, Riemannian.dcArcLength_eq_integral_dcSpeed, Riemannian.dcEnergy_eq_integral_dcSpeed_sq}
  \uses{def:dc-ch1-2-9}
  \leanok
  \dcref{ch9:2.2}
  For a curve $c:[a,b]\to M$ the \emph{speed} is $|dc/dt|=\langle dc/dt,dc/dt\rangle^{1/2}$
  and the \emph{energy} of the segment $c|[a,b]$ is
  $$
  E(c)=\int_a^b\Big|\frac{dc}{dt}\Big|^2\,dt,
  $$
  while its arc length is $L(c)=\int_a^b|dc/dt|\,dt$.
\end{definition}
\begin{lemma}[smooth curves have integrable speed and energy density]
  \label{lem:dc-ch9-2-2-smooth-speed-integrable}
  \lean{Riemannian.Variation.ContMDiffOn.continuousOn_speedSq_of_isOpen, Riemannian.Variation.ContMDiffOn.continuousOn_dcSpeed_of_isOpen, Riemannian.Variation.ContMDiffOn.intervalIntegrable_dcSpeed_of_isOpen, Riemannian.Variation.ContMDiffOn.intervalIntegrable_dcSpeed_sq_of_isOpen}
  \uses{def:dc-ch9-2-2-energy}
  \leanok
  \dcref{ch9:2.2}
  If a curve is smooth on an open neighborhood of $[a,b]$, then its speed and
  squared speed are continuous there and integrable on $[a,b]$. Consequently
  its arc length and energy on $[a,b]$ are represented by ordinary interval
  integrals.
\end{lemma}
\begin{proof}
  \leanok
  \sketch
  Near each time, choose one coordinate chart containing the curve. In that
  chart the velocity depends continuously on time, and the metric coefficients
  are smooth, so the squared speed is continuous. Its nonnegative square root
  is the speed and is also continuous. Continuous functions on the compact
  interval $[a,b]$ are integrable.
\end{proof}
\begin{lemma}[Schwarz inequality $L(c)^2\le a\,E(c)$]
  \label{lem:dc-ch9-2-2-schwarz}
  \lean{Riemannian.dcArcLength_sq_le_mul_dcEnergy, Riemannian.dcArcLength_sq_eq_iff, Riemannian.sq_integral_le_sub_mul_integral_sq, Riemannian.sq_integral_eq_iff_ae_eq_const}
  \uses{def:dc-ch9-2-2-energy}
  \leanok
  \dcref{ch9:2.2}
  For a curve $c:[a,b]\to M$ whose speed and squared speed are integrable,
  $$
  L(c)^2\le (b-a)\,E(c),
  $$
  with equality if and only if the speed $|dc/dt|$ is (almost everywhere) constant, that
  is, if and only if $t$ is proportional to arc length.

  Only integrability is required; the speed of a piecewise differentiable curve
  may jump at subdivision points.
\end{lemma}
\begin{proof}
  \leanok
  \uses{def:dc-ch9-2-2-energy}
  \sketch
  Let $f=|dc/dt|$ and let $m=\frac{1}{b-a}\int_a^b f$ be its mean. Expanding the square,
  $$
  0\le\int_a^b (f-m)^2\,dt=\int_a^b f^2\,dt-2m\int_a^b f\,dt+m^2(b-a)=E(c)-\frac{L(c)^2}{b-a},
  $$
  which is the inequality. Equality holds exactly when $\int_a^b(f-m)^2\,dt=0$, and a
  non-negative integrand has vanishing integral precisely when it vanishes almost
  everywhere; that is, exactly when $f=m$ almost everywhere.
\end{proof}
\begin{lemma}[minimizing geodesics minimize energy]
  \label{lem:dc-ch9-2-3}
  \lean{Riemannian.Geodesic.IsMinimizingGeodesicSegment.dcEnergy_le, Riemannian.Geodesic.IsMinimizingGeodesicSegment.geodesicOn_and_minimizing_of_dcEnergy_eq, Riemannian.Geodesic.IsMinimizingGeodesicSegment.dcEnergy_eq_iff_geodesicOn_and_minimizing}
  \uses{cor:dc-ch3-3-9, lem:dc-ch9-2-2-schwarz, lem:dc-ch9-2-3-inequality, lem:dc-ch9-2-3-equality-case, lem:dc-ch9-2-3-arclength-bridge, lem:dc-ch9-2-3-pointwise-arclength}
  \leanok
  \dcref{ch9:2.3}
  Let $p,q\in M$ and let $\gamma:[0,a]\to M$ be a minimizing geodesic joining $p$ to
  $q$. Then, for all curves $c:[0,a]\to M$ joining $p$ to $q$,

  $$
  E(\gamma)\le E(c)
  $$

  with equality holding if and only if $c$ is a minimizing geodesic.
\end{lemma}
\begin{proof}
  \leanok
  \sketch
  $$
  aE(\gamma)=(L(\gamma))^2\le(L(c))^2\le aE(c),
  $$

  which proves the first assertion. If equality holds, then $(L(c))^2=aE(c)$, so the
  parameter of $c$ is proportional to arc length, and $L(\gamma)=L(c)$, so $c$ is a
  minimizing geodesic (see \cref{cor:dc-ch3-3-9}). The converse is obvious.

\end{proof}
\begin{lemma}[arc length as a path length]
  \label{lem:dc-ch9-2-3-arclength-bridge}
  \lean{Riemannian.ofReal_dcArcLength_eq_pathELength, Riemannian.enorm_mfderiv_eq_ofReal_dcSpeed, Riemannian.dcArcLength_eq_toReal_pathELength, Riemannian.dcArcLength_le_of_pathELength_le}
  \uses{def:dc-ch1-2-9, def:dc-ch9-2-2-energy}
  \leanok
  \dcref{ch9:2.3}
  For a curve $c:[a,b]\to M$ whose speed is integrable, the arc length
  $L(c)=\int_a^b|dc/dt|\,dt$ agrees with the length of $c$
  measured as a path in the metric space $M$. Consequently a comparison of path lengths
  between two curves with integrable speeds is a comparison of their arc lengths $L$.

  (The integrability hypothesis is not removable: a path length is allowed to be infinite,
  while $L(c)=\int_a^b|dc/dt|\,dt$ of a non-integrable speed is not, so for such a curve the
  two genuinely differ.)
\end{lemma}
\begin{proof}
  \leanok
  \uses{def:dc-ch1-2-9}
  \sketch
  The two integrands agree pointwise: the norm of the velocity is $|dc/dt|$ by definition of
  the metric on the tangent spaces. The two notions therefore differ only in the mode of
  integration --- an integral of a real-valued function against $dt$, versus a lower integral
  of the corresponding non-negative extended-real function --- and for a non-negative
  integrable function the two agree.
\end{proof}
\begin{lemma}[a geodesic attains equality in the Schwarz comparison]
  \label{lem:dc-ch9-2-3-geodesic-equality}
  \lean{Riemannian.dcArcLength_sq_eq_mul_dcEnergy_of_isGeodesicOn, Riemannian.dcSpeed_eq_of_isGeodesicOn}
  \uses{lem:dc-ch9-2-2-schwarz, lem:dc-ch3-2-1-constspeed, def:dc-ch9-2-2-energy}
  \leanok
  \dcref{ch9:2.3}
  For a geodesic $\gamma:[a,b]\to M$ with integrable speed and squared speed,
  $$
  L(\gamma)^2 = (b-a)\,E(\gamma).
  $$
\end{lemma}
\begin{proof}
  \leanok
  \uses{lem:dc-ch9-2-2-schwarz, lem:dc-ch3-2-1-constspeed}
  \sketch
  A geodesic has constant speed (\cref{lem:dc-ch3-2-1-constspeed}), so its parameter is
  proportional to arc length, which is precisely the equality case of
  \cref{lem:dc-ch9-2-2-schwarz}.

\end{proof}
\begin{lemma}[a minimizing geodesic minimizes energy: the inequality]
  \label{lem:dc-ch9-2-3-inequality}
  \lean{Riemannian.dcEnergy_le_of_dcArcLength_le, Riemannian.Geodesic.IsMinimizingGeodesicSegment.dcEnergy_le}
  \uses{lem:dc-ch9-2-3-geodesic-equality, lem:dc-ch9-2-2-schwarz, def:dc-ch9-2-2-energy}
  \leanok
  \dcref{ch9:2.3}
  Let $\gamma:[a,b]\to M$ be a geodesic and $c:[a,b]\to M$ a curve with $L(\gamma)\le L(c)$
  (both with integrable speed and squared speed). Then $E(\gamma)\le E(c)$.

\end{lemma}
\begin{proof}
  \leanok
  \uses{lem:dc-ch9-2-3-geodesic-equality, lem:dc-ch9-2-2-schwarz}
  \sketch
  $(b-a)E(\gamma)=L(\gamma)^2\le L(c)^2\le (b-a)E(c)$, the first equality by
  \cref{lem:dc-ch9-2-3-geodesic-equality}, the middle step by hypothesis (both lengths being
  non-negative), and the last by \cref{lem:dc-ch9-2-2-schwarz}. Divide by $b-a>0$.

\end{proof}
\begin{lemma}[a smooth fixed-endpoint variation of a minimizer minimizes energy]
  \label{lem:dc-ch9-2-3-smooth-variation-energy-min}
  \lean{Riemannian.Variation.isLocalMin_dcEnergy_of_smooth_fixedEndpoint_variation}
  \uses{lem:dc-ch9-2-2-smooth-speed-integrable, lem:dc-ch9-2-3-arclength-bridge, lem:dc-ch9-2-3-inequality}
  \leanok
  \dcref{ch9:2.3}
  Let $\gamma:[a,b]\to M$ be a minimizing geodesic, and let
  $f:(-\delta,\delta)\times J\to M$ be a smooth variation on an open
  neighborhood $J$ of $[a,b]$, with $f(0,t)=\gamma(t)$ and with both endpoint
  curves fixed. Then the energy $E(s)=E(f(s,\cdot))$ has a local minimum at
  $s=0$.
\end{lemma}
\begin{proof}
  \leanok
  \sketch
  Every varied slice joins the same endpoints, so its path length is at least
  the distance between them, which is the length of $\gamma$. By
  \cref{lem:dc-ch9-2-2-smooth-speed-integrable}, all speeds and squared speeds
  involved are integrable. The path-length comparison becomes an arc-length
  comparison by \cref{lem:dc-ch9-2-3-arclength-bridge}, and
  \cref{lem:dc-ch9-2-3-inequality} gives $E(0)\le E(s)$ for every sufficiently
  small $s$.
\end{proof}
\begin{lemma}[equality in the energy comparison]
  \label{lem:dc-ch9-2-3-equality-case}
  \lean{Riemannian.dcSpeed_ae_const_of_dcEnergy_eq, Riemannian.dcArcLength_eq_of_dcEnergy_eq}
  \uses{lem:dc-ch9-2-3-geodesic-equality, lem:dc-ch9-2-2-schwarz, def:dc-ch9-2-2-energy}
  \leanok
  \dcref{ch9:2.3}
  In the situation of \cref{lem:dc-ch9-2-3-inequality}, if moreover $E(\gamma)=E(c)$, then
  \emph{(a)} the speed $|dc/dt|$ is (almost everywhere) constant, i.e. the parameter of $c$
  is proportional to arc length; and \emph{(b)} $L(\gamma)=L(c)$, i.e. $c$ is minimizing too.

\end{lemma}
\begin{proof}
  \leanok
  \uses{lem:dc-ch9-2-3-geodesic-equality, lem:dc-ch9-2-2-schwarz}
  \sketch
  Under $E(\gamma)=E(c)$ every inequality in the chain
  $(b-a)E(\gamma)=L(\gamma)^2\le L(c)^2\le (b-a)E(c)$ of
  \cref{lem:dc-ch9-2-3-inequality} collapses to an equality. From $L(c)^2=(b-a)E(c)$ and the
  equality case of \cref{lem:dc-ch9-2-2-schwarz}, the speed of $c$ is a.e. constant, which is
  (a); from $L(\gamma)^2=L(c)^2$ and non-negativity of both lengths, $L(\gamma)=L(c)$, which
  is (b).
\end{proof}
\begin{definition}[the covariant derivative along a curve, chart-independently]
  \label{def:dc-ch9-2-covariant-pair}
  \lean{Riemannian.Variation.IsCovariantDerivSolOn, Riemannian.Variation.IsCovariantDerivFieldAlongOn, Riemannian.Variation.IsCovariantDerivSolOn.transfer, Riemannian.Variation.IsCovariantDerivFieldAlongOn.isCovariantDerivSolOn_of_mem_source}
  \uses{prop:dc-ch2-2-2, lem:dc-ch2-2-2-coord}
  \leanok
  \dcref{ch9:2.4}
  Let $c:I\to M$ be a curve and let $V$, $DV$ be vector fields along $c$. Say that
  \emph{$DV$ is the covariant derivative of $V$ along $c$} on $[a,b]$ if near every
  $t_0\in[a,b]$ there is a chart containing the nearby piece of $c$ in which the
  coordinate readings satisfy formula (1) of \cref{lem:dc-ch2-2-2-coord}, namely
  $\dot V=DV-\Gamma(\dot u,V)(u)$, i.e. $\frac{DV}{dt}=DV$.

  The condition does not depend on the charts chosen: under a change of chart the
  inhomogeneous second-derivative term of the transition map produced by the product rule
  cancels against the transformation law of the Christoffel symbols, which is exactly the
  statement that $\frac{D}{dt}$ is a geometric operator. Being chart-local, the notion is
  meaningful for a curve that leaves every single chart, and on any subinterval lying in
  one chart it localizes back to the coordinate equation.

  (Only differentiability of the curve read in charts is required --- not that $c$ be a
  geodesic. The variations of \cref{def:dc-ch9-2-1} have non-geodesic transversals, so a
  geodesic hypothesis here would exclude precisely the fields this section is about.)
\end{definition}
\begin{lemma}[product rule along a curve, chart-independently]
  \label{lem:dc-ch9-2-covariant-pair-leibniz}
  \lean{Riemannian.Variation.IsCovariantDerivFieldAlongOn.hasDerivAt_metricInner}
  \uses{def:dc-ch9-2-covariant-pair, lem:dc-ch2-3-2-coord}
  \leanok
  \dcref{ch9:2.4}
  Let $V$, $W$ be vector fields along a curve $c$ with covariant derivatives $DV$, $DW$
  in the sense of \cref{def:dc-ch9-2-covariant-pair}. Then, at interior times,
  $$
  \frac{d}{dt}\langle V,W\rangle=\Big\langle\frac{DV}{dt},W\Big\rangle
    +\Big\langle V,\frac{DW}{dt}\Big\rangle.
  $$
\end{lemma}
\begin{proof}
  \leanok
  \uses{def:dc-ch9-2-covariant-pair, lem:dc-ch2-3-2-coord}
  \sketch
  The claim is local in $t$, so fix an interior time and a chart containing the nearby
  piece of $c$, as provided by \cref{def:dc-ch9-2-covariant-pair}. In that chart the
  intrinsic pairing $\langle V,W\rangle$ is the Gram pairing of the coordinate readings,
  and the coordinate metric compatibility of $\frac{D}{dt}$ (\cref{lem:dc-ch2-3-2-coord},
  the Levi-Civita case of formula (3) of \cref{prop:dc-ch2-3-2}) gives its derivative as
  $\langle\frac{DV}{dt},W\rangle+\langle V,\frac{DW}{dt}\rangle$ with $\frac{D}{dt}$ the
  coordinate covariant derivative of \cref{lem:dc-ch2-2-2-coord}. By
  \cref{def:dc-ch9-2-covariant-pair} those two coordinate covariant derivatives are the
  readings of $DV$ and $DW$, which is the assertion.
\end{proof}
\begin{lemma}[the symmetry lemma, operator form]
  \label{lem:dc-ch9-2-4-symmetry}
  \lean{Riemannian.Variation.surfaceCovariantDerivS_snd_eq_surfaceCovariantDerivT_fst, Riemannian.Variation.surfaceCovariantDerivS_snd_eq_surfaceCovariantDerivT_fst_of_forall}
  \uses{lem:dc-ch3-3-4, def:dc-ch3-3-3}
  \leanok
  \dcref{ch9:2.4}
  Let $f$ be a ($C^2$) parametrized surface (\cref{def:dc-ch3-3-3}) \emph{read in a fixed
  chart}, and let $\frac{D}{\partial s}$, $\frac{D}{\partial t}$ be the covariant
  derivatives along $f$ of that chart. Then the two exchange on the velocity fields of $f$:
  $$
  \frac{D}{\partial s}\frac{\partial f}{\partial t}
    = \frac{D}{\partial t}\frac{\partial f}{\partial s}.
  $$
  No curvature term appears. Contrast \cref{lem:dc-ch4-4-1}, which exchanges
  $\frac{D}{\partial s}$ and $\frac{D}{\partial t}$ on an \emph{arbitrary} field $V$ along
  $f$ and does pick up $R(\frac{\partial f}{\partial s},\frac{\partial f}{\partial t})V$;
  here the two sides differ only by the mixed second partial, which Schwarz kills.

\end{lemma}
\begin{proof}
  \leanok
  \uses{lem:dc-ch3-3-4}
  \sketch
  Both sides expand, in a chart, as a mixed second partial of $f$ plus a Christoffel term:
  $\frac{D}{\partial s}\frac{\partial f}{\partial t}
   = \frac{\partial^2 f}{\partial s\,\partial t}
     + \Gamma\big(\frac{\partial f}{\partial s},\frac{\partial f}{\partial t}\big)(f)$, and
  symmetrically with $s$, $t$ interchanged. The mixed partials agree by Schwarz, and the two
  Christoffel terms agree because $\Gamma$ is symmetric in its two vector slots --- which is
  exactly the hypothesis that the connection is symmetric. This is \cref{lem:dc-ch3-3-4} at
  $v=(1,0)$, $w=(0,1)$.
\end{proof}
\begin{lemma}[the $s$-derivative of the energy density]
  \label{lem:dc-ch9-2-4-energy-density}
  \lean{Riemannian.Variation.hasDerivAt_chartEnergyDensity_slice, Riemannian.Variation.hasDerivAt_chartEnergyDensity_slice_symm}
  \uses{lem:dc-ch9-2-4-symmetry, lem:dc-ch2-3-2-coord}
  \leanok
  \dcref{ch9:2.4}
  Let $f$ be a ($C^2$) parametrized surface \emph{read in a fixed chart}, as in
  \cref{lem:dc-ch9-2-4-symmetry}. Then, pointwise,
  $$
  \frac{\partial}{\partial s}\Big\langle\frac{\partial f}{\partial t},\frac{\partial f}{\partial t}\Big\rangle
    = 2\Big\langle\frac{D}{\partial s}\frac{\partial f}{\partial t},\frac{\partial f}{\partial t}\Big\rangle
    = 2\Big\langle\frac{D}{\partial t}\frac{\partial f}{\partial s},\frac{\partial f}{\partial t}\Big\rangle.
  $$
\end{lemma}
\begin{proof}
  \leanok
  \uses{lem:dc-ch9-2-4-symmetry, lem:dc-ch2-3-2-coord}
  \sketch
  The first equality is coordinate metric compatibility (\cref{lem:dc-ch2-3-2-coord})
  applied along the $s$-slice with $V=W=\frac{\partial f}{\partial t}$: the product rule's
  two terms $\langle\frac{D}{\partial s}V,W\rangle+\langle V,\frac{D}{\partial s}W\rangle$
  coincide because the pairing is symmetric, which supplies the factor $2$. The second
  equality is \cref{lem:dc-ch9-2-4-symmetry}, the symmetry of the connection, applied to
  the first factor. No curvature appears.
\end{proof}
\begin{lemma}[integration by parts along a curve]
  \label{lem:dc-ch9-2-4-parts}
  \lean{Riemannian.Variation.IsCovariantDerivFieldAlongOn.integral_metricInner_add, Riemannian.Variation.IsCovariantDerivFieldAlongOn.integral_metricInner_covariantDeriv_left, Riemannian.Variation.IsCovariantDerivFieldAlongOn.integral_metricInner_eq_neg_integral_of_proper}
  \uses{lem:dc-ch9-2-covariant-pair-leibniz, def:dc-ch9-2-covariant-pair}
  \leanok
  \dcref{ch9:2.4}
  Let $V$, $W$ be vector fields along a curve $c$ with covariant derivatives $DV$, $DW$
  in the sense of \cref{def:dc-ch9-2-covariant-pair}, on a segment $[a,b]$ carrying no
  breakpoint, and suppose the pairings $\langle\frac{DV}{dt},W\rangle$ and
  $\langle V,\frac{DW}{dt}\rangle$ are integrable on $[a,b]$. Then
  $$
  \int_a^b\Big\langle\frac{DV}{dt},W\Big\rangle dt
    = \langle V(b),W(b)\rangle-\langle V(a),W(a)\rangle
      -\int_a^b\Big\langle V,\frac{DW}{dt}\Big\rangle dt.
  $$
  If $V(a)=V(b)=0$, the identity becomes
  $\int_a^b\langle\frac{DV}{dt},W\rangle dt=-\int_a^b\langle V,\frac{DW}{dt}\rangle dt$.

  For $W=dc/dt$, summing over the segments of a piecewise differentiable curve
  telescopes the endpoint contributions into
  $-\sum_i\langle V(t_i),\frac{dc}{dt}(t_i^+)-\frac{dc}{dt}(t_i^-)\rangle$.
\end{lemma}
\begin{proof}
  \leanok
  \uses{lem:dc-ch9-2-covariant-pair-leibniz}
  \sketch
  By \cref{lem:dc-ch9-2-covariant-pair-leibniz}, $\frac{d}{dt}\langle V,W\rangle
  =\langle\frac{DV}{dt},W\rangle+\langle V,\frac{DW}{dt}\rangle$ at interior times, and
  $t\mapsto\langle V,W\rangle$ is continuous on $[a,b]$. The fundamental theorem of
  calculus (in the form requiring continuity on $[a,b]$ and a derivative on $(a,b)$ only,
  which is what \cref{lem:dc-ch9-2-covariant-pair-leibniz} supplies) integrates this to
  $\int_a^b(\langle\frac{DV}{dt},W\rangle+\langle V,\frac{DW}{dt}\rangle)dt
  =\langle V(b),W(b)\rangle-\langle V(a),W(a)\rangle$; splitting the integral and
  rearranging gives the claim. When $V(a)=V(b)=0$ both boundary terms vanish by linearity
  of the pairing in its first slot.
\end{proof}
\begin{lemma}[the $s$-derivative of the energy density, chart-free]
  \label{lem:dc-ch9-2-4-density-self}
  \lean{Riemannian.Variation.IsCovariantDerivFieldAlongOn.hasDerivAt_metricInner_self}
  \uses{lem:dc-ch9-2-covariant-pair-leibniz, def:dc-ch9-2-covariant-pair}
  \leanok
  \dcref{ch9:2.4}
  Let $V$ be a vector field along a curve $\gamma$ with covariant derivative $DV$ in the
  sense of \cref{def:dc-ch9-2-covariant-pair}. Then, at interior times,
  $$
  \frac{d}{dt}\langle V,V\rangle = 2\Big\langle\frac{DV}{dt},V\Big\rangle.
  $$

  Applied along a \emph{transversal} $\sigma\mapsto f(\sigma,t)$ of a variation
  (\cref{def:dc-ch9-2-1}) with $V=\frac{\partial f}{\partial t}$, this is the first
  equality of \cref{lem:dc-ch9-2-4-energy-density} with the chart removed: a transversal is
  a curve like any other, and \cref{def:dc-ch9-2-covariant-pair} is chart-independent.
\end{lemma}
\begin{proof}
  \leanok
  \uses{lem:dc-ch9-2-covariant-pair-leibniz}
  \sketch
  By \cref{lem:dc-ch9-2-covariant-pair-leibniz} at $W=V$, one has
  $\frac{d}{dt}\langle V,V\rangle
  =\langle\frac{DV}{dt},V\rangle+\langle V,\frac{DV}{dt}\rangle$ at interior times. The two
  summands are equal by the symmetry of the metric, which supplies the factor $2$.

\end{proof}
\begin{lemma}[symmetry of the connection along a surface, chart-free]
  \label{lem:dc-ch9-2-4-symmetry-manifold}
  \lean{Riemannian.Variation.covariantDerivS_velT_eq_covariantDerivT_velS}
  \uses{lem:dc-ch9-2-4-symmetry, def:dc-ch9-2-covariant-pair}
  \leanok
  \dcref{ch9:2.4}
  Let $f$ be a parametrized surface in $M$, let $\frac{\partial f}{\partial t}$ and
  $\frac{\partial f}{\partial s}$ be its coordinate velocity fields, and suppose
  $\big(\frac{\partial f}{\partial t},\frac{D}{\partial s}\frac{\partial f}{\partial t}\big)$
  is a covariant pair along the transversal through $(s_0,t_0)$ and
  $\big(\frac{\partial f}{\partial s},\frac{D}{\partial t}\frac{\partial f}{\partial s}\big)$
  a covariant pair along the curve in the variation through $(s_0,t_0)$, both in the sense
  of \cref{def:dc-ch9-2-covariant-pair}. If $f$ is $C^2$ near $(s_0,t_0)$ and the two slice
  curves through $(s_0,t_0)$ stay in the source of a single chart on their windows, then
  $$
  \frac{D}{\partial s}\frac{\partial f}{\partial t}(s_0,t_0)
    = \frac{D}{\partial t}\frac{\partial f}{\partial s}(s_0,t_0).
  $$

  The hypotheses concern only the two \emph{slice lines} through $(s_0,t_0)$.
\end{lemma}
\begin{proof}
  \leanok
  \uses{lem:dc-ch9-2-4-symmetry, def:dc-ch9-2-covariant-pair}
  \sketch
  Choose a chart containing the relevant pieces of both slice curves. In that chart,
  the covariant-pair hypotheses identify the intrinsic derivatives with the coordinate
  covariant derivatives. Apply \cref{lem:dc-ch9-2-4-symmetry} to the coordinate
  representation and transport the resulting equality back to $T_{f(s_0,t_0)}M$.

\end{proof}
\begin{lemma}[differentiation of the energy under the integral sign]
  \label{lem:dc-ch9-2-4-under-integral}
  \lean{Riemannian.Variation.hasDerivAt_dcEnergy_of_dominated}
  \uses{lem:dc-ch9-2-4-density-self, def:dc-ch9-2-2-energy, def:dc-ch9-2-1, def:dc-ch9-2-covariant-pair}
  \leanok
  \dcref{ch9:2.4}
  Let $f$ be a variation of $c$ and let $E$ be its energy (\cref{def:dc-ch9-2-2-energy}).
  Suppose $\big(\frac{\partial f}{\partial t},\frac{D}{\partial s}\frac{\partial f}{\partial t}\big)$
  is a covariant pair along each transversal, that the energy density is integrable at
  $s_0$, and that on a neighbourhood of $s_0$ the integrands
  $\big\langle\frac{D}{\partial s}\frac{\partial f}{\partial t},\frac{\partial f}{\partial t}\big\rangle$
  are dominated by a fixed integrable function of $t$. Then
  $$
  E'(s_0)
    = 2\int_a^b\Big\langle\frac{D}{\partial s}\frac{\partial f}{\partial t},\frac{\partial f}{\partial t}\Big\rangle\Big|_{s_0} dt .
  $$

\end{lemma}
\begin{proof}
  \leanok
  \uses{lem:dc-ch9-2-4-density-self, def:dc-ch9-2-2-energy}
  \sketch
  By \cref{def:dc-ch9-2-2-energy}, $E(s)=\int_a^b\langle\frac{\partial f}{\partial t},\frac{\partial f}{\partial t}\rangle dt$
  with the integrand evaluated along the curve $f(s,\cdot)$. For each fixed $t$ the map
  $\sigma\mapsto\langle\frac{\partial f}{\partial t},\frac{\partial f}{\partial t}\rangle(\sigma,t)$
  is, by \cref{lem:dc-ch9-2-4-density-self} applied along the transversal through $t$,
  differentiable at each interior $\sigma$ with derivative
  $2\langle\frac{D}{\partial s}\frac{\partial f}{\partial t},\frac{\partial f}{\partial t}\rangle(\sigma,t)$.
  Together with integrability at $s_0$ and the domination hypothesis, the standard theorem
  on differentiation under the integral sign applies and yields the displayed identity.

\end{proof}
\begin{lemma}[first variation on a segment without breakpoints]
  \label{lem:dc-ch9-2-4-segment}
  \lean{Riemannian.Variation.hasDerivAt_dcEnergy_eq_first_variation}
  \uses{lem:dc-ch9-2-4-under-integral, lem:dc-ch9-2-4-parts, def:dc-ch9-2-2-energy, def:dc-ch9-2-covariant-pair}
  \leanok
  \dcref{ch9:2.4}
  In the situation of \cref{lem:dc-ch9-2-4-under-integral}, let $V=\frac{\partial f}{\partial s}(s_0,\cdot)$
  carry the covariant derivative $\frac{DV}{dt}$ and let $\frac{dc}{dt}=\frac{\partial f}{\partial t}(s_0,\cdot)$
  carry $\frac{D}{dt}\frac{dc}{dt}$, both along $c=f(s_0,\cdot)$ and in the sense of
  \cref{def:dc-ch9-2-covariant-pair}. Assume the symmetry of the connection along $c$,
  $$
  \frac{D}{\partial s}\frac{\partial f}{\partial t}(s_0,t)=\frac{D}{\partial t}\frac{\partial f}{\partial s}(s_0,t)
  \qquad\text{for every interior }t\in(a,b),
  $$
  and the integrability of the two pairings of \cref{lem:dc-ch9-2-4-parts}. Then
  $$
  \frac{1}{2}E'(s_0)
    = -\int_a^b\Big\langle V,\frac{D}{dt}\frac{dc}{dt}\Big\rangle dt
      -\Big\langle V(a),\frac{dc}{dt}(a)\Big\rangle
      +\Big\langle V(b),\frac{dc}{dt}(b)\Big\rangle .
  $$

\end{lemma}
\begin{proof}
  \leanok
  \uses{lem:dc-ch9-2-4-under-integral, lem:dc-ch9-2-4-parts}
  \sketch
  By \cref{lem:dc-ch9-2-4-under-integral},
  $E'(s_0)=2\int_a^b\langle\frac{D}{\partial s}\frac{\partial f}{\partial t},\frac{\partial f}{\partial t}\rangle dt$.
  The symmetry hypothesis rewrites the integrand as
  $\langle\frac{DV}{dt},\frac{dc}{dt}\rangle$ pointwise, so
  $\frac{1}{2}E'(s_0)=\int_a^b\langle\frac{DV}{dt},\frac{dc}{dt}\rangle dt$. Applying
  \cref{lem:dc-ch9-2-4-parts} with $W=\frac{dc}{dt}$ and $DW=\frac{D}{dt}\frac{dc}{dt}$
  turns the right-hand side into
  $\langle V(b),\frac{dc}{dt}(b)\rangle-\langle V(a),\frac{dc}{dt}(a)\rangle
  -\int_a^b\langle V,\frac{D}{dt}\frac{dc}{dt}\rangle dt$, which is the claim.
\end{proof}
\begin{proposition}[formula for the first variation of the energy of a curve]
  \label{prop:dc-ch9-2-4}
  \uses{def:dc-ch9-2-1, def:dc-ch9-2-2-energy, lem:dc-ch9-2-4-symmetry, def:dc-ch9-2-covariant-pair, lem:dc-ch9-2-4-energy-density, lem:dc-ch9-2-4-parts, lem:dc-ch9-2-4-density-self, lem:dc-ch9-2-4-symmetry-manifold, lem:dc-ch9-2-4-under-integral, lem:dc-ch9-2-4-segment, lem:dc-ch9-2-4-piecewise-assembly}
  \dcref{ch9:2.4}
  Let $c:[0,a]\to M$ be a piecewise differentiable curve and let
  $f:(-\varepsilon,\varepsilon)\times[0,a]\to M$ be a variation of $c$. If
  $E:(-\varepsilon,\varepsilon)\to\mathbf{R}$ is the energy of $f$ then

  $$
  \frac{1}{2}E'(0)=-\int_0^a\Big\langle V(t),\frac{D}{dt}\frac{dc}{dt}\Big\rangle dt-\sum_{i=1}^k\Big\langle V(t_i),\frac{dc}{dt}(t_i^+)-\frac{dc}{dt}(t_i^-)\Big\rangle-\Big\langle V(0),\frac{dc}{dt}(0)\Big\rangle+\Big\langle V(a),\frac{dc}{dt}(a)\Big\rangle,\qquad(1)
  $$

  where $V(t)$ is the variational field of $f$, and
  $\frac{dc}{dt}(t_i^+)=\lim_{t\to t_i,\,t>t_i}\frac{dc}{dt}$,
  $\frac{dc}{dt}(t_i^-)=\lim_{t\to t_i,\,t<t_i}\frac{dc}{dt}$.
\end{proposition}
\begin{proof}
  \uses{def:dc-ch9-2-2-energy, lem:dc-ch9-2-4-energy-density, lem:dc-ch9-2-4-parts,
    def:dc-ch9-2-covariant-pair}
  \sketch
  By definition,
  $E(s)=\int_0^a\langle\frac{\partial f}{\partial t},\frac{\partial f}{\partial t}\rangle dt=\sum_{i=0}^k\int_{t_i}^{t_{i+1}}\langle\frac{\partial f}{\partial t},\frac{\partial f}{\partial t}\rangle dt$.
  Differentiating under the integral sign and using the symmetry of the Riemannian
  connection,

  $$
  \frac{d}{ds}\int_{t_i}^{t_{i+1}}\Big\langle\frac{\partial f}{\partial t},\frac{\partial f}{\partial t}\Big\rangle dt=2\Big\langle\frac{\partial f}{\partial s},\frac{\partial f}{\partial t}\Big\rangle\Big|_{t_i}^{t_{i+1}}-2\int_{t_i}^{t_{i+1}}\Big\langle\frac{\partial f}{\partial s},\frac{D}{dt}\frac{\partial f}{\partial t}\Big\rangle dt.
  $$

  Therefore,

  $$
  \frac{1}{2}\frac{dE}{ds}=\sum_{i=0}^k\Big\langle\frac{\partial f}{\partial s},\frac{\partial f}{\partial t}\Big\rangle\Big|_{t_i}^{t_{i+1}}-\int_0^a\Big\langle\frac{\partial f}{\partial s},\frac{D}{dt}\frac{\partial f}{\partial t}\Big\rangle dt.\qquad(2)
  $$

  Putting $s=0$ in (2) yields (1).
\end{proof}
\begin{lemma}[a geodesic is a critical point of the energy]
  \label{lem:dc-ch9-2-5-geodesic}
  \lean{Riemannian.Variation.IsCovariantDerivFieldAlongOn.integral_metricInner_covariantDeriv_eq_zero_of_geodesic}
  \uses{lem:dc-ch9-2-4-parts, def:dc-ch9-2-covariant-pair}
  \leanok
  \dcref{ch9:2.5}
  Let $W$ be a \emph{parallel} field along a curve $c$ --- $\frac{DW}{dt}=0$ --- and let
  $V$, $DV$ be a covariant-derivative pair along $c$ with $V(a)=V(b)=0$. Then on the
  segment $[a,b]$
  $$
  \int_a^b\Big\langle\frac{DV}{dt},W\Big\rangle dt=0.
  $$

  For $W=dc/dt$, parallelity is the geodesic equation; for the variational
  field of a proper variation, the endpoint condition holds. Together with
  \cref{prop:dc-ch9-2-4}, this gives the forward implication of
  \cref{prop:dc-ch9-2-5}.
\end{lemma}
\begin{proof}
  \leanok
  \uses{lem:dc-ch9-2-4-parts}
  \sketch
  Apply \cref{lem:dc-ch9-2-4-parts} with $DW=0$. The remaining integral
  $\int_a^b\langle V,\frac{DW}{dt}\rangle dt$ vanishes because its integrand is
  identically zero, and the two boundary terms vanish because $V(a)=V(b)=0$.
\end{proof}
\begin{proposition}[geodesics as critical points of energy]
  \label{prop:dc-ch9-2-5}
  \uses{prop:dc-ch9-2-4, prop:dc-ch9-2-2, lem:dc-ch9-2-5-geodesic}
  \dcref{ch9:2.5}
  A piecewise differentiable curve $c:[0,a]\to M$ is a geodesic if and only if, for
  every proper variation $f$ of $c$, we have $\frac{dE}{ds}(0)=0$.
\end{proposition}
\begin{proof}
  \uses{prop:dc-ch9-2-4, prop:dc-ch9-2-2, lem:dc-ch9-2-5-geodesic}
  \sketch
  If $c$ is a geodesic, $\frac{D}{dt}\frac{dc}{dt}=0$ and $c$ is regular;
  for $f$ proper, $V(0)=V(a)=0$, and all terms of (1) are zero
  (\cref{lem:dc-ch9-2-5-geodesic}). Conversely, suppose
  $\frac{dE}{ds}(0)=0$ for all proper variations. Let
  $V(t)=g(t)\frac{D}{dt}\frac{dc}{dt}$, where $g:[0,a]\to\mathbf{R}$ is piecewise
  differentiable with $g(t)>0$ if $t\ne t_i$ and $g(t_i)=0$, $i=0,\dots,k+1$.
  Constructing a variation
  with this $V(t)$,

  $$
  \frac{1}{2}\frac{dE}{ds}(0)=-\int_0^a g(t)\Big\langle\frac{D}{dt}\frac{dc}{dt},\frac{D}{dt}\frac{dc}{dt}\Big\rangle dt=0,
  $$

  so $\frac{D}{dt}\frac{dc}{dt}=0$ on each $(t_i,t_{i+1})$, that is, $c$ is a geodesic
  on each $(t_i,t_{i+1})$. To handle the points $t_i$, take another variational field
  $\overline V(t)$ with $\overline V(0)=\overline V(a)=0$ and, for $t\ne 0,a$,
  $\overline V(t_i)=\frac{dc}{dt}(t_i^+)-\frac{dc}{dt}(t_i^-)$. Then

  $$
  0=\frac{1}{2}\frac{dE}{ds}(0)=-\sum_{i=1}^k\Big|\frac{dc}{dt}(t_i^+)-\frac{dc}{dt}(t_i^-)\Big|^2,
  $$

  so $c$ is of class $C^1$ at each $t_i$. Since $\frac{D}{dt}\frac{dc}{dt}=0$ at
  $t_i$, $c$ satisfies the geodesic equation on $(0,a)$; by uniqueness of solutions
  to ordinary differential equations, $c\in C^\infty$ and is a geodesic.
\end{proof}
\begin{remark}[Path-space interpretation of the energy functional]
  \label{rem:dc-ch9-2-7}
  \dcref{ch9:2.7}
  For fixed $p,q\in M$, let $\Omega_{p,q}$ denote the space of piecewise
  differentiable curves from $p$ to $q$. Formally, its tangent space at a curve
  $c$ consists of piecewise differentiable vector fields along $c$ that vanish
  at both endpoints. The first variation is the directional derivative of
  $E:\Omega_{p,q}\to\mathbf{R}$, and the geodesics from $p$ to $q$ are its
  critical points. This path space is infinite-dimensional.
\end{remark}
\begin{lemma}[the second differentiation under the integral sign]
  \label{lem:dc-ch9-2-8-under-integral}
  \lean{Riemannian.Variation.hasDerivAt_dcPairing_of_dominated}
  \uses{lem:dc-ch9-2-covariant-pair-leibniz, def:dc-ch9-2-covariant-pair, lem:dc-ch9-2-4-under-integral}
  \leanok
  \dcref{ch9:2.8}
  Let $f$ be a variation of $c$, let $\frac{\partial f}{\partial t}$ carry the covariant
  $s$-derivative $\frac{D}{\partial s}\frac{\partial f}{\partial t}$ along each transversal,
  and let $\frac{D}{\partial s}\frac{\partial f}{\partial t}$ in turn carry its own covariant
  $s$-derivative $\frac{D}{\partial s}\frac{D}{\partial s}\frac{\partial f}{\partial t}$, all
  in the sense of \cref{def:dc-ch9-2-covariant-pair}; suppose the second integrand is
  dominated near $s_0$ by a fixed integrable function of $t$. Then
  $$
  \frac{d}{ds}\int_a^b\Big\langle\frac{D}{\partial s}\frac{\partial f}{\partial t},
    \frac{\partial f}{\partial t}\Big\rangle dt\Big|_{s_0}
  = \int_a^b\Big\{\Big\langle\frac{D}{\partial s}\frac{D}{\partial s}\frac{\partial f}{\partial t},
    \frac{\partial f}{\partial t}\Big\rangle
    + \Big\langle\frac{D}{\partial s}\frac{\partial f}{\partial t},
      \frac{D}{\partial s}\frac{\partial f}{\partial t}\Big\rangle\Big\}\,dt .
  $$
  The pointwise input is metric compatibility along each transversal
  (\cref{lem:dc-ch9-2-covariant-pair-leibniz}) applied to the pairs
  $\big(\frac{D}{\partial s}\frac{\partial f}{\partial t},
    \frac{D}{\partial s}\frac{D}{\partial s}\frac{\partial f}{\partial t}\big)$ and
  $\big(\frac{\partial f}{\partial t},\frac{D}{\partial s}\frac{\partial f}{\partial t}\big)$;
  the exchange of $\frac{d}{ds}$ with $\int dt$ is the same differentiation-under-the-integral
  theorem used in \cref{lem:dc-ch9-2-4-under-integral}.
\end{lemma}
\begin{proof}
  \leanok
  \uses{lem:dc-ch9-2-covariant-pair-leibniz}
  \sketch
  For each fixed $t$ the map
  $\sigma\mapsto\langle\frac{D}{\partial s}\frac{\partial f}{\partial t},\frac{\partial f}{\partial t}\rangle(\sigma,t)$
  is differentiable at interior $\sigma$ with derivative
  $\langle\frac{D}{\partial s}\frac{D}{\partial s}\frac{\partial f}{\partial t},\frac{\partial f}{\partial t}\rangle
  +\langle\frac{D}{\partial s}\frac{\partial f}{\partial t},\frac{D}{\partial s}\frac{\partial f}{\partial t}\rangle$,
  by \cref{lem:dc-ch9-2-covariant-pair-leibniz} (the product rule, not its $V=W$ collapse).
  Together with the domination hypothesis the standard theorem on differentiation under the
  integral sign yields the displayed identity.
\end{proof}
\begin{lemma}[the second variation of the energy, before the curvature substitution]
  \label{lem:dc-ch9-2-8-energy-snd-deriv}
  \lean{Riemannian.Variation.hasDerivAt_deriv_dcEnergy_second_variation}
  \uses{lem:dc-ch9-2-8-under-integral, lem:dc-ch9-2-4-under-integral, def:dc-ch9-2-2-energy, def:dc-ch9-2-covariant-pair}
  \leanok
  \dcref{ch9:2.8}
  In the situation of \cref{lem:dc-ch9-2-8-under-integral}, suppose moreover that the first
  variation $E'(\sigma)=2\int_a^b\langle\frac{D}{\partial s}\frac{\partial f}{\partial t},
  \frac{\partial f}{\partial t}\rangle\,dt$ of \cref{lem:dc-ch9-2-4-under-integral} holds at
  every $\sigma$ in a neighbourhood of $s_0$. Then
  $$
  \tfrac12 E''(s_0)
  = \int_a^b\Big\{\Big\langle\frac{D}{\partial s}\frac{D}{\partial s}\frac{\partial f}{\partial t},
    \frac{\partial f}{\partial t}\Big\rangle
    + \Big\langle\frac{D}{\partial s}\frac{\partial f}{\partial t},
      \frac{D}{\partial s}\frac{\partial f}{\partial t}\Big\rangle\Big\}\,dt .
  $$

\end{lemma}
\begin{proof}
  \leanok
  \uses{lem:dc-ch9-2-8-under-integral, lem:dc-ch9-2-4-under-integral}
  \sketch
  The first variation, supplied on a neighbourhood of $s_0$, says that $E'$ agrees there with
  $\sigma\mapsto 2\int_a^b\langle\frac{D}{\partial s}\frac{\partial f}{\partial t},\frac{\partial f}{\partial t}\rangle\,dt$;
  hence $E''(s_0)$ is the derivative of the latter at $s_0$, which
  \cref{lem:dc-ch9-2-8-under-integral} computes. Transferring the derivative along the
  eventual equality gives the claim.
\end{proof}
\begin{lemma}[the Ricci-identity substitution, metric-paired, chart-read]
  \label{lem:dc-ch9-2-8-curvature-substitution}
  \lean{Riemannian.Variation.chartMetricInner_surface_covariant_commutator_eq_curvatureFormAt, Riemannian.Variation.chartMetricInner_chartCurvatureContraction2_eq_curvatureFormAt, Riemannian.Variation.chartMetricInner_chartCurvatureContraction2_chartVectorRep_eq_curvatureFormAt}
  \uses{lem:dc-ch4-4-1, def:dc-ch4-2-1}
  \leanok
  \dcref{ch9:2.8}
  Let $f$ be a ($C^2$) parametrized surface \emph{read in a fixed chart} and $V$ a field along
  it. The commutator of the two covariant derivatives, paired against a vector $w$, is the
  intrinsic curvature $(0,4)$ form of the chart-frame realizations of $\frac{\partial f}{\partial s}$,
  $\frac{\partial f}{\partial t}$, $V$ and $w$:
  $$
  \Big\langle\frac{D}{\partial t}\frac{D}{\partial s}V - \frac{D}{\partial s}\frac{D}{\partial t}V,\ w\Big\rangle
    = \Big\langle R\Big(\frac{\partial f}{\partial s},\frac{\partial f}{\partial t}\Big)V,\ w\Big\rangle,
  $$
\end{lemma}
\begin{proof}
  \leanok
  \uses{lem:dc-ch4-4-1}
  \sketch
  The chart-level Ricci identity (\cref{lem:dc-ch4-4-1}) equates the commutator with the
  coordinate curvature contraction $\mathcal{R}_{\mathrm{chart}}(\frac{\partial f}{\partial s},\frac{\partial f}{\partial t})V$.
  Pairing against $w$ and using the coordinate-to-intrinsic curvature identity rewrites
  the result as the intrinsic curvature $(0,4)$ form
  $R(\frac{\partial f}{\partial s},\frac{\partial f}{\partial t},V,w)$.
\end{proof}
\begin{lemma}[the Ricci-identity substitution, chart-free]
  \label{lem:dc-ch9-2-8-curvature-substitution-manifold}
  \lean{Riemannian.Variation.metricInner_covariantDerivT_covariantDerivS_commutator_eq_curvatureFormAt}
  \uses{lem:dc-ch9-2-8-curvature-substitution, def:dc-ch9-2-covariant-pair, lem:dc-ch9-2-4-symmetry-manifold}
  \leanok
  \dcref{ch9:2.8}
  Let $f$ be a parametrized surface in $M$ and $V$ a field along it, and suppose
  $\frac{D}{\partial s}V$, $\frac{D}{\partial t}V$ are covariant derivatives of $V$ along the
  slices in the sense of \cref{def:dc-ch9-2-covariant-pair}, and
  $\frac{D}{\partial t}\frac{D}{\partial s}V$, $\frac{D}{\partial s}\frac{D}{\partial t}V$ the
  covariant derivatives of those, along the two slice curves through $(s_0,t_0)$. Then, at
  $(s_0,t_0)$,
  $$
  \Big\langle\frac{D}{\partial t}\frac{D}{\partial s}V - \frac{D}{\partial s}\frac{D}{\partial t}V,\ W\Big\rangle
    = \Big\langle R\Big(\frac{\partial f}{\partial s},\frac{\partial f}{\partial t}\Big)V,\ W\Big\rangle
  $$
  for every $W\in T_{f(s_0,t_0)}M$.

  The covariant derivatives are \cref{def:dc-ch9-2-covariant-pair} pairs. The hypotheses
  constrain $\frac{D}{\partial s}V$, $\frac{D}{\partial t}V$ on the two-dimensional window
  (the outer covariant derivatives differentiate them along the transverse slice), and
  $\frac{D}{\partial t}\frac{D}{\partial s}V$, $\frac{D}{\partial s}\frac{D}{\partial t}V$ only
  on their own slice line.
\end{lemma}
\begin{proof}
  \leanok
  \uses{lem:dc-ch9-2-8-curvature-substitution, lem:dc-ch9-2-4-symmetry-manifold}
  \sketch
  Choose a common chart for the two slice curves. The covariant-pair hypotheses
  identify all four intrinsic derivatives with their coordinate counterparts, so
  \cref{lem:dc-ch9-2-8-curvature-substitution} applies. The coordinate equality
  transports to the intrinsic tangent space and yields the displayed formula.
\end{proof}
\begin{lemma}[proper second variation on a segment]
  \label{lem:dc-ch9-2-8-segment}
  \lean{Riemannian.Variation.deriv_deriv_dcEnergy_eq_second_variation, Riemannian.Variation.integral_second_variation_integrand_eq, Riemannian.Variation.curvatureFormAt_swap_pairs}
  \uses{lem:dc-ch9-2-8-energy-snd-deriv, lem:dc-ch9-2-4-symmetry-manifold, lem:dc-ch9-2-8-curvature-substitution-manifold, lem:dc-ch9-2-4-parts, lem:dc-ch9-2-5-geodesic, def:dc-ch9-2-2-energy, def:dc-ch9-2-covariant-pair}
  \leanok
  \dcref{ch9:2.8}
  Let $\gamma=f(s_0,\cdot):[a,b]\to M$ be a geodesic and $f$ a variation of $\gamma$ with
  variational field $V=\frac{\partial f}{\partial s}(s_0,\cdot)$ satisfying $V(a)=V(b)=0$ and
  $\frac{D}{\partial s}\frac{\partial f}{\partial s}(s_0,a)
    =\frac{D}{\partial s}\frac{\partial f}{\partial s}(s_0,b)=0$ (all of which hold when $f$ is
  proper). Then, on the segment $[a,b]$ carrying no breakpoint,
  $$
  \frac{1}{2}E''(s_0)=-\int_a^b\Big\langle V,\frac{D^2 V}{dt^2}+R(\gamma',V)\gamma'\Big\rangle dt .
  $$

\end{lemma}
\begin{proof}
  \leanok
  \uses{lem:dc-ch9-2-8-energy-snd-deriv, lem:dc-ch9-2-4-symmetry-manifold,
    lem:dc-ch9-2-8-curvature-substitution-manifold, lem:dc-ch9-2-4-parts, lem:dc-ch9-2-5-geodesic}
  \sketch
  By \cref{lem:dc-ch9-2-8-energy-snd-deriv},
  $\frac12 E''(s_0)=\int_a^b\{\langle\frac{D}{\partial s}\frac{D}{\partial s}\frac{\partial f}{\partial t},\frac{\partial f}{\partial t}\rangle+\langle\frac{D}{\partial s}\frac{\partial f}{\partial t},\frac{D}{\partial s}\frac{\partial f}{\partial t}\rangle\}\,dt$.
  The symmetry of the connection (\cref{lem:dc-ch9-2-4-symmetry-manifold}) rewrites the second
  integrand as $\langle V',V'\rangle$, which integration by parts
  (\cref{lem:dc-ch9-2-4-parts}, the variation being proper) turns into
  $-\int_a^b\langle V,\frac{D^2V}{dt^2}\rangle\,dt$. For the first integrand, the symmetry of
  the connection followed by the Ricci identity
  (\cref{lem:dc-ch9-2-8-curvature-substitution-manifold}) gives
  $\langle\frac{D}{\partial s}\frac{D}{\partial s}\frac{\partial f}{\partial t},\frac{\partial f}{\partial t}\rangle
    =\langle\frac{D}{\partial t}\frac{D}{\partial s}\frac{\partial f}{\partial s},\frac{\partial f}{\partial t}\rangle
     -R\big(\tfrac{\partial f}{\partial s},\tfrac{\partial f}{\partial t},\tfrac{\partial f}{\partial s},\tfrac{\partial f}{\partial t}\big)$,
  whose first term integrates to $0$ because $\gamma$ is a geodesic and
  $\frac{D}{\partial s}\frac{\partial f}{\partial s}$ vanishes at the endpoints
  (\cref{lem:dc-ch9-2-5-geodesic}), and whose curvature term is $-\langle R(\gamma',V)\gamma',V\rangle$:
  at $s=0$, $\frac{\partial f}{\partial s}=V$ and $\frac{\partial f}{\partial t}=\gamma'$, so it reads
  $R(V,\gamma',V,\gamma')$, which the interchange of the two antisymmetric pairs turns into
  $R(\gamma',V,\gamma',V)=\langle R(\gamma',V)\gamma',V\rangle$.
  Summing gives formula (3).
\end{proof}
\begin{proposition}[formula for the second variation]
  \label{prop:dc-ch9-2-8}
  \uses{def:dc-ch9-2-1, def:dc-ch9-2-2-energy, def:dc-ch9-2-covariant-pair, def:dc-ch4-2-1}
  \dcref{ch9:2.8}
  Let $\gamma:[0,a]\to M$ be a geodesic and let
  $f:(-\varepsilon,\varepsilon)\times[0,a]\to M$ be a proper variation of $\gamma$. Let $E$
  be the energy function of the variation. Then

  $$
  \frac{1}{2}E''(0)=-\int_0^a\Big\langle V(t),\frac{D^2 V}{dt^2}+R\Big(\frac{d\gamma}{dt},V\Big)\frac{d\gamma}{dt}\Big\rangle dt-\sum_{i=1}^k\Big\langle V(t_i),\frac{DV}{dt}(t_i^+)-\frac{DV}{dt}(t_i^-)\Big\rangle,\qquad(3)
  $$

  where $V$ is the variational field of $f$, $R$ is the curvature of $M$ and
  $\frac{DV}{dt}(t_i^+)=\lim_{t\to t_i,\,t>t_i}\frac{DV}{dt}$,
  $\frac{DV}{dt}(t_i^-)=\lim_{t\to t_i,\,t<t_i}\frac{DV}{dt}$.
\end{proposition}
\begin{proof}
  \uses{lem:dc-ch9-2-8-energy-snd-deriv, lem:dc-ch9-2-4-symmetry-manifold, lem:dc-ch9-2-8-curvature-substitution-manifold, lem:dc-ch9-2-4-parts, lem:dc-ch9-2-5-geodesic, lem:dc-ch9-2-8-jump-sum}
  \sketch
  Taking the derivative of (2),

  $$
  \frac{1}{2}\frac{d^2 E}{ds^2}=\sum_{i=0}^k\Big\langle\frac{D}{ds}\frac{\partial f}{\partial s},\frac{\partial f}{\partial t}\Big\rangle\Big|_{t_i}^{t_{i+1}}+\sum_{i=0}^k\Big\langle\frac{\partial f}{\partial s},\frac{D}{ds}\frac{\partial f}{\partial t}\Big\rangle\Big|_{t_i}^{t_{i+1}}-\int_0^a\Big\langle\frac{D}{ds}\frac{\partial f}{\partial s},\frac{D}{dt}\frac{\partial f}{\partial t}\Big\rangle dt-\int_0^a\Big\langle\frac{\partial f}{\partial s},\frac{D}{ds}\frac{D}{dt}\frac{\partial f}{\partial t}\Big\rangle dt.
  $$

  Putting $s=0$, the first and third terms are zero, since $f$ is proper and
  $\gamma$ is a geodesic. Because

  $$
  \frac{D}{ds}\frac{D}{dt}\frac{\partial f}{\partial t}=\frac{D}{dt}\frac{D}{ds}\frac{\partial f}{\partial t}+R\Big(\frac{\partial f}{\partial t},\frac{\partial f}{\partial s}\Big)\frac{\partial f}{\partial t},
  $$

  we have, at $s=0$,
  $\frac{D}{ds}\frac{D}{dt}\frac{\partial f}{\partial t}=\frac{D^2}{dt^2}V+R(\frac{d\gamma}{dt},V)\frac{d\gamma}{dt}$.
  Further use of the fact that the variation is proper yields

  $$
  \sum_{i=0}^k\Big\langle\frac{\partial f}{\partial s},\frac{D}{ds}\frac{\partial f}{\partial t}\Big\rangle\Big|_{t_i}^{t_{i+1}}=-\sum_{i=1}^k\Big\langle V(t_i),\frac{DV}{dt}(t_i^+)-\frac{DV}{dt}(t_i^-)\Big\rangle.\qquad(4)
  $$

  Putting all these facts together, we obtain (3).
\end{proof}
\begin{remark}[nonproper second variation with boundary terms]
  \label{rem:dc-ch9-2-9}
  \uses{prop:dc-ch9-2-8}
  \lean{Riemannian.Variation.hasDerivAt_deriv_dcEnergy_eq_piecewise_second_variation_nonproper,
    Riemannian.Variation.deriv_deriv_dcEnergy_eq_piecewise_second_variation_nonproper}
  \dcref{ch9:2.9}
  If the variation is not proper, the endpoint terms in the second-variation
  formula need not vanish, and the formula becomes

  $$
  \frac{1}{2}E''(0)=-\int_0^a\Big\langle V(t),\frac{D^2 V}{dt^2}+R\Big(\frac{d\gamma}{dt},V\Big)\frac{d\gamma}{dt}\Big\rangle dt-\sum_{i=1}^k\Big\langle V(t_i),\frac{DV}{dt}(t_i^+)-\frac{DV}{dt}(t_i^-)\Big\rangle
  $$

  $$
  -\Big\langle\frac{D}{ds}\frac{\partial f}{\partial s},\frac{d\gamma}{dt}\Big\rangle(0,0)+\Big\langle\frac{D}{ds}\frac{\partial f}{\partial s},\frac{d\gamma}{dt}\Big\rangle(0,a)-\Big\langle V(0),\frac{DV}{dt}(0)\Big\rangle+\Big\langle V(a),\frac{DV}{dt}(a)\Big\rangle.\qquad(5)
  $$
\end{remark}
\begin{lemma}[nonproper second variation on a segment]
  \label{lem:dc-ch9-2-9-segment}
  \lean{Riemannian.Variation.deriv_deriv_dcEnergy_eq_second_variation_nonproper, Riemannian.Variation.integral_second_variation_integrand_eq_nonproper}
  \uses{lem:dc-ch9-2-8-energy-snd-deriv, lem:dc-ch9-2-4-symmetry-manifold, lem:dc-ch9-2-8-curvature-substitution-manifold, lem:dc-ch9-2-4-parts, def:dc-ch9-2-2-energy, def:dc-ch9-2-covariant-pair}
  \leanok
  \dcref{ch9:2.9}
  Let $\gamma=f(s_0,\cdot):[a,b]\to M$ be a geodesic and $f$ a variation of $\gamma$ that need
  \emph{not} be proper. Then, on the segment $[a,b]$ carrying no breakpoint,
  $$
  \frac{1}{2}E''(s_0)=-\int_a^b\Big\langle V,\frac{D^2 V}{dt^2}+R(\gamma',V)\gamma'\Big\rangle dt
    +\Big\langle\frac{D}{\partial s}\frac{\partial f}{\partial s},\gamma'\Big\rangle\Big|_a^b
    +\big\langle V,V'\big\rangle\Big|_a^b .
  $$

\end{lemma}
\begin{proof}
  \leanok
  \uses{lem:dc-ch9-2-8-energy-snd-deriv, lem:dc-ch9-2-4-symmetry-manifold,
    lem:dc-ch9-2-8-curvature-substitution-manifold, lem:dc-ch9-2-4-parts}
  \sketch
  Identical to the assembly of \cref{lem:dc-ch9-2-8-segment}, except that the two integrations
  by parts (\cref{lem:dc-ch9-2-4-parts}) are applied in their boundary-keeping form: the
  $\langle V',V'\rangle$ term integrates to
  $\langle V,V'\rangle|_a^b-\int_a^b\langle V,\frac{D^2V}{dt^2}\rangle$, and the geodesic term
  $\langle\frac{D}{\partial t}\frac{D}{\partial s}\frac{\partial f}{\partial s},\gamma'\rangle$
  integrates to $\langle\frac{D}{\partial s}\frac{\partial f}{\partial s},\gamma'\rangle|_a^b$
  (the remaining integral vanishes as $\gamma$ is a geodesic). The curvature term is unchanged
  from \cref{lem:dc-ch9-2-8-segment}.
\end{proof}
\begin{remark}[the index form]
  \label{rem:dc-ch9-2-10}
  \lean{Riemannian.Variation.deriv_deriv_dcEnergy_eq_piecewise_indexForm, Riemannian.Variation.deriv_deriv_dcEnergy_eq_piecewise_indexForm_of_proper}
  \uses{prop:dc-ch9-2-8, rem:dc-ch9-2-9}
  \dcref{ch9:2.10}
  On each interval where $V$ is differentiable,
  $\frac{d}{dt}\langle V,\frac{DV}{dt}\rangle=\langle V,\frac{D^2 V}{dt^2}\rangle+\langle\frac{DV}{dt},\frac{DV}{dt}\rangle$.
  For a geodesic $\gamma:[0,a]\to M$ and a partition
  $0=t_0<t_1<\cdots<t_{k+1}=a$, formula (5) is

  $$
  \frac{1}{2}E''(0)=\int_0^a\{\langle V',V'\rangle-\langle R(\gamma',V)\gamma',V\rangle\}dt-\Big\langle\frac{D}{ds}\frac{\partial f}{\partial s},\gamma'\Big\rangle(0,0)+\Big\langle\frac{D}{ds}\frac{\partial f}{\partial s},\gamma'\Big\rangle(0,a).
  $$

  Define

  $$
  \int_0^a\{\langle V',V'\rangle-\langle R(\gamma',V)\gamma',V\rangle\}dt=I_a(V,V).
  $$

  If the variation is proper, $I_a(V,V)=\frac{1}{2}E''(0)$ depends only on $V$. In
  the general case,

  $$
  \frac{1}{2}E''(0)=I_a(V,V)-\Big\langle\frac{D}{ds}\frac{\partial f}{\partial s},\gamma'\Big\rangle(0,0)+\Big\langle\frac{D}{ds}\frac{\partial f}{\partial s},\gamma'\Big\rangle(0,a).\qquad(6)
  $$
\end{remark}
\begin{lemma}[differentiating $\langle V,DV/dt\rangle$]
  \label{lem:dc-ch9-2-10-differentiation}
  \lean{Riemannian.Variation.hasDerivAt_metricInner_covariantDeriv}
  \uses{lem:dc-ch9-2-covariant-pair-leibniz}
  \leanok
  \dcref{ch9:2.10}
  On each interval where $V$ is differentiable,
  $$
  \frac{d}{dt}\Big\langle V,\frac{DV}{dt}\Big\rangle
    =\Big\langle\frac{DV}{dt},\frac{DV}{dt}\Big\rangle
     +\Big\langle V,\frac{D^2V}{dt^2}\Big\rangle,
  $$
  the identity used to pass from (5) to (6).
\end{lemma}
\begin{proof}
  \leanok
  \uses{lem:dc-ch9-2-covariant-pair-leibniz}
  \sketch
  Apply the product rule \cref{lem:dc-ch9-2-covariant-pair-leibniz} to the pair $V$, with
  covariant derivative $\frac{DV}{dt}$, against the pair $\frac{DV}{dt}$, with covariant
  derivative $\frac{D^2V}{dt^2}$.
\end{proof}
\begin{definition}[the index form $I_a$]
  \label{def:dc-ch9-2-10-index-form}
  \lean{Riemannian.Variation.indexForm, Riemannian.Variation.indexForm_def}
  \uses{def:dc-ch9-2-covariant-pair, def:dc-ch4-2-1, def:dc-ch1-2-9}
  \leanok
  \dcref{ch9:2.10}
  For a geodesic $\gamma:[0,a]\to M$ and a vector field $V$ along $\gamma$ with covariant
  derivative $V'=\frac{DV}{dt}$, the \emph{index form} is
  $$
  I_a(V,V)=\int_0^a\{\langle V',V'\rangle-\langle R(\gamma',V)\gamma',V\rangle\}\,dt,
  $$
  with $R$ the curvature tensor and $\gamma'$ the geodesic velocity. For a
  proper variation, \cref{lem:dc-ch9-2-10-formula6} gives
  $I_a(V,V)=\frac{1}{2}E''(0)$.
\end{definition}
\begin{lemma}[second variation as the index form on a segment]
  \label{lem:dc-ch9-2-10-formula6}
  \lean{Riemannian.Variation.deriv_deriv_dcEnergy_eq_indexForm, Riemannian.Variation.integral_second_variation_eq_indexForm, Riemannian.Variation.ProperSecondVariationData.deriv_deriv_dcEnergy_eq_indexForm}
  \uses{def:dc-ch9-2-10-index-form, lem:dc-ch9-2-8-energy-snd-deriv, lem:dc-ch9-2-4-symmetry-manifold, lem:dc-ch9-2-8-curvature-substitution-manifold, lem:dc-ch9-2-5-geodesic}
  \leanok
  \dcref{ch9:2.10}
  Let $\gamma=f(s_0,\cdot):[a,b]\to M$ be a geodesic and $f$ a proper variation of $\gamma$
  (in the sense of \cref{lem:dc-ch9-2-8-segment}). Then
  $$
  \frac{1}{2}E''(0)=I_a(V,V).
  $$
\end{lemma}
\begin{proof}
  \leanok
  \uses{lem:dc-ch9-2-8-energy-snd-deriv, lem:dc-ch9-2-4-symmetry-manifold,
    lem:dc-ch9-2-8-curvature-substitution-manifold, lem:dc-ch9-2-5-geodesic,
    def:dc-ch9-2-10-index-form}
  \sketch
  By \cref{lem:dc-ch9-2-8-energy-snd-deriv},
  $\frac12 E''(0)=\int_a^b\{\langle\frac{D}{\partial s}\frac{D}{\partial s}\frac{\partial f}{\partial t},\frac{\partial f}{\partial t}\rangle+\langle\frac{D}{\partial s}\frac{\partial f}{\partial t},\frac{D}{\partial s}\frac{\partial f}{\partial t}\rangle\}\,dt$.
  The symmetry of the connection (\cref{lem:dc-ch9-2-4-symmetry-manifold}) makes the second
  integrand $\langle V',V'\rangle$; the symmetry followed by the Ricci identity
  (\cref{lem:dc-ch9-2-8-curvature-substitution-manifold}), together with the geodesic
  vanishing of the boundary term (\cref{lem:dc-ch9-2-5-geodesic}), makes the first integrand
  integrate to $-\int_a^b\langle R(\gamma',V)\gamma',V\rangle\,dt$. Their sum is $I_a(V,V)$.

\end{proof}
\begin{remark}[Index form as the Hessian of energy]
  \label{rem:dc-ch9-2-12}
  \uses{rem:dc-ch9-2-7}
  \dcref{ch9:2.12}
  At a geodesic $\gamma\in\Omega_{p,q}$, the quadratic form
  $2I_a(V,V)$ is the Hessian $d^2E(V,V)$ of the energy functional in the
  direction of the endpoint-vanishing field $V$.
\end{remark}
\section{The theorems of Bonnet--Myers and of Synge--Weinstein}

\begin{lemma}[a velocity-seeded parallel orthonormal frame along a geodesic]
  \label{lem:dc-ch9-3-1-velocity-frame}
  \lean{Riemannian.Variation.exists_velocitySeededParallelOrthoFrameAlongOn}
  \leanok
  \dcref{ch9:3.1}
  Let $\gamma:[a,b]\to M$ be a geodesic with $\gamma'(a)\ne 0$. Then there are parallel fields
  $e_1,\dots,e_n$ along $\gamma$, orthonormal at every $t\in[a,b]$, and a distinguished index
  $n_0$ with $e_{n_0}(t)=\ell^{-1}\gamma'(t)$ for all $t$, where $\ell=|\gamma'|>0$ is the
  (constant) speed. In particular $e_i(t)\perp\gamma'(t)$ for $i\ne n_0$.
\end{lemma}
\begin{proof}
  \leanok
  \sketch
  Extend the unit velocity $\ell^{-1}\gamma'(a)$ to a $g$-orthonormal basis of $T_{\gamma(a)}M$
  and parallel-transport it along $\gamma$; parallel transport is a linear isometry, so the frame
  stays orthonormal. The distinguished transported member and $\ell^{-1}\gamma'$ are both parallel
  fields with the same value at $a$, so by uniqueness of parallel transport they agree on
  $[a,b]$.
\end{proof}
\begin{lemma}[the index form of a scaled parallel field]
  \label{lem:dc-ch9-3-1-index-form-scaled}
  \lean{Riemannian.Variation.indexForm_smul_eq}
  \uses{def:dc-ch9-2-10-index-form, def:dc-ch9-2-covariant-pair}
  \leanok
  \dcref{ch9:3.1}
  For a field $V(t)=\varphi(t)\,e(t)$ with covariant derivative $V'=\varphi'\,e$ (as when $e$ is
  parallel), the index form is
  $$I_a(V,V)=\int_a^b\Big((\varphi')^2\,\langle e,e\rangle
    -\varphi^2\,\langle R(\gamma',e)\gamma',e\rangle\Big)\,dt.$$
  In particular, this applies to $V_j=(\sin\pi t)e_j$ in
  \cref{thm:dc-ch9-3-1}.
\end{lemma}
\begin{proof}
  \leanok
  \sketch
  Pull $\varphi'$ out of $\langle\varphi' e,\varphi' e\rangle$ (bilinearity of the metric) and
  $\varphi$ out of each $e$-argument of $\langle R(\gamma',\varphi e)\gamma',\varphi e\rangle$
  (homogeneity of the curvature form in slots 2 and 4).
\end{proof}
\begin{lemma}[the index form of $V_j$ in the velocity-seeded frame]
  \label{lem:dc-ch9-3-1-index-form-frame}
  \lean{Riemannian.Variation.indexForm_smul_frame_eq}
  \uses{lem:dc-ch9-3-1-index-form-scaled, lem:dc-ch9-3-1-velocity-frame}
  \leanok
  \dcref{ch9:3.1}
  With the frame of \cref{lem:dc-ch9-3-1-velocity-frame} (so $e=e_j$ is a unit field and
  $\gamma'=\ell\,e_n$), for $V=\varphi\,e_j$
  $$I_a(V,V)=\int_a^b\Big((\varphi')^2-\varphi^2\,\ell^2\,\langle R(e_n,e_j)e_n,e_j\rangle\Big)\,dt.$$
  For $\varphi=\sin\pi t$ and orthonormal $e_n\perp e_j$, the curvature
  term is $\ell^2\sin^2(\pi t)K(e_n,e_j)$.
\end{lemma}
\begin{proof}
  \leanok
  \uses{lem:dc-ch9-3-1-index-form-scaled}
  \sketch
  Substitute $\langle e_j,e_j\rangle=1$ and $\gamma'=\ell\,e_n$ into
  \cref{lem:dc-ch9-3-1-index-form-scaled} and use homogeneity of the curvature form in slots
  1 and 3.
\end{proof}
\begin{lemma}[the proper sine exponential variation]
  \label{lem:dc-ch9-3-1-sine-variation}
  \lean{Riemannian.Variation.bonnetMyersSineVariation, Riemannian.Variation.bonnetMyersSineVariation_zero, Riemannian.Variation.bonnetMyersSineVariation_left, Riemannian.Variation.bonnetMyersSineVariation_right, Riemannian.Variation.hasDerivAt_bonnetMyersSineVariation_zero, Riemannian.Variation.exists_contMDiffOn_infty_bonnetMyersSineVariation_strip, Riemannian.Variation.exists_proper_sine_smoothExponentialVariation, Riemannian.Variation.exists_proper_sine_exponentialVariation}
  \leanok
  \dcref{ch9:3.1}
  Let $\gamma:\mathbb R\to M$ be a geodesic in a complete Riemannian
  manifold, and let $e$ be a parallel field along an open interval containing
  $[0,1]$. Then
  \[
    f(s,t)=\exp_{\gamma(t)}\!\bigl(s\sin(\pi t)e(t)\bigr)
  \]
  is smooth on an open product neighborhood of
  $\{0\}\times[0,1]$, satisfies $f(0,t)=\gamma(t)$, fixes both endpoints, and
  has variational field $V(t)=\sin(\pi t)e(t)$.
\end{lemma}
\begin{proof}
  \leanok
  \sketch
  Completeness makes every exponential map defined on the whole tangent
  space. Smooth dependence of geodesics on their initial data, together with
  smoothness of the parallel field, gives a common smooth product
  neighborhood of the compact zero slice. At $s=0$ the exponential of the
  zero vector is $\gamma(t)$; at $t=0,1$ the sine factor vanishes. Finally,
  differentiating $s\mapsto\exp_{\gamma(t)}(sV(t))$ at zero gives $V(t)$.
\end{proof}
\begin{lemma}[energy minimality of the sine variation]
  \label{lem:dc-ch9-3-1-sine-energy-min}
  \lean{Riemannian.Variation.isLocalMin_dcEnergy_bonnetMyersSineVariation}
  \uses{lem:dc-ch9-2-3-smooth-variation-energy-min, lem:dc-ch9-3-1-sine-variation}
  \leanok
  \dcref{ch9:3.1}
  If $\gamma:[0,1]\to M$ is a distance-minimizing geodesic and $e$ is a
  parallel field as in \cref{lem:dc-ch9-3-1-sine-variation}, then the energy
  of the sine exponential variation has a local minimum at $s=0$.
\end{lemma}
\begin{proof}
  \leanok
  \sketch
  The sine exponential variation is smooth on a product neighborhood, has
  zero slice $\gamma$, and fixes both endpoints by
  \cref{lem:dc-ch9-3-1-sine-variation}. Apply
  \cref{lem:dc-ch9-2-3-smooth-variation-energy-min}.
\end{proof}
\begin{lemma}[second-order necessary condition at a minimum]
  \label{lem:dc-ch9-3-1-second-order-min}
  \lean{Riemannian.Variation.isLocalMin_deriv_deriv_nonneg, Riemannian.Variation.indexForm_nonneg_of_energyLocalMin, Riemannian.Variation.indexForm_nonneg_of_energyLocalMin_at}
  \leanok
  \dcref{ch9:3.1}
  If $f:\mathbf{R}\to\mathbf{R}$ has a local minimum at $x$, is continuous there and is twice
  differentiable at $x$, then $f''(x)\ge 0$. This is the step by which a \emph{minimizing}
  geodesic forces $E_j''(0)\ge 0$ in \cref{thm:dc-ch9-3-1}: each curve of the variation joins
  the same endpoints, so $E_j(s)\ge E_j(0)$, i.e. $s=0$ is a minimum of $E_j$.
\end{lemma}
\begin{proof}
  \leanok
  \sketch
  If $f''(x)<0$, then for sufficiently small $h\ne0$ the symmetric second
  difference $f(x+h)+f(x-h)-2f(x)$ is negative. One of $f(x+h)$ and
  $f(x-h)$ is therefore smaller than $f(x)$, contradicting local minimality.
\end{proof}
\begin{lemma}[the summed index form is negative]
  \label{lem:dc-ch9-3-1-index-sum-neg}
  \lean{Riemannian.Variation.integral_frame_sum_lt_zero, Riemannian.Variation.sum_indexForm_smul_frame_neg, Riemannian.Variation.exists_indexForm_smul_frame_neg}
  \uses{lem:dc-ch9-3-1-velocity-frame, lem:dc-ch9-3-1-index-form-frame, lem:dc-ch9-3-3-ricci-sectional, def:dc-ch9-2-10-index-form}
  \leanok
  \dcref{ch9:3.1}
  Let $\gamma:[0,1]\to M$ be a geodesic of constant speed $\ell$ with the velocity-seeded
  parallel orthonormal frame of \cref{lem:dc-ch9-3-1-velocity-frame} (so each $e_j$, $j\ne n_0$,
  is a unit field and $\gamma'=\ell\,e_{n_0}$), and put $V_j(t)=(\sin\pi t)\,e_j(t)$. If $n\ge 2$
  (there is at least one index $j\ne n_0$), $r>0$, $\pi r<\ell$, and the raw curvature sum satisfies
  $\sum_{j\ne n_0}\langle R(e_{n_0},e_j)e_{n_0},e_j\rangle\ge (n-1)/r^2$ at every $t\in[0,1]$
  (each term continuous in $t$), then
  $$
  \sum_{j\ne n_0} I_a(V_j,V_j)<0,\qquad\text{hence}\quad I_a(V_j,V_j)<0\ \text{for some }j\ne n_0.
  $$

  By \cref{lem:dc-ch9-3-3-ricci-sectional}, the displayed curvature sum is
  the unnormalized Ricci form $Q(e_{n_0},e_{n_0})$.
\end{lemma}
\begin{proof}
  \leanok
  \uses{lem:dc-ch9-3-1-index-form-frame}
  \sketch
  For each $j\ne n_0$, \cref{lem:dc-ch9-3-1-index-form-frame} with $\varphi=\sin\pi t$
  (so $\varphi'=\pi\cos\pi t$) gives
  $I_a(V_j,V_j)=\int_0^1\big((\pi\cos\pi t)^2-\sin^2\pi t\,(\ell^2\langle R(e_{n_0},e_j)e_{n_0},e_j\rangle)\big)\,dt$.
  The integrands are continuous, so the finite sum commutes with the integral, and summing over
  the $n-1$ indices $j\ne n_0$ collapses it to
  $$
  \int_0^1\Big((n-1)(\pi\cos\pi t)^2-\sin^2\pi t\,\big(\ell^2\textstyle\sum_{j\ne n_0}\langle R(e_{n_0},e_j)e_{n_0},e_j\rangle\big)\Big)\,dt.
  $$
  Write $S(t)=\ell^2\sum_{j\ne n_0}\langle R(e_{n_0},e_j)e_{n_0},e_j\rangle$. The integrand equals
  $(n-1)\pi^2\cos 2\pi t-\sin^2\pi t\,(S(t)-(n-1)\pi^2)$; the first part integrates to $0$
  (as $\int_0^1\cos 2\pi t\,dt=0$), and the second is the integral of a function that is $\ge 0$
  on $[0,1]$ and strictly positive on $(0,1)$ --- there $\sin\pi t>0$ and, since $\pi r<\ell$
  gives $\pi^2 r^2<\ell^2$, $S(t)\ge\ell^2(n-1)/r^2>(n-1)\pi^2$. Hence the whole integral is
  $<0$. A strictly negative finite sum has a strictly negative summand.
\end{proof}
\begin{theorem}[Bonnet--Myers]
  \label{thm:dc-ch9-3-1}
  \uses{lem:dc-ch9-2-3, prop:dc-ch9-2-8, lem:dc-ch9-3-1-velocity-frame, lem:dc-ch9-3-1-index-form-frame, lem:dc-ch9-3-1-index-sum-neg, lem:dc-ch9-3-1-sine-variation, lem:dc-ch9-3-1-sine-energy-min, lem:dc-ch9-3-1-second-order-min, lem:dc-ch9-3-3-ricci-sectional}
  \notready
  \dcref{ch9:3.1}
  \notready
  Let $M^n$ be a complete Riemannian manifold. Suppose that the Ricci curvature of
  $M$ satisfies

  $$
  \mathrm{Ric}_p(v)\ge\frac{1}{r^2}>0,
  $$

  for all $p\in M$ and for all $v\in T_p(M)$. Then $M$ is compact and the diameter
  $\mathrm{diam}(M)\le\pi r$.
\end{theorem}
\begin{proof}
  \sketch
  Let $p$ and $q$ be any two points in $M$. Since $M$ is complete, there
  exists a minimizing geodesic $\gamma:[0,1]\to M$ joining $p$ to $q$. It is enough
  to show that the length $\ell(\gamma)\le\pi r$, because then $M$ is bounded and
  complete, therefore compact; and since $d(p,q)\le\pi r$ for any $p,q$, it follows
  that $\mathrm{diam}(M)\le\pi r$. Suppose, to the contrary, that
  $\ell(\gamma)=\ell>\pi r$. Consider parallel fields $e_1(t),\dots,e_{n-1}(t)$ along
  $\gamma$ which are orthonormal for each $t\in[0,1]$ and belong to the orthogonal
  complement of $\gamma'(t)$. Let $e_n(t)=\frac{\gamma'(t)}{\ell}$ and let $V_j$ be
  the vector field along $\gamma$ given by $V_j(t)=(\sin\pi t)e_j(t)$,
  $j=1,\dots,n-1$. Then $V_j(0)=V_j(1)=0$, so $V_j$ generates a proper variation of
  $\gamma$ with energy $E_j$. Using the formula for the second variation and the fact
  that $e_j$ is parallel,

  $$
  \frac{1}{2}E_j''(0)=-\int_0^1\langle V_j,V_j''+R(\gamma',V_j)\gamma'\rangle dt=\int_0^1\sin^2\pi t\,(\pi^2-\ell^2 K(e_n(t),e_j(t)))\,dt,
  $$

  where $K(e_n(t),e_j(t))$ is the sectional curvature at $\gamma(t)$ with respect to
  the plane generated by $e_n(t),e_j(t)$. Summing on $j$ and using the definition of
  Ricci curvature,

  $$
  \frac{1}{2}\sum_{j=1}^{n-1}E_j''(0)=\int_0^1\{\sin^2\pi t\,((n-1)\pi^2-(n-1)\ell^2\,\mathrm{Ric}_{\gamma(t)}(e_n(t)))\}\,dt.
  $$

  Since $\mathrm{Ric}_{\gamma(t)}(e_n(t))\ge\frac{1}{r^2}$ and $\ell>\pi r$, we have
  $(n-1)\ell^2\,\mathrm{Ric}_{\gamma(t)}(e_n(t))>(n-1)\pi^2$, hence

  $$
  \frac{1}{2}\sum_{j=1}^{n-1}E_j''(0)<\int_0^1\sin^2\pi t\,((n-1)\pi^2-(n-1)\pi^2)\,dt=0.
  $$

  Thus some $E_j''(0)<0$, which, by \cref{lem:dc-ch9-2-3},
  contradicts the fact that $\gamma$ is minimizing. Therefore $\ell\le\pi r$.
\end{proof}
\begin{corollary}[finiteness of the fundamental group]
  \label{cor:dc-ch9-3-2}
  \dcref{ch9:3.2}
  Let $M$ be a complete Riemannian manifold with $\mathrm{Ric}_p(v)\ge\delta>0$, for
  all $p\in M$ and all $v\in T_p(M)$. Then the universal cover of $M$ is compact. In
  particular, the fundamental group $\pi_1(M)$ is finite.
\end{corollary}

\begin{lemma}[finite fundamental group from a compact simply connected cover]
  \label{lem:dc-ch9-3-2-compact-cover}
  \dcref{ch9:3.2}
  \lean{Riemannian.Variation.finite_fiber_of_compact_covering,
    Riemannian.Variation.monodromy_at_injective,
    Riemannian.Variation.finite_fundamentalGroup_of_compact_simplyConnected_cover}
  \leanok
  Let $p:X\to Y$ be a surjective covering map, with $X$ compact and simply connected and
  $Y$ a $T_1$ space.  Then the fundamental group $\pi_1(Y,y)$ is finite for every basepoint
  $y\in Y$.
\end{lemma}
\begin{proof}
  \leanok
  A covering trivialization makes each fibre discrete, while compactness makes it finite.
  Simple connectedness of $X$ makes monodromy at a chosen point of the fibre injective on
  $\pi_1(Y,y)$; hence the fundamental group injects into a finite fibre.
\end{proof}

\begin{proof}
  \leanok
  \sketch
  Introducing on the universal cover $\pi:\tilde M\to M$ the covering metric
  (the metric such that $\pi$ is a local isometry), $\tilde M$ is complete and its
  Ricci curvature satisfies $\widetilde{\mathrm{Ric}}_p\ge\delta>0$. From the theorem,
  $\tilde M$ is compact. Hence the number of sheets of the covering is finite; since
  that is the number of elements in $\pi_1(M)$, we conclude that $\pi_1(M)$ is
  finite.
\end{proof}
\begin{lemma}[Ricci as a sum of sectional curvatures]
  \label{lem:dc-ch9-3-3-ricci-sectional}
  \lean{Riemannian.sectionalCurvature_eq_of_orthonormal, Riemannian.ricciForm_self_eq_sum_sectionalCurvature, Riemannian.ricciForm_self_ge_of_sectionalCurvature_ge, Riemannian.Variation.HasSectionalCurvatureLowerBound, Riemannian.Variation.hasRicciLowerBound_of_sectionalCurvatureLowerBound}
  \uses{def:dc-ch4-3-2, def:dc-ch4-4-ricci-form}
  \leanok
  \dcref{ch9:3.3}
  Let $\{e_1,\dots,e_n\}$ be an orthonormal basis of $T_pM$ and put $v=e_{i_0}$. Then
  $$
  Q(v,v)=\sum_{i\ne i_0}K(v,e_i),
  $$
  where $Q$ is the (unnormalized) Ricci form and $K$ the sectional curvature.
  Consequently, if $K(v,e_i)\ge k$ for every $i\ne i_0$, then $Q(v,v)\ge(n-1)k$.

  With the normalized Ricci curvature
  $\mathrm{Ric}_p(v)=Q(v,v)/(n-1)$, the same hypothesis gives
  $\mathrm{Ric}_p(v)\ge k$. The sum runs over $n-1$ terms because the diagonal contribution $B(v,v,v,v)$
  vanishes, and each pair $(v,e_i)$, $i\ne i_0$, being orthonormal, has $|v\wedge e_i|^2=1$,
  so its sectional curvature is the bare numerator $B(v,e_i,v,e_i)$.
\end{lemma}
\begin{proof}
  \leanok
  \uses{def:dc-ch4-3-2, def:dc-ch4-4-ricci-form}
  \sketch
  By the orthonormal-basis formula for the Ricci form, $Q(v,v)=\sum_i B(v,e_i,v,e_i)$. The
  term $i=i_0$ is $B(v,v,v,v)=0$ by the antisymmetry of $B$ in its first two slots. For
  $i\ne i_0$ the pair $(v,e_i)$ is orthonormal, so
  $|v\wedge e_i|^2=\langle v,v\rangle\langle e_i,e_i\rangle-\langle v,e_i\rangle^2=1$ and
  hence $K(v,e_i)=B(v,e_i,v,e_i)/1=B(v,e_i,v,e_i)$. Summing gives the identity. The
  inequality follows by bounding each of the $n-1$ summands below by $k$.
\end{proof}
\begin{corollary}[positive sectional curvature bound]
  \label{cor:dc-ch9-3-3}
  \uses{thm:dc-ch9-3-1, cor:dc-ch9-3-2, lem:dc-ch9-3-3-ricci-sectional}
  \dcref{ch9:3.3}
  Let $M$ be a complete Riemannian manifold with sectional curvature
  $K\ge\frac{1}{r^2}>0$. Then $M$ is compact, $\mathrm{diam}(M)\le\pi r$ and
  $\pi_1(M)$ is finite.
\end{corollary}
\begin{remark}[Positivity without a uniform lower bound]
  \label{rem:dc-ch9-3-4}
  \dcref{ch9:3.4}
  The hypothesis $K\ge\delta>0$ cannot be relaxed to $K>0$. Indeed, the paraboloid
  $\{(x,y,z)\in\mathbf{R}^3;z=x^2+y^2\}$ has curvature $K>0$, but is complete and
  non-compact.
\end{remark}
\begin{remark}[Sharpness of the Bonnet--Myers diameter bound]
  \label{rem:dc-ch9-3-6}
  \uses{thm:dc-ch9-3-1}
  \dcref{ch9:3.6}
  The estimate for the diameter given by \cref{thm:dc-ch9-3-1} cannot be improved. Indeed, the
  unit sphere $S^n\subset\mathbf{R}^{n+1}$ has constant sectional curvature equal to
  1 (therefore its Ricci curvature also is a constant equal to 1) and
  $\mathrm{diam}(S^n)=\pi$. This example is unique in the following sense: let $M^n$
  be complete with $\mathrm{Ric}_p(v)\ge 1/r^2$, for all $p\in M$ and all $v\in T_pM$;
  if $\mathrm{diam}(M)=\pi r$, then $M^n$ is isometric to the sphere $S^n$ of
  curvature $1/r^2$.
\end{remark}
\begin{theorem}[Synge--Weinstein]
  \label{thm:dc-ch9-3-7}
  \dcref{ch9:3.7}
  Let $f$ be an isometry of a compact oriented Riemannian manifold $M^n$. Suppose
  that $M$ has positive sectional curvature and that $f$ preserves the orientation of
  $M$ if $n$ is even, and reverses it if $n$ is odd. Then $f$ has a fixed point, i.e.,
  there exists $p\in M$ with $f(p)=p$.
\end{theorem}
\begin{proof}
  \proves{thm:dc-ch9-3-7}
  \uses{rem:dc-ch9-2-9, lem:dc-ch9-3-8, lem:dc-ch9-2-2-schwarz}
  \sketch
  Suppose, to the contrary, that $f(q)\ne q$ for all $q\in M$.
  Let $p\in M$ such that $d(p,f(p))$ attains a minimum. Since $M$ is complete, there
  exists a normalized minimizing geodesic $\gamma:[0,\ell]\to M$, joining $p$ to
  $f(p)$, that is, $\gamma(0)=p$, $\gamma(\ell)=f(p)$. Let
  $\tilde A=P\circ df_p:T_p(M)\to T_p(M)$, where $P$ is the parallel transport along
  $\gamma$ from $f(p)=\gamma(\ell)$ to $p$. Then $\tilde A$ is an isometry. We show
  that $(f\circ\gamma)'(0)=\gamma'(\ell)$. Consider the geodesic $f\circ\gamma$ which
  joins $f(p)$ to $f^2(p)$. Let $p'=\gamma(t')$, $t'\ne 0$, $t'\ne\ell$, and
  $f(p')=f\circ\gamma(t')$. Since $f$ is an isometry, $d(p,p')=d(f(p),f(p'))$ and,
  from the triangle inequality,

  $$
  d(p',f(p'))\le d(p',f(p))+d(f(p),f(p'))=d(p',f(p))+d(p,p')=d(p,f(p)).
  $$

  Because $d(p,f(p))$ is a minimum, $d(p',f(p'))=d(p',f(p))+d(f(p),f(p'))$. Therefore
  the curve formed by $\gamma$ and $f\circ\gamma$ is a geodesic, hence
  $(f\circ\gamma)'(0)=\gamma'(\ell)$, as asserted. It follows that $\tilde A$ leaves
  $\gamma'(0)$ fixed, since

  $$
  \tilde A(\gamma'(0))=P\circ df_p(\gamma'(0))=P((f\circ\gamma)'(0))=P(\gamma'(\ell))=\gamma'(0).
  $$

  Let $A$ be the restriction of $\tilde A$ to the orthogonal complement of
  $\gamma'(0)$. Then $A$ is orthogonal on $\mathbf{R}^{n-1}$ and, since $P$ is an
  orientation-preserving isometry,

  $$
  \det A=\det\tilde A=\det(P\circ df_p)=(-1)^n,
  $$

  where in the last equality we use the hypothesis on $f$ and the fact that $P$
  preserves orientation. By \cref{lem:dc-ch9-3-8}, $A$ leaves a vector invariant. Let $e_1(t)$
  be a unit parallel field along $\gamma$ such that, for each $t$, $e_1(t)$ belongs to
  the orthogonal complement of $\gamma'(t)$ and $e_1(0)$ is invariant by $A$. Let
  $\beta(s)$, $s\in(-\varepsilon,\varepsilon)$, be a geodesic such that $\beta(0)=p$ and
  $\beta'(0)=e_1(0)$. Because $P\circ df_p(e_1(0))=e_1(0)$, we have
  $df_p(e_1(0))=e_1(\ell)$, that is, the geodesic $f\circ\beta$ is such that
  $f\circ\beta(0)=f(p)$ and $(f\circ\beta)'(0)=e_1(\ell)$. Let $h$ be a variation of
  $\gamma$ given by $h(s,t)=\exp_{\gamma(t)}(se_1(t))$,
  $s\in(-\varepsilon,\varepsilon)$, $t\in[0,\ell]$. Since $h(s,0)=\beta(s)$, then
  $h(s,\ell)=\exp_{f(p)}(se_1(\ell))=(f\circ\beta)(s)$. Therefore,

  $$
  V(t)=\frac{\partial}{\partial s}\exp_{\gamma(t)}(se_1(t))\Big|_{s=0}=e_1(t),
  $$

  hence $\frac{D^2 V}{dt^2}=0$. The endpoint transverse curves are the
  geodesics $\beta$ and $f\circ\beta$, so their acceleration boundary terms vanish.
  Using the second-variation formula \cref{rem:dc-ch9-2-9} and
  $\frac{\partial h}{\partial s}(0,t)=e_1(t)$,

  $$
  \frac{1}{2}\frac{d^2 E}{ds^2}(0)=-\int_0^\ell\Big\langle V(t),R\Big(\frac{d\gamma}{dt},V\Big)\frac{d\gamma}{dt}\Big\rangle dt+\Big\langle\frac{d\gamma}{dt},\frac{D}{ds}\frac{\partial h}{\partial s}\Big\rangle(0,\ell)-\Big\langle\frac{d\gamma}{dt},\frac{D}{ds}\frac{\partial h}{\partial s}\Big\rangle(0,0)=-\int_0^\ell K\Big(e_1(t),\frac{d\gamma}{dt}\Big)dt.
  $$

  Because the sectional curvature is positive, $\frac{1}{2}\frac{d^2 E}{ds^2}(0)<0$,
  and therefore $\frac{dE}{ds}$ is strictly decreasing on a neighborhood of zero.
  This shows that there exists a curve $c$ in the variation such that
  $(\ell(c))^2\le\ell E(c)<\ell E(\gamma)=\ell(\gamma)^2$. Since the curves in the
  variation join $q$ to $f(q)$, we obtain a contradiction to the fact that
  $d(p,f(p))$ is a minimum.
\end{proof}
\begin{lemma}[invariant vector of an orthogonal map]
  \label{lem:dc-ch9-3-8}
  \lean{Riemannian.Variation.exists_fixed_vector_of_det_eq, Riemannian.Variation.exists_fixed_vector_of_orientation_parity}
  \leanok
  \dcref{ch9:3.8}
  Let $A$ be an orthogonal linear transformation of $\mathbf{R}^{n-1}$ and suppose
  that $\det A=(-1)^n$. Then $A$ fixes a non-zero vector of
  $\mathbf{R}^{n-1}$.
\end{lemma}
\begin{proof}
  \uses{rem:dc-ch9-2-9, lem:dc-ch9-3-7-orientation-det, lem:dc-ch9-3-8, lem:dc-ch9-2-2-schwarz}
  \sketch
  Since $A$ is orthogonal, $A-I=A(I-A^{T})$. Taking determinants in dimension
  $n-1$ gives
  $$
  \det(A-I)=\det(A)(-1)^{n-1}\det(A-I)=-\det(A-I).
  $$
  Hence $\det(A-I)=0$, so $1$ is an eigenvalue of $A$.
\end{proof}
\begin{remark}[Conformal Synge--Weinstein theorem]
  \label{rem:dc-ch9-3-9}
  \dcref{ch9:3.9}
  The conclusion of \cref{thm:dc-ch9-3-7} remains valid when $f$ is a
  conformal diffeomorphism. Extending it to arbitrary diffeomorphisms would
  preclude positive sectional curvature on $S^2\times S^2$, because the
  product of the antipodal maps is orientation-preserving and fixed-point free.
\end{remark}
\begin{corollary}[Synge]
  \label{cor:dc-ch9-3-10}
  \uses{thm:dc-ch9-3-7, cor:dc-ch9-3-3}
  \dcref{ch9:3.10}
  Let $M^n$ be a compact manifold with positive sectional curvature.
  \emph{(a)} If $M^n$ is orientable and $n$ is even, then $M$ is simply connected.
  \emph{(b)} If $n$ is odd, then $M^n$ is orientable.
\end{corollary}
\begin{proof}
  \sketch
  \emph{(a)} Equip the universal cover $\pi:\tilde M\to M$ with the
  covering metric and orientation. Compactness and positivity give a uniform
  lower curvature bound; hence \cref{cor:dc-ch9-3-3} makes $\tilde M$
  compact. Every deck transformation is an orientation-preserving isometry, so
  \cref{thm:dc-ch9-3-7} gives it a fixed point. A deck transformation with a
  fixed point is the identity, and therefore $\pi_1(M)$ is trivial.

  \emph{(b)} If $M$ were non-orientable, its orientable double cover
  $\overline M$ would be compact and carry the covering metric. Its nontrivial
  deck transformation is an orientation-reversing isometry without fixed
  points, contradicting \cref{thm:dc-ch9-3-7} when $n$ is odd.
\end{proof}
\begin{remark}[Projective-space counterexamples]
  \label{rem:dc-ch9-3-11}
  \dcref{ch9:3.11}
  The real projective space $P^2(\mathbf{R})$ of dimension two, which is compact,
  non-orientable and not simply connected, is an example which shows the necessity of
  orientability in a) and odd dimension in b). On the other hand, $P^3(\mathbf{R})$
  which is compact, orientable and not simply connected is an example which shows the
  necessity of even dimension in (a).
\end{remark}

\begin{lemma}[complete-exponential realization of a prescribed field]
  \label{lem:dc-ch9-2-2-exponential}
  \dcref{ch9:2.2}
  \lean{Riemannian.Variation.exponentialVariation,
    Riemannian.Variation.exponentialVariation_zero,
    Riemannian.Variation.exponentialVariation_isProperVariation,
    Riemannian.Variation.hasDerivAt_extChartAt_exponentialVariation_zero,
    Riemannian.Variation.variationalField_exponentialVariation,
    Riemannian.Variation.exponentialVariation_isVariationOfOrder_of_contMDiffOn,
    Riemannian.Variation.exponentialVariation_isVariation_of_contMDiffOn,
    Riemannian.Variation.exponentialVariation_isVariationOfOrder_of_piecewise_contMDiffOn,
    Riemannian.Variation.exponentialVariation_isVariation_of_piecewise_contMDiffOn,
    Riemannian.Variation.exists_variation_with_variationalField_of_contMDiffOn,
    Riemannian.Variation.exists_properVariation_with_variationalField_of_contMDiffOn,
    Riemannian.Variation.exists_properVariation_with_variationalField_of_piecewise_contMDiffOn}
  \leanok
  \uses{def:dc-ch9-2-1}
  Let $c:\mathbb R\to M$ and $V:\mathbb R\to TM$ be represented in the model
  vector space of a complete Riemannian manifold.  If, for some $a,\varepsilon>0$,
  the complete-exponential family
  $$f(s,t)=\exp_{c(t)}(sV(t))$$
  is $C^1$ on $(-\varepsilon,\varepsilon)\times[0,a]$, then $f$ is a variation of
  $c$ and its variational field is $V$.  If $V(0)=V(a)=0$, the same family is a
  proper variation.
\end{lemma}
\begin{proof}
  \leanok
  The zero-slice identity follows from $\exp_p(0)=p$.  Differentiating the radial
  exponential at the zero vector gives
  $\partial_s f(0,t)=V(t)$, so the variational field is prescribed.  The assumed
  $C^1$ regularity supplies the variation condition on the one-piece subdivision;
  when the endpoint values of $V$ vanish, the endpoint exponential identities make
  both endpoint curves constant, which is precisely properness.
\end{proof}

\begin{lemma}[equality gives pointwise proportional arc length]
  \label{lem:dc-ch9-2-3-pointwise-arclength}
  \dcref{ch9:2.3}
  \lean{Riemannian.pathELength_eq_of_dcEnergy_eq}
  \leanok
  \uses{lem:dc-ch9-2-3-equality-case, lem:dc-ch9-2-3-arclength-bridge}
  In the situation of \cref{lem:dc-ch9-2-3-equality-case}, for every $t\in[a,b]$,
  $$
  L\big(c|_{[a,t]}\big)=\frac{L(c)}{b-a}(t-a).
  $$
  Thus the parameter of $c$ is proportional to arc length pointwise, not merely almost
  everywhere.
\end{lemma}
\begin{proof}
  \leanok
  \uses{lem:dc-ch9-2-3-equality-case, lem:dc-ch9-2-3-arclength-bridge}
  By \cref{lem:dc-ch9-2-3-equality-case}, the speed of $c$ equals
  $L(c)/(b-a)$ almost everywhere on $(a,b]$. Restricting this equality to $(a,t]$ and
  integrating gives
  $L(c|_{[a,t]})=(L(c)/(b-a))(t-a)$. The identification of this integral with path length
  is \cref{lem:dc-ch9-2-3-arclength-bridge}. $\square$
\end{proof}

\begin{lemma}[finite-segment assembly of the first variation]
  \label{lem:dc-ch9-2-4-piecewise-assembly}
  \dcref{ch9:2.4}
  \lean{Riemannian.Variation.hasDerivAt_dcEnergy_eq_piecewise_first_variation, Riemannian.Variation.hasDerivAt_dcEnergy_eq_piecewise_first_variation_of_segment_formulas}
  \leanok
  \uses{def:dc-ch9-2-2-energy, lem:dc-ch9-2-4-segment}
  Let $0=t_0<\cdots<t_{k+1}=a$. Assume the energy density is integrable on every
  segment near $s_0$ and formula (1) has been established separately on each
  $[t_i,t_{i+1}]$. Then the energy on $[0,a]$ is the sum of the segment energies, its
  derivative is the sum of their derivatives, and the endpoint terms telescope to
  $$
  \left\langle V(a),\frac{dc}{dt}(a^-)\right\rangle
  -\left\langle V(0),\frac{dc}{dt}(0^+)\right\rangle
  -\sum_{i=1}^k\left\langle V(t_i),
    \frac{dc}{dt}(t_i^+)-\frac{dc}{dt}(t_i^-)\right\rangle.
  $$
  Thus the finite-additivity, differentiation, and jump-term algebra needed for the
  piecewise formula are formalized. The full proposition still has to extract the
  segment hypotheses and one-sided fields from an arbitrary variation.
\end{lemma}
\begin{proof}
  \leanok
  \uses{def:dc-ch9-2-2-energy, lem:dc-ch9-2-4-segment}
  Apply finite additivity of interval integrals, differentiate the resulting finite sum,
  and telescope consecutive endpoint terms. $\square$
\end{proof}

\begin{lemma}[proper smooth variations through a geodesic are stationary]
  \label{lem:dc-ch9-2-5-geodesic-energy}
  \dcref{ch9:2.5}
  \lean{Riemannian.Variation.hasDerivAt_dcEnergy_zero_of_geodesic, Riemannian.Variation.IsVariation.hasDerivAt_dcEnergy_zero_of_proper_geodesic, Riemannian.Variation.hasDerivAt_dcEnergy_zero_of_piecewise_geodesic_segment_formulas}
  \leanok
  \uses{lem:dc-ch9-2-4-segment, lem:dc-ch9-2-5-geodesic, def:dc-ch9-2-1}
  In the hypotheses of \cref{lem:dc-ch9-2-4-segment} on a segment $[a,b]$, suppose that
  the covariant derivative of the velocity of the base curve is zero and that the
  variational field vanishes at both endpoints. Then the energy of the variation is
  stationary at the base parameter:
  $$
  E'(s_0)=0.
  $$
\end{lemma}
\begin{proof}
  \leanok
  \uses{lem:dc-ch9-2-4-segment, lem:dc-ch9-2-5-geodesic}
  Formula (1) of \cref{lem:dc-ch9-2-4-segment} has three contributions: the integral
  against the covariant acceleration and the two endpoint pairings. The first is zero
  when the base curve is a geodesic, and the endpoint pairings are zero for a proper
  variation. Hence $E'(s_0)=0$.
\end{proof}

\begin{lemma}[telescoping the breakpoint terms]
  \label{lem:dc-ch9-2-8-jump-sum}
  \dcref{ch9:2.8}
  \lean{Riemannian.Variation.sum_segment_boundary_eq_endpoint_sub_sum_jump,
    Riemannian.Variation.sum_segment_boundary_eq_neg_sum_jump}
  \leanok
  Let $0=t_0<t_1<\cdots<t_{k+1}=a$ be a subdivision and write
  \[
    b_i^- = \left\langle V(t_i),\frac{DV}{dt}(t_i^-)\right\rangle,
    \qquad
    b_i^+ = \left\langle V(t_i),\frac{DV}{dt}(t_i^+)\right\rangle.
  \]
  Then the sum of the endpoint contributions from the $k+1$ segments is
  \[
    \sum_{i=0}^k (b_{i+1}^- - b_i^+)
      = b_{k+1}^- - b_0^+ - \sum_{i=1}^k (b_i^+ - b_i^-).
  \]
  In particular, for a proper variation the two outer terms vanish, and the
  segment boundaries give exactly the negative jump sum in formula (3).
\end{lemma}
\begin{proof}
  \leanok
  Induct on the number of internal subdivision points.  Adding the last
  segment replaces the old terminal term by the new terminal term and the
  difference $b_{k+1}^- - b_{k+1}^+$, which is the next jump contribution.
  The proper case follows by setting the two outer pairings equal to zero.
\end{proof}

\begin{lemma}[piecewise second-variation assembly under segment hypotheses]
  \label{lem:dc-ch9-2-8-piecewise-assembly}
  \dcref{ch9:2.8}
  \lean{Riemannian.Variation.hasDerivAt_deriv_dcEnergy_eq_piecewise_second_variation,
    Riemannian.Variation.hasDerivAt_deriv_dcEnergy_eq_piecewise_second_variation_of_proper,
    Riemannian.Variation.deriv_deriv_dcEnergy_eq_piecewise_second_variation_of_proper}
  \leanok
  \uses{lem:dc-ch9-2-8-segment, lem:dc-ch9-2-8-jump-sum, def:dc-ch9-2-2-energy}
  Fix a finite subdivision $\tau_0,\ldots,\tau_{k+1}$. Assume that, near the base
  parameter $s_0$, each segment energy is integrable and differentiable, and assume that
  the segment second-variation derivative is
  $$
  \frac{d}{ds}\,E_i'(s_0)=2\big((b^-_{i+1}-b^+_i)-B_i\big).
  $$
  Then finite additivity of energy and differentiation of a finite sum give the global
  derivative with endpoint terms and the internal jump sum:
  $$
  E''(s_0)=2\left(b^-_{k+1}-b^+_0-
    \sum_{i=1}^{k}(b^+_i-b^-_i)-\sum_{i=0}^{k}B_i\right).
  $$
  If the two outer pairings vanish, $b^+_0=b^-_{k+1}=0$ (as properness supplies
  in do Carmo's setting), this yields formula (3) with the subdivision jump term.
  The hypotheses are intentionally stated at the segment level:
  this lemma supplies the finite-subdivision assembly, while \cref{prop:dc-ch9-2-8} still
  needs the geometric/regularity bridge from an arbitrary variation to those hypotheses.
\end{lemma}
\begin{proof}
  \leanok
  \uses{lem:dc-ch9-2-8-segment, lem:dc-ch9-2-8-jump-sum}
  Apply finite additivity of the energy integral on the subdivision, differentiate the finite
  sum, and substitute the assumed segment identities. Reindexing the endpoint terms gives
  the displayed telescoping expression; properness removes the two outer pairings. $\square$
\end{proof}

\begin{lemma}[sectional Bonnet--Myers under explicit variation data]
  \label{lem:dc-ch9-3-3-sectional-analytic}
  \dcref{ch9:3.3}
  \lean{Riemannian.Variation.bonnetMyers_diameterBound_of_sectional_of_analytic}
  \leanok
  \uses{lem:dc-ch9-3-3-ricci-sectional, lem:dc-ch9-3-1-velocity-frame,
    lem:dc-ch9-3-1-index-sum-neg, thm:dc-ch7-2-8}
  Let $M$ be a connected, complete Riemannian manifold of dimension at least two, and let
  $r>0$. Assume that every sectional curvature is at least $1/r^2$. Suppose in addition that
  along every distance-realizing geodesic of length $\ell>\pi r$, the velocity-seeded parallel
  orthonormal frame and the sine fields $V_j(t)=\sin(\pi t)e_j(t)$ satisfy the explicit
  nonnegativity conclusion of the second-variation argument (the analytic variation data).
  Then $\operatorname{diam}(M)\le\pi r$ and $M$ is compact.

  This is the sectional-curvature specialization of the Bonnet--Myers assembly with its
  remaining variation regularity exposed; it makes no assertion about $\pi_1(M)$.
\end{lemma}
\begin{proof}
  \leanok
  By \cref{lem:dc-ch9-3-3-ricci-sectional}, the sectional lower bound gives the unnormalized
  Ricci lower bound $(n-1)/r^2$ in every orthonormal frame. If two points were farther apart
  than $\pi r$, Hopf--Rinow (\cref{thm:dc-ch7-2-8}) would provide a distance-realizing geodesic
  of length $\ell>\pi r$. The velocity-seeded frame and the sine index-form computation
  (\cref{lem:dc-ch9-3-1-velocity-frame,lem:dc-ch9-3-1-index-sum-neg}) would make the sum of
  the perpendicular index forms strictly negative, while the assumed analytic variation data
  make each summand nonnegative, a contradiction. Thus all distances are at most $\pi r$.
  Completeness and Hopf--Rinow then make the bounded manifold compact.
\end{proof}

\begin{remark}[Remark 2.11]
  \label{rem:dc-ch9-2-11}
  \dcref{ch9:2.11}
  It is possible to establish formulas for the first and second variations of the
  arc length function $L(s)$. The expressions are entirely analogous to those for the
  energy and, since they will not be used here, are not treated (see [dC 2]
  paragraph 5.4).
\end{remark}

\begin{remark}[Remark 2.6]
  \label{rem:dc-ch9-2-6}
  \dcref{ch9:2.6}
  \uses{prop:dc-ch9-2-5}
  \cref{prop:dc-ch9-2-5} furnishes a characterization of geodesics as critical points of
  the energy for all proper variations. In this sense geodesics can be thought of as
  solutions to a variational problem; such a characterization is not local but
  involves the behavior of the geodesic as a whole.
\end{remark}

\begin{remark}[Remark 3.5]
  \label{rem:dc-ch9-3-5}
  \dcref{ch9:3.5}
  In fact it is not necessary that $K$ be bounded away from zero but only that $K$
  not approach zero too fast. In this respect, see E. Calabi, "On Ricci curvatures
  and geodesics", Duke Math. J. 34 (1967), 667–676 and R. Schneider, "Konvexe
  Flächen mit langsam abnehmender Krümmung", Archiv der Math. 23 (1972), 650–654.
\end{remark}

\chapter{The Rauch Comparison Theorem}
\label{chap:dc-rauch-comparison-theorem}

\section{The Theorem of Rauch}

For a normalized geodesic $\gamma$ and a Jacobi field $J$ with
$J(0)=0$, $|J'(0)|=1$, and $J'(0)\perp\gamma'(0)$, the local expansion is
\[
 |J(t)|=t-\frac{K}{6}t^3+R,\qquad \lim_{t\to0}\frac{R}{t^3}=0,
\]
where $K$ is the sectional curvature of the plane spanned by $\gamma'(0)$ and
$J'(0)$. Thus larger curvature produces smaller transverse Jacobi fields near
$t=0$. Rauch's theorem gives a global comparison: if $\tilde\gamma$ has no
conjugate points and its sectional curvature in planes containing
$\tilde\gamma'$ is at least that of the corresponding planes along $\gamma$,
then normalized Jacobi fields with matching initial data satisfy
$|\tilde J(t)|\leq |J(t)|$. In dimension two this reduces to Sturm comparison for
\[
 f''+Kf=0,\qquad \tilde f''+\widetilde K\tilde f=0,
\]
with $f(0)=\tilde f(0)=0$ and equal positive initial derivatives.

\begin{lemma}[Factorization of a function vanishing at the origin and index form]
  \label{lem:dc-ch10-2-1}
  \dcref{ch10:2.1}
  If $h:[0,1]\to\mathbb R$ is differentiable and $h(0)=0$, there is a differentiable
  $\phi:[0,1]\to\mathbb R$ with $\phi(0)=h'(0)$ and $h(t)=t\phi(t)$. For a geodesic
  $\gamma:[0,a]\to M$ and a piecewise differentiable field $V$ along it, define
  \[
  I_{t_0}(V,V)=\int_0^{t_0}\bigl(\langle V',V'\rangle-
  \langle R(\gamma',V)\gamma',V\rangle\bigr)\,dt.
  \]
\end{lemma}

\begin{lemma}[Index comparison for fields vanishing at both endpoints]
  \label{lem:dc-ch10-2-2}
  \uses{lem:dc-ch10-2-1}
  \dcref{ch10:2.2}
  Let $\gamma:[0,a]\to M$ have no conjugate point to $\gamma(0)$ on $(0,a]$.
  If $J$ is a normal Jacobi field, $V$ a piecewise differentiable normal field,
  $J(0)=V(0)=0$, and $J(t_0)=V(t_0)$ for $t_0\in(0,a]$, then
  \[
  I_{t_0}(J,J)\leq I_{t_0}(V,V),
  \]
  with equality exactly when $V=J$ on $[0,t_0]$.
\end{lemma}
\begin{proof}
  \sketch
  Choose a basis $J_1,\ldots,J_{n-1}$ of normal Jacobi fields vanishing at $0$.
  Absence of conjugate points makes $J_i(t)$ a basis of $\gamma'(t)^\perp$ for
  $t>0$. Write $V=\sum_i f_iJ_i$; the factorization lemma extends the coefficients
  piecewise differentiably to $0$. The Wronskian identity
  \[
  \langle J_i',J_j\rangle=\langle J_i,J_j'\rangle
  \]
  follows because its derivative is zero and its value at $0$ is zero. Integration
  by parts therefore gives
  \[
  I_{t_0}(V,V)=\Big\langle\sum_i f_iJ_i,\sum_j f_jJ_j'\Big\rangle(t_0)
  +\int_0^{t_0}\Big|\sum_i f_i'J_i\Big|^2dt,
  \]
  while $I_{t_0}(J,J)$ is the same boundary term with constant coefficients
  $\alpha_i$. Since $V(t_0)=J(t_0)$, $f_i(t_0)=\alpha_i$, and hence
  \[
  I_{t_0}(V,V)=I_{t_0}(J,J)+\int_0^{t_0}\Big|\sum_i f_i'J_i\Big|^2dt.
  \]
  The inequality and its equality criterion follow from linear independence for
  $t>0$.
\end{proof}

\begin{theorem}[Rauch comparison]
  \label{thm:dc-ch10-2-3}
  \uses{lem:dc-ch10-2-2, prop:dc-ch5-3-6}
  \dcref{ch10:2.3}
  Let $\gamma:[0,a]\to M^n$ and $\tilde\gamma:[0,a]\to\widetilde M^{n+k}$ be
  geodesics with $|\gamma'|=|\tilde\gamma'|$. Let $J,\tilde J$ be Jacobi fields
  satisfying
  \[
  J(0)=\tilde J(0)=0,\quad
  \langle J'(0),\gamma'(0)\rangle=\langle\tilde J'(0),\tilde\gamma'(0)\rangle,
  \quad |J'(0)|=|\tilde J'(0)|.
  \]
  Assume that $\tilde\gamma$ has no conjugate points on $(0,a]$ and that
  \[
  \widetilde K(\tilde x,\tilde\gamma'(t))\geq K(x,\gamma'(t))
  \]
  for every $t$ and all tangent vectors $x,\tilde x$ (with the sectional
  curvatures taken in the planes they span with the geodesic velocities). Then
  $|\tilde J(t)|\leq |J(t)|$ on $[0,a]$. If equality holds at some
  $t_0\in(0,a]$, then
  \[
  \widetilde K(\tilde J(t),\tilde\gamma'(t))=K(J(t),\gamma'(t))
  \quad(0\leq t\leq t_0).
  \]
\end{theorem}
\begin{proof}
  \sketch
  By \cref{prop:dc-ch5-3-6}, the initial inner-product condition makes the
  tangential components of $J$ and $\tilde J$ equal in length; subtracting these
  parallel components reduces to the normal case. If both initial derivatives
  vanish the result is immediate. Otherwise put $v=|J|^2$ and
  $\tilde v=|\tilde J|^2$. The denominator is nonzero on $(0,a]$, and
  $v/\tilde v\to1$ at $0$. For $t_0$ with $v(t_0)\ne0$, normalize
  $U=J/\sqrt{v(t_0)}$ and $\tilde U=\tilde J/\sqrt{\tilde v(t_0)}$. The Jacobi
  equation gives
  \[
  \frac{v'(t_0)}{v(t_0)}=2I_{t_0}(U,U),\qquad
  \frac{\tilde v'(t_0)}{\tilde v(t_0)}=2I_{t_0}(\tilde U,\tilde U).
  \]
  Choose parallel orthonormal frames with first vectors tangent to the geodesics,
  and identify coefficient fields by an isometry $\phi$; then
  $\langle\phi V_1,\phi V_2\rangle=\langle V_1,V_2\rangle$ and
  $(\phi V)'=\phi(V')$. The curvature inequality gives
  $I_{t_0}(\phi U,\phi U)\leq I_{t_0}(U,U)$, while the index lemma gives
  $I_{t_0}(\tilde U,\tilde U)\leq I_{t_0}(\phi U,\phi U)$. Hence
  $(v/\tilde v)'\geq0$ and $|\tilde J|\leq|J|$. Equality at $t_0$ forces equality in
  both index comparisons on every preceding interval; the curvature inequality is
  then equality on the planes generated by the Jacobi fields, and continuity gives
  the stated interval including $0$.
\end{proof}

\begin{proposition}[Bounds on the spacing of conjugate points]
  \label{prop:dc-ch10-2-4}
  \dcref{ch10:2.4}
  If $0<L\leq K\leq H$ on $M$, then consecutive conjugate points along any geodesic
  are separated by a distance $d$ satisfying
  \[
  \frac{\pi}{\sqrt H}\leq d\leq\frac{\pi}{\sqrt L}.
  \]
\end{proposition}
\begin{proof}
  \sketch
  Compare with the sphere of constant curvature $H$. A normal Jacobi field on the
  sphere has its first zero at $\pi/\sqrt H$; Rauch comparison implies that a
  conjugate point in $M$ cannot occur earlier. Comparison with the sphere of
  curvature $L$ shows that a gap larger than $\pi/\sqrt L$ would force the sphere's
  first conjugate point to occur later, a contradiction.
\end{proof}

\begin{proposition}[Length decrease under radial curvature comparison]
  \label{prop:dc-ch10-2-5}
  \dcref{ch10:2.5}
  Let $M^n,\widetilde M^n$ satisfy $\widetilde K\geq K$ for all sectional planes.
  Fix $p\in M$, $\tilde p\in\widetilde M$, a linear isometry
  $i:T_pM\to T_{\tilde p}\widetilde M$, and $r>0$ such that $\exp_p$ is a
  diffeomorphism on $B_r(0)$ and $\exp_{\tilde p}$ is nonsingular on the corresponding
  ball. For a differentiable curve $c$ in $\exp_p(B_r(0))$, define
  \[
  \tilde c(s)=\exp_{\tilde p}\!\left(i\bigl(\exp_p^{-1}(c(s))\bigr)\right).
  \]
  Then $\ell(c)\geq\ell(\tilde c)$.
\end{proposition}
\begin{proof}
  \sketch
  Write $\bar c=\exp_p^{-1}c$ and consider the radial surface
  $f(t,s)=\exp_p(t\bar c(s))$. Its variational field $J_s$ along each radial geodesic
  has $J_s(0)=0$, $J_s(1)=c'(s)$, and initial derivative $\bar c'(s)$. The analogous
  surface in $\widetilde M$ has field $\tilde J_s$ with endpoint $\tilde c'(s)$ and
  initial derivative $i\bar c'(s)$. The initial lengths and tangential components
  agree, so Rauch gives $|\tilde c'(s)|\leq|c'(s)|$. Integrating over $s$ proves the
  length inequality.
\end{proof}

\begin{remark}[Ricci lower bound and volume comparison]
  \label{rem:dc-ch10-2-6}
  \dcref{ch10:2.6}
  If $M$ is complete with $\operatorname{Ric}_M\geq H$ and
  $\widetilde M(H)$ is the complete simply connected space of constant sectional
  curvature $H$, then normal balls satisfy
  \[
  \operatorname{vol} B_r(p)\leq \operatorname{vol}B_r(\tilde p).
  \]
  Equality implies $\operatorname{Ric}_M(\gamma'(t))=H$ for every normalized radial
  geodesic and $t<r$. The analogous upper-Ricci statement is false in general.
\end{remark}

\begin{remark}[Global volume-ratio monotonicity]
  \label{rem:dc-ch10-2-7}
  \uses{rem:dc-ch9-3-6, rem:dc-ch10-2-6}
  \dcref{ch10:2.7}
  With the notation of \cref{rem:dc-ch10-2-6}, set
  $V_r=\operatorname{vol}B_r(p)$ and $\widetilde V_r=\operatorname{vol}B_r(\tilde p)$.
  If $\operatorname{Ric}_M\geq H$, then for $R\geq r>0$
  \[
  \frac{V_r}{\widetilde V_r}\geq\frac{V_R}{\widetilde V_R}.
  \]
  For $R\leq\operatorname{diam}(M)$, equality holds exactly when
  $B_R(p)$ is isometric to $B_R(\tilde p)$.
\end{remark}

\section{Applications of the Index Lemma to immersions}

\begin{theorem}[Moore's nonimmersion theorem]
  \label{thm:dc-ch10-3-1}
  \uses{lem:dc-ch10-3-2, lem:dc-ch10-3-3}
  \dcref{ch10:3.1}
  Let $\overline M$ be complete and simply connected with
  $\overline K\leq b\leq0$, and let $M$ be compact with
  $K-\overline K\leq-b$. If $\dim\overline M<2\dim M$, no isometric immersion
  $f:M\to\overline M$ exists.
\end{theorem}
\begin{proof}
  \sketch
  Assume an immersion exists. Choose $\bar p\notin f(M)$ and $q\in M$ maximizing
  distance from $\bar p$, and let $\gamma:[0,\ell]\to\overline M$ be the normalized
  minimizing geodesic from $\bar p$ to $q$. It is perpendicular to $f(M)$ at $q$.
  For each unit $v\in T_qM$, vary $q$ in a curve tangent to $v$ and join the varied
  points radially to $\bar p$. The variational field $V$ satisfies $V(0)=0$, $V(\ell)=v$.
  Comparing the index form with the constant-curvature-$b$ model and using the index
  lemma gives $I_\ell(V,V)\geq I_\ell(\tilde J,\tilde J)$, where for a parallel field
  $\tilde w$ with $\tilde w(\ell)=\tilde v$ the model Jacobi field is
  \[
  \tilde J(t)=\frac{\sinh(t\sqrt{-b})}{\sinh(\ell\sqrt{-b})}\tilde w(t)\quad(b<0),
  \qquad \tilde J(t)=\frac{t}{\ell}\tilde w(t)\quad(b=0),
  \]
  and $I_\ell(\tilde J,\tilde J)=\sqrt{-b}\coth(\ell\sqrt{-b})>\sqrt{-b}$
  when $b<0$ (while it is positive for $b=0$). Since the
  varied radial geodesics are no longer than $\gamma$, the second variation yields
  \[
  0\geq\tfrac12E''(0)=I_\ell(V,V)+
  \langle S_{\gamma'(\ell)}v,v\rangle,
  \]
  hence
  \[
  \|B(v,v)\|>\sqrt{-b}\quad\text{for every unit }v. \tag{6}
  \]
  The Gauss equation gives, for orthonormal $v,w$,
  \[
  \langle B(v,v),B(w,w)\rangle-\|B(v,w)\|^2
  =K(v,w)-\overline K(v,w)\leq-b. \tag{7}
  \]
  Lemma \cref{lem:dc-ch10-3-3} says these inequalities force the codimension to be
  at least $\dim M$, contradicting $\dim\overline M<2\dim M$.
\end{proof}

\begin{lemma}[Second variation with the second fundamental form]
  \label{lem:dc-ch10-3-2}
  \dcref{ch10:3.2}
  If $E(s)$ is the energy of the radial variation in the proof of
  \cref{thm:dc-ch10-3-1}, then
  \[
  \tfrac12E''(0)=I_\ell(V,V)+
  \langle S_{\gamma'(\ell)}V(\ell),V(\ell)\rangle,
  \]
  where $S_{\gamma'(\ell)}$ is the shape operator of the immersion at $q$ in the
  normal direction $\gamma'(\ell)$.
\end{lemma}
\begin{proof}
  \sketch
  The second variation formula gives the index form plus
  $\langle\overline\nabla_{\overline V}\overline V,\overline N\rangle$ at $q$, for
  extensions $\overline V,\overline N$ of $V(\ell),\gamma'(\ell)$ with zero inner
  product. Metric compatibility converts this term to
  $\langle S_{\gamma'(\ell)}V(\ell),V(\ell)\rangle$.
\end{proof}

\begin{lemma}[Otsuki's bilinear-form lemma]
  \label{lem:dc-ch10-3-3}
  \dcref{ch10:3.3}
  Let $B:\mathbb R^n\times\mathbb R^n\to\mathbb R^k$ be symmetric and let $b\leq0$.
  If for every orthonormal pair $v,w$
  \[
  \langle B(v,v),B(w,w)\rangle-\|B(v,w)\|^2\leq-b,
  \qquad \|B(v,v)\|>\sqrt{-b},
  \]
  then $k\geq n$.
\end{lemma}
\begin{proof}
  \sketch
  If $k<n$, minimize $F(v)=\|B(v,v)\|^2$ on the unit sphere. At a minimizer $v$,
  differentiation along $v\cos t+w\sin t$ gives
  \[
  0=4\langle B(v,w),B(v,v)\rangle, \tag{8}
  \]
  and
  \[
  0\leq 8\|B(v,w)\|^2-4\|B(v,v)\|^2+
  4\langle B(w,w),B(v,v)\rangle. \tag{9}
  \]
  The map $w\mapsto B(v,w)$ on $v^\perp$ has image orthogonal to $B(v,v)$ and
  therefore a nonzero kernel vector $w_0$. Substituting $B(v,w_0)=0$ into (9) gives
  \[
  0\leq\langle B(w_0,w_0),B(v,v)\rangle-\|B(v,v)\|^2
  <\langle B(w_0,w_0),B(v,v)\rangle-(\sqrt{-b})^2
  \leq-b-(\sqrt{-b})^2=0,
  \]
  a contradiction.
\end{proof}

\begin{corollary}[Tompkins nonimmersion corollary]
  \label{cor:dc-ch10-3-4}
  \uses{rem:dc-ch10-2-7}
  \dcref{ch10:3.4}
  A compact flat $n$-manifold has no isometric immersion in $\mathbb R^{n+k}$ when
  $k<n$; in particular, the flat torus $T^n$ cannot be immersed in
  $\mathbb R^{2n-1}$.
\end{corollary}

\begin{corollary}[O'Neill nonimmersion corollary]
  \label{cor:dc-ch10-3-5}
  \dcref{ch10:3.5}
  If $\overline M$ is complete simply connected with $K_{\overline M}\leq0$, $M$ is
  compact with $\dim\overline M<2\dim M$, and $K_M\leq K_{\overline M}$, then no
  isometric immersion $M\to\overline M$ exists.
\end{corollary}

\begin{remark}[Necessity of compactness]
  \label{rem:dc-ch10-3-6}
  \uses{thm:dc-ch10-3-1}
  \dcref{ch10:3.6}
  Compactness in \cref{thm:dc-ch10-3-1} cannot be dropped: complete nonpositively
  curved surfaces can immerse in $\mathbb R^3$. The hyperbolic plane, however, has
  no isometric immersion in $\mathbb R^3$.
\end{remark}

\begin{remark}[Chern--Kuiper generalization]
  \label{rem:dc-ch10-3-7}
  \uses{thm:dc-ch10-3-1}
  \dcref{ch10:3.7}
  Let $M^n$ be compact and suppose each $T_pM$ contains an $m$-dimensional subspace
  $V$ on which all sectional curvatures are nonpositive. If $k<m$, there is no
  isometric immersion $M^n\to\mathbb R^{n+k}$.
\end{remark}

\begin{remark}[Failure without simple connectivity]
  \label{rem:dc-ch10-3-8}
  \uses{thm:dc-ch10-3-1, rem:dc-ch10-2-7}
  \dcref{ch10:3.8}
  The flat torus $T^{n+1}$ contains $T^n$ as a totally geodesic submanifold, so the
  simple-connectivity hypothesis on the ambient manifold in
  \cref{thm:dc-ch10-3-1} is necessary.
\end{remark}

\section{Focal points and an extension of Rauch's theorem}

Let $N\subset M$ be a submanifold and let $\gamma$ start orthogonally from
$p\in N$. Variations through geodesics whose initial points lie in $N$ and whose
initial velocities are normal to $N$ produce Jacobi fields satisfying the boundary
conditions below.

\begin{lemma}[Boundary conditions for normal variations]
  \label{lem:dc-ch10-4-1}
  \dcref{ch10:4.1}
  For such a variation, its Jacobi field $J$ satisfies
  \begin{enumerate}
  \item $J(0)\in T_pN$;
  \item $J'(0)+S_{\gamma'(0)}(J(0))\in(T_pN)^\perp$,
  \end{enumerate}
  where $S_{\gamma'(0)}$ is the shape operator of $N$. Conversely, every Jacobi
  field satisfying these two conditions is the variational field of a variation
  through geodesics normal to $N$ at their initial points.
\end{lemma}
\begin{proof}
  \sketch
  The first condition is $J(0)=\alpha'(0)$ for the initial curve $\alpha$ in $N$.
  Differentiating the normality of the initial velocity $A(s)$ gives, for every
  $v\in T_pN$,
  $\langle J'(0),v\rangle=-\langle S_{\gamma'(0)}J(0),v\rangle$, proving the
  second condition. Conversely, choose $\alpha$ in $N$ with $\alpha'(0)=J(0)$,
  extend $\gamma'(0)$ to a field $W$ along $\alpha$ with prescribed covariant
  derivative $J'(0)$, project $W$ to the normal bundle, and set
  $f(s,t)=\exp_{\alpha(s)}(tV(s))$. The boundary condition ensures that its
  variational field agrees with $J$.
\end{proof}

\begin{definition}[Focal point]
  \label{def:dc-ch10-4-2}
  \dcref{ch10:4.2}
  A point $q$ is focal to $N\subset M$ if there are $p\in N$ and a geodesic
  $\gamma:[0,\ell]\to M$ with $\gamma(0)=p$, $\gamma'(0)\perp T_pN$, and
  $\gamma(\ell)=q$, together with a nonzero Jacobi field satisfying the two
  conditions of \cref{lem:dc-ch10-4-1} and $J(\ell)=0$.
\end{definition}

\begin{example}[Focal points of an equator]
  \label{ex:dc-ch10-4-3}
  \dcref{ch10:4.3}
  The north and south poles are focal points of the equator $S^{n-1}\subset S^n$.
  More generally, for the normal bundle $T(N)^\perp\to N$, the normal exponential
  map is the restriction
  \[
  \exp^\perp:T(N)^\perp\to M,
  \]
  and $\dim T(N)^\perp=\dim M$.
\end{example}

\begin{proposition}[Focal points and the normal exponential map]
  \label{prop:dc-ch10-4-4}
  \uses{lem:dc-ch10-4-1}
  \dcref{ch10:4.4}
  A point $q$ is focal to $N$ if and only if it is a critical value of
  $\exp^\perp:T(N)^\perp\to M$.
\end{proposition}
\begin{proof}
  \sketch
  A normal geodesic variation with a focal Jacobi field gives a curve
  $w(s)=(\alpha(s),\ell A(s))$ in $T(N)^\perp$ whose image under $\exp^\perp$ has
  zero derivative at $s=0$; hence its endpoint is a critical value. Conversely,
  represent a critical tangent vector to $\exp^\perp$ by a curve
  $w(s)=(\sigma(s),\ell V(s))$, use
  $f(s,t)=\exp_{\sigma(s)}(tV(s))$, and apply
  \cref{lem:dc-ch10-4-1}. The resulting Jacobi field is nonzero, satisfies the
  boundary conditions, and vanishes at the critical endpoint.
\end{proof}

\begin{definition}[Focal set]
  \label{def:dc-ch10-4-5}
  \dcref{ch10:4.5}
  The focal set $F(N)\subset M$ is the set of focal points of $N$.
\end{definition}

\begin{example}[Focal surfaces in Euclidean space]
  \label{ex:dc-ch10-4-6}
  \uses{prop:dc-ch10-4-4}
  \dcref{ch10:4.6}
  For a regular surface $N^2\subset\mathbb R^3$, the focal set is obtained by
  moving along each unit normal a distance equal to the reciprocal of a principal
  curvature. In local coordinates with first fundamental matrix $(g_{ij})$ and
  second fundamental matrix $(h_{ij})$, the normal exponential map is
  $\mathbf x+tn$ and is singular exactly when $(g_{ij}-th_{ij})$ is singular. Thus
  $\mathbf x+tn$ is focal precisely when $1/t$ is a principal curvature. For a
  sphere the focal set is its center; for a circular cylinder it is its axis.
\end{example}

\begin{definition}[Focal-point-free geodesic]
  \label{def:dc-ch10-4-7}
  \dcref{ch10:4.7}
  A geodesic $\gamma:[0,a]\to M$ is focal-point free on $(0,a]$ if some ball
  $B_\varepsilon(0)\subset\gamma'(0)^\perp$ has the property that $\gamma$ has no
  focal point relative to $\Sigma_\varepsilon=\exp_{\gamma(0)}B_\varepsilon(0)$.
  Since $\Sigma_\varepsilon$ is geodesic at $\gamma(0)$, a Jacobi field with
  $J(0)\ne0$ and $J'(0)=0$ has $S_{\gamma'(0)}J(0)=0$.
\end{definition}

\begin{lemma}[Index comparison for focal-point-free geodesics]
  \label{lem:dc-ch10-4-8}
  \dcref{ch10:4.8}
  If $\gamma:[0,a]\to M^n$ is focal-point free on $(0,a]$, $J$ is a normal Jacobi
  field with $J'(0)=0$, and $V$ is piecewise differentiable with
  $J(t_0)=V(t_0)$, then
  \[
  I_{t_0}(J,J)\leq I_{t_0}(V,V),
  \]
  with equality exactly when $V=J$ on $[0,t_0]$.
\end{lemma}
\begin{proof}
  \sketch
  Jacobi fields normal to $\gamma$ with derivative zero at $0$ form an
  $(n-1)$-dimensional space. Focal-point freeness makes their values a basis of
  $\gamma'(t)^\perp$ for every $t>0$. Expanding $V$ in this basis and repeating the
  integration-by-parts argument of \cref{lem:dc-ch10-2-2} yields the inequality and
  the same equality criterion.
\end{proof}

\begin{theorem}[Rauch comparison with focal points]
  \label{thm:dc-ch10-4-9}
  \uses{lem:dc-ch10-4-8, thm:dc-ch10-2-3}
  \dcref{ch10:4.9}
  Let $\gamma:[0,a]\to M^n$ and $\tilde\gamma:[0,a]\to\widetilde M^{n+k}$ have
  equal speeds, and let $J,\tilde J$ be Jacobi fields satisfying
  \[
  J'(0)=\tilde J'(0)=0,\quad
  \langle J(0),\gamma'(0)\rangle=\langle\tilde J(0),\tilde\gamma'(0)\rangle,
  \quad |J(0)|=|\tilde J(0)|.
  \]
  Assume $\tilde\gamma$ is focal-point free on $(0,a]$ and
  \[
  \widetilde K(\tilde x,\tilde\gamma'(t))\geq K(x,\gamma'(t))
  \]
  for all $t$ and all tangent vectors $x,\tilde x$. Then
  $|\tilde J(t)|\leq|J(t)|$. Equality at $t_0>0$ implies
  \[
  \widetilde K(\tilde J(t),\tilde\gamma'(t))=K(J(t),\gamma'(t))
  \quad(0\leq t\leq t_0).
  \]
\end{theorem}
\begin{proof}
  \sketch
  Repeat the quotient-of-squared-lengths argument in
  \cref{thm:dc-ch10-2-3}, now with initial derivatives zero. Focal-point freeness
  ensures that $|J|^2/|\tilde J|^2$ is defined on $(0,a]$, and the focal index lemma
  reverses the comparison step to give
  $I_{t_0}(\phi(U),\phi(U))\geq I_{t_0}(\tilde U,\tilde U)$. The resulting quotient is
  nondecreasing, proving the norm inequality; equality forces equality in the index
  and curvature comparisons on every preceding interval, yielding the curvature
  equality by continuity.
\end{proof}

\chapter{The Morse Index Theorem}
\label{chap:dc-morse-index-theorem}

\section{The Index Theorem}

Let $\gamma:[0,a]\to M^n$ be a geodesic and let
\[
  \mathcal V=\mathcal V(0,a)=\{V:\text{piecewise differentiable along }\gamma:
  V(0)=V(a)=0\}.
\]

\begin{definition}[Index form]
  \label{def:dc-ch11-2-1}
  \uses{prop:dc-ch4-2-5}
  \dcref{ch11:2.1}
  The index form is the symmetric bilinear form
  \[
  I_a(V,W)=\int_0^a\bigl(\langle V',W'\rangle-
  \langle R(\gamma',V)\gamma',W\rangle\bigr)\,dt,
  \qquad V,W\in\mathcal V.
  \]
  Its index is the maximal dimension of a subspace on which $I_a(V,V)$ is
  negative definite. Its null space is
  \[
  \{V\in\mathcal V:I_a(V,W)=0\text{ for every }W\in\mathcal V\},
  \]
  and the nullity is its dimension. The form is degenerate when this dimension
  is positive. Symmetry follows from part (d) of \cref{prop:dc-ch4-2-5}.
\end{definition}

\begin{theorem}[Morse index theorem]
  \label{thm:dc-ch11-2-2}
  \dcref{ch11:2.2}
  The index of $I_a$ is finite and equals the number of points $\gamma(t)$,
  $0<t<a$, conjugate to $\gamma(0)$ along $\gamma$, with each point counted
  according to its multiplicity.
\end{theorem}
\begin{proof}
  \sketch
  For $t\in[0,a]$, let $i(t)$ be the index of the index form on
  $\gamma|_{[0,t]}$. By \cref{cor:dc-ch11-2-6}, \cref{lem:dc-ch11-2-7}, and
  \cref{lem:dc-ch11-2-8}, $i$ is zero near the origin, left-continuous, and its
  jump at $t$ equals the nullity of $I_t$. By \cref{cor:dc-ch11-2-4}, this nullity
  is the multiplicity of the conjugate point $\gamma(t)$; away from conjugate
  points it is zero. Thus $i(a)$ is the sum of the multiplicities of all conjugate
  points in $(0,a)$, proving \cref{thm:dc-ch11-2-2}.
\end{proof}

\begin{proposition}[Null space of the index form]
  \label{prop:dc-ch11-2-3}
  \dcref{ch11:2.3}
  An element $V\in\mathcal V$ lies in the null space of $I_a$ if and only if
  $V$ is a Jacobi field along $\gamma$.
\end{proposition}
\begin{proof}
  \sketch
  If $0=t_0<\cdots<t_k=a$ is a subdivision on which $V$ is differentiable, integration
  by parts gives
  \[
  I_a(V,W)=-\int_0^a\langle V''+R(\gamma',V)\gamma',W\rangle\,dt
  -\sum_{j=1}^{k-1}\left\langle V'(t_j^+)-V'(t_j^-),W(t_j)\right\rangle.\tag{1}
  \]
  A Jacobi field therefore annihilates $I_a$. Conversely, if $I_a(V,W)=0$ for all
  $W$, choose a differentiable $f$ positive away from the subdivision points and
  zero at them, and put $W=f(V''+R(\gamma',V)\gamma')$. Equation (1) yields
  \[
  \int_0^a f\,\lVert V''+R(\gamma',V)\gamma'\rVert^2\,dt=0,
  \]
  so $V$ satisfies the Jacobi equation on each open subinterval. A field $T$ with
  $T(t_j)=V'(t_j^+)-V'(t_j^-)$ then gives
  \[
  0=-\sum_{j=1}^{k-1}\lVert V'(t_j^+)-V'(t_j^-)\rVert^2,
  \]
  hence $V$ is $C^1$ and the Jacobi equation makes it smooth on all of $[0,a]$.
\end{proof}

\begin{corollary}[Degeneracy and conjugate endpoints]
  \label{cor:dc-ch11-2-4}
  \dcref{ch11:2.4}
  The form $I_a$ is degenerate exactly when $\gamma(a)$ is conjugate to
  $\gamma(0)$ along $\gamma$; then its nullity equals the multiplicity of
  $\gamma(a)$. Choose a normal subdivision
  \[
    0=t_0<t_1<\cdots<t_{k-1}<t_k=a
  \]
  for which every segment $\gamma|_{[t_{j-1},t_j]}$ lies in a totally normal
  neighborhood. Let $\mathcal V^-$ consist of fields that are Jacobi on each open
  subinterval and let $\mathcal V^+$ consist of fields vanishing at every interior
  subdivision point. The space $\mathcal V^-$ is finite-dimensional.
\end{corollary}

\begin{proposition}[Orthogonal splitting of the variation space]
  \label{prop:dc-ch11-2-5}
  \dcref{ch11:2.5}
  \mathcal V=\mathcal V^+\oplus\mathcal V^-,
  \qquad I_a(\mathcal V^+,\mathcal V^-)=0,
  \]
  and $I_a$ restricted to $\mathcal V^+$ is positive definite.
\end{proposition}
\begin{proof}
  \sketch
  For $V\in\mathcal V$, solve the Jacobi equation on each normal subinterval with
  the prescribed values $V(t_j)$; the resulting $W\in\mathcal V^-$ is unique and
  $V-W\in\mathcal V^+$. Integration by parts shows that for $X\in\mathcal V^-$ and
  $Y\in\mathcal V^+$ all boundary terms vanish, so $I_a(X,Y)=0$. Each normal
  segment minimizes energy, hence $I_a(Y,Y)\geq0$ for $Y\in\mathcal V^+$. If equality
  held for a nonzero $Y$, then for every $Z\in\mathcal V^+$ the inequality
\[
    0\leq I_a(Y+cZ,Y+cZ)=2cI_a(Y,Z)+c^2I_a(Z,Z)
  \]
  for all $c\in\mathbb R$ would force $I_a(Y,Z)=0$; orthogonality gives the same
  for $Z\in\mathcal V^-$. Thus $Y$ is a Jacobi field by \cref{prop:dc-ch11-2-3}, and
  its zeros at the subdivision points imply $Y=0$, a contradiction.
\end{proof}

\begin{corollary}[Finite-dimensional reduction]
  \label{cor:dc-ch11-2-6}
  \dcref{ch11:2.6}
  The index and nullity of $I_a$ equal those of its restriction to the finite-dimensional
  space $\mathcal V^-$. In particular, the index is finite.
\end{corollary}

For $t\in[0,a]$, write $\gamma_t=\gamma|_{[0,t]}$, let $I_t$ be its index form,
and let $i(t)$ denote the index of $I_t$. Choosing a normal subdivision with
$t\in(t_{j-1},t_j)$ identifies $\mathcal V^-(0,t)$ with
\[
  S_j=T_{\gamma(t_1)}M\oplus\cdots\oplus T_{\gamma(t_{j-1})}M.
\]
Thus the forms $I_t$ for $t\in(t_{j-1},t_j)$ act on a fixed finite-dimensional
space and vary continuously with $t$. The index is zero near $0$ and is
nondecreasing: extending a negative subspace on $[0,t]$ by zero to $[0,\bar t]$
preserves its form when $\bar t>t$.

\begin{lemma}[Left continuity of the index]
  \label{lem:dc-ch11-2-7}
  \dcref{ch11:2.7}
  If $\varepsilon>0$ is sufficiently small, then
  \[
    i(t-\varepsilon)=i(t).
  \]
  If $d$ denotes the nullity of $I_t$, then $d=0$ whenever $\gamma(t)$ is not
  conjugate to $\gamma(0)$.
\end{lemma}
\begin{proof}
  \sketch
  Monotonicity gives $i(t-\varepsilon)\leq i(t)$. On a subspace of dimension
  $i(t)$ where $I_t$ is negative definite, continuity of the family $I_s$ on the
  fixed space $S_j$ preserves negative definiteness for $s=t-\varepsilon$ close
  to $t$, giving the reverse inequality.
\end{proof}

\begin{lemma}[Index jump at a conjugate point]
  \label{lem:dc-ch11-2-8}
  \uses{thm:dc-ch11-2-2}
  \dcref{ch11:2.8}
  If $d$ is the nullity of $I_t$, then for sufficiently small $\varepsilon>0$,
  \[
    i(t+\varepsilon)=i(t)+d.
  \]
\end{lemma}
\begin{proof}
  \sketch
  Since $\dim S_j=n(j-1)$, $I_t$ is positive definite on a subspace of dimension
  $n(j-1)-i(t)-d$. Continuity preserves this subspace for $t+\varepsilon$, so
  $i(t+\varepsilon)\leq i(t)+d$. For the opposite inequality use the Index Lemma
  (\cref{lem:dc-ch10-2-2}): for a broken Jacobi field $V_{t_0}$ vanishing at
  $t_0\in(t_{j-1},t_j)$, extend it by zero on $[t_0,t_0+\varepsilon]$. The Index
  Lemma gives
  \[
    I_{t_0}(V_{t_0},V_{t_0})>
    I_{t_0+\varepsilon}(V_{t_0+\varepsilon},V_{t_0+\varepsilon}),
  \]
  with strictness because the extension is not Jacobi. Hence every vector with
  $I_t(V,V)\leq0$ becomes negative for $I_{t+\varepsilon}$; a negative subspace
  together with the null space of $I_t$ is therefore negative, yielding
  $i(t+\varepsilon)\geq i(t)+d$.
\end{proof}


\begin{corollary}[Jacobi's criterion]
  \label{cor:dc-ch11-2-9}
  \dcref{ch11:2.9}
  Let $\gamma:[0,a]\to M$ be a geodesic with $\gamma(a)$ not conjugate to
  $\gamma(0)$. Then $\gamma$ has no conjugate point in $(0,a)$ if and only if every
  proper variation has $E(s)\geq E(0)$ for all sufficiently small $|s|$. In
  particular, a minimizing geodesic has no conjugate point in its interior.
\end{corollary}

\begin{corollary}[Discreteness of conjugate points]
  \label{cor:dc-ch11-2-10}
  \dcref{ch11:2.10}
  The conjugate points along a geodesic form a discrete subset of the geodesic.
\end{corollary}

\begin{remark}[Finite-dimensional model of the path space]
  \label{rem:dc-ch11-2-12}
  \uses{rem:dc-ch9-2-7, cor:dc-ch11-2-6}
  \dcref{ch11:2.12}
  For the path space $\Omega_{p,q}$ and $c>0$, let $\Omega^c$ (respectively
  $\mathring{\Omega}{}^c$) consist of paths of energy at most $c$ (respectively
  strictly less than $c$). These paths lie in a compact set $S\subset M$. Choose
  $\delta>0$ so that points of $S$ at distance below $\delta$ have a unique
  minimizing geodesic, and subdivide $[0,1]$ by
  $|t_i-t_{i-1}|<\delta^2/c$. Let $B\subset\Omega^c$ be the broken geodesics
  whose pieces join consecutive subdivision points, and let $\mathring B$ be its
  intersection with $\mathring{\Omega}{}^c$. Since
  \[
    L^2(w_i)=(t_i-t_{i-1})E(w_i)<\delta^2,
  \]
  each piece is determined by its endpoints, and
  \[
    w\longmapsto(w(t_1),\ldots,w(t_{k-1}))
  \]
  identifies $\mathring B$ with an open subset of $M^{k-1}$. Its tangent vectors
  are broken Jacobi fields. There is a deformation retract
  $h_s:\mathring{\Omega}{}^c\to\mathring{\Omega}{}^c$ with image $\mathring B$.
  The restriction of the energy to $\mathring B$ has precisely the geodesics as
  critical points, and at a geodesic its Hessian has the same index and nullity as
  the index form by \cref{cor:dc-ch11-2-6}. Thus $\mathring B$ is a finite-dimensional
  model for the sublevel path space.
\end{remark}

\begin{remark}[Remark 2.11]
  \label{rem:dc-ch11-2-11}
  \dcref{ch11:2.11}
  The index theorem can be generalized to the case in which the points $\gamma(0)$
  and $\gamma(a)$ are replaced by submanifolds. For this see W. Ambrose [Am 2]. For
  an extremely general version of the index theorem, see S. Smale [Sm].
\end{remark}

\chapter{The Fundamental Group of Manifolds of Negative Curvature}
\label{chap:dc-fundamental-group-negative-curvature}

Let $M^n$ be complete with sectional curvature $K<0$. By the Hadamard theorem,
its universal cover is diffeomorphic to $\mathbb{R}^n$; consequently
$\pi_k(M)=0$ for $k\geq 2$. The results below constrain $\pi_1(M)$ when $M$ is
compact.

\section{Existence of closed geodesics}

Let $\pi:\tilde M\to M$ be the universal covering with the covering metric, and
let $A(\tilde M)$ be its group of covering transformations. Every nonidentity
element of $A(\tilde M)$ is a fixed-point-free isometry, and a choice of
$\tilde p\in\pi^{-1}(p)$ determines an isomorphism
$\pi_1(M;p)\cong A(\tilde M)$. Explicitly, for $g\in\pi_1(M;p)$ and
$\tilde q\in\tilde M$, choose a path $\tilde\sigma$ from $\tilde q$ to
$\tilde p$, set $\sigma=\pi\circ\tilde\sigma$ and
$\beta=\sigma g\sigma^{-1}$, and define
\[
  \alpha_{\tilde p}(\tilde q)=\tilde\beta_{\tilde q}(1),
\]
where $\tilde\beta_{\tilde q}$ is the lift of $\beta$ starting at $\tilde q$.

\begin{definition}[Free homotopy class]
  \label{def:dc-ch12-2-1}
  \dcref{ch12:2.1}
  Closed paths $f,g:I\to M$ are \emph{freely homotopic} if there is a map
  $F:I\times I\to M$ such that
  \[
    F(0,t)=f(t),\qquad F(1,t)=g(t),\qquad F(s,0)=F(s,1).
  \]
  The corresponding equivalence classes form the set $C_1(M)$. Unlike based
  homotopy classes, free homotopy classes allow the base point to vary.

  A \emph{closed geodesic} is a closed curve geodesic at every point. A
  \emph{geodesic lasso} is closed and geodesic except at one nonregular point.
\end{definition}

\begin{theorem}[Cartan's closed-geodesic theorem]
  \label{thm:dc-ch12-2-2}
  \dcref{ch12:2.2}
  If $M$ is compact and $\mathcal L\in C_1(M)$ is nonconstant, then
  $\mathcal L$ contains a closed geodesic.
\end{theorem}
\begin{proof}
  \sketch
  Let $d$ be the infimum of the lengths of piecewise differentiable curves in
  $\mathcal L$; nontriviality gives $d>0$. Choose broken geodesics
  $\gamma_j:[0,1]\to M$, parametrized proportionally to arc length, with
  $\ell(\gamma_j)\to d$. If $L=\sup_j\ell(\gamma_j)$, then
  \[
    d(\gamma_j(t_1),\gamma_j(t_2))
    \leq \int_{t_1}^{t_2}|\gamma_j'(t)|\,dt
    \leq L(t_2-t_1).
  \]
  Compactness and equicontinuity give a uniformly convergent subsequence with
  closed limit $\gamma_0$.

  Choose $0=t_0<\cdots<t_k=1$ so that each corresponding arc of $\gamma_0$
  lies in a totally normal neighborhood, and join successive points
  $\gamma_0(t_i)$ by the unique minimizing segments to obtain
  $\gamma\in\mathcal L$. If $\ell(\gamma)>d$, set
  $\varepsilon=(\ell(\gamma)-d)/(2k+1)$. For large $j$,
  \[
    \ell(\gamma_j)-d<\varepsilon,
    \qquad d(\gamma_j(t),\gamma_0(t))<\varepsilon,
  \]
  so
  \[
    \sum_{i=1}^k\bigl(\ell(\gamma_j^i)+2\varepsilon\bigr)
    <\ell(\gamma)=\sum_{i=1}^k\ell(\gamma^i).
  \]
  Some minimizing segment $\gamma^i$ would then be longer than a competitor,
  a contradiction. Thus $\ell(\gamma)=d$. After arc-length parametrization,
  $\gamma$ has no corner: a corner inside a convex ball could be shortened by
  replacing the two incident arcs with the minimizing chord, without changing
  the free homotopy class. Hence $\gamma$ is a closed geodesic.
\end{proof}

\begin{remark}[Failure without compactness]
  \label{rem:dc-ch12-2-3}
  \dcref{ch12:2.3}
  Compactness is necessary.
\end{remark}
\begin{proof}

  \sketch
  A surface of revolution generated by a curve
  asymptotic to its axis has nontrivial free homotopy classes containing curves
  of arbitrarily small length, so those classes have no length minimizer.
\end{proof}

\begin{remark}[Simply connected compact case]
  \label{rem:dc-ch12-2-4}
  \dcref{ch12:2.4}
  A compact simply connected Riemannian manifold also has a closed geodesic,
  although this does not follow from a nonconstant free homotopy class.
\end{remark}

\begin{definition}[Translation along a geodesic]
  \label{def:dc-ch12-2-5}
  \dcref{ch12:2.5}
  A fixed-point-free isometry $f:\tilde M\to\tilde M$ is a
  \emph{translation} if some geodesic $\tilde\gamma$ is invariant under $f$,
  meaning
  \[
    f\bigl(\tilde\gamma(\mathbb R)\bigr)=\tilde\gamma(\mathbb R).
  \]
  In that case $f$ is a \emph{translation along} $\tilde\gamma$.
\end{definition}

\begin{proposition}[Covering transformations are translations]
  \label{prop:dc-ch12-2-6}
  \dcref{ch12:2.6}
  Let $M$ be compact. Every nonidentity covering transformation of
  $\tilde M$ is a translation.
\end{proposition}
\begin{proof}
  \sketch
  Let $\alpha$ correspond to $g\in\pi_1(M;p)$ after choosing
  $\tilde p\in\pi^{-1}(p)$. By \cref{thm:dc-ch12-2-2}, the associated free
  homotopy class contains a closed geodesic $\gamma$. If $q\in\gamma$ and
  $\sigma$ joins $p$ to $q$, lift $\sigma$ from $\tilde p$ to an endpoint
  $\tilde q$, and lift $\gamma$ from $\tilde q$. The deck transformation
  $\alpha_{\tilde q}$ associated with $[\gamma]\in\pi_1(M;q)$ satisfies
  \[
    \alpha(\tilde q)=\alpha_{\tilde q}(\tilde q).
  \]
  Thus $\alpha\alpha_{\tilde q}^{-1}$ has a fixed point and is the identity, so
  $\alpha=\alpha_{\tilde q}$. Uniqueness of path lifting then yields
  \[
    \alpha(\tilde\gamma(s))
    =\alpha_{\tilde q}(\tilde\gamma(s))\in\tilde\gamma(\mathbb R),
  \]
  proving invariance.
\end{proof}

\section{Preissman's Theorem}

A \emph{geodesic triangle} consists of three unit-speed minimizing geodesic
segments
\[
  \gamma_1:[0,\ell_1]\to M,\qquad
  \gamma_2:[0,\ell_2]\to M,\qquad
  \gamma_3:[0,\ell_3]\to M,
\]
with $\gamma_i(\ell_i)=\gamma_{i+1}(0)$ for $i=1,2$ and
$\gamma_3(\ell_3)=\gamma_1(0)$. Its interior angles are
\[
  \sphericalangle(-\gamma_i'(\ell_i),\gamma_{i+1}'(0))\quad(i=1,2),
  \qquad
  \sphericalangle(-\gamma_3'(\ell_3),\gamma_1'(0)).
\]

\begin{lemma}[Comparison for geodesic triangles when $K\leq 0$]
  \label{lem:dc-ch12-3-1}
  \uses{prop:dc-ch10-2-5}
  \dcref{ch12:3.1}
  Let $\tilde M$ be complete, simply connected, and satisfy $K\leq0$. Points
  $a,b,c\in\tilde M$ determine a unique geodesic triangle. Let
  $\alpha,\beta,\gamma$ be the angles at $a,b,c$, and let $A,B,C$ be the
  lengths of the opposite sides. Then
  \begin{enumerate}
    \item $A^2+B^2-2AB\cos\gamma\leq C^2$;
    \item $\alpha+\beta+\gamma\leq\pi$.
  \end{enumerate}
  Both inequalities are strict when $K<0$.
\end{lemma}
\begin{proof}
  \sketch
  Let $\gamma_A,\gamma_B,\gamma_C$ be the sides of lengths $A,B,C$, with the
  first two issuing from $c$, and set
  $\Gamma_i=\exp_c^{-1}\circ\gamma_i$ for $i=A,B,C$. Let $\Gamma_0$ be the
  straight segment joining the endpoints of $\Gamma_C$. The radial sides satisfy
  \[
    A=\ell(\Gamma_A),\qquad B=\ell(\Gamma_B),
  \]
  while Euclidean geometry in $T_c\tilde M$ gives
  \[
    \ell(\Gamma_0)^2=A^2+B^2-2AB\cos\gamma.
  \]
  Since $\ell(\Gamma_0)\leq\ell(\Gamma_C)$, Rauch comparison
  (\cref{prop:dc-ch10-2-5}) gives
  \[
    A^2+B^2-2AB\cos\gamma
    \leq\ell(\Gamma_C)^2
    \leq\ell(\gamma_C)^2=C^2,
  \]
  with strict final inequality when $K<0$. This proves the first assertion.

  The side lengths satisfy the triangle inequalities, so they determine a
  Euclidean comparison triangle with opposite angles
  $\alpha',\beta',\gamma'$. Applying the first assertion at each vertex yields
  $\alpha\leq\alpha'$, $\beta\leq\beta'$, and
  $\gamma\leq\gamma'$, all strict when $K<0$. Since
  $\alpha'+\beta'+\gamma'=\pi$, the angle-sum assertion follows.
\end{proof}

For the remainder of the chapter, $M$ is complete with $K<0$, and
$\pi:\tilde M\to M$ is its universal covering with the covering metric.

\begin{theorem}[Preissman's cyclic-subgroup theorem]
  \label{thm:dc-ch12-3-2}
  \dcref{ch12:3.2}
  If $M$ is compact, every nontrivial abelian subgroup of $\pi_1(M)$ is
  infinite cyclic.
\end{theorem}
\begin{proof}
  \sketch
  Choose a nonidentity element of the abelian subgroup $H$. By
  \cref{prop:dc-ch12-2-6}, it translates along a geodesic
  $\tilde\gamma$. The uniqueness in \cref{lem:dc-ch12-3-3} and commutativity,
  through \cref{lem:dc-ch12-3-4}, imply that every element of $H$ preserves
  $\tilde\gamma$. Now \cref{lem:dc-ch12-3-5} shows that $H$ is infinite
  cyclic.
\end{proof}

\begin{lemma}[Uniqueness of an invariant geodesic]
  \label{lem:dc-ch12-3-3}
  \uses{lem:dc-ch12-3-1}
  \dcref{ch12:3.3}
  If $K<0$ and a nonidentity translation $f:\tilde M\to\tilde M$ preserves
  $\tilde\gamma$, then $\tilde\gamma$ is the unique geodesic preserved by $f$.
\end{lemma}
\begin{proof}
  \sketch
  Suppose $f$ preserves distinct geodesics $\tilde\gamma_1$ and
  $\tilde\gamma_2$. They are disjoint: if their intersection were a single
  point, it would be fixed by $f$, while two intersection points would force
  the geodesics to coincide. Join
  $\tilde p_i\in\tilde\gamma_i$ by a minimizing geodesic. The resulting
  quadrilateral with vertices
  $\tilde p_1,f(\tilde p_1),f(\tilde p_2),\tilde p_2$ has interior angles
  $\alpha,\pi-\alpha',\beta,\pi-\beta'$. Since $f$ is an isometry,
  $\alpha=\alpha'$ and $\beta=\beta'$, so their sum is $2\pi$. Dividing the
  quadrilateral into two triangles forces one triangle to have angle sum at
  least $\pi$, contradicting the strict case of
  \cref{lem:dc-ch12-3-1}.
\end{proof}

\begin{lemma}[A commuting isometry preserves the translation axis]
  \label{lem:dc-ch12-3-4}
  \dcref{ch12:3.4}
  If $K<0$, $f$ is a nonidentity translation along $\tilde\gamma$, and the
  fixed-point-free isometry $g$ commutes with $f$, then $g$ is a translation
  along $\tilde\gamma$.
\end{lemma}
\begin{proof}
  \sketch
  Commutativity gives
  \[
    f\bigl(g(\tilde\gamma)\bigr)
    =g\bigl(f(\tilde\gamma)\bigr)=g(\tilde\gamma).
  \]
  Thus $g(\tilde\gamma)$ is another $f$-invariant geodesic, so
  \cref{lem:dc-ch12-3-3} gives $g(\tilde\gamma)=\tilde\gamma$.
\end{proof}

\begin{lemma}[A subgroup preserving one geodesic is infinite cyclic]
  \label{lem:dc-ch12-3-5}
  \uses{prop:dc-ch12-2-6}
  \dcref{ch12:3.5}
  If every element of a nontrivial subgroup $H\subset\pi_1(M)$, acting on
  $\tilde M$, preserves a fixed geodesic $\tilde\gamma$, then $H$ is infinite
  cyclic.
\end{lemma}
\begin{proof}
  \sketch
  Orient $\tilde\gamma$, fix $\tilde p\in\tilde\gamma$, and assign to
  $h\in H$ the signed displacement
  \[
    \theta(h)=\pm d(\tilde p,h(\tilde p)).
  \]
  Because every $h$ acts isometrically on $\tilde\gamma$, this defines a
  homomorphism $\theta:H\to(\mathbb R,+)$. It is injective: equality of two
  images would make $h_1h_2^{-1}$ fix $\tilde p$. Hence $H$ is isomorphic to
  an additive subgroup of $\mathbb R$. Such a subgroup is either dense or
  infinite cyclic; proper discontinuity of the deck action excludes density.
\end{proof}

\begin{example}[Surface of genus two times a circle]
  \label{ex:dc-ch12-3-6}
  \dcref{ch12:3.6}
  Let $N$ be a surface of genus two. Then $N\times S^1$ admits no negatively
  curved metric.
\end{example}
\begin{proof}

  \sketch
  For any infinite cyclic subgroup $C\subset\pi_1(N)$, the subgroup
  $C\oplus\mathbb Z\subset\pi_1(N\times S^1)$ is abelian but not cyclic,
  contradicting \cref{thm:dc-ch12-3-2}.
\end{proof}

\begin{example}[The $m$-torus]
  \label{ex:dc-ch12-3-7}
  \dcref{ch12:3.7}
  For $m\geq2$, the torus $T^m=(S^1)^m$ admits no negatively curved metric.
\end{example}
\begin{proof}

  \sketch
  Its fundamental group is $\mathbb Z^m$, a noncyclic abelian group, contrary
  to \cref{thm:dc-ch12-3-2}.
\end{proof}

\begin{theorem}[Nonabelian fundamental group]
  \label{thm:dc-ch12-3-8}
  \dcref{ch12:3.8}
  If $M$ is compact and $K<0$, then $\pi_1(M)$ is not abelian.
\end{theorem}
\begin{proof}
  \sketch
  Suppose $\pi_1(M)$ were abelian. If it were trivial, every geodesic in
  $\tilde M$ would be invariant under all deck transformations, contradicting
  \cref{lem:dc-ch12-3-9}. Otherwise, choose a nonidentity deck transformation
  and its unique invariant geodesic. By \cref{lem:dc-ch12-3-4}, every deck
  transformation preserves that geodesic, again contradicting
  \cref{lem:dc-ch12-3-9} and compactness.
\end{proof}

\begin{lemma}[An invariant geodesic forces noncompactness]
  \label{lem:dc-ch12-3-9}
  \uses{lem:dc-ch12-3-1}
  \dcref{ch12:3.9}
  If $M$ is complete, $K\leq0$, and some geodesic in $\tilde M$ is invariant
  under every element of $A(\tilde M)$, then $M$ is noncompact.
\end{lemma}
\begin{proof}
  \sketch
  Let $\tilde\gamma$ be invariant, fix $\tilde p\in\tilde\gamma$, and for
  $t>0$ let $\tilde\beta:[0,t]\to\tilde M$ be the unit-speed geodesic from
  $\tilde p$ perpendicular to $\tilde\gamma$. Write
  $\beta=\pi\circ\tilde\beta$, $p=\pi(\tilde p)$, and let $\alpha_t$ be a
  minimizing geodesic from $\beta(t)$ to $p$. The lift $\tilde\alpha_t$ from
  $\tilde\beta(t)$ ends on $\tilde\gamma$ because the latter is invariant under
  every deck transformation. The comparison inequality
  \cref{lem:dc-ch12-3-1} gives
  $\ell(\tilde\alpha_t)\geq\ell(\tilde\beta)=t$, whereas
  \[
    \ell(\tilde\alpha_t)=\ell(\alpha_t)
    \leq\ell(\beta)=\ell(\tilde\beta)=t.
  \]
  Thus $d(\beta(t),p)=t$ for every $t>0$, so $M$ is unbounded and hence
  noncompact.
\end{proof}

\begin{theorem}[Byers's solvable-subgroup theorem]
  \label{thm:dc-ch12-3-10}
  \uses{thm:dc-ch12-3-8, lem:dc-ch12-3-5}
  \dcref{ch12:3.10}
  If $M$ is compact, $K<0$, and $H\neq\{e\}$ is a solvable subgroup of
  $\pi_1(M)$, then $H$ is infinite cyclic. Moreover, $\pi_1(M)$ has no cyclic
  subgroup of finite index.
\end{theorem}
\begin{proof}
  \sketch
  Choose a shortest solvable series
  \[
    H=H_0\supset H_1\supset\cdots\supset H_{k-1}\supset H_k=\{e\},
  \]
  where $H_{i+1}\triangleleft H_i$ and $H_i/H_{i+1}$ is abelian. Then
  $H_{k-1}$ is nontrivial and abelian, hence infinite cyclic by
  \cref{thm:dc-ch12-3-2}. Let $g$ generate it and let $\tilde\gamma$ be the
  $g$-invariant geodesic. If $k=1$, this already proves the first assertion.
  Assume $k\geq2$. For $a\in H_{k-2}$ and nonidentity
  $b\in H_{k-1}$, normality gives an integer $n$ such that
  \[
    a^{-1}b^{-1}ab=g^n.
  \]
  Since both $b$ and $g^n$ preserve $\tilde\gamma$, $b^{-1}$ preserves
  $a(\tilde\gamma)$. Uniqueness of the invariant geodesic gives
  $a(\tilde\gamma)=\tilde\gamma$. Hence every element of $H_{k-2}$ preserves
  $\tilde\gamma$, and \cref{lem:dc-ch12-3-5} makes $H_{k-2}$ infinite
  cyclic. Repeating upward through the series proves that $H$ is infinite
  cyclic.

  Suppose now that a cyclic subgroup $H\subset\pi_1(M)$ has finite index. If
  $H=\{e\}$, then $\pi_1(M)$ is finite; \cref{thm:dc-ch12-3-2} forces it to
  be trivial, contrary to \cref{thm:dc-ch12-3-8}. Thus
  $H=\langle g\rangle$ with $g\neq e$. It is proper, since otherwise
  $\pi_1(M)$ would be abelian. If $\tilde\gamma$ is the $g$-invariant geodesic
  and $a\notin H$, finite index gives $a^n\in H$ for some $n>0$.
  Preissman's theorem applied to $\langle a\rangle$ implies $a^n\neq e$, so
  $a^n=g^m$ for some $m\neq0$. Both $\tilde\gamma$ and
  $a(\tilde\gamma)$ are invariant under $a^n$, and uniqueness yields
  $a(\tilde\gamma)=\tilde\gamma$. Thus all of $\pi_1(M)$ preserves
  $\tilde\gamma$, so \cref{lem:dc-ch12-3-5} would make $\pi_1(M)$ infinite
  cyclic, contradicting \cref{thm:dc-ch12-3-8}.
\end{proof}

\begin{remark}[Further nonpositive-curvature structure]
  \label{rem:dc-ch12-3-11}
  \dcref{ch12:3.11}
  The results above use only part of the structure relating the geometry of a
  nonpositively curved manifold to its fundamental group.
\end{remark}

\chapter{The Sphere Theorem}
\label{chap:dc-sphere-theorem}

\begin{theorem}[Sphere Theorem]
  \label{thm:dc-ch13-1-1}
  \dcref{ch13:1.1}
  Let $M^n$ be a compact simply connected Riemannian manifold, whose sectional
  curvature $K$ satisfies

  $$
  0<hK_{\max}<K\leq K_{\max}.\qquad(1)
  $$

  Then if $h=1/4$, $M$ is homeomorphic to a sphere.

  The number $h$ is the pinching constant. Rescaling the metric gives the equivalent
  normalization

  $$
  0<h<K\leq 1.\qquad(1')
  $$

  In even dimension the non-strict bound

  $$
  0<1/4\leq K\leq1,\tag{2}
  $$
  admits simply connected examples not homeomorphic to a sphere in even dimensions.
  Under (2), the alternatives are $\operatorname{diam}(M)>\pi$ with $M$ homeomorphic
  to a sphere, or $\operatorname{diam}(M)=\pi$ with $M$ isometric to a symmetric
  space. In odd dimensions (2) still implies the conclusion. For $n=2,3$, positive
  sectional curvature suffices for a compact simply connected manifold to be
  homeomorphic to $S^n$.

  Unless otherwise stated, manifolds are complete and geodesics are normalized.
\end{theorem}
\begin{proof}
  \sketch
  Let $p,q\in M$ with $\operatorname{diam}M=d(p,q)$. Let $D_1$,
  $D_2$ be the subsets of $M$ formed by all geodesic segments $\overline{pm}$ and
  $\overline{qn}$ respectively, with $m,n$ from \cref{lem:dc-ch13-4-3}. By continuity of $m$, $D_1,D_2$ are closed; we claim $D_1\cup D_2=M$ and
  $\partial D_1=\partial D_2=D_1\cap D_2$. Indeed, for $r\in M$, by \cref{lem:dc-ch13-4-2} either
  $d(p,r)<\rho$ or $d(q,r)<\rho$; in the first case, since $d(p,C_m(p))\geq\pi>\rho$,
  there is a unique minimizing $\gamma$ through $p,r$ and a unique $m$ on $\gamma$ with
  $d(p,m)=d(q,m)<\rho$, and $r\in\overline{pm}=D_1$ if $d(p,r)<d(q,r)$, $r\in\partial D_1$
  if $d(p,r)=d(q,r)$ (then $r=m$), and $r\in D_2$ if $d(p,r)>d(q,r)$; so
  $r\in D_1\cup D_2$, and $r\in D_1\cap D_2$ forces $d(p,r)=d(q,r)$, $r=m=n$.

  Define $\varphi:S^n\to M^n$ as follows. Associate to a fixed $N\in S^n$ the point
  $p$, and to its antipode $S$ the point $q$. Choose a linear isometry
  $i:T_NS^n\to T_pM$. For each $e$ in the equator $E$ of $S^n$ relative to $N$, take
  the geodesic $\gamma(s)$ of $S^n$, $0\leq s\leq\pi$, with $\gamma(0)=N$,
  $\gamma(\pi/2)=e$; let $m$ be the point of \cref{lem:dc-ch13-4-3} on the geodesic of $M$ through
  $p$ with velocity $i\gamma'(0)$. Define

  $$
  \varphi(\gamma(s))=\exp_p\!\Big(s\tfrac{2}{\pi}d(p,m)\,(i\circ\gamma'(0))\Big),\quad 0<s\leq\tfrac{\pi}{2},
  $$

  $$
  \varphi(\gamma(s))=\exp_q\!\Big(\big(2-\tfrac{2s}{\pi}\big)d(q,m)\,V\Big),\quad\tfrac{\pi}{2}\leq s<\pi,
  $$

  where $V$ is the unit tangent at $q$ of the unique minimizing geodesic from $q$ to
  $m$. Then $\varphi$ maps the closed northern hemisphere bijectively onto $D_1$, the
  closed southern hemisphere bijectively onto $D_2$, and $E$ onto
  $\partial D_1=\partial D_2=D_1\cap D_2$. By uniqueness of $m,n$ (\cref{lem:dc-ch13-4-3}),
  $\varphi$ is continuous; since $M=D_1\cup D_2$, $\varphi$ is surjective; since
  $D_1\cap D_2=\partial D_1=\varphi(E)$, $\varphi$ is injective. As $S^n$ is compact,
  $\varphi$ is a homeomorphism.
\end{proof}

\section{The cut locus}

Let $M$ be complete, $p\in M$, and $\gamma:[0,\infty)\to M$ a normalized geodesic
with $\gamma(0)=p$. For small $t>0$, $d(\gamma(0),\gamma(t))=t$, and once a segment
ceases to be minimizing no later segment is minimizing. Thus the set of $t>0$ with
$d(\gamma(0),\gamma(t))=t$ is of the form $[0,t_0]$ or $[0,\infty)$. In the first
case $\gamma(t_0)$ is called the \emph{cut point of $p$ along $\gamma$}; in the second
no cut point exists. The \emph{cut locus of $p$}, denoted $C_m(p)$, is the union of the
cut points of $p$ along all geodesics starting from $p$. If $M$ is compact, a cut
point exists along every geodesic from $p$; the converse holds by
\cref{cor:dc-ch13-2-11}.

\begin{proposition}[Characterization of a cut point]
  \label{prop:dc-ch13-2-2}
  \uses{prop:dc-ch13-2-9}
  \dcref{ch13:2.2}
  Suppose that $\gamma(t_0)$ is the cut point of $p=\gamma(0)$ along $\gamma$. Then:
  (a) either $\gamma(t_0)$ is the first conjugate point of $\gamma(0)$ along
    $\gamma$,
  (b) or there exists a geodesic $\sigma\neq\gamma$ from $p$ to $\gamma(t_0)$ such
    that $\ell(\sigma)=\ell(\gamma)$.

  Conversely, if (a) or (b) holds, then there exists $\tilde t\in(0,t_0]$ such that
  $\gamma(\tilde t)$ is the cut point of $p$ along $\gamma$.
\end{proposition}
\begin{proof}
\sketch
  Take a sequence $\{t_0+\varepsilon_i\}$, $\varepsilon_i>0$,
  $\varepsilon_i\to 0$, and minimizing geodesics $\sigma_i$ joining $p$ to
  $\gamma(t_0+\varepsilon_i)$. The unit sphere in $T_pM$ being compact, pass to a
  subsequence with $\sigma_i'(0)\to\sigma'(0)$; by continuity $\sigma$ is a
  minimizing geodesic from $p$ to $\gamma(t_0)$, so $\ell(\sigma)=\ell(\gamma)$. If
  $\sigma\neq\gamma$, (b) holds. If $\sigma=\gamma$, we show $d\exp_p$ is singular at
  $t_0\gamma'(0)$: otherwise $\exp_p$ is a diffeomorphism on a neighborhood $U$ of
  $t_0\gamma'(0)$, and writing $\gamma(t_0+\varepsilon_j)=\sigma_j(t_0+\varepsilon_j')$
  with $\varepsilon_j'\leq\varepsilon_j$ gives
  $\exp_p(t_0+\varepsilon_j)\gamma'(0)=\exp_p(t_0+\varepsilon_j')\sigma_j'(0)$, hence
  $\gamma'(0)=\sigma_j'(0)$, contradicting the definition of $t_0$. Conversely, if
  (a) holds, since a geodesic does not minimize after the first conjugate point
  (\cref{cor:dc-ch11-2-9}), the cut point occurs at
  $\gamma(\tilde t)$, $\tilde t\leq t_0$. If (b) holds, joining $\sigma(t_0-\varepsilon)$
  to $\gamma(t_0+\varepsilon)$ by the unique minimizing geodesic $\tau$ inside a
  totally normal neighborhood of $\gamma(t_0)$ yields a curve of length strictly less
  than $t_0+\varepsilon$, so again the cut point occurs at $\gamma(\tilde t)$,
  $\tilde t\leq t_0$.
\end{proof}

\begin{example}[Cut locus of a sphere]
  \label{ex:dc-ch13-2-3}
  \dcref{ch13:2.3}
  If $M^n=S^n$ and $p\in S^n$, the cut locus of $p$ is its antipodal point. Here the
  cut locus of $p$ coincides with the locus of conjugate points to $p$.
\end{example}

\begin{example}[Cut and conjugate loci of real projective space]
  \label{ex:dc-ch13-2-4}
  \dcref{ch13:2.4}
  If $M^n=P^n(\mathbb{R})$, obtained by identifying antipodal points of $S^n$, the
  cut locus of $[p,-p]$ is the subset $P^{n-1}(\mathbb{R})\subset P^n(\mathbb{R})$
  obtained by identifying the antipodal points of the "equator" of $p$. The conjugate
  locus of $[p,-p]$ is $[p,-p]$ itself.
\end{example}

\begin{example}[Cut loci in flat torus and cylinder]
  \label{ex:dc-ch13-2-5}
  \dcref{ch13:2.5}
  On the flat torus $T$ obtained by identifying opposite sides of a quadrilateral
  $ABCD$, the cut locus of $A$ is the set formed by the medians of the segments $BA$
  and $BC$; the conjugate locus is empty. Similarly, the cut locus of a point $p$ of
  a cylinder in $\mathbb{R}^3$ is the "opposite generator" to that through $p$.
\end{example}

\begin{remark}[Cut and conjugate alternatives]
  \label{rem:dc-ch13-2-6}
  \uses{ex:dc-ch13-2-3, ex:dc-ch13-2-4, ex:dc-ch13-2-5, prop:dc-ch13-2-2}
  \dcref{ch13:2.6}
  In \cref{ex:dc-ch13-2-3} the situations (a) and (b) of \cref{prop:dc-ch13-2-2} occur simultaneously;
  in \cref{ex:dc-ch13-2-4} situation (b) occurs before (a); in \cref{ex:dc-ch13-2-5} only situation (b)
  occurs. For the ellipsoid, except for two directions, (b) occurs before (a).
\end{remark}

\begin{corollary}[Symmetry of cut points]
  \label{cor:dc-ch13-2-7}
  \uses{prop:dc-ch13-2-2}
  \dcref{ch13:2.7}
  If $q$ is a cut point of $p$ along $\gamma$, then $p$ is the cut point of $q$ along
  $-\gamma$; in particular $q\in C_m(p)$ if and only if $p\in C_m(q)$.
\end{corollary}
\begin{proof}
\sketch
  By \cref{prop:dc-ch13-2-2}, either $q$ is conjugate to $p$, or there is a geodesic
  $\sigma\neq\gamma$ from $p$ to $q$ with $\ell(\sigma)=\ell(\gamma)=d(p,q)$. In both
  cases the cut point of $q$ along $-\gamma$ does not occur before $p$; since
  $\ell(-\gamma)=d(p,q)$, $p$ is that cut point.
\end{proof}

\begin{corollary}[Injectivity radius]
  \label{cor:dc-ch13-2-8}
  \dcref{ch13:2.8}
  If $q\in M-C_m(p)$, there exists a unique minimizing geodesic joining $p$ to $q$.

  \cref{cor:dc-ch13-2-8} shows that $\exp_p$ is injective on the open ball $B_r(p)$ if and only
  if $r\leq d(p,C_m(p))$. For this reason

  $$
  i(M)=\inf_{p\in M}d(p,C_m(p))
  $$

  is called the \emph{injectivity radius of $M$}. Moreover, $M-C_m(p)$ is
  homeomorphic to an open Euclidean ball.
\end{corollary}

\begin{proposition}[Continuity of cut distance]
  \label{prop:dc-ch13-2-9}
  \dcref{ch13:2.9}
  Define $f:T_1M\to\mathbb{R}\cup\{\infty\}$ by $f(\gamma(0),\gamma'(0))=t_0$ if
  $\gamma(t_0)$ is the cut point of $\gamma(0)$ along $\gamma$, and
  $f(\gamma(0),\gamma'(0))=\infty$ if no cut point exists, where $\mathbb{R}\cup\{\infty\}$
  carries the topology with base the open intervals $(a,b)$ together with the sets
  $(a,\infty]=(a,\infty)\cup\{\infty\}$. Then $f$ is continuous.
\end{proposition}
\begin{proof}
\sketch
  Let $\gamma_i(0)\to\gamma(0)$, $\gamma_i'(0)\to\gamma'(0)$, with cut points
  $\gamma_i(t_o^i)$ and $\gamma(t_o)$. We show $\lim t_o^i=t_o$. First
  $\limsup t_o^i\leq t_o$: if $t_o<\infty$ and $\varepsilon>0$, only finitely many $j$
  can have $t_o+\varepsilon<t_o^j$, else $d(\gamma_j(0),\gamma_j(t_o+\varepsilon))=t_o+\varepsilon$
  gives in the limit $d(\gamma(0),\gamma(t_o+\varepsilon))=t_o+\varepsilon$,
  contradicting that $\gamma(t_o)$ is a cut point. Now let
  $\bar t=\liminf t_o^i\leq t_o$; suppose $\bar t<\infty$ and pass to a subsequence
  $t_o^j\to\bar t$. If for a subsequence the $\gamma_j(t_o^j)$ are conjugate to
  $\gamma_j(0)$, then (accumulation of conjugate points is conjugate) $\gamma(\bar t)$
  is conjugate to $\gamma(0)$, so $\bar t\geq t_o$. Otherwise, by \cref{prop:dc-ch13-2-2},
  geodesics $\sigma_j\neq\gamma_j$ exist with $\sigma_j(0)=\gamma_j(0)$,
  $\sigma_j(t_o^j)=\gamma_j(t_o^j)$, $\ell(\sigma_j)=\ell(\gamma_j)$; a subsequence
  $\sigma_j\to\sigma$ joins $\gamma(0)$ to $\gamma(\bar t)$. If $\sigma\neq\gamma$,
  \cref{prop:dc-ch13-2-2} gives $t_o\leq\bar t$; if $\sigma=\gamma$, a similar argument shows
  $\gamma(\bar t)$ conjugate to $\gamma(0)$, so $\bar t\geq t_o$.
\end{proof}

\begin{corollary}[Closedness of the cut locus]
  \label{cor:dc-ch13-2-10}
  \uses{prop:dc-ch13-2-9}
  \dcref{ch13:2.10}
  For all $p\in M$, $C_m(p)$ is closed; in particular, if $M$ is compact, $C_m(p)$ is
  compact.
\end{corollary}
\begin{proof}
\sketch
  One verifies $C_m(p)=\{\gamma(t);t=f(p,\gamma'(0))\}$, where $f$ is the
  function of \cref{prop:dc-ch13-2-9}. If $q$ is an accumulation point of $C_m(p)$, take
  $\gamma_j(t_j)\to q$ with $t_j=f(p,\gamma_j'(0))$ and $\gamma_j'(0)\to v$. With
  $\gamma(0)=p$, $\gamma'(0)=v$, continuity of $f$ gives
  $q=\exp_p(f(p,\gamma'(0))\gamma'(0))=\gamma(f(p,\gamma'(0)))\in C_m(p)$.
\end{proof}

\begin{corollary}[Cut point in every direction implies compactness]
  \label{cor:dc-ch13-2-11}
  \dcref{ch13:2.11}
  Suppose $M$ is complete and there exists $p\in M$ which has a cut point for every
  geodesic starting from $p$. Then $M$ is compact.
\end{corollary}
\begin{proof}
\sketch
  Since $M=\bigcup\{\gamma(t);t\leq f(p,\gamma'(0))\}$ over all geodesics
  $\gamma$ with $\gamma(0)=p$, and $f$ is continuous and (by hypothesis) bounded, $M$
  is bounded; the result follows from the Hopf–Rinow Theorem.
\end{proof}

\begin{proposition}[Geometry of a nearest cut point]
  \label{prop:dc-ch13-2-12}
  \uses{prop:dc-ch13-2-2}
  \dcref{ch13:2.12}
  Let $p\in M$ and suppose there exists $q\in C_m(p)$ realizing the distance from $p$
  to $C_m(p)$. Then:
  (a) either there exists a minimizing geodesic $\gamma$ from $p$ to $q$ along which
    $q$ is conjugate to $p$,
  (b) or there exist exactly two minimizing geodesics $\gamma$ and $\sigma$ from $p$
    to $q$, with $\gamma'(\ell)=-\sigma'(\ell)$, $\ell=d(p,q)$.
\end{proposition}
\begin{proof}
\sketch
  Let $\gamma$ be minimizing from $p$ to $q$. By \cref{prop:dc-ch13-2-2} either $q$
  is conjugate to $p$ along $\gamma$ — case (a) — or there is another minimizing
  $\sigma\neq\gamma$ with $\ell(\sigma)=\ell(\gamma)$. Suppose $q$ is not conjugate to
  $p$ along $\gamma$ or $\sigma$, and $\gamma'(\ell)\neq-\sigma'(\ell)$; we derive a
  contradiction (which also shows there can be only two such geodesics). Choose
  $V\in T_qM$ with $\langle V,\gamma'(\ell)\rangle<0$ and
  $\langle V,\sigma'(\ell)\rangle<0$. Since $q$ is not conjugate along $\gamma$, build
  a variation $\gamma_s$ of $\gamma$ with
  $\frac{d}{ds}\ell(\gamma_s)|_{s=0}=\langle V,\gamma'(\ell)\rangle<0$ by the formula
  for the first variation, and similarly $\sigma_s$ of $\sigma$ with
  $\frac{d}{ds}\ell(\sigma_s)|_{s=0}=\langle V,\sigma'(\ell)\rangle<0$ and
  $\sigma_s(\ell)=\tau(s)$. For small $s>0$, $\ell(\gamma_s)<\ell(\gamma)$ and
  $\ell(\sigma_s)<\ell(\sigma)$. If $\ell(\gamma_s)=\ell(\sigma_s)$, by \cref{prop:dc-ch13-2-2}
  $\tau(s)=\gamma_s(\ell)$ is a cut point with
  $d(p,\gamma_s(\ell))=\ell(\gamma_s)<d(p,C_m(p))$, a contradiction; if
  $\ell(\gamma_s)<\ell(\sigma_s)$, $\sigma_s$ is not minimizing so a cut point
  $\sigma_s(\tilde t)$, $\tilde t<\ell$, contradicts again; the case
  $\ell(\gamma_s)>\ell(\sigma_s)$ is similar.
\end{proof}

\begin{proposition}[Injectivity radius or a shortest closed geodesic]
  \label{prop:dc-ch13-2-13}
  \uses{prop:dc-ch13-2-9, prop:dc-ch10-2-4, prop:dc-ch13-2-12}
  \dcref{ch13:2.13}
  If the sectional curvature $K$ of a complete Riemannian manifold $M$ satisfies
  $0<K_{\min}\leq K\leq K_{\max}$, then:
  (a) $i(M)\geq\pi/\sqrt{K_{\max}}$, or
  (b) there exists a closed geodesic $\gamma$ in $M$, whose length is less than that
    of any other closed geodesic in $M$, with $i(M)=\tfrac{1}{2}\ell(\gamma)$.
\end{proposition}
\begin{proof}
\sketch
  By Bonnet--Myers, $M$ is compact. Since $T_1M$ is
  compact, \cref{prop:dc-ch13-2-9} gives $p\in M$ with $d(p,C_m(p))=\inf_{r\in M}d(r,C_m(r))$,
  and $C_m(p)$ compact gives $q\in C_m(p)$ realizing this distance. If $q$ is conjugate
  to $p$, then $d(p,q)\geq\pi/\sqrt{K_{\max}}$ by \cref{prop:dc-ch10-2-4} — case
  (a). If $q$ is not conjugate to $p$, by \cref{prop:dc-ch13-2-12} there are two minimizing
  geodesics $\mu,\sigma$ from $p$ to $q$ with $\mu'(\ell)=-\sigma'(\ell)$,
  $\ell=d(p,q)$. Since $q\in C_m(p)$ we have $p\in C_m(q)$ and $p$ realizes the
  distance from $q$ to $C_m(q)$, so $\mu'(0)=-\sigma'(0)$; thus $\mu$ and $\sigma$ form
  a closed geodesic $\gamma$ satisfying (b).
\end{proof}

\section{The estimate of the injectivity radius}

\begin{proposition}[Lower bound for the injectivity radius]
  \label{prop:dc-ch13-3-1}
  \dcref{ch13:3.1}
  Let $M^n$, $n\geq 3$, be a simply connected, compact Riemannian manifold with
  $1/4<K\leq 1$. Then $i(M)\geq\pi$.
\end{proposition}
\begin{proof}
  \sketch
  Suppose $i(M)<\pi$. By \cref{prop:dc-ch13-2-13} there is a
  closed geodesic $\gamma$ of length $\ell=\ell(\gamma)<2\pi$. By \cref{cor:dc-ch11-2-10} the conjugate points to $\gamma(0)$ along $\gamma$ are discrete; choose
  $\varepsilon>0$ with: (1) $\gamma(\ell-\varepsilon)$ not conjugate to $p=\gamma(0)$;
  (2) $\exp_p$ a diffeomorphism on $B_{2\varepsilon}(p)$;
  (3) $3\varepsilon<2\pi-\pi/\sqrt{\bar K}$, $\bar K=\inf_M K$;
  (4) $3\varepsilon<2\pi-\ell$; (5) $5\varepsilon<2\pi$. By Sard's Theorem pick a
  regular value $q\in B_\varepsilon(\gamma(\ell-\varepsilon))$ of $\exp_p$, with
  $q=\gamma_1(t)$ a geodesic from $p$, $3\varepsilon<\ell(\gamma_1)<\ell$. A
  minimizing geodesic $\gamma_o$ from $p$ to $q$ has
  $\ell(\gamma_o)\leq d(p,\gamma(\ell-\varepsilon))+d(\gamma(\ell-\varepsilon),q)\leq 2\varepsilon$,
  so $\gamma_o\neq\gamma_1$.

  Consider $\Omega_{p,q}^c$ (\cref{rem:dc-ch11-2-12}) and its finite-dimensional
  approximation $B$. Since $q$ is a regular value, all critical points in
  $\Omega_{p,q}^c$ are non-degenerate. As $M$ is simply connected, there is a homotopy
  $\gamma_s$ between $\gamma_o$ and $\gamma_1$ fixing $p,q$; deforming into $B$ and
  applying the index-zero/one variant of \cref{lem:dc-ch13-3-3}, for all
  $\delta>0$ the larger-length curve
  $\bar\gamma$ of the deformed $\bar\gamma_s$ satisfies $\ell(\bar\gamma)<a+\delta$,
  $a=\max\{\ell(\gamma_o),\ell(\gamma_1),\ell(\sigma)\}$, with $\sigma$ the largest
  geodesic of index $<2$. Put $\delta=\varepsilon$. Here $\ell(\gamma_o)\leq2\varepsilon$
  and $\ell(\gamma_1)<\ell<2\pi-3\varepsilon$ by (4). To estimate $\ell(\sigma)$, apply
  \cref{lem:dc-ch13-3-2} with $\bar M$ a sphere of curvature $\bar K=\inf_M K>1/4$, getting
  $\operatorname{index}\bar\sigma\leq\operatorname{index}\sigma<2$. Since $\bar\sigma$
  is a geodesic on $S^n$, the Morse Index Theorem gives: if
  $\ell(\sigma)>\pi/\sqrt{\bar K}$ then $\operatorname{index}\bar\sigma\geq2$;
  therefore by (3) $\ell(\sigma)\leq\pi/\sqrt{\bar K}<2\pi-3\varepsilon$. Hence by (5)
  $\ell(\bar\gamma)<a+\varepsilon<2\pi-3\varepsilon+\varepsilon=2\pi-2\varepsilon$.
  On the other hand, the closed-geodesic estimate gives in the
  homotopy a curve $\bar\gamma_{s_o}$ with $\ell(\gamma_o)+\ell(\bar\gamma_{s_o})\geq2\pi$,
  so $\ell(\bar\gamma)\geq2\pi-\ell(\gamma_o)\geq2\pi-2\varepsilon$, contradicting
  $\ell(\bar\gamma)<2\pi-2\varepsilon$. Thus $i(M)<\pi$ is untenable.
\end{proof}

\begin{lemma}[Index comparison]
  \label{lem:dc-ch13-3-2}
  \dcref{ch13:3.2}
  Let $M$ and $\bar M$ be Riemannian manifolds of the same dimension with sectional
  curvatures satisfying $\sup\bar K\leq\inf K$. Consider a geodesic
  $\sigma:[0,\ell]\to M$, $\sigma(0)=p$, fix $\bar p\in\bar M$, let
  $i:T_pM\to T_{\bar p}\bar M$ be a linear isometry, and set
  $\bar\sigma(t)=\exp_{\bar p}t(i\sigma'(0))$. Then
  $\operatorname{index}\sigma\geq\operatorname{index}\bar\sigma$.
\end{lemma}
\begin{proof}
\sketch
  For a piecewise differentiable field $\bar W$ along $\bar\sigma$, define
  $W$ along $\sigma$ by $W(t)=P_t\circ i^{-1}\circ\bar P_t^{-1}(\bar W(t))$ (parallel
  transports). Then $\langle W,\sigma'\rangle=\langle\bar W,\bar\sigma'\rangle$,
  $|W(t)|=|\bar W(t)|$, $\langle W',W'\rangle=\langle\bar W',\bar W'\rangle$. Since
  $\sup\bar K\leq\inf K$, $I(W,W)\leq I(\bar W,\bar W)$; so $I(\bar W,\bar W)<0$ forces
  $I(W,W)<0$.
\end{proof}

For the next lemma we need the following facts. Let $f:M^n\to\mathbb{R}$ be a
differentiable function on a differentiable manifold $M$. A point $p\in M$ is a
critical point of $f$ if $df(p)=0$; $f(p)$ is then called a critical value of $f$.
If $p$ is a critical point and $(x_1,\dots,x_n)$ is a coordinate system on $M$
around $p$, the hessian of $f$, defined by the matrix
$(\partial^2f/\partial x_i\partial x_j)(p)$, represents a symmetric bilinear form
on $T_pM$ that does not depend on the chosen system of coordinates. We say that
the critical point is non-degenerate if the determinant of this matrix is
different from zero. Such a critical point is isolated, and it is possible to
choose a coordinate system $(x_1,\dots,x_{n-\lambda},y_1,\dots,y_\lambda)$ in such
a way that, in the coordinate neighborhood considered,

$$
f(x_1,\dots,x_{n-\lambda},y_1,\dots,y_\lambda)
=f(p)+x_1^2+\cdots+x_{n-\lambda}^2-y_1^2-\cdots-y_\lambda^2.
$$

The integer $\lambda$ is the index of the non-degenerate critical point $p$.

\begin{lemma}[Morse-theoretic deformation]
  \label{lem:dc-ch13-3-3}
  \uses{prop:dc-ch13-3-1, prop:dc-ch13-2-13, prop:dc-ch13-2-12, lem:dc-ch13-3-2, cor:dc-ch13-2-10}
  \dcref{ch13:3.3}
  Let $M^n$ be a differentiable manifold and $f:M^n\to\mathbb{R}$ a function whose
  critical points are all non-degenerate. Given $p,q\in M$ and a differentiable curve
  $\gamma:[0,1]\to M$ with $\gamma(0)=p$, $\gamma(1)=q$, let $a=\max\{f(p),f(q)\}$,
  $M^a=\{x\in M;f(x)\leq a\}$, and $b=\max_{t\in[0,1]}(f\circ\gamma(t))$. Assume
  $f^{-1}([a,b])$ is compact and contains no critical points of index zero or one.
  Then for all $\delta>0$, $\gamma$ is homotopic, keeping endpoints fixed, to a curve
  $\bar\gamma$ with $\bar\gamma([0,1])\subset M^{a+\delta}$.
\end{lemma}
\begin{proof}
\sketch
  If $f^{-1}([a_1,a_2])$ is compact and
  critical-point-free, the gradient flow of $f$ gives a deformation retract
  $M^{a_2}\to M^{a_1}$. Let $p_1,\dots,p_k$ be the critical points in $M^b-M^a$ with
  distinct values $b\geq c_1>\dots>c_k>a$. At a critical value $c_1$ with index
  $\lambda\geq 2$, write $f=c_1+x_1^2+\dots+x_{n-\lambda}^2-y_1^2-\dots-y_\lambda^2$ on
  a neighborhood $H$, retract $M^{c_1+\varepsilon}\to M^{c_1-\varepsilon}\cup H$, and
  let $L=\{y_1=\dots=y_\lambda=0\}$. Since
  $\dim\bar\gamma_1'+\dim L\leq 1+n-2<\dim M$, transversality lets one perturb
  $\bar\gamma_1$ off $L$ within $H$ fixing endpoints; the gradient flow then retracts
  into $M^{c_1-\varepsilon}$. Proceeding inductively yields
  $\bar\gamma\subset M^{a+\delta}$.

  If $f^{-1}([a,b])$ contains critical points of index zero or one and $c$ is the
  largest value among them, the same argument gives
  $\bar\gamma([0,1])\subset M^{c+\delta}$.
\end{proof}

\begin{proposition}[Even-dimensional injectivity-radius estimate]
  \label{prop:dc-ch13-3-4}
  \uses{prop:dc-ch13-2-12, ex:dc-ch13-2-4}
  \dcref{ch13:3.4}
  If the sectional curvature $K$ of a compact orientable even-dimensional Riemannian
  manifold $M^n$ satisfies $0<K\leq 1$, then $i(M)\geq\pi$.
\end{proposition}
\begin{proof}
\sketch
  By compactness there are $p,q\in M$ with $q\in C_m(p)$ and $d(p,q)=i(M)$.
  Suppose $d(p,q)<\pi$. If $p$ is conjugate to $q$, then $d(p,q)\geq\pi$ by
  \cref{prop:dc-ch10-2-4}, a contradiction. So $p$ is not conjugate to $q$; by
  \cref{prop:dc-ch13-2-12,prop:dc-ch13-2-13} there is a closed geodesic $\gamma$ through $p=\gamma(0)$
  and $q$ with $\ell(\gamma)<2\pi$. Parallel transport along $\gamma$ leaves invariant
  a vector $V$ orthogonal to $\gamma$; computing the second
  variation $E_V''(0)$ for this $V$, it is strictly negative (here even dimension is
  crucial). Hence there is a variation through regular closed curves $\gamma_s$ with
  $\ell(\gamma_s)<\ell(\gamma)$ for $s\neq0$. Let $q_s$ be a point of $\gamma_s$ at
  maximum distance from $\gamma_s(0)$; the minimizing geodesic $\sigma_s$ from
  $q_s=\sigma_s(0)$ to $\gamma_s(0)$ has $\sigma_s(0)\to q$, and an accumulation point
  $w$ of $\sigma_s'(0)$ gives a minimizing geodesic $\sigma(t)=\exp_q tw$ from $q$ to
  $p$. By the first-variation formula $\sigma_s'(0)\perp\gamma_s'$ at $q_s$, so
  $\sigma'(0)\perp\gamma'$ at $q$. This yields three minimizing geodesics from $p$ to
  $q$; since $p$ is not conjugate to $q$, this contradicts \cref{prop:dc-ch13-2-12}.
\end{proof}

\begin{remark}[Orientability in even dimensions]
  \label{rem:dc-ch13-3-5}
  \uses{prop:dc-ch13-3-4}
  \dcref{ch13:3.5}
  By Synge's theorem the manifolds satisfying the hypotheses of
  \cref{prop:dc-ch13-3-4} are simply connected. Conversely a simply connected manifold is
  orientable (otherwise its orientation double cover would contradict simple
  connectivity). Hence the orientability hypothesis in
  \cref{prop:dc-ch13-3-4} is equivalent to assuming $M$ simply connected.
\end{remark}

\section{The Sphere Theorem}

In what follows, whether $\dim M$ is even or odd plays no role.

\begin{lemma}[Initial direction to a diameter endpoint]
  \label{lem:dc-ch13-4-1}
  \dcref{ch13:4.1}
  Let $M$ be a compact Riemannian manifold and $p,q\in M$ with
  $d(p,q)=\operatorname{diam}M$. Then for all $W\in T_pM$ there exists a minimizing
  geodesic $\gamma$ from $p=\gamma(0)$ to $q$ with $\langle\gamma'(0),W\rangle\geq0$.
\end{lemma}
\begin{proof}
\sketch
  Let $\lambda(t)=\exp_p tW$ and let $\gamma_t$ be a minimizing geodesic with
  $\gamma_t(0)=\lambda(t)$, $\gamma_t(\ell(\gamma_t))=q$. If for each $n$ there is
  $t_n\in[0,1/n]$ with $\langle\gamma_{t_n}'(0),\lambda'(t_n)\rangle\geq0$, a
  subsequence $\gamma_{t_n}\to\gamma$ gives
  $0\leq\langle\gamma'(0),\lambda'(0)\rangle=\langle\gamma'(0),W\rangle$. Otherwise
  some $n$ has $\langle\gamma_t'(0),\lambda'(t)\rangle<0$ for all $t\in[0,1/n]$. Taking
  a totally normal neighborhood $U$ of $\lambda(t)$, $q_o\in U$ on $\gamma_t$, and
  $\sigma_s$ minimizing from $q_o$ to $\lambda(s)$, the first variation of energy gives
  $\tfrac{1}{2}\frac{d}{ds}E(\sigma_s)|_{s=t}=-\langle\gamma_t'(0),\lambda'(t)\rangle>0$,
  so for $s<t$, $d(q_o,\lambda(s))<d(q_o,\lambda(t))$, hence
  $d(q,\lambda(s))<d(q,\lambda(t))$; but $p$ is at maximum distance from $q$, so
  $d(q,\lambda(t))\leq d(q,\lambda(0))$ — a contradiction.
\end{proof}

\begin{lemma}[Covering by two balls]
  \label{lem:dc-ch13-4-2}
  \uses{ex:dc-ch13-2-5}
  \dcref{ch13:4.2}
  Let $M^n$ be a compact, simply connected Riemannian manifold whose sectional
  curvature satisfies $\tfrac14<\delta\leq K\leq 1$, and let $p,q\in M$ with
  $d(p,q)=\operatorname{diam}(M)$. Then

  $$
  M=B_\rho(p)\cup B_\rho(q),
  $$

  where $B_r(p)$ denotes the open geodesic ball of radius $r$ and center $p$, and
  $\rho$ is such that $\pi/(2\sqrt\delta)<\rho<\pi$.
\end{lemma}
\begin{proof}
\sketch
  By \cref{prop:dc-ch13-3-1,prop:dc-ch13-3-4}, $B_\rho(p)$ contains no
  points of $C_m(p)$, so it is diffeomorphic via $\exp$ to a Euclidean ball, and
  likewise $B_\rho(q)$. (Here $i(M)\geq\pi$ enters crucially.) Suppose some $r\in M$
  has $d(p,r)\geq\rho$ and $d(q,r)\geq\rho$, with $d(p,r)\geq d(q,r)\geq\rho$. A
  minimizing geodesic from $q$ to $r$ meets $\partial B_\rho(q)$ at a point
  $q'\notin B_\rho(p)$ (else $d(r,q')>d(r,B_\rho(p))\geq d(r,B_\rho(q))=d(r,q')$,
  absurd). By Bonnet–Myers, $\operatorname{diam}(M)\leq\pi/\sqrt\delta<2\rho$, so the
  intersection $q''$ of the minimizing geodesic from $q$ to $p$ with $\partial B_\rho(q)$
  satisfies $d(p,q'')=d(p,q)-d(q,q'')<2\rho-\rho=\rho$, i.e. $q''\in B_\rho(p)$. As
  $\partial B_\rho(q)$ is path connected, $\partial B_\rho(p)\cap\partial B_\rho(q)\neq\emptyset$,
  giving $r_o$ with $d(r_o,p)=d(r_o,q)=\rho$. Take $\lambda$ minimizing from $p$ to
  $r_o$; by \cref{lem:dc-ch13-4-1} there is $\gamma$ minimizing from $p$ to $q$ with
  $\langle\gamma'(0),\lambda'(0)\rangle\geq0$. Let $s\in\gamma$ with $d(p,s)=\rho$.
  Applying Rauch comparison (\cref{prop:dc-ch10-2-5}), comparing $M^n$ with a
  sphere $S^n$ of curvature $\delta$: since the angle
  $\sphericalangle r_o p s\leq\pi/2$, $d(r_o,s)\leq\pi/(2\sqrt\delta)$. As
  $d(r_o,p)=d(r_o,q)=\rho$ and some $s$ has $d(r_o,s)<\rho$, the distance from $r_o$ to
  $\gamma$ is realized by an interior point $s_o$, with the minimizing geodesic from
  $r_o$ to $s_o$ orthogonal to $\gamma$ and
  $d(r_o,\gamma)=d(r_o,s_o)\leq\pi/(2\sqrt\delta)$. Since $d(p,q)\leq\pi/\sqrt\delta$,
  either $d(p,s_o)\leq\pi/(2\sqrt\delta)$ or $d(q,s_o)\leq\pi/(2\sqrt\delta)$; in the
  former case, since $d(r_o,s_o)\leq\pi/(2\sqrt\delta)$ and
  $\sphericalangle p s_o r_o=\pi/2$, Rauch gives $d(p,r_o)\leq\pi/(2\sqrt\delta)<\rho$,
  contradicting $d(p,r_o)=\rho$ (the other case is analogous).
\end{proof}

\begin{lemma}[Equator between two diameter endpoints]
  \label{lem:dc-ch13-4-3}
  \uses{lem:dc-ch13-4-2}
  \dcref{ch13:4.3}
  Under the conditions of \cref{lem:dc-ch13-4-2}, on each geodesic of length $\rho$ starting from
  $p$ there exists a unique point $m$ such that $d(p,m)=d(q,m)<\rho$. Similarly, on
  each geodesic of length $\rho$ starting from $q$ there exists a unique point $n$
  equidistant from $p$ and $q$.
\end{lemma}
\begin{proof}
\sketch
  For a geodesic $\gamma$ with $\gamma(0)=p$, set
  $f(s)=d(q,\gamma(s))-d(p,\gamma(s))$. Then $f$ is continuous and $f(0)=d(q,p)>0$. If
  $\gamma(s_o)$ is the cut point of $p$ along $\gamma$,
  \cref{prop:dc-ch13-3-1,prop:dc-ch13-3-4} gives
  $d(p,\gamma(s_o))\geq\pi>\rho$, and
  \cref{lem:dc-ch13-4-2} gives $d(q,\gamma(s_o))<\rho$, so $f(s_o)<0$. Hence $f(s_1)=0$ for some
  $s_1\in(0,s_o)$, giving $m=\gamma(s_1)$. For uniqueness, suppose $m_1\neq m_2$ both
  equidistant with $m_1$ between $p$ and $m_2$; then
  $d(q,m_2)=d(p,m_2)=d(p,m_1)+d(m_1,m_2)=d(q,m_1)+d(m_1,m_2)$, so the unique minimizing
  geodesic $\sigma_1$ from $q$ to $m_1$ coincides with $\gamma$, forcing $q\in\gamma$
  and then $p=q$, a contradiction.
\end{proof}

\begin{remark}[Continuity and sphere gluing]
  \label{rem:dc-ch13-4-4}
  \uses{thm:dc-ch13-1-1, lem:dc-ch13-4-2, lem:dc-ch13-4-3}
  \dcref{ch13:4.4}
  By its uniqueness, the point $m$ depends continuously on the (initial direction of
  the) geodesic which contains it.
\end{remark}

\begin{remark}[Remark 2.1]
  \label{rem:dc-ch13-2-1}
  \dcref{ch13:2.1}
  The cut locus, for surfaces, was introduced as the "ligne de partage" by
  H. Poincaré in 1905 [Po]. For Riemannian manifolds it was introduced by
  J. H. C. Whitehead in 1935 [Wh], who proved several of the properties presented
  here (pp. 700–704 of [Wh]). Interest revived in 1959 when Klingenberg [Kℓ 1]
  showed its usefulness in improving the "pinching" of the sphere theorem.
\end{remark}

% Formalization support lemmas.
%
% NOT book material. These declarations were invented in this workspace to
% carry a Lean formalization and have no counterpart in the source text --
% they carry no \dcref. They live here, outside the book chapters, so the
% book blueprint stays timeless mathematics while their \leanok marks and
% Lean anchors are preserved.
%
% Do not add book mathematics to this file, and do not add formalization
% notes to the book chapters.

\begin{corollary}[Positive constant-curvature endpoints lie in the conjugate locus]
  \label{cor:dc-ch5-3-3-conjugate-locus-endpoint}
  \lean{Riemannian.Jacobi.mem_conjugateLocus_constCurvature_pos}
  \leanok
  \uses{cor:dc-ch5-3-3-first-conjugate,def:dc-ch5-3-4}
  Under the hypotheses of the preceding corollary, if $p=\gamma(0)$ then
  $\gamma(\pi/\sqrt{K_0})\in C(p)$. Thus the intrinsic formalized consequence
  of the sphere calculation is that every such first endpoint belongs to the
  conjugate locus; identifying it with the antipode, and proving
  $C(p)=\{-p\}$ for the round sphere, requires the additional global sphere
  geometry.
\end{corollary}

\begin{corollary}[First conjugate point in positive constant curvature]
  \label{cor:dc-ch5-3-3-first-conjugate}
  \lean{Riemannian.Jacobi.isFirstConjugatePointAt_constCurvature_pos}
  \leanok
  \uses{ex:dc-ch5-3-3,def:dc-ch5-3-1,lem:dc-ch8-2-1-no-conjugate}
  Let $M^n$ have constant sectional curvature $K_0>0$, with $n\geq 2$, and let
  $\gamma\colon[0,\pi/\sqrt{K_0}]\to M$ be a unit-speed geodesic. Then
  $\gamma(\pi/\sqrt{K_0})$ is the first conjugate point to $\gamma(0)$ along
  $\gamma$. Indeed, the sine Jacobi fields of \cref{ex:dc-ch5-3-3} give
  conjugacy at the endpoint, while the constant-curvature non-conjugacy
  criterion rules out every $0<t<\pi/\sqrt{K_0}$.
\end{corollary}

\begin{lemma}[pointwise Riemannian volume form]
  \label{lem:dc-ch1-2-11-pointwise}
  \uses{def:dc-ch1-2-1}
  \lean{Riemannian.DCPointwiseOrientation, Riemannian.DCVolumeFormAt,
    Riemannian.DCVolumeFormAt_eq_det}
  \leanok
  Let $(M,g)$ be an $n$-dimensional Riemannian manifold and choose an orientation
  $o_p$ of each tangent space.  At every $p\in M$ there is a canonical alternating
  $n$-form $\nu_p$ determined by $g_p$ and $o_p$.  If $(e_1,\ldots,e_n)$ is a
  positively oriented orthonormal basis of $T_pM$, then
  $$
  \nu_p(v_1,\ldots,v_n)=\det\bigl(\langle v_i,e_j\rangle_p\bigr)_{i,j}.
  $$
\end{lemma}
\begin{proof}
  \leanok
  The oriented inner-product space $(T_pM,g_p,o_p)$ has its canonical top
  alternating form.  By its defining invariance, this form is the determinant
  form of every positively oriented orthonormal basis, which gives the displayed
  formula.
\end{proof}

\begin{lemma}[synchronized finite chart partition for two curves]
  \label{lem:dc-ch8-2-1-chart-partition}
  \lean{Riemannian.Jacobi.exists_pair_curve_chart_partition}
  \leanok
  If $\gamma:[0,1]\to M$ and $\tilde\gamma:[0,1]\to\tilde M$ are continuous, then there are
  a monotone partition $0=\tau_0\leq\tau_1\leq\cdots$ which is eventually equal to $1$, and
  charts $\alpha_i$ on $M$ and $\tilde\alpha_i$ on $\tilde M$, such that each closed piece
  $[\tau_i,\tau_{i+1}]$ is mapped by $\gamma$ into the source of $\alpha_i$ and by
  $\tilde\gamma$ into the source of $\tilde\alpha_i$.  Thus both curves have a common finite
  subdivision by one-chart pieces.
\end{lemma}
\begin{proof}
  \leanok
  Apply the compactness lemma to the open cover of $[0,1]$ by the intersections of the pullbacks
  of chart sources along $\gamma$ and $\tilde\gamma$.
\end{proof}

\begin{lemma}[a common chart window for the Cartan transfer]
  \label{lem:dc-ch8-2-1-window-transfer}
  \lean{Riemannian.Jacobi.exists_common_chart_window_metricInner_mfderiv_eq_of_semiconjugacy_of_curvatureFormAt_of_isLocalDiffeomorphAt}
  \leanok
  \uses{lem:dc-ch8-2-1-window,
    lem:dc-ch8-2-1-exp-norm-transfer-general-discharged}
  In the setting of \cref{lem:dc-ch8-2-1-exp-norm-transfer-general-discharged}, suppose the two
  radial geodesics lie in fixed chart sources on $[0,1]$. Then there is one common interval
  $[a',b']$, with $a'<0<1<b'$, on which both chart-source conditions hold, and Cartan's
  curvature match for $\phi_t$ on that interval implies that the semiconjugacy preserves the
  metric at $q=\exp_p(v)$.
\end{lemma}
\begin{proof}
  \leanok
  \uses{lem:dc-ch8-2-1-window,
    lem:dc-ch8-2-1-exp-norm-transfer-general-discharged}
  Apply \cref{lem:dc-ch8-2-1-window} to each radial geodesic. If the resulting intervals are
  $[a_1,b_1]$ and $[a_2,b_2]$, take
  $$
  a'=\max\{a_1,a_2\},\qquad b'=\min\{b_1,b_2\}.
  $$
  Both strict endpoint inequalities persist, and the common interval lies in both original
  intervals. The two chart-source conditions therefore hold there, so the curvature match and
  \cref{lem:dc-ch8-2-1-exp-norm-transfer-general-discharged} give the metric identity. $\square$
\end{proof}

\begin{lemma}[cyclic polarization of the differentiated conformality equations]
  \label{lem:dc-ch8-5-differential-core}
  \uses{def:dc-ch8-5-conformal}
  \lean{Riemannian.liouville_secondDerivative_formula,
    Riemannian.liouville_secondDerivative_formula_of_inner_eq_zero}
  \leanok
  Let $E$ be a finite-dimensional real inner-product space. Let $L:E\to E$ be a
  conformal linear map with nonzero coefficient $\lambda$, let $B:E\times E\to E$
  be symmetric bilinear, and let $d\lambda:E\to\mathbb R$ have gradient
  $\nabla\lambda$. Suppose the differentiated conformality identity is
  \[
    \langle B(u,v),Lw\rangle+\langle Lv,B(u,w)\rangle
      =2\lambda\,d\lambda(u)\,\langle v,w\rangle
  \]
  for all $u,v,w$. Then
  \[
    B(u,v)=\frac1\lambda\bigl(d\lambda(u)Lv+d\lambda(v)Lu
      -\langle u,v\rangle L(\nabla\lambda)\bigr).
  \]
  In particular, if $u\perp v$, then
  \[
    B(u,v)=\frac1\lambda\bigl(d\lambda(u)Lv+d\lambda(v)Lu\bigr),
  \]
  which is equation (7) in the proof of Liouville's theorem.
\end{lemma}
\begin{proof}
  \leanok
  Write the differentiated identity for $(u,v,w)$, $(v,w,u)$, and $(w,u,v)$.
  Symmetry of $B$ and symmetry of the inner product identify the three repeated
  terms. Adding the first two equations and subtracting the third gives
  \[
    \langle B(u,v),Lw\rangle
      =\lambda\bigl(d\lambda(u)\langle v,w\rangle
        +d\lambda(v)\langle u,w\rangle
        -d\lambda(w)\langle u,v\rangle\bigr).
  \]
  The conformality identity for $L$ and
  $d\lambda(w)=\langle\nabla\lambda,w\rangle$ show that the right-hand side
  is the inner product of $Lw$ with the displayed formula for $B(u,v)$.
  Since $\lambda\ne0$, $L$ is injective and hence surjective in finite
  dimension; equality against every $Lw$ therefore gives equality of vectors.
  The orthogonal case is obtained by setting $\langle u,v\rangle=0$.
\end{proof}

\begin{lemma}[orientation determines the determinant of an orthogonal map]
  \label{lem:dc-ch9-3-7-orientation-det}
  \lean{Riemannian.Variation.det_sq_eq_one_of_isometry,
    Riemannian.Variation.det_eq_one_of_orientation_preserving,
    Riemannian.Variation.det_eq_neg_one_of_orientation_reversing}
  \leanok
  Let $V$ be a finite-dimensional oriented real inner-product space and let $A:V\to V$
  be an orthogonal transformation. Then $(\det A)^2=1$. If $A$ preserves the orientation,
  then $\det A=1$, while if $A$ reverses the orientation, then $\det A=-1$.
\end{lemma}
\begin{proof}
  \leanok
  In an orthonormal basis, the matrix of $A$ satisfies $A^{T}A=I$. Taking determinants gives
  $(\det A)^2=1$. Preservation of orientation is equivalent to $\det A>0$, and reversal of
  orientation is equivalent to $\det A<0$, which selects the corresponding root.
\end{proof}

